\documentclass[UTF8,zihao=-4]{ctexart}
\usepackage[a4paper,margin=2.35cm]{geometry}
\usepackage{fontspec}
\usepackage{xeCJK}
\usepackage{graphicx}
\usepackage{caption}
\usepackage{float}
\usepackage{longtable}
\usepackage{booktabs}
\usepackage{array}
\usepackage{amsmath,amssymb}
\usepackage{hyperref}
\usepackage{xcolor}
\usepackage{etoolbox}
\usepackage{fancyhdr}

\setmainfont{Times New Roman}
\IfFontExistsTF{SimSun}{\setCJKmainfont{SimSun}}{\setCJKmainfont{Noto Serif CJK SC}}
\graphicspath{{images/}{E:/大学/书/与自然对话/vlm/images/}}
\linespread{1.12}
\setlength{\parindent}{2em}
\setlength{\parskip}{0.45em}
\hypersetup{colorlinks=true,linkcolor=black,urlcolor=blue}
\pagestyle{fancy}
\fancyhf{}
\fancyfoot[C]{\scriptsize 译者：智乃的脚小小的}
\renewcommand{\headrulewidth}{0pt}
\renewcommand{\footrulewidth}{0pt}
\fancypagestyle{plain}{%
  \fancyhf{}%
  \fancyfoot[C]{\scriptsize 译者：智乃的脚小小的}%
  \renewcommand{\headrulewidth}{0pt}%
  \renewcommand{\footrulewidth}{0pt}%
}

\newcommand{\BodyFont}{\fontsize{10.5pt}{15.5pt}\selectfont}
\newcommand{\origlabel}{\textbf{原文}\quad}
\newcommand{\translabel}{\textbf{译文}\quad}
\newenvironment{enblock}{\par\noindent\begingroup\BodyFont\origlabel}{\par\endgroup}
\newenvironment{zhblock}{\par\noindent\begingroup\BodyFont\translabel}{\par\endgroup\vspace{0.35em}}
\renewenvironment{quote}{\par\noindent\begingroup\BodyFont}{\par\endgroup}
\newcommand{\en}[1]{\begin{enblock}#1\end{enblock}}
\newcommand{\zh}[1]{\begin{zhblock}#1\end{zhblock}}
\newcommand{\pair}[2]{\begin{enblock}#1\end{enblock}\begin{zhblock}#2\end{zhblock}}
\newcommand{\speaker}[2]{\par\medskip\noindent\textbf{#1 / #2}\par}
\newcommand{\safegraphics}[3]{%
  \IfFileExists{images/#1}{\includegraphics[width=#2,height=#3,keepaspectratio]{images/#1}}{%
  \IfFileExists{E:/大学/书/与自然对话/vlm/images/#1}{\includegraphics[width=#2,height=#3,keepaspectratio]{E:/大学/书/与自然对话/vlm/images/#1}}{%
    \fbox{\parbox{0.72\linewidth}{\centering 图片引用保留：\texttt{\detokenize{#1}}}}%
  }}%
}
\newcommand{\safeincludeimage}[1]{\begin{center}\safegraphics{#1}{0.86\linewidth}{0.46\textheight}\end{center}}
\newcommand{\safeincludeimagepair}[2]{\begin{center}\safegraphics{#1}{0.45\linewidth}{0.34\textheight}\hfill\safegraphics{#2}{0.45\linewidth}{0.34\textheight}\end{center}}
\newcommand{\preservedimage}[1]{\safeincludeimage{#1}}
\newcommand{\preserveimage}[1]{\safeincludeimage{#1}}
\newcommand{\editorialnote}[1]{\par\smallskip\noindent{\footnotesize\color{gray}\textbf{编者按：}#1}\par\smallskip}
\newcommand{\orig}{\par\noindent\begingroup\BodyFont\origlabel}
\newcommand{\zhline}[1]{\par\endgroup\noindent\begingroup\BodyFont\translabel #1\par\endgroup\vspace{0.35em}}

\begin{document}
\sloppy
\BodyFont
% ===== 00_front_intro.tex =====
\begin{figure}[p]
\centering
\safeincludeimage{a6efff7140e16fd202e094f504efbaea737deea9ccb83d4abcff42c748ab01a2.jpg}
\end{figure}

\begin{center}
{\large 香港中文大學}\\
The Chinese University of Hong Kong

\vspace{1em}

{\Huge 與自然對話}\\
{\Huge IN DIALOGUE WITH NATURE}

\vspace{1em}

通識教育基礎課程讀本\\
Textbook for General Education Foundation Programme

\vspace{1em}

Second edition, 2012
\end{center}

\begin{figure}[p]
\centering
\safeincludeimage{8784b1a80b63978c6fc27b4602b394f6fdeabc934a2093867ec476486f1f88da.jpg}
\end{figure}

\section*{香港中文大學 大學通識教育部}
\begin{center}
Office of University General Education\\
The Chinese University of Hong Kong

\vspace{1em}

{\LARGE In Dialogue with Nature}\\
{\LARGE 與自然對話}
\end{center}

\noindent \textit{In Dialogue with Nature}（《與自然對話》）是香港中文大學通識教育基礎課程兩門課程之一「UGFN 1000 In Dialogue with Nature」的讀本。

\vspace{1em}

\noindent © Office of University General Education, The Chinese University of Hong Kong, 2011, 2012\\
Second edition 2012

\vspace{1em}

\noindent 版權所有。未經香港中文大學大學通識教育部事先書面許可，本出版物任何部分不得以任何形式複製，亦不得提供予他人使用。

\noindent 本書所選篇章均經版權持有人許可重印，或取自公有領域。除公有領域材料外，任何選文均不得在未取得「鳴謝」所列版權持有人事先許可的情況下複製。

\section*{Editorial Board}
Chan Chi Wang\\
Szeto Wai Man\\
Wong Wing Hung

\vspace{1em}

\noindent Managing editor / Ng Hiu Chun Emily\\
Editorial assistants / Lam Yee Ki Dora, Ng Tsz Ying Tracy\\
Cover and layout design / Ng Hiu Chun Emily

\vspace{1em}

\noindent ISBN: 978-988-16660-1-7

\vspace{1em}

\noindent Published by Office of University General Education\\
The Chinese University of Hong Kong\\
Shatin, Hong Kong

\vspace{1em}

\noindent http://www.cuhk.edu.hk/oge/gef\\
gef@cuhk.edu.hk

\vspace{1em}

\noindent Printed in Hong Kong by Wellfit Production \& Printing Company（唯美印刷製作公司）

\vspace{1em}

\noindent 本書以不經漂白和增白的原漿紙印刷。\\
This book is printed on unbleached, with no brightening agents, original pulp paper.

\section*{鳴謝}
\noindent 本書所選篇章均經版權持有人許可重印，或取自公有領域。除公有領域材料外，任何選文均不得在未取得下列版權持有人事先許可的情況下複製。

\vspace{0.5em}

\noindent pp. 5--10: From Plato, \textit{Republic}, translated by C.D.C. Reeve. Copyright © 2004 by Hackett Publishing Company, Inc. Reprinted by permission of Hackett Publishing Company, Inc. All rights reserved.

\noindent pp. 11--48: From David C. Lindberg, \textit{The Beginnings of Western Science: The European Scientific Tradition in Philosophical, Religious, and Institutional Context, Prehistory to A.D. 1450} (Second Edition). Published by The University of Chicago Press. Copyright © 1992, 2007 by The University of Chicago. All rights reserved.

\noindent pp. 49--62: From I. Bernard Cohen, \textit{The Birth of a New Physics: Revised and Updated}. Copyright © 1985 by I. Bernard Cohen. Copyright © 1960 by Educational Services Incorporated. Used by permission of W. W. Norton \& Company, Inc.

\noindent pp. 63--70: From Isaac Newton, \textit{The Principia: Mathematical Principles of Natural Philosophy}. Copyright © 2012 by University of California Press. Reproduced with permission of University of California Press via Copyright Clearance Center.

\noindent pp. 73--96: From Charles Darwin, \textit{On the Origin of Species by Means of Natural Selection} (First Edition) (1859). Reproduced from the public domain.

\noindent pp. 97--142: From James D. Watson and with Andrew Berry, \textit{DNA: The Secret of Life}. Copyright © 2003 by DNA Show LLC. Used by permission of Alfred A. Knopf, a division of Random House, Inc.

\noindent pp. 143--158: From Rachel Carson, \textit{Silent Spring: 40th Anniversary Edition}. Copyright © 1962 by Rachel L. Carson. Copyright renewed 1990 by Roger Christie. This selection originally appeared in \textit{SILENT SPRING: 40th Anniversary Edition} by Rachel Carson. This edition of ``pages 65 to pages 83 from \textit{SILENT SPRING: 40th Anniversary Edition} by Rachel Carson'' is published by Office of University General Education, The Chinese University of Hong Kong.

\noindent Published by arrangement with Frances Collin, Trustee u-w-o Rachel Carson through Bardon-Chinese Media Agency.

\noindent ALL RIGHTS RESERVED.

\noindent pp. 161--178: From Henri Poincaré, \textit{The Value of Science: Essential Writings of Henri Poincaré} (2001). Published by Modern Library, reproduced from the public domain.

\noindent pp. 179--194: From Eric Kandel, \textit{In Search of Memory: The Emergence of a New Science of Mind}. Copyright © 2006 by Eric R. Kandel. Used by permission of W. W. Norton \& Company, Inc.

\noindent pp. 195--218: From Joseph Needham and Colin A. Ronan, \textit{The Shorter Science and Civilisation in China: An Abridgement of Joseph Needham's Original Text}, Vol. 1 (1978). © Cambridge University Press, reproduced with permission.

\noindent pp. 219--244: Nathan Sivin, ``Why the Scientific Revolution Did Not Take Place in China -- or Didn't It?'' originally published in \textit{Chinese Science}, 1982, 5: 45--66, revised 2005.8.24, used by permission of the author.

\noindent pp. 245--252: 文言文篇錄自沈括，《新校正夢溪筆談》（1975），由中華書局（香港）有限公司授權翻印。詮譯及註釋由香港中文大學大學通識教育部提供。

\noindent pp. 253--258 (except for pp. 255--256 ``The Rainbow''): From Joseph Needham, \textit{Science and Civilisation in China}, Vol. 3 (1959), Vol. 4 (1962), Vol. 5 (1985). © Cambridge University Press, reproduced with permission.

\noindent pp. 255--256 (``The Rainbow''): Nathan Sivin, ``The Rainbow,'' unpublished, used by permission of the author.

\noindent pp. 259--274: From William Dunham, \textit{The Mathematical Universe: An Alphabetical Journey Through the Great Proofs, Problems, and Personalities}. Copyright © 1994 by John Wiley \& Sons, Inc. Reproduced with permission of John Wiley \& Sons, Inc.

\noindent pp. 275--290: From Euclid, \textit{Elements}. The text used in this anthology is adopted from a digitized version of Euclid, \textit{The Thirteen Books of Euclid's Elements}, translated by Thomas L. Heath (1956), made available by the Perseus Digital Library Project. Ed. Gregory R. Crane. Tufts University. (Date of access: Feb 2012) <http://www.perseus.tufts.edu>. The illustrations are reproduced by this publication.

\section*{目錄}
\noindent Letter to Students \dotfill vii\\
Notes on editing \dotfill xi\\
Introduction to ``In Dialogue with Nature'' \dotfill 1

\section*{Part I: Human Exploration of the Physical Universe}
\noindent Introduction \dotfill 3\\
Text 1a from \textit{Republic} / Plato \dotfill 5\\
Text 1b from \textit{The Beginnings of Western Science} / David C. Lindberg \dotfill 11\\
Text 2 from \textit{The Beginnings of Western Science} / David C. Lindberg \dotfill 17\\
Text 3a from \textit{The Birth of a New Physics} / I. Bernard Cohen \dotfill 49\\
Text 3b from \textit{The Principia: Mathematical Principles of Natural Philosophy} / Isaac Newton \dotfill 63

\section*{Part II: Human Exploration of the World of Life}
\noindent Introduction \dotfill 71\\
Text 4 from \textit{On the Origin of Species} / Charles Darwin \dotfill 73\\
Text 5 from \textit{DNA: The Secret of Life} / James D. Watson \dotfill 97\\
Text 6 from \textit{Silent Spring} / Rachel Carson \dotfill 143

\section*{Part III: Our Understanding of Human Understanding}
\noindent Introduction \dotfill 159\\
Text 7 from \textit{Science and Method} / Henri Poincaré \dotfill 161\\
Text 8 from \textit{In Search of Memory: The Emergence of a New Science of Mind} / Eric R. Kandel \dotfill 179\\
Text 9 from \textit{The Shorter Science and Civilisation in China} / Joseph Needham \dotfill 195\\
Text 10a from ``Why the Scientific Revolution Did Not Take Place in China -- or Didn't It?'' / Nathan Sivin \dotfill 219\\
Text 10b from 《夢溪筆談》 / 沈括 \dotfill 245\\
Brush Talks from Dream Brook / Shen Kua \dotfill 253\\
Text 11a from \textit{The Mathematical Universe} / William Dunham \dotfill 259\\
Text 11b from \textit{Elements} / Euclid \dotfill 275



\section*{致同學}
\noindent 親愛的同學：

\noindent 歡迎加入通識教育基礎課程（General Education Foundation Programme, GEF）的學習共同體。我們即將展開一段充滿挑戰的旅程：穿行於不同文化與學科的流水之間，進而航向遼闊的知識海洋。GEF 開設的兩門課程，即「與人文對話」（\textit{In Dialogue with Humanity}）與「與自然對話」（\textit{In Dialogue with Nature}），旨在拓闊視野，展示彼此不同的世界。我們將反思至關重要的議題與關懷，並在此過程中努力更完整地理解自身，以及我們所身處的宇宙。

\noindent 多數人進入大學，是為了在主修學科中獲得專門知識或一套技能，以此通向專業或學術生涯。然而，在當前這個技術快速變化、通訊高度整合、經濟波動不定的社會環境，也就是我們所稱的知識經濟之中，傳統的就業資歷已不再足夠。理解全球化世界的複雜性，需要寬廣而多元的視角。為了解決日益複雜的問題，往往必須同時而協調地調動多個知識領域。

\noindent 從根本上影響我們在世生活的，不只是我們擁有什麼知識或技能，也包括我們成為怎樣的人。因此，通識教育的終極目標，是培養負責任的知識分子與世界公民所需的素質。我們必須意識到認知世界的不同方式，回應人類存在的共同關切，對文化多樣性保持敏感；而最重要的是，我們必須培養面對陌生事物、並就其作出明智判斷的智性能力與心志。

\noindent 「與人文對話」與「與自然對話」構成香港中文大學通識教育課程的核心部分。這兩門課是以主題為本、由經典文本閱讀引導的研討課，也可從幾個方面被視為「核心」。

\noindent 就內容而言，GEF 課程涵蓋人類知識與實踐的兩個關鍵或核心領域：對作為人的意義與價值的追尋，以及人類理解自然的成就與限制。我們將閱讀的經典所包含的思想，觸及人類存在與知識中的重大問題，並界定了與當代生活仍然相關、歷久不衰的信念與價值。

\noindent 此外，兩門課也促進重要能力的發展，包括主動投入閱讀、討論、反思與寫作。這些能力既是終身學習者不可或缺的素質，也是大學生活以至往後人生得以充實展開所需的智性工具。

\noindent 更重要的是，GEF 課程提供了一個空間，讓核心態度得以開始形成。首先，是面對智性挑戰時的無畏精神。我們將直接與來自迥異傳統、文化與學科的原典對話，並努力理解它們。

\noindent 同樣重要的，是開放心靈的精神。當我們與文本、與彼此展開對話時，會發現各自的觀點常常互有差異。我們的信念與價值也可能受到挑戰。我們或會逐漸明白，對許多重要問題而言，並不存在確定或最終的答案；人類的知識與實踐是一個永無止境的尋求與再尋求的過程。也正是在這一過程中，我們將為負責任而積極地參與知識與福祉的追尋奠定基礎。

\noindent GEF 的參與者都是主動學習者。這一過程不應，也不會，被一個字母等級所概括。GEF 的學習成果包括：對文化差異更敏銳的覺察，以及對價值多樣性更深的敏感；對普遍關懷更強的意識，以及對人類知識的範圍與限制更深的敏感；面對艱深文本與複雜思想時的信心；以及以口頭或書面方式表達有根據論點的能力。要使這趟智性旅程成功，你們的主動參與正是關鍵。

\noindent 準備好吧！

\vspace{1em}

\noindent Leung Mei Yee\\
Director\\
General Education Foundation Programme

\section*{译者注}

\noindent 感谢强大的AI技术，让我可以毫无压力的完成对这本书的翻译。之前也有这本书的中文译本存在过，但那本书的翻译又臭又长还不准确。这本书略微的优化了翻译，但是也没好多少（笑）。本书完全由ChatGPT 5.5 进行翻译，译者本人没有进行过任何校对，但由于保留了英文原文，读者应该能很方便的进行对照和理解。希望这本书能对你的学习有所帮助。

\noindent 译者本人不对此书中的任何内容负责。

\noindent 另外，此译本完全免费发行，如果花钱你就被骗了。

\hfill 译者：智乃的脚小小的




\section*{編輯說明}
\begin{enumerate}
\item 本選集作為「與自然對話」課程的讀本出版；該課程是一個以主題為取向、以經典文本為本的核心課程的一部分。因此，對原典篇章的選取與刪節，均以使課程主題更為清晰集中為目標。
\item 本選集中重印的文本，除校正排印錯誤、為編輯一致性而調整排版樣式，以及刪去或調整原文本中對頁碼或附錄的引用外，均按原樣重印。
\item 書邊所有段落編號均由本書加上，以便課堂討論時引用。編號依照段落在原文本中出現的順序排列。因此，凡有段落被刪去，編號亦會跳過。
\item 所有腳註及方括號內文字均出自原文本；標明為「Ed.」者除外。
\item 原文本中的省略以方括號省略號「[...]」標示。凡如此標示的省略，均由本書作出。
\item 第 10b 篇「Brush Talks from Dream Brook 夢溪筆談」所有標題均由本書加上，並非原文本所有。中文文本中亦列出條目號碼，以標明相關條目在原文本中的出現次序。
\item 第 11b 篇「Elements」有以下特別說明：
\end{enumerate}

\begin{itemize}
\item 本選集只收入 Heath 的 \textit{Elements} 英譯，並刪去他對希臘文本的 ``scholia'' 評注以及譯註。
\item 在原出版物中，Heath 使用不同類型的括號，標示他為使譯文清楚而補入、但希臘原文中並非逐字具備的文字。為方便一般讀者閱讀，本書刪去這些括號而保留其中文字。唯一例外，是 Heath 用來指引讀者參照 \textit{Elements} 相關部分的方括號，例如「[Def. 10]」指定義 10，「[Post. 1]」指公設 1，「[C.N. 3]」指公理 3，「[I.4]」指第一卷命題 4；這些均予保留。
\end{itemize}

\section*{In Dialogue with Nature}
\section*{導論}
\noindent 人類長久以來都對自然懷有好奇。早在文明曙光出現以前，我們便一直努力理解世界所呈現的繁多現象。宇宙的某些運行似乎有序，例如日出日落、四季更替、天體旋轉；另一些則似乎混沌而隨機，例如風暴與火山爆發。或許，在這些模式或秩序背後，存在某些自然法則。為了理解這些觀察，人們在世界各地不同的具體境況與文化中，發展並確立了極為不同的理解自然的方法，並以之作為古代文化的基礎。

\noindent 早期自然哲學家及其後繼者，即今日的現代科學家，致力發現自然法則；這一努力一直是一場自然與我們之間的對話。這場對話透過經驗與推理進行。它始於約公元前 500 年的古希臘；當時哲學家開始把邏輯與理性的智性工具運用於對自然現象的觀察，以嘗試解釋存在。他們對自然的態度與探究方法，奠定了我們今日現代西方經驗科學及其所支撐技術的基礎。其後，更大量、更系統的觀察逐步修正了我們對自然的理解。現代科學家則從這些觀察中概括出自然定律。

\noindent 然而，在這一過程中，我們對自然的理解經歷了許多調整，甚至出現了一些極為深刻的根本變化。某些新的思維模式偶爾甚至挑戰了「理解」這一概念本身，迫使人類反思知識的意義。

\noindent 與此同時，在大致相同的時期，中國對理解自然的追尋走向了十分不同的答案。中國學者發展出陰陽、五行等概念，形成了一套與希臘同行迥然有別的理解。直到近代早期，西方科學的概念才開始滲入中國的文化精神；然而，中國在許多世紀中即使沒有被認為由科學帶來的那些益處，仍是古代世界公認的強國。顯然，東西文化之間持續不斷的相遇，對我們理解自然仍大有可為。

\noindent 本課程邀請同學追溯前人在這場求知歷程中的思路，並與他們的著作展開對話。沿著偉大心靈的足跡前行，同學將能對自然以及人類與自然的互動形成有根據的見解。

\section*{課程結構}
\noindent 本課程將帶領同學踏上一段智性旅程。我們首先走向外在世界，考察的對象是物質宇宙與生命世界。這兩項考察構成本課程的前兩部分：第一部分「人類對物質宇宙的探索」（Human Exploration of the Physical Universe）與第二部分「人類對生命世界的探索」（Human Exploration of the World of Life）。我們將認識西方科學若干重大發現：宇宙受物理定律支配；物種正在經歷自然選擇；DNA 是生命的密碼。

\noindent 接著，我們將回到自身，考察我們的心靈與理解。這是本課程第三部分「我們對人類理解的理解」（Our Understanding of Human Understanding），也是一段反思之旅。

\noindent 首先，我們將反思自己對自然的理解，以及科學知識的有效性與限制。我們將嘗試理解具備理性思考能力的人類心靈的本質。其次，透過比較西方科學與中國哲學，我們將嘗試理解為什麼現代科學沒有在中國出現；並且在中國哲學中尋找一些視角，幫助我們反思現代科學處理人與自然關係的方式。

\vspace{1em}

\noindent Wong Wing Hung\\
Deputy Director\\
General Education Foundation Programme

\section*{第一部分：人類對物質宇宙的探索}
\noindent 這是旅程的第一部分。我們將從柏拉圖的《理想國》（\textit{Republic}）開始。在通常被稱為「洞穴寓言」（Allegory of the Cave）的選段（Text 1a）中，柏拉圖透過蘇格拉底與格勞孔的對話，區分了感知與實在。這一區分奠定了西方科學的基本基調。

\noindent 依柏拉圖之見，我們在「可見領域」（visible realm）中所見的只是影子。只有在「可知領域」（intelligible realm）中，人才能找到真正實在的「善的理念」（Form of the Good），它是「一切正確而美好事物的原因」。這是真正追求理解者的終極關切。對現代科學家而言，也存在一個由自然法則構成的「可知領域」。科學家的終極關切，不是對自然現象作全面描述，而是理解其背後的法則。David Lindberg 是研究中世紀與近代早期科學的著名史家。在他廣受使用的教科書 \textit{The Beginnings of Western Science} 的選段（Text 1b）中，他闡明了柏拉圖兩個領域的觀念及其與科學的關係。

\noindent 第二篇文本（Text 2）同樣選自 \textit{The Beginnings of Western Science}。它介紹亞里士多德物理學，並以簡明的歷史敘述，說明中世紀科學家後來如何發展運動概念，從而為牛頓物理學的出現奠定基礎。

\noindent Text 3a 選自 \textit{The Birth of a New Physics}。這本書由國際知名的牛頓研究者 I. Bernard Cohen 面向一般讀者撰寫。在這一選段中，Cohen 闡述了牛頓物理學的背景，以及偉大著作 \textit{The Principia: Mathematical Principles of Natural Philosophy}（《自然哲學的數學原理》）背後的主要思想；該書長久以來被視為人類歷史上最重要的著作之一。Text 3b 是這部巨著開首數頁；牛頓在其中建構了一個關於世界的數學體系。讀者將能體會牛頓如何努力為質量、力、運動等抽象概念作出精確定義，也能看到他試圖在數學與哲學之間劃定界線。


% ===== 01_plato_republic.tex =====
\section*{Socrates' Narration Continues / 苏格拉底继续叙述}

\speaker{SOCRATES}{苏格拉底}
\en{Next, then, compare the effect of education and that of the lack of it on our nature to an experience like this. Imagine human beings living in an underground, cavelike dwelling, with an entrance a long way up that is open to the light and as wide as the cave itself. They have been there since childhood, with their necks and legs fettered, so that they are fixed in the same place, able to see only in front of them, because their fetter prevents them from turning their heads around. Light is provided by a fire burning far above and behind them. Between the prisoners and the fire, there is an elevated road stretching. Imagine that along this road a low wall has been built--like the screen in front of people that is provided by puppeteers, and above which they show their puppets.}
\zh{那么接下来，请把教育以及缺乏教育对我们本性的影响，同下面这种经历相比较。设想有一些人住在地下洞穴般的居所里，洞口在远处上方，向光敞开，宽度与洞穴本身相当。他们自幼就在那里，颈项和腿都被镣铐束缚，因此固定在原处，只能看见面前的东西，因为束缚使他们不能转头。光来自他们身后高处燃烧的一团火。在囚徒和火之间，有一条高起的道路延伸过去。再设想沿着这条路筑起了一道矮墙--就像傀儡戏艺人置于观众面前、并在其上方展示傀儡的屏幕。}

\speaker{GLAUCON}{格劳孔}
\en{I am imagining it.}
\zh{我想象到了。}

\speaker{SOCRATES}{苏格拉底}
\en{Also imagine, then, that there are people alongside the wall carrying multifarious artifacts that project above it--statues of people and other animals, made of stone, wood, and every material. And as you would expect, some of the carriers are talking and some are silent.}
\zh{那么也请设想，墙边还有一些人，携带着各种伸出墙头的器物--用石头、木头以及各类材料制成的人像和其他动物像。正如你会预料的那样，搬运者有些在说话，有些沉默不语。}

\speaker{GLAUCON}{格劳孔}
\en{It is a strange image you are describing, and strange prisoners.}
\zh{你描绘的是一个奇特的图像，囚徒也奇特。}

\speaker{SOCRATES}{苏格拉底}
\en{They are like us. I mean, in the first place, do you think these prisoners have ever seen anything of themselves and one another besides the shadows that the fire casts on the wall of the cave in front of them?}
\zh{他们就像我们。我的意思是，首先，你认为这些囚徒除了火光投在他们面前洞壁上的自己和彼此的影子之外，还曾看见过自己或同伴的任何东西吗？}

\speaker{GLAUCON}{格劳孔}
\en{How could they, if they have to keep their heads motionless throughout life?}
\zh{如果他们一生都不得不让头保持不动，又怎么可能看见呢？}

\speaker{SOCRATES}{苏格拉底}
\en{What about the things carried along the wall? Isn't the same true where they are concerned?}
\zh{那么沿墙被搬运的那些东西呢？关于它们，情形不也是一样吗？}

\speaker{GLAUCON}{格劳孔}
\en{Of course.}
\zh{当然。}

\speaker{SOCRATES}{苏格拉底}
\en{And if they could engage in discussion with one another, don't you think they would assume that the words they used applied to the things they see passing in front of them?}
\zh{如果他们能够彼此交谈，你不认为他们会以为自己所用的词语，就是指他们眼前经过的那些东西吗？}

\speaker{GLAUCON}{格劳孔}
\en{They would have to.}
\zh{必然如此。}

\speaker{SOCRATES}{苏格拉底}
\en{What if their prison also had an echo from the wall facing them? When one of the carriers passing along the wall spoke, do you think they would believe that anything other than the shadow passing in front of them was speaking?}
\zh{如果他们的囚室还会从对面的墙壁传来回声，又会怎样？当一个沿墙走过的搬运者说话时，你认为他们会相信说话的不是眼前经过的影子，而是别的什么东西吗？}

\speaker{GLAUCON}{格劳孔}
\en{I do not, by Zeus.}
\zh{宙斯为证，我不认为会。}

\speaker{SOCRATES}{苏格拉底}
\en{All in all, then, what the prisoners would take for true reality is nothing other than the shadows of those artifacts.}
\zh{总而言之，囚徒们所当作真正实在的，不会是别的，只会是那些器物的影子。}

\speaker{GLAUCON}{格劳孔}
\en{That's entirely inevitable.}
\zh{这完全不可避免。}

\speaker{SOCRATES}{苏格拉底}
\en{Consider, then, what being released from their bonds and cured of their foolishness would naturally be like, if something like this should happen to them. When one was freed and suddenly compelled to stand up, turn his neck around, walk, and look up toward the light, he would be pained by doing all these things and be unable to see the things whose shadows he had seen before, because of the flashing lights. What do you think he would say if we told him that what he had seen before was silly nonsense, but that now--because he is a bit closer to what is, and is turned toward things that are more--he sees more correctly? And in particular, if we pointed to each of the things passing by and compelled him to answer what each of them is, don't you think he would be puzzled and believe that the things he saw earlier were more truly real than the ones he was being shown?}
\zh{那么请考虑一下：如果这样的事发生在他们身上，他们摆脱束缚、治愈愚昧，按自然情形会是什么样子。当其中一人被释放，并突然被迫站起、转动脖子、行走、抬头看向光明时，做这些事都会使他痛苦；由于光芒刺眼，他也无法看见那些过去只见其影子的东西。如果我们告诉他，他从前所见只是无谓的幻象，而现在--因为他稍微更接近存在，并转向更具有实在性的事物--他看得更正确，你认为他会怎么说？尤其是，如果我们指着每一个经过的东西，强迫他回答它是什么，你不认为他会感到困惑，并相信他早先所见比现在展示给他的东西更真实地存在吗？}

\speaker{GLAUCON}{格劳孔}
\en{Much more so.}
\zh{肯定会更这样认为。}

\speaker{SOCRATES}{苏格拉底}
\en{And if he were compelled to look at the light itself, wouldn't his eyes be pained and wouldn't he turn around and flee toward the things he is able to see, and believe that they are really clearer than the ones he is being shown?}
\zh{如果他被迫去看光本身，他的眼睛岂不会疼痛？他岂不会转身逃向那些自己能够看见的东西，并相信它们实际上比正在展示给他的那些东西更清楚吗？}

\speaker{GLAUCON}{格劳孔}
\en{He would.}
\zh{他会的。}

\speaker{SOCRATES}{苏格拉底}
\en{And if someone dragged him by force away from there, along the rough, steep, upward path, and did not let him go until he had dragged him into the light of the sun, wouldn't he be pained and angry at being treated that way? And when he came into the light, wouldn't he have his eyes filled with sunlight and be unable to see a single one of the things now said to be truly real?}
\zh{如果有人用强力把他从那里拖走，沿着崎岖、陡峭、向上的道路前行，直到把他拖进太阳之光中才放手，他岂不会因为受到这样的对待而痛苦和愤怒？当他来到光中，眼里充满阳光时，岂不会连现在所谓真正实在之物中的任何一个都看不见吗？}

\speaker{GLAUCON}{格劳孔}
\en{No, he would not be able to--at least not right away.}
\zh{是的，他不能，至少不会立刻能。}

\speaker{SOCRATES}{苏格拉底}
\en{He would need time to get adjusted, I suppose, if he is going to see the things in the world above. At first, he would see shadows most easily, then images of men and other things in water, then the things themselves. From these, it would be easier for him to go on to look at the things in the sky and the sky itself at night, gazing at the light of the stars and the moon, than during the day, gazing at the sun and the light of the sun.}
\zh{我想，如果他要看见上方世界中的事物，就需要时间来适应。起初，他最容易看见影子；接着，看见水中人和其他事物的映像；然后，才看见事物本身。由此再进一步，他在夜间观看天上的事物和天空本身，凝视星月之光，会比在白昼凝视太阳及太阳之光更容易。}

\speaker{GLAUCON}{格劳孔}
\en{Of course.}
\zh{当然。}

\speaker{SOCRATES}{苏格拉底}
\en{Finally, I suppose, he would be able to see the sun--not reflections of it in water or some alien place, but the sun just by itself in its own place--and be able to look at it and see what it is like.}
\zh{最后，我想，他就能够看见太阳--不是它在水中或异处的倒影，而是太阳在自身所在之处、就其自身而言的样子--并且能够直视它，看清它是什么样。}

\speaker{GLAUCON}{格劳孔}
\en{Necessarily.}
\zh{必然如此。}

\speaker{SOCRATES}{苏格拉底}
\en{After that, he would already be able to conclude about it that it provides the seasons and the years, governs everything in the visible world, and is in some way the cause of all the things that he and his fellows used to see.}
\zh{在那之后，他便已经能够关于太阳得出结论：它提供季节和岁月，统摄可见世界中的一切，并且在某种意义上，是他和同伴们过去所见一切的原因。}

\speaker{GLAUCON}{格劳孔}
\en{That would clearly be his next step.}
\zh{显然，这会是他的下一步。}

\speaker{SOCRATES}{苏格拉底}
\en{What about when he reminds himself of his first dwelling place, what passed for wisdom there, and his fellow prisoners? Don't you think he would count himself happy for the change and pity the others?}
\zh{那么，当他回想自己最初的住处、那里所谓的智慧以及那些同伴囚徒时，又会怎样？你不认为他会为自己的改变而庆幸，并怜悯其他人吗？}

\speaker{GLAUCON}{格劳孔}
\en{Certainly.}
\zh{当然会。}

\speaker{SOCRATES}{苏格拉底}
\en{And if there had been honors, praises, or prizes among them for the one who was sharpest at identifying the shadows as they passed by; and was best able to remember which usually came earlier, which later, and which simultaneously; and who was thus best able to prophesy the future, do you think that our man would desire these rewards or envy those among the prisoners who were honored and held power? Or do you think he would feel with Homer that he would much prefer ``to work the earth as a serf for another man, a man without possessions of his own,'' and go through any sufferings, rather than share their beliefs and live as they do?}
\zh{如果他们中间曾有荣誉、赞扬或奖赏，颁给那个最敏锐地辨认经过影子的人；颁给那个最能记住哪些影子通常在先、哪些在后、哪些同时出现的人；也就是颁给那个最能预言未来的人，那么你认为我们这个人会渴望这些奖赏，或嫉妒那些在囚徒中受尊荣、掌权的人吗？还是说，你认为他会同荷马一起觉得，自己宁愿“替另一个人，一个没有自己财产的人，像雇工那样耕作土地”，并忍受任何痛苦，也不愿分享他们的信念、过他们那样的生活？}

\speaker{GLAUCON}{格劳孔}
\en{Yes, I think he would rather suffer anything than live like that.}
\zh{是的，我认为他宁愿忍受任何事，也不愿那样生活。}

\speaker{SOCRATES}{苏格拉底}
\en{Consider this too, then. If this man went back down into the cave and sat down in his same seat, wouldn't his eyes be filled with darkness, coming suddenly out of the sun like that?}
\zh{那么也请考虑这一点。如果这个人再下到洞穴中，坐回原来的座位，他这样突然从太阳下回来，眼睛岂不会充满黑暗吗？}

\speaker{GLAUCON}{格劳孔}
\en{Certainly.}
\zh{当然会。}

\speaker{SOCRATES}{苏格拉底}
\en{Now, if he had to compete once again with the perpetual prisoners in recognizing the shadows, while his sight was still dim and before his eyes had recovered, and if the time required for readjustment was not short, wouldn't he provoke ridicule? Wouldn't it be said of him that he had returned from his upward journey with his eyes ruined, and that it is not worthwhile even to try to travel upward? And as for anyone who tried to free the prisoners and lead them upward, if they could somehow get their hands on him, wouldn't they kill him?}
\zh{现在，如果他在视力仍然昏暗、眼睛还没有恢复的时候，又不得不再次和那些终身囚徒比赛辨认影子，而且重新适应所需的时间并不短，他岂不会招来嘲笑？人们岂不会说，他完成上行之旅回来，眼睛已经坏了，甚至连尝试上行都不值得？至于任何试图释放囚徒并把他们带向上方的人，如果他们能设法抓住他，岂不会杀了他吗？}

\speaker{GLAUCON}{格劳孔}
\en{They certainly would.}
\zh{他们一定会。}

\speaker{SOCRATES}{苏格拉底}
\en{This image, my dear Glaucon, must be fitted together as a whole with what we said before. The realm revealed through sight should be likened to the prison dwelling, and the light of the fire inside it to the sun's power. And if you think of the upward journey and the seeing of things above as the upward journey of the soul to the intelligible realm, you won't mistake my intention--since it is what you wanted to hear about. Only the god knows whether it is true. But this is how these phenomena seem to me: in the knowable realm, the last thing to be seen is the form of the good, and it is seen only with toil and trouble. Once one has seen it, however, one must infer that it is the cause of all that is correct and beautiful in anything, that in the visible realm it produces both light and its source, and that in the intelligible realm it controls and provides truth and understanding; and that anyone who is to act sensibly in private or public must see it.}
\zh{亲爱的格劳孔，这个图像必须作为一个整体，同我们先前所说的东西配合起来理解。由视觉显现的领域，应当比作囚居之所；其中火的光，应当比作太阳的力量。而如果你把上行之旅以及观看上方事物，理解为灵魂上升到可知领域的旅程，你就不会误解我的意思--因为这正是你想听的。至于它是否为真，只有神知道。但在我看来，情形是这样：在可知领域中，最后被看见的是善之理念（善的理型），而且只有经过辛劳困苦才能看见。然而，一旦看见它，就必须推知：它是一切事物中正当者与美好者的原因；在可见领域中，它产生光及光源；在可知领域中，它主宰并提供真理和理解；并且，任何想要在私人事务或公共事务中明智行事的人，都必须看见它。}

\speaker{GLAUCON}{格劳孔}
\en{I agree, so far as I am able.}
\zh{就我所能理解的而言，我同意。}

\speaker{SOCRATES}{苏格拉底}
\en{Come on, then, and join me in this further thought: you should not be surprised that the ones who get to this point are not willing to occupy themselves with human affairs, but that, on the contrary, their souls are always eager to spend their time above. I mean, that is surely what we would expect, if indeed the image I described before is also accurate here.}
\zh{那么，请继续同我一起思考这一点：那些达到这一境地的人不愿忙于人间事务，反而其灵魂总是渴望在上方停留，对此你不应感到惊讶。我的意思是，如果我先前描述的图像在这里也确实恰当，这当然正是我们会预期的情形。}

\speaker{GLAUCON}{格劳孔}
\en{It is what we would expect.}
\zh{这正是我们会预期的。}

\speaker{SOCRATES}{苏格拉底}
\en{What about when someone, coming from looking at divine things, looks to the evils of human life? Do you think it is surprising that he behaves awkwardly and appears completely ridiculous, if--while his sight is still dim and he has not yet become accustomed to the darkness around him--he is compelled, either in the courts or elsewhere, to compete about the shadows of justice, or about the statues of which they are the shadows; and to dispute the way these things are understood by people who have never seen justice itself?}
\zh{那么，当一个人刚观看过神圣事物，转而察看人类生活的诸种恶时，又会怎样？如果--在他的视力仍然昏暗、还没有习惯周围黑暗的时候--他被迫在法庭或别处，就正义的影子，或就投下那些影子的雕像展开争竞，并且同那些从未见过正义本身的人争论这些东西应如何理解，你认为他举止笨拙、显得十分可笑，会令人惊讶吗？}

\speaker{GLAUCON}{格劳孔}
\en{It is not surprising at all.}
\zh{这根本不令人惊讶。}

\begin{center}
\textit{[\dots]}
\end{center}

\safeincludeimage{c71e19475dd43f2902a91e328edeb2062a34b8a6f56720da0301812d3489866e.jpg}
\safeincludeimage{aaa4581cac55b98fb3bd196278730238d2964fc259e31265903b352a51812148.jpg}
\safeincludeimage{2509989365e4ca64c8315dcf70650facf7be0779ef7a983d4ed5c3e65ad553a6.jpg}


% ===== 02_lindberg_beginnings.tex =====
\section*{The Beginnings of Western Science / 《西方科学的起源》}
\textbf{David C. Lindberg / 戴维 C. 林德伯格}

\section*{Chapter Two: The Greeks and the Cosmos / 第二章：希腊人与宇宙}

\subsection*{Plato's World of Forms / 柏拉图的形式世界}

\orig The death of Socrates in 399 B.C., coming as it did around the turn of the century (not on their calendar, of course, but on ours), has made it a convenient point of demarcation in the history of Greek philosophy. Thus Socrates' predecessors of the sixth and fifth centuries (the philosophers who have occupied us until now in this chapter) are commonly called the ``pre-Socratic philosophers.'' But Socrates' prominence is more than an accident of the calendar, for Socrates represents a shift in emphasis within Greek philosophy, away from the cosmological concerns of the sixth and fifth centuries toward political and ethical matters. Nonetheless, the shift was not so dramatic as to preclude continuing attention to the major problems of pre-Socratic philosophy. We find both the new and the old in the work of Socrates' younger friend and disciple, Plato (fig. 2.4).

\zhline{苏格拉底死于公元前399年，恰好接近一个世纪之交（当然不是按他们的历法，而是按我们的历法），因此成了划分希腊哲学史的一个方便界标。于是，公元前六、五世纪那些早于苏格拉底的哲学家（也就是本章迄今讨论的那些哲学家）通常被称为“前苏格拉底哲学家”。但是，苏格拉底的重要性并不只是历法上的偶然；他代表了希腊哲学重心的一次转移：从公元前六、五世纪的宇宙论关怀，转向政治和伦理问题。尽管如此，这种转移并没有剧烈到使人们不再关注前苏格拉底哲学的主要问题。在苏格拉底较年轻的朋友和弟子柏拉图的著作中，我们既能看到新的关怀，也能看到旧的问题（图2.4）。}

\begin{figure}[htbp]
\centering
\safeincludeimage{76f2553df76f729a90b5a54f134b9c18710e62ab8ac46e3c53b24a76fd0f8cbb.jpg}
\caption{Fig. 2.4. Plato (1st c. A.D. copy). Museo Vaticano, Vatican City. Alinari/Art Resource N.Y. / 图2.4 柏拉图（公元1世纪摹本）。梵蒂冈城，梵蒂冈博物馆。Alinari/Art Resource N.Y.}
\end{figure}

\orig Plato (427--348/47 B.C.) was born into a distinguished Athenian family, active in affairs of state; he was undoubtedly a close observer of the political events that led up to Socrates' execution. After Socrates' death, Plato left Athens and visited Italy and Sicily, where he seems to have come into contact with Pythagorean philosophers. In 388 Plato returned to Athens and founded a school of his own, the Academy, where young men could pursue advanced studies (see fig. 4.1). Plato's literary output appears to have consisted almost entirely of dialogues, the majority of which have survived. We will find it necessary to be highly selective in our examination of Plato's philosophy; let us begin with his quest for the underlying reality.

\zhline{柏拉图（公元前427--348/47年）出生于一个显赫的雅典家族，这个家族积极参与城邦政治；毫无疑问，他密切目睹了导致苏格拉底被处死的一系列政治事件。苏格拉底死后，柏拉图离开雅典，游历意大利和西西里；在那里，他似乎接触到了毕达哥拉斯学派的哲学家。公元前388年，柏拉图回到雅典，创办了自己的学校“学园”（Academy），年轻人在那里可以从事高等研究（见图4.1）。柏拉图的著作似乎几乎全由对话构成，其中大部分保存至今。我们在考察柏拉图哲学时必须高度取舍；先从他对根本实在的追寻说起。}

\begin{figure}[htbp]
\centering
\safeincludeimage{5dec5d9ba0cfe21df1e401825ccfb025e409a982dea1eaad5bbd1158e8374444.jpg}
\caption{Fig. 4.1. The schools of Hellenistic Athens. \copyright{} Candace H. Smith. / 图4.1 希腊化时期雅典的各学派。\copyright{} Candace H. Smith.}
\end{figure}

\orig In a passage in one of his dialogues, the \textit{Republic}, Plato reflected on the relationship between the actual tables constructed by a carpenter and the idea or definition of a table in the carpenter's mind. The carpenter replicates the mental idea as closely as possible in each table he makes, but always imperfectly. No two manufactured tables are alike down to the smallest detail, and limitations in the material (a knot here, a warped board there) ensure that none will fully measure up to the ideal.

\zhline{在其对话《理想国》的一段文字中，柏拉图思考了木匠实际制作的桌子与木匠心中“桌子”的理念或定义之间的关系。木匠在每一张桌子上都尽可能复制心中的理念，但这种复制总是不完美。没有两张制成的桌子在最细微之处完全相同，材料的局限（这里有个木节，那里有块弯曲的木板）也保证了没有哪张桌子能完全达到理想。}

\orig Now, Plato argued, there is a divine craftsman who bears the same relationship to the cosmos as the carpenter bears to his tables. The divine craftsman (the Demiurge) constructed the cosmos according to an idea or plan, so that the cosmos and everything in it are replicas of eternal ideas or forms---but always imperfect replicas because of limitations inherent in the materials available to the Demiurge. In short, there are two realms: a realm of forms or ideas, containing the perfect form of everything; and the material realm in which these forms or ideas are imperfectly replicated.

\zhline{柏拉图于是主张，有一位神圣工匠与宇宙的关系，正如木匠与其桌子的关系一样。这位神圣工匠（德谟革，Demiurge）按照某种理念或计划构造了宇宙；因此，宇宙以及其中的一切都是永恒理念或形式的摹本。但由于德谟革所能使用的材料本身有限，这些摹本总是不完美。简言之，存在两个领域：一个是形式或理念的领域，包含万物的完美形式；另一个是质料领域，在其中，这些形式或理念被不完美地复制出来。}

\orig Plato's notion of two distinct realms will seem strange to many people, and we must therefore stress several points of importance. The forms are incorporeal, intangible, and insensible; they have always existed, sharing the property of eternity with the Demiurge; and they are absolutely changeless. They include the form, the perfect idea, of everything in the material world. One does not speak of their location, since they are incorporeal and therefore not spatial. Although incorporeal and imperceptible by the senses, they objectively exist; indeed, true reality (reality in its fullness) is located only in the world of forms. The sensible, corporeal world, by contrast, is imperfect and transitory. It is less real in the sense that the corporeal object is a replica of, and therefore dependent for its existence upon, the form. The form has primary existence, its corporeal replica secondary existence.

\zhline{柏拉图关于两个不同领域的观念，对许多人来说会显得陌生，因此我们必须强调几个要点。形式是无形体、不可触摸、不可感知的；它们始终存在，与德谟革一样具有永恒性；并且绝对不变。它们包括质料世界中一切事物的形式，也就是其完美理念。我们不能谈论它们的位置，因为它们没有形体，因而也不占据空间。尽管形式无形体且不能由感官感知，它们却客观存在；事实上，真正的实在（充分意义上的实在）只位于形式世界。相比之下，可感的、有形的世界是不完美且短暂的。它之所以较少真实，是因为有形事物只是形式的摹本，因此其存在依赖于形式。形式具有首要存在，它的有形摹本则具有次级存在。}

\orig Plato illustrated this conception of reality in his famous ``allegory of the cave,'' found in book VII of the \textit{Republic}. Men are imprisoned within a deep cave, chained so as to be incapable of moving their heads. Behind them is a wall, and beyond that a fire. People walk back and forth behind the wall, holding above it various objects, including statues of humans and animals; the objects cast shadows on the wall that is visible to the prisoners. The prisoners see only the shadows cast by these objects; and, having lived in the cave from childhood, they no longer recall any other reality. They do not suspect that these shadows are but imperfect images of objects that they cannot see; and consequently they mistake the shadows for the real.

\zhline{柏拉图在《理想国》第七卷著名的“洞穴寓言”中说明了这种实在观。人们被囚禁在一个深洞之中，身上缚有锁链，头也无法转动。他们身后有一堵墙，墙后有一团火。有人在墙后往来行走，高举着各种物体，包括人和动物的雕像；这些物体把影子投到囚徒能够看见的墙面上。囚徒只能看见这些物体投下的影子；由于他们自幼便生活在洞中，早已不记得任何别的实在。他们并不怀疑这些影子只是他们看不见的物体的不完美影像；于是，他们误把影子当成真实。}

\orig So it is with all of us, says Plato. We are souls imprisoned in bodies. The shadows of the allegory represent the world of sense experience. The soul, peering out from its prison, is able to perceive only these flickering shadows, and the ignorant claim that this is all there is to reality. However, there do exist the statues and other objects of which the shadows are feeble representations and also the humans and animals of which the statues are imperfect replicas. To gain access to these higher realities, we must escape the bondage of sense experience and climb out of the cave, until we find ourselves able, finally, to gaze on the eternal realities, thereby entering the realm of true knowledge.

\zhline{柏拉图说，我们所有人都是如此。我们是被囚禁在身体中的灵魂。寓言中的影子代表感官经验的世界。灵魂从它的牢笼中向外窥视，只能感知这些闪烁不定的影子；无知者便声称这就是实在的全部。然而，确实存在着那些雕像和其他物体，影子只是它们微弱的表象；也确实存在着人和动物，雕像只是它们不完美的摹本。若要接近这些更高的实在，我们必须摆脱感官经验的束缚，爬出洞穴，直到最终能够凝视那些永恒实在，从而进入真正知识的领域。}

\orig What are the implications of these views for the concerns of the pre-Socratic philosophers? First, Plato equated his forms with the underlying reality, while assigning derivative or secondary existence to the corporeal world of sensible things. Second, Plato has made room for both change and stability by assigning each to a different level of reality: the corporeal realm is the scene of imperfection and change, while the realm of forms is characterized by eternal, changeless perfection. Both change and stability are therefore genuine; each characterizes something; but changelessness belongs to the forms and thus shares their fuller reality.

\zhline{这些观点对于前苏格拉底哲学家的关切意味着什么？第一，柏拉图把形式等同于根本实在，同时把可感事物的有形世界置于派生的或次级的存在地位。第二，柏拉图通过把变化和稳定分别安置在不同层次的实在中，为二者都留出了位置：有形领域是不完美和变化的场所，而形式领域则以永恒不变的完美为特征。因此，变化和稳定都是真实的；二者各自表征某种事物；但不变性属于形式，因而分享形式那种更充分的实在性。}

\orig Third, as we have seen, Plato addressed epistemological questions, placing observation and true knowledge (or understanding) in opposition. Far from leading upward to knowledge or understanding, the senses are chains that tie us down; the route to knowledge is through philosophical reflection. This is explicit in the \textit{Phaedo}, where Plato maintains the uselessness of the senses for the acquisition of truth and points out that when the soul attempts to employ them it is inevitably deceived.

\zhline{第三，如我们所见，柏拉图处理了认识论问题，把观察与真正的知识（或理解）置于对立之中。感官远非引导我们上升到知识或理解的途径，而是束缚我们的锁链；通向知识的道路在于哲学反思。这一点在《斐多篇》中说得很明确：柏拉图主张感官对于获得真理毫无用处，并指出灵魂一旦试图运用感官，就不可避免地受到欺骗。}

\orig Now the short account of Plato's epistemology frequently ends here; but there are important qualifications that it would be a serious mistake to omit. Plato did not, in fact, dismiss the senses altogether, as Parmenides had done and as the passage from the \textit{Phaedo} might suggest Plato did. Sense experience, in Plato's view, served various useful functions. First, sense experience may provide wholesome recreation. Second, observation of certain sensible objects (especially those with geometrical properties) may serve to direct the soul toward nobler objects in the realm of forms. Plato used this argument as justification for the pursuit of astronomy. Third, Plato argued (in his theory of reminiscence) that sense experience may actually stir the memory and remind the soul of forms that it knew in a prior existence, thus stimulating a process of recollection that will lead to actual knowledge of the forms.

\zhline{对柏拉图认识论的简短说明常常到这里就结束了；但还有一些重要限定，若省略就会造成严重误解。事实上，柏拉图并没有像巴门尼德那样完全排斥感官，尽管《斐多篇》的那段文字可能让人以为他如此。按柏拉图的看法，感官经验具有若干有益功能。第一，感官经验可以提供健康的消遣。第二，观察某些可感对象（尤其是具有几何性质的对象）可以引导灵魂转向形式界中更高贵的对象。柏拉图用这一论证为研究天文学辩护。第三，柏拉图在其“回忆说”中主张，感官经验实际上可以激发记忆，使灵魂想起它在前世曾经认识的形式，从而推动一种回忆过程，并最终导致对形式的真实知识。}

\orig Finally, although Plato firmly believed that knowledge of the eternal forms (the highest, and perhaps the only true, form of knowledge) is obtainable only through the exercise of reason, the changeable realm of matter is also an acceptable object of study. Such studies serve the purpose of supplying examples of the operation of reason in the cosmos. If this is what interests us (as it sometimes did Plato), the best method of exploring it is surely to observe it. The legitimacy and utility of sense experience are clearly implied in the \textit{Republic}, where Plato acknowledged that a prisoner emerging from the cave first employs his sense of sight to apprehend living creatures, the stars, and finally the most noble of visible (material) things, the sun. But if he aspires to apprehend ``the essential reality,'' he must proceed ``through the discourse of reason unaided by any of the senses.'' Both reason and sense are thus instruments worth having; which one we employ on a particular occasion will depend on the object of study.

\zhline{最后，尽管柏拉图坚定地相信，对永恒形式的知识（最高级、或许也是唯一真正的知识形式）只能通过理性的运用来获得，但变化着的质料领域同样可以成为正当的研究对象。这类研究的目的，是为宇宙中理性运作的方式提供例证。如果我们关心的正是这一点（柏拉图有时确实如此），那么探索它的最佳方法无疑就是观察它。《理想国》清楚地暗示了感官经验的正当性和有用性：柏拉图承认，走出洞穴的囚徒首先运用视觉来把握生物、星辰，最后把握可见（质料性）事物中最高贵者，即太阳。但是，如果他想把握“本质的实在”，就必须“通过不借助任何感官的理性论辩”来前进。因此，理性和感官都是值得拥有的工具；在特定场合使用哪一个，取决于研究对象。}

\orig There is another way of expressing all of this, which may shed light on Plato's achievement. When Plato assigned reality to the forms, he was, in fact, identifying reality with the properties that classes of things have in common. The bearer of true reality is not (for example) this dog with the droopy left ear or that one with the menacing bark, but the idealized form of a dog shared (imperfectly, to be sure) by every individual dog---those characteristics by virtue of which we are able to classify all of them as dogs. Therefore, to gain true knowledge, we must set aside all characteristics peculiar to things as individuals and seek the shared characteristics that define them into classes. Now stated in this modest fashion, Plato's view has a distinctly modern ring. Idealization is a prominent feature of a great deal of modern science; we develop models or laws that overlook the incidental in favor of the essential. However, Plato went beyond this, maintaining not merely that true reality is to be found in the common properties of classes of things, but also that this common property (the idea or form) has objective, independent, and indeed prior existence.

\zhline{还可以用另一种方式来表述这一切，这或许能显示柏拉图的成就。柏拉图把实在性赋予形式时，实际上是在把实在等同于各类事物共同具有的性质。真正实在的承载者并不是（例如）这条左耳下垂的狗，或那条叫声凶狠的狗，而是所有个别的狗共同分享（诚然，是不完美地分享）的狗的理想化形式，也就是那些使我们能够把它们统统归为狗的特征。因此，若要获得真正知识，我们必须撇开事物作为个体所特有的一切特征，寻求那些把它们界定为某一类别的共同特征。若以这种较为温和的方式表述，柏拉图的观点便颇具现代意味。理想化是许多现代科学的显著特征；我们建立模型或规律时，会忽略偶然因素而偏重本质。不过，柏拉图走得更远：他不仅主张真正实在存在于各类事物的共同性质中，而且主张这种共同性质（理念或形式）具有客观、独立、并且确实先在的存在。}

\begin{figure}[htbp]
\centering
\safeincludeimagepair{9519836f4685aaff51f66711e7f2b7597d7605556dc3b35ffcbfe2d575d961d5.jpg}{aa581e862cc9d1691c97016b9b4021f9589afd23337d79179d3b7e7d6cbce561.jpg}\hfill
\safeincludeimage{8ea6db286861b81aca334bb0da5448aee2e09ebaf1d510659a5b0cd9f6654773.jpg}
\caption{Uncaptioned image-only material preserved from the OCR source between Chapter 2 and Chapter 3. / OCR 源在第二章与第三章之间保留的无可靠图注图像，按原顺序保留。}
\end{figure}

\section*{Chapter Three: Aristotle's Philosophy of Nature / 第三章：亚里士多德的自然哲学}

\subsection*{Life and Works / 生平与著作}

\orig Aristotle (fig. 3.1) was born in 384 B.C. in the northern Greek town of Stagira, into a privileged family. His father was personal physician to the Macedonian king, Amyntas II (grandfather of Alexander the Great). Aristotle had the advantage of an exceptional education: at age seventeen, he was sent to Athens to study with Plato. He remained in Athens as a member of Plato's Academy for twenty years, until Plato's death about 347. Aristotle then spent several years in travel and study, crossing the Aegean Sea to Asia Minor and its coastal islands. During this period he undertook biological studies, and he encountered Theophrastus (from the island of Lesbos), who was to become his pupil and lifetime colleague. He returned to Macedonia in 342 to become the tutor of the young Alexander (later ``the Great''). In 335, when Athens fell under Macedonian rule, Aristotle returned to the city and began to teach in the Lyceum, a public garden frequented by teachers. He remained there, establishing an informal school, until shortly before his death in 322.

\zhline{亚里士多德（图3.1）于公元前384年出生在希腊北部城镇斯塔吉拉，出身于一个有特权的家庭。他的父亲是马其顿国王阿明塔斯二世（亚历山大大帝的祖父）的私人医生。亚里士多德享有卓越教育的优势：十七岁时，他被送到雅典跟随柏拉图学习。他作为柏拉图学园的一员在雅典停留了二十年，直到柏拉图约于公元前347年去世。此后，亚里士多德用了数年时间旅行和研究，渡过爱琴海前往小亚细亚及其沿岸岛屿。在这一时期，他从事生物学研究，并遇到了泰奥弗拉斯托斯（来自莱斯博斯岛），后者后来成为他的学生和终身同事。公元前342年，他回到马其顿，担任年轻的亚历山大（后来被称为“大帝”）的导师。公元前335年，雅典处于马其顿统治之下时，亚里士多德回到这座城市，并开始在吕克昂（Lyceum）授课；吕克昂是一处教师们常去的公共园地。他在那里建立了一所非正式学校，并一直留到公元前322年去世前不久。}

\orig In the course of his long career as student and teacher, Aristotle systematically and comprehensively addressed the major philosophical issues of his day. He is credited with more than 150 treatises, approximately 30 of which have come down to us. The surviving works appear to consist mainly of lecture notes or unfinished treatises not intended for wide circulation; whatever their exact origin, they were obviously directed to other philosophers, including advanced students. In modern translation, they occupy well over a foot of bookshelf, and they present a philosophical system overwhelming in power and scope. It is out of the question for us to survey the whole of Aristotle's philosophy, and we must be content with examining the fundamentals of his philosophy of nature---beginning with his response to positions taken by the pre-Socratics and Plato.

\zhline{在其漫长的求学和教学经历中，亚里士多德系统而全面地处理了当时的主要哲学问题。传世目录归于他名下的论著超过150部，其中约30部流传至今。现存著作主要看来是讲义或未完成的论著，本不打算广泛流通；无论其确切来源如何，它们显然是写给其他哲学家看的，其中也包括高阶学生。按现代译本计算，它们占据书架的长度远不止一英尺，并呈现出一个力量和范围都令人震撼的哲学体系。我们不可能综览亚里士多德哲学的全貌，只能满足于考察其自然哲学的基本方面，先从他对前苏格拉底哲学家和柏拉图立场的回应开始。}

\begin{figure}[htbp]
\centering
\safeincludeimage{88ca2aa3a48f069b56deb949a5f0e30c086ed3d94423489f429999f388b54bf7.jpg}
\caption{Fig. 3.1. Aristotle. Museo Nazionale, Rome. Alinari/Art Resource N.Y. / 图3.1 亚里士多德。罗马，国立博物馆。Alinari/Art Resource N.Y.}
\end{figure}

\subsection*{Metaphysics and Epistemology / 形而上学与认识论}

\orig Through his long association with Plato, Aristotle had, of course, become thoroughly versed in Plato's theory of forms. Plato had drastically diminished (without totally rejecting) the reality of the material world observed by the senses. Reality in its perfect fullness, Plato argued, is found only in the eternal forms, which are dependent on nothing else for their existence. The objects that make up the sensible world, by contrast, derive their characteristics and their very being from the forms; it follows that sensible objects exist only derivatively or dependently.

\zhline{由于长期与柏拉图相处，亚里士多德当然已经非常熟悉柏拉图的形式论。柏拉图大幅削弱了（但并未完全否定）感官所观察到的质料世界的实在性。柏拉图主张，充分而完美的实在只存在于永恒形式之中，而这些形式的存在不依赖任何其他东西。相比之下，构成可感世界的对象，其特征乃至其存在本身都来自形式；由此可见，可感对象只能以派生的或依赖的方式存在。}

\orig Aristotle refused to accept this diminished, dependent status that Plato assigned to sensible objects. They must exist fully and independently, for in Aristotle's view they were what make up the real world. Moreover, the traits that give an individual object its character do not, Aristotle argued, have a prior and separate existence in a world of forms, but belong to the object itself. There is no perfect form of a dog, for example, existing independently in the world of forms and replicated imperfectly in individual dogs, imparting to them their attributes. For Aristotle, there were just individual dogs. These dogs certainly shared a set of attributes---for otherwise we would not be entitled to call them ``dogs''---but these attributes exist in, and belong to, individual dogs.

\zhline{亚里士多德拒绝接受柏拉图赋予可感对象的这种被削弱的、依赖性的地位。可感对象必须充分且独立地存在，因为在亚里士多德看来，它们正是构成真实世界的东西。此外，亚里士多德论证说，赋予个别对象以特征的那些性状，并不在形式世界中具有先在而分离的存在，而是属于对象本身。比如，并不存在某个狗的完美形式，它独立存在于形式世界中，然后在个别的狗身上被不完美地复制，并把属性赋予它们。对亚里士多德来说，存在的只是个别的狗。这些狗当然共享一组属性，否则我们便无权称它们为“狗”；但这些属性存在于、并属于个别的狗。}

\orig Perhaps this way of viewing the world has a familiar ring. Making individual sensible objects the primary realities (``substances,'' Aristotle called them) will seem like good common sense to most readers of this book, and probably struck Aristotle's contemporaries the same way. But if it makes good common sense, can it also be good philosophy? That is, can it deal successfully, or at least plausibly, with the difficult philosophical issues raised by the pre-Socratics and Plato---the nature of the fundamental reality, epistemological concerns, and the problem of change and stability? Let us take up these problems one by one.

\zhline{这种看待世界的方式也许听起来很熟悉。把个别可感对象视为首要实在（亚里士多德称其为“实体”），对本书大多数读者来说会显得很合乎常识；对亚里士多德的同时代人来说，大概也是如此。但是，合乎常识是否也能成为好的哲学？也就是说，它能否成功地、或至少看似合理地处理前苏格拉底哲学家和柏拉图提出的那些困难哲学问题：根本实在的性质、认识论上的关切，以及变化与稳定的问题？让我们逐一讨论。}

\orig The decision to locate reality in sensible, corporeal objects does not yet tell us very much about reality---only that we should look for it in the sensible world. Already in Aristotle's day, any philosopher would demand to know more: one thing he would demand to know was whether the corporeal materials of daily experience (wood, water, air, stone, metal, flesh, etc.) are themselves the fundamental, irreducible constituents of things, or whether they are composites of still more fundamental stuff. Aristotle addressed this question by drawing a distinction between properties and their subjects. He maintained (as most of us would) that a property has to be the property of something; we call that something its ``subject.'' To be a property is to belong to a subject; properties cannot exist independently.

\zhline{决定把实在安置在可感的、有形的对象中，并没有告诉我们太多关于实在的内容；它只告诉我们应当在可感世界中寻找实在。早在亚里士多德时代，任何哲学家都会要求知道更多：其中一个问题便是，日常经验中的有形材料（木、水、气、石、金属、肉体等）本身是否就是事物的基本且不可还原的构成成分，还是说它们也是由更基本的东西复合而成。亚里士多德通过区分属性及其主体来处理这个问题。他主张（我们大多数人也会这样认为），属性必须是某物的属性；我们把这个“某物”称为属性的“主体”。作为属性就是属于某个主体；属性不能独立存在。}

\orig Individual corporeal objects, then, have both properties (color, weight, texture, and the like) and something other than properties to serve as their subject. These two roles are played by ``form'' and ``matter,'' respectively. Corporeal objects are ``composites'' of form and matter---form consisting of the properties that make the thing what it is, matter serving as the subject or substratum for the form. A white rock, for example, is white, hard, heavy, and so forth, by virtue of its form; but matter must also be present, to serve as subject for the form, and this matter brings no properties of its own to its union with form. (Aristotle's doctrine will be further discussed in chap. 12, below, in connection with medieval attempts to clarify and extend it.)

\zhline{因此，个别有形对象既具有属性（颜色、重量、质地等），也具有某种不是属性、但可以作为属性主体的东西。这两种角色分别由“形式”和“质料”承担。有形对象是形式与质料的“复合物”：形式由使事物成为其所是的那些属性构成，质料则作为形式的主体或基底。例如，一块白色的石头之所以是白的、硬的、重的等等，是凭借它的形式；但质料也必须在场，作为形式的主体，而这种质料在与形式结合时并不带来任何它自身的属性。（亚里士多德这一学说将在下文第12章中结合中世纪对它的澄清和扩展再作讨论。）}

\orig We can never, in actuality, separate form and matter; they are presented to us only as a unitary composite. If they were separable, we should be able to put the properties (no longer the properties of anything) in one pile, the matter (absolutely propertyless) in another---an obvious impossibility. But if form and matter can never be separated, is it not meaningless to speak of them as the real constituents of things? Isn't this a purely logical distinction, existing in our minds, but not in the external world? Surely not for Aristotle, and perhaps not for us; most of us would think twice before denying the real existence of cold or red, although we can never collect a bucket of either one. In short, Aristotle once again surprises us by using commonsense notions to build a persuasive philosophical edifice.

\zhline{实际上，我们永远不能把形式和质料分离开来；它们只作为一个统一复合物呈现在我们面前。如果二者可以分离，我们就应当能够把属性（不再是任何东西的属性）堆成一堆，把质料（绝对没有属性）堆成另一堆；这显然不可能。但是，如果形式和质料永远不能分离，那么把它们称为事物的真实构成要素，难道不是毫无意义吗？这难道不是一种纯粹逻辑上的区分，只存在于我们的心中，而不存在于外部世界吗？至少对亚里士多德来说绝非如此，也许对我们来说也不是；我们大多数人在否认“冷”或“红”的真实存在之前，都会三思，尽管我们永远不能收集一桶“冷”或一桶“红”。简言之，亚里士多德再次让我们惊讶：他用常识性的观念建构出一座有说服力的哲学大厦。}

\orig Aristotle's claim that the primary realities are concrete individuals surely has epistemological implications, since true knowledge must be knowledge of truly real things. By this criterion, Plato's attention was naturally directed toward the eternal forms, knowable through reason or philosophical reflection. Aristotle's metaphysics of concrete individuals, by contrast, directed his quest for knowledge toward the material world of individuals, of nature, and of change---a world encountered through the senses.

\zhline{亚里士多德关于首要实在是具体个体的主张，当然具有认识论含义，因为真正知识必须是关于真正实在之物的知识。按照这一标准，柏拉图的注意力自然转向通过理性或哲学反思而可知的永恒形式。相比之下，亚里士多德关于具体个体的形而上学，则把他对知识的追求引向由个体、自然和变化构成的质料世界，也就是通过感官遭遇到的世界。}

\orig Aristotle's epistemology is complex and sophisticated. It must suffice here to indicate that the process of acquiring knowledge begins with sense experience. From repeated sense experience follows memory; and from memory, by a process of ``intuition'' or insight, the experienced investigator is able to discern the universal features of things. By the repeated observation of dogs, for example, an experienced dog breeder comes to know what a dog really is; that is, he comes to understand the form or definition of a dog, the crucial traits without which an animal cannot be a dog. Note that Aristotle, no less than Plato, was determined to grasp the universal traits or properties of things; but, unlike his teacher, Aristotle argued that one must start with the individual material thing. Once we grasp the universal properties or definition, we can put it to use as the premise of deductive demonstrations.

\zhline{亚里士多德的认识论复杂而精细。这里只需指出，获得知识的过程始于感官经验。反复的感官经验产生记忆；由记忆出发，经过一种“直观”或洞见的过程，有经验的研究者便能辨识事物的普遍特征。举例说，通过反复观察狗，一位有经验的养犬者会认识到狗究竟是什么；也就是说，他会理解狗的形式或定义，即那些缺少了便不能称为狗的关键特征。需要注意，亚里士多德和柏拉图一样，决心把握事物的普遍性状或属性；但是，与其老师不同，亚里士多德主张必须从个别的质料事物出发。一旦我们把握了普遍属性或定义，就可以把它作为演绎证明的前提来使用。}

\orig Knowledge is thus gained by a process that begins with experience (a term broad enough, in some contexts, to include common opinion or the reports of distant observers). In that sense knowledge is empirical; nothing can be known apart from such experience. But what we learn by this ``inductive'' process does not acquire the status of true knowledge until put into deductive form; the end product is a deductive demonstration (nicely illustrated in a Euclidean proof) beginning from universal definitions as premises. Although Aristotle discussed both the inductive and deductive phases (the latter far more than the former) in the acquisition of knowledge, he stopped considerably short of later methodologists, especially in the analysis of induction.

\zhline{因此，知识是通过一个始于经验的过程获得的（“经验”一词在某些语境下足够宽泛，可以包括普通意见或远方观察者的报告）。在这个意义上，知识是经验性的；离开这种经验，什么也无法被认识。但是，我们通过这种“归纳”过程学到的东西，只有在被置入演绎形式之后，才获得真正知识的地位；其最终产物是从普遍定义作为前提出发的演绎证明（欧几里得式证明很好地说明了这一点）。虽然亚里士多德讨论了获得知识过程中的归纳阶段和演绎阶段（后者远多于前者），但他距离后来的方法论者仍相当遥远，尤其是在对归纳的分析方面。}

\orig This is the theory of knowledge outlined by Aristotle in the abstract. Is it also the method actually employed in Aristotle's own scientific investigations? Probably not---with perhaps an occasional exception. Like modern scientists, Aristotle did not proceed by following a methodological recipe book, but rather by rough and ready methods, familiar procedures that had proved themselves in practice. Somebody has defined science as ``doing your damnedest, no holds barred''; when it came (for example) to his extensive biological researches, this is exactly what Aristotle did. It is not a surprise, and certainly no character defect, that Aristotle should, in the course of thinking about the nature and the foundations of knowledge, formulate a theoretical scheme (an epistemology) not perfectly consistent with his own scientific practice.

\zhline{这就是亚里士多德在抽象层面勾勒出的知识理论。那么，这也是亚里士多德本人在科学研究中实际采用的方法吗？大概不是，或许偶有例外。像现代科学家一样，亚里士多德并不是照着一本方法论食谱做事，而是采用粗略但可行的方法，以及那些在实践中证明有效的熟悉程序。有人曾把科学定义为“竭尽全力，不设禁区”；就他广泛的生物学研究而言，亚里士多德正是这样做的。亚里士多德在思考知识的性质和基础时，提出一个与其自身科学实践并不完全一致的理论方案（认识论），这并不令人惊讶，当然也不构成性格缺陷。}

\subsection*{Nature and Change / 自然与变化}

\orig The problem of change had become a celebrated philosophical issue (within the quite small community of philosophers) in the fifth century B.C. In the fourth century, Plato had dealt with it by restricting change to the imperfect material replica of the changeless world of forms. For Aristotle, a distinguished naturalist who was philosophically committed to the full reality of the changeable individuals that make up the sensible world, the problem of change was a most pressing one.

\zhline{到公元前五世纪，变化问题已经成为一个著名的哲学议题（当然是在人数很少的哲学家共同体内部）。到公元前四世纪，柏拉图通过把变化限制在不变形式世界的不完美质料摹本中来处理这一问题。对亚里士多德这位杰出的自然研究者来说，变化问题尤其迫切，因为他在哲学上承认可感世界中那些变化着的个体具有充分实在性。}

\orig Aristotle's starting point was the commonsense assumption that change is genuine. But this does not, by itself, get us very far; it remains to be demonstrated that the idea of change can withstand philosophical scrutiny; it must also be shown how change can be explained. Aristotle had various weapons in his arsenal by which to achieve these ends. The first was his doctrine of form and matter. If every object is constituted of form and matter, then Aristotle could make room for both change and stability by arguing that when an object undergoes change, its form changes (by a process of replacement, the new form replacing the old one) while its matter remains unchanged. Aristotle went on to argue that change in form takes place between a pair of opposites or contraries, one of which is the form to be achieved, the other its privation or absence. When the dry becomes wet, or the cold becomes hot, this is change from privation (dry or cold) to the intended form (wet or hot). Change, for Aristotle, is thus never random, but confined to the narrow corridor connecting pairs of contrary qualities; order is thus discernible even in the midst of change.

\zhline{亚里士多德的出发点，是一个常识性假设：变化是真实的。但仅凭这一点并不能走得很远；还必须证明变化观念经得起哲学审查，也必须说明变化如何能够得到解释。为达到这些目的，亚里士多德有多种理论武器。第一种便是他的形式与质料学说。如果每个对象都由形式和质料构成，那么亚里士多德就可以通过如下论证为变化和稳定同时留出位置：当对象发生变化时，它的形式改变了（通过替换过程，新形式取代旧形式），而它的质料保持不变。亚里士多德进一步论证说，形式的变化发生在一对相反者或对立性质之间，其中一个是将要实现的形式，另一个则是它的缺失或不在。当干变成湿，或冷变成热时，这就是从缺失（干或冷）到预期形式（湿或热）的变化。因此，在亚里士多德看来，变化从不是随机的，而是被限制在连接成对对立性质的狭窄通道之中；即便在变化之中，秩序也清晰可辨。}

\orig A determined Parmenidean might protest that to this point the analysis does nothing to escape Parmenides' objection to all change on the ground that inevitably it calls for the emergence of something out of nothing. Aristotle's reply is found in his doctrine of potentiality and actuality. Aristotle would undoubtedly have granted that if the only two possibilities are being and nonbeing---that is, if things either exist or do not exist---then the transition from non-hot to hot would indeed involve passage from nonbeing to being (the nonexistence of hot to the existence of hot) and would thus be vulnerable to Parmenides' objection. But Aristotle believed that the objection could be successfully circumvented by supposing that there are three categories associated with being instead of two: not just being and nonbeing, but (1) nonbeing, (2) potential being, and (3) actual being. If such is the state of things, then change can occur between potential being and actual being without nonbeing ever entering the picture. What Aristotle has in mind is perhaps most easily illustrated by examples from the biological realm. An acorn is potentially, but not actually, an oak tree. In becoming an oak tree, it becomes actually what it originally was only potentially. The change thus involves passage from potentiality to actuality---not from nonbeing to being, but from one kind or degree of being to another. Or for a pair of nonbiological examples, a heavy body held above the earth falls in order to fulfill its potential of being situated with other heavy things about the center of the universe. And a sculptor, with mallet and hammer, reveals in actuality a shape that existed potentially within the original block of marble.

\zhline{一个坚定的巴门尼德信徒可能会抗议说，到目前为止，这一分析并没有避开巴门尼德对一切变化的反对：变化不可避免地要求某物从无中产生。亚里士多德的回答体现在他的潜能与现实学说中。亚里士多德无疑会承认，如果只有两种可能，即存在和非存在，或者说事物要么存在、要么不存在，那么从“非热”到“热”的转变确实会涉及从非存在到存在的过渡（从热的不存在到热的存在），因而会受到巴门尼德反驳的攻击。但是，亚里士多德认为，只要假定与存在相关的范畴不是两个而是三个，即不只是存在和非存在，而是（1）非存在，（2）潜在存在，（3）现实存在，就可以成功绕开这一反驳。若事情如此，变化就可以发生在潜在存在与现实存在之间，而无需让非存在进入图景。亚里士多德的意思，也许最容易用生物领域的例子来说明。一颗橡子潜在地是橡树，但现实地还不是橡树。在成为橡树时，它现实地成为了它原先只是潜在地所是的东西。因此，这种变化涉及从潜能到现实的过渡，不是从非存在到存在，而是从一种或一个程度的存在到另一种或另一个程度的存在。再举两个非生物的例子：悬在地面上方的重物下落，是为了实现它与其他重物一起位于宇宙中心附近的潜能；雕刻家用槌和凿，把原本潜在存在于大理石块中的形状现实地显现出来。}

\orig If these arguments allow us to escape the logical dilemmas associated with the idea of change, and therefore to believe in its possibility, they do not yet tell us anything about the cause of change. Why should an acorn move from the status of potential oak tree to that of actual oak tree, or an object change from black to white, rather than remaining in its original state? Aristotle answered with an intricate, subtle, and not always consistent, theory of nature and causation. Given these difficulties, we will spare ourselves the pain of an exhaustive account and treat ourselves to the short version.

\zhline{如果这些论证使我们得以摆脱变化观念所伴随的逻辑困境，从而相信变化是可能的，它们却尚未告诉我们变化的原因是什么。为什么橡子会从潜在橡树的状态转向现实橡树的状态？为什么一个对象会由黑变白，而不是停留在原有状态？亚里士多德用一套复杂、细致、且并非总是一贯的自然与因果理论来回答。鉴于这些困难，我们不必忍受详尽说明的痛苦，采用简短版本即可。}

\orig The world we inhabit is an orderly one, in which things generally behave in predictable ways, Aristotle argued, because every natural object has a ``nature''---an attribute (associated primarily with form) that makes the object behave in its customary fashion, provided no insurmountable obstacle intervenes---or, as a modern commentator has put it, ``that within a thing which determines basically what that thing does when it is being itself.'' For Aristotle, a brilliant zoologist, the growth and development of biological organisms were easily explained by the activity of such an inner driving force. An acorn becomes an oak tree because its nature is to do so. But the theory was applicable beyond biological growth and, indeed, beyond the biological realm altogether. Dogs bark, rocks fall, and marble yields to the hammer and chisel of the sculptor because of their respective natures. Ultimately, Aristotle argued, all change and motion in the universe can be traced back to the natures of things. For the natural philosopher, who by definition is interested in change and things capable of undergoing change, these natures are the central object of study.

\zhline{亚里士多德主张，我们所居住的世界是一个有秩序的世界，其中事物通常以可预期的方式行动，因为每个自然对象都有一种“自然本性”：这是一种属性（主要与形式相关），在没有不可克服的障碍介入时，使对象以其惯常方式行动；或者用一位现代评论者的话说，它是“事物内部那种基本决定该事物在成为其自身时会做什么的东西”。对杰出的动物学家亚里士多德来说，生物有机体的生长和发展，很容易由这种内在驱动力的活动来解释。橡子之所以成为橡树，是因为它的自然本性就是如此。但这一理论不仅适用于生物生长，事实上也不仅适用于生物领域。狗会吠叫，石头会下落，大理石会屈从于雕刻家的锤凿，都是由于它们各自的自然本性。亚里士多德最终主张，宇宙中所有变化和运动都可以追溯到事物的自然本性。对于自然哲学家而言，这些本性就是研究的中心对象；因为自然哲学家按定义所关心的，正是变化以及能够发生变化的事物。}

\orig To this general statement of Aristotle's theory of ``nature,'' we need to add a qualification---namely, that an artificially produced object is a special case, for such an object possesses no nature other than the natures of its ingredients. If a chariot is constructed of wood and iron, the nature of wood and nature of iron do not yield to a composite ``nature of a chariot.'' By contrast, in the organic world the natures of the organs and tissues that make up an organism yield to the nature of the organism as a whole. The nature of the human body is not the sum of the natures of its various tissues and organs, but a unique nature characteristic of that living human as an organic whole.

\zhline{对亚里士多德“自然”理论的这一概括性说明，还需要加上一项限定：人工制成的对象是特殊情况，因为这种对象除了其组成成分的自然本性之外，并不具有别的自然本性。如果一辆战车由木和铁构成，木的本性和铁的本性并不会让位于某种复合的“战车的本性”。相比之下，在有机世界中，组成有机体的器官和组织的本性，会让位于作为整体的有机体的本性。人体的本性不是其各种组织和器官本性的总和，而是作为有机整体的活生生的人所特有的一种独特本性。}

\orig With this theory of nature in mind, we can understand a feature of Aristotle's scientific practice that has puzzled and distressed modern commentators and critics---namely, the absence from his work of anything resembling controlled experimentation. Unfortunately, such criticism overlooks Aristotle's aims, which drastically limited his methodological options. If, as Aristotle believed, the nature of a thing is to be discovered through the behavior of that thing in its natural, unfettered state, then artificial constraints will merely interfere and corrupt. If, despite interference, the object behaves in its customary fashion, we have troubled ourselves for no purpose. If we set up conditions that prevent the nature of an object from revealing itself, all we have learned is that it can be interfered with to the point of remaining concealed. Contrived experimentation violates, rather than reveals, the natures of things. Aristotle's scientific practice is not to be explained, therefore, as a result of stupidity or deficiency on his part---failure to perceive an obvious procedural improvement---but as a method compatible with the world as he perceived it and suited to the questions that interested him. Experimental science emerged not when, at long last, the human race produced somebody clever enough to perceive that artificial conditions would assist in the exploration of nature, but when a rich variety of conditions were fulfilled---including the emergence of questions to which such a procedure promised to provide answers.

\zhline{牢记这一自然理论，我们就能理解亚里士多德科学实践中一个令现代评论者和批评者困惑并不满的特征：他的著作中缺少任何类似受控实验的东西。遗憾的是，这类批评忽视了亚里士多德的目标，而这些目标极大地限制了他可采用的方法。倘若如亚里士多德所信，事物的自然本性应当通过该事物在自然、无束缚状态下的行为来发现，那么人为约束只会干扰并破坏它。如果在干扰之下，对象仍按其惯常方式行动，我们便白费了力气；如果我们设置的条件阻止对象的自然本性显现，那么我们学到的也只是它可以被干扰到保持隐藏的程度。人为设计的实验不是揭示事物的自然本性，而是违背它们。因此，亚里士多德的科学实践不能解释为他的愚蠢或缺陷，也不是他未能看出某个显而易见的程序改进；它应被理解为一种与他所感知的世界相容、并适合他所关心问题的方法。实验科学之所以出现，并不是因为人类终于产生了一个聪明到足以看出人工条件有助于探索自然的人，而是因为一系列丰富条件得到满足，其中包括出现了这样一些问题：这种程序有望为它们提供答案。}

\orig To complete our analysis of Aristotle's theory of change, we must briefly consider the celebrated four Aristotelian causes. To understand a change or the production of an artifact is to know its causes (perhaps best translated ``explanatory conditions and factors''). There are four of these: the form of a thing; the matter underlying that form, which persists through the change; the agency that brings about the change; and the purpose served by the change. These are called, respectively, formal cause, material cause, efficient cause, and final cause. To take an extremely simple example---the production of a statue---the formal cause is the shape given the marble, the material cause is the marble that receives this shape, the efficient cause is the sculptor, and the final cause is the purpose for which the statue is produced (perhaps the beautification of Athens or the celebration of one of its heroes). There are cases in which identifying one or another of the causes is difficult, or in which one or more causes merge, but Aristotle was convinced that his four causes provided an analytical scheme of general applicability.

\zhline{为完成我们对亚里士多德变化理论的分析，必须简要考察著名的亚里士多德四因。理解一种变化或一件人工制品的产生，就是认识它的原因（也许最好译作“解释性条件和因素”）。这些原因有四种：事物的形式；承载该形式并在变化中持续存在的质料；引起变化的作用者；以及变化所服务的目的。它们分别称为形式因、质料因、动力因和目的因。举一个极为简单的例子：制作一尊雕像。形式因是赋予大理石的形状；质料因是接受这一形状的大理石；动力因是雕刻家；目的因则是制作这尊雕像的目的（也许是美化雅典，或纪念它的一位英雄）。在某些情况下，辨认某一原因会很困难，或者一个或多个原因会相互合并；但亚里士多德坚信，他的四因提供了一种具有普遍适用性的分析框架。}

\orig We have said enough about the form-matter distinction to make clear what was meant by ``formal'' and ``material'' causes, and ``efficient'' cause is close enough to modern notions of causation to require no further comment; but ``final'' cause requires explanation. In the first place, the expression ``final cause'' is an English cognate derived from the Latin word \textit{finis}, meaning ``goal,'' ``purpose,'' or ``end,'' and it has nothing to do with the fact that it often appears last in the list of Aristotelian causes. Aristotle argued, quite rightly, that many things cannot be understood without knowledge of purpose or function. To explain the arrangement of teeth in the mouth, for example, we must understand their functions (sharp teeth in front for tearing, molars in back for grinding). Or to take an example from the inorganic realm, it is not possible to grasp why a saw is made as it is without knowing the function the saw is meant to serve. Aristotle went so far as to give final cause priority over material cause, noting that the purpose of the saw determines the material (iron) of which it must be made, whereas the fact that we possess a piece of iron does nothing to determine that we will make it into a saw.

\zhline{关于形式与质料的区分，我们已经说得足够多，足以说明“形式因”和“质料因”的意思；“动力因”又与现代因果观念相当接近，无需进一步说明。但是，“目的因”需要解释。首先，“final cause”（目的因）这个英文表达源自拉丁词 \textit{finis}，意为“目标”“目的”或“终点”，它与目的因常常出现在亚里士多德诸因列表的最后这一事实没有关系。亚里士多德相当正确地论证说，许多事物若不了解其目的或功能，就无法理解。比如，要解释口腔中牙齿的排列，就必须理解它们的功能（前面的尖牙用于撕裂，后面的臼齿用于研磨）。再举一个无机领域的例子，如果不知道锯子要服务的功能，就不可能把握锯子为什么被制作成这个样子。亚里士多德甚至把目的因置于质料因之上，指出锯子的目的决定了它必须由什么材料（铁）制成；相反，我们拥有一块铁这一事实，并不能决定我们会把它做成锯子。}

\orig Perhaps the most important point to be made about final cause is its clear illustration of the role of purpose (the more technical term is ``teleology'') in Aristotle's universe. The world of Aristotle is not the inert, mechanistic world of the atomists, in which the individual atom pursues its own course mindless of all others. Aristotle's world is not a world of chance and coincidence, but an orderly, organized world, a world of purpose, in which things develop toward ends determined by their natures. It would be unfair and pointless to judge Aristotle's success by the degree to which he anticipated modern science (as though his goal was to answer our questions, rather than his own); it is nonetheless worth noting that the emphasis on functional explanation to which Aristotle's teleology leads would prove to be of profound significance for all of the sciences and remains to this day a dominant mode of explanation within the biological sciences.

\zhline{关于目的因，也许最重要的一点是：它清楚说明了目的（更技术性的术语是“目的论”）在亚里士多德宇宙中的作用。亚里士多德的世界并不是原子论者那种惰性的、机械论的世界；在原子论世界中，个别原子各行其道，对其他一切毫不关心。亚里士多德的世界不是偶然和巧合的世界，而是一个有秩序、有组织的世界，一个有目的的世界；在其中，事物依照其自然本性所规定的终点发展。以亚里士多德在多大程度上预示了现代科学来评判他的成功，是不公平且无意义的（仿佛他的目标是回答我们的问题，而不是他自己的问题）。尽管如此，仍值得注意的是，亚里士多德目的论所导向的功能性解释，对所有科学都将具有深远意义，并且至今仍是生物科学中的一种主导解释方式。}

\subsection*{Cosmology / 宇宙论}

\orig Aristotle not only devised methods and principles by which to investigate and understand the world: form and matter, nature, potentiality and actuality, and the four causes. In the process, he also developed detailed and influential theories regarding an enormous range of natural phenomena, from the heavens above to the earth and its inhabitants below.

\zhline{亚里士多德不仅提出了研究和理解世界的方法与原则：形式和质料、自然本性、潜能与现实，以及四因。在这一过程中，他还就范围极广的自然现象提出了详细而影响深远的理论，涵盖上方的天界以及下方的大地和其中的居民。}

\orig Let us start with the question of origins. Aristotle adamantly denied the possibility of a beginning, insisting that the universe must be eternal. The alternative---that the universe came into being at some point in time---he regarded as unthinkable, violating (among other things) Parmenidean strictures about something coming from nothing. Aristotle's position on this question would prove troublesome for medieval Christian Aristotelians.

\zhline{让我们从起源问题说起。亚里士多德坚决否认宇宙有开端的可能性，坚持认为宇宙必定是永恒的。另一种选择，即宇宙在时间中的某一点开始存在，被他视为不可思议；它违背了（除此之外还包括）巴门尼德关于某物不能从无中产生的禁令。亚里士多德在这一问题上的立场，后来会给中世纪基督教亚里士多德主义者带来麻烦。}

\orig Aristotle considered this eternal universe to be a great sphere, divided into an upper and a lower region by the spherical shell in which the moon is situated. Above the moon is the celestial region; below is the terrestrial region; the moon, spatially intermediate, is also of intermediate nature. The terrestrial or sublunar region is characterized by birth, death, and transient change of all kinds; the celestial or supralunar region, by contrast, is a region of eternally unchanging cycles. That this scheme had its origin in observation would seem clear enough; in his \textit{On the Heavens}, Aristotle noted that ``in the whole range of time past, so far as our inherited records reach, no change appears to have taken place either in the whole scheme of the outermost heaven or in any of its proper parts.'' If in the heavens we observe eternally unvarying circular motion, he continued, we can infer that the heavens are not made of the terrestrial elements, the nature of which (observation reveals) is to rise or fall in transient rectilinear motions. The heavens must consist of an incorruptible fifth element (there are four terrestrial elements): the quintessence (literally, the fifth essence) or aether. The celestial region is completely filled with this quintessence (no void space) and divided, as we shall see, into concentric spherical shells bearing the planets. It had, for Aristotle, a superior, quasi-divine status.

\zhline{亚里士多德认为，这个永恒宇宙是一个巨大的球体，并由月球所在的球壳分成上、下两个区域。月球之上是天界；月球之下是地界；月球在空间上居于中间，其性质也居于中间。地界或月下区域以生成、死亡以及各种短暂变化为特征；相比之下，天界或月上区域则是永恒不变循环的区域。这个图式源自观察，这一点似乎相当清楚；在《论天》中，亚里士多德指出：“在过去的整个时间范围内，就我们传承下来的记录所及，似乎无论在最外层天的整体结构中，还是在其任何固有部分中，都没有发生过变化。”他继续说，如果我们在天上观察到永恒不变的圆周运动，就可以推断天体不是由地上的元素构成的；因为观察显示，地上元素的本性是在短暂的直线运动中上升或下落。天体必定由一种不可朽坏的第五元素构成（地上元素有四种）：第五精华（字面意思即第五种本质），或以太。天界完全由这种第五元素充满（没有虚空），并且如我们将看到的，被划分为承载行星的同心球壳。对亚里士多德来说，天界具有一种较高的、近乎神圣的地位。}

\orig The sublunar region is the scene of generation, corruption, and impermanence. Aristotle, like his predecessors, inquired into the basic element or elements to which the multitude of substances found in the terrestrial region can be reduced. He accepted the four elements originally proposed by Empedocles and subsequently adopted by Plato---earth, water, air, and fire. He agreed with Plato that these elements are in fact reducible to something even more fundamental; but he did not share Plato's mathematical inclination and therefore refused to accept Plato's regular solids and their constituent triangles. Instead, he expressed his own commitment to the reality of the world of sense experience by choosing sensible qualities as the ultimate building blocks. Two pairs of qualities are crucial: hot-cold and wet-dry. These combine in four pairs, each of which yields one of the elements (see fig. 3.2). Notice the use made once again of contraries. There is nothing to forbid any of the four qualities being replaced by its contrary, as the result of outside influence. If water is heated, so that the cold of water yields to hot, the water is transformed into air. Such a process easily explains changes of state (from solid to liquid to vapor, and conversely), but also more general transmutation of one substance into another. On such a theory as this, alchemists could easily build.

\zhline{月下区域是生成、消灭和无常的场所。像他的前辈一样，亚里士多德探究了地界中众多物质可以被还原为何种基本元素。他接受了恩培多克勒最初提出、随后为柏拉图采纳的四元素：土、水、气、火。他同意柏拉图的看法，认为这些元素实际上还能还原为更基本的东西；但是，他并不分享柏拉图的数学倾向，因此拒绝接受柏拉图的正多面体及其构成三角形。相反，他选择可感性质作为最终构成要素，从而表达了自己对感官经验世界之实在性的承诺。两对性质至关重要：热--冷和湿--干。它们组合成四对，每一对产生一种元素（见图3.2）。请注意，对立性质在这里再次发挥作用。外部影响导致四种性质中的任何一种被其相反性质取代，这并无不可。若水被加热，使水的冷让位于热，水就转化为空气。这样的过程不仅容易解释状态变化（从固体到液体到气体，以及反向变化），也容易解释一种物质向另一种物质的更一般的嬗变。炼金术士很容易在这样一种理论上建立自己的学说。}

\orig cold and dry = earth; cold and wet = water; hot and wet = air; hot and dry = fire.

\zhline{冷而干 = 土；冷而湿 = 水；热而湿 = 气；热而干 = 火。}

\begin{figure}[htbp]
\centering
\safeincludeimage{749fd116e948c6da12cfa6db546ce82a7d13b303bbd503ba4bce425b8ec8e85a.jpg}
\caption{Fig. 3.2. Square of opposition of the Aristotelian elements and qualities. For a medieval (9th c.) version of this diagram, see John E. Murdoch, \textit{Album of Science: Antiquity and the Middle Ages}, p. 352. / 图3.2 亚里士多德元素与性质的对立方阵。关于此图的一个中世纪（9世纪）版本，见 John E. Murdoch, \textit{Album of Science: Antiquity and the Middle Ages}, p. 352。}
\end{figure}

\orig The various substances that make up the cosmos totally fill it, leaving no empty space. To appreciate Aristotle's view, we must lay aside our almost automatic inclination to think atomistically; we must conceive material things not as aggregates of tiny particles but as continuous wholes. If it is obvious that, say, a loaf of bread is composed of crumbs separated by small spaces, there is no reason not to suppose that those spaces are filled by some finer substance, such as air or water. And there is certainly no simple way of demonstrating, nor indeed any obvious reason for believing, that water and air are anything but continuous. Similar reasoning, applied to the whole of the universe, led Aristotle to the conclusion that the universe is full, a plenum, containing no void space. This claim would be attacked by medieval scholars.

\zhline{构成宇宙的各种物质把宇宙完全充满，不留下任何空处。为了理解亚里士多德的观点，我们必须暂时放下那种几乎自动出现的原子式思维倾向；必须把质料事物设想为连续整体，而不是微小粒子的聚合体。比如，如果一条面包显然由许多面包屑组成，且面包屑之间隔着细小空隙，那么没有理由不设想这些空隙被某种更精细的物质，如空气或水，所填充。而且，也确实没有简单方法可以证明水和空气不是连续的；事实上，也没有明显理由相信它们不是连续的。类似推理应用到整个宇宙，便使亚里士多德得出结论：宇宙是充满的，是一个充盈体，没有虚空。这个主张后来会受到中世纪学者的攻击。}

\orig Aristotle defended this conclusion with a variety of arguments, such as the following. The speed of a falling body is dependent on the density of the medium through which it falls---the less the density, the swifter the motion of the falling body. It follows that in a void space (density zero), there is nothing to slow the descent of the body, from which we would be forced to conclude that the body would fall with infinite speed---a nonsensical notion, since it implies that the body could be at two places at the same time. Critics have frequently noted that this argument can just as well be taken to prove that the absence of resistance does not entail infinite speed as to prove that void does not exist. The point is, of course, well taken. However, we need to understand that Aristotle's denial of the void did not rest on this single piece of reasoning. In fact, this was but one small part of a lengthy campaign against the atomists, in which Aristotle battled the notion of void space (or void place) with a variety of arguments, some more and some less persuasive.

\zhline{亚里士多德用多种论证来捍卫这一结论，下面便是一例。落体的速度取决于它下落时穿过的介质密度；密度越小，落体运动越快。由此可推得，在虚空中（密度为零），没有任何东西减缓物体下落，于是我们就会被迫得出结论：物体将以无限速度下落。这个观念荒谬无稽，因为它意味着物体可以同时位于两个地方。批评者常常指出，这一论证同样可以被用来证明“没有阻力并不导致无限速度”，而不必证明虚空不存在。当然，这个批评切中要害。不过，我们需要理解，亚里士多德对虚空的否定并不依赖这一条推理。事实上，这只是他长期反对原子论者的一小部分；在这场论战中，亚里士多德用多种论证反击虚空空间（或虚空位置）的观念，其中有的较有说服力，有的较弱。}

\orig In addition to being hot or cold and wet or dry, each of the elements is also heavy or light. Earth and water are heavy, but earth is the heavier of the two. Air and fire are light, fire being the lighter of the two. In assigning levity to two of the elements, Aristotle did not mean (as we might, if we were making the claim) simply that they are less heavy, but that they are light in an absolute sense; levity is not a weaker version of gravity, but its contrary. Because earth and water are heavy, it is their nature to descend toward the center of the universe; because air and fire are light, it is their nature to ascend toward the periphery (that is, the periphery of the terrestrial region, the spherical shell that contains the moon). If there were no hindrances, therefore, earth and water would collect at the center; because of its greater heaviness, earth would achieve a lower position, forming a sphere at the very center of the universe; water would collect in a concentric spherical shell just outside it. Air and fire naturally ascend, but fire, owing to its greater levity, occupies the outermost region, with air as a concentric sphere just inside it. In the ideal case (in which there are no mixed bodies and nothing prevents the natures of the four elements from fulfilling themselves), the elements would thus form a set of concentric spheres: fire on the outside, followed by air and water, and finally earth at the center (see fig. 3.3). But in reality, the world is composed largely of mixed bodies, one always interfering with another, and the ideal is never attained. Nonetheless, the ideal arrangement defines the natural place of each of the elements; the natural place of earth is at the center of the universe, of fire just inside the sphere of the moon, and so forth.

\zhline{除了热或冷、湿或干之外，每一种元素也都是重的或轻的。土和水是重的，但土在二者中更重；气和火是轻的，火在二者中更轻。亚里士多德把轻性赋予两种元素时，并不是说（如我们若提出这一说法时可能会说的那样）它们只是较不重，而是说它们在绝对意义上是轻的；轻性不是重性的较弱版本，而是其相反者。由于土和水是重的，它们的自然本性便是下降到宇宙中心；由于气和火是轻的，它们的自然本性便是上升到外围（也就是地界的外围，包含月球的球壳）。因此，如果没有阻碍，土和水会聚集在中心；由于土更重，它会取得更低的位置，在宇宙最中心形成一个球体；水则会聚集在它外侧的同心球壳中。气和火自然上升；但火由于轻性更大，占据最外层区域，气则作为一个同心球层位于其内侧。在理想情况下（没有混合物体，也没有任何东西阻止四元素的本性实现自身），各元素便会形成一组同心球：外层是火，其次是气和水，最后是中心的土（见图3.3）。但在现实中，世界主要由混合物体构成，彼此总是相互干扰，理想状态从未实现。尽管如此，这种理想排列界定了每种元素的自然位置：土的自然位置在宇宙中心，火的自然位置就在月球天球内侧，等等。}

\orig It must be emphasized that the arrangement of the elements is spherical. Earth collects at the center to form the earth, and it too is spherical. Aristotle defended this belief with a variety of arguments. Arguing from his natural philosophy, he pointed out that since the natural tendency of earth is to move toward the center of the universe, it must arrange itself symmetrically about that point. But he also called attention to observational evidence, including the circular shadow cast by the earth during a lunar eclipse and the fact that north-south motion by an observer on the surface of the earth alters the apparent position of the stars. Aristotle even reported an estimate by mathematicians of the earth's circumference (400,000 stades = about 45,000 miles, roughly 1.8 times the modern value). The sphericity of the earth, thus defended by Aristotle, would never be forgotten or seriously questioned. The widespread myth that medieval people believed in a flat earth is of modern origin.

\zhline{必须强调，元素的排列是球形的。土聚集在中心形成大地，而大地本身也是球形的。亚里士多德用多种论证捍卫这一信念。从他的自然哲学出发，他指出，由于土的自然倾向是朝宇宙中心运动，它必定围绕这一点对称地排列自身。但他也提醒人们注意观察证据，包括月食时地球投下的圆形影子，以及地表观察者向南北方向移动会改变恒星的视位置这一事实。亚里士多德甚至报告了数学家对地球周长的一个估计（400,000斯塔德，约45,000英里，约为现代数值的1.8倍）。亚里士多德如此捍卫的地球球形说，后来从未被遗忘或受到严肃质疑。中世纪人相信地平说这一广泛流传的神话，乃是现代的产物。}

\begin{figure}[htbp]
\centering
\safeincludeimage{2800b314c989309a31b367a2dcd79da8bd2197dacae4678ce6e7f2442cd1e18d.jpg}
\caption{Fig. 3.3. The Aristotelian cosmos. / 图3.3 亚里士多德宇宙。}
\end{figure}

\orig Finally, we must note one of the implications of this cosmology, namely that space, instead of being a neutral, homogeneous backdrop (analogous to our modern notion of geometrical space) against which events occur, has properties. Or to express the point more precisely, ours is a world of space, whereas Aristotle's was a world of place. Heavy bodies move toward their place at the center of the universe not because of a tendency to unite with other heavy bodies located there, but simply because it is their nature to seek that central place; if by some miracle the center happened to be vacant (a physical impossibility in an Aristotelian universe, but an interesting imaginary state of affairs), it would remain the destination of every heavy body.

\zhline{最后，我们必须注意这一宇宙论的一个含义：空间并不是事件发生时一个中性的、同质的背景（类似我们现代的几何空间观念），而是具有属性。或者更准确地说，我们的世界是空间的世界，而亚里士多德的世界则是位置的世界。重物朝向宇宙中心的自然位置运动，并不是因为它们倾向于与那里其他重物结合，而只是因为它们的自然本性就是寻求那个中心位置；即使奇迹般地中心恰好空着（这在亚里士多德宇宙中是物理上不可能的，但作为想象情形颇有趣），它仍会是每一个重物的目的地。}

\subsection*{Motion, Terrestrial and Celestial / 地上运动与天体运动}

\orig We can best understand Aristotle's theory of motion by grasping its two most fundamental claims. The first is that motion is never spontaneous; there is no motion without a mover. The second is the distinction between two types of motion: motion toward the natural place of the moving body is ``natural'' motion; motion in any other direction occurs only under coercion from an outside force and is therefore a ``forced'' or ``violent'' motion.

\zhline{要最好地理解亚里士多德的运动理论，需把握其两个最基本的主张。第一，运动从不是自发的；没有推动者，就没有运动。第二，运动分为两类：朝向运动物体自然位置的运动，是“自然”运动；任何其他方向的运动，都只能在外力强迫下发生，因此是“强迫的”或“暴力的”运动。}

\orig The mover in the case of natural motion is the nature of the body, which is responsible for its tendency to move toward its natural place as defined by the ideal spherical arrangement of the elements. Mixed bodies have a directional tendency that depends on the proportion of the various elements in their composition. When a body undergoing natural motion reaches its natural place, its motion ceases. The mover in the case of forced motion is an external force, which compels the body to violate its natural tendency and move in a direction or manner other than straight-line motion toward its natural place. Such motion ceases when the external force is withdrawn.

\zhline{在自然运动的情况下，推动者是物体的自然本性；这种本性负责使物体倾向于朝向它在元素理想球形排列中所规定的自然位置运动。混合物体的方向倾向取决于其构成中各种元素的比例。当一个进行自然运动的物体到达其自然位置时，它的运动便停止。在强迫运动的情况下，推动者是外在力量；它迫使物体违背其自然倾向，以不同于朝向自然位置的直线运动的方向或方式运动。当外力撤去时，这种运动便停止。}

\orig So far, this seems sensible. One obvious difficulty, however, is to explain why a projectile hurled horizontally, and therefore undergoing forced motion, does not come to an immediate halt when it loses contact with whatever propelled it. Aristotle's answer was that the medium takes over as mover. When we project an object, we also act on the surrounding medium (air, for instance), imparting to it the power to move objects; this power is communicated from part to part, in such a way that the projectile is always in contact with a portion of the medium capable of keeping it in motion. If this seems implausible, consider the greater implausibility (from Aristotle's standpoint) of the alternative---that a projectile, which is inclined by nature to move toward the center of the universe, moves horizontally or upward despite the fact that there is no longer anything causing it to do so.

\zhline{到目前为止，这似乎合情合理。然而，一个明显困难是解释为什么一个被水平抛出的投射物，既然处于强迫运动之中，在失去与推动它的东西的接触后并不会立即停下。亚里士多德的回答是，介质接替成为推动者。当我们投射一个物体时，也作用于周围介质（例如空气），把移动物体的能力赋予介质；这种能力从一个部分传递到另一个部分，使投射物总是与某部分能够维持其运动的介质相接触。如果这看起来难以置信，那么请考虑另一种选择从亚里士多德立场看更加难以置信：一个自然上倾向于朝宇宙中心运动的投射物，竟然在已经没有任何东西使它这样做的情况下，仍然水平或向上运动。}

\orig Force is not the only determinant of motion. In all real cases of motion in the terrestrial realm, there will also be a resistance or opposing force. And it seemed clear to Aristotle that the quickness of motion must depend on these two determining factors---the motive force and the resistance. The question arose: what is the relationship between force, resistance, and speed? Although it probably did not occur to Aristotle that there might be a quantitative law of universal applicability, he was not without interest in the question and did make several forays into quantitative territory. In reference to natural motion in his \textit{On the Heavens} and again in his \textit{Physics}, Aristotle claimed that when two bodies of differing weight descend, the times required to cover a given distance will be inversely proportional to the weights. (A body twice as heavy will require half the time.) In the same chapter of the \textit{Physics}, Aristotle introduced resistance into the analysis of natural motion, arguing that if bodies of equal weight move through media of different densities, the times required to traverse a given distance are proportional to the densities of the respective media; that is, the greater the resistance the slower the body moves. Finally, Aristotle also dealt with forced motion in his \textit{Physics}, claiming that if a given force moves a given weight (against its nature) for a given distance in a given time, the same force will move half that weight twice the distance in that same time (or the same distance in half that time); alternatively, half the force will move half the weight the same distance in the same time.

\zhline{力并不是运动的唯一决定因素。在地界中一切真实运动的情形里，也会存在阻力或反向力。在亚里士多德看来，运动快慢显然必须取决于这两个决定因素：推动力和阻力。问题随之出现：力、阻力和速度之间的关系是什么？虽然亚里士多德大概没有想到可能存在一个普遍适用的定量规律，但他并非对此问题毫无兴趣，并且确实多次涉足定量领域。关于自然运动，亚里士多德在《论天》中、并且又在《物理学》中主张，当两个重量不同的物体下落时，它们通过给定距离所需时间与其重量成反比。（一个重两倍的物体将需要一半时间。）在《物理学》的同一章中，亚里士多德把阻力引入自然运动分析，论证说，如果同等重量的物体穿过不同密度的介质运动，那么它们通过给定距离所需时间与相应介质的密度成正比；也就是说，阻力越大，物体运动越慢。最后，亚里士多德在《物理学》中也处理了强迫运动，声称如果某一给定力能在给定时间内使某一给定重量（违反其本性）移动给定距离，那么同样的力就能在同样时间内使一半重量移动两倍距离（或在一半时间内移动同样距离）；或者，半个力能使一半重量在同样时间内移动同样距离。}

\orig From such statements, some of Aristotle's successors have made a determined effort to extract a general law. This law is customarily stated as:
\[
\mathrm{v} \propto \mathrm{F}/\mathrm{R}.
\]
That is, velocity (v) is proportional to the motive force (F) and inversely proportional to the resistance (R). For the special case of the natural descent of a heavy body, the motive force is the weight (W) of the body, and the relationship then becomes:
\[
\mathrm{v} \propto \mathrm{W}/\mathrm{R}.
\]

\zhline{从这些陈述中，亚里士多德的一些后继者努力抽取出一个普遍规律。这个规律通常表述为：}
\[
\mathrm{v} \propto \mathrm{F}/\mathrm{R}.
\]
也就是说，速度（v）与推动力（F）成正比，并与阻力（R）成反比。对于重物自然下落这一特殊情况，推动力就是物体的重量（W），于是关系变为：
\[
\mathrm{v} \propto \mathrm{W}/\mathrm{R}.
\]

\orig Such relationships probably do no great violence to Aristotle's intent for most cases of motion; however, giving them mathematical form, as we have done, suggests that they hold for all values of v, F (or W), and R---a claim that Aristotle would certainly have denied. He stated explicitly, for example, that a resistance equal to the motive force will prevent motion altogether, whereas the formula above offers no such result. Moreover, the appearance of velocity in these relationships seriously misrepresents Aristotle's conceptual framework, which contained no concept of velocity as a quantifiable measure of motion, but described motion only in terms of distances and times. Velocity as a technical scientific term to which numerical values might be assigned was a contribution of the Middle Ages (see below, chap. 12).

\zhline{对于大多数运动情形来说，这些关系大概并没有严重违背亚里士多德的意图；然而，像我们这样赋予它们数学形式，便暗示这些关系对 v、F（或 W）和 R 的所有数值都成立，而亚里士多德肯定会否认这一点。例如，他明确说过，若阻力等于推动力，运动就会完全被阻止；而上述公式并不给出这种结果。此外，速度出现在这些关系中，也严重误表了亚里士多德的概念框架；他的框架中并没有把速度作为运动的可量化尺度这一概念，而只是用距离和时间来描述运动。速度作为可以赋予数值的技术性科学术语，是中世纪的贡献（见下文第12章）。}

\orig Aristotle has been severely criticized for this theory of motion, on the assumption that any sensible person should have recognized its fatal flaws. Is such criticism justified? In the first place, our goal is to understand the behavior, beliefs, and achievements of historical actors against the background of the culture in which they lived, rather than to assess credit or blame according to the degree to which those historical actors resemble us. In short, historians must always contextualize their subjects. Second, some of the criticisms of Aristotle's theories of motion apply only to theories foisted onto Aristotle by followers and critics, rather than to his own. Third, the theory in its genuinely Aristotelian (and properly contextualized) version makes quite good sense today and would surely have made good sense in the fourth century B.C. For example, various surveys have shown that the majority of modern, university-educated people are prepared to assent to many of the basics of Aristotle's theory of motion. Fourth, the relatively modest level of quantitative content in Aristotle's theory is easily explained as the outcome of his larger philosophy of nature. His primary goal was to understand essential natures, not to explore quantitative relationships between such incidental factors as the space-time (or place-time) coordinates applicable to a moving body; even an exhaustive investigation of the latter gives us no useful information about the former. You may criticize Aristotle, if you like, for not being interested in whatever interests modern scientists, but we do not thereby learn anything significant about Aristotle.

\zhline{亚里士多德因其运动理论受到严厉批评，批评者假定任何明智的人都应当看出其中致命缺陷。这种批评合理吗？首先，我们的目标是在历史行动者所生活的文化背景中理解他们的行为、信念和成就，而不是按照他们在多大程度上像我们来评判功过。简言之，历史学家必须始终把研究对象放回语境中。第二，对亚里士多德运动理论的一些批评，只适用于后来追随者和批评者强加给亚里士多德的理论，而不适用于他自己的理论。第三，在真正亚里士多德式的（并且经过恰当语境化的）版本中，这一理论今天也相当说得通，在公元前四世纪当然更会显得合情合理。例如，多项调查显示，现代受过大学教育的人中，大多数愿意赞同亚里士多德运动理论的许多基本点。第四，亚里士多德理论中定量内容相对有限，很容易解释为他更大范围自然哲学的结果。他的主要目标是理解本质性的自然本性，而不是探索运动物体所适用的时空（或位置--时间）坐标等偶然因素之间的定量关系；即使把后者研究得再详尽，也不会给我们关于前者的有用信息。如果你愿意，可以批评亚里士多德不关心现代科学家所关心的东西，但我们并不能因此学到任何关于亚里士多德的重要内容。}

\orig Motion in the celestial sphere is an altogether different sort of phenomenon. The heavens, composed of the incorruptible quintessence, possess no contraries and are therefore incapable of qualitative change. It might seem fitting for such a region to be absolutely motionless, but this hypothesis is defeated by the most casual observation of the heavens. Aristotle therefore assigned to the heavens the most perfect of motions---continuous uniform circular motion. Besides being the most perfect of motions, uniform circular motion appears to have the capability of explaining the observed celestial cycles.

\zhline{天球中的运动则是完全不同的一类现象。天界由不可朽坏的第五元素构成，没有相反者，因此不能发生性质变化。这样一个区域似乎应当绝对静止，但对天空最随意的观察都会推翻这一假设。因此，亚里士多德把最完美的运动赋予天界，即连续的匀速圆周运动。匀速圆周运动除了是最完美的运动之外，似乎还能够解释观察到的天体周期。}

\orig By Aristotle's day, these cycles had been an object of study for centuries in the Greek world and for millennia in its predecessor civilizations. It was understood that the ``fixed'' stars move with perfect uniformity, as though fixed to a uniformly rotating sphere, with a period of rotation of approximately one day. But there were seven stars, the wandering stars or planets, that displayed a more intricate motion, apparently crawling around on the stellar sphere as it went through its daily rotation. These seven were the Sun, Moon, Mercury, Venus, Mars, Jupiter, and Saturn. The sun crawls slowly (about $1^\circ$/day), west to east with small variations in speed, through the sphere of fixed stars along a path called the ecliptic, which passes through the center of the zodiac. The moon follows approximately the same course, but at the more rapid rate of about $12^\circ$/day. The remaining planets also move along the ecliptic (or in its vicinity) with variable speed and with an occasional reversal of direction.

\zhline{到亚里士多德时代，这些周期在希腊世界已经被研究了数百年，在希腊之前的文明中则已被研究了数千年。人们理解到，“恒定的”星辰以完美的均匀性运动，仿佛固定在一个匀速旋转的天球上，旋转周期约为一日。但是，有七颗星，即游星或行星，表现出更复杂的运动：它们显然在恒星天球进行每日旋转时，于其上缓慢爬行。这七颗星是太阳、月亮、水星、金星、火星、木星和土星。太阳沿着一条称为黄道、穿过黄道带中心的路径，在恒星天球上由西向东缓慢移动（约每天 $1^\circ$），速度有小幅变化。月亮大致沿同一路径运行，但速度更快，约每天 $12^\circ$。其余行星也沿黄道（或其附近）运动，速度可变，并且偶尔逆行。}

\orig Are such complex motions compatible with the requirement of uniform circular motion in the heavens? Eudoxus, a generation before Aristotle, had already shown that they are. I will return to this subject in chap. 5; for the moment, it will be sufficient to point out that Eudoxus treated each complex planetary motion as a composite of a series of simple uniform circular movements. He did this by assigning to each planet a set of concentric spheres, and to each sphere one component of the complex planetary motion. Aristotle took over this scheme, with various modifications. When he was finished, he had produced an intricate piece of celestial machinery, consisting of fifty-five planetary spheres plus the sphere of the fixed stars.

\zhline{这种复杂运动与天界中匀速圆周运动的要求相容吗？比亚里士多德早一代的欧多克索斯已经表明二者相容。第5章将回到这个主题；眼下只需指出，欧多克索斯把每一种复杂的行星运动处理为一系列简单匀速圆周运动的复合。他的做法是为每颗行星分配一组同心球，并把复杂行星运动的一个分量分配给每个球。亚里士多德接过这一方案，并作了若干修改。完成时，他已经构造出一套精巧的天体机械，由五十五个行星天球加上恒星天球构成。}

\orig What is the cause of movement in the heavens? Aristotle's natural philosophy would not allow such a question to go unasked. The celestial spheres are composed, of course, of the quintessence; their motion, being eternal, must be natural rather than forced. The cause of this eternal motion must itself be unmoved, for if we do not postulate an unmoved mover, we quickly find ourselves trapped in an infinite regress: a moving mover must have acquired its motion from yet another moving mover, and so on. Aristotle identified the unmoved mover for the planetary spheres as the ``Prime Mover,'' a living deity representing the highest good, wholly actualized, totally absorbed in self-contemplation, nonspatial, separated from the spheres it (or he or she) moves, and not at all like the traditional anthropomorphic Greek gods. How, then, does the Prime Mover or Unmoved Mover cause motion in the heavens? Not as efficient cause, for that would require contact between the mover and the moved, but as final cause. That is, the Prime Mover is the object of desire for the celestial spheres, which endeavor to imitate its changeless perfection by assuming eternal, uniform circular motions. Any reader who has followed this much of Aristotle's discussion would be justified in assuming that there is a single Unmoved Mover for the entire cosmos; it comes as something of a surprise, therefore, when Aristotle announces that, in fact, each of the celestial spheres has its own Unmoved Mover, the object of its affection and the final cause of its motion.

\zhline{天上的运动原因是什么？亚里士多德的自然哲学不会允许这个问题不被提出。天球当然由第五元素构成；它们的运动既然是永恒的，就必定是自然运动而不是强迫运动。这种永恒运动的原因本身必须是不动的，因为如果不设定一个不动的推动者，我们很快就会陷入无穷倒退：一个运动着的推动者必定从另一个运动着的推动者那里获得其运动，如此无穷。亚里士多德把行星天球的不动推动者认定为“第一推动者”：这是一位活的神祇，代表最高善，完全现实化，完全沉浸于自我观照，没有空间位置，与它（或他、或她）所推动的天球相分离，并且完全不像传统拟人化的希腊诸神。那么，第一推动者或不动的推动者如何造成天上的运动？不是作为动力因，因为那要求推动者与被推动者之间有接触；而是作为目的因。也就是说，第一推动者是天球欲求的对象；天球试图通过采取永恒、匀速的圆周运动来模仿其不变的完美。任何读到亚里士多德这一部分讨论的读者，都有理由假定整个宇宙只有一个不动的推动者；因此，当亚里士多德宣布事实上每一个天球都有自己的不动推动者，即它所钟爱的对象和其运动的目的因时，这多少令人惊讶。}


% ===== 03_cohen_newton.tex =====
% Processed from chunks/03_cohen_newton.md.
% Bilingual LaTeX fragment. Compile with a Chinese-capable engine such as XeLaTeX/LuaLaTeX.
% Images are preserved with project-root-relative paths and require graphicx in the parent document.

\section*{Chapter 7: The Grand Design -- A New Physics\\第七章：宏伟构想：一种新物理学}
\begin{flushright}
I. Bernard Cohen
\end{flushright}

\noindent\textbf{原文} The publication of Isaac Newton's \emph{Principia} in 1687 was one of the most notable events in the whole history of physical science. In it one may find the culmination of thousands of years of striving to comprehend the system of the world, the principles of force and of motion, and the physics of bodies moving in different media. It is no small testimony to the vitality of Newton's scientific genius that although the physics of the \emph{Principia} has been altered, improved, and challenged ever since, we still set about solving most problems of celestial mechanics and the physics of gross bodies by proceeding essentially as Newton did some 300 years ago. Newtonian principles of celestial mechanics guide our artificial satellites, our space shuttles, and every spacecraft we launch to explore the vast reaches of our solar system. And if this is not enough to satisfy the canons of greatness, Newton was equally great as a pure mathematician. He invented the differential and integral calculus (produced simultaneously and independently by the German philosopher Gottfried Wilhelm Leibniz), which is the language of physics; he developed the binomial theorem and various properties of infinite series; and he laid the foundations for the calculus of variations.
\par\noindent\textbf{译文} 艾萨克·牛顿的《自然哲学的数学原理》（\emph{Principia}）于1687年出版，这是整个物理科学史上最重要的事件之一。在这部著作中，人们可以看到数千年来理解世界体系、力与运动的原理，以及不同介质中物体运动之物理学的努力达到了顶点。此后，《原理》中的物理学不断被修改、完善并受到挑战；然而直到今天，我们处理天体力学和宏观物体物理学的大多数问题时，实质上仍沿着牛顿约三百年前开辟的道路前进，这充分显示了牛顿科学天才的生命力。牛顿的天体力学原理指引着人造卫星、航天飞机，以及我们发射去探索太阳系广阔空间的每一艘航天器。若这些还不足以满足伟大的标准，那么牛顿作为纯数学家同样伟大。他发明了微分和积分学（德国哲学家戈特弗里德·威廉·莱布尼茨也同时且独立地创立了它），而微积分正是物理学的语言；他发展了二项式定理和无穷级数的多种性质；并为变分法奠定了基础。
\medskip

\noindent\textbf{原文} In optics, Newton began the experimental study of the analysis and composition of light, showing that white light is a mixture of light of many colors, each having a characteristic index of refraction. Upon these researches have risen the science of spectroscopy and the methods of color analysis. Newton invented a reflecting telescope and so showed astronomers how to transcend the limitations of telescopes built of lenses. All in all, his was a fantastic scientific achievement--of a kind that has never been equaled and may never be equaled again.
\par\noindent\textbf{译文} 在光学方面，牛顿开创了对光的分析与合成的实验研究，证明白光是许多种颜色光的混合，而每一种颜色都有其特定的折射率。在这些研究的基础上，光谱学和颜色分析方法发展起来。牛顿发明了反射望远镜，由此向天文学家展示了如何超越透镜望远镜的局限。总而言之，他的科学成就是惊人的，其性质前无古人，或许也后无来者。
\medskip

\noindent\textbf{原文} In this book we shall deal exclusively with Newton's system of dynamics and gravitation, the central problems for which the preceding chapters have been a preparation. If you have read them carefully, you have in mind all but one of the major ingredients requisite to an understanding of the Newtonian system of the world. But even if that one were to be given--the analysis of uniform circular motion--the guiding hand of Newton would still be required to put the ingredients together. It took genius to supply the new concept of universal gravitation. Let us see what Newton actually did.
\par\noindent\textbf{译文} 本书将专门讨论牛顿的动力学与引力体系，前面各章正是为这些核心问题所作的准备。如果你已仔细读过前文，那么理解牛顿世界体系所需的主要要素，除了一个之外，都已在你心中具备。可是即便再给出那一个要素，即匀速圆周运动的分析，仍然需要牛顿的引导之手将这些要素结合起来。提出万有引力这一新概念，需要天才。让我们看看牛顿实际做了什么。
\medskip

\noindent\textbf{原文} First of all, it must be understood that Galileo himself never attempted to display any scheme of forces that would account for the movement of the planets, or of their satellites. As for Copernicus, the \emph{De revolutionibus} contains no important insight into a celestial mechanics. Kepler had tried to supply a celestial mechanism, but the result was never a very happy one. He held that the \emph{anima motrix} emanating from the sun would cause planets to revolve about the sun in circles. He further supposed that magnetic interactions of the sun and a planet would shift the planet during an otherwise circular revolution into an elliptical orbit. Others who contemplated the problems of planetary motion proposed systems of mechanics containing certain features that were later to appear in Newtonian dynamics. One of these was Robert Hooke, who quite understandably thought that Newton should have given him more credit than a mere passing reference for having anticipated parts of the laws of dynamics and gravitation.
\par\noindent\textbf{译文} 首先必须明白，伽利略本人从未试图提出一套力的方案，以解释行星或其卫星的运动。至于哥白尼，《天体运行论》中并没有关于天体力学的重要洞见。开普勒曾试图提供一种天体机制，但结果并不十分成功。他认为从太阳发出的\emph{anima motrix}（运动灵魂）会使行星围绕太阳作圆周运动。他还进一步设想，太阳与行星之间的磁相互作用，会把本来作圆周运行的行星推移到椭圆轨道上。其他思考行星运动问题的人也提出过一些力学体系，其中含有后来会出现在牛顿动力学中的若干特征。罗伯特·胡克就是其中之一；他认为自己预见了动力学与引力定律的一部分，牛顿本应给予他比一句顺带提及更多的承认，这种想法是可以理解的。
\medskip

\subsection*{Newtonian Anticipations\\牛顿思想的先声}

\noindent\textbf{原文} The climactic chapter in the discovery of the mechanics of the universe starts with a pretty story. By the third quarter of the seventeenth century, a group of men had become so eager to advance the new mathematical experimental sciences that they banded together to perform experiments in concert, to present problems for solution to one another, and to report on their own researches and on those of others as revealed by correspondence, books, and pamphlets. Thus it came about that Robert Hooke, Edmond Halley, and Sir Christopher Wren, England's foremost architect, met to discuss the question, Under what law of force would a planet follow an elliptical orbit? From Kepler's laws--especially the third or harmonic law, but also the second or law of areas--it was clear that the sun somehow or other must control or at least affect the motion of a planet in accordance with the relative proximity of the planet to the sun. Even if the particular mechanisms proposed by Kepler (an \emph{anima motrix} and a magnetic force) had to be rejected, there could be no doubt that some kind of planet-sun interaction keeps the planets in their courses. Furthermore, a more acute intuition than Kepler's would sense that any force emanating from the sun must spread out in all directions from that body, presumably diminishing according to the inverse of the square of its distance from the sun--as the intensity of light diminishes in relation to distance. But to say this much is a very different thing from proving it mathematically. For to prove it would require a complete physics with mathematical methods for solving all the attendant and consequent problems. When Newton declined to credit authors who tossed off general statements without being able to prove them mathematically or fit them into a valid framework of dynamics, he was quite justified in saying, as he did of Hooke's claims: ``Now is not this very fine? Mathematicians that find out, settle, and do all the business must content themselves with being nothing but dry calculators and drudges; and another, that does nothing but pretend and grasp at all things, must carry away all the invention, as well of those that were to follow him as of those that went before.'' (See, further, Supplement 11.)
\par\noindent\textbf{译文} 宇宙力学发现史中最高潮的一章，是从一个颇动人的故事开始的。到17世纪第三个四分之一时期，一批人已急切地想推进新的数学实验科学，他们结成团体，共同做实验，彼此提出待解的问题，并通过通信、书籍和小册子报告自己的研究以及他人的研究。于是，罗伯特·胡克、埃德蒙·哈雷和英国首屈一指的建筑师克里斯托弗·雷恩爵士聚在一起，讨论这样一个问题：在什么样的力的定律下，行星会沿椭圆轨道运行？根据开普勒定律，尤其是第三定律即调和定律，同时也包括第二定律即面积定律，显然太阳必定以某种方式按照行星与太阳的相对距离来控制，或至少影响行星的运动。即使必须抛弃开普勒所提出的具体机制（运动灵魂和磁力），也毫无疑问，某种行星--太阳相互作用使行星保持在各自的轨道上。此外，比开普勒更敏锐的直觉会意识到，从太阳发出的任何力都必定从太阳向四面八方扩散，并且大概按距太阳距离平方的倒数减弱，正如光强随距离而减弱一样。但是，说到这一点与用数学证明这一点，是完全不同的两回事。要证明它，就需要一套完整的物理学，以及解决所有相关和后续问题的数学方法。因此，当有些作者只能随口提出一般性断言，却既不能作数学证明，也不能把这些断言纳入有效的动力学框架时，牛顿拒绝给予他们荣誉是有充分理由的。对于胡克的主张，牛顿曾说：“这岂不是妙极了吗？数学家们发现、确定并做完全部工作，却只能满足于不过是枯燥的计算者和苦工；而另一个人什么也不做，只是自称什么都懂、什么都要抓到手，却必须把所有发明都归于自己，既包括后来者的，也包括前人已有的。”（另见补编11。）
\medskip

\noindent\textbf{原文} In any event, by January 1684 Halley had concluded that the force acting on planets to keep them in their orbits ``decreased in the proportion of the squares of the distances reciprocally,''
\[
F \propto \frac{1}{D^{2}},
\]
but he was not able to deduce from that hypothesis the observed motions of the celestial bodies. When Wren and Hooke met later in the month, they agreed with Halley's supposition of a solar force. Hooke boasted ``that upon that principle all the laws of the celestial motions were to be [i.e., could be] demonstrated, and that he himself had done it.'' But despite repeated urgings and Wren's offer of a considerable monetary prize, Hooke did not--and presumably could not--produce a solution. Six months later, in August 1684, Halley decided to go to Cambridge to consult Isaac Newton. On his arrival he learned the ``good news'' that Newton ``had brought this demonstration to perfection.'' Here is De Moivre's almost contemporaneous account of that visit:
\par\noindent\textbf{译文} 无论如何，到1684年1月，哈雷已经得出结论：作用在行星上、使其保持在轨道中的力，“按距离平方的倒数比例减小”，
\[
F \propto \frac{1}{D^{2}},
\]
但他无法从这一假设推出天体的观测运动。那个月稍晚些时候，雷恩和胡克会面时，他们同意哈雷关于太阳力的设想。胡克夸口说，“依照那个原理，天体运动的一切定律都将[也就是能够]得到证明，而且他本人已经做到了”。可是，尽管众人一再催促，雷恩还提出一笔可观的奖金，胡克并没有拿出解答，而且大概也拿不出来。六个月后，即1684年8月，哈雷决定去剑桥请教艾萨克·牛顿。他到达后得知了一个“好消息”：牛顿“已经把这个证明完成得尽善尽美”。以下是德·莫阿弗对那次访问所作的几乎同时代的叙述：
\medskip

\begin{quote}
\noindent\textbf{原文} After they had been some time together, the Dr. [Halley] asked him what he thought the curve would be that would be described by the planets supposing the force of attraction towards the sun to be reciprocal to the square of their distance from it. Sir Isaac replied immediately that it would be an ellipsis. The Doctor, struck with joy and amazement, asked him how he knew it. Why, saith he, I have calculated it. Whereupon Dr. Halley asked him for his calculation without any further delay. Sir Isaac looked among his papers but could not find it, but he promised him to renew it and then to send it him. Sir Isaac, in order to make good his promise, fell to work again, but he could not come to that conclusion which he thought he had before examined with care. However, he attempted a new way which, though longer than the first, brought him again to his former conclusion. Then he examined carefully what might be the reason why the calculation he had undertaken before did not prove right, and he found that, having drawn an ellipsis coarsely with his own hand, he had drawn the two axes of the curve, instead of drawing two diameters somewhat inclined to one another, whereby he might have fixed his imagination to any two conjugate diameters, which was requisite he should do. That being perceived, he made both his calculations agree together.
\par\noindent\textbf{译文} 两人相处了一段时间之后，博士[哈雷]问他：假定指向太阳的吸引力与行星距太阳距离的平方成反比，那么行星将描画出什么曲线？艾萨克爵士立刻回答说，会是一个椭圆。博士又惊又喜，问他是怎么知道的。他说，哦，我算过。于是哈雷博士不再迟疑，请他拿出计算。艾萨克爵士在自己的论文中寻找，却找不到；但他答应重新作一次，再寄给哈雷。为了履行诺言，艾萨克爵士重新动手，可是他无法得到自己以为先前已经仔细考察过的结论。不过，他尝试了另一种方法；这种方法虽比第一种更长，却又把他带回到原先的结论。随后他仔细检查，想找出先前所作计算为何不正确；他发现，自己用手粗略画椭圆时，画的是曲线的两条轴，而不是两条彼此略有倾斜的直径。要使想象固定在任意两条共轭直径上，后者才是必须画的。意识到这一点后，他使两种计算彼此一致了。
\end{quote}

\noindent\textbf{原文} Spurred on by Halley's visit, Newton resumed work on a subject that had commanded his attention in his twenties when he had laid the foundations of his other great scientific discoveries: the nature of white light and color and the differential and integral calculus. He now put his investigations in order, made great progress, and in the fall term of the year, discussed his research in a series of lectures on dynamics that he gave at Cambridge University, as required by his professorship. Eventually, with Halley's encouragement, a draft of these lectures, \emph{De motu corporum}, grew into one of the greatest and most influential books any man has yet conceived. Many a scientist has echoed the sentiment that Halley expressed in the ode he wrote as a preface to Newton's \emph{Principia} (or, to give Newton's masterpiece its full title, \emph{Philosophiae naturalis principia mathematica}, \emph{Mathematical Principles of Natural Philosophy}, London, 1687):
\par\noindent\textbf{译文} 在哈雷来访的推动下，牛顿重新研究起一个在他二十多岁时就曾吸引其注意力的课题；那时他已经为另外两项伟大的科学发现奠定了基础，即白光与颜色的本性，以及微分和积分学。现在他整理自己的研究，取得重大进展，并在当年秋季学期按教授职务的要求，在剑桥大学作了一系列关于动力学的讲座，讨论这些研究。最终，在哈雷的鼓励下，这些讲稿的草稿《物体运动论》（\emph{De motu corporum}）发展成为人类迄今构思出的最伟大、最有影响力的著作之一。许多科学家都呼应过哈雷在为牛顿《原理》所写序诗中表达的感情（若写出牛顿这部杰作的全名，则是\emph{Philosophiae naturalis principia mathematica}，《自然哲学的数学原理》，伦敦，1687年）：
\medskip

\begin{verse}
Then ye who now on heavenly nectar fare,\\
Come celebrate with me in song the name\\
Of Newton, to the Muses dear; for he\\
Unlocked the hidden treasuries of Truth:\\
So richly through his mind had Phoebus cast\\
The radiance of his own divinity.\\
Nearer the gods no mortal may approach.
\end{verse}

\begin{verse}
那么，你们这些如今享用天上琼浆的人，\\
请与我同歌牛顿之名，\\
他为缪斯所珍爱；因为他\\
开启了真理深藏的宝库：\\
福玻斯曾如此丰沛地\\
把自身神性的光辉投射进他的心灵。\\
凡人再也不能更接近诸神。
\end{verse}

\subsection*{The Principia\\《原理》}

\noindent\textbf{原文} The \emph{Principia} is divided into three parts or books; we shall concentrate on the first and third. In Book One Newton develops the general principles of the dynamics of moving bodies, and in Book Three he applies the principles to the mechanism of the universe. Book Two deals with a facet of fluid mechanics, the theory of waves, and other aspects of physics.
\par\noindent\textbf{译文} 《原理》分为三部分，或三卷；我们将集中讨论第一卷和第三卷。在第一卷中，牛顿发展了运动物体的动力学一般原理；在第三卷中，他把这些原理应用于宇宙的机制。第二卷处理流体力学的一个方面、波的理论以及物理学的其他方面。
\medskip

\noindent\textbf{原文} In Book One, following the preface, a set of definitions, and a discussion of the nature of time and space, Newton presented the ``axioms, or laws of motion.''
\par\noindent\textbf{译文} 在第一卷中，牛顿在序言、一组定义以及关于时间与空间性质的讨论之后，提出了“公理，或运动定律”。
\medskip

\subsubsection*{Law I\\定律一}
\noindent\textbf{原文} Every body perseveres in its state of being at rest or of moving uniformly straight forward, except insofar as it is compelled to change its state by forces impressed upon it.
\par\noindent\textbf{译文} 每一个物体都会保持其静止状态，或沿直线作匀速运动的状态，除非它被施加在其上的力迫使改变这种状态。
\medskip

\subsubsection*{Law II\\定律二}
\noindent\textbf{原文} A change in motion is proportional to the motive force impressed and takes place in the direction of the straight line along which that force is impressed. [See Supplementary Note on p. 184.]
\par\noindent\textbf{译文} 运动的改变与所施加的动力成正比，并沿着施加该力的直线方向发生。[见第184页补注。]
\medskip

\begin{center}
\safeincludeimage{0f043290c0aee6258d3b9dc29864191fc41e15048c067bc65b960e2596f4c37d.jpg}
\end{center}

\noindent\textbf{Plate VI.} A ball fired with a spring gun in the smokestack of a moving toy locomotive describes a parabola and lands on the locomotive instead of going straight up and down as it does when the locomotive is standing still. These stroboscopic pictures, with exposures at intervals of one-thirtieth of a second, vividly illustrate one of Galileo's arguments on the behavior of falling bodies and settle the ancient debate about bodies dropped from the masts of moving ships. If the speed of the locomotive were absolutely uniform and if the ball met no air resistance, it would land in the smokestack. (In fact, even under the imperfect conditions of the experiment, the ball hits the smokestack more often than not.) Note that the ball attains the same height whether the locomotive is at rest or moving. Notice, too, that in the picture where the locomotive is standing still, the distances traveled by the ball in the intervals of exposure correspond almost exactly. On the ascent, gravity slows it down; on the descent, gravity speeds it up. Photographs by Berenice Abbott.
\par\noindent\textbf{图版VI。} 在一辆运动中的玩具机车烟囱里，用弹簧枪射出一个小球；小球描出一条抛物线，并落回机车上，而不是像机车静止时那样笔直上升又下降。这些频闪照片以三十分之一秒为曝光间隔，生动说明了伽利略关于落体行为的一个论证，也解决了关于从运动船只桅杆上落下物体的古老争论。如果机车速度绝对均匀，并且小球不受空气阻力，小球就会落回烟囱中。（事实上，即使在实验条件并不完美的情况下，小球击中烟囱的次数也多于未击中的次数。）注意，无论机车静止还是运动，小球达到的高度都相同。还应注意，在机车静止的照片中，小球在各曝光间隔内通过的距离几乎正好对应。上升时，重力使它减速；下降时，重力使它加速。照片由贝伦妮丝·阿博特拍摄。
\medskip

\noindent\textbf{原文} Observe that if a body is in uniform motion in a straight line, a force at right angles to the direction of motion of the body will not affect the forward motion. This follows from the fact that the acceleration is always in the same direction as the force producing it, so that the acceleration in this case is at right angles to the direction of motion. Thus in the toy train experiment of Chapter 5, the chief force acting is the downward force of gravity, producing a vertical acceleration. The ball, whether moving forward or at rest, is thus caused to slow down in its upward motion until it comes to rest, and then be speeded up or accelerated on the way down.
\par\noindent\textbf{译文} 请注意，如果一个物体正沿直线作匀速运动，那么与该物体运动方向成直角的力不会影响其前进运动。这是因为加速度总是与产生它的力同方向，因此在这种情况下，加速度与运动方向成直角。所以，在第五章的玩具火车实验中，起主要作用的力是向下的重力，它产生竖直加速度。于是，无论小球本来是否向前运动，它在向上运动时都会减速，直到停止；随后在下降途中又会加速。
\medskip

\noindent\textbf{原文} A comparison of the two sets of photographs (p. 83) shows that the upward and downward motions are exactly the same whether the train is at rest or in uniform motion. In the forward direction there is no effect of weight or gravity, since this acts only in a downward direction. The only force in the forward or horizontal direction is the small amount of air friction, which is almost negligible; so one may say that in the horizontal direction there is no force acting. According to Newton's first law of motion, the ball will continue to move in the forward direction with uniform motion in a straight line just as the train does--a fact you can check by inspecting the photograph. The ball remains above the locomotive whether the train is at rest or in uniform motion in a straight line. This law of motion is sometimes called the principle of inertia, and the property that material bodies have of continuing in a state of rest or of uniform motion in a straight line is sometimes known as the bodies' inertia.\textsuperscript{1}
\par\noindent\textbf{译文} 比较两组照片（第83页）可见，无论火车静止还是作匀速运动，小球的上升和下降运动都完全相同。在前进方向上，重量或重力不起作用，因为重力只沿向下方向作用。前进方向或水平方向上唯一的力，是很小的一点空气摩擦，几乎可以忽略；因此可以说，水平方向上没有力在作用。根据牛顿第一运动定律，小球将像火车一样继续沿前进方向作直线匀速运动；这一点可以通过观察照片来核对。无论火车静止还是沿直线匀速运动，小球都保持在机车上方。这条运动定律有时称为惯性原理，而物质物体继续保持静止状态或直线匀速运动状态的性质，有时称为物体的惯性。\textsuperscript{1}
\medskip

\noindent\textbf{原文} Newton illustrated Law I by reference to projectiles that continue in their forward motions ``so far as they are not retarded by the resistance of the air, or impelled downward by the force of gravity,'' and he referred also to ``the greater bodies of planets and comets.'' (On the inertial aspect of the motion of ``greater bodies'' such as ``planets and comets,'' see Supplement 12.) At this one stroke Newton postulated the opposite view of Aristotelian physics. In the latter, no celestial body could move uniformly in a straight line in the absence of a force, because this would be a ``violent'' motion and so contrary to its nature. Nor could a terrestrial object, as we have seen, move along its ``natural'' straight line without an external mover or an internal motive force. Newton, presenting a physics that applies simultaneously to both terrestrial and celestial objects, stated that in the absence of a force bodies do not necessarily stand still or come to rest as Aristotle supposed, but they may move at constant rectilinear speed. This ``indifference'' of all sorts of bodies to rest or uniform straight-line motion in the absence of a force clearly is an advanced form of Galileo's statement in his book on sunspots (p. 88), the difference being that in that work Galileo was writing about uniform motion along a great spherical surface concentric with the earth.
\par\noindent\textbf{译文} 牛顿用抛射体来说明定律一：只要抛射体“未被空气阻力所阻滞，或未被重力向下推动”，它们就会继续向前运动；他还提到“行星和彗星这些较大的物体”。（关于“行星和彗星”等“较大物体”运动中的惯性方面，见补编12。）仅此一举，牛顿就提出了与亚里士多德物理学相反的观点。在亚里士多德物理学中，没有力时任何天体都不可能沿直线匀速运动，因为这会是“强迫”运动，因而违背其本性。同样，正如我们已经看到的，地上物体如果没有外在推动者或内在动力，也不能沿其“自然”直线运动。牛顿提出了一种同时适用于地上物体和天上物体的物理学，指出在没有力的情况下，物体并不一定像亚里士多德所设想的那样静止不动或趋于静止，而是可以以恒定速度作直线运动。在没有力时，各类物体对于静止和直线匀速运动表现出的这种“无差别”，显然是伽利略在其太阳黑子著作（第88页）中说法的进一步形式；不同的是，伽利略在那里谈的是沿着与地球同心的巨大球面作匀速运动。
\medskip

\noindent\textbf{原文} Newton said of the laws of motion that they were ``such principles as have been received by mathematicians, and \ldots{} confirmed by [an] abundance of experiments. By the first two Laws and the first two Corollaries, Galileo discovered that the descent of bodies varies as the square of the time and that the motion of projectiles is in the curve of a parabola, experience agreeing with both, unless so far as these motions are a little retarded by the resistance of the air.'' The ``two Corollaries'' deal with methods used by Galileo and many of his predecessors to combine two different forces or two independent motions. Fifty years after the publication of Galileo's \emph{Two New Sciences} it was difficult for Newton, who had already established an inertial physics, to conceive that Galileo could have come as close as he had to the concept of inertia without having taken full leave of circularity and having known the true principle of linear inertia.
\par\noindent\textbf{译文} 牛顿论及运动定律时说，它们是“数学家们已经接受，并且……由大量实验所证实的原理。凭借前两条定律和前两个推论，伽利略发现物体下落与时间的平方成比例，抛射体的运动沿抛物线进行；经验与这两点一致，只是这些运动会因空气阻力而略有迟滞”。所谓“两个推论”，涉及伽利略及其许多前人用来合成两个不同的力或两个独立运动的方法。在伽利略《两门新科学》出版五十年后，对于已经建立惯性物理学的牛顿来说，很难想象伽利略竟能如此接近惯性概念，却没有完全摆脱圆形运动，也没有认识真正的线性惯性原理。
\medskip

\noindent\textbf{原文} Newton was being very generous to Galileo because, however it may be argued that Galileo ``really did'' have the law of inertia or Newton's Law I, a great stretch of the imagination is required to assign any credit to Galileo for Law II. This law has two parts. In the second half of Newton's statement of Law II, the ``change in motion'' produced by an ``impressed'' or ``motive'' force--whether that is a change in the speed with which a body moves or a change in the direction in which it is moving--is said to be ``in the direction of the straight line along which that force is impressed.'' This much is certainly implied in Galileo's analysis of projectile motion because Galileo assumed that in the forward direction there is no acceleration because there is no horizontal force, except the negligible action of air friction; but in the vertical direction there is an acceleration or continual increase of downward speed, because of the downward-acting weight force. But the first part of Law II--that the change in the magnitude of the motion is related to the motive force--is something else again; only a Newton could have seen it in Galileo's studies of falling bodies. This part of the law says that if an object were to be acted on first by one force $F_{1}$ and then by some other force $F_{2}$, the accelerations or changes in speed produced, $A_{1}$ and $A_{2}$, would be proportional to the forces, or that
\[
\frac{F_{1}}{F_{2}} = \frac{A_{1}}{A_{2}}, \qquad \mathrm{or} \qquad
\frac{F_{1}}{A_{1}} = \frac{F_{2}}{A_{2}}.
\]
\par\noindent\textbf{译文} 牛顿对伽利略已经十分慷慨，因为无论人们如何主张伽利略“确实”已经具有惯性定律或牛顿定律一的思想，要把定律二的功劳也归给伽利略，都需要极大的想象力。这条定律有两个部分。在牛顿对定律二的表述后半部分中，由“施加的”力或“动力”所产生的“运动的改变”——无论是物体运动速度的改变，还是运动方向的改变——都被说成“沿着施加该力的直线方向”发生。这一点当然隐含在伽利略对抛射体运动的分析中，因为伽利略假定，在前进方向上除几乎可忽略的空气摩擦外没有水平力，所以没有加速度；而在竖直方向上，由于有向下作用的重量力，存在加速度，或者说向下速度不断增加。但是定律二的前半部分，即运动量大小的改变与动力有关，却是另一回事；只有牛顿才能在伽利略的落体研究中看出这一点。定律的这一部分说，如果一个物体先受一个力$F_{1}$作用，然后又受另一个力$F_{2}$作用，那么所产生的加速度或速度改变$A_{1}$和$A_{2}$将与这些力成正比，或者说
\[
\frac{F_{1}}{F_{2}} = \frac{A_{1}}{A_{2}}, \qquad \hbox{或} \qquad
\frac{F_{1}}{A_{1}} = \frac{F_{2}}{A_{2}}.
\]
\medskip

\noindent\textbf{原文} But in analyzing falling, Galileo was dealing with a situation in which only one force acted on each body, its weight $W$, and the acceleration it produced was $g$, the acceleration of a freely falling body. (For the two forms of Newton's Law II, see p. 184.)
\par\noindent\textbf{译文} 可是在分析下落时，伽利略处理的是每个物体只受一个力作用的情形，即它的重量$W$；这个力所产生的加速度是$g$，即自由落体的加速度。（关于牛顿定律二的两种形式，见第184页。）
\medskip

\noindent\textbf{原文} Where Aristotle had said that a given force gives an object a certain characteristic speed, Newton now said that a given force always produces in that body a definite acceleration $A$. To find the speed $V$, we must know how long a time $T$ the force has acted, or how long the object has been accelerated, so that Galileo's law
\[
V = AT
\]
may be applied.
\par\noindent\textbf{译文} 亚里士多德曾说，给定的力会给予物体某种特定速度；而牛顿现在说，给定的力总是在该物体中产生确定的加速度$A$。要找到速度$V$，我们必须知道力作用了多长时间$T$，也就是物体被加速了多长时间，这样才可以应用伽利略定律
\[
V = AT.
\]
\medskip

\noindent\textbf{原文} At this point let us try a thought experiment, in which we assume we have two cubes of aluminum, one just twice the volume of the other. (Incidentally, to ``duplicate'' a cube--or make a cube having exactly twice the volume as some given cube--is as impossible within the framework of Euclidean geometry as to trisect an angle or to square a circle.) We now subject the smaller cube to a series of forces $F_{1}, F_{2}, F_{3}, \ldots$ and determine the corresponding accelerations $A_{1}, A_{2}, A_{3}, \ldots$. In accordance with Law II, we would find that there is a certain constant value of the ratio of force to acceleration
\[
\frac{F_{1}}{A_{1}} = \frac{F_{2}}{A_{2}} = \frac{F_{3}}{A_{3}} = \dots = m_{s},
\]
which for this object we may call $m_s$.
\par\noindent\textbf{译文} 此时让我们做一个思想实验。假定我们有两个铝立方体，其中一个的体积正好是另一个的两倍。（顺便说一句，在欧几里得几何框架内，“倍立方”——也就是作出一个体积正好为某给定立方体两倍的立方体——与三等分任意角或化圆为方一样不可能。）现在我们让较小的立方体依次受到一系列力$F_{1}, F_{2}, F_{3}, \ldots$，并测定相应的加速度$A_{1}, A_{2}, A_{3}, \ldots$。根据定律二，我们会发现力与加速度之比有某个恒定值：
\[
\frac{F_{1}}{A_{1}} = \frac{F_{2}}{A_{2}} = \frac{F_{3}}{A_{3}} = \dots = m_{s},
\]
对于这个物体，我们可以把它称为$m_s$。
\medskip

\noindent\textbf{原文} We now repeat the operations with the larger cube and find that the same set of forces $F_1, F_2, F_3, \ldots$ respectively produces another set of accelerations $a_1, a_2, a_3, \ldots$. In accordance with Newton's second law, the force-acceleration ratio is again a constant which for this object we may call $m_l$:
\[
\frac{F_{1}}{a_{1}} = \frac{F_{2}}{a_{2}} = \frac{F_{3}}{a_{3}} = \dots = m_{l}.
\]
\par\noindent\textbf{译文} 现在我们对较大的立方体重复这些操作，发现同一组力$F_1, F_2, F_3, \ldots$分别产生另一组加速度$a_1, a_2, a_3, \ldots$。根据牛顿第二定律，力与加速度之比再次是一个常数；对于这个物体，我们可以称之为$m_l$：
\[
\frac{F_{1}}{a_{1}} = \frac{F_{2}}{a_{2}} = \frac{F_{3}}{a_{3}} = \dots = m_{l}.
\]
\medskip

\noindent\textbf{原文} For the larger object the constant proves to be just twice as large as the constant obtained for the smaller one and, in general, so long as we deal with a single variety of matter like pure aluminum, this constant is proportional to the volume and so is a measure of the amount of aluminum in any sample. This particular constant is a measure of an object's resistance to acceleration, or a measure of the tendency of that object to stay as it is--either at rest, or in motion in a straight line. For observe that $m_l$ was twice $m_s$; to give both objects the same acceleration or change in motion the force required for the larger object is just twice what it must be for the smaller. The tendency of any object to continue in its state of motion (at constant speed in a straight line) or its state of rest is called its inertia; hence, Newton's Law I is also called the principle of inertia. The constant determined by finding the constant force-acceleration ratio for any given body may thus be called the body's inertia. But for our aluminum blocks this same constant is also a measure of the ``quantity of matter'' in the object, which is called its mass. We now make precise the condition that two objects of different material--say one of brass and the other of wood--shall have the same ``quantity of matter'': it is that they have the same mass as determined by the force-acceleration ratio, or the same inertia.
\par\noindent\textbf{译文} 对较大的物体而言，这个常数恰好是较小物体所得常数的两倍；一般地说，只要我们处理的是像纯铝这样单一类型的物质，这个常数就与体积成正比，因此也是任何样品中铝的多少的量度。这个特定常数是物体抵抗加速的量度，或者说是物体保持原状的倾向的量度——无论原状是静止，还是沿直线运动。请注意，$m_l$是$m_s$的两倍；若要使两个物体获得相同的加速度或运动改变，较大物体所需的力正好是较小物体的两倍。任何物体继续保持其运动状态（以恒定速度沿直线运动）或静止状态的倾向称为它的惯性；因此，牛顿定律一也称为惯性原理。对于任一给定物体，通过求出恒定的力--加速度比所确定的常数，也可称为该物体的惯性。但对我们的铝块来说，这同一个常数也是物体中“物质多少”的量度，而这种量度称为它的质量。现在我们可以精确说明，不同材料的两个物体——比如一个黄铜物体和一个木头物体——具有相同“物质多少”的条件是什么：它们必须具有由力--加速度比确定的相同质量，或者说具有相同惯性。
\medskip

\noindent\textbf{原文} In ordinary life, we do not compare the ``quantity of matter'' in objects in terms of their inertias, but in terms of their weight. Newtonian physics makes it clear why we can, and through its clarification we are able to understand why at any place on the earth two unequal weights in a vacuum fall at the same rate. But we may observe that in at least one common situation we always compare the inertias of objects rather than their weights. This happens when a person hefts two objects to find which is heavier, or has the greater mass. He does not hold them out to see which pulls down more on his arm; instead, he moves them up and down to find which is easier to move. In this way he determines which has the greater resistance to a change in its state of motion in a straight line or of rest--that is, which has the greater inertia. (On Newton's concept of inertia, see Supplement 15.)
\par\noindent\textbf{译文} 在日常生活中，我们并不根据物体的惯性来比较其中的“物质多少”，而是根据其重量来比较。牛顿物理学说明了我们为什么可以这样做；通过这种说明，我们也能够理解为什么在地球上任一地点，真空中两个重量不同的物体会以相同速率下落。不过我们可以注意到，至少在一种常见情形中，我们总是在比较物体的惯性而不是重量。这发生在一个人掂量两个物体，以判断哪一个更重或质量更大时。他并不是把它们伸出去，看哪一个把手臂向下拉得更多；相反，他上下移动它们，看哪一个更容易移动。通过这种方式，他判断哪一个对其直线运动状态或静止状态的改变有更大的抵抗，也就是哪一个有更大的惯性。（关于牛顿的惯性概念，见补编15。）
\medskip

\subsection*{Final Formulation of the Law of Inertia\\惯性定律的最终表述}

\noindent\textbf{原文} At one point in his \emph{Discourses and Demonstrations Concerning Two New Sciences}, Galileo imagined a ball to be rolling along a plane and noted that ``equable motion on this plane would be perpetual if the plane were of infinite extent.'' A plane without limit is all right for a pure mathematician, who is a Platonist in any case. But Galileo was a man who combined just such a Platonism with a concern for applications to the real world of sensory experience. In the \emph{Two New Sciences}, Galileo was not interested only in abstractions as such, but in the analysis of real motions on or near the earth. So we understand that having talked about a plane without limit, he does not continue with such a fancy, but asks what would happen on such a plane if it were a real earthly plane, which for him means that it is ``ended, and [situated] on high.'' The ball, in the real world of physics, falls off the plane and begins to fall to the ground. In this case,
\par\noindent\textbf{译文} 在《两门新科学的对话与数学证明》的某处，伽利略想象一个球沿平面滚动，并指出：“如果这个平面无限延伸，其上的匀速运动就会永远持续。”一个没有边界的平面对于纯数学家来说当然可以接受，因为纯数学家无论如何都是柏拉图主义者。但伽利略把这种柏拉图主义与对感官经验中现实世界之应用的关切结合在一起。在《两门新科学》中，伽利略并不只关心抽象本身，而是关心地球上或地球附近真实运动的分析。因此我们理解，在谈到一个无边界平面之后，他并没有继续这种想象，而是问：如果这样的平面是一个真实的地上平面，会发生什么？对他来说，这意味着该平面是“有尽头的，并且[位于]高处”。在物理学的现实世界中，球会滚出平面，开始向地面下落。在这种情况下，
\medskip

\begin{quote}
\noindent\textbf{原文} the movable (which I conceive of as being endowed with heaviness), driven to the end of this plane and going on further, adds on to its previous equable and indelible motion that downward tendency which it has from its own heaviness. Thus there emerges a certain motion, compounded from equable horizontal and from naturally accelerated downward [motion], which I call ``projection.''
\par\noindent\textbf{译文} 可动体（我设想它具有重性）被推到这个平面的尽头并继续前进时，会在其先前均匀且不可磨灭的运动之上，增添由自身重性产生的向下倾向。于是便出现一种运动，它由均匀的水平运动和自然加速的向下[运动]复合而成，我称之为“抛射”。
\end{quote}

\noindent\textbf{原文} Unlike Galileo, Newton made a clear separation between the world of abstract mathematics and the world of physics, which he still called philosophy. Thus the \emph{Principia} included both ``mathematical principles'' as such and those that could be applied in ``natural philosophy,'' but Galileo's \emph{Two New Sciences} included only those mathematical conditions exemplified in nature. For instance, Newton plainly knew that the attractive force exerted by the sun on a planet varies as the inverse square of the distance,
\[
F \propto \frac{1}{D^{2}},
\]
but in Book One of the \emph{Principia} he explored the consequences not only of this particular force but of others with quite different dependence on the distance, including
\[
F \propto D.
\]
\par\noindent\textbf{译文} 与伽利略不同，牛顿在抽象数学世界和物理世界之间作出了清楚区分；他仍把物理世界称为哲学。因此，《原理》既包括作为数学本身的“数学原理”，也包括可以应用于“自然哲学”的原理；而伽利略的《两门新科学》只包括自然中所体现的那些数学条件。例如，牛顿清楚知道太阳对行星施加的吸引力随距离平方的倒数而变化，
\[
F \propto \frac{1}{D^{2}},
\]
但在《原理》第一卷中，他不仅考察了这种特定力的结果，也考察了其他与距离有截然不同依赖关系的力的结果，包括
\[
F \propto D.
\]
\medskip

\subsection*{``The System of the World''\\“世界体系”}

\noindent\textbf{原文} At the beginning of Book Three, which was devoted to ``The System of the World,'' Newton explained how it differed from the preceding two, which had been dealing with ``The Motion of Bodies'':
\par\noindent\textbf{译文} 在第三卷开头，牛顿把该卷用于讨论“世界体系”，并说明它与前两卷的不同；前两卷讨论的是“物体运动”：
\medskip

\begin{quote}
\noindent\textbf{原文} In the preceding Books I have laid down \ldots{} principles not philosophical [pertaining to physics] but mathematical: such, namely, as we may build our reasonings upon in philosophical inquiries. These principles are laws and conditions of certain motions, and powers or forces, which chiefly have respect to philosophy; but, lest they should have appeared of themselves dry and barren, I have illustrated them here and there with some philosophical scholiums, giving an account of such things as are of a more general nature, and which philosophy seems chiefly to be founded on: such as the density and the resistance of bodies, spaces void of all bodies, and the motion of light and sounds. It remains that, from the same principles, I now demonstrate the structure of the System of the World.
\par\noindent\textbf{译文} 在前几卷中，我提出了……并非哲学的[即属于物理学的]、而是数学的原理，也就是我们在哲学探究中可以据以推理的原理。这些原理是某些运动的定律和条件，以及能力或力；它们主要关系到哲学。不过，为免它们本身显得枯燥贫瘠，我在各处用一些哲学性释义加以说明，论述那些更具一般性、而哲学似乎主要建立于其上的事物，例如物体的密度和阻力、没有任何物体的空间，以及光和声音的运动。现在剩下的，就是从这些同一原理出发，证明世界体系的结构。
\end{quote}

\noindent\textbf{原文} I believe it fair to say that it was the freedom to consider problems either in a purely mathematical way or in a ``philosophical'' (or physical) way that enabled Newton to express the first law and to develop a complete inertial physics. After all, physics as a science may be developed in a mathematical way but it always must rest on experience--and experience never shows us pure inertial motion. Even in the limited examples of linear inertia discussed by Galileo, there was always some air friction and the motion ceased almost at once, as when a projectile strikes the ground. In the whole range of physics explored by Galileo there is no example of a physical object that has even a component of pure inertial motion for more than a very short time. It was perhaps for this reason that Galileo never framed a general law of inertia. He was too much a physicist.
\par\noindent\textbf{译文} 我认为可以公允地说，正是这种自由——既可以用纯数学方式，也可以用“哲学的”（或物理的）方式来思考问题——使牛顿能够表述第一定律，并发展出完整的惯性物理学。毕竟，物理学作为一门科学可以以数学方式发展，但它始终必须建立在经验之上；而经验从未向我们显示纯粹的惯性运动。即使在伽利略讨论的线性惯性的有限例子中，也总有一些空气摩擦，而且运动几乎立刻终止，例如抛射体撞到地面时。在伽利略所探索的整个物理范围内，没有任何物理物体能在超过很短时间内具有哪怕一个纯粹惯性运动的分量。也许正因为如此，伽利略从未形成一条普遍的惯性定律。他太像一位物理学家了。
\medskip

\noindent\textbf{原文} But as a mathematician Newton could easily conceive of a body's moving along a straight line at constant speed forever. The concept ``forever,'' which implies an infinite universe, held no terror for him. Observe that his statement of the law of inertia, that it is the natural condition for bodies to move in straight lines at a constant speed, occurs in Book One of the \emph{Principia}, the portion said by him to be mathematical rather than physical. Now, if it is the natural condition of motion for bodies to move uniformly in straight lines, then this kind of inertial motion must characterize the planets. The planets, however, do not move in straight lines, but rather along ellipses. Using a kind of Galilean approach to this single problem, Newton could say that the planets must therefore be subject to two motions: one inertial (along a straight line at constant speed) and one always at right angles to that straight line, drawing each planet toward its orbit. (See, further, Supplements 11 and 12.)
\par\noindent\textbf{译文} 但是作为数学家，牛顿很容易设想一个物体以恒定速度沿直线永远运动。“永远”这个概念暗含无限宇宙，却并不令他畏惧。请注意，他对惯性定律的表述——物体以恒定速度沿直线运动是其自然状态——出现在《原理》第一卷中，而牛顿说这一卷是数学的，而非物理的。现在，如果物体的自然运动状态是沿直线匀速运动，那么这种惯性运动必定也是行星的特征。然而，行星并不沿直线运动，而是沿椭圆运动。牛顿用一种类似伽利略的方法处理这个单一问题，可以说行星因此必定受制于两种运动：一种是惯性运动（以恒定速度沿直线运动），另一种总是与这条直线成直角，把每颗行星拉向其轨道。（另见补编11和12。）
\medskip

\noindent\textbf{原文} Though not moving in a straight line, each planet nevertheless represents the best example of inertial motion observable in the universe. Were it not for that component of inertial motion, the force that continually draws the planet away from the straight line would draw the planet in toward the sun until the two bodies collided. Newton once used this argument to prove the existence of God. If the planets had not received a push to give them an inertial (or tangential) component of motion, he said, the solar attractive force would not draw them into an orbit but instead would move each planet in a straight line toward the sun itself. Hence the universe could not be explained in terms of matter alone.
\par\noindent\textbf{译文} 尽管行星并不沿直线运动，每颗行星仍然代表着宇宙中可观测到的惯性运动的最佳例子。如果没有那一惯性运动分量，那么持续把行星从直线拉开的力，就会把行星向太阳拉去，直到两个物体相撞。牛顿曾用这个论证来证明上帝的存在。他说，如果行星没有受到一次推动，从而获得惯性（或切向）运动分量，那么太阳的吸引力不会把它们带入轨道，而会使每颗行星沿直线朝太阳本身运动。因此，宇宙不能仅仅用物质来解释。
\medskip

\noindent\textbf{原文} For Galileo pure circular motion could still be inertial, as in the example of an object on or near the surface of the earth. But for Newton pure circular motion was not inertial; it was accelerated and required a force for its continuance. Thus it was Newton who finally shattered the bonds of ``circularity'' which still had held Galileo in thrall. And so we may understand that it was Newton who showed how to build a celestial mechanics based on the laws of motion, since the elliptical (or almost circular) orbital motion of planets is not purely inertial, but requires additionally the constant action of a force, which turns out to be the force of universal gravitation.
\par\noindent\textbf{译文} 对伽利略来说，纯圆周运动仍然可以是惯性的，例如在地球表面或其附近的物体就是如此。但对牛顿来说，纯圆周运动不是惯性的；它是加速运动，并需要一个力来维持。因而，正是牛顿最终打碎了仍束缚着伽利略的“圆形性”枷锁。于是我们可以理解，正是牛顿展示了如何在运动定律基础上建立天体力学，因为行星的椭圆（或近乎圆形）轨道运动并非纯粹惯性运动，还需要一个力的持续作用；这个力最终证明就是万有引力。
\medskip

\noindent\textbf{原文} Thus Newton, again unlike Galileo, set out to ``demonstrate the structure of the System of the World,'' or--as we would say today--to show how the general laws of terrestrial motion may be applied to the planets and to their satellites.
\par\noindent\textbf{译文} 因此，牛顿再次不同于伽利略，他着手“证明世界体系的结构”，或者用今天的话说，展示地上运动的一般定律如何能够应用于行星及其卫星。
\medskip

\begin{center}
\ldots
\end{center}

\noindent\textbf{原文} Isaac Newton's system of mechanics came to symbolize the rational order of the world, functioning under the ``rule of nature.'' Not only could Newtonian science account for present and past phenomena; the principles could be applied to the prediction of future events. In the \emph{Principia} Newton proved that comets are like the planets, moving in great orbits that must (according to Newtonian rules) be conic sections. Some comets move in ellipses and these must return periodically from far out in space to the visible regions of our solar system, whereas others will visit our solar system and never return. Edmond Halley applied these Newtonian results to an analysis of cometary records of the past and found--among others--a comet with a period of some seventy-five and a half years. He made a bold Newtonian prediction that this comet would reappear in 1758. When it did so, right on schedule, though long after Halley and Newton were dead, men and women everywhere experienced a new feeling of awe for the powers of human reason abetted by mathematics. This new respect for science was expressed by such adjectives as ``amazing,'' ``phenomenal,'' or ``extraordinary.'' This successful prediction of a future event symbolized the force of the new science: the perfection of the mathematical understanding of nature, realized in the ability to make reliable predictions of the future. Not surprisingly, men and women everywhere saw a promise that all of human knowledge and the regulation of human affairs would yield to a similar rational system of deduction and mathematical inference coupled with experiment and critical observation. The eighteenth century not only was the Enlightenment, but became ``preeminently the age of faith in science.'' Newton became the symbol of successful science, the ideal for all thought--in philosophy, psychology, government, and the science of society.
\par\noindent\textbf{译文} 艾萨克·牛顿的力学体系逐渐成为世界理性秩序的象征，这一秩序在“自然法则”支配下运行。牛顿科学不仅能够解释现在和过去的现象；其原理还可以用于预测未来事件。在《原理》中，牛顿证明彗星与行星相似，它们沿巨大的轨道运动，而这些轨道按牛顿规则必定是圆锥曲线。有些彗星沿椭圆运动，必定会从遥远太空周期性返回太阳系的可见区域；另一些则会造访我们的太阳系后永不返回。埃德蒙·哈雷把这些牛顿结果应用于对过去彗星记录的分析，并在其中发现了一颗周期约为七十五年半的彗星。他作出一个大胆的牛顿式预言：这颗彗星将在1758年重现。后来它果然准时重现，尽管那时哈雷和牛顿都已去世很久，世界各地的人们仍对借助数学的人类理性力量产生了一种新的敬畏。这种对科学的新尊重，体现在“惊人”“非凡”或“卓越”等形容词中。对未来事件的成功预测，象征着新科学的力量：对自然的数学理解已臻完善，并表现为可靠预测未来的能力。不足为奇的是，世界各地的人们都看到了一种希望：人类全部知识以及人类事务的管理，都将服从一种类似的理性体系，这一体系把演绎和数学推理同实验及批判性观察结合起来。18世纪不仅是启蒙时代，而且成了“最突出地信仰科学的时代”。牛顿成为成功科学的象征，成为一切思想的理想典范——无论在哲学、心理学、政府，还是社会科学中。
\medskip

\noindent\textbf{原文} Newton's genius enables us to see the full significance of both Galilean mechanics and Kepler's laws of planetary motion as manifested in the development of the inertial principles required for the Copernican-Keplerian universe. A great French mathematician, Joseph Louis Lagrange (1736--1813), best defined Newton's achievement. There is only one law of the universe, he said, and Newton discovered it. Newton did not develop modern dynamics all by himself but depended heavily on certain of his predecessors; the debt in no way lessens the magnitude of his achievement. It only emphasizes the importance of such men as Galileo and Kepler, and Descartes, Hooke, and Huygens, who were great enough to make significant contributions to the Newtonian enterprise. Above all, we may see in Newton's work the degree to which science is a collective and a cumulative activity and we may find in it the measure of the influence of an individual genius on the future of a cooperative scientific effort. In Newton's achievement, we see how science advances by heroic exercises of the imagination, rather than by patient collecting and sorting of myriads of individual facts. Who, after studying Newton's magnificent contribution to thought, could deny that pure science exemplifies the creative accomplishment of the human spirit at its pinnacle?
\par\noindent\textbf{译文} 牛顿的天才使我们能够看出伽利略力学和开普勒行星运动定律的全部意义，因为它们体现在哥白尼--开普勒宇宙所需要的惯性原理的发展之中。伟大的法国数学家约瑟夫·路易·拉格朗日（1736--1813）最恰当地概括了牛顿的成就。他说，宇宙只有一条定律，而牛顿发现了它。牛顿并不是独自发展出现代动力学的，他在很大程度上依赖某些前辈；这种债务丝毫不减损其成就的伟大。它只是强调了伽利略、开普勒、笛卡尔、胡克和惠更斯等人的重要性，他们足够伟大，能够为牛顿事业作出重要贡献。最重要的是，我们可以在牛顿的工作中看到，科学在多大程度上是一种集体的、累积的活动；也可以从中衡量一个个人天才对合作性科学努力之未来的影响。在牛顿的成就中，我们看到科学并不是靠耐心收集和分类无数个别事实而前进，而是靠想象力的英雄式运用而前进。研究过牛顿对思想的宏伟贡献之后，谁还能否认纯科学体现了人类精神创造性成就的顶峰？
\medskip

\section*{Text 3b: From \emph{The Principia: Mathematical Principles of Natural Philosophy}\\文本3b：选自《自然哲学的数学原理》}
\begin{flushright}
Isaac Newton
\end{flushright}

\subsection*{Definitions\\定义}

\subsubsection*{Definition 1\\定义一}
\noindent\textbf{原文} Quantity of matter is a measure of matter that arises from its density and volume jointly.
\par\noindent\textbf{译文} 物质的量（即质量）是由密度和体积共同决定的物质量度。
\medskip

\noindent\textbf{原文} If the density of air is doubled in a space that is also doubled, there is four times as much air, and there is six times as much if the space is tripled. The case is the same for snow and powders condensed by compression or liquefaction, and also for all bodies that are condensed in various ways by any causes whatsoever. For the present, I am not taking into account any medium, if there should be any, freely pervading the interstices between the parts of bodies. Furthermore, I mean this quantity whenever I use the term ``body'' or ``mass'' in the following pages. It can always be known from a body's weight, for--by making very accurate experiments with pendulums--I have found it to be proportional to the weight, as will be shown below.
\par\noindent\textbf{译文} 如果在一个也扩大为两倍的空间中，空气密度加倍，那么空气的量就是四倍；若空间扩大为三倍，则空气的量就是六倍。被压缩或液化而凝聚的雪和粉末也是如此；一切由于任何原因以各种方式被凝聚的物体也是如此。目前我并不考虑是否存在某种介质，能够自由贯穿物体各部分之间的空隙。此外，在以下各页中，我使用“物体”或“质量”一词时，指的就是这个量。这个量总是可以由物体的重量得知；因为通过非常精确的摆实验，我发现它与重量成正比，这将在后文说明。
\medskip

\subsubsection*{Definition 2\\定义二}
\noindent\textbf{原文} Quantity of motion is a measure of motion that arises from the velocity and the quantity of matter jointly.
\par\noindent\textbf{译文} 动量是由速度和物质的量（质量）共同决定的运动量度。
\medskip

\noindent\textbf{原文} The motion of a whole is the sum of the motions of the individual parts, and thus if a body is twice as large as another and has equal velocity there is twice as much motion, and if it has twice the velocity there is four times as much motion.
\par\noindent\textbf{译文} 整体的运动是各个部分运动的总和；因此，如果一个物体是另一个物体的两倍大且速度相同，它就有两倍的动量；如果它又有两倍的速度，它就有四倍的动量。
\medskip

\subsubsection*{Definition 3\\定义三}
\noindent\textbf{原文} Inherent force of matter is the power of resisting by which every body, so far as it is able, perseveres in its state either of resting or of moving uniformly straight forward.
\par\noindent\textbf{译文} 物质的固有力，是每个物体借以抵抗的能力；凭借这种能力，物体尽其所能保持其静止状态，或沿直线作匀速运动的状态。
\medskip

\noindent\textbf{原文} This force is always proportional to the body and does not differ in any way from the inertia of the mass except in the manner in which it is conceived. Because of the inertia of matter, every body is only with difficulty put out of its state either of resting or of moving. Consequently, inherent force may also be called by the very significant name of force of inertia. Moreover, a body exerts this force only during a change of its state, caused by another force impressed upon it, and this exercise of force is, depending on the viewpoint, both resistance and impetus: resistance insofar as the body, in order to maintain its state, strives against the impressed force, and impetus insofar as the same body, yielding only with difficulty to the force of a resisting obstacle, endeavors to change the state of that obstacle. Resistance is commonly attributed to resting bodies and impetus to moving bodies; but motion and rest, in the popular sense of the terms, are distinguished from each other only by point of view, and bodies commonly regarded as being at rest are not always truly at rest.
\par\noindent\textbf{译文} 这种力总是与物体成正比；除了理解方式不同之外，它与质量的惯性没有任何差别。由于物质具有惯性，每个物体都只能很困难地脱离其静止状态或运动状态。因此，固有力也可以用一个非常有意义的名称来称呼，即惯性力。此外，只有当另一个施加在物体上的力导致其状态改变时，物体才发挥这种力；这种力的发挥，依观察角度而定，既是抵抗，也是冲力：说它是抵抗，是因为物体为了保持其状态而努力反抗所施加的力；说它是冲力，是因为同一物体只是在困难地屈从于阻碍物的力时，努力改变那个阻碍物的状态。抵抗通常归于静止物体，冲力通常归于运动物体；但按通俗意义来说，运动和静止仅仅因观察角度不同而彼此区别，那些通常被认为静止的物体并不总是真正静止。
\medskip

\subsubsection*{Definition 4\\定义四}
\noindent\textbf{原文} Impressed force is the action exerted on a body to change its state either of resting or of moving uniformly straight forward.
\par\noindent\textbf{译文} 施加力，是作用于物体、使其静止状态或沿直线匀速运动状态发生改变的作用。
\medskip

\noindent\textbf{原文} This force consists solely in the action and does not remain in a body after the action has ceased. For a body perseveres in any new state solely by the force of inertia. Moreover, there are various sources of impressed force, such as percussion, pressure, or centripetal force.
\par\noindent\textbf{译文} 这种力仅存在于作用之中；作用停止后，它并不留存在物体内。因为物体完全凭借惯性力而保持在任何新的状态中。此外，施加力有多种来源，例如撞击、压力或向心力。
\medskip

\subsubsection*{Definition 5\\定义五}
\noindent\textbf{原文} Centripetal force is the force by which bodies are drawn from all sides, are impelled, or in any way tend, toward some point as to a center.
\par\noindent\textbf{译文} 向心力，是物体从各个方向被吸引、被推动，或以任何方式趋向某一点如同趋向中心的力。
\medskip

\noindent\textbf{原文} One force of this kind is gravity, by which bodies tend toward the center of the earth; another is magnetic force, by which iron seeks a lodestone; and yet another is that force, whatever it may be, by which the planets are continually drawn back from rectilinear motions and compelled to revolve in curved lines.
\par\noindent\textbf{译文} 这种力的一种是重力，物体借此趋向地球中心；另一种是磁力，铁借此趋向磁石；还有一种是使行星不断从直线运动中被拉回并被迫沿曲线运行的力，无论这种力究竟是什么。
\medskip

\noindent\textbf{原文} A stone whirled in a sling endeavors to leave the hand that is whirling it, and by its endeavor it stretches the sling, doing so the more strongly the more swiftly it revolves; and as soon as it is released, it flies away. The force opposed to that endeavor, that is, the force by which the sling continually draws the stone back toward the hand and keeps it in an orbit, I call centripetal, since it is directed toward the hand as toward the center of an orbit. And the same applies to all bodies that are made to move in orbits. They all endeavor to recede from the centers of their orbits, and unless some force opposed to that endeavor is present, restraining them and keeping them in orbits and hence called by me centripetal, they will go off in straight lines with uniform motion. If a projectile were deprived of the force of gravity, it would not be deflected toward the earth but would go off in a straight line into the heavens and do so with uniform motion, provided that the resistance of the air were removed. The projectile, by its gravity, is drawn back from a rectilinear course and continually deflected toward the earth, and this is so to a greater or lesser degree in proportion to its gravity and its velocity of motion. The less its gravity in proportion to its quantity of matter, or the greater the velocity with which it is projected, the less it will deviate from a rectilinear course and the farther it will go. If a lead ball were projected with a given velocity along a horizontal line from the top of some mountain by the force of gunpowder and went in a curved line for a distance of two miles before falling to the earth, then the same ball projected with twice the velocity would go about twice as far and with ten times the velocity about ten times as far, provided that the resistance of the air were removed.
\par\noindent\textbf{译文} 用投石索旋转的石块努力离开旋转它的手；由于这种努力，它拉紧投石索，而且旋转得越快，拉得越强；一旦被释放，它就飞走。与这种努力相反的力，也就是投石索不断把石块拉回手的方向并使它保持在轨道中的力，我称之为向心力，因为它指向手，如同指向轨道的中心。所有被迫在轨道中运动的物体也都是如此。它们都努力离开自身轨道的中心；如果没有某种与这种努力相反的力存在，没有某种约束它们、使它们保持在轨道中并因而被我称为向心的力，它们就会沿直线作匀速运动而离去。如果一个抛射体失去重力，而且空气阻力被消除，它就不会向地球偏转，而会沿直线匀速飞向天空。抛射体由于自身重力而从直线路径被拉回，并不断向地球偏转；偏转程度的大小，与其重力和运动速度有关。相对于其物质的量（质量）而言重力越小，或者抛出时速度越大，它偏离直线路径就越少，飞得就越远。假如一个铅球从某座山顶被火药之力以给定速度沿水平线发射，在落到地面前沿曲线前进了两英里，那么在去除空气阻力的条件下，同一个球以两倍速度发射就会走大约两倍远，以十倍速度发射就会走大约十倍远。
\medskip

\noindent\textbf{原文} And by increasing the velocity, the distance to which it would be projected could be increased at will and the curvature of the line that it would describe could be decreased, in such a way that it would finally fall at a distance of 10 or 30 or 90 degrees or even go around the whole earth or, lastly, go off into the heavens and continue indefinitely in this motion. And in the same way that a projectile could, by the force of gravity, be deflected into an orbit and go around the whole earth, so too the moon, whether by the force of gravity--if it has gravity--or by any other force by which it may be urged toward the earth, can always be drawn back toward the earth from a rectilinear course and deflected into its orbit; and without such a force the moon cannot be kept in its orbit. If this force were too small, it would not deflect the moon sufficiently from a rectilinear course; if it were too great, it would deflect the moon excessively and draw it down from its orbit toward the earth. In fact, it must be of just the right magnitude, and mathematicians have the task of finding the force by which a body can be kept exactly in any given orbit with a given velocity and, alternatively, to find the curvilinear path into which a body leaving any given place with a given velocity is deflected by a given force.
\par\noindent\textbf{译文} 通过增加速度，它被抛射到的距离可以任意增加，它所描绘曲线的曲率可以减小，以至于它最终会在相距10度、30度或90度处落下，甚至绕整个地球运行；最后，它也可能飞向天空，并无限期地继续这种运动。正如抛射体可以因重力而被偏转进入轨道并环绕整个地球一样，月球也可以——无论是由于重力（如果它有重力），还是由于任何其他使它趋向地球的力——总是从直线路径被拉回地球方向，并被偏转进入自己的轨道；没有这样的力，月球就不能保持在其轨道中。如果这个力太小，它就不能使月球充分偏离直线路径；如果太大，它就会使月球过度偏转，并把月球从轨道上拉向地球。事实上，这个力必须大小正合适；数学家的任务就是找出能够使物体以给定速度恰好保持在任一给定轨道中的力，反过来也要找出物体从任一给定地点以给定速度出发时，在给定力作用下被偏转成的曲线路径。
\medskip

\noindent\textbf{原文} The quantity of centripetal force is of three kinds: absolute, accelerative, and motive.
\par\noindent\textbf{译文} 向心力的量有三种：绝对量、加速量和动力量。
\medskip

\begin{center}
\ldots
\end{center}

\noindent\textbf{原文} Further, it is in this same sense that I call attractions and impulses accelerative and motive. Moreover, I use interchangeably and indiscriminately words signifying attraction, impulse, or any sort of propensity toward a center, considering these forces not from a physical but only from a mathematical point of view. Therefore, let the reader beware of thinking that by words of this kind I am anywhere defining a species or mode of action or a physical cause or reason, or that I am attributing forces in a true and physical sense to centers (which are mathematical points) if I happen to say that centers attract or that centers have forces.
\par\noindent\textbf{译文} 此外，也正是在同一意义上，我把吸引和推动称为加速性的和动力性的。并且，我可以互换而不加区别地使用表示吸引、推动或任何趋向中心之倾向的词，因为我并不是从物理观点，而只是从数学观点来考察这些力。因此，读者切勿以为我用这类词是在任何地方定义某种作用的种类或方式，或某种物理原因或理由；也不要以为当我偶尔说中心会吸引或中心具有力时，我是在真实的物理意义上把力归于中心（中心只是数学点）。
\medskip

\subsection*{Axioms, or the Laws of Motion\\公理，或运动定律}

\subsubsection*{Law 1\\定律一}
\noindent\textbf{原文} Every body perseveres in its state of being at rest or of moving uniformly straight forward, except insofar as it is compelled to change its state by forces impressed.
\par\noindent\textbf{译文} 每一个物体都会保持其静止状态，或沿直线作匀速运动的状态，除非它被所施加的力迫使改变这种状态。
\medskip

\noindent\textbf{原文} Projectiles persevere in their motions, except insofar as they are retarded by the resistance of the air and are impelled downward by the force of gravity. A spinning hoop, which has parts that by their cohesion continually draw one another back from rectilinear motions, does not cease to rotate, except insofar as it is retarded by the air. And larger bodies--planets and comets--preserve for a longer time both their progressive and their circular motions, which take place in spaces having less resistance.
\par\noindent\textbf{译文} 抛射体会保持其运动，除非它们受到空气阻力而迟滞，或被重力向下推动。一个旋转的环，其各部分由于凝聚力而不断把彼此从直线运动中拉回；除非受到空气迟滞，它不会停止旋转。更大的物体——行星和彗星——则在阻力较小的空间中，更长久地保持其前进运动和圆周运动。
\medskip

\subsubsection*{Law 2\\定律二}
\noindent\textbf{原文} A change in motion is proportional to the motive force impressed and takes place along the straight line in which that force is impressed.
\par\noindent\textbf{译文} 动量的改变与所施加的动力成正比，并沿着施加该力的直线发生。
\medskip

\noindent\textbf{原文} If some force generates any motion, twice the force will generate twice the motion, and three times the force will generate three times the motion, whether the force is impressed all at once or successively by degrees. And if the body was previously moving, the new motion (since motion is always in the same direction as the generative force) is added to the original motion if that motion was in the same direction or is subtracted from the original motion if it was in the opposite direction or, if it was in an oblique direction, is combined obliquely and compounded with it according to the directions of both motions.
\par\noindent\textbf{译文} 如果某个力产生某种动量，两倍的力就会产生两倍的动量，三倍的力就会产生三倍的动量；无论该力是一次施加，还是逐渐连续施加，都是如此。如果物体先前已经在运动，新的动量（由于动量总是与产生它的力同方向）若与原动量同方向，就加到原动量上；若与原动量相反，就从原动量中减去；若成斜向，则按照两个动量的方向作斜向结合并合成。
\medskip

\subsubsection*{Law 3\\定律三}
\noindent\textbf{原文} To any action there is always an opposite and equal reaction; in other words, the actions of two bodies upon each other are always equal and always opposite in direction.
\par\noindent\textbf{译文} 对于任何作用，总有一个相反且相等的反作用；换言之，两个物体相互之间的作用总是大小相等，方向相反。
\medskip

\noindent\textbf{原文} Whatever presses or draws something else is pressed or drawn just as much by it. If anyone presses a stone with a finger, the finger is also pressed by the stone. If a horse draws a stone tied to a rope, the horse will (so to speak) also be drawn back equally toward the stone, for the rope, stretched out at both ends, will urge the horse toward the stone and the stone toward the horse by one and the same endeavor to go slack and will impede the forward motion of the one as much as it promotes the forward motion of the other. If some body impinging upon another body changes the motion of that body in any way by its own force, then, by the force of the other body (because of the equality of their mutual pressure), it also will in turn undergo the same change in its own motion in the opposite direction. By means of these actions, equal changes occur in the motions, not in the velocities--that is, of course, if the bodies are not impeded by anything else. For the changes in velocities that likewise occur in opposite directions are inversely proportional to the bodies because the motions are changed equally. This law is valid also for attractions, as will be proved in the next scholium.
\par\noindent\textbf{译文} 凡是压迫或牵引别物者，也被别物以同样大小压迫或牵引。如果有人用手指压一块石头，手指也被石头压。如果一匹马牵引一块系在绳上的石头，那么这匹马（可以这样说）也同样被向石头方向拉回；因为绳子在两端被拉紧时，会由于同一个趋于松弛的努力，把马推向石头，也把石头推向马，并且它阻碍一方前进运动的程度，正如它促进另一方前进运动的程度一样。如果某个物体撞击另一个物体，并以自身的力以某种方式改变那个物体的动量，那么由于另一物体的力（因为二者相互压力相等），它自身也会反过来在相反方向上经历同样的动量改变。通过这些作用，发生相等的动量改变，而不是相等的速度改变——当然，前提是这些物体不受其他事物阻碍。因为既然动量改变相等，那么同样在相反方向发生的速度改变就与物体的质量成反比。这条定律也适用于吸引作用，这将在下一条释义中得到证明。
\medskip

\subsubsection*{Corollary 1\\推论一}
\noindent\textbf{原文} A body acted on by [two] forces acting jointly describes the diagonal of a parallelogram in the same time in which it would describe the sides if the forces were acting separately.
\par\noindent\textbf{译文} 一个物体在[两个]力共同作用下所描绘的平行四边形对角线，所用时间与这些力分别作用时它描绘各边所用的时间相同。
\medskip

\begin{center}
\safeincludeimage{86e81927fb72b006792bc4605366d6baa2c5963aedd08e9e51d23ec99fb2611f.jpg}
\end{center}

\noindent\textbf{原文} Let a body in a given time, by force $M$ alone impressed in $A$, be carried with uniform motion from $A$ to $B$, and, by force $N$ alone impressed in the same place, be carried from $A$ to $C$; then complete the parallelogram $ABDC$, and by both forces the body will be carried in the same time along the diagonal from $A$ to $D$. For, since force $N$ acts along the line $AC$ parallel to $BD$, this force, by law 2, will make no change at all in the velocity toward the line $BD$ which is generated by the other force. Therefore, the body will reach the line $BD$ in the same time whether force $N$ is impressed or not, and so at the end of that time will be found somewhere on the line $BD$. By the same argument, at the end of the same time it will be found somewhere on the line $CD$, and accordingly it is necessarily found at the intersection $D$ of both lines. And, by law 1, it will go with [uniform] rectilinear motion from $A$ to $D$.
\par\noindent\textbf{译文} 设一个物体在给定时间内，若只受在$A$处施加的力$M$作用，便以匀速运动从$A$被带到$B$；若只受同一地点施加的力$N$作用，便从$A$被带到$C$。现在补成平行四边形$ABDC$，则在两个力共同作用下，物体将在同一时间内沿对角线从$A$被带到$D$。因为力$N$沿$AC$线作用，而$AC$平行于$BD$；根据定律二，这个力不会使另一个力所产生的、朝向$BD$线的速度发生任何改变。因此，无论是否施加力$N$，物体都会在同一时间内到达$BD$线，所以在该时间结束时，它将位于$BD$线上的某处。用同样的论证，在同一时间结束时，它也将位于$CD$线上的某处；于是它必定位于两条线的交点$D$。并且，根据定律一，它将以[匀速]直线运动从$A$到$D$。
\medskip

\begin{center}
\ldots
\end{center}

\begin{center}
\safeincludeimage{96a7884762717ddc050668629e01a03990e49f38be99f72b422c072fe378ebb4.jpg}
\end{center}

\begin{center}
\safeincludeimage{453ff7a961a91fd9458260fc2cb186d44f83f69bb271740c7c73792307573ad1.jpg}
\end{center}

\begin{center}
\safeincludeimage{a57a6587536014598149960fe6cfc7bd98e69910545d5bb0e3e380f76308e9ba.jpg}
\end{center}


% ===== 04_darwin.tex =====
% Cleaned bilingual LaTeX from chunks/04_darwin.md.
% Compile with XeLaTeX or LuaLaTeX for Chinese text. Requires graphicx for preserved images.

\safeincludeimage{5f8522bcc13764df00ac41eb9c256db35f9c3982afaeab440b049ff6ba14d2ff.jpg}
\safeincludeimage{04890d618174d03e224376ed406b2bf95249a4704840eb328a547de8c5b8d66a.jpg}
\safeincludeimage{298c02d16853568398e971ddca4f30f6ce8faa50bad1e7e8b79389d63480f4ed.jpg}

\section*{On the Origin of Species\\《物种起源》}
\begin{center}
Charles Darwin\\
查尔斯·达尔文
\end{center}

\subsection*{Chapter IV: Natural Selection\\第四章：自然选择}

\pair{Natural Selection—its power compared with man's selection—its power on characters of trifling importance—its power at all ages and on both sexes—Sexual Selection—On the generality of intercrosses between individuals of the same species—Circumstances favourable and unfavourable to Natural Selection, namely, intercrossing, isolation, number of individuals—Slow action—Extinction caused by Natural Selection—Divergence of Character, related to the diversity of inhabitants of any small area, and to naturalisation—Action of Natural Selection, through Divergence of Character and Extinction, on the descendants from a common parent—Explains the Grouping of all organic beings.}{自然选择——它同人工选择相比的力量——它对极细微性状的作用——它在各个年龄和两性中的作用——性选择——同种个体之间杂交的普遍性——有利于和不利于自然选择的条件，即杂交、隔离和个体数量——缓慢的作用——自然选择所造成的灭绝——性状分歧，它同任何小区域内生物的多样性以及归化有关——自然选择通过性状分歧和灭绝，对共同祖先的后裔所发生的作用——解释一切生物的类群排列。}

\pair{How will the struggle for existence, discussed too briefly in the last chapter, act in regard to variation? Can the principle of selection, which we have seen is so potent in the hands of man, apply in nature? I think we shall see that it can act most effectually. Let it be borne in mind in what an endless number of strange peculiarities our domestic productions, and, in a lesser degree, those under nature, vary; and how strong the hereditary tendency is. Under domestication, it may be truly said that the whole organisation becomes in some degree plastic. Let it be borne in mind how infinitely complex and close-fitting are the mutual relations of all organic beings to each other and to their physical conditions of life. Can it, then, be thought improbable, seeing that variations useful to man have undoubtedly occurred, that other variations useful in some way to each being in the great and complex battle of life should sometimes occur in the course of thousands of generations? If such do occur, can we doubt, remembering that many more individuals are born than can possibly survive, that individuals having any advantage, however slight, over others would have the best chance of surviving and of procreating their kind? On the other hand, we may feel sure that any variation in the least degree injurious would be rigidly destroyed. This preservation of favourable variations and the rejection of injurious variations, I call Natural Selection. Variations neither useful nor injurious would not be affected by natural selection, and would be left a fluctuating element, as perhaps we see in the species called polymorphic.}{上一章过于简略地讨论过的生存斗争，将怎样作用于变异呢？我们已经看到，选择原则在人的手中具有巨大力量；它能否应用于自然界呢？我想我们将会看到，它能够极其有效地发生作用。请记住，我们的家养生物，以及程度较低的自然状态下的生物，会在无穷无尽的奇特性状上发生变异；还要记住遗传倾向是多么强。在家养状态下，确实可以说整个机体在某种程度上都变得具有可塑性。再请记住，一切生物彼此之间以及它们同物理生活条件之间的相互关系，是何等无限复杂而又严密契合。那么，既然毫无疑问已经出现过对人有用的变异，难道在成千上万代的过程中，有时不会出现某些以某种方式有利于每一种生物在复杂生命大战中生存的变异吗？如果这类变异确会发生，那么，只要想到出生的个体远比可能存活的个体多，我们还能怀疑那些比其他个体具有任何优势、哪怕极其微小优势的个体，最有机会存活并繁殖其同类吗？另一方面，我们也可以确信，任何哪怕只有最轻微有害的变异，都会被严格消灭。保存有利变异、排斥有害变异，我称之为自然选择。既无益也无害的变异不会受到自然选择的影响，而会作为一种浮动因素保留下来，也许正如我们在所谓多型种中所见到的那样。}

\pair{We shall best understand the probable course of natural selection by taking the case of a country undergoing some physical change, for instance, of climate. The proportional numbers of its inhabitants would almost immediately undergo a change, and some species might become extinct. We may conclude, from what we have seen of the intimate and complex manner in which the inhabitants of each country are bound together, that any change in the numerical proportions of some of the inhabitants, independently of the change of climate itself, would most seriously affect many of the others. If the country were open on its borders, new forms would certainly immigrate, and this also would seriously disturb the relations of some of the former inhabitants. Let it be remembered how powerful the influence of a single introduced tree or mammal has been shown to be. But in the case of an island, or of a country partly surrounded by barriers, into which new and better adapted forms could not freely enter, we should then have places in the economy of nature which would assuredly be better filled up if some of the original inhabitants were in some manner modified; for, had the area been open to immigration, these same places would have been seized on by intruders. In such case, every slight modification which in the course of ages chanced to arise, and which in any way favoured the individuals of any of the species by better adapting them to their altered conditions, would tend to be preserved; and natural selection would thus have free scope for the work of improvement.}{要最清楚地理解自然选择可能的进程，我们不妨设想一个国家正在经历某种物理变化，例如气候变化。它的生物居民在数量比例上几乎会立即发生改变，有些种还可能灭绝。我们已经看到，每个地区的生物居民以亲密而复杂的方式彼此联系；因此可以推断，某些居民数量比例的任何变化，即使撇开气候变化本身不论，也会极其严重地影响许多其他居民。如果这个地区的边界是开放的，新的类型必定会迁入，而这也会严重扰乱原有某些居民之间的关系。请记住，单单引入一种树木或一种哺乳动物，就已显示出何等强大的影响。可是，对于岛屿，或部分被障碍包围、使新的并且更适应的类型不能自由进入的地区来说，自然经济中就会有一些位置；倘若某些原有居民以某种方式被改造，这些位置必定会被更好地填补。因为如果该区域向移入开放，同样的位置就会被外来的入侵者占据。在这种情况下，漫长岁月中偶然发生的每一种轻微改变，只要能以任何方式有利于某一种个体，使它们更好地适应改变了的条件，就会倾向于被保存下来；自然选择由此便有充分余地去进行改进的工作。}

\pair{We have reason to believe, as stated in the first chapter, that a change in the conditions of life, by specially acting on the reproductive system, causes or increases variability; and in the foregoing case the conditions of life are supposed to have undergone a change, and this would manifestly be favourable to natural selection, by giving a better chance of profitable variations occurring; and unless profitable variations do occur, natural selection can do nothing. Not that, as I believe, any extreme amount of variability is necessary; as man can certainly produce great results by adding up in any given direction mere individual differences, so could Nature, but far more easily, from having incomparably longer time at her disposal. Nor do I believe that any great physical change, as of climate, or any unusual degree of isolation to check immigration, is actually necessary to produce new and unoccupied places for natural selection to fill up by modifying and improving some of the varying inhabitants. For as all the inhabitants of each country are struggling together with nicely balanced forces, extremely slight modifications in the structure or habits of one inhabitant would often give it an advantage over others; and still further modifications of the same kind would often still further increase the advantage. No country can be named in which all the native inhabitants are now so perfectly adapted to each other and to the physical conditions under which they live, that none of them could anyhow be improved; for in all countries, the natives have been so far conquered by naturalised productions that they have allowed foreigners to take firm possession of the land. And as foreigners have thus everywhere beaten some of the natives, we may safely conclude that the natives might have been modified with advantage, so as to have better resisted such intruders.}{正如第一章所说，我们有理由相信，生活条件的改变由于特别作用于生殖系统，会引起或增加变异性；在上述例子中，我们假定生活条件已经改变，这显然有利于自然选择，因为它使有益变异的发生有更好的机会；而如果没有有益变异发生，自然选择便无能为力。并不是说，我认为必须有极大量的变异性；人能够仅仅把某一方向上的个体差异累加起来而确实产生巨大结果，自然也能如此，而且由于它拥有无可比拟地更长的时间，就更加容易。我也不认为必须有任何重大的物理变化，例如气候变化，或必须有任何异常程度的隔离来阻止迁入，才会产生新的、空缺的位置，供自然选择通过改造和改进某些变异中的居民来填补。因为每个地区的一切居民都以精细平衡的力量共同斗争，某一居民在构造或习性上的极其轻微改变，往往就会使它比其他居民占有优势；而同类的进一步改变，又往往会进一步增加这种优势。没有任何一个国家可以说，它的一切本地居民现在已经如此完美地彼此适应，并适应其所处的物理条件，以至于它们无论如何都不能再被改进；因为在所有国家，本地生物都已经在某种程度上被归化生物所征服，允许外来者牢固地占据土地。既然外来者在各处都曾击败某些本地生物，我们便可以稳妥地断定：本地生物本来可以有利地被改造，从而更好地抵抗这些入侵者。}

\pair{As man can produce and certainly has produced a great result by his methodical and unconscious means of selection, what may not nature effect? Man can act only on external and visible characters: nature cares nothing for appearances, except in so far as they may be useful to any being. She can act on every internal organ, on every shade of constitutional difference, on the whole machinery of life. Man selects only for his own good; Nature only for that of the being which she tends. Every selected character is fully exercised by her; and the being is placed under well-suited conditions of life. Man keeps the natives of many climates in the same country; he seldom exercises each selected character in some peculiar and fitting manner; he feeds a long- and a short-beaked pigeon on the same food; he does not exercise a long-backed or long-legged quadruped in any peculiar manner; he exposes sheep with long and short wool to the same climate. He does not allow the most vigorous males to struggle for the females. He does not rigidly destroy all inferior animals, but protects during each varying season, as far as lies in his power, all his productions. He often begins his selection by some half-monstrous form; or at least by some modification prominent enough to catch his eye, or to be plainly useful to him. Under nature, the slightest difference of structure or constitution may well turn the nicely-balanced scale in the struggle for life, and so be preserved. How fleeting are the wishes and efforts of man! how short his time! and consequently how poor will his products be, compared with those accumulated by nature during whole geological periods. Can we wonder, then, that nature's productions should be far “truer” in character than man's productions; that they should be infinitely better adapted to the most complex conditions of life, and should plainly bear the stamp of far higher workmanship?}{既然人能够并且确实已经通过有计划的和无意识的选择方法产生巨大结果，那么自然又有什么不能做到呢？人只能作用于外部的、可见的性状；自然并不关心外表，除非外表对某种生物有用。自然能够作用于每一个内部器官，作用于体质差异的每一种细微层次，作用于整个生命机器。人只是为自己的利益而选择；自然只为它所照管的生物的利益而选择。每一种被选择的性状都由自然充分加以运用；生物也被置于适宜的生活条件之下。人把许多气候下的本地生物养在同一个国家；他很少以某种特殊而合适的方式锻炼每一种被选择的性状；他用同样的食物喂养长喙和短喙的鸽子；他并不以任何特殊方式锻炼长背或长腿的四足动物；他让长毛和短毛的绵羊暴露于同一种气候之下。他不允许最强健的雄性为争夺雌性而斗争。他并不严格消灭一切劣等动物，而是在每个变化的季节里尽其所能保护他的所有产物。他的选择往往从某种半畸形的类型开始；或至少从某种显著到足以引起他注意、或明显对他有用的改变开始。在自然状态下，构造或体质上最轻微的差异，完全可能在生存斗争中拨动那架精密平衡的天平，从而被保存下来。人的愿望和努力多么短暂！他的时间多么短促！因此，同自然在整个地质时期中累积起来的产物相比，人的产物将会何等贫弱。既然如此，自然的产物在性状上远比人的产物更“真实”，更能无限良好地适应最复杂的生活条件，并且明显带有远为高超的工艺印记，我们又有什么可惊奇的呢？}

\pair{It may be said that natural selection is daily and hourly scrutinising, throughout the world, every variation, even the slightest; rejecting that which is bad, preserving and adding up all that is good; silently and insensibly working, whenever and wherever opportunity offers, at the improvement of each organic being in relation to its organic and inorganic conditions of life. We see nothing of these slow changes in progress, until the hand of time has marked the long lapse of ages, and then so imperfect is our view into long past geological ages that we only see that the forms of life are now different from what they formerly were.}{可以说，自然选择在全世界每时每刻都在审察每一个变异，哪怕最轻微的变异也不放过；它排斥坏的，保存并累加一切好的；只要机会出现，无论何时何地，它都在静默而不知不觉地工作，使每一种生物在同其有机和无机生活条件的关系中得到改进。我们看不见这些缓慢变化正在进行，直到时间之手标记出漫长岁月的流逝；而那时我们对于久远地质年代的观察又如此不完善，以至于只能看出现在的生命类型不同于从前。}

\pair{Although natural selection can act only through and for the good of each being, yet characters and structures which we are apt to consider as of very trifling importance may thus be acted on. When we see leaf-eating insects green, and bark-feeders mottled grey; the alpine ptarmigan white in winter, the red grouse the colour of heather, and the black grouse that of peaty earth, we must believe that these tints are of service to these birds and insects in preserving them from danger. Grouse, if not destroyed at some period of their lives, would increase in countless numbers; they are known to suffer largely from birds of prey; and hawks are guided by eyesight to their prey, so much so that on parts of the Continent persons are warned not to keep white pigeons, as being the most liable to destruction. Hence I can see no reason to doubt that natural selection might be most effective in giving the proper colour to each kind of grouse, and in keeping that colour, when once acquired, true and constant. Nor ought we to think that the occasional destruction of an animal of any particular colour would produce little effect: we should remember how essential it is in a flock of white sheep to destroy every lamb with the faintest trace of black. In plants the down on the fruit and the colour of the flesh are considered by botanists as characters of the most trifling importance: yet we hear from an excellent horticulturist, Downing, that in the United States smooth-skinned fruits suffer far more from a beetle, a curculio, than those with down; that purple plums suffer far more from a certain disease than yellow plums; whereas another disease attacks yellow-fleshed peaches far more than those with other coloured flesh. If, with all the aids of art, these slight differences make a great difference in cultivating the several varieties, assuredly, in a state of nature, where the trees would have to struggle with other trees and with a host of enemies, such differences would effectually settle which variety, whether a smooth or downy, a yellow- or purple-fleshed fruit, should succeed.}{尽管自然选择只能通过并为了每一种生物的利益而发生作用，但我们容易认为极其微不足道的性状和构造，也可能因此受到作用。当我们看到食叶昆虫呈绿色，食树皮的昆虫呈斑驳的灰色；高山雷鸟冬季为白色，红松鸡呈石南花的颜色，黑松鸡呈泥炭土的颜色时，我们必须相信，这些色调有助于这些鸟类和昆虫免遭危险。松鸡若在生命的某个时期不被消灭，就会增加到不可计数；众所周知，它们大量遭受猛禽的危害；而鹰隼凭视觉寻找猎物，甚至在欧洲大陆的某些地方，人们会被告诫不要饲养白鸽，因为白鸽最容易被捕杀。因此，我看不出有什么理由怀疑，自然选择能够极其有效地赋予每一种松鸡以适当的颜色，并在这种颜色一经获得后，使它真实而恒定地保持下去。我们也不应以为偶尔消灭某种特定颜色的动物不会产生多大效果；我们应当记住，在一群白羊中，消灭每一只带有最轻微黑色痕迹的羔羊是多么必要。在植物中，果实上的茸毛和果肉的颜色被植物学家看作极其细微的性状；可是我们从杰出的园艺家唐宁那里听说，在美国，光滑果皮的果实比有茸毛的果实更容易遭受一种象甲危害；紫色李子比黄色李子更容易遭受某种病害；而另一种病害又比起其他颜色果肉的桃子，更容易侵袭黄肉桃。即使借助所有人工手段，这些细微差异在栽培若干品种时也会造成巨大差别；那么在自然状态下，树木必须同其他树木以及大批敌害斗争时，这些差异必定会有效决定哪一种品种能够成功，是光滑的还是有茸毛的果实，是黄色果肉还是紫色果肉的果实。}

\begin{center}
[\ldots]
\end{center}

\pair{Natural selection will modify the structure of the young in relation to the parent, and of the parent in relation to the young. In social animals it will adapt the structure of each individual for the benefit of the community, if each in consequence profits by the selected change. What natural selection cannot do is to modify the structure of one species, without giving it any advantage, for the good of another species; and though statements to this effect may be found in works of natural history, I cannot find one case which will bear investigation. A structure used only once in an animal's whole life, if of high importance to it, might be modified to any extent by natural selection; for instance, the great jaws possessed by certain insects, and used exclusively for opening the cocoon, or the hard tip to the beak of nestling birds, used for breaking the egg. It has been asserted that of the best short-beaked tumbler-pigeons more perish in the egg than are able to get out of it; so that fanciers assist in the act of hatching. Now, if nature had to make the beak of a full-grown pigeon very short for the bird's own advantage, the process of modification would be very slow, and there would be simultaneously the most rigorous selection of the young birds within the egg which had the most powerful and hardest beaks, for all with weak beaks would inevitably perish; or, more delicate and more easily broken shells might be selected, the thickness of the shell being known to vary like every other structure.}{自然选择会改变幼体相对于亲体的构造，也会改变亲体相对于幼体的构造。在社会性动物中，如果每一个个体都因此从所选择的改变中获益，自然选择就会为了群体的利益而使每个个体的构造得到适应。自然选择所不能做的，是不赋予某一种任何优势，却为了另一种的利益而改变前一种的构造；尽管自然史著作中可以找到这类说法，我却找不到一个经得起考察的例子。某种构造即使在动物一生中只使用一次，只要对它极其重要，也可以经由自然选择被改变到任何程度；例如某些昆虫拥有的大颚，专门用于打开茧；又如雏鸟喙端的坚硬突起，用于破壳而出。有人断言，在最优良的短喙翻飞鸽中，死在卵内的比能出壳的还多；因此养鸽者会协助孵化。现在，如果自然为了成鸽本身的利益而不得不使鸽子的喙变得很短，改变过程将会非常缓慢；同时，卵内幼鸟中那些拥有最有力、最坚硬喙的个体会受到最严格的选择，因为喙软弱的个体必然死亡；或者，更薄、更容易破裂的蛋壳可能被选择，因为壳的厚度也同其他一切构造一样会发生变异。}

\subsection*{Sexual Selection\\性选择}

\pair{Inasmuch as peculiarities often appear under domestication in one sex and become hereditarily attached to that sex, the same fact probably occurs under nature; and if so, natural selection will be able to modify one sex in its functional relations to the other sex, or in relation to wholly different habits of life in the two sexes, as is sometimes the case with insects. And this leads me to say a few words on what I call Sexual Selection. This depends not on a struggle for existence, but on a struggle between the males for possession of the females; the result is not death to the unsuccessful competitor, but few or no offspring. Sexual selection is, therefore, less rigorous than natural selection. Generally, the most vigorous males, those which are best fitted for their places in nature, will leave most progeny. But in many cases victory will depend not on general vigour, but on having special weapons, confined to the male sex. A hornless stag or spurless cock would have a poor chance of leaving offspring. Sexual selection, by always allowing the victor to breed, might surely give indomitable courage, length to the spur, and strength to the wing to strike in the spurred leg, as well as the brutal cock-fighter, who knows well that he can improve his breed by careful selection of the best cocks. How low in the scale of nature this law of battle descends I know not; male alligators have been described as fighting, bellowing, and whirling round, like Indians in a war dance, for the possession of the females; male salmons have been seen fighting all day long; male stag-beetles often bear wounds from the huge mandibles of other males. The war is, perhaps, severest between the males of polygamous animals, and these seem oftenest provided with special weapons. The males of carnivorous animals are already well armed; though to them and to others, special means of defence may be given through means of sexual selection, as the mane to the lion, the shoulder-pad to the boar, and the hooked jaw to the male salmon; for the shield may be as important for victory as the sword or spear.}{既然在家养状态下，特殊性状常常出现在某一性别中，并以遗传方式附着于这一性别，那么同样的事实大概也发生在自然状态下；如果是这样，自然选择就能够按照一性同另一性之间的功能关系来改变某一性，或者按照两性在生活习性上的完全不同来改变某一性，正如昆虫中有时所见的那样。这使我想谈几句我称为性选择的东西。性选择并不依赖于生存斗争，而依赖于雄性之间为占有雌性而进行的斗争；其结果并不是失败竞争者死亡，而是留下很少后代或没有后代。因此，性选择没有自然选择那样严酷。一般说来，最强健的雄性，即最适合其在自然中地位的雄性，会留下最多后裔。但在许多情况下，胜利并不取决于一般的强健，而取决于仅限于雄性的特殊武器。无角的雄鹿或无距的公鸡，留下后代的机会就很小。性选择总是让胜利者繁殖，当然能够赋予不可征服的勇气、增长鸡距，并增强翅膀击打带距腿部的力量；残酷的斗鸡者也深知，他能通过仔细选择最好的公鸡来改良自己的品种。我不知道这一战斗法则在自然阶梯中向下延伸到何等低级的地方；据描述，雄性短吻鳄会为了占有雌性而争斗、吼叫、旋转，像印第安人跳战舞一样；人们见过雄性鲑鱼整天争斗；雄性锹甲常带着其他雄性巨大上颚造成的伤痕。战争也许在多配偶动物的雄性之间最为激烈，而它们似乎也最常装备特殊武器。食肉动物的雄性本来已经武装良好；不过对它们以及其他动物来说，通过性选择仍可赋予特殊的防御手段，例如狮子的鬃毛、野猪的肩垫、雄鲑鱼的钩状颌；因为盾牌对于胜利可能同剑或矛一样重要。}

\pair{Amongst birds, the contest is often of a more peaceful character. All those who have attended to the subject believe that there is the severest rivalry between the males of many species to attract the females by singing. The rock-thrush of Guiana, birds of Paradise, and some others, congregate; and successive males display their gorgeous plumage and perform strange antics before the females, which, standing by as spectators, at last choose the most attractive partner. Those who have closely attended to birds in confinement well know that they often take individual preferences and dislikes: thus Sir R. Heron has described how one pied peacock was eminently attractive to all his hen birds. It may appear childish to attribute any effect to such apparently weak means: I cannot here enter on the details necessary to support this view; but if man can in a short time give elegant carriage and beauty to his bantams, according to his standard of beauty, I can see no good reason to doubt that female birds, by selecting during thousands of generations the most melodious or beautiful males according to their standard of beauty, might produce a marked effect. I strongly suspect that some well-known laws with respect to the plumage of male and female birds, in comparison with the plumage of the young, can be explained on the view of plumage having been chiefly modified by sexual selection, acting when the birds have come to the breeding age or during the breeding season; the modifications thus produced being inherited at corresponding ages or seasons, either by the males alone, or by the males and females; but I have not space here to enter on this subject.}{在鸟类中，争斗常常具有较为和平的性质。凡是注意过这一问题的人都相信，许多种的雄鸟之间存在最激烈的竞争，以歌声吸引雌鸟。圭亚那的岩鸫、极乐鸟以及其他一些鸟类会聚集起来；一只只雄鸟依次在雌鸟面前展示华美的羽衣，表演奇异的动作；雌鸟则站在一旁观看，最后选择最有吸引力的伴侣。仔细观察过笼养鸟类的人都很清楚，它们往往有个体性的偏好和厌恶；例如 R. Heron 爵士曾描述过，一只花斑孔雀怎样对他所有雌孔雀都具有格外的吸引力。把任何效果归因于这种看似微弱的手段，或许显得幼稚；我在这里不能详述支持这一观点所需的细节。但是，如果人能够在短时间内按照自己的美的标准，赋予矮脚鸡以优雅姿态和美丽，那么我看不出有什么充分理由怀疑，雌鸟按照自己的美的标准，在成千上万代中选择歌声最悦耳或最美丽的雄鸟，也可能产生显著效果。我强烈怀疑，关于雌雄鸟羽衣同幼鸟羽衣相比较的一些众所周知的规律，可以用这种观点来解释：羽衣主要是由性选择所改变，而性选择在鸟类达到繁殖年龄或处于繁殖季节时发生作用；由此产生的改变在相应年龄或季节遗传下来，或仅由雄性遗传，或由雌雄双方遗传。不过我在这里没有篇幅深入讨论这一问题。}

\pair{Thus it is, as I believe, that when the males and females of any animal have the same general habits of life, but differ in structure, colour, or ornament, such differences have been mainly caused by sexual selection; that is, individual males have had, in successive generations, some slight advantage over other males in their weapons, means of defence, or charms; and have transmitted these advantages to their male offspring. Yet I would not wish to attribute all such sexual differences to this agency: for we see peculiarities arising and becoming attached to the male sex in our domestic animals, as the wattle in male carriers, horn-like protuberances in the cocks of certain fowls, \&c., which we cannot believe to be either useful to the males in battle or attractive to the females. We see analogous cases under nature; for instance, the tuft of hair on the breast of the turkey-cock, which can hardly be either useful or ornamental to this bird. Indeed, had the tuft appeared under domestication, it would have been called a monstrosity.}{因此，我相信，当任何动物的雌雄具有相同的一般生活习性，却在构造、颜色或装饰上有所不同，这类差异主要是由性选择造成的；也就是说，在连续世代中，个别雄性在武器、防御手段或魅力方面比其他雄性稍有优势，并把这些优势传给了它们的雄性后代。不过，我并不愿意把所有这类性别差异都归因于这一作用；因为我们在家养动物中看到一些特征发生并附着于雄性，例如雄性邮鸽的肉垂、某些家鸡公鸡的角状突起等等，而我们不能相信这些特征对雄性战斗有用，或对雌性有吸引力。在自然状态下我们也看到类似情形；例如火鸡公鸡胸前的一簇毛，对这种鸟几乎既无实用也无装饰作用。事实上，如果这簇毛出现在家养状态下，它就会被称为畸形。}

\subsection*{Illustrations of the Action of Natural Selection\\自然选择作用的例证}

\pair{In order to make it clear how, as I believe, natural selection acts, I must beg permission to give one or two imaginary illustrations. Let us take the case of a wolf, which preys on various animals, securing some by craft, some by strength, and some by fleetness; and let us suppose that the fleetest prey, a deer for instance, had from any change in the country increased in numbers, or that other prey had decreased in numbers during that season of the year when the wolf is hardest pressed for food. I can under such circumstances see no reason to doubt that the swiftest and slimmest wolves would have the best chance of surviving, and so be preserved or selected, provided always that they retained strength to master their prey at this or at some other period of the year, when they might be compelled to prey on other animals. I can see no more reason to doubt this than that man can improve the fleetness of his greyhounds by careful and methodical selection, or by that unconscious selection which results from each man trying to keep the best dogs without any thought of modifying the breed.}{为了说明我所相信的自然选择是怎样发生作用的，请允许我举出一两个假想的例子。让我们以狼为例。狼捕食各种动物，有的靠狡诈取得，有的靠力量取得，有的靠速度取得；再设想最快捷的猎物，例如鹿，由于该地区发生某种变化而数量增加，或者其他猎物在一年中狼最缺食物的季节数量减少。在这种情况下，我看不出有什么理由怀疑，最快、最瘦长的狼会有最好的存活机会，并因此被保存或被选择；当然前提是它们仍保有力量，能在这一时期或一年中其他时期制服猎物，因为那时它们可能被迫捕食其他动物。我看不出这有什么理由比下列事实更值得怀疑：人能够通过仔细而有计划的选择改进灵缇犬的速度，或者通过无意识的选择来改进，即每个人都力求保留最好的犬，却并没有任何改造品种的念头。}

\pair{Even without any change in the proportional numbers of the animals on which our wolf preyed, a cub might be born with an innate tendency to pursue certain kinds of prey. Nor can this be thought very improbable; for we often observe great differences in the natural tendencies of our domestic animals: one cat, for instance, taking to catch rats, another mice; one cat, according to Mr. St. John, bringing home winged game, another hares or rabbits, and another hunting on marshy ground and almost nightly catching woodcocks or snipes. The tendency to catch rats rather than mice is known to be inherited. Now, if any slight innate change of habit or of structure benefited an individual wolf, it would have the best chance of surviving and of leaving offspring. Some of its young would probably inherit the same habits or structure, and by the repetition of this process a new variety might be formed which would either supplant or coexist with the parent-form of wolf. Or, again, the wolves inhabiting a mountainous district, and those frequenting the lowlands, would naturally be forced to hunt different prey; and from the continued preservation of the individuals best fitted for the two sites, two varieties might slowly be formed. These varieties would cross and blend where they met; but to this subject of intercrossing we shall soon have to return. I may add that, according to Mr. Pierce, there are two varieties of the wolf inhabiting the Catskill Mountains in the United States, one with a light greyhound-like form, which pursues deer, and the other more bulky, with shorter legs, which more frequently attacks the shepherd's flocks.}{即使我们所说的狼所捕食的动物在数量比例上没有任何变化，也可能有一只幼狼生来就具有追逐某些猎物的内在倾向。这也不能被认为很不可能；因为我们常在家养动物的自然倾向中观察到巨大差异。例如，一只猫喜欢捕大鼠，另一只喜欢捕小鼠；据 St. John 先生说，一只猫把有翼猎禽带回家，另一只带回野兔或家兔，还有一只在沼泽地捕猎，几乎每夜都捕到丘鹬或鹬。捕大鼠而不是小鼠的倾向，已知是会遗传的。现在，如果习性或构造上的任何轻微内在改变有利于某一只狼，它就会有最好的存活并留下后代的机会。它的一些幼仔很可能遗传同样的习性或构造；通过这一过程的反复，一种新品种可能形成，它或者取代狼的亲本类型，或者与其共存。再者，栖居山区的狼和常到低地活动的狼，自然会被迫猎取不同的猎物；持续保存最适合这两种地点的个体，便可能慢慢形成两个品种。这些品种在相遇处会杂交并融合；但关于杂交这个问题，我们很快还要再谈。我还可以补充说，据 Pierce 先生所述，美国卡茨基尔山中有两个狼的品种：一种体形轻捷，像灵缇犬，追逐鹿；另一种较粗壮，腿较短，更常袭击牧人的羊群。}

\pair{Let us now take a more complex case. Certain plants excrete a sweet juice, apparently for the sake of eliminating something injurious from their sap: this is effected by glands at the base of the stipules in some Leguminosae, and at the back of the leaf of the common laurel. This juice, though small in quantity, is greedily sought by insects. Let us now suppose a little sweet juice or nectar to be excreted by the inner bases of the petals of a flower. In this case insects in seeking the nectar would get dusted with pollen, and would certainly often transport the pollen from one flower to the stigma of another flower. The flowers of two distinct individuals of the same species would thus get crossed; and the act of crossing, we have good reason to believe, as will hereafter be more fully alluded to, would produce very vigorous seedlings, which consequently would have the best chance of flourishing and surviving. Some of these seedlings would probably inherit the nectar-excreting power. Those individual flowers which had the largest glands or nectaries, and which excreted most nectar, would be oftenest visited by insects and would be oftenest crossed; and so in the long run would gain the upper hand. Those flowers also which had their stamens and pistils placed, in relation to the size and habits of the particular insects which visited them, so as to favour in any degree the transportal of their pollen from flower to flower, would likewise be favoured or selected. We might have taken the case of insects visiting flowers for the sake of collecting pollen instead of nectar; and as pollen is formed for the sole object of fertilisation, its destruction appears a simple loss to the plant; yet if a little pollen were carried, at first occasionally and then habitually, by the pollen-devouring insects from flower to flower, and a cross thus effected, although nine-tenths of the pollen were destroyed, it might still be a great gain to the plant; and those individuals which produced more and more pollen, and had larger and larger anthers, would be selected.}{现在让我们举一个更复杂的例子。某些植物分泌甜汁，显然是为了从其汁液中排除某种有害物；在一些豆科植物中，这是由托叶基部的腺体完成的，在普通月桂中则由叶背完成。这种汁液数量虽少，却受到昆虫贪婪寻求。现在设想一朵花的花瓣内基部分泌少量甜汁或花蜜。在这种情况下，昆虫寻求花蜜时会沾上花粉，并且必定常把花粉从一朵花带到另一朵花的柱头上。这样，同种两个不同个体的花就会发生杂交；而杂交这一行为，我们有充分理由相信，正如以后还会更详细提到的，会产生非常强健的幼苗，因而这些幼苗最有机会繁盛并存活。其中一些幼苗大概会遗传分泌花蜜的能力。那些腺体或蜜腺最大的、分泌花蜜最多的个体花，会最常受到昆虫访问，也会最常发生杂交；因此从长远来看会占上风。那些雄蕊和雌蕊的位置，同访问它们的特定昆虫的大小和习性相适合，从而在任何程度上有利于花粉在花与花之间转移的花，也同样会受到有利对待或被选择。我们本来也可以设想昆虫不是为了花蜜而是为了采集花粉来访花；花粉形成的唯一目的在于受精，因此它被破坏似乎只是植物的损失；然而，如果食花粉的昆虫先是偶尔、后来习惯性地把少量花粉从一朵花带到另一朵花，从而实现杂交，那么即使十分之九的花粉被破坏，对植物仍可能是巨大收益；那些产生越来越多花粉、花药越来越大的个体，便会被选择。}

\pair{When our plant, by this process of the continued preservation or natural selection of more and more attractive flowers, had been rendered highly attractive to insects, they would, unintentionally on their part, regularly carry pollen from flower to flower; and that they can most effectually do this I could easily show by many striking instances. I will give only one, not as a very striking case, but as likewise illustrating one step in the separation of the sexes of plants, presently to be alluded to. Some holly-trees bear only male flowers, which have four stamens producing rather a small quantity of pollen, and a rudimentary pistil; other holly-trees bear only female flowers; these have a full-sized pistil, and four stamens with shrivelled anthers, in which not a grain of pollen can be detected. Having found a female tree exactly sixty yards from a male tree, I put the stigmas of twenty flowers, taken from different branches, under the microscope, and on all, without exception, there were pollen grains, and on some a profusion of pollen. As the wind had set for several days from the female to the male tree, the pollen could not thus have been carried. The weather had been cold and boisterous, and therefore not favourable to bees; nevertheless every female flower which I examined had been effectually fertilised by the bees, accidentally dusted with pollen, having flown from tree to tree in search of nectar. But to return to our imaginary case: as soon as the plant had been rendered so highly attractive to insects that pollen was regularly carried from flower to flower, another process might commence. No naturalist doubts the advantage of what has been called the “physiological division of labour”; hence we may believe that it would be advantageous to a plant to produce stamens alone in one flower or on one whole plant, and pistils alone in another flower or on another plant. In plants under culture and placed under new conditions of life, sometimes the male organs and sometimes the female organs become more or less impotent; now if we suppose this to occur in ever so slight a degree under nature, then, as pollen is already carried regularly from flower to flower, and as a more complete separation of the sexes of our plant would be advantageous on the principle of the division of labour, individuals with this tendency more and more increased would be continually favoured or selected, until at last a complete separation of the sexes would be effected.}{当我们的植物通过持续保存或自然选择越来越能吸引昆虫的花这一过程，已经变得对昆虫具有很强吸引力时，昆虫就会无意中经常把花粉从一朵花带到另一朵花；它们能够极其有效地做到这一点，我本可以轻易举出许多醒目的例子。我只举一个例子，不是因为它特别醒目，而是因为它也说明了植物性别分离的一个步骤，这一点马上还要提到。有些冬青树只开雄花，花中有四枚雄蕊，产生相当少量的花粉，并有一个退化的雌蕊；另一些冬青树只开雌花，具有发育完全的雌蕊和四枚花药萎缩的雄蕊，在其中检测不到一粒花粉。我发现一株雌树距离一株雄树正好六十码，便把取自不同枝条的二十朵花的柱头放在显微镜下观察；所有柱头无一例外都有花粉粒，有些还沾有大量花粉。由于风已经连续几天从雌树吹向雄树，花粉不可能由风这样带来。天气寒冷而多风，因此并不适合蜜蜂活动；然而我检查的每一朵雌花，都已由蜜蜂有效受精，蜜蜂在寻找花蜜、从树飞到树时，偶然沾上了花粉。但回到我们的假想例子：植物一旦变得如此吸引昆虫，以致花粉经常被从花带到花，另一个过程便可能开始。没有自然学家怀疑所谓“生理分工”的好处；因此我们可以相信，一株植物在一朵花或整株植物上只产生雄蕊，而在另一朵花或另一株植物上只产生雌蕊，对它是有利的。在栽培并置于新生活条件下的植物中，有时雄性器官、有时雌性器官会或多或少失去功能；现在，如果我们假定这种情形在自然状态下哪怕以极轻微程度发生，那么既然花粉已经经常从花传到花，而且根据分工原则，我们的植物性别更完全的分离会是有利的，具有这一倾向并逐渐增强的个体就会不断受到有利对待或被选择，直到最后实现性别的完全分离。}

\pair{Let us now turn to the nectar-feeding insects in our imaginary case: we may suppose the plant of which we have been slowly increasing the nectar by continued selection to be a common plant, and that certain insects depended in main part on its nectar for food. I could give many facts showing how anxious bees are to save time; for instance, their habit of cutting holes and sucking the nectar at the bases of certain flowers, which they can, with a very little more trouble, enter by the mouth. Bearing such facts in mind, I can see no reason to doubt that an accidental deviation in the size and form of the body, or in the curvature and length of the proboscis, \&c., far too slight to be appreciated by us, might profit a bee or other insect, so that an individual so characterised would be able to obtain its food more quickly, and so have a better chance of living and leaving descendants. Its descendants would probably inherit a tendency to a similar slight deviation of structure. The tubes of the corollas of the common red and incarnate clovers, Trifolium pratense and incarnatum, do not on a hasty glance appear to differ in length; yet the hive-bee can easily suck the nectar out of the incarnate clover, but not out of the common red clover, which is visited by humble-bees alone; so that whole fields of the red clover offer in vain an abundant supply of precious nectar to the hive-bee. Thus it might be a great advantage to the hive-bee to have a slightly longer or differently constructed proboscis. On the other hand, I have found by experiment that the fertility of clover greatly depends on bees visiting and moving parts of the corolla so as to push the pollen on to the stigmatic surface. Hence, again, if humble-bees were to become rare in any country, it might be a great advantage to the red clover to have a shorter or more deeply divided tube to its corolla, so that the hive-bee could visit its flowers. Thus I can understand how a flower and a bee might slowly become, either simultaneously or one after the other, modified and adapted in the most perfect manner to each other by the continued preservation of individuals presenting mutual and slightly favourable deviations of structure.}{现在让我们转向假想例子中以花蜜为食的昆虫：我们可以设想，那种花蜜通过连续选择而被我们逐渐增多的植物，是一种普通植物；某些昆虫主要依赖它的花蜜为食。我可以举出许多事实，说明蜜蜂多么急于节省时间；例如，它们习惯在某些花的基部咬洞吸取花蜜，而这些花其实只要多费一点点力气就能从花口进入。牢记这些事实，我看不出有什么理由怀疑，在身体大小和形状上，或在吻的弯曲度和长度等方面，偶然出现的、细微到我们难以察觉的偏差，可能有利于一只蜜蜂或其他昆虫，使具有这种性状的个体能够更快获得食物，从而有更好的存活和留下后代的机会。它的后裔很可能遗传类似轻微构造偏差的倾向。普通红三叶草和绛三叶草（Trifolium pratense 和 incarnatum）的花冠管，匆匆一看似乎长度没有差别；然而家蜂能够容易地吸取绛三叶草的花蜜，却不能吸取普通红三叶草的花蜜，后者只由熊蜂访问；因此整片红三叶草田虽向家蜂提供大量珍贵花蜜，却是徒然的。这样，家蜂具有稍长的或构造略有不同的吻，可能就是一项巨大优势。另一方面，我通过实验证明，三叶草的结实率在很大程度上依赖蜜蜂访问并移动花冠的某些部分，从而把花粉推到柱头表面。因此，如果熊蜂在某个国家变得稀少，红三叶草具有更短或分裂更深的花冠管，使家蜂能够访问它的花，就可能是一项巨大优势。由此我能够理解，通过持续保存那些呈现相互且稍微有利的构造偏差的个体，一朵花和一只蜜蜂怎样会缓慢地、或同时或先后地被改变并以最完善的方式相互适应。}

\pair{I am well aware that this doctrine of natural selection, exemplified in the above imaginary instances, is open to the same objections which were at first urged against Sir Charles Lyell's noble views on “the modern changes of the earth, as illustrative of geology”; but we now very seldom hear the action, for instance, of the coast waves called a trifling and insignificant cause, when applied to the excavation of gigantic valleys or to the formation of the longest lines of inland cliffs. Natural selection can act only by the preservation and accumulation of infinitesimally small inherited modifications, each profitable to the preserved being; and as modern geology has almost banished such views as the excavation of a great valley by a single diluvial wave, so will natural selection, if it be a true principle, banish the belief in the continued creation of new organic beings, or in any great and sudden modification in their structure.}{我很清楚，上述假想例子所说明的自然选择学说，容易遭到最初针对查尔斯·莱尔爵士关于“用地球的现代变化说明地质”的高明见解所提出的同类反对；但是现在，当把海岸波浪的作用用于解释巨大山谷的开凿，或最长内陆峭壁线的形成时，我们已经很少再听见有人称这类作用为微不足道、无足轻重的原因。自然选择只能通过保存并累积极其微小的遗传改变而发生作用，而每一种改变都对被保存的生物有利；正如现代地质学几乎已经驱逐了由一次洪水波开凿巨大山谷这类观点一样，如果自然选择是一个真实原则，它也将驱逐那种相信新生物不断被创造出来、或其构造发生任何巨大而突然改变的信念。}

\begin{center}
[\ldots]
\end{center}

\subsection*{Extinction\\灭绝}

\pair{This subject will be more fully discussed in our chapter on Geology; but it must be here alluded to from being intimately connected with natural selection. Natural selection acts solely through the preservation of variations in some way advantageous, which consequently endure. But as, from the high geometrical powers of increase of all organic beings, each area is already fully stocked with inhabitants, it follows that as each selected and favoured form increases in number, so will the less favoured forms decrease and become rare. Rarity, as geology tells us, is the precursor to extinction. We can also see that any form represented by few individuals will, during fluctuations in the seasons or in the number of its enemies, run a good chance of utter extinction. But we may go further than this; for as new forms are continually and slowly being produced, unless we believe that the number of specific forms goes on perpetually and almost indefinitely increasing, numbers inevitably must become extinct. That the number of specific forms has not indefinitely increased, geology shows us plainly; and indeed we can see reason why they should not have thus increased, for the number of places in the polity of nature is not indefinitely great, not that we have any means of knowing that any one region has as yet got its maximum of species. Probably no region is as yet fully stocked, for at the Cape of Good Hope, where more species of plants are crowded together than in any other quarter of the world, some foreign plants have become naturalised without causing, as far as we know, the extinction of any natives.}{这一问题将在地质学一章中作更充分的讨论；但由于它同自然选择关系密切，这里必须提到。自然选择只通过保存某些方面有利的变异而发生作用，而这些变异因此得以持续。可是，由于一切生物都具有高度的几何级数增长能力，每一个区域已经住满了居民；于是，随着每一种被选择和受有利对待的类型在数量上增加，较少受到有利对待的类型就会减少并变得稀有。地质学告诉我们，稀有是灭绝的先导。我们还可以看到，任何只由少数个体代表的类型，在季节波动或敌害数量波动期间，都很有可能完全灭绝。但我们还可以更进一步；因为新的类型不断而缓慢地产生，除非我们相信种的类型数会永远并几乎无限地增加，否则必定有许多类型灭绝。地质学清楚地向我们表明，种的类型数并没有无限增加；事实上，我们也可以看出它们为什么不应这样增加，因为自然体系中的位置数量并不是无限大的。不过，我们并没有办法知道任何一个地区是否已经达到其种数的最大值。大概没有任何地区已经完全住满；因为在好望角，那里挤聚在一起的植物种比世界任何其他地方都多，而一些外来植物已经归化，据我们所知，并没有造成任何本地植物灭绝。}

\pair{Furthermore, the species which are most numerous in individuals will have the best chance of producing within any given period favourable variations. We have evidence of this in the facts given in the second chapter, showing that it is the common species which afford the greatest number of recorded varieties, or incipient species. Hence rare species will be less quickly modified or improved within any given period, and they will consequently be beaten in the race for life by the modified descendants of the commoner species.}{此外，个体数量最多的种，在任何给定时期内最有机会产生有利变异。第二章所列举的事实为此提供了证据，表明正是普通种提供了最多有记录的品种，或初成种。因此，稀有种在任何给定时期内被改变或改进的速度都会较慢，于是它们将在生命竞赛中被较普通种的已改变后裔所击败。}

\pair{From these several considerations I think it inevitably follows that, as new species in the course of time are formed through natural selection, others will become rarer and rarer, and finally extinct. The forms which stand in closest competition with those undergoing modification and improvement will naturally suffer most. And we have seen in the chapter on the Struggle for Existence that it is the most closely allied forms—varieties of the same species, and species of the same genus or of related genera—which, from having nearly the same structure, constitution, and habits, generally come into the severest competition with each other. Consequently, each new variety or species, during the progress of its formation, will generally press hardest on its nearest kindred, and tend to exterminate them. We see the same process of extermination amongst our domesticated productions, through the selection of improved forms by man. Many curious instances could be given showing how quickly new breeds of cattle, sheep, and other animals, and varieties of flowers, take the place of older and inferior kinds. In Yorkshire, it is historically known that the ancient black cattle were displaced by the long-horns, and that these “were swept away by the short-horns” (I quote the words of an agricultural writer) “as if by some murderous pestilence.”}{根据这几方面的考虑，我认为必然可以得出这样的结论：随着新的种在时间进程中通过自然选择形成，其他种会变得越来越稀有，最后灭绝。那些同正在改变和改进的类型处于最密切竞争中的类型，自然受害最重。我们在关于生存斗争的一章中已经看到，正是关系最接近的类型，即同一种的品种、同属的种或相关属的种，因为具有几乎相同的构造、体质和习性，通常彼此进入最激烈的竞争。因此，每一个新品种或新种在形成过程中，通常会最严厉地压迫其最亲近的亲族，并倾向于消灭它们。通过人对改良类型的选择，我们在家养生物中也看到同样的消灭过程。可以举出许多有趣例子，说明牛、羊和其他动物的新品种，以及花卉品种，怎样迅速取代较老而较劣的种类。在约克郡，历史上知道古老的黑牛曾被长角牛取代，而长角牛又“被短角牛扫荡一空”（我引用一位农业作者的话），“仿佛遭遇某种杀戮性的瘟疫”。}

\subsection*{Divergence of Character\\性状分歧}

\pair{The principle which I have designated by this term is of high importance on my theory, and explains, as I believe, several important facts. In the first place, varieties, even strongly marked ones, though having somewhat of the character of species—as is shown by the hopeless doubts in many cases how to rank them—yet certainly differ from each other far less than do good and distinct species. Nevertheless, according to my view, varieties are species in the process of formation, or are, as I have called them, incipient species. How, then, does the lesser difference between varieties become augmented into the greater difference between species? That this does habitually happen, we must infer from most of the innumerable species throughout nature presenting well-marked differences; whereas varieties, the supposed prototypes and parents of future well-marked species, present slight and ill-defined differences. Mere chance, as we may call it, might cause one variety to differ in some character from its parents, and the offspring of this variety again to differ from its parent in the very same character and in a greater degree; but this alone would never account for so habitual and large an amount of difference as that between varieties of the same species and species of the same genus.}{我用这一术语所指称的原则，在我的理论中具有高度重要性；我相信，它解释了若干重要事实。首先，品种，即使是特征很强的品种，虽然多少具有种的性状，正如许多情况下怎样给它们定级令人无望地疑惑所显示的那样，但它们彼此之间的差异，肯定远小于良好而明确的种之间的差异。然而，按照我的观点，品种就是正在形成中的种，或者如我所称，是初成种。那么，品种之间较小的差异，怎样会增大为种之间较大的差异呢？我们必须从自然界无数种中的大多数都呈现明显差异这一事实推断，这种情况是经常发生的；而品种，即未来明显种的假定原型和亲本，却呈现轻微而界限不明的差异。我们所谓的单纯偶然，可能使一个品种在某一性状上不同于其亲本，而该品种的后代又在同一性状上以更大程度不同于其亲本；但仅凭这一点，绝不能解释同一种的品种之间和同属的种之间那种如此经常且巨大的差异量。}

\pair{As has always been my practice, let us seek light on this head from our domestic productions. We shall here find something analogous. A fancier is struck by a pigeon having a slightly shorter beak; another fancier is struck by a pigeon having a rather longer beak; and on the acknowledged principle that “fanciers do not and will not admire a medium standard, but like extremes,” they both go on, as has actually occurred with tumbler-pigeons, choosing and breeding from birds with longer and longer beaks, or with shorter and shorter beaks. Again, we may suppose that at an early period one man preferred swifter horses, another stronger and more bulky horses. The early differences would be very slight; in the course of time, from the continued selection of swifter horses by some breeders, and of stronger ones by others, the differences would become greater, and would be noted as forming two sub-breeds; finally, after the lapse of centuries, the sub-breeds would become converted into two well-established and distinct breeds. As the differences slowly become greater, the inferior animals with intermediate characters, being neither very swift nor very strong, will have been neglected and will have tended to disappear. Here, then, we see in man's productions the action of what may be called the principle of divergence, causing differences, at first barely appreciable, steadily to increase, and the breeds to diverge in character both from each other and from their common parent.}{像我一贯的做法一样，让我们从家养生物中寻求这一问题的线索。我们将在这里发现类似情形。一个爱好者被喙稍短的鸽子打动；另一个爱好者被喙较长的鸽子打动；依据公认的原则，即“爱好者并不欣赏、也不会欣赏中等标准，而是喜欢极端”，他们两人都继续选择并繁殖喙越来越长或越来越短的鸟，翻飞鸽中实际已经发生过这种情形。再者，我们可以设想，在早期，一个人偏爱更快的马，另一个人偏爱更强壮、更粗大的马。最初的差异会非常轻微；随着时间推移，由于一些育种者持续选择更快的马，另一些持续选择更强壮的马，差异会变得更大，并会被注意到形成两个亚品种；最后，经过数世纪，亚品种会转化为两个确立而不同的品种。随着差异慢慢增大，那些具有中间性状、既不很快也不很强壮的劣等动物，就会被忽视，并倾向于消失。因此，在人的产物中，我们看到了可以称为分歧原则的作用：起初几乎无法察觉的差异，稳步增加，品种在性状上既彼此分歧，也同它们的共同亲本分歧。}

\pair{But how, it may be asked, can any analogous principle apply in nature? I believe it can and does apply most efficiently, from the simple circumstance that the more diversified the descendants from any one species become in structure, constitution, and habits, by so much will they be better enabled to seize on many and widely diversified places in the polity of nature, and so be enabled to increase in numbers.}{但是，或许有人会问，任何类似的原则怎样能应用于自然界呢？我相信它能够，而且确实极其有效地应用于自然界，理由只是这样一个简单情形：任何一个种的后裔在构造、体质和习性上越是多样化，就越能占据自然体系中许多广泛多样的位置，从而能够增加数量。}

\pair{We can clearly see this in the case of animals with simple habits. Take the case of a carnivorous quadruped, of which the number that can be supported in any country has long ago arrived at its full average. If its natural powers of increase be allowed to act, it can succeed in increasing, the country not undergoing any change in its conditions, only by its varying descendants seizing on places at present occupied by other animals: some of them, for instance, being enabled to feed on new kinds of prey, either dead or alive; some inhabiting new stations, climbing trees, frequenting water, and some perhaps becoming less carnivorous. The more diversified in habits and structure the descendants of our carnivorous animal became, the more places they would be enabled to occupy. What applies to one animal will apply throughout all time to all animals, that is, if they vary; for otherwise natural selection can do nothing. So it will be with plants. It has been experimentally proved that if a plot of ground be sown with one species of grass, and a similar plot be sown with several distinct genera of grasses, a greater number of plants and a greater weight of dry herbage can thus be raised. The same has been found to hold good when first one variety and then several mixed varieties of wheat have been sown on equal spaces of ground. Hence, if any one species of grass were to go on varying, and those varieties were continually selected which differed from each other in at all the same manner as distinct species and genera of grasses differ from each other, a greater number of individual plants of this species of grass, including its modified descendants, would succeed in living on the same piece of ground. And we well know that each species and each variety of grass is annually sowing almost countless seeds; and thus, as it may be said, is striving its utmost to increase its numbers. Consequently, I cannot doubt that in the course of many thousands of generations, the most distinct varieties of any one species of grass would always have the best chance of succeeding and of increasing in numbers, and thus of supplanting the less distinct varieties; and varieties, when rendered very distinct from each other, take the rank of species.}{在习性简单的动物中，我们可以清楚看到这一点。以一种食肉四足动物为例，任何国家所能供养的这种动物数量，早已达到其完全平均值。如果允许它的自然增长能力发生作用，而该国家的条件又没有发生任何变化，那么它只有通过其发生变异的后裔占据目前由其他动物占据的位置，才能成功增加数量：例如，其中一些能够以新的猎物为食，无论是死的还是活的；一些居住于新的场所，爬树，常到水边；还有一些也许变得不那么食肉。我们的食肉动物的后裔在习性和构造上越是多样化，就越能占据更多位置。适用于一种动物的情形，也会在一切时期适用于一切动物，也就是说，只要它们发生变异；否则自然选择便无能为力。植物也会如此。实验证明，如果一块地播种一种草，而另一块类似的地播种几个不同属的草，那么后一种方式能够产生更多植株和更大重量的干草。先在同等面积上播种一种小麦品种，再播种若干混合品种时，也发现同样规律成立。因此，如果任何一种草继续变异，而那些彼此之间的差异多少像不同种和不同属的草彼此差异那样的品种不断被选择，那么这种草的更多个体植株，包括它的已改变后裔，就会成功地生活在同一块土地上。我们也很清楚，每一种和每一个品种的草每年都播下几乎不可计数的种子；也可以说，它们正尽最大努力增加自己的数量。因此，我毫不怀疑，在许多千代的过程中，任何一种草中最为不同的品种，总会最有成功并增加数量的机会，从而取代较不相同的品种；而品种一旦彼此变得十分不同，便取得种的地位。}

\pair{The truth of the principle, that the greatest amount of life can be supported by great diversification of structure, is seen under many natural circumstances. In an extremely small area, especially if freely open to immigration, and where the contest between individual and individual must be severe, we always find great diversity in its inhabitants. For instance, I found that a piece of turf, three feet by four in size, which had been exposed for many years to exactly the same conditions, supported twenty species of plants, and these belonged to eighteen genera and to eight orders, which shows how much these plants differed from each other. So it is with the plants and insects on small and uniform islets; and so in small ponds of fresh water. Farmers find that they can raise most food by a rotation of plants belonging to the most different orders: nature follows what may be called a simultaneous rotation. Most of the animals and plants which live close round any small piece of ground could live on it, supposing it not to be in any way peculiar in its nature, and may be said to be striving to the utmost to live there; but it is seen that where they come into the closest competition with each other, the advantages of diversification of structure, with the accompanying differences of habit and constitution, determine that the inhabitants which thus jostle each other most closely shall, as a general rule, belong to what we call different genera and orders.}{构造高度多样化能够支持最大生命量这一原则的真实性，在许多自然情形中都可以看到。在极小的区域内，尤其是当它自由向迁入开放、个体与个体之间的竞争必定激烈时，我们总会发现其居民具有很大多样性。例如，我发现一块三英尺乘四英尺、许多年来暴露于完全相同条件下的草皮，供养着二十种植物，而这些植物分属十八个属、八个目，这表明这些植物彼此差异多么大。小而均一的岛屿上的植物和昆虫如此，小淡水池中的情形也如此。农民发现，通过轮作最不同目的植物，他们能够生产最多食物；自然遵循着一种可以称为同时轮作的方式。生活在任何一小块土地周围的大多数动植物，假定这块土地在性质上没有任何特殊之处，本都能生活在其上，并且可以说都在尽力争取生活在那里；但是可以看到，在它们彼此进入最密切竞争的地方，构造多样化所带来的优势以及随之而来的习性和体质差异，决定了这些最紧密相互挤压的居民，通常属于我们所谓的不同属和不同目。}

\begin{center}
[\ldots]
\end{center}

\pair{After the foregoing discussion, which ought to have been much amplified, we may, I think, assume that the modified descendants of any one species will succeed by so much the better as they become more diversified in structure, and are thus enabled to encroach on places occupied by other beings. Now let us see how this principle of great benefit being derived from divergence of character, combined with the principles of natural selection and of extinction, will tend to act.}{经过前面的讨论，尽管本应大为扩展，我想我们可以假定：任何一个种的已改变后裔，在构造上越是多样化，从而越能侵占其他生物所占据的位置，就越能更好地成功。现在让我们看看，从性状分歧中获得巨大好处这一原则，同自然选择和灭绝原则结合起来时，将会怎样发生作用。}

\pair{The accompanying diagram will aid us in understanding this rather perplexing subject. Let A to L represent the species of a genus large in its own country; these species are supposed to resemble each other in unequal degrees, as is so generally the case in nature, and as is represented in the diagram by the letters standing at unequal distances. I have said a large genus, because we have seen in the second chapter that, on an average, more of the species of large genera vary than of small genera; and the varying species of the large genera present a greater number of varieties. We have also seen that the species which are the commonest and the most widely diffused vary more than rare species with restricted ranges. Let (A) be a common, widely diffused, and varying species belonging to a genus large in its own country. The little fan of diverging dotted lines of unequal lengths proceeding from (A) may represent its varying offspring. The variations are supposed to be extremely slight, but of the most diversified nature; they are not supposed all to appear simultaneously, but often after long intervals of time; nor are they all supposed to endure for equal periods. Only those variations which are in some way profitable will be preserved or naturally selected. And here the importance of the principle of benefit being derived from divergence of character comes in; for this will generally lead to the most different or divergent variations, represented by the outer dotted lines, being preserved and accumulated by natural selection. When a dotted line reaches one of the horizontal lines, and is there marked by a small numbered letter, a sufficient amount of variation is supposed to have been accumulated to have formed a fairly well-marked variety, such as would be thought worthy of record in a systematic work.}{附图将帮助我们理解这个相当令人困惑的问题。让 A 到 L 代表某一属的种，这个属在其本国是一个大属；这些种彼此相似的程度被假定为不相等，这在自然界极为普遍，图中字母间距离不等也表示了这一点。我说的是一个大属，因为我们在第二章已经看到，平均说来，大属中发生变异的种比小属更多；而大属中发生变异的种，呈现出更多品种。我们也已经看到，最普通、分布最广的种，比稀有且分布范围受限的种变异更多。设 (A) 是一个普通的、分布广泛的、发生变异的种，属于其本国的一个大属。从 (A) 发出的那一小束长短不等、分歧的虚线，可以代表它发生变异的后代。这里假定变异极其轻微，但性质最为多样；并不假定它们同时出现，而是常在很长间隔之后出现；也不假定它们都持续同等时期。只有那些以某种方式有益的变异会被保存，或被自然选择。这里，性状分歧带来利益这一原则的重要性便显现出来；因为这一原则通常会导致最不同或最分歧的变异，即由外侧虚线表示的变异，被自然选择保存并累积。当一条虚线到达一条水平线，并在该处用一个带小数字的字母标出时，就假定已经累积了足够数量的变异，形成了相当明显的品种，足以被认为值得在分类学著作中记录。}

\begin{figure}[htbp]
\centering
\IfFileExists{../images/bbe4172e667b7677be63c8bcf78a3ef3c8b0cb1900750fa6e01b6f2f0383ecbc.jpg}{%
  \safeincludeimage{bbe4172e667b7677be63c8bcf78a3ef3c8b0cb1900750fa6e01b6f2f0383ecbc.jpg}%
}{%
  \IfFileExists{images/bbe4172e667b7677be63c8bcf78a3ef3c8b0cb1900750fa6e01b6f2f0383ecbc.jpg}{%
    \safeincludeimage{bbe4172e667b7677be63c8bcf78a3ef3c8b0cb1900750fa6e01b6f2f0383ecbc.jpg}%
  }{%
    \fbox{\parbox{0.82\linewidth}{\centering Missing image reference preserved:\\\texttt{images/bbe4172e667b7677be63c8bcf78a3ef3c8b0cb1900750fa6e01b6f2f0383ecbc.jpg}}}%
  }%
}
\caption{达尔文用于说明性状分歧、自然选择与灭绝的谱系图。}
\end{figure}

\pair{The intervals between the horizontal lines in the diagram may represent each a thousand generations; but it would have been better if each had represented ten thousand generations. After a thousand generations, species (A) is supposed to have produced two fairly well-marked varieties, namely $a^{1}$ and $m^{1}$. These two varieties will generally continue to be exposed to the same conditions which made their parents variable, and the tendency to variability is in itself hereditary; consequently they will tend to vary, and generally to vary in nearly the same manner as their parents varied. Moreover, these two varieties, being only slightly modified forms, will tend to inherit those advantages which made their common parent (A) more numerous than most of the other inhabitants of the same country; they will likewise partake of those more general advantages which made the genus to which the parent species belonged a large genus in its own country. And these circumstances we know to be favourable to the production of new varieties.}{图中水平线之间的间隔，每一段可以代表一千代；不过如果每一段代表一万代会更好。一千代以后，假定种 (A) 已经产生了两个相当明显的品种，即 $a^{1}$ 和 $m^{1}$。这两个品种通常会继续暴露于使其亲本发生变异的同样条件之下，而变异倾向本身是可遗传的；因此它们会倾向于变异，而且通常会以几乎同其亲本相同的方式变异。此外，这两个品种只是稍有改变的类型，因而会倾向于遗传那些使其共同亲本 (A) 比同一国家大多数其他居民数量更多的优势；它们也会分享那些更一般的优势，正是这些优势使亲本种所属的属在其本国成为大属。我们知道，这些情形有利于新品种的产生。}

\pair{If, then, these two varieties be variable, the most divergent of their variations will generally be preserved during the next thousand generations. And after this interval, variety $a^{1}$ is supposed in the diagram to have produced variety $a^{2}$, which will, owing to the principle of divergence, differ more from (A) than did variety $a^{1}$. Variety $m^{1}$ is supposed to have produced two varieties, namely $m^{2}$ and $s^{2}$, differing from each other, and more considerably from their common parent (A). We may continue the process by similar steps for any length of time; some of the varieties, after each thousand generations, producing only a single variety but in a more and more modified condition, some producing two or three varieties, and some failing to produce any. Thus the varieties or modified descendants proceeding from the common parent (A) will generally go on increasing in number and diverging in character. In the diagram the process is represented up to the ten-thousandth generation, and under a condensed and simplified form up to the fourteen-thousandth generation.}{那么，如果这两个品种是可变的，在接下来的一千代中，它们的变异中最分歧的那些通常会被保存下来。在这段间隔之后，图中假定品种 $a^{1}$ 产生了品种 $a^{2}$；由于分歧原则，$a^{2}$ 将比品种 $a^{1}$ 更不同于 (A)。假定品种 $m^{1}$ 产生了两个品种，即 $m^{2}$ 和 $s^{2}$，它们彼此不同，并且同其共同亲本 (A) 的差异更大。我们可以按类似步骤把这一过程继续任意长时间；有些品种在每一千代之后只产生一个品种，但处于越来越改变的状态；有些产生两个或三个品种；有些则没有产生任何品种。因此，出自共同亲本 (A) 的品种或已改变后裔，通常会在数量上继续增加，并在性状上继续分歧。图中这一过程表示到第一万代，并以压缩和简化形式表示到第一万四千代。}

\pair{But I must here remark that I do not suppose that the process ever goes on so regularly as is represented in the diagram, though in itself made somewhat irregular. I am far from thinking that the most divergent varieties will invariably prevail and multiply: a medium form may often long endure, and may or may not produce more than one modified descendant; for natural selection will always act according to the nature of the places which are either unoccupied or not perfectly occupied by other beings, and this will depend on infinitely complex relations. But as a general rule, the more diversified in structure the descendants from any one species can be rendered, the more places they will be enabled to seize on, and the more their modified progeny will be increased. In our diagram the line of succession is broken at regular intervals by small numbered letters marking the successive forms which have become sufficiently distinct to be recorded as varieties. But these breaks are imaginary, and might have been inserted anywhere, after intervals long enough to have allowed the accumulation of a considerable amount of divergent variation.}{不过我在这里必须说明，我并不假定这一过程会像图中表示的那样规则，尽管图本身已经画得有些不规则。我完全不认为最分歧的品种总会不变地占优势并增多：一种中间类型常常可以长期持续，并且可能产生也可能不产生一个以上的已改变后裔；因为自然选择总会按照那些尚未被其他生物占据或尚未被完全占据的位置的性质而发生作用，而这又取决于无限复杂的关系。但作为一般规则，任何一个种的后裔在构造上能够被造成得越多样化，它们就越能占据更多位置，它们的已改变子代也就会增加得越多。在我们的图中，继承线在规则间隔处被带小数字的字母打断，用以标示那些已经变得足够不同、可作为品种记录的连续类型。但是这些中断是假想的，本可以在任何地方插入，只要其间隔足够长，足以允许相当数量的分歧变异累积起来。}

\pair{As all the modified descendants from a common and widely diffused species, belonging to a large genus, will tend to partake of the same advantages which made their parent successful in life, they will generally go on multiplying in number as well as diverging in character: this is represented in the diagram by the several divergent branches proceeding from (A). The modified offspring from the later and more highly improved branches in the lines of descent will, it is probable, often take the place of, and so destroy, the earlier and less improved branches: this is represented in the diagram by some of the lower branches not reaching to the upper horizontal lines. In some cases I do not doubt that the process of modification will be confined to a single line of descent, and the number of the descendants will not be increased; although the amount of divergent modification may have been increased in the successive generations. This case would be represented in the diagram if all the lines proceeding from (A) were removed, excepting that from $a^{1}$ to $a^{10}$. In the same way, for instance, the English race-horse and English pointer have apparently both gone on slowly diverging in character from their original stocks, without either having given off any fresh branches or races.}{由于出自一个普通而分布广泛、又属于大属的种的一切已改变后裔，都会倾向于分享使其亲本在生活中成功的同样优势，它们通常会在数量上继续增多，同时在性状上继续分歧；图中由 (A) 发出的几条分歧枝表示了这一点。血统线上较晚且改进程度更高的分枝所产生的已改变后代，很可能常常取代并因此消灭较早且改进较少的分枝；图中一些较低分枝没有达到较高水平线，便表示了这一点。在某些情况下，我毫不怀疑改变过程会限于单一血统线，后裔的数量不会增加；不过在连续世代中，分歧改变的数量可能已经增加。如果去掉图中从 (A) 发出的所有线，只留下从 $a^{1}$ 到 $a^{10}$ 的那一条线，就可以表示这种情形。同样，例如英国赛马和英国指示犬显然都从其原始种群出发，在性状上缓慢分歧，而二者都没有产生任何新的分枝或品系。}

\pair{After ten thousand generations, species (A) is supposed to have produced three forms, $a^{10}$, $f^{10}$, and $m^{10}$, which, from having diverged in character during the successive generations, will have come to differ largely, but perhaps unequally, from each other and from their common parent. If we suppose the amount of change between each horizontal line in our diagram to be excessively small, these three forms may still be only well-marked varieties; or they may have arrived at the doubtful category of subspecies; but we have only to suppose the steps in the process of modification to be more numerous or greater in amount to convert these three forms into well-defined species: thus the diagram illustrates the steps by which the small differences distinguishing varieties are increased into the larger differences distinguishing species. By continuing the same process for a greater number of generations, as shown in the diagram in a condensed and simplified manner, we get eight species, marked by the letters between $a^{14}$ and $m^{14}$, all descended from (A). Thus, as I believe, species are multiplied and genera are formed.}{一万代之后，假定种 (A) 已经产生了三个类型，即 $a^{10}$、$f^{10}$ 和 $m^{10}$；由于它们在连续世代中发生了性状分歧，它们彼此之间以及同其共同亲本之间已经产生了很大但也许并不相等的差异。如果我们假定图中每两条水平线之间所代表的改变量极其微小，那么这三个类型可能仍只是明显的品种；或者它们可能已经达到亚种这一可疑范畴；但是只要假定改变过程中的步骤更多或数量更大，就能把这三个类型转变为界限明确的种。因此，该图说明了区分品种的小差异怎样逐步增大，成为区分种的大差异。把同一过程继续到更多世代，如图中以压缩和简化方式所示，我们得到八个种，由 $a^{14}$ 到 $m^{14}$ 之间的字母标出，它们全都出自 (A)。因此，我相信，种就是这样增多，属就是这样形成的。}

\pair{In a large genus it is probable that more than one species would vary. In the diagram I have assumed that a second species (I) has produced, by analogous steps, after ten thousand generations, either two well-marked varieties, $w^{10}$ and $z^{10}$, or two species, according to the amount of change supposed to be represented between the horizontal lines. After fourteen thousand generations, six new species, marked by the letters $n^{14}$ to $z^{14}$, are supposed to have been produced. In each genus, the species which are already extremely different in character will generally tend to produce the greatest number of modified descendants; for these will have the best chance of filling new and widely different places in the polity of nature: hence in the diagram I have chosen the extreme species (A), and the nearly extreme species (I), as those which have largely varied and have given rise to new varieties and species. The other nine species, marked by capital letters, of our original genus may for a long period continue transmitting unaltered descendants; and this is shown in the diagram by the dotted lines not prolonged far upwards from want of space.}{在一个大属中，很可能不止一个种会发生变异。在图中，我假定第二个种 (I) 通过类似步骤，在一万代之后产生了两个明显品种，即 $w^{10}$ 和 $z^{10}$，或者产生了两个种，这取决于假定水平线之间代表的改变量。第一万四千代以后，假定已经产生了六个新种，由 $n^{14}$ 到 $z^{14}$ 的字母标出。在每一个属中，那些在性状上已经极为不同的种，通常会倾向于产生最大数量的已改变后裔；因为这些后裔最有机会填补自然体系中新的、广泛不同的位置。因此在图中，我选择极端的种 (A) 和接近极端的种 (I)，作为那些发生大量变异并产生新品种和新种的种。我们原来这个属的其他九个种，即用大写字母标出的那些，可能在很长时期内继续传递未改变的后裔；图中由于空间不足，虚线没有向上延伸很远，表示了这一点。}

\pair{But during the process of modification represented in the diagram, another of our principles, namely that of extinction, will have played an important part. As in each fully stocked country natural selection necessarily acts by the selected form having some advantage in the struggle for life over other forms, there will be a constant tendency in the improved descendants of any one species to supplant and exterminate in each stage of descent their predecessors and their original parent. For it should be remembered that the competition will generally be most severe between those forms which are most nearly related to each other in habits, constitution, and structure. Hence all the intermediate forms between the earlier and later states, that is between the less and more improved state of a species, as well as the original parent species itself, will generally tend to become extinct. So it probably will be with many whole collateral lines of descent, which will be conquered by later and improved lines of descent. If, however, the modified offspring of a species get into some distinct country, or become quickly adapted to some quite new station in which child and parent do not come into competition, both may continue to exist.}{但是，在图中所表示的改变过程中，我们另一个原则，即灭绝原则，将起到重要作用。由于在每一个已经住满的国家，自然选择必然通过被选择的类型在生存斗争中比其他类型具有某种优势而发生作用，所以任何一个种的改进后裔，在每一继承阶段中，都将不断倾向于取代并消灭它们的前身和原始亲本。因为应当记住，竞争通常在习性、体质和构造上彼此最接近的类型之间最为激烈。因此，早期状态和晚期状态之间的一切中间类型，也就是一个种较少改进和更多改进状态之间的一切中间类型，以及原始亲本种本身，通常都会倾向于灭绝。许多完整的旁支血统很可能也是如此，它们会被较晚且改进的血统所征服。不过，如果某个种的已改变后代进入某个不同地区，或迅速适应某个全新的位置，使子代和亲本不发生竞争，二者就可能继续存在。}

\pair{If then our diagram be assumed to represent a considerable amount of modification, species (A) and all the earlier varieties will have become extinct, having been replaced by eight new species, $a^{14}$ to $m^{14}$; and (I) will have been replaced by six, $n^{14}$ to $z^{14}$, new species.}{那么，如果假定我们的图代表相当数量的改变，种 (A) 以及所有较早的品种都将已经灭绝，被八个新种，即 $a^{14}$ 到 $m^{14}$ 所取代；而 (I) 将被六个新种，即 $n^{14}$ 到 $z^{14}$ 所取代。}

\pair{But we may go further than this. The original species of our genus were supposed to resemble each other in unequal degrees, as is so generally the case in nature; species (A) being more nearly related to B, C, and D than to the other species, and species (I) more to G, H, K, L than to the others. These two species, (A) and (I), were also supposed to be very common and widely diffused species, so that they must originally have had some advantage over most of the other species of the genus. Their modified descendants, fourteen in number at the fourteen-thousandth generation, will probably have inherited some of the same advantages: they have also been modified and improved in a diversified manner at each stage of descent, so as to have become adapted to many related places in the natural economy of their country. It seems, therefore, to me extremely probable that they will have taken the places of, and thus exterminated, not only their parents (A) and (I), but likewise some of the original species which were most nearly related to their parents. Hence very few of the original species will have transmitted offspring to the fourteen-thousandth generation. We may suppose that only one (F), of the two species which were least closely related to the other nine original species, has transmitted descendants to this late stage of descent.}{但是我们还可以更进一步。我们假定本属的原始种彼此相似程度不等，这在自然界是极其普遍的；种 (A) 同 B、C、D 的关系，比同其他种的关系更近；种 (I) 同 G、H、K、L 的关系，比同其他种的关系更近。我们还假定这两个种，即 (A) 和 (I)，非常普通并且分布很广，因此它们原本一定比本属大多数其他种具有某些优势。到第一万四千代时，它们的已改变后裔共有十四个，大概会遗传某些同样的优势；它们还在每一继承阶段以多样方式被改变和改进，以致已经适应本国自然经济中许多相关位置。因此，在我看来，它们极有可能不仅已经取代并因而消灭了它们的亲本 (A) 和 (I)，而且也消灭了某些同其亲本关系最接近的原始种。因此，只有极少数原始种会把后代传到第一万四千代。我们可以假定，在同其他九个原始种关系最不密切的两个种中，只有一个 (F) 把后裔传到了这一较晚的继承阶段。}

\pair{The new species in our diagram descended from the original eleven species will now be fifteen in number. Owing to the divergent tendency of natural selection, the extreme amount of difference in character between species $a^{14}$ and $z^{14}$ will be much greater than that between the most different of the original eleven species. The new species, moreover, will be allied to each other in a widely different manner. Of the eight descendants from (A), the three marked $a^{14}$, $q^{14}$, $p^{14}$ will be nearly related from having recently branched off from $a^{10}$; $b^{14}$ and $f^{14}$, from having diverged at an earlier period from $a^{5}$, will be in some degree distinct from the three first-named species; and lastly, $o^{14}$, $e^{14}$, and $m^{14}$ will be nearly related one to the other, but, from having diverged at the first commencement of the process of modification, will be widely different from the other five species, and may constitute a subgenus or even a distinct genus.}{在我们的图中，出自原始十一个种的新种现在将有十五个。由于自然选择的分歧倾向，种 $a^{14}$ 和 $z^{14}$ 之间性状差异的极端程度，将远大于原始十一个种中差异最大的两个种之间的差异。此外，新种之间的亲缘关系方式也会大不相同。在 (A) 的八个后裔中，标为 $a^{14}$、$q^{14}$、$p^{14}$ 的三个，由于最近从 $a^{10}$ 分出，关系将很近；$b^{14}$ 和 $f^{14}$ 由于在较早时期从 $a^{5}$ 分歧出来，将同前三个种有某种程度的区别；最后，$o^{14}$、$e^{14}$ 和 $m^{14}$ 彼此关系很近，但由于在改变过程一开始就分歧出来，将同其他五个种大不相同，并可能构成一个亚属，甚至一个不同的属。}

\pair{The six descendants from (I) will form two subgenera or even genera. But as the original species (I) differed largely from (A), standing nearly at the extreme points of the original genus, the six descendants from (I) will, owing to inheritance, differ considerably from the eight descendants from (A); the two groups, moreover, are supposed to have gone on diverging in different directions. The intermediate species also, and this is a very important consideration, which connected the original species (A) and (I), have all become, excepting (F), extinct, and have left no descendants. Hence the six new species descended from (I), and the eight descended from (A), will have to be ranked as very distinct genera, or even as distinct subfamilies.}{出自 (I) 的六个后裔将形成两个亚属，甚至两个属。但是，由于原始种 (I) 同 (A) 差异很大，几乎处于原始属的两个极端点，出自 (I) 的六个后裔由于遗传，将同出自 (A) 的八个后裔有相当大的差异；此外，假定这两个类群已经继续向不同方向分歧。连接原始种 (A) 和 (I) 的中间种也都已经灭绝，并且没有留下后裔，除了 (F) 以外；这是一个非常重要的考虑。因此，出自 (I) 的六个新种和出自 (A) 的八个新种，必须被列为十分不同的属，甚至不同的亚科。}

\pair{Thus it is, as I believe, that two or more genera are produced by descent, with modification, from two or more species of the same genus. And the two or more parent species are supposed to have descended from some one species of an earlier genus. In our diagram this is indicated by the broken lines beneath the capital letters, converging in subbranches downwards towards a single point; this point representing a single species, the supposed single parent of our several new subgenera and genera.}{我相信，两个或更多属就是这样通过伴随改变的传衍，从同一属的两个或更多种产生出来的。而这两个或更多亲本种，又被假定是从较早某一属的某一个种传衍而来的。在我们的图中，这一点由大写字母下方的断线表示；这些断线以小分枝向下汇聚到一个点，这个点代表一个单一的种，即我们若干新亚属和新属的假定单一亲本。}

\begin{center}
[\ldots]
\end{center}

\subsection*{Summary of Chapter\\本章小结}

\pair{If during the long course of ages and under varying conditions of life, organic beings vary at all in the several parts of their organisation, and I think this cannot be disputed; if there be, owing to the high geometrical powers of increase of each species, at some age, season, or year, a severe struggle for life, and this certainly cannot be disputed; then, considering the infinite complexity of the relations of all organic beings to each other and to their conditions of existence, causing an infinite diversity in structure, constitution, and habits to be advantageous to them, I think it would be a most extraordinary fact if no variation had ever occurred useful to each being's own welfare, in the same way as so many variations have occurred useful to man. But if variations useful to any organic being do occur, assuredly individuals thus characterised will have the best chance of being preserved in the struggle for life; and from the strong principle of inheritance they will tend to produce offspring similarly characterised. This principle of preservation I have called, for the sake of brevity, Natural Selection. Natural selection, on the principle of qualities being inherited at corresponding ages, can modify the egg, seed, or young as easily as the adult. Amongst many animals, sexual selection will give its aid to ordinary selection by assuring to the most vigorous and best adapted males the greatest number of offspring. Sexual selection will also give characters useful to the males alone in their struggles with other males.}{如果在漫长年代过程中，并在变化的生活条件下，生物在其机体的各个部分确实会发生变异，我想这一点不能争辩；如果由于每一种具有高度的几何级数增长能力，在某一年龄、季节或年份中存在严酷的生存斗争，而这一点当然也不能争辩；那么，考虑到一切生物彼此之间以及它们同生存条件之间关系的无限复杂性，并由此使构造、体质和习性上的无限多样性对它们有利，我认为，若竟然从未发生过任何有利于每一种生物自身福利的变异，就像已经发生过那么多有利于人的变异那样，那将是最异常的事实。可是，如果确实发生了对任何生物有用的变异，那么具有这些性状的个体在生存斗争中必定最有机会被保存下来；并且由于强大的遗传原则，它们会倾向于产生具有相似性状的后代。为简洁起见，我把这种保存原则称为自然选择。按照性状在相应年龄遗传的原则，自然选择能够像改变成体一样容易地改变卵、种子或幼体。在许多动物中，性选择将帮助普通选择，因为它保证最强健、最适应的雄性留下最大数量的后代。性选择还会赋予只对雄性有用的性状，用于它们同其他雄性的斗争。}

\pair{Whether natural selection has really thus acted in nature, in modifying and adapting the various forms of life to their several conditions and stations, must be judged of by the general tenor and balance of evidence given in the following chapters. But we already see how it entails extinction, and how largely extinction has acted in the world's history geology plainly declares. Natural selection also leads to divergence of character; for more living beings can be supported on the same area the more they diverge in structure, habits, and constitution, of which we see proof by looking at the inhabitants of any small spot or at naturalised productions. Therefore, during the modification of the descendants of any one species, and during the incessant struggle of all species to increase in numbers, the more diversified these descendants become, the better will be their chance of succeeding in the battle of life. Thus the small differences distinguishing varieties of the same species will steadily tend to increase till they come to equal the greater differences between species of the same genus, or even of distinct genera.}{自然选择是否真的这样在自然界发生作用，改变各种生命类型并使之适应各自的条件和位置，必须依据以下各章所提供证据的一般趋势和权衡来判断。但是我们已经看到，它怎样导致灭绝；而地质学也清楚表明，灭绝在世界历史中曾起过多么巨大的作用。自然选择还会导致性状分歧；因为在同一面积上，生物在构造、习性和体质上越是分歧，就能供养越多生物。我们只要观察任何小地点的居民，或观察归化生物，便能看到这一点的证据。因此，在任何一个种的后裔发生改变的过程中，以及在一切种为增加数量而不断进行的斗争中，这些后裔变得越多样化，它们在生命战斗中成功的机会就越好。于是，区分同一种品种的小差异，将稳定地倾向于增大，直到等于同属种之间、甚至不同属之间较大的差异。}

\pair{We have seen that it is the common, the widely diffused, and widely ranging species, belonging to the larger genera, which vary most; and these will tend to transmit to their modified offspring that superiority which now makes them dominant in their own countries. Natural selection, as has just been remarked, leads to divergence of character and to much extinction of the less improved and intermediate forms of life. On these principles, I believe, the nature of the affinities of all organic beings may be explained. It is a truly wonderful fact—the wonder of which we are apt to overlook from familiarity—that all animals and all plants throughout all time and space should be related to each other in group subordinate to group, in the manner which we everywhere behold: namely, varieties of the same species most closely related together, species of the same genus less closely and unequally related together, forming sections and subgenera, species of distinct genera much less closely related, and genera related in different degrees, forming subfamilies, families, orders, subclasses, and classes. The several subordinate groups in any class cannot be ranked in a single file, but seem rather to be clustered round points, and these round other points, and so on in almost endless cycles. On the view that each species has been independently created, I can see no explanation of this great fact in the classification of all organic beings; but, to the best of my judgment, it is explained through inheritance and the complex action of natural selection, entailing extinction and divergence of character, as we have seen illustrated in the diagram.}{我们已经看到，属于较大属的普通、分布广泛、范围辽阔的种，变异最多；这些种会倾向于把现在使它们在本国占优势的优越性传给它们已改变的后代。正如刚才所说，自然选择导致性状分歧，并导致较少改进和中间生命类型的大量灭绝。依据这些原则，我相信，一切生物亲缘关系的性质都可以得到解释。一个真正奇妙的事实——只是由于我们太熟悉，往往忽略了它的奇妙——就是所有动物和所有植物，在全部时间和空间中，都以我们处处看到的那种类群隶属于类群的方式相互关联：也就是说，同一种的品种彼此关系最密切；同属的种关系较不密切且不相等，形成组和亚属；不同属的种关系远为疏远；属之间又以不同程度相关，形成亚科、科、目、亚纲和纲。任何一纲中的若干从属类群，不能排成单一行列，反而似乎是围绕一些点聚集，而这些点又围绕其他点，如此几乎无穷循环。如果认为每一个种都是独立创造的，我看不出怎样解释生物分类中的这一伟大事实；但据我最好的判断，它可以通过遗传和自然选择的复杂作用来解释，而自然选择导致灭绝和性状分歧，正如我们在图中已经说明的那样。}

\pair{The affinities of all the beings of the same class have sometimes been represented by a great tree. I believe this simile largely speaks the truth. The green and budding twigs may represent existing species; and those produced during each former year may represent the long succession of extinct species. At each period of growth all the growing twigs have tried to branch out on all sides, and to overtop and kill the surrounding twigs and branches, in the same manner as species and groups of species have tried to overmaster other species in the great battle for life. The limbs divided into great branches, and these into lesser and lesser branches, were themselves once, when the tree was small, budding twigs; and this connection of the former and present buds by ramifying branches may well represent the classification of all extinct and living species in groups subordinate to groups. Of the many twigs which flourished when the tree was a mere bush, only two or three, now grown into great branches, yet survive and bear all the other branches; so with the species which lived during long-past geological periods, very few now have living and modified descendants. From the first growth of the tree, many a limb and branch has decayed and dropped off; and these lost branches of various sizes may represent those whole orders, families, and genera which now have no living representatives, and which are known to us only from having been found in a fossil state. As we here and there see a thin straggling branch springing from a fork low down in a tree, and which by some chance has been favoured and is still alive on its summit, so we occasionally see an animal like the Ornithorhynchus or Lepidosiren, which in some small degree connects by its affinities two large branches of life, and which has apparently been saved from fatal competition by having inhabited a protected station. As buds give rise by growth to fresh buds, and these, if vigorous, branch out and overtop on all sides many a feebler branch, so by generation I believe it has been with the great Tree of Life, which fills with its dead and broken branches the crust of the earth, and covers the surface with its ever-branching and beautiful ramifications.}{同一纲中一切生物的亲缘关系，有时被比作一棵大树。我相信这个比喻在很大程度上说明了真理。绿色而发芽的小枝可以代表现存的种；以前每一年所产生的小枝，可以代表已经灭绝的种的漫长连续。在每一个生长期，所有生长中的小枝都试图向四面八方分枝，并越过和杀死周围的小枝和枝条，正如种和种的类群在生命的伟大战斗中试图压倒其他种一样。主干分成大枝，大枝又分成越来越小的枝条；而这些枝条本身，在树还很小时，也曾是发芽的小枝。过去和现在的芽通过分枝相连，这很可以代表一切已经灭绝和现存的种在类群隶属于类群中的分类。当这棵树还只是灌木时曾繁茂的许多小枝中，只有两三枝如今长成大枝，仍然存活并承载所有其他枝条；生活在久远地质时期的种也是这样，现在只有极少数具有活着的并已改变的后裔。从这棵树最初生长以来，许多主枝和枝条已经枯朽并脱落；这些大小不一的失落枝条，可以代表那些如今没有任何现生代表、我们只因其以化石状态被发现才知道的整个目、科和属。正如我们有时看到一根细弱零散的枝条从树下部的分叉处生出，并因某种偶然机缘受到有利对待，仍在树顶存活一样，我们也偶尔看到鸭嘴兽或肺鱼这样的动物，它们通过亲缘关系在某种小程度上连接生命的两大分枝，并且显然由于栖居在一个受保护的位置而免遭致命竞争。正如芽通过生长产生新芽，而这些新芽如果强健，就会向四面分枝并高过许多较弱枝条，我相信伟大的生命之树在世代更替中也是如此：它用已经死亡和断裂的枝条充满地壳，并以不断分枝而美丽的枝叶覆盖地表。}


% ===== 05_watson_dna.tex =====
% Processed OCR chunk: chunks/05_watson_dna.md
% Scope: Watson/Berry DNA selection only. Compile with a CJK-capable LaTeX engine.
% The parent document should load graphicx.


\section*{Beginnings of Genetics: From Mendel to Hitler}
\begin{center}
\large 遗传学的开端：从孟德尔到希特勒
\end{center}

\begin{quote}
\noindent\textbf{原文} My mother, Bonnie Jean, believed in genes. She was proud of her father's Scottish origins, and saw in him the traditional Scottish virtues of honesty, hard work, and thriftiness. She, too, possessed these qualities and felt that they must have been passed down to her from him. His tragic early death meant that her only nongenetic legacy was a set of tiny little girl's kilts he had ordered for her from Glasgow. Perhaps therefore it is not surprising that she valued her father's biological legacy over his material one.

\smallskip
\noindent\textbf{译文} 我的母亲邦妮·琼相信基因。她为父亲的苏格兰出身感到自豪，也在他身上看到了诚实、勤劳、节俭这些传统的苏格兰美德。她自己也具备这些品质，并觉得它们一定是父亲传给她的。父亲英年早逝，留给她的唯一非遗传遗产，是他从格拉斯哥为她订的一套小女孩穿的小格子裙。也许正因如此，她看重父亲留给她的生物学遗产，胜过物质遗产。
\end{quote}

\begin{quote}
\noindent\textbf{原文} Growing up, I had endless arguments with Mother about the relative roles played by nature and nurture in shaping us. By choosing nurture over nature, I was effectively subscribing to the belief that I could make myself into whatever I wanted to be. I did not want to accept that my genes mattered that much, preferring to attribute my Watson grandmother's extreme fatness to her having overeaten. If her shape was the product of her genes, then I too might have a hefty future. However, even as a teenager, I would not have disputed the evident basics of inheritance, that like begets like. My arguments with my mother concerned complex characteristics like aspects of personality, not the simple attributes that, even as an obstinate adolescent, I could see were passed down over the generations, resulting in "family likeness." My nose is my mother's and now belongs to my son Duncan.

\smallskip
\noindent\textbf{译文} 在成长过程中，我和母亲无休止地争论：在塑造一个人这件事上，先天与后天各自起多大作用。我站在后天一边，实际上就是认同这样一种信念：我可以把自己塑造成任何想成为的人。我不愿承认自己的基因有那么重要，更愿意把沃森家祖母异常肥胖的体形归因于吃得过多。如果她的身材是基因的产物，那么我自己的未来也可能相当沉重。不过，即便在十几岁时，我也不会否认遗传最显而易见的基本事实：种瓜得瓜，种豆得豆。我和母亲争论的，是性格等复杂特征；至于那些简单属性，即使是固执的少年，我也能看出它们会一代代传下去，造成所谓“家族相似”。我的鼻子来自母亲，如今又长在我儿子邓肯脸上。
\end{quote}

\begin{quote}
\noindent\textbf{原文} Sometimes characteristics come and go within a few generations, but sometimes they persist over many. One of the most famous examples of a long-lived trait is known as the "Hapsburg Lip." This distinctive elongation of the jaw and droopiness to the lower lip---which made the Hapsburg rulers of Europe such a nightmare assignment for generations of court portrait painters---was passed down intact over at least twenty-three generations.

\smallskip
\noindent\textbf{译文} 有些特征会在短短几代之间出现又消失，有些却能延续许多代。一个最著名的长期存在的性状，就是所谓的“哈布斯堡唇”。这种下颌明显拉长、下唇下垂的特征，使欧洲哈布斯堡统治者成为一代又一代宫廷肖像画家的噩梦；它完整地传递了至少二十三代。
\end{quote}

\begin{figure}[htbp]
\centering
\safeincludeimage{3e4cb4b042e59d2bcb9a019553b80b5387727b74271918f84ee6aa8dac26102c.jpg}
\caption{At age eleven, with my sister Elizabeth and my father, James.\\十一岁时，我和姐姐伊丽莎白以及父亲詹姆斯在一起。}
\end{figure}

\begin{quote}
\noindent\textbf{原文} The Hapsburgs added to their genetic woes by intermarrying. Arranging marriages between different branches of the Hapsburg clan and often among close relatives may have made political sense as a way of building alliances and ensuring dynastic succession, but it was anything but astute in genetic terms. Inbreeding of this kind can result in genetic disease, as the Hapsburgs found out to their cost. Charles II, the last of the Hapsburg monarchs in Spain, not only boasted a prize-worthy example of the family lip---he could not even chew his own food---but was also a complete invalid, and incapable, despite two marriages, of producing children.

\smallskip
\noindent\textbf{译文} 哈布斯堡家族又通过族内通婚加重了自己的遗传困境。安排哈布斯堡家族不同分支之间，且常常是近亲之间的婚姻，作为建立联盟、确保王朝继承的手段，在政治上也许说得通，但从遗传学上看绝不明智。这样的近亲繁殖可能导致遗传病，哈布斯堡家族为此付出了代价。西班牙哈布斯堡王朝最后一位君主查理二世，不仅拥有堪称“获奖级”的家族唇，甚至连自己的食物都无法咀嚼；他还体弱多病，尽管结过两次婚，却不能生育子女。
\end{quote}

\begin{quote}
\noindent\textbf{原文} Genetic disease has long stalked humanity. In some cases, such as Charles II's, it has had a direct impact on history. Retrospective diagnosis has suggested that George III, the English king whose principal claim to fame is to have lost the American colonies in the Revolutionary War, suffered from an inherited disease, porphyria, which causes periodic bouts of madness. Some historians---mainly British ones---have argued that it was the distraction caused by George's illness that permitted the Americans' against-the-odds military success. While most hereditary diseases have no such geopolitical impact, they nevertheless have brutal and often tragic consequences for the afflicted families, sometimes for many generations. Understanding genetics is not just about understanding why we look like our parents. It is also about coming to grips with some of humankind's oldest enemies: the flaws in our genes that cause genetic disease.

\smallskip
\noindent\textbf{译文} 遗传病长期纠缠着人类。在某些情况下，例如查理二世，它还直接影响了历史。后来的诊断推测认为，英王乔治三世患有一种遗传病卟啉症，这种病会引发周期性的精神失常；而乔治三世最出名的事，大概就是在美国独立战争中失去了北美殖民地。一些历史学家，主要是英国历史学家，认为正是乔治的疾病造成的分心，使美国人取得了原本胜算不大的军事胜利。多数遗传病并没有这样的地缘政治影响，但它们仍会给患病家庭带来残酷而常常悲惨的后果，有时会延续许多代。理解遗传学，不只是理解我们为什么长得像父母；它还意味着直面人类最古老的敌人之一：导致遗传病的基因缺陷。
\end{quote}

\begin{quote}
\noindent\textbf{原文} Our ancestors must have wondered about the workings of heredity as soon as evolution endowed them with brains capable of formulating the right kind of question. And the readily observable principle that close relatives tend to be similar can carry you a long way if, like our ancestors, your concern with the application of genetics is limited to practical matters like improving domesticated animals, for, say, milk yield in cattle, and plants, for, say, the size of fruit. Generations of careful selection---breeding initially to domesticate appropriate species, and then breeding only from the most productive cows and from the trees with the largest fruit---resulted in animals and plants tailor-made for human purposes. Underlying this enormous unrecorded effort is that simple rule of thumb: that the most productive cows will produce highly productive offspring and from the seeds of trees with large fruit large-fruited trees will grow. Thus, despite the extraordinary advances of the past hundred years or so, the twentieth and twenty-first centuries by no means have a monopoly on genetic insight. Although it was not until 1909 that the British biologist William Bateson gave the science of inheritance a name, genetics, and although the DNA revolution has opened up new and extraordinary vistas of potential progress, in fact the single greatest application of genetics to human well-being was carried out eons ago by anonymous ancient farmers. Almost everything we eat---cereals, fruit, meat, dairy products---is the legacy of that earliest and most far-reaching application of genetic manipulations to human problems.

\smallskip
\noindent\textbf{译文} 一旦进化赋予我们的祖先足以提出恰当问题的大脑，他们一定就开始思考遗传是如何运作的。近亲往往相似，这个很容易观察到的原则，如果像我们的祖先那样只把遗传学应用于实际事务，譬如改良家畜以提高牛奶产量，或改良植物以增大果实，就已经足以走得很远。一代又一代谨慎的选择，先是繁育适合驯化的物种，随后只让产奶最多的牛和果实最大的树参与繁殖，最终产生了专为人类目的而存在的动植物。在这项庞大而无文字记录的努力背后，是一条朴素经验法则：产量最高的母牛会生出高产后代，大果树的种子会长成结大果的树。因此，尽管过去一百多年进步惊人，二十世纪和二十一世纪绝非独占遗传学洞见。虽然直到1909年，英国生物学家威廉·贝特森才把遗传科学命名为 genetics（遗传学），虽然DNA革命已经打开了前所未有的进步前景，但事实上，对人类福祉影响最大的遗传学应用，是在远古时期由无名的古代农民完成的。我们吃的几乎一切东西，谷物、水果、肉类、乳制品，都是那次最早、也最深远的遗传操控应用于人类问题的遗产。
\end{quote}

\begin{quote}
\noindent\textbf{原文} An understanding of the actual mechanics of genetics proved a tougher nut to crack. Gregor Mendel (1822--1884) published his famous paper on the subject in 1866, and it was ignored by the scientific community for another thirty-four years. Why did it take so long? After all, heredity is a major aspect of the natural world, and, more important, it is readily and universally observable: a dog owner sees how a cross between a brown and black dog turns out, and all parents consciously or subconsciously track the appearance of their own characteristics in their children. One simple reason is that genetic mechanisms turn out to be complicated. Mendel's solution to the problem is not intuitively obvious: children are not, after all, simply a blend of their parents' characteristics. Perhaps most important was the failure by early biologists to distinguish between two fundamentally different processes, heredity and development. Today we understand that a fertilized egg contains the genetic information, contributed by both parents, that determines whether someone will be afflicted with, say, porphyria. That is heredity. The subsequent process, the development of a new individual from that humble starting point of a single cell, the fertilized egg, involves implementing that information. Broken down in terms of academic disciplines, genetics focuses on the information and developmental biology focuses on the use of that information. Lumping heredity and development together into a single phenomenon, early scientists never asked the questions that might have steered them toward the secret of heredity. Nevertheless, the effort had been under way in some form since the dawn of Western history.

\smallskip
\noindent\textbf{译文} 真正理解遗传学的实际机制，就困难得多。格雷戈尔·孟德尔（1822--1884）在1866年发表了关于这一主题的著名论文，但此后整整三十四年，科学共同体都对它置之不理。为什么要这么久？毕竟遗传是自然界的一个重要方面，更重要的是，它显而易见、人人可见：养狗的人会看到棕狗和黑狗杂交会生出怎样的后代；所有父母也都会有意无意地留意自己的特征在孩子身上出现。一个简单原因是，遗传机制确实复杂。孟德尔给出的解决方案并不符合直觉：孩子终究不是父母特征的简单混合。也许最重要的是，早期生物学家未能区分两个根本不同的过程：遗传与发育。今天我们知道，受精卵含有来自父母双方的遗传信息，这些信息决定了一个人是否会患上卟啉症之类疾病。这就是遗传。随后，从受精卵这一单细胞的谦卑起点发育成一个新个体的过程，则涉及对这些信息的执行。若按学科划分，遗传学关注信息本身，发育生物学关注这些信息如何被使用。早期科学家把遗传和发育混为一个现象，因而从未提出那些本可把他们引向遗传秘密的问题。不过，从西方历史的黎明起，相关努力就已经以某种形式展开了。
\end{quote}

\begin{quote}
\noindent\textbf{原文} The Greeks, including Hippocrates, pondered heredity. They devised a theory of "pangenesis," which claimed that sex involved the transfer of miniaturized body parts: "Hairs, nails, veins, arteries, tendons and their bones, albeit invisible as their particles are so small. While growing, they gradually separate from each other." This idea enjoyed a brief renaissance when Charles Darwin, desperate to support his theory of evolution by natural selection with a viable hypothesis of inheritance, put forward a modified version of pangenesis in the second half of the nineteenth century. In Darwin's scheme, each organ---eyes, kidneys, bones---contributed circulating "gemmules" that accumulated in the sex organs, and were ultimately exchanged in the course of sexual reproduction. Because these gemmules were produced throughout an organism's lifetime, Darwin argued any change that occurred in the individual after birth, like the stretch of a giraffe's neck imparted by craning for the highest foliage, could be passed on to the next generation. Ironically, then, to buttress his theory of natural selection Darwin came to champion aspects of Jean-Baptiste Lamarck's theory of inheritance of acquired characteristics---the very theory that his evolutionary ideas did so much to discredit. Darwin was invoking only Lamarck's theory of inheritance; he continued to believe that natural selection was the driving force behind evolution, but supposed that natural selection operated on the variation produced by pangenesis. Had Darwin known about Mendel's work, although Mendel published his results shortly after \emph{The Origin of Species} appeared and Darwin was never aware of them, he might have been spared the embarrassment of this late-career endorsement of some of Lamarck's ideas.

\smallskip
\noindent\textbf{译文} 希腊人，包括希波克拉底，也思考过遗传。他们提出了“泛生论”，主张性行为涉及微小化身体部件的转移：“毛发、指甲、静脉、动脉、肌腱以及它们的骨骼，虽然因颗粒极小而不可见；在生长过程中，它们逐渐彼此分离。”十九世纪下半叶，查尔斯·达尔文急切地想为自己的自然选择进化论找到一种可行的遗传假说，于是提出了修订版泛生论，这个观念因此短暂复兴。在达尔文的方案中，每个器官，眼睛、肾脏、骨骼，都会贡献循环于体内的“胚芽粒”；这些胚芽粒聚集到性器官中，最终在有性生殖过程中交换。由于这些胚芽粒在生物个体的一生中持续产生，达尔文认为，出生后个体发生的任何变化，例如长颈鹿为够到最高处树叶而伸长脖子，都可以传给下一代。具有讽刺意味的是，为了支撑自然选择理论，达尔文竟然支持了让-巴蒂斯特·拉马克获得性遗传理论的若干方面，而他的进化思想本来恰恰最有力地削弱了这一理论。达尔文只是借用了拉马克的遗传理论；他仍然相信自然选择是进化的驱动力，只是假设自然选择作用于泛生论所产生的变异。如果达尔文知道孟德尔的工作，虽然孟德尔的成果在《物种起源》发表后不久就出版了，达尔文却从未得知，他也许就不必在晚年因认可拉马克某些观点而尴尬了。
\end{quote}

\begin{quote}
\noindent\textbf{原文} Whereas pangenesis supposed that embryos were assembled from a set of minuscule components, another approach, "preformationism," avoided the assembly step altogether: either the egg or the sperm, exactly which was a contentious issue, contained a complete preformed individual called a homunculus. Development was therefore merely a matter of enlarging this into a fully formed being. In the days of preformationism, what we now recognize as genetic disease was variously interpreted: sometimes as a manifestation of the wrath of God or the mischief of demons and devils; sometimes as evidence of either an excess of or a deficit of the father's "seed"; sometimes as the result of "wicked thoughts" on the part of the mother during pregnancy. On the premise that fetal malformation can result when a pregnant mother's desires are thwarted, leaving her feeling stressed and frustrated, Napoleon passed a law permitting expectant mothers to shoplift. None of these notions, needless to say, did much to advance our understanding of genetic disease.

\smallskip
\noindent\textbf{译文} 泛生论认为胚胎由一套微小部件组装而成；另一种观点“预成论”则干脆省去了组装步骤：卵子或精子中，究竟是哪一个曾引发争论，已经包含了一个完整预成的微型个体，称为“小人”。发育因此不过是把这个小人放大成完整生命。在预成论盛行的年代，我们今天称为遗传病的现象有各种解释：有时被视为上帝震怒或恶魔作祟的表现；有时被看作父亲“种子”过多或不足的证据；有时则被归因于母亲怀孕期间的“邪念”。拿破仑甚至基于这样一种前提，通过了一项允许孕妇顺手牵羊的法律：如果孕妇的欲望受挫，使她感到压力和沮丧，胎儿可能因此畸形。不用说，这些观念都没有多少助益于我们理解遗传病。
\end{quote}

\begin{quote}
\noindent\textbf{原文} By the early nineteenth century, better microscopes had defeated preformationism. Look as hard as you like, you will never see a tiny homunculus curled up inside a sperm or egg cell. Pangenesis, though an earlier misconception, lasted rather longer---the argument would persist that the gemmules were simply too small to visualize---but was eventually laid to rest by August Weismann, who argued that inheritance depended on the continuity of germ plasm between generations and thus changes to the body over an individual's lifetime could not be transmitted to subsequent generations. His simple experiment involved cutting the tails off several generations of mice. According to Darwin's pangenesis, tailless mice would produce gemmules signifying "no tail," and so their offspring should develop a severely stunted hind appendage or none at all. When Weismann showed that the tail kept appearing after many generations of amputees, pangenesis bit the dust.

\smallskip
\noindent\textbf{译文} 到十九世纪初，更好的显微镜击败了预成论。无论你多么用力观察，都不可能在精子或卵细胞里看到一个蜷缩着的小人。泛生论虽然是更早的误解，却存活得更久一些，因为支持者可以坚持说胚芽粒实在太小，无法看见。最终奥古斯特·魏斯曼终结了它。他认为，遗传依赖的是代际之间生殖质的连续性，因此个体一生中身体发生的变化不能传给后代。他的简单实验，是连续几代剪掉小鼠的尾巴。按照达尔文的泛生论，无尾小鼠会产生表示“无尾”的胚芽粒，其后代就应当长出极度退化的后肢附属物，或者根本没有尾巴。当魏斯曼证明，许多代被截尾小鼠之后，尾巴仍然照常出现时，泛生论便寿终正寝。
\end{quote}

\begin{figure}[htbp]
\centering
\safeincludeimage{dbdb3499f1e7a9fdf01b306b63c809b52b17d23cc88f3666e7b6956cc8d0f509.jpg}
\caption{Genetics before Mendel: a homunculus, a preformed miniature person imagined to exist in the head of a sperm cell.\\孟德尔之前的遗传学：所谓“小人”，即想象中预先成形、存在于精子头部的微型人。}
\end{figure}

\begin{quote}
\noindent\textbf{原文} Gregor Mendel was the one who got it right. By any standards, however, he was an unlikely candidate for scientific superstardom. Born to a farming family in what is now the Czech Republic, he excelled at the village school and, at twenty-one, entered the Augustinian monastery at Brünn. After proving a disaster as a parish priest---his response to the ministry was a nervous breakdown---he tried his hand at teaching. By all accounts he was a good teacher, but in order to qualify to teach a full range of subjects, he had to take an exam. He failed it. Mendel's father superior, Abbot Napp, then dispatched him to the University of Vienna, where he was to bone up full-time for the retesting. Despite apparently doing well in physics at Vienna, Mendel again failed the exam, and so never rose above the rank of substitute teacher.

\smallskip
\noindent\textbf{译文} 真正说对的是格雷戈尔·孟德尔。不过无论按什么标准看，他都不像会成为科学巨星的人。孟德尔出生于今天捷克境内的一个农民家庭，在村校成绩出众，二十一岁时进入布尔诺的奥古斯丁修道院。作为堂区神父，他表现得一塌糊涂，面对神职工作的反应是精神崩溃；随后他尝试教书。众人都说他是个好教师，但为了取得教授完整课程的资格，他必须参加考试。他失败了。孟德尔的上级、修道院院长纳普随后派他去维也纳大学，让他全职补习准备重考。尽管他在维也纳的物理学方面似乎学得不错，孟德尔还是再次考试失败，因此从未升到代课教师以上的职位。
\end{quote}

\begin{quote}
\noindent\textbf{原文} Around 1856, at Abbot Napp's suggestion, Mendel undertook some scientific experiments on heredity. He chose to study a number of characteristics of the pea plants he grew in his own patch of the monastery garden. In 1865 he presented his results to the local natural history society in two lectures, and, a year later, published them in the society's journal. The work was a tour de force: the experiments were brilliantly designed and painstakingly executed, and his analysis of the results was insightful and deft. It seems that his training in physics contributed to his breakthrough because, unlike other biologists of that time, he approached the problem quantitatively. Rather than simply noting that crossbreeding of red and white flowers resulted in some red and some white offspring, Mendel actually counted them, realizing that the ratios of red to white progeny might be significant---as indeed they are. Despite sending copies of his article to various prominent scientists, Mendel found himself completely ignored by the scientific community. His attempt to draw attention to his results merely backfired. He wrote to his one contact among the ranking scientists of the day, botanist Karl Nägeli in Munich, asking him to replicate the experiments, and he duly sent off 140 carefully labeled packets of seeds. He should not have bothered. Nägeli believed that the obscure monk should be of service to him, rather than the other way around, so he sent Mendel seeds of his own favorite plant, hawkweed, challenging the monk to re-create his results with a different species. Sad to say, for various reasons, hawkweed is not well suited to breeding experiments such as those Mendel had performed on the peas. The entire exercise was a waste of his time.

\smallskip
\noindent\textbf{译文} 大约在1856年，在纳普院长的建议下，孟德尔开始进行一些关于遗传的科学实验。他选择研究自己在修道院花园一小块地里种植的豌豆的若干性状。1865年，他在当地自然历史学会用两场讲座报告了结果；一年后，又在学会期刊上发表。那项工作堪称杰作：实验设计精妙、执行细致，结果分析敏锐而娴熟。他在物理学方面的训练似乎促成了突破，因为与当时其他生物学家不同，他以定量方式处理问题。孟德尔并不只是注意到红花和白花杂交会产生一些红花后代和一些白花后代，他真正去数了它们，并意识到红白后代的比例可能重要，事实也正是如此。尽管他把论文副本寄给多位著名科学家，科学共同体却完全无视他。他试图引起注意，反而适得其反。他写信给自己在当时重要科学家中唯一的联系人，慕尼黑植物学家卡尔·内格利，请他重复实验，并寄去140包仔细贴好标签的种子。他本不该费这个劲。内格利认为，应当是这个无名修士为他服务，而不是反过来，于是寄给孟德尔自己最喜欢的植物山柳菊种子，要求这位修士用另一个物种重现结果。遗憾的是，出于多种原因，山柳菊并不适合孟德尔用豌豆所做的那类育种实验。整件事完全浪费了他的时间。
\end{quote}

\begin{quote}
\noindent\textbf{原文} Mendel's low-profile existence as monk-teacher-researcher ended abruptly in 1868 when, on Napp's death, he was elected abbot of the monastery. Although he continued his research---increasingly on bees and the weather---administrative duties were a burden, especially as the monastery became embroiled in a messy dispute over back taxes. Other factors, too, hampered him as a scientist. Portliness eventually curtailed his fieldwork: as he wrote, hill climbing had become "very difficult for me in a world where universal gravitation prevails." His doctors prescribed tobacco to keep his weight in check, and he obliged them by smoking twenty cigars a day, as many as Winston Churchill. It was not his lungs, however, that let him down: in 1884, at the age of sixty-one, Mendel succumbed to a combination of heart and kidney disease.

\smallskip
\noindent\textbf{译文} 孟德尔作为修士、教师和研究者的低调生活，在1868年戛然而止：纳普去世后，他被选为修道院院长。虽然他继续研究，且越来越多地研究蜜蜂和天气，但行政事务成为沉重负担，尤其是修道院卷入一场混乱的补税争端。其他因素也妨碍了他作为科学家的工作。肥胖最终限制了他的野外活动：正如他所写，在一个“万有引力占统治地位的世界里”，爬山对他已变得“非常困难”。医生开给他烟草来控制体重，他照办了，每天抽二十支雪茄，数量与温斯顿·丘吉尔相当。不过，最终拖垮他的并不是肺。1884年，六十一岁的孟德尔死于心脏病和肾病的共同作用。
\end{quote}

\begin{quote}
\noindent\textbf{原文} Not only were Mendel's results buried in an obscure journal, but they would have been unintelligible to most scientists of the era. He was far ahead of his time with his combination of careful experiment and sophisticated quantitative analysis. Little wonder, perhaps, that it was not until 1900 that the scientific community caught up with him. The rediscovery of Mendel's work, by three plant geneticists interested in similar problems, provoked a revolution in biology. At last the scientific world was ready for the monk's peas.

\smallskip
\noindent\textbf{译文} 孟德尔的结果不仅埋没在一份不起眼的期刊里，而且对当时大多数科学家来说也难以理解。他把严谨实验与精细的定量分析结合起来，远远超前于时代。因此，直到1900年科学共同体才追上他，也许并不奇怪。三位研究类似问题的植物遗传学家重新发现了孟德尔的工作，引发了一场生物学革命。科学世界终于准备好迎接这位修士的豌豆了。
\end{quote}

\begin{quote}
\noindent\textbf{原文} Mendel realized that there are specific factors---later to be called "genes"---that are passed from parent to offspring. He worked out that these factors come in pairs and that the offspring receives one from each parent.

\smallskip
\noindent\textbf{译文} 孟德尔认识到，有一些特定因子会从亲代传给子代；这些因子后来被称为“基因”。他弄清楚了这些因子成对存在，而后代从父母双方各获得一个。
\end{quote}

\begin{quote}
\noindent\textbf{原文} Noticing that peas came in two distinct colors, green and yellow, he deduced that there were two versions of the pea-color gene. A pea has to have two copies of the G version if it is to become green, in which case we say that it is GG for the pea-color gene. It must therefore have received a G pea-color gene from both of its parents. However, yellow peas can result both from YY and YG combinations. Having only one copy of the Y version is sufficient to produce yellow peas. Y trumps G. Because in the YG case the Y signal dominates the G signal, we call Y "dominant." The subordinate G version of the pea-color gene is called "recessive." Each parent pea plant has two copies of the pea-color gene, yet it contributes only one copy to each offspring; the other copy is furnished by the other parent. In plants, pollen grains contain sperm cells---the male contribution to the next generation---and each sperm cell contains just one copy of the pea-color gene. A parent pea plant with a YG combination will produce sperm that contain either a Y version or a G one. Mendel discovered that the process is random: 50\% of the sperm produced by that plant will have a Y and 50\% will have a G.

\smallskip
\noindent\textbf{译文} 孟德尔注意到豌豆有绿色和黄色两种截然不同的颜色，于是推断豌豆颜色基因有两种版本，也就是两种等位基因。一个豌豆若要成为绿色，就必须拥有两个G版本；在这种情况下，我们说它在豌豆颜色基因上是GG。因此，它必然分别从父母双方各获得了一个G版本的豌豆颜色基因。然而，黄色豌豆既可能来自YY组合，也可能来自YG组合。只要有一个Y版本，就足以产生黄色豌豆。Y压过G。因为在YG情况下，Y信号支配G信号，我们称Y为“显性”；而从属的G版本豌豆颜色基因则称为“隐性”。每株亲代豌豆都有两个豌豆颜色基因副本，却只给每个后代贡献一个；另一个副本由另一亲本提供。在植物中，花粉粒含有精细胞，这是给下一代的雄性贡献；每个精细胞只含有豌豆颜色基因的一个副本。带有YG组合的亲代豌豆，会产生含有Y版本或G版本的精子。孟德尔发现这个过程是随机的：该植物产生的精子有50\%带Y，另50\%带G。
\end{quote}

\begin{quote}
\noindent\textbf{原文} Suddenly many of the mysteries of heredity made sense. Characteristics, like the Hapsburg Lip, that are transmitted with a high probability, actually 50\%, from generation to generation are dominant. Other characteristics that appear in family trees much more sporadically, often skipping generations, may be recessive. When a gene is recessive an individual has to have two copies of it for the corresponding trait to be expressed. Those with one copy of the gene are carriers: they do not themselves exhibit the characteristic, but they can pass the gene on. Albinism, in which the body fails to produce pigment so the skin and hair are strikingly white, is an example of a recessive characteristic that is transmitted in this way. Therefore, to be albino you have to have two copies of the gene, one from each parent. This was the case with the Reverend Dr. William Archibald Spooner, who was also---perhaps only by coincidence---prone to a peculiar form of linguistic confusion whereby, for example, "a well-oiled bicycle" might become "a well-boiled icicle." Such reversals would come to be termed "spoonerisms" in his honor. Your parents, meanwhile, may have shown no sign of the gene at all. If, as is often the case, each has only one copy, then they are both carriers. The trait has skipped at least one generation.

\smallskip
\noindent\textbf{译文} 一下子，遗传的许多谜团都有了道理。像哈布斯堡唇这样以高概率，实际上是50\%，代代相传的性状，是显性性状。其他在家谱中出现得零散得多、常常隔代出现的性状，可能是隐性性状。当一个基因是隐性时，个体必须拥有它的两个副本，相应性状才会表现出来。只拥有一个副本的人是携带者：他们自己不表现这个特征，却可以把该基因传下去。白化病就是这样传递的隐性特征之一；患者身体不能产生色素，因此皮肤和头发呈醒目的白色。所以，要成为白化病患者，你必须拥有该基因的两个副本，一个来自父亲，一个来自母亲。威廉·阿奇博尔德·斯普纳牧师博士就是如此；他还容易出现一种特殊的语言混乱，也许只是巧合，例如把“a well-oiled bicycle”（一辆润滑良好的自行车）说成“a well-boiled icicle”（一根煮得很好的冰柱）。这种音位调换后来便以他的名字称为“斯普纳式口误”。与此同时，你的父母可能完全没有显示出该基因的迹象。如果像常见情况那样，他们各自只有一个副本，那么他们都是携带者。这个性状至少跳过了一代。
\end{quote}

\begin{quote}
\noindent\textbf{原文} Mendel's results implied that things---material objects---were transmitted from generation to generation. But what was the nature of these things?

\smallskip
\noindent\textbf{译文} 孟德尔的结果暗示，有某些东西，某些物质实体，在一代又一代之间传递。但这些东西的本质是什么？
\end{quote}

\begin{quote}
\noindent\textbf{原文} At about the time of Mendel's death in 1884, scientists using ever-improving optics to study the minute architecture of cells coined the term "chromosome" to describe the long stringy bodies in the cell nucleus. But it was not until 1902 that Mendel and chromosomes came together.

\smallskip
\noindent\textbf{译文} 大约在孟德尔1884年去世前后，科学家利用日益改进的光学设备研究细胞的微细结构，并创造了“染色体”一词，用来描述细胞核中长而丝状的结构。但直到1902年，孟德尔的因子才与染色体联系到一起。
\end{quote}

\begin{figure}[htbp]
\centering
\safeincludeimage{441c947adb700e4e5337fe0936f94205f10859932393b7b8c9ca50eb6e986016.jpg}
\caption{The human X chromosome, as seen with an electron microscope.\\电子显微镜下的人类X染色体。}
\end{figure}

\begin{quote}
\noindent\textbf{原文} A medical student at Columbia University, Walter Sutton, realized that chromosomes had a lot in common with Mendel's mysterious factors. Studying grasshopper chromosomes, Sutton noticed that most of the time they are doubled up---just like Mendel's paired factors. But Sutton also identified one type of cell in which chromosomes were not paired: the sex cells. Grasshopper sperm have only a single set of chromosomes, not a double set. This was exactly what Mendel had described: his pea plant sperm cells also only carried a single copy of each of his factors. It was clear that Mendel's factors, now called genes, must be on the chromosomes.

\smallskip
\noindent\textbf{译文} 哥伦比亚大学医学生沃尔特·萨顿意识到，染色体与孟德尔那些神秘因子有许多共同点。萨顿研究蝗虫染色体时注意到，它们大多数时候成对存在，正像孟德尔的成对因子一样。但萨顿还发现了一类染色体不成对的细胞：性细胞。蝗虫精子只有一套染色体，而不是两套。这正是孟德尔描述过的情况：他的豌豆精细胞也只携带每个因子的一个副本。显然，孟德尔的因子，如今称为基因，一定在染色体上。
\end{quote}

\begin{quote}
\noindent\textbf{原文} In Germany Theodor Boveri independently came to the same conclusions as Sutton, and so the biological revolution their work had precipitated came to be called the Sutton--Boveri chromosome theory of inheritance. Suddenly genes were real. They were on chromosomes, and you could actually see chromosomes through the microscope.

\smallskip
\noindent\textbf{译文} 在德国，西奥多·鲍维里独立得出了与萨顿相同的结论。因此，他们的工作所触发的生物学革命，被称为萨顿--鲍维里染色体遗传理论。突然之间，基因变成了真实存在的东西。它们位于染色体上，而染色体确实可以通过显微镜看到。
\end{quote}

\begin{quote}
\noindent\textbf{原文} Not everyone bought the Sutton--Boveri theory. One skeptic was Thomas Hunt Morgan, also at Columbia. Looking down the microscope at those stringy chromosomes, he could not see how they could account for all the changes that occur from one generation to the next. If all the genes were arranged along chromosomes, and all chromosomes were transmitted intact from one generation to the next, then surely many characteristics would be inherited together. But since empirical evidence showed this not to be the case, the chromosomal theory seemed insufficient to explain the variation observed in nature. Being an astute experimentalist, however, Morgan had an idea how he might resolve such discrepancies. He turned to the fruit fly, \emph{Drosophila melanogaster}, the drab little beast that, ever since Morgan, has been so beloved by geneticists.

\smallskip
\noindent\textbf{译文} 并非所有人都接受萨顿--鲍维里理论。怀疑者之一是同在哥伦比亚大学的托马斯·亨特·摩尔根。他透过显微镜看着那些丝状染色体，却看不出它们如何能够解释一代到下一代之间发生的所有变化。如果所有基因都沿染色体排列，而所有染色体又完整地从一代传到下一代，那么许多性状理应一起遗传。但经验证据显示事实并非如此，所以染色体理论似乎不足以解释自然界中观察到的变异。不过，摩尔根是一位敏锐的实验学家，他想到也许可以解决这些矛盾。他转向了果蝇，\emph{Drosophila melanogaster}，这只不起眼的小生物自摩尔根以来一直深受遗传学家喜爱。
\end{quote}

\begin{figure}[htbp]
\centering
\safeincludeimage{58fbae006993552b21bb6d1e4a5d19607609d27f70da86419fe2eebeb43e61d7.jpg}
\caption{Notoriously camera-shy T. H. Morgan was photographed surreptitiously while at work in the fly room.\\以怕照相闻名的T. H. 摩尔根，在果蝇室工作时被人偷偷拍了下来。}
\end{figure}

\begin{quote}
\noindent\textbf{原文} In fact, Morgan was not the first to use the fruit fly in breeding experiments---that distinction belonged to a lab at Harvard that first put the critter to work in 1901---but it was Morgan's work that put the fly on the scientific map. \emph{Drosophila} is a good choice for genetic experiments. It is easy to find, as anyone who has left out a bunch of overripe bananas during the summer well knows; it is easy to raise, bananas will do as feed; you can accommodate hundreds of flies in a single milk bottle, and Morgan's students had no difficulty acquiring milk bottles, pinching them at dawn from doorsteps in their Manhattan neighborhood; and it breeds and breeds and breeds, a whole generation taking about ten days, with each female laying several hundred eggs. Starting in 1907 in a famously squalid, cockroach-infested, banana-stinking lab that came to be known affectionately as the "fly room," Morgan and his students, "Morgan's boys" as they were called, set to work on fruit flies.

\smallskip
\noindent\textbf{译文} 事实上，摩尔根并不是第一个把果蝇用于育种实验的人；这一荣誉属于哈佛的一个实验室，那里在1901年率先让这种小虫投入工作。但正是摩尔根的研究让果蝇登上了科学地图。果蝇是遗传实验的好材料。它容易找到，夏天把一串过熟香蕉放在外面的人都知道这一点；它容易饲养，香蕉就能当饲料；一个牛奶瓶里可以容纳数百只果蝇，而摩尔根的学生弄到牛奶瓶毫不费力，他们清晨从曼哈顿邻里门前顺走即可；它还不停繁殖，一代约十天，每只雌蝇产卵数百枚。1907年起，在一个出了名肮脏、蟑螂横行、香蕉味刺鼻，后来被亲切称为“果蝇室”的实验室里，摩尔根和他的学生们，人称“摩尔根的男孩们”，开始研究果蝇。
\end{quote}

\begin{quote}
\noindent\textbf{原文} Unlike Mendel, who could rely on the variant strains isolated over the years by farmers and gardeners---yellow peas as opposed to green ones, wrinkled skin as opposed to smooth---Morgan had no menu of established genetic differences in the fruit fly to draw upon. And you cannot do genetics until you have isolated some distinct characteristics to track through the generations. Morgan's first goal therefore was to find "mutants," the fruit fly equivalents of yellow or wrinkled peas. He was looking for genetic novelties, random variations that somehow simply appeared in the population.

\smallskip
\noindent\textbf{译文} 孟德尔可以依靠农民和园丁多年分离出来的变异品系，例如黄豌豆相对于绿豌豆、皱皮相对于光滑；摩尔根却没有一份果蝇既有遗传差异的菜单可用。而在分离出某些明确特征并能跨代追踪之前，你无法做遗传学。因此，摩尔根的首要目标是寻找“突变体”，也就是果蝇中相当于黄豌豆或皱豌豆的东西。他寻找的是遗传新奇性，是种群中不知怎么就随机出现的变异。
\end{quote}

\begin{quote}
\noindent\textbf{原文} One of the first mutants Morgan observed turned out to be one of the most instructive. While normal fruit flies have red eyes, these had white ones. And he noticed that the white-eyed flies were typically male. It was known that the sex of a fruit fly---or, for that matter, the sex of a human---is determined chromosomally: females have two copies of the X chromosome, whereas males have one copy of the X and one copy of the much smaller Y. In light of this information, the white-eye result suddenly made sense: the eye-color gene is located on the X chromosome and the white-eye mutation, W, is recessive. Because males have only a single X chromosome, even recessive genes, in the absence of a dominant counterpart to suppress them, are automatically expressed. White-eyed females were relatively rare because they typically had only one copy of W, so they expressed the dominant red eye color. By correlating a gene---the one for eye color---with a chromosome, the X, Morgan, despite his initial reservations, had effectively proved the Sutton--Boveri theory. He had also found an example of "sex-linkage," in which a particular characteristic is disproportionately represented in one sex.

\smallskip
\noindent\textbf{译文} 摩尔根最早观察到的突变体之一，后来证明极富启发性。正常果蝇有红眼，而这些果蝇是白眼。他还注意到，白眼果蝇通常是雄性。人们已经知道，果蝇的性别，事实上人类的性别也是如此，由染色体决定：雌性有两条X染色体，雄性有一条X染色体和一条小得多的Y染色体。有了这一信息，白眼结果立刻变得有道理：眼色基因位于X染色体上，白眼突变W是隐性的。由于雄性只有一条X染色体，即便是隐性基因，只要没有显性对应基因压制，也会自动表现出来。白眼雌蝇相对罕见，因为它们通常只有一个W副本，因此表现显性红眼。摩尔根把一个基因，也就是眼色基因，与一条染色体X联系起来，尽管他最初有所保留，却实际上证明了萨顿--鲍维里理论。他还发现了“伴性遗传”的例子，即某一特定性状在某一性别中不成比例地出现。
\end{quote}

\begin{quote}
\noindent\textbf{原文} Like Morgan's fruit flies, Queen Victoria provides a famous example of sex-linkage. On one of her X chromosomes, she had a mutated gene for hemophilia, the "bleeding disease" in whose victims proper blood clotting fails to occur. Because her other copy was normal, and the hemophilia gene is recessive, she herself did not have the disease. But she was a carrier. Her daughters did not have the disease either; evidently each possessed at least one copy of the normal version. But Victoria's sons were not all so lucky. Like all males, fruit fly males included, each had only one X chromosome; this was necessarily derived from Victoria, since a Y chromosome could have come only from Prince Albert, Victoria's husband. Because Victoria had one mutated copy and one normal copy, each of her sons had a 50--50 chance of having the disease. Prince Leopold drew the short straw: he developed hemophilia, and died at thirty-one, bleeding to death after a minor fall. Two of Victoria's daughters, Princesses Alice and Beatrice, were carriers, having inherited the mutated gene from their mother. They each produced carrier daughters and sons with hemophilia. Alice's grandson Alexis, heir to the Russian throne, had hemophilia, and would doubtless have died young had the Bolsheviks not gotten to him first.

\smallskip
\noindent\textbf{译文} 与摩尔根的果蝇一样，维多利亚女王也提供了一个著名的伴性遗传例子。她的一条X染色体上带有血友病的突变基因；血友病又称“出血病”，患者无法正常凝血。由于她的另一个副本正常，而血友病基因是隐性的，她本人并未患病。但她是携带者。她的女儿们也没有患病，显然每个人至少有一个正常版本。可维多利亚的儿子们并非都如此幸运。像所有雄性一样，包括雄性果蝇在内，每个男性只有一条X染色体；这条X必然来自维多利亚，因为Y染色体只能来自维多利亚的丈夫阿尔伯特亲王。由于维多利亚有一个突变副本和一个正常副本，她每个儿子患病的机会都是一半对一半。利奥波德王子抽中了下签：他患上血友病，三十一岁时因一次小摔倒后流血不止而死。维多利亚的两个女儿，爱丽丝公主和比阿特丽斯公主，也是携带者，她们从母亲那里继承了突变基因。她们各自生下携带者女儿和血友病儿子。爱丽丝的孙子、俄国王位继承人阿列克谢患有血友病；若不是布尔什维克先下了手，他无疑也会早夭。
\end{quote}

\begin{quote}
\noindent\textbf{原文} Morgan's fruit flies had other secrets to reveal. In the course of studying genes located on the same chromosome, Morgan and his students found that chromosomes actually break apart and re-form during the production of sperm and egg cells. This meant that Morgan's original objections to the Sutton--Boveri theory were unwarranted: the breaking and re-forming---"recombination," in modern genetic parlance---shuffles gene copies between members of a chromosome pair. This means that, say, the copy of chromosome 12 I got from my mother, the other of course comes from my father, is in fact a mix of my mother's two copies of chromosome 12, one of which came from her mother and one from her father. Her two 12s recombined---exchanged material---during the production of the egg cell that eventually turned into me. Thus my maternally derived chromosome 12 can be viewed as a mosaic of my grandparents' 12s. Of course, my mother's maternally derived 12 was itself a mosaic of her grandparents' 12s, and so on.

\smallskip
\noindent\textbf{译文} 摩尔根的果蝇还有其他秘密可揭示。在研究位于同一染色体上的基因时，摩尔根和学生发现，染色体在精子和卵细胞形成过程中实际上会断裂并重新连接。这意味着摩尔根最初对萨顿--鲍维里理论的反对并不成立：断裂和重组，用现代遗传学术语说就是“重组”，会在一对染色体成员之间洗牌式地交换基因副本。也就是说，比如我从母亲那里得到的第12号染色体副本，另一条当然来自父亲，事实上是母亲两条12号染色体的混合体；其中一条来自她的母亲，另一条来自她的父亲。在形成最终变成我的那枚卵细胞时，她的两条12号染色体发生了重组，也就是交换了物质。因此，我母源性的12号染色体可以看作外祖父母两条12号染色体的镶嵌体。当然，我母亲从她母亲那里得到的12号染色体本身，又是她的祖父母两条12号染色体的镶嵌体，如此上溯。
\end{quote}

\begin{quote}
\noindent\textbf{原文} Recombination permitted Morgan and his students to map out the positions of particular genes along a given chromosome. Recombination involves breaking and re-forming chromosomes. Because genes are arranged like beads along a chromosome string, a break is statistically much more likely to occur between two genes that are far apart, with more potential break points intervening, on the chromosome than between two genes that are close together. If, therefore, we see a lot of reshuffling for any two genes on a single chromosome, we can conclude that they are a long way apart; the rarer the reshuffling, the closer the genes likely are. This basic and immensely powerful principle underlies all of genetic mapping. One of the primary tools of scientists involved in the Human Genome Project and of researchers at the forefront of the battle against genetic disease was thus developed all those years ago in the filthy, cluttered Columbia fly room. Each new headline in the science section of the newspaper these days along the lines of "Gene for Something Located" is a tribute to the pioneering work of Morgan and his boys.

\smallskip
\noindent\textbf{译文} 重组使摩尔根和学生能够绘制特定基因在给定染色体上的位置图。重组涉及染色体的断裂与重新连接。由于基因像珠子一样沿染色体这根线排列，两个相距很远的基因之间有更多潜在断点，因此断裂在它们之间发生的统计概率，远高于在两个相邻基因之间发生的概率。所以，如果我们看到同一条染色体上任意两个基因发生大量重新组合，就可以断定它们相距很远；重组越少，基因就越可能靠得近。这个基本而极其有力的原则，是全部遗传作图的基础。后来人类基因组计划科学家以及抗击遗传病前沿研究者的一项主要工具，竟是在多年前哥伦比亚大学那间肮脏杂乱的果蝇室里发展出来的。如今报纸科学版每出现一条类似“某某基因定位成功”的新标题，都是对摩尔根和他的男孩们开创性工作的致敬。
\end{quote}

\begin{quote}
\noindent\textbf{原文} The rediscovery of Mendel's work, and the breakthroughs that followed it, sparked a surge of interest in the social significance of genetics. While scientists had been grappling with the precise mechanisms of heredity through the eighteenth and nineteenth centuries, public concern had been mounting about the burden placed on society by what came to be called the "degenerate classes"---the inhabitants of poorhouses, workhouses, and insane asylums. What could be done with these people? It remained a matter of controversy whether they should be treated charitably---which, the less charitably inclined claimed, ensured such folk would never exert themselves and would therefore remain forever dependent on the largesse of the state or of private institutions---or whether they should be simply ignored, which, according to the charitably inclined, would result only in perpetuating the inability of the unfortunate to extricate themselves from their blighted circumstances.

\smallskip
\noindent\textbf{译文} 孟德尔工作的重新发现以及随后的突破，引发了人们对遗传学社会意义的浓厚兴趣。十八、十九世纪期间，科学家一直在努力解决遗传的精确机制；与此同时，公众越来越担心所谓“退化阶层”给社会带来的负担，即济贫院、劳役所和精神病院中的居民。该如何对待这些人？一种主张是慈善救助；不那么有慈善心的人认为，这只会保证这些人永远不肯努力，从而永远依赖国家或私人机构的慷慨施舍。另一种主张是干脆不管；而有慈善心的人则认为，这只会让不幸者无法摆脱悲惨处境的状态永久化。争议一直存在。
\end{quote}

\begin{quote}
\noindent\textbf{原文} The publication of Darwin's \emph{Origin of Species} in 1859 brought these issues into sharp focus. Although Darwin carefully omitted to mention human evolution, fearing that to do so would only further inflame an already raging controversy, it required no great leap of imagination to apply his idea of natural selection to humans. Natural selection is the force that determines the fate of all genetic variations in nature---mutations like the one Morgan found in the fruit fly eye-color gene, but also perhaps differences in the abilities of human individuals to fend for themselves.

\smallskip
\noindent\textbf{译文} 1859年达尔文《物种起源》的出版，使这些问题变得格外突出。达尔文担心提及人类进化只会进一步激化已经沸腾的争论，因此谨慎地避开了这一点；但把他的自然选择思想应用于人类，并不需要多大的想象力跳跃。自然选择是决定自然界一切遗传变异命运的力量，既包括摩尔根在果蝇眼色基因中发现的那类突变，也许也包括人类个体自谋生存能力上的差异。
\end{quote}

\begin{quote}
\noindent\textbf{原文} Natural populations have an enormous reproductive potential. Take fruit flies, with their generation time of just ten days, and females that produce some three hundred eggs apiece, half of which will be female: starting with a single fruit fly couple, after a month, that is, three generations later, you will have \(150 \times 150 \times 150\) fruit flies on your hands---that is more than 3 million flies, all of them derived from just one pair in just one month. Darwin made the point by choosing a species from the other end of the reproductive spectrum:

\smallskip
\noindent\textbf{译文} 自然种群具有巨大的繁殖潜力。以果蝇为例，它们一代只需十天，每只雌蝇产卵约三百枚，其中一半会是雌性。从一对果蝇开始，一个月后，也就是三代以后，你手上会有 \(150 \times 150 \times 150\) 只果蝇，超过三百万只，而且全部在一个月内由最初一对繁衍而来。达尔文为了说明这一点，选择了繁殖谱另一端的物种：
\end{quote}

\begin{quote}
\noindent\textbf{English quotation.} The elephant is reckoned to be the slowest breeder of all known animals, and I have taken some pains to estimate its probable minimum rate of natural increase: it will be under the mark to assume that it breeds when thirty years old, and goes on breeding till ninety years old, bringing forth three pairs of young in this interval; if this be so, at the end of the fifth century there would be alive fifteen million elephants, descended from the first pair.

\smallskip
\noindent\textbf{中文引文。} 大象被认为是所有已知动物中繁殖最慢的；我曾费心估算它可能的最低自然增长率：假设它三十岁开始繁殖，并持续到九十岁，在此期间生下三对子代，这个估计还会偏低；若果真如此，到第五个世纪末，将会有一千五百万头大象存活，全都源自最初那一对。
\end{quote}

\begin{quote}
\noindent\textbf{原文} All these calculations assume that all the baby fruit flies and all the baby elephants make it successfully to adulthood. In theory, therefore, there must be an infinitely large supply of food and water to sustain this kind of reproductive overdrive. In reality, of course, those resources are limited, and not all baby fruit flies or baby elephants make it. There is competition among individuals within a species for those resources. What determines who wins the struggle for access to the resources? Darwin pointed out genetic variation means that some individuals have advantages in what he called "the struggle for existence." To take the famous example of Darwin's finches from the Galápagos Islands, those individuals with genetic advantages---like the right size of beak for eating the most abundant seeds---are more likely to survive and reproduce. So the advantageous genetic variant---having a bill the right size---tends to be passed on to the next generation. The result is that natural selection enriches the next generation with the beneficial mutation so that eventually, over enough generations, every member of the species ends up with that characteristic.

\smallskip
\noindent\textbf{译文} 所有这些计算都假定，所有果蝇幼体和所有小象都能成功长到成年。因此在理论上，必须有无限多的食物和水来支撑这种繁殖狂飙。当然现实中这些资源是有限的，并不是所有果蝇幼体或小象都能活下来。同一物种的个体之间会为这些资源竞争。什么决定了谁在获取资源的斗争中获胜？达尔文指出，遗传变异意味着有些个体在他所谓的“生存斗争”中具有优势。以加拉帕戈斯群岛上著名的达尔文雀为例，具有遗传优势的个体，例如喙的大小正适合吃最丰富种子者，更有可能存活并繁殖。因此，有利的遗传变异，也就是拥有合适大小的喙，往往会传给下一代。结果是，自然选择让下一代中有益突变更加丰富；经过足够多代后，该物种每个成员最终都会拥有这一特征。
\end{quote}

\begin{quote}
\noindent\textbf{原文} The Victorians applied the same logic to humans. They looked around and were alarmed by what they saw. The decent, moral, hardworking middle classes were being massively out-reproduced by the dirty, immoral, lazy lower classes. The Victorians assumed that the virtues of decency, morality, and hard work ran in families just as the vices of filth, wantonness, and indolence did. Such characteristics must then be hereditary; thus, to the Victorians, morality and immorality were merely two of Darwin's genetic variants. And if the great unwashed were out-reproducing the respectable classes, then the "bad" genes would be increasing in the human population. The species was doomed! Humans would gradually become more and more depraved as the "immorality" gene became more and more common.

\smallskip
\noindent\textbf{译文} 维多利亚时代的人把同样逻辑应用到人类身上。他们环顾四周，并对所见之事感到惊恐。体面、道德、勤劳的中产阶级，在繁殖数量上被肮脏、不道德、懒惰的下层阶级远远超过。维多利亚人假定，体面、道德和勤劳这些美德像污秽、放荡和懒惰这些恶习一样，会在家族中流传。那么这些特征必然是遗传的；于是对维多利亚人来说，道德和不道德不过是达尔文式遗传变异的两种。如果“不洗澡的大众”比体面阶级繁殖更多，那么“坏”基因就会在人类群体中增加。这个物种注定完了！随着“不道德”基因越来越常见，人类会逐渐变得越来越堕落。
\end{quote}

\begin{quote}
\noindent\textbf{原文} Francis Galton had good reason to pay special attention to Darwin's book, as the author was his cousin and friend. Darwin, some thirteen years older, had provided guidance during Galton's rather rocky college experience. But it was \emph{The Origin of Species} that would inspire Galton to start a social and genetic crusade that would ultimately have disastrous consequences. In 1883, a year after his cousin's death, Galton gave the movement a name: eugenics.

\smallskip
\noindent\textbf{译文} 弗朗西斯·高尔顿有充分理由格外关注达尔文的书，因为作者既是他的表兄，也是他的朋友。年长约十三岁的达尔文曾在高尔顿颇为坎坷的大学经历中给予指导。但正是《物种起源》激励高尔顿发起了一场社会和遗传学十字军运动，其后果最终将十分灾难。1883年，也就是表兄去世一年后，高尔顿给这场运动取了一个名字：优生学。
\end{quote}

\begin{quote}
\noindent\textbf{原文} Eugenics was only one of Galton's many interests; Galton enthusiasts refer to him as a polymath, detractors as a dilettante. In fact, he made significant contributions to geography, anthropology, psychology, genetics, meteorology, statistics, and, by setting fingerprint analysis on a sound scientific footing, to criminology. Born in 1822 into a prosperous family, his education---partly in medicine and partly in mathematics---was mostly a chronicle of defeated expectations. The death of his father when he was twenty-one simultaneously freed him from paternal restraint and yielded a handsome inheritance; the young man duly took advantage of both. After a full six years of being what might be described today as a trust-fund dropout, however, Galton settled down to become a productive member of the Victorian establishment. He made his name leading an expedition to a then little-known region of southwest Africa in 1850--1852. In his account of his explorations, we encounter the first instance of the one strand that connects his many varied interests: he counted and measured everything. Galton was only happy when he could reduce a phenomenon to a set of numbers.

\smallskip
\noindent\textbf{译文} 优生学只是高尔顿众多兴趣之一；崇拜他的人称他为博学多才者，批评者则称他为浅尝辄止的业余爱好者。事实上，他对地理学、人类学、心理学、遗传学、气象学、统计学都有重要贡献，并且通过把指纹分析建立在稳固的科学基础上，也贡献于犯罪学。高尔顿1822年出生于富裕家庭，他的教育部分是医学、部分是数学，却多半是一连串落空期望的记录。他二十一岁时父亲去世，这既使他摆脱了父亲约束，也给了他一笔可观遗产；年轻的高尔顿自然充分利用了二者。不过，在整整六年过着今天可称为“信托基金辍学者”的生活后，他安定下来，成为维多利亚体制中富有产出的一员。1850--1852年，他率领探险队前往当时鲜为人知的非洲西南部，并由此成名。在他对探险的记述中，我们首次看到贯穿他各种兴趣的一条线索：他对一切都计数和测量。只有当他能把一个现象化约为一组数字时，高尔顿才会高兴。
\end{quote}

\begin{quote}
\noindent\textbf{原文} At a missionary station he encountered a striking specimen of steatopygia---a condition of particularly protuberant buttocks, common among the indigenous Nama women of the region---and realized that this woman was naturally endowed with the figure that was then fashionable in Europe. The only difference was that it required enormous and costly ingenuity on the part of European dressmakers to create the desired "look" for their clients.

\smallskip
\noindent\textbf{译文} 在一个传教站，他遇到了一位突出的脂臀体征者。脂臀是臀部特别隆起的状态，在当地土著纳马女性中很常见。高尔顿意识到，这位女性天生拥有当时欧洲流行的体形。唯一不同的是，欧洲裁缝若要为顾客制造这种理想“外观”，需要付出巨大而昂贵的巧思。
\end{quote}

\begin{quote}
\noindent\textbf{English quotation.} I profess to be a scientific man, and was exceedingly anxious to obtain accurate measurements of her shape; but there was a difficulty in doing this. I did not know a word of Hottentot [the Dutch name for the Nama], and could never therefore have explained to the lady what the object of my footrule could be; and I really dared not ask my worthy missionary host to interpret for me. I therefore felt in a dilemma as I gazed at her form, that gift of bounteous nature to this favoured race, which no mantua-maker, with all her crinoline and stuffing, can do otherwise than humbly imitate. The object of my admiration stood under a tree, and was turning herself about to all points of the compass, as ladies who wish to be admired usually do. Of a sudden my eye fell upon my sextant; the bright thought struck me, and I took a series of observations upon her figure in every direction, up and down, crossways, diagonally, and so forth, and I registered them carefully upon an outline drawing for fear of any mistake; this being done, I boldly pulled out my measuring tape, and measured the distance from where I was to the place she stood, and having thus obtained both base and angles, I worked out the results by trigonometry and logarithms.

\smallskip
\noindent\textbf{中文引文。} 我自称是科学人，极想取得她体形的准确测量；但这样做有困难。我一句霍屯督语（纳马人的荷兰称呼）也不会，因此无从向这位女士解释我的英尺尺究竟有什么目的；我也实在不敢请我可敬的传教士主人替我翻译。于是我望着她的形体陷入两难：那是慷慨自然赐予这个受惠种族的礼物，任何女装裁缝凭借裙撑和填充物，都只能谦卑模仿。我的赞赏对象站在一棵树下，像希望被欣赏的女士通常会做的那样，向罗盘的各个方向转身。突然我的目光落在六分仪上；一个绝妙念头击中了我，于是我从各个方向对她的身形做了一系列观测，上下、横向、斜向，等等，并小心地记录在轮廓图上，以免出错。做完这些，我大胆拿出卷尺，测量自己所在位置到她站立处的距离；这样取得底边和角度后，我便用三角学和对数算出了结果。
\end{quote}

\begin{figure}[htbp]
\centering
\safeincludeimage{ab6c7a6dc9c3960917faef512bf665aaa8b9de182031bc523ecdd38196dd229c.jpg}
\caption{A nineteenth-century exaggerated view of a Nama woman.\\十九世纪对一名纳马女性的夸张描绘。}
\end{figure}

\begin{quote}
\noindent\textbf{原文} Galton's passion for quantification resulted in his developing many of the fundamental principles of modern statistics. It also yielded some clever observations. For example, he tested the efficacy of prayer. He figured that if prayer worked, those most prayed for should be at an advantage; to test the hypothesis he studied the longevity of British monarchs. Every Sunday, congregations in the Church of England following the \emph{Book of Common Prayer} beseeched God to "Endue the king/queen plenteously with heavenly gifts; Grant him/her in health and wealth long to live." Surely, Galton reasoned, the cumulative effect of all those prayers should be beneficial. In fact, prayer seemed ineffectual: he found that on average the monarchs died somewhat younger than other members of the British aristocracy.

\smallskip
\noindent\textbf{译文} 高尔顿对量化的热情，使他发展出现代统计学许多基本原则。它也产生了一些聪明观察。比如，他检验了祈祷的效力。他推想，如果祈祷有效，那么被祈祷最多的人应当有优势；为检验这一假说，他研究了英国君主的寿命。每个星期天，英国国教会众按照《公祷书》恳求上帝：“把丰盛的天赐礼物赋予国王/女王；赐予他/她健康与财富，使其长寿。”高尔顿推理道，所有这些祈祷的累积效果当然应该有益。事实上，祈祷似乎无效：他发现君主平均死亡年龄比英国贵族其他成员还略低。
\end{quote}

\begin{quote}
\noindent\textbf{原文} Because of the Darwin connection---their common grandfather, Erasmus Darwin, too, was one of the intellectual giants of his day---Galton was especially sensitive to the way in which certain lineages seemed to spawn disproportionately large numbers of prominent and successful people. In 1869 he published what would become the underpinning of all his ideas on eugenics, a treatise called \emph{Hereditary Genius: An Inquiry into Its Laws and Consequences}. In it he purported to show that talent, like simple genetic traits such as the Hapsburg Lip, does indeed run in families; he recounted, for example, how some families had produced generation after generation of judges. His analysis largely neglected to take into account the effect of the environment: the son of a prominent judge is, after all, rather more likely to become a judge---by virtue of his father's connections, if nothing else---than the son of a peasant farmer. Galton did not, however, completely overlook the effect of the environment, and it was he who first referred to the "nature/nurture" dichotomy, possibly in reference to Shakespeare's irredeemable villain, Caliban, "a devil, a born devil, on whose nature/Nurture can never stick."

\smallskip
\noindent\textbf{译文} 由于与达尔文的关系，他们共同的祖父伊拉斯谟·达尔文也曾是当时代的思想巨人之一，高尔顿对某些家系似乎不成比例地产生大量杰出成功人物这一现象格外敏感。1869年，他发表了后来成为其全部优生学思想基础的著作《遗传的天才：其规律及后果研究》。在书中，他声称证明了才能确实像哈布斯堡唇这类简单遗传性状一样，在家族中流传；例如，他讲述某些家族如何一代又一代产生法官。他的分析在很大程度上忽略了环境影响：毕竟，一位显赫法官的儿子仅凭父亲的人脉，就比农民的儿子更可能成为法官。不过，高尔顿并没有完全忽略环境作用，也是他首先提出了“先天/后天”二分，也许是在回应莎士比亚笔下不可救赎的恶人凯列班：“一个魔鬼，天生的魔鬼，教养永远无法附着于其本性。”
\end{quote}

\begin{quote}
\noindent\textbf{原文} The results of his analysis, however, left no doubt in Galton's mind.

\smallskip
\noindent\textbf{译文} 然而，他的分析结果在高尔顿心中没有留下任何疑问。
\end{quote}

\begin{quote}
\noindent\textbf{English quotation.} I have no patience with the hypothesis occasionally expressed, and often implied, especially in tales written to teach children to be good, that babies are born pretty much alike, and that the sole agencies in creating differences between boy and boy, and man and man, are steady application and moral effort. It is in the most unqualified manner that I object to pretensions of natural equality.

\smallskip
\noindent\textbf{中文引文。} 我不能容忍这样一种假说，它偶尔被明说，常常被暗示，尤其出现在教孩子做好人的故事中：婴儿出生时大体相同，而造成男孩与男孩、男人与男人之间差异的唯一因素，是持续努力和道德奋斗。我以最毫无保留的方式反对所谓自然平等的主张。
\end{quote}

\begin{quote}
\noindent\textbf{原文} A corollary of his conviction that these traits are genetically determined, he argued, was that it would be possible to "improve" the human stock by preferentially breeding gifted individuals, and preventing the less gifted from reproducing.

\smallskip
\noindent\textbf{译文} 他认为，既然这些特征由基因决定，一个推论就是：通过优先让有天赋者繁殖，并阻止天赋较低者繁殖，就可能“改良”人类种群。
\end{quote}

\begin{quote}
\noindent\textbf{English quotation.} It is easy \(\ldots\) to obtain by careful selection a permanent breed of dogs or horses gifted with peculiar powers of running, or of doing anything else, so it would be quite practicable to produce a highly gifted race of men by judicious marriages during several consecutive generations.

\smallskip
\noindent\textbf{中文引文。} 通过仔细选择，很容易获得一个具有特殊奔跑能力或其他能力的稳定犬种或马种；因此，在连续若干代中通过明智婚配，培育出一个高度天赋的人类种族，也完全可行。
\end{quote}

\begin{quote}
\noindent\textbf{原文} Galton introduced the term eugenics, literally "good in birth," to describe this application of the basic principle of agricultural breeding to humans. In time, eugenics came to refer to "self-directed human evolution": by making conscious choices about who should have children, eugenicists believed that they could head off the "eugenic crisis" precipitated in the Victorian imagination by the high rates of reproduction of inferior stock coupled with the typically small families of the superior middle classes.

\smallskip
\noindent\textbf{译文} 高尔顿引入“优生学”一词，字面意思是“出生良好”，用来描述把农业育种基本原则应用于人类。随着时间推移，优生学开始指“自我导向的人类进化”：通过有意识地选择谁应当生育，优生学者相信，他们能够阻止维多利亚想象中由劣等血统高繁殖率与优秀中产阶级家庭规模通常较小共同造成的“优生危机”。
\end{quote}

\begin{center}
\([\ldots]\)
\end{center}

\begin{figure}[htbp]
\centering
\safeincludeimage{5f5c2f9beaea0e59a190e44bd43fb9b305d148cb9108bddcf414f9951c876fb4.jpg}
\caption{Eugenics as it was perceived during the first part of the twentieth century: an opportunity for humans to control their own evolutionary destiny.\\二十世纪上半叶人们眼中的优生学：人类控制自身进化命运的机会。}
\end{figure}

\section*{The Double Helix: This Is Life}
\begin{center}
\large 双螺旋：这就是生命
\end{center}

\begin{quote}
\noindent\textbf{原文} I got hooked on the gene during my third year at the University of Chicago. Until then, I had planned to be a naturalist and looked forward to a career far removed from the urban bustle of Chicago's South Side, where I grew up. My change of heart was inspired not by an unforgettable teacher but a little book that appeared in 1944, \emph{What Is Life?}, by the Austrian-born father of wave mechanics, Erwin Schrödinger. It grew out of several lectures he had given the year before at the Institute for Advanced Study in Dublin. That a great physicist had taken the time to write about biology caught my fancy. In those days, like most people, I considered chemistry and physics to be the "real" sciences, and theoretical physicists were science's top dogs.

\smallskip
\noindent\textbf{译文} 我在芝加哥大学三年级时迷上了基因。在那之前，我原计划成为一名博物学家，期待一种远离芝加哥南区都市喧嚣的职业生涯；我正是在那里长大的。让我改变心意的不是一位令人难忘的老师，而是一本1944年出版的小书《生命是什么？》，作者是奥地利出生的波动力学之父埃尔温·薛定谔。此书源自他前一年在都柏林高等研究院所作的几场讲座。一位伟大的物理学家愿意花时间写生物学，这件事吸引了我。那时候，像大多数人一样，我认为化学和物理才是“真正的”科学，而理论物理学家则是科学界的顶尖人物。
\end{quote}

\begin{quote}
\noindent\textbf{原文} Schrödinger argued that life could be thought of in terms of storing and passing on biological information. Chromosomes were thus simply information bearers. Because so much information had to be packed into every cell, it must be compressed into what Schrödinger called a "hereditary code-script" embedded in the molecular fabric of chromosomes. To understand life, then, we would have to identify these molecules, and crack their code. He even speculated that understanding life---which would involve finding the gene---might take us beyond the laws of physics as we then understood them. Schrödinger's book was tremendously influential. Many of those who would become major players in Act 1 of molecular biology's great drama, including Francis Crick, a former physicist himself, had, like me, read \emph{What Is Life?} and been impressed.

\smallskip
\noindent\textbf{译文} 薛定谔认为，可以从储存和传递生物信息的角度来理解生命。因此，染色体不过是信息载体。由于每个细胞必须容纳大量信息，这些信息必然被压缩进薛定谔所谓的“遗传密码脚本”中，并嵌入染色体的分子结构里。要理解生命，我们就必须识别这些分子，并破解它们的密码。他甚至推测，理解生命，也就是寻找基因，可能会把我们带出现有物理定律的范围。薛定谔的书影响极大。后来成为分子生物学伟大戏剧第一幕中主要人物的许多人，包括弗朗西斯·克里克这位前物理学家，都像我一样读过《生命是什么？》并深受触动。
\end{quote}

\begin{quote}
\noindent\textbf{原文} In my own case, Schrödinger struck a chord because I too was intrigued by the essence of life. A small minority of scientists still thought life depended upon a vital force emanating from an all-powerful god. But like most of my teachers, I disdained the very idea of vitalism. If such a "vital" force were calling the shots in nature's game, there was little hope life would ever be understood through the methods of science. On the other hand, the notion that life might be perpetuated by means of an instruction book inscribed in a secret code appealed to me. What sort of molecular code could be so elaborate as to convey all the multitudinous wonder of the living world? And what sort of molecular trick could ensure that the code is exactly copied every time a chromosome duplicates?

\smallskip
\noindent\textbf{译文} 对我来说，薛定谔之所以拨动心弦，是因为我也被生命的本质所吸引。少数科学家仍认为，生命依赖一种来自全能上帝的生命力。但像我的多数老师一样，我轻视这种活力论观念。如果这样一种“生命”力量主宰自然游戏，那么通过科学方法理解生命就几乎没有希望。另一方面，生命也许依靠一本用秘密密码写成的说明书得以延续，这个想法吸引了我。什么样的分子密码能够复杂到传达生物世界如此繁多的奇妙？又是什么样的分子技巧，能够保证每次染色体复制时，密码都被精确复制？
\end{quote}

\begin{quote}
\noindent\textbf{原文} At the time of Schrödinger's Dublin lectures, most biologists supposed that proteins would eventually be identified as the primary bearers of genetic instruction. Proteins are molecular chains built up from twenty different building blocks, the amino acids. Because permutations in the order of amino acids along the chain are virtually infinite, proteins could, in principle, readily encode the information underpinning life's extraordinary diversity. DNA then was not considered a serious candidate for the bearer of code-scripts, even though it was exclusively located on chromosomes and had been known about for some seventy-five years. In 1869, Friedrich Miescher, a Swiss biochemist working in Germany, had isolated from pus-soaked bandages supplied by a local hospital a substance he called "nuclein." Because pus consists largely of white blood cells, which, unlike red blood cells, have nuclei and therefore DNA-containing chromosomes, Miescher had stumbled on a good source of DNA. When he later discovered that "nuclein" was to be found in chromosomes alone, Miescher understood that his discovery was indeed a big one. In 1893, he wrote: "Inheritance insures a continuity in form from generation to generation that lies even deeper than the chemical molecule. It lies in the structuring atomic groups. In this sense, I am a supporter of the chemical heredity theory."

\smallskip
\noindent\textbf{译文} 在薛定谔都柏林讲座的时代，多数生物学家认为，蛋白质最终会被确认是遗传指令的主要承载者。蛋白质是由二十种不同构件，即氨基酸，组成的分子链。由于氨基酸在链上排列顺序的组合实际上无穷无尽，原则上蛋白质很容易编码支撑生命非凡多样性的信息。那时，DNA尚未被认真视为密码脚本的承载候选者，尽管它专门位于染色体上，并且已为人所知约七十五年。1869年，在德国工作的瑞士生物化学家弗里德里希·米歇尔，从当地医院提供的浸满脓液的绷带中分离出一种他称为“核素”的物质。脓液主要由白细胞构成；白细胞不同于红细胞，有细胞核，因此含有带DNA的染色体。米歇尔偶然找到了一个很好的DNA来源。后来他发现“核素”只存在于染色体中，便明白自己的发现确实意义重大。1893年，他写道：“遗传保证了形态从一代到下一代的连续性，这种连续性比化学分子更深；它存在于构成结构的原子团之中。从这个意义上说，我是化学遗传理论的支持者。”
\end{quote}

\begin{figure}[htbp]
\centering
\safeincludeimage{283af47520fa3108f7a0662a80b6c7e627dddf8a7e270d4eec97eebd5d0967f8.jpg}
\caption{The physicist Erwin Schrödinger, whose book \emph{What Is Life?} turned me on to the gene.\\物理学家埃尔温·薛定谔，他的《生命是什么？》让我迷上了基因。}
\end{figure}

\begin{quote}
\noindent\textbf{原文} Nevertheless, for decades afterward, chemistry would remain unequal to the task of analyzing the immense size and complexity of the DNA molecule. Only in the 1930s was DNA shown to be a long molecule containing four different chemical bases: adenine (A), guanine (G), thymine (T), and cytosine (C). But at the time of Schrödinger's lectures, it was still unclear just how the subunits, called deoxynucleotides, of the molecule were chemically linked. Nor was it known whether DNA molecules might vary in their sequences of the four different bases. If DNA were indeed Schrödinger's code-script, then the molecule would have to be capable of existing in an immense number of different forms. But back then it was still considered a possibility that one simple sequence like AGTC might be repeated over and over along the entire length of DNA chains.

\smallskip
\noindent\textbf{译文} 尽管如此，此后数十年，化学仍无力分析DNA分子巨大的尺寸和复杂性。直到二十世纪三十年代，人们才证明DNA是一种长分子，含有四种不同化学碱基：腺嘌呤（A）、鸟嘌呤（G）、胸腺嘧啶（T）和胞嘧啶（C）。但在薛定谔讲座时，人们仍不清楚这种分子的亚单位，即脱氧核苷酸，究竟如何以化学方式连接。人们也不知道DNA分子在四种不同碱基的序列上是否可能变化。如果DNA确实是薛定谔所谓的密码脚本，那么这种分子就必须能够以数量巨大的不同形式存在。但当时仍有人认为，像AGTC这样简单的序列也许会沿DNA链全长一遍遍重复。
\end{quote}

\begin{quote}
\noindent\textbf{原文} DNA did not move into the genetic limelight until 1944, when Oswald Avery's lab at the Rockefeller Institute in New York City reported that the composition of the surface coats of pneumonia bacteria could be changed. This was not the result he and his junior colleagues, Colin MacLeod and Maclyn McCarty, expected.

\smallskip
\noindent\textbf{译文} 直到1944年，DNA才进入遗传学聚光灯。当时纽约市洛克菲勒研究所奥斯瓦尔德·艾弗里的实验室报告说，肺炎细菌表面荚膜的组成可以被改变。这并不是他和年轻同事科林·麦克劳德、麦克林·麦卡蒂预期的结果。
\end{quote}

\begin{quote}
\noindent\textbf{原文} For more than a decade Avery's group had been following up on another most unexpected observation made in 1928 by Fred Griffith, a scientist in the British Ministry of Health. Griffith was interested in pneumonia and studied its bacterial agent, \emph{Pneumococcus}. It was known that there were two strains, designated "smooth" (S) and "rough" (R) according to their appearance under the microscope. These strains differed not only visually but also in their virulence. Inject S bacteria into a mouse, and within a few days the mouse dies; inject R bacteria and the mouse remains healthy. It turns out that S bacterial cells have a coating that prevents the mouse's immune system from recognizing the invader. The R cells have no such coating and are therefore readily attacked by the mouse's immune defenses.

\smallskip
\noindent\textbf{译文} 十多年来，艾弗里团队一直在追踪英国卫生部科学家弗雷德·格里菲思1928年的另一个极其意外的观察。格里菲思关注肺炎，并研究其致病细菌肺炎球菌。人们已知这种细菌有两个菌株，根据显微镜下外观分别称为“光滑型”（S）和“粗糙型”（R）。这些菌株不仅外观不同，毒力也不同。把S型细菌注入小鼠，几天内小鼠死亡；注入R型细菌，小鼠仍保持健康。原来S型细胞有一层外衣，可以阻止小鼠免疫系统识别入侵者。R型细胞没有这层外衣，因此很容易被小鼠免疫防御攻击。
\end{quote}

\begin{figure}[htbp]
\centering
\safeincludeimage{0edbfbcce31b2e08b7233f24242c1aaad9c5441d31125ba58e18200cb469a047.jpg}
\caption{A view through the microscope of blood cells treated with a chemical that stains DNA. In order to maximize their oxygen-transporting capacity, red blood cells have no nucleus and therefore no DNA. But white blood cells, which patrol the bloodstream in search of intruders, have a nucleus containing chromosomes.\\显微镜下用DNA染色化学物质处理过的血细胞。为了最大限度提高输氧能力，红细胞没有细胞核，因此没有DNA。但在血流中巡逻、寻找入侵者的白细胞有细胞核，其中含有染色体。}
\end{figure}

\begin{quote}
\noindent\textbf{原文} Through his involvement with public health, Griffith knew that multiple strains had sometimes been isolated from a single patient, and so he was curious about how different strains might interact in his unfortunate mice. With one combination, he made a remarkable discovery: when he injected heat-killed S bacteria, harmless, and normal R bacteria, also harmless, the mouse died. How could two harmless forms of bacteria conspire to become lethal? The clue came when he isolated the \emph{Pneumococcus} bacteria retrieved from the dead mice and discovered living S bacteria. It appeared the living innocuous R bacteria had acquired something from the dead S variant; whatever it was, that something had allowed the R in the presence of the heat-killed S bacteria to transform itself into a living killer S strain. Griffith confirmed that this change was for real by culturing the S bacteria from the dead mouse over several generations: the bacteria bred true for the S type, just as any regular S strain would. A genetic change had indeed occurred to the R bacteria injected into the mouse.

\smallskip
\noindent\textbf{译文} 由于参与公共卫生工作，格里菲思知道有时会从同一名患者体内分离出多个菌株，因此他很好奇不同菌株在他那些不幸的小鼠体内会如何相互作用。在一种组合中，他有了惊人发现：当他注入热杀死的S型细菌（无害）和正常R型细菌（也无害）时，小鼠却死了。两种无害细菌怎么会合谋致命？线索来自他从死亡小鼠体内分离肺炎球菌时发现了活的S型细菌。看来，活的无害R型细菌从死去的S型变体那里获得了某种东西；不管那是什么，它使R型在热杀死的S型细菌存在时，转化成了活的致命S型菌株。格里菲思通过把死鼠体内的S型细菌连续培养几代，确认这种改变是真实的：这些细菌像任何普通S型菌株一样稳定繁殖。注入小鼠的R型细菌确实发生了遗传改变。
\end{quote}

\begin{quote}
\noindent\textbf{原文} Though this transformation phenomenon seemed to defy all understanding, Griffith's observations at first created little stir in the scientific world. This was partly because Griffith was intensely private and so averse to large gatherings that he seldom attended scientific conferences. Once, he had to be virtually forced to give a lecture. Bundled into a taxi and escorted to the hall by colleagues, he discoursed in a mumbled monotone, emphasizing an obscure corner of his microbiological work but making no mention of bacterial transformation. Luckily, however, not everyone overlooked Griffith's breakthrough.

\smallskip
\noindent\textbf{译文} 虽然这种转化现象看似完全不可理解，格里菲思的观察最初并未在科学界引起多大轰动。部分原因是格里菲思极其重视隐私，又非常厌恶大型聚会，几乎不参加科学会议。有一次，他几乎是被迫去做演讲的。同事把他塞进出租车并护送到会场后，他用含糊单调的声音讲述自己微生物学工作中一个冷僻角落，却丝毫没有提及细菌转化。不过幸运的是，并非所有人都忽视了格里菲思的突破。
\end{quote}

\begin{quote}
\noindent\textbf{原文} Oswald Avery was also interested in the sugarlike coats of the \emph{Pneumococcus}. He set out to duplicate Griffith's experiment in order to isolate and characterize whatever it was that had caused those R cells to change to the S type. In 1944 Avery, MacLeod, and McCarty published their results: an exquisite set of experiments showing unequivocally that DNA was the transforming principle. Culturing the bacteria in the test tube rather than in mice made it much easier to search for the chemical identity of the transforming factor in the heat-killed S cells. Methodically destroying one by one the biochemical components of the heat-treated S cells, Avery and his group looked to see whether transformation was prevented. First they degraded the sugarlike coat of the S bacteria. Transformation still occurred: the coat was not the transforming principle. Next they used a mixture of two protein-destroying enzymes, trypsin and chymotrypsin, to degrade virtually all the proteins in the S cells. To their surprise, transformation was again unaffected. Next they tried an enzyme, RNase, that breaks down RNA, ribonucleic acid, a second class of nucleic acids similar to DNA and possibly involved in protein synthesis. Again transformation occurred. Finally, they came to DNA, exposing the S bacterial extracts to the DNA-destroying enzyme, DNase. This time they hit a home run. All S-inducing activity ceased completely. The transforming factor was DNA.

\smallskip
\noindent\textbf{译文} 奥斯瓦尔德·艾弗里也对肺炎球菌的糖样荚膜感兴趣。他开始重复格里菲思的实验，以分离并表征究竟是什么导致R型细胞变成S型。1944年，艾弗里、麦克劳德和麦卡蒂发表结果：一组精巧实验明确显示，DNA就是转化原则。与其在小鼠体内培养细菌，不如在试管中培养，这使他们更容易寻找热杀死S型细胞中转化因子的化学身份。艾弗里团队有条不紊地逐一破坏热处理S型细胞的生化成分，再观察转化是否被阻止。首先，他们降解S型细菌的糖样荚膜。转化仍然发生：荚膜不是转化原则。接着，他们使用两种破坏蛋白质的酶，胰蛋白酶和胰凝乳蛋白酶，几乎降解了S型细胞中所有蛋白质。令他们惊讶的是，转化仍不受影响。随后他们尝试RNase，这是一种分解RNA（核糖核酸）的酶；RNA是类似DNA的第二类核酸，可能参与蛋白质合成。转化又一次发生。最后，他们轮到DNA，把S型细菌提取物暴露给破坏DNA的酶DNase。这一次他们击中了靶心。所有诱导S型的活性完全停止。转化因子就是DNA。
\end{quote}

\begin{quote}
\noindent\textbf{原文} In part because of its bombshell implications, the resulting February 1944 paper by Avery, MacLeod, and McCarty met with a mixed response. Many geneticists accepted their conclusions. After all, DNA was found on every chromosome; why should it not be the genetic material? By contrast, however, most biochemists expressed doubt that DNA was a complex enough molecule to act as the repository of such a vast quantity of biological information. They continued to believe that proteins, the other component of chromosomes, would prove to be the hereditary substance. In principle, as the biochemists rightly noted, it would be much easier to encode a vast body of complex information using the twenty-letter amino-acid alphabet of proteins than the four-letter nucleotide alphabet of DNA. Particularly vitriolic in his rejection of DNA as the genetic substance was Avery's own colleague at the Rockefeller Institute, the protein chemist Alfred Mirsky. By then, however, Avery was no longer scientifically active. The Rockefeller Institute had mandatorily retired him at age sixty-five.

\smallskip
\noindent\textbf{译文} 部分由于其含义像炸弹一样震撼，艾弗里、麦克劳德和麦卡蒂1944年2月发表的论文反响不一。许多遗传学家接受了他们的结论。毕竟，每条染色体上都有DNA；为什么它不能是遗传物质？相比之下，多数生物化学家怀疑DNA分子是否足够复杂，能作为如此巨大生物信息量的储存库。他们继续相信，染色体的另一种成分蛋白质才会被证明是遗传物质。原则上，正如生物化学家正确指出的，用蛋白质二十个字母的氨基酸字母表编码大量复杂信息，要比用DNA四个字母的核苷酸字母表容易得多。在否认DNA是遗传物质方面，尤其尖刻的是艾弗里在洛克菲勒研究所的同事、蛋白质化学家阿尔弗雷德·米尔斯基。然而，到那时艾弗里已经不再活跃于科学研究。洛克菲勒研究所强制他在六十五岁退休。
\end{quote}

\begin{quote}
\noindent\textbf{原文} Avery missed out on more than the opportunity to defend his work against the attacks of his colleagues: he was never awarded the Nobel Prize, which was certainly his due, for identifying DNA as the transforming principle. Because the Nobel committee makes its records public fifty years following each award, we now know that Avery's candidacy was blocked by the Swedish physical chemist Einar Hammarsten. Though Hammarsten's reputation was based largely on his having produced DNA samples of unprecedented high quality, he still believed genes to be an undiscovered class of proteins. In fact, even after the double helix was found, Hammarsten continued to insist that Avery should not receive the prize until after the mechanism of DNA transformation had been completely worked out. Avery died in 1955; had he lived only a few more years, he would almost certainly have gotten the prize.

\smallskip
\noindent\textbf{译文} 艾弗里错过的不只是为自己的工作抵御同事攻击的机会。由于识别出DNA是转化原则，他理应获得诺贝尔奖，却从未获奖。诺贝尔委员会会在每次授奖五十年后公开档案，因此我们现在知道，艾弗里的候选资格被瑞典物理化学家埃纳尔·哈马斯滕阻止了。哈马斯滕的声誉主要建立在他制备出前所未有高质量DNA样品之上，但他仍相信基因是一类尚未发现的蛋白质。事实上，即便双螺旋被发现以后，哈马斯滕仍坚持认为，在DNA转化机制被完全弄清楚之前，不应把奖授予艾弗里。艾弗里1955年去世；如果他再多活几年，几乎肯定会获奖。
\end{quote}

\begin{quote}
\noindent\textbf{原文} When I arrived at Indiana University in the fall of 1947 with plans to pursue the gene for my Ph.D. thesis, Avery's paper came up over and over in conversations. By then, no one doubted the reproducibility of his results, and more recent work coming out of the Rockefeller Institute made it all the less likely that proteins would prove to be the genetic actors in bacterial transformation. DNA had at last become an important objective for chemists setting their sights on the next breakthrough. In Cambridge, England, the canny Scottish chemist Alexander Todd rose to the challenge of identifying the chemical bonds that linked together nucleotides in DNA. By early 1951, his lab had proved that these links were always the same, such that the backbone of the DNA molecule was very regular. During the same period, the Austrian-born refugee Erwin Chargaff, at the College of Physicians and Surgeons of Columbia University, used the new technique of paper chromatography to measure the relative amounts of the four DNA bases in DNA samples extracted from a variety of vertebrates and bacteria. While some species had DNA in which adenine and thymine predominated, others had DNA with more guanine and cytosine. The possibility thus presented itself that no two DNA molecules had the same composition.

\smallskip
\noindent\textbf{译文} 1947年秋天，我来到印第安纳大学，计划以基因为博士论文主题，艾弗里的论文在谈话中一次又一次被提起。到那时，已经没有人怀疑他的结果可以重复；洛克菲勒研究所新近的工作也使蛋白质会被证明是细菌转化中遗传行动者的可能性越来越小。DNA终于成为化学家瞄准下一次突破时的重要目标。在英国剑桥，精明的苏格兰化学家亚历山大·托德迎接了挑战：识别把DNA中核苷酸连接起来的化学键。到1951年初，他的实验室已经证明这些连接始终相同，因此DNA分子的骨架非常规则。同一时期，出生于奥地利的难民埃尔温·查伽夫在哥伦比亚大学内外科医学院，利用纸层析新技术，测量从多种脊椎动物和细菌中提取的DNA样品里四种DNA碱基的相对含量。有些物种的DNA中腺嘌呤和胸腺嘧啶占优势，另一些物种的DNA则含有更多鸟嘌呤和胞嘧啶。由此出现了一种可能：没有两个DNA分子的组成完全相同。
\end{quote}

\begin{quote}
\noindent\textbf{原文} At Indiana I joined a small group of visionary scientists, mostly physicists and chemists, studying the reproductive process of the viruses that attack bacteria, bacteriophages---"phages" for short. The Phage Group was born when my Ph.D. supervisor, the Italian-trained medic Salvador Luria, and his close friend, the German-born theoretical physicist Max Delbrück, teamed up with the American physical chemist Alfred Hershey. During World War II both Luria and Delbrück were considered enemy aliens, and thus ineligible to serve in the war effort of American science, even though Luria, a Jew, had been forced to leave France for New York City and Delbrück had fled Germany as an objector to Nazism. Thus excluded, they continued to work in their respective university labs---Luria at Indiana and Delbrück at Vanderbilt---and collaborated on phage experiments during successive summers at Cold Spring Harbor. In 1943, they joined forces with the brilliant but taciturn Hershey, then doing phage research of his own at Washington University in St. Louis.

\smallskip
\noindent\textbf{译文} 在印第安纳，我加入了一小群有远见的科学家，他们多半是物理学家和化学家，研究攻击细菌的病毒，即噬菌体，简称“噬菌体”的繁殖过程。噬菌体小组诞生于我的博士导师、受意大利训练的医生萨尔瓦多·卢里亚，与他的密友、德国出生的理论物理学家马克斯·德尔布吕克，以及美国物理化学家阿尔弗雷德·赫尔希联合之时。第二次世界大战期间，卢里亚和德尔布吕克都被视为敌国侨民，因此没有资格服务于美国科学的战时努力，尽管身为犹太人的卢里亚被迫从法国逃往纽约，德尔布吕克则因反对纳粹主义而逃离德国。于是，被排除在外的他们继续在各自大学实验室工作，卢里亚在印第安纳，德尔布吕克在范德比尔特，并在连续几个夏天于冷泉港合作开展噬菌体实验。1943年，他们与才华横溢但沉默寡言的赫尔希联手；赫尔希当时正在圣路易斯华盛顿大学从事自己的噬菌体研究。
\end{quote}

\begin{quote}
\noindent\textbf{原文} The Phage Group's program was based on its belief that phages, like all viruses, were in effect naked genes. This concept had first been proposed in 1922 by the imaginative American geneticist Herman J. Muller, who three years later demonstrated that X rays cause mutations. His belated Nobel Prize came in 1946, just after he joined the faculty of Indiana University. It was his presence, in fact, that led me to Indiana. Having started his career under T. H. Morgan, Muller knew better than anyone else how genetics had evolved during the first half of the twentieth century, and I was enthralled by his lectures during my first term. His work on fruit flies, \emph{Drosophila}, however, seemed to me to belong more to the past than to the future, and I only briefly considered doing thesis research under his supervision. I opted instead for Luria's phages, an even speedier experimental subject than \emph{Drosophila}: genetic crosses of phages done one day could be analyzed the next.

\smallskip
\noindent\textbf{译文} 噬菌体小组的计划建立在这样一种信念上：噬菌体像所有病毒一样，实际上是裸露的基因。这个概念最早由富有想象力的美国遗传学家赫尔曼·J. 穆勒在1922年提出；三年后，他证明X射线会导致突变。他迟来的诺贝尔奖在1946年到来，就在他加入印第安纳大学教职之后。事实上，正是他的存在把我带到了印第安纳。穆勒在T. H. 摩尔根门下开始职业生涯，他比任何人都更清楚遗传学在二十世纪上半叶如何发展；我第一学期听他的课时深深着迷。不过，他关于果蝇的工作在我看来更属于过去而不是未来，我只是短暂考虑过在他指导下做论文研究。我最终选择了卢里亚的噬菌体，它比果蝇还是更快的实验对象：一天完成的噬菌体遗传杂交，第二天就能分析。
\end{quote}

\begin{quote}
\noindent\textbf{原文} For my Ph.D. thesis research, Luria had me follow in his footsteps by studying how X rays killed phage particles. Initially I had hoped to show that viral death was caused by damage to phage DNA. Reluctantly, however, I eventually had to concede that my experimental approach could never give unambiguous answers at the chemical level. I could draw only biological conclusions. Even though phages were indeed effectively naked genes, I realized that the deep answers the Phage Group was seeking could be arrived at only through advanced chemistry. DNA somehow had to transcend its status as an acronym; it had to be understood as a molecular structure in all its chemical detail.

\smallskip
\noindent\textbf{译文} 我的博士论文研究中，卢里亚让我沿着他的道路，研究X射线如何杀死噬菌体颗粒。起初我希望证明，病毒死亡是由噬菌体DNA受损造成的。然而，我最终不得不勉强承认，我的实验方法永远无法在化学层面给出明确答案。我只能得出生物学结论。尽管噬菌体确实等同于裸露基因，我意识到，噬菌体小组寻求的深层答案只能通过高级化学获得。DNA必须以某种方式超越其缩写词地位；它必须被理解为一个具有全部化学细节的分子结构。
\end{quote}

\begin{quote}
\noindent\textbf{原文} Upon finishing my thesis, I saw no alternative but to move to a lab where I could study DNA chemistry. Unfortunately, however, knowing almost no pure chemistry, I would have been out of my depth in any lab attempting difficult experiments in organic or physical chemistry. I therefore took a postdoctoral fellowship in the Copenhagen lab of the biochemist Herman Kalckar in the fall of 1950. He was studying the synthesis of the small molecules that make up DNA, but I figured out quickly that his biochemical approach would never lead to an understanding of the essence of the gene. Every day spent in his lab would be one more day's delay in learning how DNA carried genetic information.

\smallskip
\noindent\textbf{译文} 完成论文后，我看不到别的选择，只能转到一个可以研究DNA化学的实验室。不幸的是，我几乎不懂纯化学，因此在任何尝试进行艰难有机化学或物理化学实验的实验室里都会力不从心。于是1950年秋，我接受了生物化学家赫尔曼·卡尔卡在哥本哈根实验室的博士后职位。他研究构成DNA的小分子的合成，但我很快看出，他的生化方法永远无法引向对基因本质的理解。在他的实验室里每多待一天，就会让学习DNA如何携带遗传信息的时间再推迟一天。
\end{quote}

\begin{quote}
\noindent\textbf{原文} My Copenhagen year nonetheless ended productively. To escape the cold Danish spring, I went to the Zoological Station at Naples during April and May. During my last week there, I attended a small conference on X-ray diffraction methods for determining the 3-D structure of molecules. X-ray diffraction is a way of studying the atomic structure of any molecule that can be crystallized. The crystal is bombarded with X rays, which bounce off its atoms and are scattered. The scatter pattern gives information about the structure of the molecule but, taken alone, is not enough to solve the structure. The additional information needed is the "phase assignment," which deals with the wave properties of the molecule. Solving the phase problem was not easy, and at that time only the most audacious scientists were willing to take it on. Most of the successes of the diffraction method had been achieved with relatively simple molecules.

\smallskip
\noindent\textbf{译文} 尽管如此，我在哥本哈根的那一年还是以富有成效的方式结束了。为了逃离丹麦寒冷的春天，我在四月和五月去了那不勒斯动物学站。在那里的最后一周，我参加了一个关于用X射线衍射方法测定分子三维结构的小型会议。X射线衍射是一种研究任何可结晶分子原子结构的方法。晶体受到X射线轰击，X射线从其原子上反弹并散射。散射图案提供分子结构信息，但单凭它还不足以解出结构。还需要额外信息，即“相位指定”，它涉及分子的波动性质。解决相位问题并不容易，在当时只有最大胆的科学家愿意承担。衍射法的大多数成功，都发生在相对简单的分子上。
\end{quote}

\begin{quote}
\noindent\textbf{原文} My expectations for the conference were low. I believed that a three-dimensional understanding of protein structure, or for that matter of DNA, was more than a decade away. Disappointing earlier X-ray photos suggested that DNA was particularly unlikely to yield up its secrets via the X-ray approach. These results were not surprising since the exact sequences of DNA were expected to differ from one individual molecule to another. The resulting irregularity of surface configurations would understandably prevent the long thin DNA chains from lying neatly side by side in the regular repeating patterns required for X-ray analysis to be successful.

\smallskip
\noindent\textbf{译文} 我对这次会议的期待很低。我认为，对蛋白质结构的三维理解，甚至对DNA的三维理解，至少还要十多年。早先令人失望的X射线照片提示，DNA尤其不可能通过X射线方法交出秘密。这些结果并不意外，因为人们预期不同DNA分子的精确序列会彼此不同。由此造成的表面构型不规则，可以理解地会阻止细长DNA链整齐并排排列成X射线分析成功所需的规则重复图案。
\end{quote}

\begin{quote}
\noindent\textbf{原文} It was therefore a surprise and a delight to hear the last-minute talk on DNA by a thirty-four-year-old Englishman named Maurice Wilkins from the Biophysics Lab of King's College, London. Wilkins was a physicist who during the war had worked on the Manhattan Project. For him, as for many of the other scientists involved, the actual deployment of the bomb on Hiroshima and Nagasaki, supposedly the culmination of all their work, was profoundly disillusioning.

\smallskip
\noindent\textbf{译文} 因此，听到伦敦国王学院生物物理实验室三十四岁的英国人莫里斯·威尔金斯临时作关于DNA的报告，既出乎意料又令人高兴。威尔金斯是物理学家，战争期间曾参与曼哈顿计划。对他来说，如同对许多参与其中的科学家一样，原本被视为全部工作巅峰的炸弹实际投向广岛和长崎，令人深感幻灭。
\end{quote}

\begin{figure}[htbp]
\centering
\safeincludeimage{a27a6630519567029c486e0a8709a2111b58555b107bf1f6aa5c6e55aa3f1623.jpg}
\caption{Maurice Wilkins in his lab at King's College, London.\\莫里斯·威尔金斯在伦敦国王学院的实验室。}
\end{figure}

\begin{quote}
\noindent\textbf{原文} He considered forsaking science altogether to become a painter in Paris, but biology intervened. He too had read Schrödinger's book, and was now tackling DNA with X-ray diffraction. He displayed a photograph of an X-ray diffraction pattern he had recently obtained, and its many precise reflections indicated a highly regular crystalline packing. DNA, one had to conclude, must have a regular structure, the elucidation of which might well reveal the nature of the gene. Instantly I saw myself moving to London to help Wilkins find the structure. My attempts to converse with him after his talk, however, went nowhere. All I got for my efforts was a declaration of his conviction that much hard work lay ahead.

\smallskip
\noindent\textbf{译文} 他曾考虑彻底放弃科学，去巴黎当画家，但生物学介入了。他也读过薛定谔的书，如今正用X射线衍射研究DNA。他展示了一张自己最近获得的X射线衍射图，许多精确反射表明其中存在高度规则的晶体堆积。人们不得不下结论：DNA一定有规则结构，而阐明这种结构很可能揭示基因的本质。我立刻想象自己搬到伦敦，帮助威尔金斯寻找结构。然而，他报告后我试图与他交谈，却毫无进展。我努力所得，只是他坚称前方还有大量艰苦工作。
\end{quote}

\begin{quote}
\noindent\textbf{原文} While I was hitting consecutive dead ends, back in America the world's preeminent chemist, Caltech's Linus Pauling, announced a major triumph: he had found the exact arrangement in which chains of amino acids, called polypeptides, fold up in proteins, and called his structure the \(\alpha\)-helix. That it was Pauling who made this breakthrough was no surprise: he was a scientific superstar. His book \emph{The Nature of the Chemical Bond} essentially laid the foundation of modern chemistry, and, for chemists of the day, it was the Bible. Pauling had been a precocious child. When he was nine, his father, a druggist in Oregon, wrote to the \emph{Oregonian} newspaper requesting suggestions of reading matter for his bookish son, adding that he had already read the Bible and Darwin's \emph{Origin of Species}. But the early death of Pauling's father, which brought the family to financial ruin, makes it remarkable that the promising young man managed to get an education at all.

\smallskip
\noindent\textbf{译文} 当我接连碰壁时，在美国，世界最杰出的化学家、加州理工学院的莱纳斯·鲍林宣布了一项重大胜利：他找到了氨基酸链，也称多肽链，在蛋白质中折叠的确切排列方式，并把这一结构称为 \(\alpha\)-螺旋。取得这一突破的人是鲍林并不令人意外：他是科学巨星。他的《化学键的本质》基本奠定了现代化学基础，对当时的化学家来说就是圣经。鲍林童年早慧。九岁时，他在俄勒冈当药剂师的父亲写信给《俄勒冈人报》，请他们为爱读书的儿子推荐读物，并补充说孩子已经读过《圣经》和达尔文的《物种起源》。但鲍林父亲早逝，使家庭陷入经济破产，这也让这个前途光明的年轻人竟能接受教育一事显得格外不易。
\end{quote}

\begin{quote}
\noindent\textbf{原文} As soon as I returned to Copenhagen I read about Pauling's \(\alpha\)-helix. To my surprise, his model was not based on a deductive leap from experimental X-ray diffraction data. Instead, it was Pauling's long experience as a structural chemist that had emboldened him to infer which type of helical fold would be most compatible with the underlying chemical features of the polypeptide chain. Pauling made scale models of the different parts of the protein molecule, working out plausible schemes in three dimensions. He had reduced the problem to a kind of three-dimensional jigsaw puzzle in a way that was simple yet brilliant.

\smallskip
\noindent\textbf{译文} 我一回到哥本哈根，就读到了鲍林的 \(\alpha\)-螺旋。令我惊讶的是，他的模型并不是基于从X射线衍射实验数据做出的演绎飞跃。相反，是鲍林作为结构化学家的长期经验，使他敢于推断哪一种螺旋折叠最能与多肽链底层化学特征相容。鲍林制作了蛋白质分子不同部分的比例模型，在三维空间中推敲可行方案。他把问题化约成一种三维拼图，方式简单却高明。
\end{quote}

\begin{quote}
\noindent\textbf{原文} Whether the \(\alpha\)-helix was correct---in addition to being pretty---was now the question. Only a week later, I got the answer. Sir Lawrence Bragg, the English inventor of X-ray crystallography and 1915 Nobel laureate in Physics, came to Copenhagen and excitedly reported that his junior colleague, the Austrian-born chemist Max Perutz, had ingeniously used synthetic polypeptides to confirm the correctness of Pauling's \(\alpha\)-helix. It was a bittersweet triumph for Bragg's Cavendish Laboratory. The year before, they had completely missed the boat in their paper outlining possible helical folds for polypeptide chains.

\smallskip
\noindent\textbf{译文} \(\alpha\)-螺旋除了漂亮之外是否正确，成了问题。仅一周后，我得到了答案。X射线晶体学的英国发明者、1915年诺贝尔物理学奖得主劳伦斯·布拉格爵士来到哥本哈根，兴奋地报告说，他的年轻同事、奥地利出生的化学家马克斯·佩鲁茨巧妙地利用合成多肽，证实了鲍林 \(\alpha\)-螺旋的正确性。对布拉格的卡文迪许实验室来说，这是一次苦乐参半的胜利。就在前一年，他们在一篇概述多肽链可能螺旋折叠的论文中完全错失了良机。
\end{quote}

\begin{figure}[htbp]
\centering
\safeincludeimage{214fa4b29d64e67c2fdb9a065099f912ba77e0f50d9c2c3553ba6b0c32ff357d.jpg}
\caption{Lawrence Bragg (left) with Linus Pauling, who is carrying a model of the \(\alpha\)-helix.\\劳伦斯·布拉格（左）与莱纳斯·鲍林；鲍林手中拿着 \(\alpha\)-螺旋模型。}
\end{figure}

\begin{quote}
\noindent\textbf{原文} By then Salvador Luria had tentatively arranged for me to take up a research position at the Cavendish. Located at Cambridge University, this was the most famous laboratory in all of science. Here Ernest Rutherford first described the structure of the atom. Now it was Bragg's own domain, and I was to work as apprentice to the English chemist John Kendrew, who was interested in determining the 3-D structure of the protein myoglobin. Luria advised me to visit the Cavendish as soon as possible. With Kendrew in the States, Max Perutz would check me out. Together, Kendrew and Perutz had earlier established the Medical Research Council (MRC) Unit for the Study of the Structure of Biological Systems.

\smallskip
\noindent\textbf{译文} 到那时，萨尔瓦多·卢里亚已初步安排我去卡文迪许实验室任研究职位。卡文迪许位于剑桥大学，是全科学界最著名的实验室。欧内斯特·卢瑟福最早就是在这里描述原子结构。如今这里是布拉格的地盘，而我将作为英国化学家约翰·肯德鲁的学徒工作；肯德鲁关心的是确定肌红蛋白的三维结构。卢里亚建议我尽快访问卡文迪许。肯德鲁在美国，所以马克斯·佩鲁茨会来考察我。肯德鲁和佩鲁茨早些时候共同建立了医学研究委员会（MRC）生物系统结构研究单元。
\end{quote}

\begin{quote}
\noindent\textbf{原文} A month later in Cambridge, Perutz assured me that I could quickly master the necessary X-ray diffraction theory and should have no difficulty fitting in with the others in their tiny MRC Unit. To my relief, he was not put off by my biology background. Nor was Lawrence Bragg, who briefly came down from his office to look me over.

\smallskip
\noindent\textbf{译文} 一个月后在剑桥，佩鲁茨向我保证，我可以很快掌握必要的X射线衍射理论，并且在他们那个小小的MRC单元里融入其他人不会有困难。令我松了一口气的是，他并没有被我的生物学背景吓退。劳伦斯·布拉格也没有；他短暂地从办公室下来打量了我一番。
\end{quote}

\begin{quote}
\noindent\textbf{原文} I was twenty-three when I arrived back at the MRC Unit in Cambridge in early October. I found myself sharing space in the biochemistry room with a thirty-five-year-old ex-physicist, Francis Crick, who had spent the war working on magnetic mines for the Admiralty. When the war ended, Crick had planned to stay on in military research, but, on reading Schrödinger's \emph{What Is Life?}, he had moved toward biology. Now he was at the Cavendish to pursue the 3-D structure of proteins for his Ph.D.

\smallskip
\noindent\textbf{译文} 十月初我回到剑桥MRC单元时二十三岁。我发现自己在生化室里与一位三十五岁的前物理学家弗朗西斯·克里克共用空间；战争期间，他为海军部研究磁性水雷。战争结束时，克里克本打算继续从事军事研究，但读了薛定谔的《生命是什么？》后，他转向了生物学。如今他在卡文迪许攻读博士，研究蛋白质三维结构。
\end{quote}

\begin{quote}
\noindent\textbf{原文} Crick was always fascinated by the intricacies of important problems. His endless questions as a child compelled his weary parents to buy him a children's encyclopedia, hoping that it would satisfy his curiosity. But it only made him insecure: he confided to his mother his fear that everything would have been discovered by the time he grew up, leaving him nothing to do. His mother reassured him, correctly, as it happened, that there would still be a thing or two for him to figure out.

\smallskip
\noindent\textbf{译文} 克里克总是被重要问题的复杂细节所吸引。小时候，他无穷无尽的问题迫使疲惫的父母给他买了一套儿童百科全书，希望满足他的好奇心。但这只让他更不安：他向母亲吐露担忧，害怕自己长大时一切都已被发现，留给他无事可做。他母亲安慰他说，仍会有一两件事等着他弄明白；事实证明她是对的。
\end{quote}

\begin{quote}
\noindent\textbf{原文} A great talker, Crick was invariably the center of attention in any gathering. His booming laugh was forever echoing down the hallways of the Cavendish. As the MRC Unit's resident theoretician, he used to come up with a novel insight at least once a month, and he would explain his latest idea at great length to anyone willing to listen. The morning we met he lit up when he learned that my objective in coming to Cambridge was to learn enough crystallography to have a go at the DNA structure. Soon I was asking Crick's opinion about using Pauling's model-building approach to go directly for the structure. Would we need many more years of diffraction experimentation before modeling would be practicable? To bring us up to speed on the status of DNA structural studies, Crick invited Maurice Wilkins, a friend since the end of the war, up from London for Sunday lunch. Then we could learn what progress Wilkins had made since his talk in Naples.

\smallskip
\noindent\textbf{译文} 克里克极其健谈，在任何聚会中总是注意力中心。他洪亮的笑声永远回荡在卡文迪许走廊中。作为MRC单元的常驻理论家，他过去每月至少会想出一个新见解，并向任何愿意听的人详尽解释自己的最新想法。我们相遇的那个早晨，当他得知我来剑桥的目标是学习足够的晶体学，以便尝试DNA结构时，他立刻来了精神。很快，我就询问克里克是否可以采用鲍林的模型构建方法，直接奔向结构。在建模可行之前，我们还需要许多年衍射实验吗？为了让我们了解DNA结构研究的现状，克里克邀请战后以来的朋友莫里斯·威尔金斯从伦敦来吃星期日午餐。这样我们就可以知道威尔金斯自那不勒斯报告以来取得了什么进展。
\end{quote}

\begin{figure}[htbp]
\centering
\safeincludeimage{3544e5837f922455f23bb7338a782d2a404c84d9e43c0ad947e0ab5094d455f8.jpg}
\caption{Francis Crick with the Cavendish X-ray tube.\\弗朗西斯·克里克与卡文迪许的X射线管。}
\end{figure}

\begin{quote}
\noindent\textbf{原文} Wilkins expressed his belief that DNA's structure was a helix, formed by several chains of linked nucleotides twisted around each other. All that remained to be settled was the number of chains. At the time, Wilkins favored three on the basis of his density measurements of DNA fibers. He was keen to start model-building, but he had run into a roadblock in the form of a new addition to the King's College Biophysics Unit, Rosalind Franklin.

\smallskip
\noindent\textbf{译文} 威尔金斯表示，他相信DNA结构是螺旋，由几条相连核苷酸链相互缠绕形成。剩下需要确定的只是链的数量。当时，基于对DNA纤维的密度测量，威尔金斯倾向于三链。他急于开始建模，但遇到了一道障碍：国王学院生物物理单元新来的人，罗莎琳德·富兰克林。
\end{quote}

\begin{quote}
\noindent\textbf{原文} A thirty-one-year-old Cambridge-trained physical chemist, Franklin was an obsessively professional scientist; for her twenty-ninth birthday all she requested was her own subscription to her field's technical journal, \emph{Acta Crystallographica}. Logical and precise, she was impatient with those who acted otherwise. And she was given to strong opinions, once describing her Ph.D. thesis adviser, Ronald Norrish, a future Nobel laureate, as "stupid, bigoted, deceitful, ill-mannered and tyrannical." Outside the laboratory, she was a determined and gutsy mountaineer, and, coming from the upper echelons of London society, she belonged to a more rarefied social world than most scientists. At the end of a hard day at the bench, she would occasionally change out of her lab coat into an elegant evening gown and disappear into the night.

\smallskip
\noindent\textbf{译文} 富兰克林三十一岁，是受剑桥训练的物理化学家，也是一位近乎执着的职业科学家；她二十九岁生日时唯一要求的礼物，是订阅自己领域的技术期刊《晶体学报》。她逻辑严密、精确，对行事不如此的人缺乏耐心。她观点强烈，曾把自己的博士论文导师、后来的诺贝尔奖得主罗纳德·诺里什形容为“愚蠢、偏执、虚伪、粗鲁而专横”。在实验室之外，她是坚定而勇敢的登山者；出身伦敦上流社会的她，属于比多数科学家更精致的社交世界。辛苦工作一天后，她偶尔会脱下实验服，换上优雅晚礼服，消失在夜色中。
\end{quote}

\begin{figure}[htbp]
\centering
\safeincludeimage{10e89428031ee62695360770c2642e8814912404728b22a692edeafa9def4bed.jpg}
\caption{Rosalind Franklin on one of the mountain hiking vacations she loved.\\罗莎琳德·富兰克林在她热爱的山地徒步假期中。}
\end{figure}

\begin{quote}
\noindent\textbf{原文} Just back from a four-year X-ray crystallographic investigation of graphite in Paris, Franklin had been assigned to the DNA project while Wilkins was away from King's. Unfortunately, the pair soon proved incompatible. Franklin, direct and data-focused, and Wilkins, retiring and speculative, were destined never to collaborate. Shortly before Wilkins accepted our lunch invitation, the two had had a big blowup in which Franklin had insisted that no model-building could commence before she collected much more extensive diffraction data. Now they effectively did not communicate, and Wilkins would have no chance to learn of her progress until Franklin presented her lab seminar scheduled for the beginning of November. If we wanted to listen, Crick and I were welcome to go as Wilkins's guests.

\smallskip
\noindent\textbf{译文} 富兰克林刚从巴黎回来，完成了为期四年的石墨X射线晶体学研究；在威尔金斯离开国王学院期间，她被分配到DNA项目。不幸的是，两人很快证明彼此不合。富兰克林直截了当、专注数据；威尔金斯内向、倾向推测，他们注定无法合作。就在威尔金斯接受我们的午餐邀请前不久，两人大吵了一架；富兰克林坚持认为，在她收集到更多广泛的衍射数据之前，不能开始任何模型构建。现在他们实际上不再交流，而威尔金斯也无从了解她的进展，直到富兰克林按计划在十一月初作实验室报告。如果我们想听，克里克和我可以作为威尔金斯的客人去。
\end{quote}

\begin{quote}
\noindent\textbf{原文} Crick was unable to make the seminar, so I attended alone and briefed him later on what I believed to be its key take-home messages on crystalline DNA. In particular, I described from memory Franklin's measurements of the crystallographic repeats and the water content. This prompted Crick to begin sketching helical grids on a sheet of paper, explaining that the new helical X-ray theory he had devised with Bill Cochran and Vladimir Vand would permit even me, a former bird-watcher, to predict correctly the diffraction patterns expected from the molecular models we would soon be building at the Cavendish.

\smallskip
\noindent\textbf{译文} 克里克无法参加报告，所以我独自前往，后来向他转述我认为其中关于晶态DNA的关键要点。尤其是，我凭记忆描述了富兰克林对晶体学重复周期和含水量的测量。这促使克里克在纸上开始画螺旋网格，并解释说，他与比尔·科克伦和弗拉基米尔·范德共同提出的新螺旋X射线理论，甚至能让我这个前观鸟者，正确预测我们很快将在卡文迪许搭建的分子模型应产生的衍射图案。
\end{quote}

\begin{quote}
\noindent\textbf{原文} As soon as we got back to Cambridge, I arranged for the Cavendish machine shop to construct the phosphorus atom models needed for short sections of the sugar-phosphate backbone found in DNA. Once these became available, we tested different ways the backbones might twist around each other in the center of the DNA molecule. Their regular repeating atomic structure should allow the atoms to come together in a consistent, repeated conformation. Following Wilkins's hunch, we focused on three-chain models. When one of these appeared to be almost plausible, Crick made a phone call to Wilkins to announce we had a model we thought might be DNA.

\smallskip
\noindent\textbf{译文} 我们一回到剑桥，我就安排卡文迪许机修车间制作DNA中糖-磷酸骨架短段所需的磷原子模型。模型做好后，我们测试骨架在DNA分子中心彼此缠绕的不同方式。它们规则重复的原子结构应当使原子以一致、重复的构象聚在一起。按照威尔金斯的直觉，我们集中研究三链模型。当其中一个模型看上去几乎可行时，克里克给威尔金斯打电话，宣布我们有了一个自认为可能是DNA的模型。
\end{quote}

\begin{quote}
\noindent\textbf{原文} The next day both Wilkins and Franklin came up to see what we had done. The threat of unanticipated competition briefly united them in common purpose. Franklin wasted no time in faulting our basic concept. My memory was that she had reported almost no water present in crystalline DNA. In fact, the opposite was true. Being a crystallographic novice, I had confused the terms "unit cell" and "asymmetric unit." Crystalline DNA was in fact water-rich. Consequently, Franklin pointed out, the backbone had to be on the outside and not, as we had it, in the center, if only to accommodate all the water molecules she had observed in her crystals.

\smallskip
\noindent\textbf{译文} 第二天，威尔金斯和富兰克林都来查看我们做了什么。意外竞争的威胁，短暂地把他们团结在共同目标上。富兰克林立刻指出我们基本概念的错误。我记得她报告过晶态DNA几乎不含水。事实上正相反。作为晶体学新手，我把“晶胞”和“不对称单元”混淆了。晶态DNA实际上富含水分。因此，富兰克林指出，骨架必须在外侧，而不能像我们的模型那样在中心，哪怕只是为了容纳她在晶体中观察到的所有水分子。
\end{quote}

\begin{quote}
\noindent\textbf{原文} That unfortunate November day cast a very long shadow. Franklin's opposition to model-building was reinforced. Doing experiments, not playing with Tinkertoy representations of atoms, was the way she intended to proceed. Even worse, Sir Lawrence Bragg passed down the word that Crick and I should desist from all further attempts at building a DNA model. It was further decreed that DNA research should be left to the King's lab, with Cambridge continuing to focus solely on proteins. There was no sense in two MRC-funded labs competing against each other. With no more bright ideas up our sleeves, Crick and I were reluctantly forced to back off, at least for the time being.

\smallskip
\noindent\textbf{译文} 那个不幸的十一月日子投下了很长的阴影。富兰克林反对建模的立场得到强化。她打算继续的方式，是做实验，而不是玩用儿童拼插玩具代表原子的游戏。更糟的是，劳伦斯·布拉格爵士传话下来，要求克里克和我停止一切进一步构建DNA模型的尝试。还进一步规定，DNA研究应交给国王学院实验室，剑桥继续只专注蛋白质。两个由MRC资助的实验室相互竞争没有意义。我们袖子里也没有更多妙计，克里克和我只好不情愿地退后一步，至少暂时如此。
\end{quote}

\begin{quote}
\noindent\textbf{原文} It was not a good moment to be condemned to the DNA sidelines. Linus Pauling had written Wilkins to request a copy of the crystalline DNA diffraction pattern. Though Wilkins had declined, saying he wanted more time to interpret it himself, Pauling was hardly obliged to depend upon data from King's. If he wished, he could easily start serious X-ray diffraction studies at Caltech.

\smallskip
\noindent\textbf{译文} 被判到DNA研究边线旁，这不是一个好时机。莱纳斯·鲍林已经写信给威尔金斯，请求一份晶态DNA衍射图样。威尔金斯拒绝了，说自己还需要更多时间解释它，但鲍林几乎不必依赖国王学院的数据。只要愿意，他完全可以在加州理工轻易开展严肃的X射线衍射研究。
\end{quote}

\begin{quote}
\noindent\textbf{原文} The following spring, I duly turned away from DNA and set about extending prewar studies on the pencil-shaped tobacco mosaic virus using the Cavendish's powerful new X-ray beam. This light experimental workload gave me plenty of time to wander through various Cambridge libraries. In the zoology building, I read Erwin Chargaff's paper describing his finding that the DNA bases adenine and thymine occurred in roughly equal amounts, as did the bases guanine and cytosine. Hearing of these one-to-one ratios Crick wondered whether, during DNA duplication, adenine residues might be attracted to thymine and vice versa, and whether a corresponding attraction might exist between guanine and cytosine. If so, base sequences on the "parental" chains, for example ATGC, would have to be complementary to those on "daughter" strands, yielding in this case TACG.

\smallskip
\noindent\textbf{译文} 第二年春天，我如命令所要求的那样转离DNA，开始使用卡文迪许强大的新X射线束，扩展战前对铅笔状烟草花叶病毒的研究。这个实验工作量不重，使我有大量时间在剑桥各个图书馆游荡。在动物学楼里，我读到了埃尔温·查伽夫的论文，文中描述了他的发现：DNA碱基腺嘌呤和胸腺嘧啶的含量大致相等，鸟嘌呤和胞嘧啶的含量也大致相等。听到这些一比一比例后，克里克想知道，在DNA复制过程中，腺嘌呤残基是否可能吸引胸腺嘧啶，反之亦然；鸟嘌呤和胞嘧啶之间是否也可能存在相应吸引。如果是这样，“亲本”链上的碱基序列，例如ATGC，就必须与“子代”链上的序列互补，在这个例子中产生TACG。
\end{quote}

\begin{quote}
\noindent\textbf{原文} These remained idle thoughts until Erwin Chargaff came through Cambridge in the summer of 1952 on his way to the International Biochemical Congress in Paris. Chargaff expressed annoyance that neither Crick nor I saw the need to know the chemical structures of the four bases. He was even more upset when we told him that we could simply look up the structures in textbooks as the need arose. I was left hoping that Chargaff's data would prove irrelevant. Crick, however, was energized to do several experiments looking for molecular "sandwiches" that might form when adenine and thymine, or alternatively guanine and cytosine, were mixed together in solution. But his experiments went nowhere.

\smallskip
\noindent\textbf{译文} 这些想法一直只是闲想，直到1952年夏天埃尔温·查伽夫途经剑桥，前往巴黎参加国际生物化学大会。查伽夫对克里克和我都不认为有必要知道四种碱基的化学结构感到恼火。当我们告诉他，若有需要可以直接去教科书上查结构时，他更加不悦。我只希望查伽夫的数据最终会被证明无关紧要。克里克却因此振奋起来，做了几项实验，寻找腺嘌呤和胸腺嘧啶，或者鸟嘌呤和胞嘧啶，在溶液中混合时可能形成的分子“三明治”。但他的实验没有进展。
\end{quote}

\begin{quote}
\noindent\textbf{原文} Like Chargaff, Linus Pauling also attended the International Biochemical Congress, where the big news was the latest result from the Phage Group. Alfred Hershey and Martha Chase at Cold Spring Harbor had just confirmed Avery's transforming principle: DNA was the hereditary material! Hershey and Chase proved that only the DNA of the phage virus enters bacterial cells; its protein coat remains on the outside. It was more obvious than ever that DNA must be understood at the molecular level if we were to uncover the essence of the gene. With Hershey and Chase's result the talk of the town, I was sure that Pauling would now bring his formidable intellect and chemical wisdom to bear on the problem of DNA.

\smallskip
\noindent\textbf{译文} 和查伽夫一样，莱纳斯·鲍林也参加了国际生物化学大会。大会上的大新闻是噬菌体小组的最新结果。冷泉港的阿尔弗雷德·赫尔希和玛莎·蔡斯刚刚证实了艾弗里的转化原则：DNA就是遗传物质！赫尔希和蔡斯证明，只有噬菌体病毒的DNA进入细菌细胞；其蛋白质外壳留在外面。若要揭示基因的本质，就必须在分子层面理解DNA，这比以往任何时候都更明显。赫尔希和蔡斯的结果成为全城热议，我确信鲍林如今会把他强大的智力和化学智慧用于DNA问题。
\end{quote}

\begin{quote}
\noindent\textbf{原文} Early in 1953, Pauling did indeed publish a paper outlining the structure of DNA. Reading it anxiously I saw that he was proposing a three-chain model with sugar-phosphate backbones forming a dense central core. Superficially it was similar to our botched model of fifteen months earlier. But instead of using positively charged atoms, for example \(Mg^{2+}\), to stabilize the negatively charged backbones, Pauling made the unorthodox suggestion that the phosphates were held together by hydrogen bonds. But it seemed to me, the biologist, that such hydrogen bonds required extremely acidic conditions never found in cells. With a mad dash to Alexander Todd's nearby organic chemistry lab my belief was confirmed: the impossible had happened. The world's best-known, if not best, chemist had gotten his chemistry wrong. In effect, Pauling had knocked the A off of DNA. Our quarry was deoxyribonucleic acid, but the structure he was proposing was not even acidic.

\smallskip
\noindent\textbf{译文} 1953年初，鲍林确实发表了一篇概述DNA结构的论文。我焦急地读着，看到他提出的是三链模型，糖-磷酸骨架形成致密中央核心。从表面看，这与我们十五个月前搞砸的模型相似。但鲍林没有使用带正电的原子，例如 \(Mg^{2+}\)，来稳定带负电的骨架，而是提出一个非正统建议：磷酸基团由氢键维系在一起。可在我这个生物学家看来，这样的氢键需要细胞内从不存在的极端酸性条件。我狂奔到附近亚历山大·托德的有机化学实验室，我的判断得到证实：不可能发生的事发生了。世界上最有名，即使不是最好的，化学家把化学弄错了。实际上，鲍林把DNA中的“A”去掉了。我们的猎物是脱氧核糖核酸，但他提出的结构甚至不是酸性的。
\end{quote}

\begin{quote}
\noindent\textbf{原文} Hurriedly I took the manuscript to London to inform Wilkins and Franklin they were still in the game. Convinced that DNA was not a helix, Franklin had no wish even to read the article and deal with the distraction of Pauling's helical ideas, even when I offered Crick's arguments for helices. Wilkins, however, was very interested indeed in the news I brought; he was now more certain than ever that DNA was helical. To prove the point, he showed me a photograph obtained more than six months earlier by Franklin's graduate student Raymond Gosling, who had X-rayed the so-called B form of DNA. Until that moment, I did not know a B form even existed. Franklin had put this picture aside, preferring to concentrate on the A form, which she thought would more likely yield useful data. The X-ray pattern of this B form was a distinct cross. Since Crick and others had already deduced that such a pattern of reflections would be created by a helix, this evidence made it clear that DNA had to be a helix! In fact, despite Franklin's reservations, this was no surprise. Geometry itself suggested that a helix was the most logical arrangement for a long string of repeating units such as the nucleotides of DNA. But we still did not know what that helix looked like, nor how many chains it contained.

\smallskip
\noindent\textbf{译文} 我匆忙把手稿带到伦敦，告诉威尔金斯和富兰克林他们还在游戏中。富兰克林确信DNA不是螺旋，甚至不愿读这篇文章，不愿应付鲍林螺旋思想的干扰，即使我提出克里克支持螺旋的论据也无济于事。然而，威尔金斯对我带来的消息确实非常感兴趣；他现在比以往更确信DNA是螺旋。为了证明这一点，他给我看了一张六个多月前由富兰克林研究生雷蒙德·戈斯林获得的照片，戈斯林曾对所谓DNA B型进行X射线照射。在那一刻之前，我甚至不知道B型存在。富兰克林把这张照片搁在一旁，更愿意集中研究A型，因为她认为A型更可能产生有用数据。这个B型的X射线图案呈现清晰十字。克里克和其他人已经推导出，螺旋会产生这种反射图案，因此这一证据清楚表明DNA必定是螺旋！事实上，尽管富兰克林有所保留，这并不令人意外。几何学本身提示，对DNA核苷酸这样长串重复单元来说，螺旋是最合乎逻辑的排列。但我们仍不知道这个螺旋是什么样子，也不知道它包含几条链。
\end{quote}

\begin{quote}
\noindent\textbf{原文} The time had come to resume building helical models of DNA. Pauling was bound to realize soon enough that his brainchild was wrong. I urged Wilkins to waste no time. But he wanted to wait until Franklin had completed her scheduled departure for another lab later that spring. She had decided to move on to avoid the unpleasantness at King's. Before leaving, she had been ordered to stop further work with DNA and had already passed on many of her diffraction images to Wilkins.

\smallskip
\noindent\textbf{译文} 重新构建DNA螺旋模型的时刻到了。鲍林迟早会意识到自己的心血结晶是错的。我敦促威尔金斯不要浪费时间。但他想等到富兰克林按计划在那年晚春离开，前往另一间实验室之后。她已经决定离开，以避开国王学院的不愉快。离开前，她被命令停止进一步研究DNA，并已把许多衍射图像交给威尔金斯。
\end{quote}

\begin{figure}[htbp]
\centering
\safeincludeimagepair{19c9e38243b212a63b0c2c3bdaff52a3070510c4e38cefa3383690dda5682970.jpg}{e490980fa1fb6465bd4c9013ed309954bdf81cfd15b9d02dcee06b0de3b45061.jpg}
\caption{X-ray photos of the A and B forms of DNA from, respectively, Maurice Wilkins and Rosalind Franklin. The differences in molecular structure are caused by differences in the amount of water associated with each DNA molecule.\\DNA的A型和B型X射线照片，分别来自莫里斯·威尔金斯和罗莎琳德·富兰克林。分子结构差异由每个DNA分子结合水量的差异造成。}
\end{figure}

\begin{quote}
\noindent\textbf{原文} When I returned to Cambridge and broke the news of the DNA B form, Bragg no longer saw any reason for Crick and me to avoid DNA. He very much wanted the DNA structure to be found on his side of the Atlantic. So we went back to model-building, looking for a way the known basic components of DNA---the backbone of the molecule and the four different bases, adenine, thymine, guanine, and cytosine---could fit together to make a helix. I commissioned the shop at the Cavendish to make us a set of tin bases, but they could not produce them fast enough for me: I ended up cutting out rough approximations from stiff cardboard.

\smallskip
\noindent\textbf{译文} 我回到剑桥并透露DNA B型的消息后，布拉格再也看不到克里克和我避开DNA的理由。他非常希望DNA结构能在大西洋这一侧被发现。于是我们重新开始建模，寻找一种方式，让DNA已知的基本组成部分，即分子骨架和四种不同碱基腺嘌呤、胸腺嘧啶、鸟嘌呤、胞嘧啶，能够组合成螺旋。我委托卡文迪许车间为我们制作一套锡制碱基，但他们做得不够快；最后我用硬纸板剪出粗略近似物。
\end{quote}

\begin{quote}
\noindent\textbf{原文} By this time I realized the DNA density-measurement evidence actually slightly favored a two-chain, rather than three-chain, model. So I decided to search out plausible double helices. As a biologist, I preferred the idea of a genetic molecule made of two, rather than three, components. After all, chromosomes, like cells, increase in number by duplicating, not triplicating.

\smallskip
\noindent\textbf{译文} 到这时，我意识到DNA密度测量证据实际上稍微支持两链模型，而不是三链模型。因此，我决定寻找可行的双螺旋。作为生物学家，我更喜欢遗传分子由两个组成部分而非三个组成部分构成的想法。毕竟，染色体像细胞一样，是通过复制而不是三倍化来增加数目的。
\end{quote}

\begin{figure}[htbp]
\centering
\safeincludeimage{c879c31c90a3b1b3c2b93ae578d9b3f6990843106b896ca71172ac4499e54282.jpg}
\caption{The chemical backbone of DNA.\\DNA的化学骨架。}
\end{figure}

\begin{quote}
\noindent\textbf{原文} I knew that our previous model with the backbone on the inside and the bases hanging out was wrong. Chemical evidence from the University of Nottingham, which I had too long ignored, indicated that the bases must be hydrogen-bonded to each other. They could only form bonds like this in the regular manner implied by the X-ray diffraction data if they were in the center of the molecule. But how could they come together in pairs? For two weeks I got nowhere, misled by an error in my nucleic acid chemistry textbook. Happily, on February 27, Jerry Donohue, a theoretical chemist visiting the Cavendish from Caltech, pointed out that the textbook was wrong. So I changed the locations of the hydrogen atoms on my cardboard cutouts of the molecules.

\smallskip
\noindent\textbf{译文} 我知道我们先前那个骨架在内、碱基朝外的模型是错的。我长期忽视的诺丁汉大学化学证据表明，碱基之间必须彼此形成氢键。只有当碱基位于分子中心时，它们才能以X射线衍射数据所暗示的规则方式形成这种键。可是它们如何成对聚合？两个星期里我毫无进展，被核酸化学教科书中的一个错误误导了。幸运的是，2月27日，从加州理工来卡文迪许访问的理论化学家杰里·多诺休指出教科书错了。于是我改变了分子硬纸板剪片上氢原子的位置。
\end{quote}

\begin{quote}
\noindent\textbf{原文} The next morning, February 28, 1953, the key features of the DNA model all fell into place. The two chains were held together by strong hydrogen bonds between adenine--thymine and guanine--cytosine base pairs. The inferences Crick had drawn the year before based on Chargaff's research had indeed been correct. Adenine does bond to thymine and guanine does bond to cytosine, but not through flat surfaces to form molecular sandwiches. When Crick arrived, he took it all in rapidly, and gave my base-pairing scheme his blessing. He realized right away that it would result in the two strands of the double helix running in opposite directions.

\smallskip
\noindent\textbf{译文} 第二天早晨，1953年2月28日，DNA模型的关键特征全部就位。两条链由腺嘌呤--胸腺嘧啶和鸟嘌呤--胞嘧啶碱基对之间的强氢键维系。克里克前一年基于查伽夫研究作出的推论确实正确。腺嘌呤会与胸腺嘧啶结合，鸟嘌呤会与胞嘧啶结合，但不是通过平面表面形成分子三明治。克里克到来后，很快理解了一切，并认可了我的碱基配对方案。他立刻意识到，这将导致双螺旋的两条链以相反方向运行。
\end{quote}

\begin{quote}
\noindent\textbf{原文} It was quite a moment. We felt sure that this was it. Anything that simple, that elegant just had to be right. What got us most excited was the complementarity of the base sequences along the two chains. If you knew the sequence---the order of bases---along one chain, you automatically knew the sequence along the other. It was immediately apparent that this must be how the genetic messages of genes are copied so exactly when chromosomes duplicate prior to cell division. The molecule would "unzip" to form two separate strands. Each separate strand then could serve as the template for the synthesis of a new strand, one double helix becoming two.

\smallskip
\noindent\textbf{译文} 那真是一个了不起的时刻。我们确信就是它了。如此简单、如此优雅的东西，必然是正确的。最令我们兴奋的是两条链上碱基序列的互补性。如果你知道一条链上的序列，也就是碱基顺序，就会自动知道另一条链上的序列。立刻显而易见的是，这一定就是染色体在细胞分裂前复制时，基因的遗传信息能够如此精确复制的方式。这个分子会像拉链一样“拉开”，形成两条单独链。每条单链随后都可以作为模板合成一条新链，一个双螺旋变成两个。
\end{quote}

\begin{quote}
\noindent\textbf{原文} In \emph{What Is Life?} Schrödinger had suggested that the language of life might be like Morse code, a series of dots and dashes. He was not far off. The language of DNA is a linear series of As, Ts, Gs, and Cs. And just as transcribing a page out of a book can result in the odd typo, the rare mistake creeps in when all these As, Ts, Gs, and Cs are being copied along a chromosome. These errors are the mutations geneticists had talked about for almost fifty years. Change an "i" to an "a" and "Jim" becomes "Jam" in English; change a T to a C and "ATG" becomes "ACG" in DNA.

\smallskip
\noindent\textbf{译文} 在《生命是什么？》中，薛定谔曾提出，生命的语言也许像莫尔斯电码，是一系列点和划。他离真相并不远。DNA的语言是由A、T、G、C组成的线性序列。正如抄写书页会偶尔出现错别字，当所有这些A、T、G、C沿染色体被复制时，罕见错误也会潜入。这些错误就是遗传学家谈论了近五十年的突变。在英语中，把“Jim”里的“i”改成“a”，“Jim”就变成“Jam”；在DNA中，把T改成C，“ATG”就变成“ACG”。
\end{quote}

\begin{figure}[htbp]
\centering
\safeincludeimagepair{9ed2d5675b2303884242823021aa3e3ba5819dec62e874dbeaa1906ab110f706.jpg}{c99acaa9355ddb2b918c14b2bb9ebb29ed0780f6348debfe1bfe251cb90b9209.jpg}
\caption{The insight that made it all come together: complementary pairing of the bases.\\让一切贯通的洞见：碱基的互补配对。}
\end{figure}

\begin{figure}[htbp]
\centering
\safeincludeimagepair{89cf17fe3c783f0129e6decdcfd90773809f800d82fac3712a7aa030f95c657c.jpg}{556ed40ecd12f173f0a6361a38347c46ec70df5becf9e667c1cb7ac95a6ed493.jpg}
\caption{Bases and backbone in place: the double helix. (A) is a schematic showing the system of base-pairing that binds the two strands together. (B) is a "space-filling" model showing, to scale, the atomic detail of the molecule.\\碱基和骨架各就各位：双螺旋。（A）为示意图，显示把两条链结合在一起的碱基配对系统。（B）为“空间填充”模型，按比例展示分子的原子细节。}
\end{figure}

\begin{quote}
\noindent\textbf{原文} The double helix made sense chemically and it made sense biologically. Now there was no need to be concerned about Schrödinger's suggestion that new laws of physics might be necessary for an understanding of how the hereditary code-script is duplicated: genes in fact were no different from the rest of chemistry. Later that day, during lunch at the Eagle, the pub virtually adjacent to the Cavendish Lab, Crick, ever the talker, could not help but tell everyone we had just found the "secret of life." I myself, though no less electrified by the thought, would have waited until we had a pretty three-dimensional model to show off.

\smallskip
\noindent\textbf{译文} 双螺旋在化学上说得通，在生物学上也说得通。现在无需再担心薛定谔关于理解遗传密码脚本复制也许需要新物理定律的说法：基因事实上并不不同于其他化学。当天晚些时候，在几乎紧挨卡文迪许实验室的“鹰”酒馆吃午饭时，一贯健谈的克里克忍不住告诉所有人，我们刚刚发现了“生命的秘密”。至于我自己，虽然同样为这个想法激动不已，却会等到我们有了漂亮的三维模型可展示之后再说。
\end{quote}

\begin{quote}
\noindent\textbf{原文} Among the first to see our demonstration model was the chemist Alexander Todd. That the nature of the gene was so simple both surprised and pleased him.

\smallskip
\noindent\textbf{译文} 最早看到我们演示模型的人之一，是化学家亚历山大·托德。基因的本质如此简单，既让他惊讶，也让他高兴。
\end{quote}

\begin{quote}
\noindent\textbf{原文} Later, however, he must have asked himself why his own lab, having established the general chemical structure of DNA chains, had not moved on to asking how the chains folded up in three dimensions. Instead the essence of the molecule was left to be discovered by a two-man team, a biologist and a physicist, neither of whom possessed a detailed command even of undergraduate chemistry. But paradoxically, this was, at least in part, the key to our success: Crick and I arrived at the double helix first precisely because most chemists at that time thought DNA too big a molecule to understand by chemical analysis.

\smallskip
\noindent\textbf{译文} 然而后来，他一定问过自己：既然自己的实验室已经建立了DNA链的一般化学结构，为什么没有进一步追问这些链在三维空间中如何折叠？相反，这个分子的本质留给了两人小组来发现，一个生物学家和一个物理学家，而两人甚至都没有详细掌握本科化学。但矛盾的是，这至少部分正是我们成功的关键：克里克和我之所以最先抵达双螺旋，恰恰因为当时多数化学家认为DNA分子太大，无法用化学分析来理解。
\end{quote}

\begin{quote}
\noindent\textbf{原文} At the same time, the only two chemists with the vision to seek DNA's 3-D structure made major tactical mistakes: Rosalind Franklin's was her resistance to model-building; Linus Pauling's was a matter of simply neglecting to read the existing literature on DNA, particularly the data on its base composition published by Chargaff. Ironically, Pauling and Chargaff sailed across the Atlantic on the same ship following the Paris Biochemical Congress in 1952, but failed to hit it off. Pauling was long accustomed to being right. And he believed there was no chemical problem he could not work out from first principles by himself. Usually this confidence was not misplaced. During the Cold War, as a prominent critic of the American nuclear weapons development program, he was questioned by the FBI after giving a talk. How did he know how much plutonium there is in an atomic bomb? Pauling's response was "Nobody told me. I figured it out."

\smallskip
\noindent\textbf{译文} 与此同时，唯有两个有远见去寻找DNA三维结构的化学家，都犯下了重大战术错误：罗莎琳德·富兰克林的错误在于抵制模型构建；莱纳斯·鲍林的错误则只是没有阅读已有DNA文献，尤其是查伽夫发表的碱基组成数据。具有讽刺意味的是，1952年巴黎生物化学大会后，鲍林和查伽夫乘同一艘船横渡大西洋，却没能投缘。鲍林早已习惯于自己是对的。他相信没有任何化学问题不能由自己从第一原理推出来。通常，这种自信并没有错置。冷战期间，他作为美国核武器发展计划的著名批评者，在一次演讲后受到FBI询问。他怎么知道一枚原子弹里有多少钚？鲍林回答说：“没人告诉我。我算出来的。”
\end{quote}

\begin{quote}
\noindent\textbf{原文} Over the next several months Crick and, to a lesser extent, I relished showing off our model to an endless stream of curious scientists. However, the Cambridge biochemists did not invite us to give a formal talk in the biochemistry building. They started to refer to it as the "WC," punning our initials with those used in Britain for the toilet or water closet. That we had found the double helix without doing experiments irked them.

\smallskip
\noindent\textbf{译文} 在接下来的几个月里，克里克和我，后者程度稍弱，都非常享受向源源不断的好奇科学家展示我们的模型。然而，剑桥的生物化学家并没有邀请我们到生化楼作正式报告。他们开始把它称作“WC”，用我们的姓名首字母与英国表示厕所或盥洗室的缩写双关。我们没有做实验却找到了双螺旋，这让他们恼火。
\end{quote}

\begin{quote}
\noindent\textbf{原文} The manuscript that we submitted to \emph{Nature} in early April was published just over three weeks later, on April 25, 1953. Accompanying it were two longer papers by Franklin and Wilkins, both supporting the general correctness of our model. In June, I gave the first presentation of our model at the Cold Spring Harbor symposium on viruses. Max Delbrück saw to it that I was offered, at the last minute, an invitation to speak. To this intellectually high-powered meeting I brought a three-dimensional model built in the Cavendish, the adenine--thymine base pairs in red and the guanine--cytosine base pairs in green.

\smallskip
\noindent\textbf{译文} 我们四月初提交给《自然》的手稿，在三周多后，即1953年4月25日发表。与之同期刊出的还有富兰克林和威尔金斯的两篇更长论文，二者都支持我们模型的大体正确性。六月，我在冷泉港病毒专题讨论会上首次介绍了我们的模型。马克斯·德尔布吕克确保我在最后时刻获得了发言邀请。面对这场智力强度极高的会议，我带来了一个在卡文迪许制作的三维模型，其中腺嘌呤--胸腺嘧啶碱基对为红色，鸟嘌呤--胞嘧啶碱基对为绿色。
\end{quote}

\section*{Molecular Structure of Nucleic Acids}
\begin{center}
\large A Structure for Deoxyribose Nucleic Acid\\
\large 核酸的分子结构：脱氧核糖核酸的一种结构
\end{center}

\begin{quote}
\noindent\textbf{原文} We wish to suggest a structure for the salt of deoxyribose nucleic acid (D.N.A.). This structure has novel features which are of considerable biological interest.

\smallskip
\noindent\textbf{译文} 我们希望提出脱氧核糖核酸（D.N.A.）盐的一种结构。这个结构具有若干新颖特征，具有相当大的生物学意义。
\end{quote}

\begin{quote}
\noindent\textbf{原文} A structure for nucleic acid has already been proposed by Pauling and Corey\textsuperscript{1}. They kindly made their manuscript available to us in advance of publication. Their model consists of three intertwined chains, with the phosphates near the fibre axis, and the bases on the outside. In our opinion, this structure is unsatisfactory for two reasons: (1) We believe that the material which gives the X-ray diagrams is the salt, not the free acid. Without the acidic hydrogen atoms it is not clear what forces would hold the structure together, especially as the negatively charged phosphates near the axis will repel each other. (2) Some of the van der Waals distances appear to be too small.

\smallskip
\noindent\textbf{译文} 鲍林和科里已经提出过一种核酸结构\textsuperscript{1}。他们友好地在发表前把手稿提供给我们。他们的模型由三条相互缠绕的链组成，磷酸基靠近纤维轴，碱基在外侧。我们认为这个结构有两个方面不能令人满意：（1）我们相信产生X射线图样的材料是盐，而不是游离酸。没有酸性氢原子，就不清楚是什么力量把结构维系在一起，尤其是靠近轴心的带负电磷酸基会彼此排斥。（2）某些范德华距离似乎太小。
\end{quote}

\begin{quote}
\noindent\textbf{原文} Another three-chain structure has also been suggested by Fraser, in the press. In his model the phosphates are on the outside and the bases on the inside, linked together by hydrogen bonds. This structure as described is rather ill-defined, and for this reason we shall not comment on it.

\smallskip
\noindent\textbf{译文} 弗雷泽也提出了另一种三链结构，论文即将发表。在他的模型中，磷酸基在外侧，碱基在内侧，并由氢键连接在一起。按其描述，这一结构相当不明确，因此我们不作评论。
\end{quote}

\begin{figure}[htbp]
\centering
\safeincludeimage{a43c7f1bb5bbda79e9e814375d12af675f324f9ef0f2b73d1f00ab1a62b64b57.jpg}
\caption{This figure is purely diagrammatic. The two ribbons symbolize the two phosphate--sugar chains, and the horizontal rods the pairs of bases holding the chains together. The vertical line marks the fibre axis.\\此图纯属示意。两条带状物表示两条磷酸--糖链，水平短杆表示把链维系在一起的碱基对。垂直线标示纤维轴。}
\end{figure}

\begin{quote}
\noindent\textbf{原文} We wish to put forward a radically different structure for the salt of deoxyribose nucleic acid. This structure has two helical chains each coiled round the same axis (see diagram). We have made the usual chemical assumptions, namely, that each chain consists of phosphate diester groups joining \(\beta\)-D-deoxyribofuranose residues with 3',5' linkages. The two chains, but not their bases, are related by a dyad perpendicular to the fibre axis. Both chains follow right-handed helices, but owing to the dyad the sequences of the atoms in the two chains run in opposite directions. Each chain loosely resembles Furberg's model No. 1; that is, the bases are on the inside of the helix and the phosphates on the outside. The configuration of the sugar and the atoms near it is close to Furberg's "standard configuration," the sugar being roughly perpendicular to the attached base. There is a residue on each chain every 3.4 \AA{} in the \(z\)-direction. We have assumed an angle of \(36^\circ\) between adjacent residues in the same chain, so that the structure repeats after 10 residues on each chain, that is, after 34 \AA{}. The distance of a phosphorus atom from the fibre axis is 10 \AA{}. As the phosphates are on the outside, cations have easy access to them.

\smallskip
\noindent\textbf{译文} 我们希望为脱氧核糖核酸盐提出一种根本不同的结构。这个结构有两条螺旋链，各自围绕同一轴缠绕（见图）。我们采用通常的化学假设，即每条链由磷酸二酯基团组成，这些基团通过3',5'键连接 \(\beta\)-D-脱氧核糖呋喃糖残基。两条链（但不包括它们的碱基）由一条垂直于纤维轴的二重轴相关。两条链都遵循右手螺旋，但由于二重轴，两条链中原子的序列方向相反。每条链大致类似富伯格的1号模型；也就是说，碱基在螺旋内部，磷酸基在外部。糖及其附近原子的构型接近富伯格的“标准构型”，糖大致垂直于所连接的碱基。沿 \(z\) 方向，每条链每隔3.4埃有一个残基。我们假定同一链上相邻残基之间的角度为 \(36^\circ\)，因此每条链上10个残基后结构重复，即34埃后重复。磷原子距纤维轴10埃。由于磷酸基在外侧，阳离子很容易接近它们。
\end{quote}

\begin{quote}
\noindent\textbf{原文} The structure is an open one, and its water content is rather high. At lower water contents we would expect the bases to tilt so that the structure could become more compact.

\smallskip
\noindent\textbf{译文} 该结构是开放的，含水量相当高。在含水量较低时，我们预期碱基会倾斜，使结构变得更加紧凑。
\end{quote}

\begin{quote}
\noindent\textbf{原文} The novel feature of the structure is the manner in which the two chains are held together by the purine and pyrimidine bases. The planes of the bases are perpendicular to the fibre axis. They are joined together in pairs, a single base from one chain being hydrogen-bonded to a single base from the other chain, so that the two lie side by side with identical \(z\)-coordinates. One of the pair must be a purine and the other a pyrimidine for bonding to occur. The hydrogen bonds are made as follows: purine position 1 to pyrimidine position 1; purine position 6 to pyrimidine position 6.

\smallskip
\noindent\textbf{译文} 这一结构的新颖特征，在于两条链由嘌呤和嘧啶碱基维系在一起的方式。碱基平面垂直于纤维轴。它们成对连接，一条链上的一个碱基与另一条链上的一个碱基形成氢键，使二者并排处在相同的 \(z\) 坐标上。要发生成键，一对中一个必须是嘌呤，另一个必须是嘧啶。氢键形成方式如下：嘌呤1位连到嘧啶1位；嘌呤6位连到嘧啶6位。
\end{quote}

\begin{quote}
\noindent\textbf{原文} If it is assumed that the bases only occur in the structure in the most plausible tautomeric forms, that is, with the keto rather than the enol configurations, it is found that only specific pairs of bases can bond together. These pairs are: adenine, a purine, with thymine, a pyrimidine; and guanine, a purine, with cytosine, a pyrimidine.

\smallskip
\noindent\textbf{译文} 如果假定碱基在结构中只以最合理的互变异构形式出现，也就是以酮式而不是烯醇式构型出现，就会发现只有特定碱基对能够彼此成键。这些配对是：腺嘌呤（嘌呤）与胸腺嘧啶（嘧啶）；鸟嘌呤（嘌呤）与胞嘧啶（嘧啶）。
\end{quote}

\begin{quote}
\noindent\textbf{原文} In other words, if an adenine forms one member of a pair, on either chain, then on these assumptions the other member must be thymine; similarly for guanine and cytosine. The sequence of bases on a single chain does not appear to be restricted in any way. However, if only specific pairs of bases can be formed, it follows that if the sequence of bases on one chain is given, then the sequence on the other chain is automatically determined.

\smallskip
\noindent\textbf{译文} 换言之，如果腺嘌呤在任一条链上构成一对中的一个成员，那么在这些假设下，另一成员必定是胸腺嘧啶；鸟嘌呤与胞嘧啶同理。单条链上的碱基序列似乎不受任何限制。然而，如果只有特定碱基对能够形成，那么只要给定一条链上的碱基序列，另一条链上的序列就会自动确定。
\end{quote}

\begin{quote}
\noindent\textbf{原文} It has been found experimentally\textsuperscript{2,4} that the ratio of the amounts of adenine to thymine, and the ratio of guanine to cytosine, are always very close to unity for deoxyribose nucleic acid.

\smallskip
\noindent\textbf{译文} 实验已经发现\textsuperscript{2,4}，在脱氧核糖核酸中，腺嘌呤与胸腺嘧啶的含量比，以及鸟嘌呤与胞嘧啶的含量比，总是非常接近1。
\end{quote}

\begin{quote}
\noindent\textbf{原文} It is probably impossible to build this structure with a ribose sugar in place of the deoxyribose, as the extra oxygen atom would make too close a van der Waals contact.

\smallskip
\noindent\textbf{译文} 如果用核糖替代脱氧核糖，可能无法构建这一结构，因为额外的氧原子会造成过近的范德华接触。
\end{quote}

\begin{quote}
\noindent\textbf{原文} The previously published X-ray data\textsuperscript{5,6} on deoxyribose nucleic acid are insufficient for a rigorous test of our structure. So far as we can tell, it is roughly compatible with the experimental data, but it must be regarded as unproved until it has been checked against more exact results. Some of these are given in the following communications. We were not aware of the details of the results presented there when we devised our structure, which rests mainly though not entirely on published experimental data and stereochemical arguments.

\smallskip
\noindent\textbf{译文} 先前发表的关于脱氧核糖核酸的X射线数据\textsuperscript{5,6}，不足以对我们的结构作严格检验。就我们所能判断，它大体上与实验数据相容，但在与更精确结果核对之前，必须视为尚未证明。随后的论文给出了一些这样的结果。我们设计该结构时，并不知道那些结果的细节；我们的结构主要但并非完全建立在已发表实验数据和立体化学论证之上。
\end{quote}

\begin{quote}
\noindent\textbf{原文} It has not escaped our notice that the specific pairing we have postulated immediately suggests a possible copying mechanism for the genetic material.

\smallskip
\noindent\textbf{译文} 我们注意到，我们所假定的特异性配对会立即提示遗传物质一种可能的复制机制。
\end{quote}

\begin{quote}
\noindent\textbf{原文} Full details of the structure, including the conditions assumed in building it, together with a set of co-ordinates for the atoms, will be published elsewhere.

\smallskip
\noindent\textbf{译文} 该结构的完整细节，包括构建时假定的条件以及一组原子坐标，将另文发表。
\end{quote}

\begin{quote}
\noindent\textbf{原文} We are much indebted to Dr. Jerry Donohue for constant advice and criticism, especially on interatomic distances. We have also been stimulated by a knowledge of the general nature of the unpublished experimental results and ideas of Dr. M. H. F. Wilkins, Dr. R. E. Franklin and their co-workers at King's College, London. One of us (J. D. W.) has been aided by a fellowship from the National Foundation for Infantile Paralysis.

\smallskip
\noindent\textbf{译文} 我们非常感谢杰里·多诺休博士持续不断的建议和批评，尤其是在原子间距离方面。我们还受到伦敦国王学院M. H. F. 威尔金斯博士、R. E. 富兰克林博士及其同事未发表实验结果和想法的一般性质的启发。我们中的一人（J. D. W.）获得了国家小儿麻痹症基金会奖学金资助。
\end{quote}

\begin{flushright}
J. D. Watson\\
F. H. C. Crick
\end{flushright}

\begin{quote}
\noindent Medical Research Council Unit for the Study of the Molecular Structure of Biological Systems, Cavendish Laboratory, Cambridge. April 2.

\smallskip
\noindent 医学研究委员会生物系统分子结构研究单元，剑桥卡文迪许实验室。4月2日。
\end{quote}

\begin{enumerate}
\item Pauling, L., and Corey, R. B., \emph{Nature}, 171, 346 (1953); \emph{Proc. U.S. Nat. Acad. Sci.}, 39, 84 (1953).
\item Furberg, S., \emph{Acta Chem. Scand.}, 6, 634 (1952).
\item Chargaff, E., for references see Zamenhof, S., Brawerman, G., and Chargaff, E., \emph{Biochim. et Biophys. Acta}, 9, 402 (1952).
\item Wyatt, G. R., \emph{J. Gen. Physiol.}, 36, 201 (1952).
\item Astbury, W. T., \emph{Symp. Soc. Exp. Biol.} 1, \emph{Nucleic Acid}, 66 (Camb. Univ. Press, 1947).
\item Wilkins, M. H. F., and Randall, J. T., \emph{Biochim. et Biophys. Acta}, 10, 192 (1953).
\end{enumerate}

\begin{quote}
\noindent\textbf{原文} Short and sweet: our \emph{Nature} paper announcing the discovery. The same issue also carried longer articles by Rosalind Franklin and Maurice Wilkins.

\smallskip
\noindent\textbf{译文} 简短而漂亮：这就是我们宣布发现的《自然》论文。同一期还刊登了罗莎琳德·富兰克林和莫里斯·威尔金斯的较长论文。
\end{quote}

\begin{figure}[htbp]
\centering
\safeincludeimage{5f5c9a86c1556b15bc87624993ecaf2a5350fa6baa71c3325042a808c1935d14.jpg}
\caption{Unveiling the double helix: my lecture at Cold Spring Harbor Laboratory, June 1953.\\揭开双螺旋：1953年6月，我在冷泉港实验室作报告。}
\end{figure}

\begin{quote}
\noindent\textbf{原文} In the audience was Seymour Benzer, yet another ex-physicist who had heeded the clarion call of Schrödinger's book. He immediately understood what our breakthrough meant for his studies of mutations in viruses. He realized that he could now do for a short stretch of bacteriophage DNA what Morgan's boys had done forty years earlier for fruit fly chromosomes: he would map mutations---determine their order---along a gene, just as the fruit fly pioneers had mapped genes along a chromosome. Like Morgan, Benzer would have to depend on recombination to generate new genetic combinations, but, whereas Morgan had the advantage of a ready mechanism of recombination---the production of sex cells in a fruit fly---Benzer had to induce recombination by simultaneously infecting a single bacterial host cell with two different strains of bacteriophage, which differed by one or more mutations in the region of interest. Within the bacterial cell, recombination---the exchange of segments of molecules---would occasionally occur between the different viral DNA molecules, producing new permutations of mutations, so-called "recombinants." Within a single astonishingly productive year in his Purdue University lab, Benzer produced a map of a single bacteriophage gene, rII, showing how a series of mutations---all errors in the genetic script---were laid out linearly along the virus DNA. The language was simple and linear, just like a line of text on the written page.

\smallskip
\noindent\textbf{译文} 听众中有西摩·本泽，他也是一位响应薛定谔之书号角召唤的前物理学家。他立刻理解我们的突破对自己研究病毒突变意味着什么。他意识到，现在他可以对一小段噬菌体DNA做四十年前“摩尔根的男孩们”对果蝇染色体所做的事：他将沿一个基因绘制突变图谱，确定它们的顺序，正如果蝇先驱者沿染色体绘制基因图谱一样。和摩尔根一样，本泽必须依赖重组来产生新的遗传组合；但摩尔根拥有现成的重组机制，即果蝇性细胞的产生，而本泽必须通过用两种不同噬菌体菌株同时感染同一个细菌宿主细胞来诱导重组，这两种菌株在所关注区域相差一个或多个突变。在细菌细胞内，不同病毒DNA分子之间偶尔会发生重组，即分子片段交换，产生突变的新排列，也就是所谓“重组体”。在普渡大学实验室里惊人高产的一年中，本泽绘制出单个噬菌体基因rII的图谱，显示一系列突变，即遗传脚本中的错误，如何沿病毒DNA线性排列。这种语言简单而线性，就像书页上一行文字。
\end{quote}

\begin{quote}
\noindent\textbf{原文} The response of the Hungarian physicist Leo Szilard to my Cold Spring Harbor talk on the double helix was less academic. His question was, "Can you patent it?" At one time Szilard's main source of income had been a patent that he held with Einstein, and he had later tried unsuccessfully to patent with Enrico Fermi the nuclear reactor they built at the University of Chicago in 1942. But then as now patents were given only for useful inventions and at the time no one could conceive of a practical use for DNA. Perhaps then, Szilard suggested, we should copyright it.

\smallskip
\noindent\textbf{译文} 匈牙利物理学家利奥·西拉德对我在冷泉港关于双螺旋报告的反应就不那么学术了。他的问题是：“你们能给它申请专利吗？”西拉德曾有一段时间主要收入来自他与爱因斯坦共同持有的一项专利，后来他还曾试图与恩里科·费米一起为他们1942年在芝加哥大学建造的核反应堆申请专利，但没有成功。然而无论当时还是现在，专利只授予有用发明，而当时无人能想象DNA有什么实际用途。那么，西拉德建议，也许我们应该给它登记版权。
\end{quote}

\begin{figure}[htbp]
\centering
\safeincludeimage{a026b9883dbbda8eac990954a994fc9eaad9b41e2e1663ff203da9d998f6ce50.jpg}
\caption{DNA replication: the double helix is unzipped and each strand copied.\\DNA复制：双螺旋被拉开，每条链都被复制。}
\end{figure}

\begin{quote}
\noindent\textbf{原文} There remained, however, a single missing piece in the double helical jigsaw puzzle: our unzipping idea for DNA replication had yet to be experimentally verified. Max Delbrück, for example, was unconvinced. Though he liked the double helix as a model, he worried that unzipping it might generate horrible knots. Five years later, a former student of Pauling's, Matt Meselson, and the equally bright young phage worker Frank Stahl put to rest such fears when they published the results of a single elegant experiment.

\smallskip
\noindent\textbf{译文} 然而，双螺旋拼图中仍缺少一块：我们关于DNA复制时“拉开拉链”的想法尚未得到实验证实。例如马克斯·德尔布吕克并不信服。虽然他喜欢双螺旋作为模型，却担心拉开它可能产生可怕的结。五年后，鲍林昔日学生马特·梅塞尔森和同样聪明的年轻噬菌体研究者弗兰克·斯塔尔，发表了一项优雅实验的结果，平息了这种担忧。
\end{quote}

\begin{quote}
\noindent\textbf{原文} They had met in the summer of 1954 at the Marine Biological Laboratory at Woods Hole, Massachusetts, where I was then lecturing, and agreed---over a good many gin martinis---that they should get together to do some science. The result of their collaboration has been described as "the most beautiful experiment in biology."

\smallskip
\noindent\textbf{译文} 他们1954年夏天在马萨诸塞州伍兹霍尔海洋生物学实验室相遇，当时我在那里讲课。两人在许多杯金酒马提尼之后达成共识：他们应该合作做点科学。他们合作的结果后来被称为“生物学中最美丽的实验”。
\end{quote}

\begin{figure}[htbp]
\centering
\safeincludeimage{72946d8f3811ea6bdd57f651161bc59556e5a1d901a9dae3b0247b825b8790ac.jpg}
\caption{The Meselson--Stahl experiment.\\梅塞尔森--斯塔尔实验。}
\end{figure}

\begin{quote}
\noindent\textbf{原文} They used a centrifugation technique that allowed them to sort molecules according to slight differences in weight; following a centrifugal spin, heavier molecules end up nearer the bottom of the test tube than lighter ones. Because nitrogen atoms (N) are a component of DNA, and because they exist in two distinct forms, one light and one heavy, Meselson and Stahl were able to tag segments of DNA and thereby track the process of its replication in bacteria. Initially all the bacteria were raised in a medium containing heavy N, which was thus incorporated in both strands of the DNA. From this culture they took a sample, transferring it to a medium containing only light N, ensuring that the next round of DNA replication would have to make use of light N. If, as Crick and I had predicted, DNA replication involves unzipping the double helix and copying each strand, the resultant two "daughter" DNA molecules in the experiment would be hybrids, each consisting of one heavy N strand, the template strand derived from the "parent" molecule, and one light N strand, the one newly fabricated from the new medium. Meselson and Stahl's centrifugation procedure bore out these expectations precisely. They found three discrete bands in their centrifuge tubes, with the heavy-then-light sample halfway between the heavy-heavy and light-light samples. DNA replication works just as our model supposed it would.

\smallskip
\noindent\textbf{译文} 他们使用一种离心技术，可以按照重量的细微差异对分子排序；离心旋转后，较重分子最终会比轻分子更靠近试管底部。氮原子（N）是DNA组成部分，而氮存在两种不同形式，一轻一重，因此梅塞尔森和斯塔尔能够标记DNA片段，从而追踪细菌中DNA复制过程。最初，所有细菌都在含重氮的培养基中培养，因此重氮被纳入DNA两条链。随后他们从该培养物中取样，转移到只含轻氮的培养基中，确保下一轮DNA复制必须使用轻氮。如果如克里克和我所预言的那样，DNA复制涉及拉开双螺旋并复制每条链，那么实验中产生的两个“子代”DNA分子将是杂合体，每个都由一条重氮链，即来自“亲本”分子的模板链，和一条轻氮链，即用新培养基新合成的链，组成。梅塞尔森和斯塔尔的离心程序精确证实了这些预期。他们在离心管中发现三条离散条带，其中先重后轻的样品位于重-重样品与轻-轻样品之间。DNA复制正如我们的模型所设想的那样运作。
\end{quote}

\begin{figure}[htbp]
\centering
\safeincludeimage{b247dd189a7c3454e931f15463bd16d36aa515299235b26261281887ab7814de.jpg}
\caption{Matt Meselson beside an ultracentrifuge, the hardware at the heart of "the most beautiful experiment in biology."\\马特·梅塞尔森站在超速离心机旁；这台设备是“生物学中最美丽的实验”的核心硬件。}
\end{figure}

\begin{quote}
\noindent\textbf{原文} The biochemical nuts and bolts of DNA replication were being analyzed at around the same time in Arthur Kornberg's laboratory at Washington University in St. Louis. By developing a new, "cell-free" system for DNA synthesis, Kornberg discovered an enzyme, DNA polymerase, that links the DNA components and makes the chemical bonds of the DNA backbone. Kornberg's enzymatic synthesis of DNA was such an unanticipated and important event that he was awarded the 1959 Nobel Prize in Physiology or Medicine, less than two years after the key experiments. After his prize was announced, Kornberg was photographed holding a copy of the double helix model I had taken to Cold Spring Harbor in 1953.

\smallskip
\noindent\textbf{译文} 大约同一时期，圣路易斯华盛顿大学阿瑟·科恩伯格实验室正在分析DNA复制的生化细节。通过开发一种新的“无细胞”DNA合成系统，科恩伯格发现了一种酶，即DNA聚合酶，它把DNA组分连接起来，并形成DNA骨架的化学键。科恩伯格的DNA酶促合成是如此出人意料而重要，以至于关键实验后不到两年，他就获得了1959年诺贝尔生理学或医学奖。获奖公布后，科恩伯格被拍到手持一份我1953年带到冷泉港的双螺旋模型。
\end{quote}

\begin{quote}
\noindent\textbf{原文} It was not until 1962 that Francis Crick, Maurice Wilkins, and I were to receive our own Nobel Prize in Physiology or Medicine. Four years earlier, Rosalind Franklin had died of ovarian cancer at the tragically young age of thirty-seven. Before then Crick had become a close colleague and a real friend of Franklin's. Following the two operations that would fail to stem the advance of her cancer, Franklin convalesced with Crick and his wife, Odile, in Cambridge.

\smallskip
\noindent\textbf{译文} 直到1962年，弗朗西斯·克里克、莫里斯·威尔金斯和我才获得我们自己的诺贝尔生理学或医学奖。四年前，罗莎琳德·富兰克林因卵巢癌去世，年仅三十七岁，令人痛惜。在那之前，克里克已成为富兰克林亲近的同事和真正的朋友。在两次未能阻止癌症进展的手术之后，富兰克林在剑桥与克里克及其妻子奥迪尔一起休养。
\end{quote}

\begin{figure}[htbp]
\centering
\safeincludeimage{64e65b4d42b8972cb02eb887296fc9edba05b01aeb760bb5593aa27a5d31f6b8.jpg}
\caption{Arthur Kornberg at the time of winning his Nobel Prize.\\阿瑟·科恩伯格获诺贝尔奖时。}
\end{figure}

\begin{quote}
\noindent\textbf{原文} It was and remains a long-standing rule of the Nobel Committee never to split a single prize more than three ways. Had Franklin lived, the problem would have arisen whether to bestow the award upon her or Maurice Wilkins. The Swedes might have resolved the dilemma by awarding them both the Nobel Prize in Chemistry that year. Instead, it went to Max Perutz and John Kendrew, who had elucidated the three-dimensional structures of hemoglobin and myoglobin respectively.

\smallskip
\noindent\textbf{译文} 诺贝尔委员会长期以来一直有一条规则，且至今如此：单项奖不能分给超过三个人。如果富兰克林还活着，就会出现该把奖授予她还是授予莫里斯·威尔金斯的问题。瑞典人或许可以通过把那年的诺贝尔化学奖同时授予他们二人来解决困境。但事实上，化学奖授予了马克斯·佩鲁茨和约翰·肯德鲁，他们分别阐明了血红蛋白和肌红蛋白的三维结构。
\end{quote}

\begin{quote}
\noindent\textbf{原文} The discovery of the double helix sounded the death knell for vitalism. Serious scientists, even those religiously inclined, realized that a complete understanding of life would not require the revelation of new laws of nature. Life was just a matter of physics and chemistry, albeit exquisitely organized physics and chemistry. The immediate task ahead would be to figure out how the DNA-encoded script of life went about its work. How does the molecular machinery of cells read the messages of DNA molecules? As the next chapter will reveal, the unexpected complexity of the reading mechanism led to profound insights into how life first came about.

\smallskip
\noindent\textbf{译文} 双螺旋的发现为活力论敲响了丧钟。严肃科学家，即使是有宗教倾向的科学家，也认识到，对生命的完整理解并不需要揭示新的自然法则。生命只是物理和化学的问题，尽管是组织得极其精巧的物理和化学。眼前的任务是弄清由DNA编码的生命脚本如何发挥作用。细胞的分子机器怎样读取DNA分子的信息？正如下章将揭示的，读取机制出人意料的复杂性，带来了关于生命最初如何出现的深刻洞见。
\end{quote}

\begin{figure}[htbp]
\centering
\safeincludeimagepair{dc1e2620c73dd956aa24b0dee5de0c025523e58fa15db48333e2b1f5e14dd38a.jpg}{28ba4827aae4f32c1a21d821087caa028bb2cd3bbbe88b6ca80cb819202d9c1f.jpg}
\hfill
\safeincludeimage{d48c5f92ded01c4d69e1782a28c8bfeeb9b00a2845d59755a7d38fec885397bf.jpg}
\caption{Uncaptioned terminal image references retained from the Watson/Berry source boundary.\\保留Watson/Berry源文本边界处未配图注的篇末图像引用。}
\end{figure}


% ===== 06_carson.tex =====
\section*{Earth's Green Mantle / 大地的绿色披覆}
\begin{flushright}
Rachel Carson, \textit{Silent Spring}
\end{flushright}

\noindent\textbf{原文} \ldots{} the booming ``weed killer'' business of the present day, in which soaring sales and expanding uses mark the production of plant-killing chemicals.
\par
\noindent\textbf{译文} ……当今蓬勃发展的“除草剂”产业：销量节节攀升，用途不断扩大，推动着各种杀灭植物的化学制剂大量生产。
\par\medskip

\noindent\textbf{原文} One of the most tragic examples of our unthinking bludgeoning of the landscape is to be seen in the sagebrush lands of the West, where a vast campaign is on to destroy the sage and to substitute grasslands. If ever an enterprise needed to be illuminated with a sense of the history and meaning of the landscape, it is this. For here the natural landscape is eloquent of the interplay of forces that have created it. It is spread before us like the pages of an open book in which we can read why the land is what it is, and why we should preserve its integrity. But the pages lie unread.
\par
\noindent\textbf{译文} 我们不假思索地重击大地景观，其中最悲惨的例子之一，就出现在西部的蒿灌丛地带：一场大规模运动正在进行，要消灭蒿灌木，代之以草地。倘若有哪项事业尤其需要以对景观历史与意义的理解来照亮，那正是这一项。因为这里的自然景观清楚讲述了塑造它的诸力如何相互作用。它像一本摊开的书展现在我们面前，使我们能够读出土地为何如此，也读出为什么应当保存它的完整性。可是这些书页无人阅读。
\par\medskip

\noindent\textbf{原文} The land of the sage is the land of the high western plains and the lower slopes of the mountains that rise above them, a land born of the great uplift of the Rocky Mountain system many millions of years ago. It is a place of harsh extremes of climate: of long winters when blizzards drive down from the mountains and snow lies deep on the plains, of summers whose heat is relieved by only scanty rains, with drought biting deep into the soil, and drying winds stealing moisture from leaf and stem.
\par
\noindent\textbf{译文} 蒿的土地，就是西部高原和从其上拔起的山脉低坡；这片土地诞生于数百万年前落基山系的巨大隆升。这里气候严酷而极端：漫长的冬季，暴风雪从山中压下，平原积雪深厚；夏季炎热，只有稀少雨水稍作缓解，干旱深深咬入土壤，干风从叶片和茎秆中偷走水分。
\par\medskip

\noindent\textbf{原文} As the landscape evolved, there must have been a long period of trial and error in which plants attempted the colonization of this high and windswept land. One after another must have failed. At last one group of plants evolved which combined all the qualities needed to survive. The sage---low-growing and shrubby---could hold its place on the mountain slopes and on the plains, and within its small gray leaves it could hold moisture enough to defy the thieving winds. It was no accident, but rather the result of long ages of experimentation by nature, that the great plains of the West became the land of the sage.
\par
\noindent\textbf{译文} 在景观演化的过程中，植物试图定居这片高寒、多风的土地，必定经历过漫长的试错。一种又一种植物想必都失败了。最后，一类植物演化出来，兼具生存所需的一切品质。蒿灌木低矮而丛生，能够在山坡和平原上站稳脚跟，并能在小小的灰色叶片里保存足够的水分，抵抗偷水的风。西部大平原成为蒿的土地，并非偶然，而是自然经历漫长年代实验的结果。
\par\medskip

\noindent\textbf{原文} Along with the plants, animal life, too, was evolving in harmony with the searching requirements of the land. In time there were two as perfectly adjusted to their habitat as the sage. One was a mammal, the fleet and graceful pronghorn antelope. The other was a bird, the sage grouse---the ``cock of the plains'' of Lewis and Clark.
\par
\noindent\textbf{译文} 与植物一道，动物生命也在演化，并与这片土地苛刻的要求相协调。久而久之，有两种动物像蒿一样完美地适应了它们的栖息地。一种是哺乳动物，迅捷而优雅的叉角羚；另一种是鸟类，蒿原松鸡，也就是刘易斯和克拉克所称的“平原雄鸡”。
\par\medskip

\noindent\textbf{原文} The sage and the grouse seem made for each other. The original range of the bird coincided with the range of the sage, and as the sagelands have been reduced, so the populations of grouse have dwindled. The sage is all things to these birds of the plains. The low sage of the foothill ranges shelters their nests and their young; the denser growths are loafing and roosting areas; at all times the sage provides the staple food of the grouse. Yet it is a two-way relationship. The spectacular courtship displays of the cocks help loosen the soil beneath and around the sage, aiding invasion by grasses which grow in the shelter of sagebrush.
\par
\noindent\textbf{译文} 蒿和松鸡仿佛天生相属。蒿原松鸡最初的分布范围与蒿灌丛的范围相吻合；随着蒿地缩减，松鸡数量也随之下降。对这些平原之鸟来说，蒿就是一切。山麓地带低矮的蒿庇护它们的巢和幼雏；更浓密的蒿丛是它们闲栖和夜宿之所；无论何时，蒿都是松鸡的主食。然而，这也是一种双向关系。雄鸟壮观的求偶炫耀，会松动蒿下及其周围的土壤，帮助草类侵入并在蒿灌木庇护下生长。
\par\medskip

\noindent\textbf{原文} The antelope, too, have adjusted their lives to the sage. They are primarily animals of the plains, and in winter when the first snows come those that have summered in the mountains move down to the lower elevations. There the sage provides the food that tides them over the winter. Where all other plants have shed their leaves, the sage remains evergreen, the gray-green leaves---bitter, aromatic, rich in proteins, fats, and needed minerals---clinging to the stems of the dense and shrubby plants. Though the snows pile up, the tops of the sage remain exposed, or can be reached by the sharp, pawing hoofs of the antelope. Then grouse feed on them too, finding them on bare and windswept ledges or following the antelope to feed where they have scratched away the snow.
\par
\noindent\textbf{译文} 叉角羚也把自己的生活调适到蒿的节律之中。它们主要是平原动物；冬天初雪降临时，在山地度过夏季的羚群便下移到较低海拔。那里，蒿提供食物，帮助它们熬过冬季。其他植物都已落叶，蒿却依然常绿，灰绿色的叶片苦涩、芳香，富含蛋白质、脂肪和必需矿物质，紧附在浓密丛生的茎上。即使积雪堆深，蒿梢仍会露出，或可由叉角羚锋利、刨挖的蹄子触及。于是松鸡也以其为食，在裸露而风扫的岩脊上寻找，或跟随叉角羚，在它们刨开雪的地方觅食。
\par\medskip

\noindent\textbf{原文} And other life looks to the sage. Mule deer often feed on it. Sage may mean survival for winter-grazing livestock. Sheep graze many winter ranges where the big sagebrush forms almost pure stands. For half the year it is their principal forage, a plant of higher energy value than even alfalfa hay.
\par
\noindent\textbf{译文} 还有其他生命仰赖蒿。骡鹿常以它为食。对于冬季放牧的牲畜来说，蒿可能意味着生存。许多冬季牧场上，大蒿几乎形成纯林，羊群就在其中采食。一年中有半年，蒿是它们的主要饲草，其能量价值甚至高于苜蓿干草。
\par\medskip

\noindent\textbf{原文} The bitter upland plains, the purple wastes of sage, the wild, swift antelope, and the grouse are then a natural system in perfect balance. Are? The verb must be changed---at least in those already vast and growing areas where man is attempting to improve on nature's way. In the name of progress the land management agencies have set about to satisfy the insatiable demands of the cattlemen for more grazing land. By this they mean grassland---grass without sage. So in a land which nature found suited to grass growing mixed with and under the shelter of sage, it is now proposed to eliminate the sage and create unbroken grassland. Few seem to have asked whether grasslands are a stable and desirable goal in this region. Certainly nature's own answer was otherwise. The annual precipitation in this land where the rains seldom fall is not enough to support good sod-forming grass; it favors rather the perennial bunch grass that grows in the shelter of the sage.
\par
\noindent\textbf{译文} 苦寒的高地平原、紫色的蒿海、野性而迅疾的叉角羚以及松鸡，原本构成一个完美平衡的自然系统。原本？这个动词必须改掉，至少在那些已经广阔且仍在扩大的地区是如此，因为人在那里试图改进自然之道。以进步之名，土地管理机构开始满足牧牛人对更多牧场无止境的要求。他们所谓的牧场，指的是草地，只有草、没有蒿。于是，在自然认为适宜草与蒿混生、并在蒿的庇护下生长的土地上，如今有人提议消灭蒿，创造连绵不断的草地。似乎很少有人问过，在这一地区，草地是否真是稳定而可取的目标。自然自身的回答显然并非如此。在这片少雨之地，年降水量不足以支撑结成良好草皮的禾草；它更有利于在蒿庇护下生长的多年生丛生禾草。
\par\medskip

\noindent\textbf{原文} Yet the program of sage eradication has been under way for a number of years. Several government agencies are active in it; industry has joined with enthusiasm to promote and encourage an enterprise which creates expanded markets not only for grass seed but for a large assortment of machines for cutting and plowing and seeding. The newest addition to the weapons is the use of chemical sprays. Now millions of acres of sagebrush lands are sprayed each year.
\par
\noindent\textbf{译文} 然而，根除蒿的计划已经推行多年。若干政府机构积极参与；工业界也热情加入，推动和鼓励这项事业，因为它不仅为草籽，也为各种割灌、犁耕和播种机械开辟了更大的市场。最新加入武器库的是化学喷洒。如今，每年都有数百万英亩蒿灌丛地被喷洒药剂。
\par\medskip

\noindent\textbf{原文} What are the results? The eventual effects of eliminating sage and seeding with grass are largely conjectural. Men of long experience with the ways of the land say that in this country there is better growth of grass between and under the sage than can possibly be had in pure stands, once the moisture-holding sage is gone.
\par
\noindent\textbf{译文} 结果如何？消灭蒿并播种牧草的最终影响，在很大程度上仍属推测。那些长期熟悉这片土地规律的人说，在这个地区，草在蒿丛之间和蒿下生长得更好；一旦能够保持水分的蒿消失，纯草地绝不可能达到同样的长势。
\par\medskip

\noindent\textbf{原文} But even if the program succeeds in its immediate objective, it is clear that the whole closely knit fabric of life has been ripped apart. The antelope and the grouse will disappear along with the sage. The deer will suffer, too, and the land will be poorer for the destruction of the wild things that belong to it. Even the livestock which are the intended beneficiaries will suffer; no amount of lush green grass in summer can help the sheep starving in the winter storms for lack of the sage and bitterbrush and other wild vegetation of the plains.
\par
\noindent\textbf{译文} 但是，即使这项计划实现了眼前目标，也很清楚，整个紧密编织的生命之网已经被撕裂。叉角羚和松鸡将随蒿一起消失。鹿也会受害；那些原本属于这片土地的野生生命被毁灭，土地也将变得更贫乏。即便是本应受益的牲畜也会遭殃；夏季再茂盛的青草，也无法帮助羊群在冬季风雪中挨饿，因为它们失去了蒿、苦灌木以及平原上的其他野生植被。
\par\medskip

\noindent\textbf{原文} These are the first and obvious effects. The second is of a kind that is always associated with the shotgun approach to nature: the spraying also eliminates a great many plants that were not its intended target. Justice William O. Douglas, in his recent book \textit{My Wilderness: East to Katahdin}, has told of an appalling example of ecological destruction wrought by the United States Forest Service in the Bridger National Forest in Wyoming. Some 10,000 acres of sagelands were sprayed by the Service, yielding to pressure of cattlemen for more grasslands. The sage was killed, as intended. But so was the green, life-giving ribbon of willows that traced its way across these plains, following the meandering streams. Moose had lived in these willow thickets, for willow is to the moose what sage is to the antelope. Beavers had lived there, too, feeding on the willows, felling them and making a strong dam across the tiny stream. Through the labor of the beavers, a lake backed up. Trout in the mountain streams seldom were more than six inches long; in the lake they thrived so prodigiously that many grew to five pounds. Waterfowl were attracted to the lake, also. Merely because of the presence of the willows and the beavers that depended on them, the region was an attractive recreational area with excellent fishing and hunting.
\par
\noindent\textbf{译文} 这些是最初且显而易见的影响。第二类影响，则总是与对自然采取霰弹枪式的粗暴办法相伴而来：喷洒还会消灭许多并非目标的植物。威廉·O·道格拉斯大法官在近作《我的荒野：从东部到卡塔丁》中，讲述了美国林务局在怀俄明州布里杰国家森林造成生态破坏的一个惊人例子。林务局屈从牧牛人要求更多草地的压力，喷洒了约一万英亩蒿地。蒿按计划被杀死了。但那条沿蜿蜒溪流横贯平原、带来生命的绿色柳带也一同死去。驼鹿曾生活在这些柳丛中，因为柳树之于驼鹿，正如蒿之于叉角羚。河狸也曾在那里生活，以柳为食，伐倒柳树，在小溪上筑起坚固堤坝。由于河狸的劳动，水积成了湖。山溪中的鳟鱼很少超过六英寸；在湖中，它们却繁盛异常，许多长到五磅。水禽也被湖吸引而来。仅仅因为有了柳树和依赖柳树的河狸，这一地区便成为有优良垂钓与狩猎条件、富有吸引力的休闲地。
\par\medskip

\noindent\textbf{原文} But with the ``improvement'' instituted by the Forest Service, the willows went the way of the sagebrush, killed by the same impartial spray. When Justice Douglas visited the area in 1959, the year of the spraying, he was shocked to see the shriveled and dying willows---the ``vast, incredible damage.'' What would become of the moose? Of the beavers and the little world they had constructed? A year later he returned to read the answers in the devastated landscape. The moose were gone and so were the beavers. Their principal dam had gone out for want of attention by its skilled architects, and the lake had drained away. None of the large trout were left. None could live in the tiny creek that remained, threading its way through a bare, hot land where no shade remained. The living world was shattered.
\par
\noindent\textbf{译文} 然而，随着林务局推行的“改良”，柳树也步了蒿灌木的后尘，被同一种不分青红皂白的喷剂杀死。道格拉斯大法官在1959年，也就是喷洒当年访问该地时，看到枯缩将死的柳树，看到那种“巨大而难以置信的损害”，深感震惊。驼鹿会怎样？河狸和它们构筑的那个小世界又会怎样？一年后，他重返此地，在被毁坏的景观中读到了答案。驼鹿不见了，河狸也不见了。它们的主坝因缺少熟练建筑师的维护而垮掉，湖水也流干。大型鳟鱼一条也没有留下。残存的小溪穿过光秃、炎热、毫无遮阴的土地，没有一条大鳟鱼能在那里活下去。这个有生命的世界被粉碎了。
\par\medskip

\noindent\textbf{原文} Besides the more than four million acres of rangelands sprayed each year, tremendous areas of other types of land are also potential or actual recipients of chemical treatments for weed control. For example, an area larger than all of New England---some 50 million acres---is under management by utility corporations and much of it is routinely treated for ``brush control.'' In the Southwest an estimated 75 million acres of mesquite lands require management by some means, and chemical spraying is the method most actively pushed. An unknown but very large acreage of timber-producing lands is now aerially sprayed in order to ``weed out'' the hardwoods from the more spray-resistant conifers. Treatment of agricultural lands with herbicides doubled in the decade following 1949, totaling 53 million acres in 1959. And the combined acreage of private lawns, parks, and golf courses now being treated must reach an astronomical figure.
\par
\noindent\textbf{译文} 除了每年被喷洒的四百多万英亩牧地之外，其他类型的巨大面积土地，也正在或可能正在接受化学除草处理。比如，由公用事业公司管理的土地面积超过整个新英格兰，约五千万英亩，其中很大一部分被例行用于“灌丛控制”。在西南部，估计有七千五百万英亩牧豆树地需要以某种方式管理，而化学喷洒是被最积极推行的方法。另有面积不明但极为广阔的产材林地，如今正通过空中喷洒来从较耐喷剂的针叶树中“剔除”阔叶树。1949年以后的十年间，农地使用除草剂处理的面积翻了一番，到1959年总计达五千三百万英亩。至于如今接受处理的私人草坪、公园和高尔夫球场，合计面积必定达到天文数字。
\par\medskip

\noindent\textbf{原文} The chemical weed killers are a bright new toy. They work in a spectacular way; they give a giddy sense of power over nature to those who wield them, and as for the long-range and less obvious effects---these are easily brushed aside as the baseless imaginings of pessimists. The ``agricultural engineers'' speak blithely of ``chemical plowing'' in a world that is urged to beat its plowshares into spray guns. The town fathers of a thousand communities lend willing ears to the chemical salesman and the eager contractors who will rid the roadsides of ``brush''---for a price. It is cheaper than mowing, is the cry. So, perhaps, it appears in the neat rows of figures in the official books; but were the true costs entered, the costs not only in dollars but in the many equally valid debits we shall presently consider, the wholesale broadcasting of chemicals would be seen to be more costly in dollars as well as infinitely damaging to the long-range health of the landscape and to all the varied interests that depend on it.
\par
\noindent\textbf{译文} 化学除草剂是一件闪亮的新玩具。它们起效的方式令人瞩目；使用者会因此产生一种对自然拥有权力的眩晕感。至于长期而不那么明显的后果，则很容易被斥为悲观者毫无根据的想象。“农业工程师”轻快地谈论“化学耕作”，仿佛世界正被敦促把犁铧改铸成喷枪。成千上万社区的镇政要人乐于听信化学品推销员和急切承包商的话；这些人愿意替道路两旁清除“灌丛”，只要付钱即可。呼声说，这比割草便宜。也许在官方账册整齐排列的数字中，看起来确是如此；可是如果把真实成本列入账内，不仅是美元成本，也包括我们马上要考量的那些同样有效的借项，那么，大规模撒布化学品就会显得在金钱上更昂贵，并且对景观的长期健康、对依赖景观的各种利益造成无穷损害。
\par\medskip

\noindent\textbf{原文} Take, for instance, that commodity prized by every chamber of commerce throughout the land---the good will of vacationing tourists. There is a steadily growing chorus of outraged protest about the disfigurement of once beautiful roadsides by chemical sprays, which substitute a sere expanse of brown, withered vegetation for the beauty of fern and wildflower, of native shrubs adorned with blossom or berry. ``We are making a dirty, brown, dying-looking mess along the sides of our roads,'' a New England woman wrote angrily to her newspaper. ``This is not what the tourists expect, with all the money we are spending advertising the beautiful scenery.''
\par
\noindent\textbf{译文} 例如，全国每个商会都珍视的一项资产：度假游客的好感。对化学喷洒毁坏昔日美丽路边景观的愤怒抗议，正日益汇成合唱；喷洒以一片焦枯、褐色、萎败的植被，取代了蕨类和野花之美，取代了开花或结莓的本地灌木之美。一位新英格兰妇女愤怒地写信给报纸说：“我们正在把道路两旁弄成肮脏、褐色、垂死模样的一团糟。我们花了那么多钱宣传美丽风景，游客期待看到的可不是这个。”
\par\medskip

\noindent\textbf{原文} In the summer of 1960 conservationists from many states converged on a peaceful Maine island to witness its presentation to the National Audubon Society by its owner, Millicent Todd Bingham. The focus that day was on the preservation of the natural landscape and of the intricate web of life whose interwoven strands lead from microbes to man. But in the background of all the conversations among the visitors to the island was indignation at the despoiling of the roads they had traveled. Once it had been a joy to follow those roads through the evergreen forests, roads lined with bayberry and sweet fern, alder and huckleberry. Now all was brown desolation. One of the conservationists wrote of that August pilgrimage to a Maine island: ``I returned \ldots{} angry at the desecration of the Maine roadsides. Where, in previous years, the highways were bordered with wildflowers and attractive shrubs, there were only the scars of dead vegetation for mile after mile \ldots{}. As an economic proposition, can Maine afford the loss of tourist goodwill that such sights induce?''
\par
\noindent\textbf{译文} 1960年夏，来自许多州的自然保护人士聚集到缅因州一座宁静小岛，见证岛主米利森特·托德·宾厄姆将其赠予全国奥杜邦学会。那一天的焦点，是保存自然景观，保存那张从微生物一直连到人类、千丝万缕相互交织的生命之网。然而，访客们在岛上所有交谈的背景里，都有一种对来路遭到破坏的愤慨。过去，沿着那些道路穿过常绿森林曾是一种喜悦，道旁生长着杨梅属灌木、香蕨木、桤木和越橘。如今，一切都是褐色荒凉。一位保护人士写到那年八月前往缅因小岛的旅程：“我回来时……为缅因路边遭到亵渎而愤怒。往年，公路边满是野花和迷人的灌木；如今一英里又一英里，只剩死去植被留下的伤疤……从经济角度看，缅因能承受这种景象导致的游客好感流失吗？”
\par\medskip

\noindent\textbf{原文} Maine roadsides are merely one example, though a particularly sad one for those of us who have a deep love for the beauty of that state, of the senseless destruction that is going on in the name of roadside brush control throughout the nation.
\par
\noindent\textbf{译文} 缅因的路边不过是一个例子；对于我们这些深深爱着该州之美的人来说，这是一个尤其令人悲伤的例子。全国各地以道路灌丛控制之名进行的无谓破坏，正在不断发生。
\par\medskip

\noindent\textbf{原文} Botanists at the Connecticut Arboretum declare that the elimination of beautiful native shrubs and wildflowers has reached the proportions of a ``roadside crisis.'' Azaleas, mountain laurel, blueberries, huckleberries, viburnums, dogwood, bayberry, sweet fern, low shadbush, winterberry, chokecherry, and wild plum are dying before the chemical barrage. So are the daisies, black-eyed Susans, Queen Anne's lace, goldenrods, and fall asters which lend grace and beauty to the landscape.
\par
\noindent\textbf{译文} 康涅狄格植物园的植物学家宣称，消灭美丽本地灌木和野花的做法，已经发展到“路边危机”的程度。杜鹃、山月桂、蓝莓、越橘、荚蒾、山茱萸、杨梅属灌木、香蕨木、矮唐棣、冬青、稠李和野李，都在化学弹幕前死去。雏菊、黑心金光菊、野胡萝卜、秋麒麟草和秋紫苑也同样如此；它们本来为景观增添优雅与美丽。
\par\medskip

\noindent\textbf{原文} The spraying is not only improperly planned but studded with abuses such as these. In a southern New England town one contractor finished his work with some chemical remaining in his tank. He discharged this along woodland roadsides where no spraying had been authorized. As a result the community lost the blue and golden beauty of its autumn roads, where asters and goldenrod would have made a display worth traveling far to see. In another New England community a contractor changed the state specifications for town spraying without the knowledge of the highway department and sprayed roadside vegetation to a height of eight feet instead of the specified maximum of four feet, leaving a broad, disfiguring, brown swath. In a Massachusetts community the town officials purchased a weed killer from a zealous chemical salesman, unaware that it contained arsenic. One result of the subsequent roadside spraying was the death of a dozen cows from arsenic poisoning.
\par
\noindent\textbf{译文} 喷洒不仅规划不当，还充斥着这类滥用。在新英格兰南部一座城镇，一名承包商完工后罐中还剩一些药剂，便把它沿林地路边排放，而那里并未获准喷洒。结果，这个社区失去了秋日道路的蓝金之美；原本紫苑和秋麒麟草会在那里形成一场值得远道而来的展示。在另一个新英格兰社区，一名承包商在公路部门不知情的情况下改变了州里关于城镇喷洒的规格，把路边植被喷到八英尺高，而规定的最高限度是四英尺，留下了一条宽阔、丑陋的褐色带状伤痕。在马萨诸塞州一个社区，镇官员从一名热心的化学品推销员那里购买除草剂，却不知道其中含砷。随后进行的路边喷洒，造成十几头牛死于砷中毒。
\par\medskip

\noindent\textbf{原文} Trees within the Connecticut Arboretum Natural Area were seriously injured when the town of Waterford sprayed the roadsides with chemical weed killers in 1957. Even large trees not directly sprayed were affected. The leaves of the oaks began to curl and turn brown, although it was the season for spring growth. Then new shoots began to be put forth and grew with abnormal rapidity, giving a weeping appearance to the trees. Two seasons later, large branches on these trees had died, others were without leaves, and the deformed, weeping effect of whole trees persisted.
\par
\noindent\textbf{译文} 1957年，沃特福德镇在路边喷洒化学除草剂时，康涅狄格植物园自然区内的树木受到严重伤害。即使没有被直接喷到的大树也受到影响。橡树叶片开始卷曲、变褐，尽管那正是春季生长时节。随后，新梢开始抽出，并以异常速度生长，使树呈现垂泣般的外观。两个生长季之后，这些树上的大枝已经死亡，另一些枝条没有叶片，整株树畸形下垂的状态仍然存在。
\par\medskip

\noindent\textbf{原文} I know well a stretch of road where nature's own landscaping has provided a border of alder, viburnum, sweet fern, and juniper with seasonally changing accents of bright flowers, or of fruits hanging in jeweled clusters in the fall. The road had no heavy load of traffic to support; there were few sharp curves or intersections where brush could obstruct the driver's vision. But the sprayers took over and the miles along that road became something to be traversed quickly, a sight to be endured with one's mind closed to thoughts of the sterile and hideous world we are letting our technicians make. But here and there authority had somehow faltered and by an unaccountable oversight there were oases of beauty in the midst of austere and regimented control---oases that made the desecration of the greater part of the road the more unbearable. In such places my spirit lifted to the sight of the drifts of white clover or the clouds of purple vetch with here and there the flaming cup of a wood lily.
\par
\noindent\textbf{译文} 我很熟悉一段道路，在那里，自然自己的造景沿路铺设了一道桤木、荚蒾、香蕨木和刺柏的边界，季节变换时，明亮的花朵点缀其间；到了秋天，成串果实像珠宝般悬挂。那条路并不承载繁重交通；急弯和交叉路口也很少，灌丛几乎不会遮挡驾驶者视线。可是喷洒者接管之后，沿途数英里变成了人们只想尽快穿过的地方，成为一种必须忍受的景象；只有闭上心，不去想我们正让技术人员制造怎样一个贫瘠而丑陋的世界，才能承受。然而，这里那里，权威不知怎的失手了，由于某种无法解释的疏忽，在严苛而整齐划一的控制中留下了美的绿洲；这些绿洲反倒使道路大部分遭受的亵渎更加难以忍受。在这样的地方，看到一片片白三叶草，或一团团紫色野豌豆云雾，其间偶尔燃起一盏林百合的火焰杯盏，我的精神便随之振奋。
\par\medskip

\noindent\textbf{原文} Such plants are ``weeds'' only to those who make a business of selling and applying chemicals. In a volume of \textit{Proceedings} of one of the weed-control conferences that are now regular institutions, I once read an extraordinary statement of a weed killer's philosophy. The author defended the killing of good plants ``simply because they are in bad company.'' Those who complain about killing wildflowers along roadsides reminded him, he said, of antivivisectionists ``to whom, if one were to judge by their actions, the life of a stray dog is more sacred than the lives of children.''
\par
\noindent\textbf{译文} 这些植物之所以是“杂草”，只是在那些以销售和施用化学品为业的人眼里如此。我曾在一本除草会议《论文集》中读到一段非同寻常的除草哲学宣言；这类会议如今已成为常设制度。作者为杀死好植物辩护，说理由“只不过是它们交了坏朋友”。他说，那些抱怨路边野花遭杀灭的人，让他想到反对动物活体实验者；“如果按他们的行为来判断，在他们眼中，一条流浪狗的生命比儿童的生命还神圣。”
\par\medskip

\noindent\textbf{原文} To the author of this paper, many of us would unquestionably be suspect, convicted of some deep perversion of character because we prefer the sight of the vetch and the clover and the wood lily in all their delicate and transient beauty to that of roadsides scorched as by fire, the shrubs brown and brittle, the bracken that once lifted high its proud lacework now withered and drooping. We would seem deplorably weak that we can tolerate the sight of such ``weeds,'' that we do not rejoice in their eradication, that we are not filled with exultation that man has once more triumphed over miscreant nature.
\par
\noindent\textbf{译文} 在这篇论文的作者看来，我们许多人无疑都很可疑，必定会因某种深层性格败坏而被定罪：因为我们宁愿看见野豌豆、白三叶草和林百合那纤细而短暂的美，也不愿看见道路两旁像被火烧过一样，灌木褐黄而脆裂，曾经高举骄傲蕾丝花纹的蕨草如今枯萎垂落。我们似乎软弱得可悲，竟能容忍这些“杂草”的景象，竟不为它们被根除而欢欣，竟不会因人类又一次战胜作恶的自然而充满狂喜。
\par\medskip

\noindent\textbf{原文} Justice Douglas tells of attending a meeting of federal field men who were discussing protests by citizens against plans for the spraying of sagebrush that I mentioned earlier in this chapter. These men considered it hilariously funny that an old lady had opposed the plan because the wildflowers would be destroyed. ``Yet, was not her right to search out a banded cup or a tiger lily as inalienable as the right of stockmen to search out grass or of a lumberman to claim a tree?'' asks this humane and perceptive jurist. ``The esthetic values of the wilderness are as much our inheritance as the veins of copper and gold in our hills and the forests in our mountains.''
\par
\noindent\textbf{译文} 道格拉斯大法官讲述过，他参加过一次联邦外勤人员会议，会上讨论公民对本章前面提到的喷洒蒿灌木计划的抗议。那些人觉得，一位老妇人因为野花会被毁灭而反对该计划，这件事滑稽极了。这位富有人道精神且洞察敏锐的法学家问道：“然而，她寻找斑纹杯花或虎百合的权利，难道不和牧场主寻找牧草的权利、木材商要求一棵树的权利同样不可剥夺吗？”“荒野的审美价值，正如山中铜金矿脉和山上森林一样，都是我们的遗产。”
\par\medskip

\noindent\textbf{原文} There is of course more to the wish to preserve our roadside vegetation than even such esthetic considerations. In the economy of nature the natural vegetation has its essential place. Hedgerows along country roads and bordering fields provide food, cover, and nesting areas for birds and homes for many small animals. Of some 70 species of shrubs and vines that are typical roadside species in the eastern states alone, about 65 are important to wildlife as food.
\par
\noindent\textbf{译文} 当然，保存路边植被的愿望并不只出于这些审美考虑。在自然的经济体系中，天然植被有其不可或缺的位置。乡间道路和田地边缘的篱丛，为鸟类提供食物、隐蔽处和筑巢地，也为许多小动物提供家园。仅在东部各州，典型路边灌木和藤本约有七十种，其中约六十五种是野生动物的重要食物。
\par\medskip

\noindent\textbf{原文} Such vegetation is also the habitat of wild bees and other pollinating insects. Man is more dependent on these wild pollinators than he usually realizes. Even the farmer himself seldom understands the value of wild bees and often participates in the very measures that rob him of their services. Some agricultural crops and many wild plants are partly or wholly dependent on the services of the native pollinating insects. Several hundred species of wild bees take part in the pollination of cultivated crops---100 species visiting the flowers of alfalfa alone. Without insect pollination, most of the soil-holding and soil-enriching plants of uncultivated areas would die out, with far-reaching consequences to the ecology of the whole region. Many herbs, shrubs, and trees of forests and range depend on native insects for their reproduction; without these plants many wild animals and range stock would find little food. Now clean cultivation and the chemical destruction of hedgerows and weeds are eliminating the last sanctuaries of these pollinating insects and breaking the threads that bind life to life.
\par
\noindent\textbf{译文} 这类植被也是野蜂和其他传粉昆虫的栖息地。人类对这些野生传粉者的依赖，远超通常所知。即使农民本人也很少理解野蜂的价值，并且常常参与那些恰恰剥夺其服务的措施。一些农作物和许多野生植物，部分或完全依赖本地传粉昆虫。数百种野蜂参与栽培作物的授粉；仅访问苜蓿花的就有一百种。没有昆虫授粉，未开垦地区多数固土、肥土植物都会消失，对整个区域的生态造成深远后果。森林和牧地中的许多草本、灌木和乔木，都依赖本地昆虫完成繁殖；没有这些植物，许多野生动物和牧场牲畜就几乎找不到食物。如今，清耕以及对篱丛和杂草的化学摧毁，正在消灭这些传粉昆虫最后的避难所，并切断生命彼此相系的丝线。
\par\medskip

\noindent\textbf{原文} These insects, so essential to our agriculture and indeed to our landscape as we know it, deserve something better from us than the senseless destruction of their habitat. Honeybees and wild bees depend heavily on such ``weeds'' as goldenrod, mustard, and dandelions for pollen that serves as the food of their young. Vetch furnishes essential spring forage for bees before the alfalfa is in bloom, tiding them over this early season so that they are ready to pollinate the alfalfa. In the fall they depend on goldenrod at a season when no other food is available, to stock up for the winter. By the precise and delicate timing that is nature's own, the emergence of one species of wild bees takes place on the very day of the opening of the willow blossoms. There is no dearth of men who understand these things, but these are not the men who order the wholesale drenching of the landscape with chemicals.
\par
\noindent\textbf{译文} 这些昆虫对我们的农业，乃至对我们所熟知的整个景观都至关重要；我们给它们的，不应是对其栖息地的无谓毁灭。蜜蜂和野蜂高度依赖秋麒麟草、芥菜、蒲公英这类“杂草”的花粉，以此喂养幼虫。苜蓿开花前，野豌豆为蜂类提供必不可少的春季饲源，帮助它们度过早春，使其准备好为苜蓿授粉。秋天，在别无食物可得的季节，它们依赖秋麒麟草为越冬储备。按照自然自身精确而细腻的时间安排，某一种野蜂的羽化，正发生在柳花开放的那一天。理解这些事的人并不少；可是下令用化学品大规模浇灌景观的人，并不是他们。
\par\medskip

\noindent\textbf{原文} And where are the men who supposedly understand the value of proper habitat for the preservation of wildlife? Too many of them are to be found defending herbicides as ``harmless'' to wildlife because they are thought to be less toxic than insecticides. Therefore, it is said, no harm is done. But as the herbicides rain down on forest and field, on marsh and rangeland, they are bringing about marked changes and even permanent destruction of wildlife habitat. To destroy the homes and the food of wildlife is perhaps worse in the long run than direct killing.
\par
\noindent\textbf{译文} 那些本应理解适当栖息地对保护野生动物有何价值的人，又在哪里？太多这样的人在为除草剂辩护，声称它们对野生动物“无害”，理由是除草剂被认为不如杀虫剂毒性强。因此，据说并未造成伤害。可是，当除草剂像雨一样落到森林和田野、沼泽和牧地上时，它们正在引起野生动物栖息地的显著改变，甚至造成永久性毁坏。摧毁野生动物的家园和食物，从长远看也许比直接杀戮更糟。
\par\medskip

\noindent\textbf{原文} The irony of this all-out chemical assault on roadsides and utility rights-of-way is twofold. It is perpetuating the problem it seeks to correct, for as experience has clearly shown, the blanket application of herbicides does not permanently control roadside ``brush'' and the spraying has to be repeated year after year. And as a further irony, we persist in doing this despite the fact that a perfectly sound method of selective spraying is known, which can achieve long-term vegetational control and eliminate repeated spraying in most types of vegetation.
\par
\noindent\textbf{译文} 这种对路边和公用设施通道发动的全面化学攻击，具有双重讽刺。它延续了自己试图纠正的问题；经验已经清楚表明，除草剂的全面覆盖施用并不能永久控制路边“灌丛”，喷洒不得不年复一年重复。更讽刺的是，尽管我们明知一种完全稳妥的选择性喷洒方法存在，能够实现长期植被控制，并在大多数植被类型中免除反复喷洒，我们仍继续这样做。
\par\medskip

\noindent\textbf{原文} The object of brush control along roads and rights-of-way is not to sweep the land clear of everything but grass; it is, rather, to eliminate plants ultimately tall enough to present an obstruction to drivers' vision or interference with wires on rights-of-way. This means, in general, trees. Most shrubs are low enough to present no hazard; so, certainly, are ferns and wildflowers.
\par
\noindent\textbf{译文} 道路和通道沿线灌丛控制的目标，并不是把土地扫荡到只剩青草；相反，它是要清除那些最终会长到足以遮挡驾驶者视线，或干扰通道上方线缆的植物。一般说来，这意味着树木。大多数灌木足够低，不构成危险；蕨类和野花当然更是如此。
\par\medskip

\noindent\textbf{原文} Selective spraying was developed by Dr. Frank Egler during a period of years at the American Museum of Natural History as director of a Committee for Brush Control Recommendations for Rights-of-Way. It took advantage of the inherent stability of nature, building on the fact that most communities of shrubs are strongly resistant to invasion by trees. By comparison, grasslands are easily invaded by tree seedlings. The object of selective spraying is not to produce grass on roadsides and rights-of-way but to eliminate the tall woody plants by direct treatment and to preserve all other vegetation. One treatment may be sufficient, with a possible follow-up for extremely resistant species; thereafter the shrubs assert control and the trees do not return. The best and cheapest controls for vegetation are not chemicals but other plants.
\par
\noindent\textbf{译文} 选择性喷洒是弗兰克·埃格勒博士在美国自然历史博物馆任通道灌丛控制建议委员会主任期间，经多年发展出来的方法。它利用自然内在的稳定性，依据的事实是：大多数灌木生物群落强烈抵抗树木入侵。相比之下，草地很容易被树苗侵入。选择性喷洒的目标不是在路边和通道上生产草，而是通过直接处理清除高大的木本植物，并保存所有其他植被。一次处理也许就足够，对特别顽强的物种可作后续处理；此后灌木便取得控制，树木不会返回。控制植被最好、最便宜的办法不是化学品，而是其他植物。
\par\medskip

\noindent\textbf{原文} The method has been tested in research areas scattered throughout the eastern United States. Results show that once properly treated, an area becomes stabilized, requiring no re-spraying for at least 20 years. The spraying can often be done by men on foot, using knapsack sprayers, and having complete control over their material. Sometimes compressor pumps and material can be mounted on truck chassis, but there is no blanket spraying. Treatment is directed only to trees and any exceptionally tall shrubs that must be eliminated. The integrity of the environment is thereby preserved, the enormous value of the wildlife habitat remains intact, and the beauty of shrub and fern and wildflower has not been sacrificed.
\par
\noindent\textbf{译文} 这种方法已在美国东部各地分散的研究区域中进行测试。结果表明，一片区域一旦经过适当处理，就会稳定下来，至少二十年不需要重新喷洒。喷洒常可由步行人员使用背负式喷雾器完成，从而完全控制药剂。有时也可把压缩泵和药剂装在卡车底盘上，但绝不进行覆盖式喷洒。处理只针对必须清除的树木和任何特别高大的灌木。环境完整性因而得以保存，野生动物栖息地的巨大价值保持无损，灌木、蕨类和野花之美也没有被牺牲。
\par\medskip

\noindent\textbf{原文} Here and there the method of vegetation management by selective spraying has been adopted. For the most part, entrenched custom dies hard and blanket spraying continues to thrive, to exact its heavy annual costs from the taxpayer, and to inflict its damage on the ecological web of life. It thrives, surely, only because the facts are not known. When taxpayers understand that the bill for spraying the town roads should come due only once a generation instead of once a year, they will surely rise up and demand a change of method.
\par
\noindent\textbf{译文} 在一些地方，选择性喷洒这种植被管理方法已经被采用。但在大多数地方，根深蒂固的习惯难以消亡，全面喷洒仍然兴盛，年复一年向纳税人索取沉重成本，并把损害加诸生命的生态之网。它之所以兴盛，当然只是因为事实还不为人知。当纳税人明白，城镇道路喷洒的账单本应一代人才到期一次，而不是一年一次，他们必定会起来要求改变方法。
\par\medskip

\noindent\textbf{原文} Among the many advantages of selective spraying is the fact that it minimizes the amount of chemical applied to the landscape. There is no broadcasting of material but, rather, concentrated application to the base of the trees. The potential harm to wildlife is therefore kept to a minimum.
\par
\noindent\textbf{译文} 选择性喷洒有许多优点，其中之一是它把施加到景观上的化学品数量降至最低。药剂不是被广泛撒布，而是集中施用于树基。因此，对野生动物的潜在伤害也被控制在最低限度。
\par\medskip

\noindent\textbf{原文} The most widely used herbicides are 2,4-D, 2,4,5-T, and related compounds. Whether or not these are actually toxic is a matter of controversy. People spraying their lawns with 2,4-D and becoming wet with spray have occasionally developed severe neuritis and even paralysis. Although such incidents are apparently uncommon, medical authorities advise caution in use of such compounds. Other hazards, more obscure, may also attend the use of 2,4-D. It has been shown experimentally to disturb the basic physiological process of respiration in the cell, and to imitate X-rays in damaging the chromosomes. Some very recent work indicates that reproduction of birds may be adversely affected by these and certain other herbicides at levels far below those that cause death.
\par
\noindent\textbf{译文} 使用最广泛的除草剂是2,4-D、2,4,5-T以及相关化合物。它们是否真正有毒，仍有争议。有些人在用2,4-D喷洒自家草坪、并被喷雾淋湿后，偶尔发生严重神经炎，甚至瘫痪。尽管这类事件显然并不常见，医学权威仍建议谨慎使用这些化合物。使用2,4-D还可能伴随其他更隐蔽的危险。实验显示，它会干扰细胞呼吸这一基本生理过程，并会像X射线一样损伤染色体。最近的一些研究表明，在远低于致死剂量的水平上，这些和某些其他除草剂可能已经对鸟类繁殖造成不利影响。
\par\medskip

\noindent\textbf{原文} Apart from any directly toxic effects, curious indirect results follow the use of certain herbicides. It has been found that animals, both wild herbivores and livestock, are sometimes strangely attracted to a plant that has been sprayed, even though it is not one of their natural foods. If a highly poisonous herbicide such as arsenic has been used, this intense desire to reach the wilting vegetation inevitably has disastrous results. Fatal results may follow, also, from less toxic herbicides if the plant itself happens to be poisonous or perhaps to possess thorns or burs. Poisonous range weeds, for example, have suddenly become attractive to livestock after spraying, and the animals have died from indulging this unnatural appetite. The literature of veterinary medicine abounds in similar examples: swine eating sprayed cockleburs with consequent severe illness, lambs eating sprayed thistles, bees poisoned by pasturing on mustard sprayed after it came into bloom. Wild cherry, the leaves of which are highly poisonous, has exerted a fatal attraction for cattle once its foliage has been sprayed with 2,4-D. Apparently the wilting that follows spraying (or cutting) makes the plant attractive. Ragwort has provided other examples. Livestock ordinarily avoid this plant unless forced to turn to it in late winter and early spring by lack of other forage. However, the animals eagerly feed on it after its foliage has been sprayed with 2,4-D.
\par
\noindent\textbf{译文} 除了任何直接毒性影响之外，某些除草剂的使用还会带来奇特的间接结果。人们发现，无论野生食草动物还是家畜，有时都会奇怪地被喷洒过的植物吸引，哪怕那并不是它们的天然食物。如果使用的是砷这类剧毒除草剂，这种接近萎蔫植被的强烈欲望，必然造成灾难性后果。即使除草剂毒性较低，只要植物本身恰好有毒，或者也许带刺或带芒刺，也可能导致死亡。例如，有毒的牧场杂草在喷洒后会突然变得吸引牲畜，动物因纵容这种反常食欲而死去。兽医学文献中充满类似例子：猪吃了喷洒过的苍耳而重病；羔羊吃了喷洒过的蓟；蜜蜂在芥菜开花后取食被喷洒的芥菜而中毒。野樱桃的叶子毒性很强；一旦其叶片被2,4-D喷洒，就会对牛产生致命吸引。显然，喷洒或切割之后发生的萎蔫，使植物变得有吸引力。千里光也提供了其他例子。牲畜通常避开这种植物，除非在晚冬和早春因缺乏其他饲草而被迫转向它。然而，它的叶片一经2,4-D喷洒，动物便会急切取食。
\par\medskip

\noindent\textbf{原文} The explanation of this peculiar behavior sometimes appears to lie in the changes which the chemical brings about in the metabolism of the plant itself. There is temporarily a marked increase in sugar content, making the plant more attractive to many animals.
\par
\noindent\textbf{译文} 对这种奇特行为的解释，有时似乎在于化学品对植物自身代谢造成的变化。植物含糖量会暂时显著增加，从而对许多动物更具吸引力。
\par\medskip

\noindent\textbf{原文} Another curious effect of 2,4-D has important effects for livestock, wildlife, and apparently for men as well. Experiments carried out about a decade ago showed that after treatment with this chemical there is a sharp increase in the nitrate content of corn and of sugar beets. The same effect was suspected in sorghum, sunflower, spiderwort, lamb's quarters, pigweed, and smartweed. Some of these are normally ignored by cattle, but are eaten with relish after treatment with 2,4-D. A number of deaths among cattle have been traced to sprayed weeds, according to some agricultural specialists. The danger lies in the increase in nitrates, for the peculiar physiology of the ruminant at once poses a critical problem. Most such animals have a digestive system of extraordinary complexity, including a stomach divided into four chambers. The digestion of cellulose is accomplished through the action of microorganisms (rumen bacteria) in one of the chambers. When the animal feeds on vegetation containing an abnormally high level of nitrates, the microorganisms in the rumen act on the nitrates to change them into highly toxic nitrites. Thereafter a fatal chain of events ensues: the nitrites act on the blood pigment to form a chocolate-brown substance in which the oxygen is so firmly held that it cannot take part in respiration, hence oxygen is not transferred from the lungs to the tissues. Death occurs within a few hours from anoxia, or lack of oxygen. The various reports of livestock losses after grazing on certain weeds treated with 2,4-D therefore have a logical explanation. The same danger exists for wild animals belonging to the group of ruminants, such as deer, antelope, sheep, and goats.
\par
\noindent\textbf{译文} 2,4-D的另一种奇特影响，对牲畜、野生动物，显然也对人具有重要后果。约十年前进行的实验显示，经这种化学品处理后，玉米和甜菜中的硝酸盐含量急剧上升。同样的影响也被怀疑出现在高粱、向日葵、紫露草、藜、苋和蓼等植物中。其中一些植物通常被牛忽视，但经2,4-D处理后却会被津津有味地吃掉。一些农业专家认为，已有若干牛死亡事件可追溯到喷洒过的杂草。危险在于硝酸盐增加，因为反刍动物独特的生理立即带来关键问题。多数这类动物拥有极其复杂的消化系统，包括分成四个室的胃。纤维素消化是在其中一个胃室里通过微生物（瘤胃细菌）的作用完成的。当动物取食含有异常高水平硝酸盐的植被时，瘤胃中的微生物作用于硝酸盐，将其转化为剧毒的亚硝酸盐。随后，一连串致命事件发生：亚硝酸盐作用于血液色素，形成一种巧克力褐色物质，氧被牢牢结合其中，以致不能参与呼吸；因此，氧不能从肺转移到组织。数小时内，动物便因缺氧而死亡。因此，关于牲畜在取食经2,4-D处理的某些杂草后死亡的各种报告，有了合乎逻辑的解释。属于反刍类的野生动物，如鹿、叉角羚、绵羊和山羊，也面临同样危险。
\par\medskip

\noindent\textbf{原文} Although various factors (such as exceptionally dry weather) can cause an increase in nitrate content, the effect of the soaring sales and applications of 2,4-D cannot be ignored. The situation was considered important enough by the University of Wisconsin Agricultural Experiment Station to justify a warning in 1957 that ``plants killed by 2,4-D may contain large amounts of nitrate.'' The hazard extends to human beings as well as animals and may help to explain the recent mysterious increase in ``silo deaths.'' When corn, oats, or sorghum containing large amounts of nitrates are ensiled they release poisonous nitrogen oxide gases, creating a deadly hazard to anyone entering the silo. Only a few breaths of one of these gases can cause a diffuse chemical pneumonia. In a series of such cases studied by the University of Minnesota Medical School all but one terminated fatally.
\par
\noindent\textbf{译文} 虽然多种因素（例如异常干旱天气）都能导致硝酸盐含量上升，但2,4-D销量和施用量飙升的影响不能被忽视。威斯康星大学农业试验站认为情况足够重要，因而在1957年发出警告：“被2,4-D杀死的植物可能含有大量硝酸盐。”这种危险不仅延伸到动物，也延伸到人类，并可能有助于解释近来“青贮窖死亡”事件神秘增加的原因。当含有大量硝酸盐的玉米、燕麦或高粱被青贮时，会释放有毒的氮氧化物气体，对任何进入青贮窖的人构成致命危险。只要吸入几口其中一种气体，就可能造成弥漫性化学性肺炎。明尼苏达大学医学院研究的一系列此类病例中，除一例外全部致死。
\par\medskip

\noindent\textbf{原文} ``Once again we are walking in nature like an elephant in the china cabinet.'' So C. J. Briejer, a Dutch scientist of rare understanding, sums up our use of weed killers. ``In my opinion too much is taken for granted. We do not know whether all weeds in crops are harmful or whether some of them are useful,'' says Dr. Briejer.
\par
\noindent\textbf{译文} “我们又一次像大象走进瓷器柜那样行走在自然中。”荷兰科学家C·J·布里耶尔以罕见的理解力，这样总结我们对除草剂的使用。布里耶尔博士说：“在我看来，人们把太多事情视为理所当然。我们并不知道作物中的所有杂草是否都有害，也不知道其中有些是否有益。”
\par\medskip

\noindent\textbf{原文} Seldom is the question asked: What is the relation between the weed and the soil? Perhaps, even from our narrow standpoint of direct self-interest, the relation is a useful one. As we have seen, soil and the living things in and upon it exist in a relation of interdependence and mutual benefit. Presumably the weed is taking something from the soil; perhaps it is also contributing something to it. A practical example was provided recently by the parks in a city in Holland. The roses were doing badly. Soil samples showed heavy infestations by tiny nematode worms. Scientists of the Dutch Plant Protection Service did not recommend chemical sprays or soil treatments; instead, they suggested that marigolds be planted among the roses. This plant, which the purist would doubtless consider a weed in any rosebed, releases an excretion from its roots that kills the soil nematodes. The advice was taken; some beds were planted with marigolds, some left without as controls. The results were striking. With the aid of the marigolds the roses flourished; in the control beds they were sickly and drooping. Marigolds are now used in many places for combating nematodes.
\par
\noindent\textbf{译文} 人们很少提出这样的问题：杂草与土壤之间有什么关系？也许，即使从我们狭隘的直接自利角度看，这种关系也是有用的。正如我们已经看到的，土壤以及生活在其中、其上的生命，存在相互依赖、互利共生的关系。可以推想，杂草从土壤中取走某些东西；也许它也向土壤贡献某些东西。荷兰一座城市的公园最近提供了一个实际例子。玫瑰长势很差。土样显示微小线虫严重侵染。荷兰植物保护局的科学家没有建议化学喷洒或土壤处理；相反，他们建议在玫瑰之间种植万寿菊。这种植物，洁癖式的园艺家无疑会认为它在任何玫瑰花坛里都是杂草，却会从根部分泌一种物质，杀死土壤线虫。建议被采纳了；一些花坛种上万寿菊，另一些不种，作为对照。结果非常显著。在万寿菊帮助下，玫瑰繁茂生长；对照花坛中的玫瑰则病弱下垂。如今，万寿菊已在许多地方用于防治线虫。
\par\medskip

\noindent\textbf{原文} In the same way, and perhaps quite unknown to us, other plants that we ruthlessly eradicate may be performing a function that is necessary to the health of the soil. One very useful function of natural plant communities---now pretty generally stigmatized as ``weeds''---is to serve as an indicator of the condition of the soil. This useful function is of course lost where chemical weed killers have been used.
\par
\noindent\textbf{译文} 同样地，也许我们全然不知，那些被我们无情根除的其他植物，可能正在履行一种维持土壤健康所必需的功能。天然植物生物群落如今已普遍被污名化为“杂草”，但它们有一项非常有用的功能：指示土壤状况。当然，在使用化学除草剂的地方，这种有用功能便丧失了。
\par\medskip

\noindent\textbf{原文} Those who find an answer to all problems in spraying also overlook a matter of great scientific importance---the need to preserve some natural plant communities. We need these as a standard against which we can measure the changes our own activities bring about. We need them as wild habitats in which original populations of insects and other organisms can be maintained, for, as will be explained in Chapter 16, the development of resistance to insecticides is changing the genetic factors of insects and perhaps other organisms. One scientist has even suggested that some sort of ``zoo'' should be established to preserve insects, mites, and the like, before their genetic composition is further changed.
\par
\noindent\textbf{译文} 那些认为喷洒能回答一切问题的人，也忽视了一件具有重大科学意义的事：需要保存某些天然植物生物群落。我们需要它们作为标准，用来衡量我们自身活动带来的变化。我们还需要它们作为野生栖息地，使昆虫和其他生物的原初种群得以维持；因为正如第十六章将说明的，对杀虫剂抗性的产生，正在改变昆虫，也许还有其他生物的遗传因素。一位科学家甚至建议，在昆虫、螨类等的遗传组成进一步改变之前，应建立某种“动物园”来保存它们。
\par\medskip

\noindent\textbf{原文} Some experts warn of subtle but far-reaching vegetational shifts as a result of the growing use of herbicides. The chemical 2,4-D, by killing out the broad-leaved plants, allows the grasses to thrive in the reduced competition---now some of the grasses themselves have become ``weeds,'' presenting a new problem in control and giving the cycle another turn. This strange situation is acknowledged in a recent issue of a journal devoted to crop problems: ``With the widespread use of 2,4-D to control broadleaved weeds, grass weeds in particular have increasingly become a threat to corn and soybean yields.''
\par
\noindent\textbf{译文} 一些专家警告，除草剂使用日增，可能导致微妙却影响深远的植被转变。2,4-D这种化学品通过杀灭阔叶植物，减少竞争，使禾草繁盛；如今，一些禾草本身又变成了“杂草”，提出新的控制问题，使循环又转了一圈。最近一期专门讨论作物问题的期刊承认了这种奇怪状况：“随着2,4-D被广泛用于控制阔叶杂草，禾本科杂草尤其日益成为对玉米和大豆产量的威胁。”
\par\medskip

\noindent\textbf{原文} Ragweed, the bane of hay fever sufferers, offers an interesting example of the way efforts to control nature sometimes boomerang. Many thousands of gallons of chemicals have been discharged along roadsides in the name of ragweed control. But the unfortunate truth is that blanket spraying is resulting in more ragweed, not less. Ragweed is an annual; its seedlings require open soil to become established each year. Our best protection against this plant is therefore the maintenance of dense shrubs, ferns, and other perennial vegetation. Spraying frequently destroys this protective vegetation and creates open, barren areas which the ragweed hastens to fill. It is probable, moreover, that the pollen content of the atmosphere is not related to roadside ragweed, but to the ragweed of city lots and fallow fields.
\par
\noindent\textbf{译文} 豚草是花粉热患者的祸害，它提供了一个有趣例子，说明控制自然的努力有时会反噬。以控制豚草为名，成千上万加仑化学品被排放到道路两旁。但不幸的事实是，覆盖式喷洒带来的豚草不是更少，而是更多。豚草是一年生植物；它的幼苗每年都需要裸露土壤才能定居。因此，我们抵御这种植物的最好保护，是维持浓密灌木、蕨类和其他多年生植被。喷洒往往破坏这种保护性植被，制造开阔贫瘠的区域，而豚草会迅速填补。况且，大气中的花粉含量很可能并不与路边豚草有关，而是与城市空地和休耕地上的豚草有关。
\par\medskip

\noindent\textbf{原文} The booming sales of chemical crabgrass killers are another example of how readily unsound methods catch on. There is a cheaper and better way to remove crabgrass than to attempt year after year to kill it out with chemicals. This is to give it competition of a kind it cannot survive, the competition of other grass. Crabgrass exists only in an unhealthy lawn. It is a symptom, not a disease in itself. By providing a fertile soil and giving the desired grasses a good start, it is possible to create an environment in which crabgrass cannot grow, for it requires open space in which it can start from seed year after year.
\par
\noindent\textbf{译文} 化学马唐除草剂销量激增，是不健全方法多么容易流行的另一个例子。要清除马唐，有一种比年复一年试图用化学品杀死它更便宜、更好的办法：给予它无法承受的竞争，也就是其他草类的竞争。马唐只存在于不健康的草坪中。它是一种症状，本身并不是疾病。只要提供肥沃土壤，并让所需草类有良好开端，就可以创造一种马唐不能生长的环境，因为它需要开阔空间，才能年复一年从种子开始生长。
\par\medskip

\noindent\textbf{原文} Instead of treating the basic condition, suburbanites---advised by nurserymen who in turn have been advised by the chemical manufacturers---continue to apply truly astonishing amounts of crabgrass killers to their lawns each year. Marketed under trade names which give no hint of their nature, many of these preparations contain such poisons as mercury, arsenic, and chlordane. Application at the recommended rates leaves tremendous amounts of these chemicals on the lawn. Users of one product, for example, apply 60 pounds of technical chlordane to the acre if they follow directions. If they use another of the many available products, they are applying 175 pounds of metallic arsenic to the acre. The toll of dead birds, as we shall see in Chapter 8, is distressing. How lethal these lawns may be for human beings is unknown.
\par
\noindent\textbf{译文} 郊区居民没有处理根本状况，而是在苗圃业者建议下继续每年向草坪施用数量惊人的马唐除草剂；而这些苗圃业者又接受了化学品制造商的建议。许多这类制剂以商标名称销售，名称丝毫不透露其本质，却含有汞、砷、氯丹等毒物。按推荐剂量施用，会在草坪上留下大量化学品。例如，某一产品的用户若按说明操作，每英亩便施用了六十磅工业氯丹。若使用另一些可选产品中的一种，他们每英亩施用的则是一百七十五磅金属砷。正如我们将在第八章看到的，死亡鸟类的代价令人痛心。这些草坪对人类可能有多致命，目前尚无人知晓。
\par\medskip

\noindent\textbf{原文} The success of selective spraying for roadside and right-of-way vegetation, where it has been practiced, offers hope that equally sound ecological methods may be developed for other vegetation programs for farms, forests, and ranges---methods aimed not at destroying a particular species but at managing vegetation as a living community.
\par
\noindent\textbf{译文} 在已经实践的地方，路边和通道植被选择性喷洒的成功，为农场、森林和牧地的其他植被计划提供了希望：同样健全的生态方法也许能够发展出来。这些方法的目标不是毁灭某个特定物种，而是把植被作为一个有生命的生物群落来管理。
\par\medskip

\noindent\textbf{原文} Other solid achievements show what can be done. Biological control has achieved some of its most spectacular successes in the area of curbing unwanted vegetation. Nature herself has met many of the problems that now beset us, and she has usually solved them in her own successful way. Where man has been intelligent enough to observe and to emulate Nature he, too, is often rewarded with success.
\par
\noindent\textbf{译文} 其他扎实成就显示了可以做到什么。生物防治在遏制不受欢迎植被的领域，取得了一些最引人注目的成功。自然本身已经遭遇过许多如今困扰我们的问题，并且通常以自己的成功方式解决了它们。当人类足够聪明，去观察并效法自然时，他也常常得到成功的回报。
\par\medskip

\noindent\textbf{原文} An outstanding example in the field of controlling unwanted plants is the handling of the Klamath-weed problem in California. Although the Klamath weed, or goatweed, is a native of Europe (where it is called St. Johnswort), it accompanied man in his westward migrations, first appearing in the United States in 1793 near Lancaster, Pennsylvania. By 1900 it had reached California in the vicinity of the Klamath River, hence the name locally given to it. By 1929 it had occupied about 100,000 acres of rangeland, and by 1952 it had invaded some two and one half million acres.
\par
\noindent\textbf{译文} 在控制不受欢迎植物的领域，加利福尼亚处理克拉马斯草问题，是一个杰出例子。克拉马斯草，又称山羊草，原产欧洲（在那里称为圣约翰草或贯叶连翘），却伴随人类西迁，1793年首先出现在美国宾夕法尼亚州兰开斯特附近。到1900年，它已抵达加利福尼亚克拉马斯河一带，因此获得当地名称。到1929年，它占据了约十万英亩牧地；到1952年，入侵面积已达约二百五十万英亩。
\par\medskip

\noindent\textbf{原文} Klamath weed, quite unlike such native plants as sagebrush, has no place in the ecology of the region, and no animals or other plants require its presence. On the contrary, wherever it appeared livestock became ``scabby, sore-mouthed, and unthrifty'' from feeding on this toxic plant. Land values declined accordingly, for the Klamath weed was considered to hold the first mortgage.
\par
\noindent\textbf{译文} 克拉马斯草与蒿灌木这类本地植物截然不同，在该地区生态中没有位置，也没有动物或其他植物需要它存在。相反，它出现在哪里，牲畜因取食这种有毒植物就会变得“结痂、口腔溃痛、发育不良”。土地价值也相应下降，因为克拉马斯草被认为握有第一抵押权。
\par\medskip

\noindent\textbf{原文} In Europe the Klamath weed, or St. Johnswort, has never become a problem because along with the plant there have developed various species of insects; these feed on it so extensively that its abundance is severely limited. In particular, two species of beetles in southern France, pea-sized and of metallic color, have their whole beings so adapted to the presence of the weed that they feed and reproduce only upon it.
\par
\noindent\textbf{译文} 在欧洲，克拉马斯草，也就是圣约翰草，从未成为问题，因为与这种植物一起演化出了多种昆虫；这些昆虫大量取食它，使其丰度受到严格限制。尤其是在法国南部，有两种豌豆大小、带金属色泽的甲虫，其整个生命都如此适应这种杂草的存在，以至于只在它身上取食和繁殖。
\par\medskip

\noindent\textbf{原文} It was an event of historic importance when the first shipments of these beetles were brought to the United States in 1944, for this was the first attempt in North America to control a plant with a plant-eating insect. By 1948 both species had become so well established that no further importations were needed. Their spread was accomplished by collecting beetles from the original colonies and redistributing them at the rate of millions a year. Within small areas the beetles accomplish their own dispersion, moving on as soon as the Klamath weed dies out and locating new stands with great precision. And as the beetles thin out the weed, desirable range plants that have been crowded out are able to return.
\par
\noindent\textbf{译文} 1944年，第一批这类甲虫被运到美国，这是一个具有历史意义的事件，因为这是北美首次尝试用食植昆虫控制一种植物。到1948年，两种甲虫都已建立稳固种群，不再需要进一步进口。它们的扩散，是通过从最初群落中采集甲虫，并以每年数百万只的速度重新分布来完成的。在小范围内，甲虫会自行扩散；克拉马斯草一旦死去，它们便迁往别处，并以极高精度找到新的草丛。随着甲虫削弱这种杂草，那些被挤出的有益牧地植物便能够返回。
\par\medskip

\noindent\textbf{原文} A ten-year survey completed in 1959 showed that control of the Klamath weed had been ``more effective than hoped for even by enthusiasts,'' with the weed reduced to a mere 1 per cent of its former abundance. This token infestation is harmless and is actually needed in order to maintain a population of beetles as protection against a future increase in the weed.
\par
\noindent\textbf{译文} 1959年完成的一项十年调查显示，对克拉马斯草的控制“比热心者所希望的还要有效”，这种杂草已减少到原来丰度的仅百分之一。这种象征性的侵染并无害处，实际上还为维持甲虫种群所必需，以防该杂草未来再次增加。
\par\medskip

\noindent\textbf{原文} Another extraordinarily successful and economical example of weed control may be found in Australia. With the colonists' usual taste for carrying plants or animals into a new country, a Captain Arthur Phillip had brought various species of cactus into Australia about 1787, intending to use them in culturing cochineal insects for dye. Some of the cacti or prickly pears escaped from his gardens and by 1925 about 20 species could be found growing wild. Having no natural controls in this new territory, they spread prodigiously, eventually occupying about 60 million acres. At least half of this land was so densely covered as to be useless.
\par
\noindent\textbf{译文} 另一个极其成功且经济的杂草控制例子，可在澳大利亚找到。殖民者惯常喜欢把植物或动物带入新国土；约在1787年，亚瑟·菲利普船长把多种仙人掌带到澳大利亚，打算用它们培养胭脂虫制取染料。其中一些仙人掌或梨果仙人掌从他的花园逃逸；到1925年，约二十个种已在野外生长。在这片新领地缺乏自然控制因素，它们便异常扩散，最终占据约六千万英亩土地。其中至少一半覆盖过密，已无法利用。
\par\medskip

\noindent\textbf{原文} In 1920 Australian entomologists were sent to North and South America to study insect enemies of the prickly pears in their native habitat. After trials of several species, 3 billion eggs of an Argentine moth were released in Australia in 1930. Seven years later the last dense growth of the prickly pear had been destroyed and the once uninhabitable areas reopened to settlement and grazing. The whole operation had cost less than a penny per acre. In contrast, the unsatisfactory attempts at chemical control in earlier years had cost about \pounds 10 per acre.
\par
\noindent\textbf{译文} 1920年，澳大利亚昆虫学家被派往南、北美洲，研究梨果仙人掌在原生栖息地中的昆虫天敌。试验若干物种之后，1930年，三十亿枚阿根廷蛾卵在澳大利亚释放。七年后，最后一片茂密的梨果仙人掌丛被摧毁，曾经不宜居住的地区重新向定居和放牧开放。整个行动每英亩花费不到一便士。相比之下，早年不令人满意的化学控制尝试，每英亩花费约十英镑。
\par\medskip

\noindent\textbf{原文} Both of these examples suggest that extremely effective control of many kinds of unwanted vegetation might be achieved by paying more attention to the role of plant-eating insects. The science of range management has largely ignored this possibility, although these insects are perhaps the most selective of all grazers and their highly restricted diets could easily be turned to man's advantage.
\par
\noindent\textbf{译文} 这两个例子都表明，只要更加重视食植昆虫的作用，许多种不受欢迎植被或许都能得到极其有效的控制。牧地管理科学在很大程度上忽视了这种可能性，尽管这些昆虫也许是所有取食者中选择性最高的，其高度受限的食谱很容易转化为有利于人类的工具。
\par


% ===== 07_poincare_kandel.tex =====
\part*{Henri Poincaré, \emph{Science and Method}\\亨利·庞加莱《科学与方法》}

\section*{I. The Selection of Facts\\一、事实的选择}

\begin{enblock}
Tolstoi explains somewhere in his writings why, in his opinion, “Science for Science's sake” is an absurd conception. We cannot know all the facts, since they are practically infinite in number. We must make a selection; and that being so, can this selection be governed by the mere caprice of our curiosity? Is it not better to be guided by utility, by our practical, and more especially our moral, necessities? Have we not some better occupation than counting the number of ladybirds in existence on this planet?
\end{enblock}
\begin{zhblock}
托尔斯泰在某处著作中解释过，为什么在他看来，“为科学而科学”是一种荒谬的观念。我们不可能认识所有事实，因为事实的数量实际上是无限的。我们必须作出选择；既然如此，这种选择能只由好奇心的一时任性来支配吗？由功用，由我们的实践需要，尤其由我们的道德需要来引导，不是更好吗？难道我们没有比清点地球上有多少只瓢虫更值得做的事吗？\end{zhblock}

\begin{enblock}
It is clear that for him the word utility has not the meaning assigned to it by business men, and, after them, by the greater number of our contemporaries. He cares but little for the industrial applications of science, for the marvels of electricity or of automobilism, which he regards rather as hindrances to moral progress. For him the useful is exclusively what is capable of making men better.
\end{enblock}
\begin{zhblock}
显然，对他来说，“功用”一词并不是商人以及随后多数当代人赋予它的含义。他几乎不关心科学的工业应用，不关心电力或汽车技术的奇迹；在他看来，这些倒更像是道德进步的障碍。对他而言，有用的东西只限于能够使人变得更好的东西。\end{zhblock}

\begin{enblock}
It is hardly necessary for me to state that, for my part, I could not be satisfied with either of these ideals. I have no liking either for a greedy and narrow plutocracy, or for a virtuous unaspiring democracy, solely occupied in turning the other cheek, in which we should find good people devoid of curiosity, who, avoiding all excesses, would not die of any disease--save boredom. But it is all a matter of taste, and that is not the point I wish to discuss.
\end{enblock}
\begin{zhblock}
几乎无须说明，就我而言，我不能满足于这两种理想中的任何一种。我既不喜欢贪婪而狭隘的富豪政治，也不喜欢一种有德却无所追求的民主社会，在那里，人们只忙着“把另一边脸也转过去”；那里有善良却没有好奇心的人，他们避开一切过度，除了无聊之外不会死于任何疾病。不过这全是趣味问题，并不是我想讨论的要点。\end{zhblock}

\begin{enblock}
Nonetheless the question remains, and it claims our attention. If our selection is only determined by caprice or by immediate necessity, there can be no science for science's sake, and consequently no science. Is this true? There is no disputing the fact that a selection must be made: however great our activity, facts outstrip us, and we can never overtake them; while the scientist is discovering one fact, millions and millions are produced in every cubic inch of his body. Trying to make science contain nature is like trying to make the part contain the whole.
\end{enblock}
\begin{zhblock}
然而问题依然存在，而且需要我们认真对待。如果我们的选择只由任性或眼前需要决定，那么就不会有“为科学而科学”，因而也不会有科学。事实真是如此吗？选择必须作出，这是无可争辩的：无论我们多么勤奋，事实总是跑在我们前面，我们永远追不上它们；科学家发现一个事实的同时，他身体每一立方英寸里又产生了成千上万的事实。想让科学容纳自然，就像想让部分容纳整体。\end{zhblock}

\begin{enblock}
But scientists believe that there is a hierarchy of facts, and that a judicious selection can be made. They are right, for otherwise there would be no science, and science does exist. One has only to open one's eyes to see that the triumphs of industry, which have enriched so many practical men, would never have seen the light if only these practical men had existed, and if they had not been preceded by disinterested fools who died poor, who never thought of the useful, and yet had a guide that was not their own caprice.
\end{enblock}
\begin{zhblock}
但是科学家相信，事实之间有一种层级，并且可以作出明智的选择。他们是对的；否则科学就不会存在，而科学确实存在。只要睁开眼睛就能看见，那些使许多实际人物致富的工业胜利，如果世界上只有这些实际人物，如果在他们之前没有那些无私的傻子，便绝不会出现。这些人穷困而死，从未想着“有用”，但他们仍有某种指引，而那并不是他们自己的任性。\end{zhblock}

\begin{enblock}
What these fools did, as Mach has said, was to save their successors the trouble of thinking. If they had worked solely in view of an immediate application, they would have left nothing behind them, and in face of a new requirement, all would have had to be done again. Now the majority of men do not like thinking, and this is perhaps a good thing, since instinct guides them, and very often better than reason would guide a pure intelligence, at least whenever they are pursuing an end that is immediate and always the same. But instinct is routine, and if it were not fertilized by thought, it would advance no further with man than with the bee or the ant. It is necessary, therefore, to think for those who do not like thinking, and as they are many, each one of our thoughts must be useful in as many circumstances as possible. For this reason, the more general a law is, the greater is its value.
\end{enblock}
\begin{zhblock}
正如马赫所说，这些傻子所做的，是替后来者省去了思考的麻烦。如果他们只为了眼前的应用而工作，他们就不会留下任何东西；一旦出现新的需求，一切都得从头再来。多数人并不喜欢思考，这也许还是好事，因为本能会引导他们；至少当他们追求的是眼前而且总是相同的目标时，本能常常比纯粹理智更能引导他们。但是本能只是惯例；如果没有思想使它受孕，它在人那里也不会比在蜜蜂或蚂蚁那里走得更远。因此，必须为那些不喜欢思考的人思考；而这样的人很多，所以我们的每一个思想都应当在尽可能多的情形中有用。正因为如此，一条定律越普遍，它的价值就越大。\end{zhblock}

\begin{enblock}
This shows us how our selection should be made. The most interesting facts are those which can be used several times, those which have a chance of recurring. We have been fortunate enough to be born in a world where there are such facts. Suppose that instead of eighty chemical elements we had eighty million, and that they were not some common and others rare, but uniformly distributed. Then each time we picked up a new pebble there would be a strong probability that it was composed of some unknown substance. Nothing that we knew of other pebbles would tell us anything about it. Before each new object we should be like a newborn child; like him we could but obey our caprices or our necessities. In such a world there would be no science; perhaps thought and even life would be impossible, since evolution could not have developed the instincts of self-preservation. Providentially it is not so; but this blessing, like all those to which we are accustomed, is not appreciated at its true value. The biologist would be equally embarrassed if there were only individuals and no species, and if heredity did not make children resemble their parents.
\end{enblock}
\begin{zhblock}
这就向我们说明，应当怎样作出选择。最有意思的事实，是那些能够被多次使用的事实，是那些有可能重复出现的事实。我们足够幸运，出生在一个确有这类事实的世界里。设想一下，如果化学元素不是八十种，而是八千万种；而且它们并非有的常见、有的稀少，而是均匀分布。那么我们每捡起一块新的卵石，它很可能都由某种未知物质组成。我们关于其他卵石所知道的一切，都不会告诉我们关于它的任何事情。面对每一个新对象，我们都会像新生儿一样；像他一样，我们只能听从自己的任性或需要。在这样的世界里就不会有科学；也许思想甚至生命都不可能存在，因为进化无法发展出自我保存的本能。幸而事实并非如此；但这种恩惠也和一切我们习以为常的恩惠一样，并没有得到应有的珍视。如果只有个体而没有物种，如果遗传不能使子女与父母相似，生物学家也会同样无所适从。\end{zhblock}

\begin{enblock}
Which, then, are the facts that have a chance of recurring? In the first place, simple facts. It is evident that in a complex fact many circumstances are united by chance, and that only a still more improbable chance could ever so unite them again. But are there such things as simple facts? and if there are, how are we to recognize them? Who can tell that what we believe to be simple does not conceal an alarming complexity? All that we can say is that we must prefer facts which appear simple, to those in which our rude vision detects dissimilar elements. Then only two alternatives are possible; either this simplicity is real, or else the elements are so intimately mingled that they do not admit of being distinguished. In the first case we have a chance of meeting the same simple fact again, either in all its purity, or itself entering as an element into some complex whole. In the second case the intimate mixture has similarly a greater chance of being reproduced than a heterogeneous assemblage. Chance can mingle, but it cannot unmingle, and a combination of various elements in a well-ordered edifice in which something can be distinguished, can only be made deliberately. There is, therefore, but little chance that an assemblage in which different things can be distinguished should ever be reproduced. On the other hand, there is great probability that a mixture which appears homogeneous at first sight will be reproduced several times. Accordingly facts which appear simple, even if they are not so in reality, will be more easily brought about again by chance.
\end{enblock}
\begin{zhblock}
那么，哪些事实有可能重复出现呢？首先是简单事实。显然，在一个复杂事实中，许多情形是由偶然性结合在一起的；只有一种更不可能的偶然性，才会再次以同样方式把它们结合起来。但是，真有简单事实这种东西吗？如果有，我们又如何识别它们？谁能断言，我们以为简单的东西没有隐藏着惊人的复杂性？我们所能说的只是：我们必须偏爱那些看起来简单的事实，而不是那些在我们粗糙的眼光中显出异质成分的事实。于是只有两种可能：这种简单性是真实的；或者，各个成分混合得如此紧密，以致无法加以区分。在第一种情形中，我们有机会再次遇到同一个简单事实，或者以其纯粹形态出现，或者作为一个成分进入某个复杂整体。在第二种情形中，这种紧密混合也比异质拼合更有可能重现。偶然性能混合，却不能解混；由多种元素构成、秩序井然而且其中某些东西可被区分的组合，只能有意造成。因此，一个其中可以区分出不同事物的集合几乎没有机会被再次重现。相反，一个乍看之下同质的混合物，多次重现的概率却很大。所以，那些看起来简单的事实，即使实际上并不简单，也更容易被偶然性再次带来。\end{zhblock}

\begin{enblock}
It is this that justifies the method instinctively adopted by scientists, and what perhaps justifies it still better is that facts which occur frequently appear to us simple just because we are accustomed to them.
\end{enblock}
\begin{zhblock}
正是这一点为科学家本能采用的方法提供了根据；也许更能为它辩护的是，频繁发生的事实之所以在我们看来简单，正是因为我们已经习惯了它们。\end{zhblock}

\begin{enblock}
But where is the simple fact? Scientists have tried to find it in the two extremes, in the infinitely great and in the infinitely small. The astronomer has found it because the distances of the stars are immense, so great that each of them appears only as a point and qualitative differences disappear, and because a point is simpler than a body which has shape and qualities. The physicist, on the other hand, has sought the elementary phenomenon in an imaginary division of bodies into infinitely small atoms, because the conditions of the problem, which undergo slow and continuous variations as we pass from one point of the body to another, may be regarded as constant within each of these little atoms. Similarly the biologist has been led instinctively to regard the cell as more interesting than the whole animal, and the event has proved him right, since cells belonging to the most diverse organisms have greater resemblances, for those who can recognize them, than the organisms themselves. The sociologist is in a more embarrassing position. The elements, which for him are men, are too dissimilar, too variable, too capricious, in a word, too complex themselves. Furthermore, history does not repeat itself; how, then, is he to select the interesting fact, the fact which is repeated? Method is precisely the selection of facts, and accordingly our first care must be to devise a method. Many have been devised because none holds the field undisputed. Nearly every sociological thesis proposes a new method, which, however, its author is very careful not to apply, so that sociology is the science with the greatest number of methods and the least results.
\end{enblock}
\begin{zhblock}
但是，简单事实在哪里呢？科学家曾试图在两个极端中寻找它：在无限大和无限小之中。天文学家找到了它，因为恒星之间的距离极其巨大，大到每一颗星看起来都只是一个点，性质上的差别消失了；而一个点比有形状、有性质的物体更简单。另一方面，物理学家则在把物体设想为无限小原子的划分中寻找基本现象，因为问题的条件在我们从物体一点移到另一点时缓慢而连续地变化，在每一个这样的小原子内部却可以被看作恒定。同样，生物学家也本能地被引导，把细胞看得比整个动物更有意思；结果证明他是对的，因为在能够识别这些相似性的人看来，属于最不同有机体的细胞之间，比这些有机体本身之间有更多相似之处。社会学家的处境则更尴尬。对他来说，组成元素是人；而人太不相同，太多变，太任性，简言之，本身就太复杂。此外，历史不会重复；那么，他怎样选择有意思的事实，怎样选择会重复的事实呢？方法恰恰就是事实的选择；因此我们的首要任务就是设计一种方法。人们已经设计出许多方法，因为没有一种能毫无争议地占据主导。几乎每一篇社会学论文都会提出一种新方法，不过作者总是十分谨慎地不去运用它；于是社会学便成了方法最多、成果最少的科学。\end{zhblock}

\begin{enblock}
It is with regular facts, therefore, that we ought to begin; but as soon as the rule is well established, as soon as it is no longer in doubt, the facts which are in complete conformity with it lose their interest, since they can teach us nothing new. Then it is the exception which becomes important. We cease to look for resemblances, and apply ourselves before all else to differences, and of these differences we select first those that are most accentuated, not only because they are the most striking, but because they will be the most instructive. This will be best explained by a simple example. Suppose we are seeking to determine a curve by observing some of the points on it. The practical man who looked only to immediate utility would merely observe the points he required for some special object; these points would be badly distributed on the curve, they would be crowded together in certain parts and scarce in others, so that it would be impossible to connect them by a continuous line, and they would be useless for any other application. The scientist would proceed in a different manner. Since he wishes to study the curve for itself, he will distribute the points to be observed regularly, and as soon as he knows some of them, he will join them by a regular line, and he will then have the complete curve. But how is he to accomplish this? If he has determined one extreme point on the curve, he will not remain close to this extremity, but will move to the other end. After the two extremities, the central point is the most instructive, and so on.
\end{enblock}
\begin{zhblock}
因此，我们应当从有规律的事实开始；但是，一旦规则已经确立，一旦它不再可疑，与它完全符合的事实就失去了兴趣，因为它们不能教给我们任何新东西。于是例外变得重要。我们不再寻找相似，而首先致力于差异；在这些差异中，我们又首先选择最显著的差异，不仅因为它们最醒目，还因为它们最有教益。一个简单例子最能说明这一点。假定我们要通过观察曲线上的若干点来确定一条曲线。只看眼前功用的实际人物，只会观察为某个特殊目的所需的点；这些点在曲线上分布很差，有的部分挤在一起，有的部分稀疏，以致无法把它们连成一条连续曲线，而且对任何其他应用都毫无用处。科学家的做法不同。既然他想研究曲线本身，他就会有规则地分布要观察的点；一旦知道其中一些点，他就会用一条有规则的线把它们连接起来，从而得到完整曲线。但他怎样做到这一点呢？如果他已经确定曲线的一个端点，他不会停留在这个端点附近，而会移向另一端。在两个端点之后，最有教益的是中心点，如此继续下去。\end{zhblock}

\begin{enblock}
Thus when a rule has been established, we have first to look for the cases in which the rule stands the best chance of being found in fault. This is one of many reasons for the interest of astronomical facts and of geological ages. By making long excursions in space or in time, we may find our ordinary rules completely upset, and these great upsettings will give us a clearer view and better comprehension of such small changes as may occur nearer us, in the small corner of the world in which we are called to live and move. We shall know this corner better for the journey we have taken into distant lands where we had no concern.
\end{enblock}
\begin{zhblock}
因此，当一条规则已经确立时，我们首先要寻找那些最可能让这条规则出错的情形。这是天文学事实和地质年代之所以令人感兴趣的许多理由之一。通过在空间或时间中作长距离远行，我们也许会发现我们的通常规则被完全颠覆；这些巨大的颠覆会使我们更清楚地看见、更好地理解那些可能发生在我们近旁的小变化，发生在我们被召唤来生活和活动的世界小角落里的变化。正因为我们曾经到那些与我们无关的远方去旅行，我们会更好地认识这个角落。\end{zhblock}

\begin{enblock}
But what we must aim at is not so much to ascertain resemblances and differences, as to discover similarities hidden under apparent discrepancies. The individual rules appear at first discordant, but on looking closer we can generally detect a resemblance; though differing in matter, they approximate in form and in the order of their parts. When we examine them from this point of view, we shall see them widen and tend to embrace everything. This is what gives a value to certain facts that come to complete a whole, and show that it is the faithful image of other known wholes.
\end{enblock}
\begin{zhblock}
不过，我们的目标与其说是确认相似和差异，不如说是发现隐藏在表面分歧之下的相似性。各个规则起初看起来并不协调，但细看之下，我们通常能够发现一种相似；它们在材料上不同，却在形式以及各部分的秩序上相近。当我们从这个角度考察它们时，就会看到它们扩展并趋向于包容一切。某些事实之所以有价值，正是因为它们来补全一个整体，并显示这个整体是其他已知整体的忠实映像。\end{zhblock}

\begin{enblock}
I cannot dwell further on this point, but these few words will suffice to show that the scientist does not make a random selection of the facts to be observed. He does not count ladybirds, as Tolstoi says, because the number of these insects, interesting as they are, is subject to capricious variations. He tries to condense a great deal of experience and a great deal of thought into a small volume, and that is why a little book on physics contains so many past experiments, and a thousand times as many possible ones, whose results are known in advance.
\end{enblock}
\begin{zhblock}
我不能在这一点上继续展开，但这几句话足以说明，科学家并不是随意选择要观察的事实。他并不像托尔斯泰所说的那样去数瓢虫，因为这些昆虫固然有趣，其数量却受任意变化支配。他试图把大量经验和大量思想压缩进一个小体积之中；正因为如此，一本小小的物理学书里包含了那么多过去的实验，也包含了多出千倍的可能实验，而这些实验的结果已经可以预先知道。\end{zhblock}

\begin{enblock}
But so far we have only considered one side of the question. The scientist does not study nature because it is useful to do so. He studies it because he takes pleasure in it, and he takes pleasure in it because it is beautiful. If nature were not beautiful it would not be worth knowing, and life would not be worth living. I am not speaking, of course, of that beauty which strikes the senses, of the beauty of qualities and appearances. I am far from despising this, but it has nothing to do with science. What I mean is that more intimate beauty which comes from the harmonious order of its parts, and which a pure intelligence can grasp. It is this that gives a body, so to speak, a skeleton to the shimmering visions that flatter our senses, and without this support the beauty of these fleeting dreams would be imperfect, because it would be indefinite and ever elusive. Intellectual beauty, on the contrary, is self-sufficing, and it is for it, more perhaps than for the future good of humanity, that the scientist condemns himself to long and painful labours.
\end{enblock}
\begin{zhblock}
但是到目前为止，我们只考察了问题的一面。科学家研究自然，并不是因为这样做有用。他研究自然，是因为他从中得到乐趣；而他从中得到乐趣，是因为自然是美的。如果自然不美，它就不值得被认识，生活也不值得过下去。当然，我说的并不是打动感官的那种美，不是性质和外观的美。我绝不轻视这种美，但它与科学无关。我所说的是一种更内在的美，它来自各部分和谐的秩序，纯粹理智能够把握它。可以说，正是这种美给那些取悦我们感官的闪烁幻象赋予了骨架；没有这种支撑，这些转瞬即逝的梦幻之美就是不完善的，因为它会是不确定的、永远逃逸的。相反，理智之美自足自立；科学家甘愿承受漫长而痛苦的劳动，也许更多是为了它，而不是为了人类未来的利益。\end{zhblock}

\begin{enblock}
It is, then, the search for this special beauty, the sense of the harmony of the world, that makes us select the facts best suited to contribute to this harmony; just as the artist selects those features of his sitter which complete the portrait and give it character and life. And there is no fear that this instinctive and unacknowledged preoccupation will divert the scientist from the search for truth. We may dream of a harmonious world, but how far it will fall short of the real world! The Greeks, the greatest artists that ever were, constructed a heaven for themselves; how poor a thing it is beside the heaven as we know it!
\end{enblock}
\begin{zhblock}
因此，正是对这种特殊美的追求，正是对世界和谐的感受，使我们选择那些最适合促成这种和谐的事实；这就像艺术家会选择模特身上那些能完成肖像、赋予肖像性格和生命的特征一样。无需担心这种本能的、未被承认的关切会使科学家偏离对真理的追求。我们也许会梦想一个和谐的世界，但它离真实世界会相差多远！希腊人是有史以来最伟大的艺术家，他们为自己建构了一个天空；与我们今日所知的天空相比，那是何等贫乏！
\end{zhblock}

\begin{enblock}
It is because simplicity and vastness are both beautiful that we seek by preference simple facts and vast facts; that we take delight, now in following the giant courses of the stars, now in scrutinizing with a microscope that prodigious smallness which is also a vastness, and now in seeking in geological ages the traces of a past that attracts us because of its remoteness.
\end{enblock}
\begin{zhblock}
正因为简单性和广大性都是美的，我们才优先寻找简单事实和宏大事实；我们时而乐于追随群星巨大的轨道，时而乐于用显微镜审视那同样也是一种广大的惊人微小，时而又在地质年代中寻找遥远过去的痕迹，而正是这种遥远吸引着我们。\end{zhblock}

\begin{enblock}
Thus we see that care for the beautiful leads us to the same selection as care for the useful. Similarly economy of thought, that economy of effort which, according to Mach, is the constant tendency of science, is a source of beauty as well as a practical advantage. The buildings we admire are those in which the architect has succeeded in proportioning the means to the end, in which the columns seem to carry the burdens imposed on them lightly and without effort, like the graceful caryatids of the Erechtheum.
\end{enblock}
\begin{zhblock}
由此可见，对美的关切会把我们引向与对有用性的关切相同的选择。同样，思想的经济性，即马赫所说的科学恒常倾向于追求的努力经济性，既是实践上的优点，也是美的源泉。我们所赞赏的建筑，是建筑师成功地使手段与目的相称的建筑；在那里，柱子似乎轻盈而毫不费力地承受着加在它们身上的负荷，就像厄瑞克忒翁神庙中优雅的女像柱。\end{zhblock}

\begin{enblock}
Whence comes this concordance? Is it merely that things which seem to us beautiful are those which are best adapted to our intelligence, and that consequently they are at the same time the tools that intelligence knows best how to handle? Or is it due rather to evolution and natural selection? Have the peoples whose ideal conformed best to their own interests, properly understood, exterminated the others and taken their place? One and all pursued their ideal without considering the consequences, but while this pursuit led some to their destruction, it gave empire to others. We are tempted to believe this, for if the Greeks triumphed over the barbarians, and if Europe, heir of the thought of the Greeks, dominates the world, it is due to the fact that the savages loved garish colours and the blatant noise of the drum, which appealed to their senses, while the Greeks loved the intellectual beauty hidden behind sensible beauty, and that it is this beauty which gives certainty and strength to the intelligence.
\end{enblock}
\begin{zhblock}
这种一致从何而来？是否只是因为那些在我们看来美的东西，正是最适合我们理智的东西，因而同时也是理智最懂得如何运用的工具？或者，它更应归因于进化和自然选择？那些理想最符合自身真正利益的民族，是否消灭了其他民族并取而代之？所有民族都追求自己的理想，并不考虑后果；但这种追求有的把人引向毁灭，有的却赋予人以统治力。我们很容易相信这一点，因为如果希腊人战胜了蛮族，如果继承希腊思想的欧洲支配了世界，那是因为野蛮人喜爱刺眼的色彩和喧嚣的鼓声，这些诉诸感官；而希腊人热爱隐藏在感性美之后的理智之美，正是这种美给予理智以确定性和力量。\end{zhblock}

\begin{enblock}
No doubt Tolstoi would be horrified at such a triumph, and he would refuse to admit that it could be truly useful. But this disinterested pursuit of truth for its own beauty is also wholesome, and can make men better. I know very well there are disappointments, that the thinker does not always find the serenity he should, and even that some scientists have thoroughly bad tempers.
\end{enblock}
\begin{zhblock}
托尔斯泰无疑会对这样的胜利感到震惊，并拒绝承认它可能真正有用。但是，这种为了真理自身之美而无私追求真理的活动，也同样有益健康，并能使人变得更好。我当然知道，其中会有失望；思想者并不总能获得他本应拥有的宁静，甚至有些科学家的脾气还糟糕透顶。\end{zhblock}

\begin{enblock}
Must we therefore say that science should be abandoned, and morality alone be studied? Does anyone suppose that moralists themselves are entirely above reproach when they have come down from the pulpit?
\end{enblock}
\begin{zhblock}
因此我们就必须说应当放弃科学，只研究道德吗？难道有人以为，道德家一旦从讲坛上走下来，他们自己就完全无可指摘吗？
\end{zhblock}

\begin{center}
[\dots]
\end{center}

\section*{III. Mathematical Discovery\\三、数学发现}

\begin{enblock}
The genesis of mathematical discovery is a problem which must inspire the psychologist with the keenest interest. For this is the process in which the human mind seems to borrow least from the exterior world, in which it acts, or appears to act, only by itself and on itself, so that by studying the process of geometric thought we may hope to arrive at what is most essential in the human mind.
\end{enblock}
\begin{zhblock}
数学发现的发生，是一个必定会引起心理学家最强烈兴趣的问题。因为在这一过程中，人类心智似乎最少向外部世界借取东西；它只凭自身并作用于自身，或者看起来如此。因此，通过研究几何思想的过程，我们也许有希望抵达人类心智中最本质的东西。\end{zhblock}

\begin{enblock}
This has long been understood, and a few months ago a review called \emph{L'Enseignement mathématique}, edited by MM. Laisant and Fehr, instituted an enquiry into the habits of mind and methods of work of different mathematicians. I had outlined the principal features of this article when the results of the enquiry were published, so that I have hardly been able to make any use of them, and I will content myself with saying that the majority of the evidence confirms my conclusions. I do not say there is unanimity, for on an appeal to universal suffrage we cannot hope to obtain unanimity.
\end{enblock}
\begin{zhblock}
这一点早已为人所理解。几个月前，由莱桑和费尔二位先生编辑的《数学教育》杂志，发起了一项调查，研究不同数学家的思维习惯和工作方法。当调查结果发表时，本文的主要轮廓我已经拟好，所以几乎没有能够利用这些材料；我只满足于说，大多数证据都证实了我的结论。我并不说意见完全一致，因为诉诸普遍投票时，我们不能指望得到全体一致。\end{zhblock}

\begin{enblock}
One first fact must astonish us, or rather would astonish us if we were not too much accustomed to it. How does it happen that there are people who do not understand mathematics? If the science invokes only the rules of logic, those accepted by all well-formed minds, if its evidence is founded on principles that are common to all men, and that none but a madman would attempt to deny, how does it happen that there are so many people who are entirely impervious to it?
\end{enblock}
\begin{zhblock}
首先有一个事实必定使我们惊讶，或者说，如果我们不是已经太习惯于它，它本来会使我们惊讶：为什么会有人不理解数学？如果这门科学所诉诸的只是逻辑规则，只是所有健全心智都承认的规则；如果它的证据建立在所有人共有、除了疯子无人会试图否认的原则之上，那么为什么竟有那么多人对它完全不能领会？
\end{zhblock}

\begin{enblock}
There is nothing mysterious in the fact that everyone is not capable of discovery. That everyone should not be able to retain a demonstration he has once learnt is still comprehensible. But what does seem most surprising, when we consider it, is that anyone should be unable to understand a mathematical argument at the very moment it is stated to him. And yet those who can only follow the argument with difficulty are in a majority; this is incontestable, and the experience of teachers of secondary education will certainly not contradict me.
\end{enblock}
\begin{zhblock}
并非每个人都能作出发现，这并不神秘。一个人学过一个证明之后不能把它记住，这仍然可以理解。但是认真想来，最令人惊讶的似乎是：有人竟然在一个数学论证正向他陈述的当下就不能理解它。然而，只有困难地跟上论证的人占多数；这一点无可争辩，中学教师的经验当然也不会反驳我。\end{zhblock}

\begin{enblock}
And still further, how is error possible in mathematics? A healthy intellect should not be guilty of any error in logic, and yet there are very keen minds which will not make a false step in a short argument such as those we have to make in the ordinary actions of life, which yet are incapable of following or repeating without error the demonstrations of mathematics which are longer, but which are, after all, only accumulations of short arguments exactly analogous to those they make so easily. Is it necessary to add that mathematicians themselves are not infallible?
\end{enblock}
\begin{zhblock}
再进一步说，数学中的错误怎么可能发生？一个健全的理智不应当犯任何逻辑错误；然而，有些非常敏锐的头脑，在日常生活行动中所需的短论证里不会走错一步，却无法无误地跟随或复述较长的数学证明，尽管这些证明归根到底只是许多短论证的累积，而这些短论证恰好与他们轻易作出的那些完全类似。还需要补充说，数学家本人也并非不会犯错吗？\end{zhblock}

\begin{enblock}
The answer appears to me obvious. Imagine a long series of syllogisms in which the conclusions of those that precede form the premises of those that follow. We shall be capable of grasping each of the syllogisms, and it is not in the passage from premises to conclusion that we are in danger of going astray. But between the moment when we meet a proposition for the first time as the conclusion of one syllogism, and the moment when we find it once more as the premise of another syllogism, much time will sometimes have elapsed, and we shall have unfolded many links of the chain; accordingly it may well happen that we shall have forgotten it, or, what is more serious, forgotten its meaning. So we may chance to replace it by a somewhat different proposition, or to preserve the same statement but give it a slightly different meaning, and thus we are in danger of falling into error.
\end{enblock}
\begin{zhblock}
在我看来，答案是显而易见的。设想一长串三段论，前面三段论的结论构成后面三段论的前提。我们能够把握每一个三段论；危险并不在于从前提过渡到结论的那一步。但是，从我们第一次把某个命题作为一个三段论的结论遇到它，到我们再次把它作为另一个三段论的前提发现它，这中间有时会经过很长时间，我们也已经展开了链条中的许多环节。因此，很可能发生的情况是，我们已经忘记了它，或者更严重地说，忘记了它的意义。于是我们可能偶然用一个略有不同的命题替代它，或者保留同一个表述，却赋予它稍稍不同的意义；这样我们就有陷入错误的危险。\end{zhblock}

\begin{enblock}
A mathematician must often use a rule, and, naturally, he begins by demonstrating the rule. At the moment the demonstration is quite fresh in his memory he understands perfectly its meaning and significance, and he is in no danger of changing it. But later on he commits it to memory, and only applies it in a mechanical way, and then, if his memory fails him, he may apply it wrongly. It is thus, to take a simple and almost vulgar example, that we sometimes make mistakes in calculation, because we have forgotten our multiplication table.
\end{enblock}
\begin{zhblock}
数学家常常必须使用一条规则；自然，他一开始会证明这条规则。当证明在记忆中还十分新鲜时，他完全理解它的含义和意义，也没有改变它的危险。但后来他把它交给记忆，只是以机械方式加以应用；这时如果记忆失灵，他就可能用错。举一个简单得几乎粗俗的例子，我们有时会在计算中出错，正是因为忘记了乘法表。\end{zhblock}

\begin{enblock}
On this view special aptitude for mathematics would be due to nothing but a very certain memory or a tremendous power of attention. It would be a quality analogous to that of the whist player who can remember the cards played, or, to rise a step higher, to that of the chess player who can picture a very great number of combinations and retain them in his memory. Every good mathematician should also be a good chess player and vice versa, and similarly he should be a good numerical calculator. Certainly this sometimes happens, and thus Gauss was at once a geometrician of genius and a very precocious and very certain calculator.
\end{enblock}
\begin{zhblock}
按这种看法，数学的特殊才能只不过来自非常可靠的记忆力或巨大的注意力。它会类似于惠斯特牌手能记住已经打出的牌，或者再高一步，类似于棋手能够想象大量组合并把它们保存在记忆中。每个优秀数学家也都应当是优秀棋手，反之亦然；同样，他也应当是优秀的数字计算者。当然，这种情况有时确实发生；高斯既是天才几何学家，也是一位早熟而极其可靠的计算者。\end{zhblock}

\begin{enblock}
But there are exceptions, or rather I am wrong, for I cannot call them exceptions, otherwise the exceptions would be more numerous than the cases of conformity with the rule. On the contrary, it was Gauss who was an exception. As for myself, I must confess I am absolutely incapable of doing an addition sum without a mistake. Similarly I should be a very bad chess player. I could easily calculate that by playing in a certain way I should be exposed to such and such a danger; I should then review many other moves, which I should reject for other reasons, and I should end by making the move I first examined, having forgotten in the interval the danger I had foreseen.
\end{enblock}
\begin{zhblock}
但也有例外；或者更确切地说，我说错了，因为我不能把它们称为例外，否则例外会比符合规则的情形更多。相反，高斯才是例外。至于我自己，必须承认，我完全不能保证做一个加法题而不出错。同样，我也会是一个很糟糕的棋手。我可以轻易算出，以某种方式走棋会使我暴露于某种危险；随后我会考察许多其他走法，又因别的理由把它们否定；最后我会走最初考察的那一步，却已经在这中间忘记了我曾预见的危险。\end{zhblock}

\begin{enblock}
In a word, my memory is not bad, but it would be insufficient to make me a good chess player. Why, then, does it not fail me in a difficult mathematical argument in which the majority of chess players would be lost? Clearly because it is guided by the general trend of the argument. A mathematical demonstration is not a simple juxtaposition of syllogisms; it consists of syllogisms placed in a certain order, and the order in which these elements are placed is much more important than the elements themselves. If I have the feeling, so to speak the intuition, of this order, so that I can perceive the whole of the argument at a glance, I need no longer be afraid of forgetting one of the elements; each of them will place itself naturally in the position prepared for it, without my having to make any effort of memory.
\end{enblock}
\begin{zhblock}
总之，我的记忆并不差，但它不足以使我成为一个好棋手。那么，为什么在一个困难的数学论证中，在大多数棋手都会迷失的地方，它却没有辜负我呢？显然，是因为它受到论证总体趋向的引导。数学证明并不是三段论的简单并列；它由按某种秩序安置的三段论构成，而这些元素被安置的秩序比元素本身重要得多。如果我对这种秩序有一种感觉，或者说一种直觉，以至于能够一眼看见论证的整体，我就不必再害怕忘记其中某个元素；每个元素都会自然地安放到为它准备的位置上，而无需我作任何记忆努力。\end{zhblock}

\begin{enblock}
It seems to me, then, as I repeat an argument I have learnt, that I could have discovered it. This is often only an illusion; but even then, even if I am not clever enough to create for myself, I rediscover it myself as I repeat it.
\end{enblock}
\begin{zhblock}
因此，当我复述一个已经学过的论证时，我似乎觉得自己本来也能够发现它。这常常只是一种幻觉；但即便如此，即便我并不聪明到能够自行创造，我在复述它时也仍是在亲自重新发现它。\end{zhblock}

\begin{enblock}
We can understand that this feeling, this intuition of mathematical order, which enables us to guess hidden harmonies and relations, cannot belong to everyone. Some have neither this delicate feeling that is difficult to define, nor a power of memory and attention above the common, and so they are absolutely incapable of understanding even the first steps of higher mathematics. This applies to the majority of people. Others have the feeling only in a slight degree, but they are gifted with an uncommon memory and a great capacity for attention. They learn the details one after the other by heart, they can understand mathematics and sometimes apply them, but they are not in a condition to create. Lastly, others possess the special intuition I have spoken of more or less highly developed, and they cannot only understand mathematics, even though their memory is in no way extraordinary, but they can become creators, and seek to make discovery with more or less chance of success, according as their intuition is more or less developed.
\end{enblock}
\begin{zhblock}
我们可以理解，这种感觉，这种对数学秩序的直觉，使我们能够猜出隐藏的和谐与关系，但它不可能属于每一个人。有些人既没有这种难以界定的微妙感觉，也没有超出常人的记忆力和注意力，所以他们甚至完全无法理解高等数学的最初几步。多数人就是如此。另一些人只在很轻微的程度上具有这种感觉，却拥有非同寻常的记忆力和很强的注意力。他们把细节一个接一个地背下来，能够理解数学，有时也能应用数学，但他们并不具备创造的条件。最后，还有一些人或多或少高度发展了我所说的特殊直觉；即便他们的记忆并无任何非凡之处，他们不仅能理解数学，而且能够成为创造者，并依其直觉发展程度的高低，在或大或小的成功机会中寻求发现。\end{zhblock}

\begin{enblock}
What, in fact, is mathematical discovery? It does not consist in making new combinations with mathematical entities that are already known. That can be done by anyone, and the combinations that could be so formed would be infinite in number, and the greater part of them would be absolutely devoid of interest. Discovery consists precisely in not constructing useless combinations, but in constructing those that are useful, which are an infinitely small minority. Discovery is discernment, selection.
\end{enblock}
\begin{zhblock}
那么，数学发现究竟是什么？它并不在于用已经知道的数学对象作出新的组合。任何人都能这样做，而能够这样形成的组合数量是无限的，其中绝大部分毫无兴趣。发现恰恰在于不构造无用组合，而构造那些有用的组合；后者只是无限小的少数。发现就是辨别，就是选择。\end{zhblock}

\begin{enblock}
How this selection is to be made I have explained above. Mathematical facts worthy of being studied are those which, by their analogy with other facts, are capable of conducting us to the knowledge of a mathematical law, in the same way that experimental facts conduct us to the knowledge of a physical law. They are those which reveal unsuspected relations between other facts, long since known, but wrongly believed to be unrelated to each other.
\end{enblock}
\begin{zhblock}
这种选择应当如何作出，我在上面已经说明。值得研究的数学事实，是那些凭借它们与其他事实的类比，能够引导我们认识一条数学定律的事实；这与实验事实引导我们认识一条物理定律相同。它们揭示出那些早已为人所知、却被错误地认为彼此无关的事实之间未曾料到的关系。\end{zhblock}

\begin{enblock}
Among the combinations we choose, the most fruitful are often those which are formed of elements borrowed from widely separated domains. I do not mean to say that for discovery it is sufficient to bring together objects that are as incongruous as possible. The greater part of the combinations so formed would be entirely fruitless, but some among them, though very rare, are the most fruitful of all.
\end{enblock}
\begin{zhblock}
在我们选择的组合中，最富成果的往往是那些由来自相距甚远领域的元素构成的组合。我并不是说，为了发现，只要把尽可能不相干的对象放在一起就够了。这样形成的大多数组合会完全没有结果；但其中少数虽然极为罕见，却是最富成果的。\end{zhblock}

\begin{enblock}
Discovery, as I have said, is selection. But this is perhaps not quite the right word. It suggests a purchaser who has been shown a large number of samples, and examines them one after the other in order to make his selection. In our case the samples would be so numerous that a whole life would not give sufficient time to examine them. Things do not happen in this way. Unfruitful combinations do not so much as present themselves to the mind of the discoverer. In the field of his consciousness there never appear any but really useful combinations, and some that he rejects, which, however, partake to some extent of the character of useful combinations. Everything happens as if the discoverer were a secondary examiner who had only to interrogate candidates declared eligible after passing a preliminary test.
\end{enblock}
\begin{zhblock}
正如我所说，发现就是选择。但这也许并不是完全恰当的词。它使人想到一个买主，面前摆着大量样品，他逐一检查，以便作出选择。在我们的情形中，样品数量如此巨大，即使用一生也没有足够时间去检查。事情并不是这样发生的。没有成果的组合甚至不会呈现在发现者的心中。在他的意识领域中，出现的从来只是确有用处的组合，以及某些被他拒绝的组合，不过这些被拒绝者在某种程度上也带有有用组合的性质。一切仿佛发现者只是第二级考官，只需询问那些已经通过初试、被宣布合格的候选者。\end{zhblock}

\begin{enblock}
But what I have said up to now is only what can be observed or inferred by reading the works of geometricians, provided they are read with some reflection.
\end{enblock}
\begin{zhblock}
但是，到目前为止我所说的，只是只要带着思考阅读几何学家的著作，便能够观察到或推断出的东西。\end{zhblock}

\begin{enblock}
It is time to penetrate further, and to see what happens in the very soul of the mathematician. For this purpose I think I cannot do better than recount my personal recollections. Only I am going to confine myself to relating how I wrote my first treatise on Fuchsian functions. I must apologize, for I am going to introduce some technical expressions, but they need not alarm the reader, for he has no need to understand them. I shall say, for instance, that I found the demonstration of such and such a theorem under such and such circumstances; the theorem will have a barbarous name that many will not know, but that is of no importance. What is interesting for the psychologist is not the theorem but the circumstances.
\end{enblock}
\begin{zhblock}
现在是进一步深入的时候了，要看看数学家的灵魂深处究竟发生了什么。为此，我想没有比叙述自己的个人回忆更好的办法。不过，我只打算讲述我怎样写出第一篇关于富克斯函数的论文。我必须先致歉，因为我要引入一些技术性表述；但它们不必使读者惊慌，因为读者并不需要理解它们。比如，我会说我在某种情形下找到了某某定理的证明；这个定理会有一个许多人不知道的生僻名称，但这并不重要。对心理学家来说，有趣的不是定理，而是情形。\end{zhblock}

\begin{enblock}
For a fortnight I had been attempting to prove that there could not be any function analogous to what I have since called Fuchsian functions. I was at that time very ignorant. Every day I sat down at my table and spent an hour or two trying a great number of combinations, and I arrived at no result. One night I took some black coffee, contrary to my custom, and was unable to sleep. A host of ideas kept surging in my head; I could almost feel them jostling one another, until two of them coalesced, so to speak, to form a stable combination. When morning came, I had established the existence of one class of Fuchsian functions, those that are derived from the hypergeometric series. I had only to verify the results, which only took a few hours.
\end{enblock}
\begin{zhblock}
有两周时间，我一直试图证明，不可能存在任何类似于我后来称为富克斯函数的函数。当时我非常无知。每天我坐到桌前，花一两个小时尝试大量组合，却得不到任何结果。一天夜里，我一反常态喝了一些黑咖啡，于是无法入睡。许多想法不断在我头脑中涌起；我几乎能感觉到它们彼此碰撞，直到其中两个可以说是合并起来，形成了一个稳定组合。到早晨时，我已经确定了一类富克斯函数的存在，即那些由超几何级数导出的函数。我只需要验证结果，而这只花了几个小时。\end{zhblock}

\begin{enblock}
Then I wished to represent these functions by the quotient of two series. This idea was perfectly conscious and deliberate; I was guided by the analogy with elliptic functions. I asked myself what must be the properties of these series, if they existed, and I succeeded without difficulty in forming the series that I have called Theta-Fuchsian.
\end{enblock}
\begin{zhblock}
随后我想用两个级数的商来表示这些函数。这个想法完全是有意识且有意为之的；我受到与椭圆函数类比的引导。我问自己，如果这些级数存在，它们必须具有什么性质；于是我毫不困难地构造出了我称为 Θ-富克斯级数的级数。\end{zhblock}

\begin{enblock}
At this moment I left Caen, where I was then living, to take part in a geological conference arranged by the School of Mines. The incidents of the journey made me forget my mathematical work. When we arrived at Coutances, we got into a brake to go for a drive, and, just as I put my foot on the step, the idea came to me, though nothing in my former thoughts seemed to have prepared me for it, that the transformations I had used to define Fuchsian functions were identical with those of non-Euclidean geometry. I made no verification, and had no time to do so, since I took up the conversation again as soon as I had sat down in the brake, but I felt absolute certainty at once. When I got back to Caen I verified the result at my leisure to satisfy my conscience.
\end{enblock}
\begin{zhblock}
就在这时，我离开当时居住的卡昂，去参加矿业学院安排的一次地质会议。旅途中的种种事情使我忘记了数学工作。我们抵达库唐斯后，登上一辆马车外出兜风；就在我把脚踏上踏板的瞬间，一个想法来到我心中，尽管我此前的任何思考似乎都没有为它作准备：我用来定义富克斯函数的变换，与非欧几何中的变换是同一回事。我没有作验证，也没有时间验证，因为一在马车里坐下，我就重新开始谈话；但我立刻感到绝对确定。回到卡昂后，我从容地验证了这个结果，以安慰自己的良心。\end{zhblock}

\begin{enblock}
I then began to study arithmetical questions without any great apparent result, and without suspecting that they could have the least connection with my previous researches. Disgusted at my want of success, I went away to spend a few days at the seaside, and thought of entirely different things. One day, as I was walking on the cliff, the idea came to me, again with the same characteristics of conciseness, suddenness, and immediate certainty, that arithmetical transformations of indefinite ternary quadratic forms are identical with those of non-Euclidean geometry.
\end{enblock}
\begin{zhblock}
随后我开始研究算术问题，表面上没有取得什么重大结果，也没有怀疑它们会与我此前的研究有哪怕最轻微的联系。由于厌恶自己的失败，我离开去海边住了几天，思考完全不同的事情。有一天，我在悬崖上散步时，一个想法又以同样简洁、突然、立刻确定的特征来到我心中：不定三元二次型的算术变换，与非欧几何中的变换是同一回事。\end{zhblock}

\begin{enblock}
Returning to Caen, I reflected on this result and deduced its consequences. The example of quadratic forms showed me that there are Fuchsian groups other than those which correspond with the hypergeometric series; I saw that I could apply to them the theory of the Theta-Fuchsian series, and that, consequently, there are Fuchsian functions other than those which are derived from the hypergeometric series, the only ones I knew up to that time. Naturally, I proposed to form all these functions. I laid siege to them systematically and captured all the outworks one after the other. There was one, however, which still held out, whose fall would carry with it that of the central fortress. But all my efforts were of no avail at first, except to make me better understand the difficulty, which was already something. All this work was perfectly conscious.
\end{enblock}
\begin{zhblock}
回到卡昂后，我思考这个结果并推出它的后果。二次型的例子使我看到，除了与超几何级数相对应的富克斯群之外，还有其他富克斯群；我看出可以把 Θ-富克斯级数理论应用到它们身上，因此，除了由超几何级数导出的那些富克斯函数，也就是我此前唯一知道的那些之外，还存在其他富克斯函数。自然，我打算构造所有这些函数。我系统地围攻它们，一个接一个攻克外围工事。然而还有一个仍在坚守；一旦它陷落，中心堡垒也会随之陷落。但起初我所有努力都无济于事，只是让我更好地理解困难，而这已经算有所收获。所有这些工作都是完全有意识的。\end{zhblock}

\begin{enblock}
Thereupon I left for Mont-Valérien, where I had to serve my time in the army, and so my mind was preoccupied with very different matters. One day, as I was crossing the street, the solution of the difficulty which had brought me to a standstill came to me all at once. I did not try to fathom it immediately, and it was only after my service was finished that I returned to the question. I had all the elements, and had only to assemble and arrange them. Accordingly I composed my definitive treatise at a sitting and without any difficulty.
\end{enblock}
\begin{zhblock}
随后我前往蒙瓦莱里安，在那里我必须服兵役，于是我的心思被完全不同的事情占据。有一天，我正在过街，使我陷入停顿的那个困难的解法突然一下子来到我心中。我没有立即试图深究它；直到服役结束后，我才回到这个问题。我已经拥有所有元素，只需把它们集合并安排起来。于是我一气呵成写出了定稿，没有任何困难。\end{zhblock}

\begin{enblock}
It is useless to multiply examples, and I will content myself with this one alone. As regards my other researches, the accounts I should give would be exactly similar, and the observations related by other mathematicians in the enquiry of \emph{L'Enseignement mathématique} would only confirm them.
\end{enblock}
\begin{zhblock}
再增加例子没有必要，我满足于这一个例子。至于我的其他研究，我所能给出的叙述会完全类似；而《数学教育》调查中其他数学家讲述的观察，也只会证实这些叙述。\end{zhblock}

\begin{enblock}
One is at once struck by these appearances of sudden illumination, obvious indications of a long course of previous unconscious work. The part played by this unconscious work in mathematical discovery seems to me indisputable, and we shall find traces of it in other cases where it is less evident. Often when a man is working at a difficult question, he accomplishes nothing the first time he sets to work. Then he takes more or less of a rest, and sits down again at his table. During the first half-hour he still finds nothing, and then all at once the decisive idea presents itself to his mind. We might say that the conscious work proved more fruitful because it was interrupted and the rest restored force and freshness to the mind. But it is more probable that the rest was occupied with unconscious work, and that the result of this work was afterwards revealed to the geometrician exactly as in the cases I have quoted, except that the revelation, instead of coming to light during a walk or a journey, came during a period of conscious work, but independently of that work, which at most only performs the unlocking process, as if it were the spur that excited into conscious form the results already acquired during the rest, which till then remained unconscious.
\end{enblock}
\begin{zhblock}
这些突然照明的出现立刻令人震动；它们显然表明此前曾有漫长的无意识工作。无意识工作在数学发现中所起的作用，在我看来无可争辩；在其他不那么明显的情形中，我们也会发现它的痕迹。一个人研究一个困难问题时，常常第一次着手工作毫无所得。随后他或多或少休息一下，再坐回桌前。最初半小时他仍然什么也找不到；然后决定性的想法突然呈现在心中。我们也许会说，有意识工作之所以更有成果，是因为它曾被中断，而休息恢复了心智的力量和新鲜感。但更可能的是，休息期间充满了无意识工作，而这项工作的结果随后向几何学家显现，正如我所举的那些情形一样；只是这种显现并不是在散步或旅行中出现，而是在一段有意识工作中出现，却独立于那项工作。后者最多只是起了解锁作用，仿佛一个刺激，把休息期间已经获得、但此前仍停留在无意识中的结果激发成有意识的形式。\end{zhblock}

\begin{enblock}
There is another remark to be made regarding the conditions of this unconscious work, which is, that it is not possible, or in any case not fruitful, unless it is first preceded and then followed by a period of conscious work. These sudden inspirations are never produced (and this is sufficiently proved already by the examples I have quoted) except after some days of voluntary efforts which appeared absolutely fruitless, in which one thought one had accomplished nothing, and seemed to be on a totally wrong track. These efforts, however, were not as barren as one thought; they set the unconscious machine in motion, and without them it would not have worked at all, and would not have produced anything.
\end{enblock}
\begin{zhblock}
关于这种无意识工作的条件，还应作另一点说明：除非它先有一段有意识工作在前，后又有一段有意识工作在后，否则它是不可能的，至少不会有成果。这些突然的灵感从不产生，除非在若干天自愿努力之后；我前面举的例子已经充分证明了这一点。这些努力看起来完全没有成果，人们以为自己什么也没有完成，似乎走在完全错误的路上。然而，这些努力并不像人们想的那样贫瘠；它们启动了无意识机器，没有它们，机器根本不会运转，也不会产生任何东西。\end{zhblock}

\begin{enblock}
The necessity for the second period of conscious work can be even more readily understood. It is necessary to work out the results of the inspiration, to deduce the immediate consequences and put them in order and to set out the demonstrations; but, above all, it is necessary to verify them. I have spoken of the feeling of absolute certainty which accompanies the inspiration; in the cases quoted this feeling was not deceptive, and more often than not this will be the case. But we must beware of thinking that this is a rule without exceptions. Often the feeling deceives us without being any less distinct on that account, and we only detect it when we attempt to establish the demonstration. I have observed this fact most notably with regard to ideas that have come to me in the morning or at night when I have been in bed in a semi-somnolent condition.
\end{enblock}
\begin{zhblock}
第二段有意识工作的必要性更容易理解。我们必须展开灵感所得的结果，推出直接后果并把它们整理好，还要陈述证明；但最重要的是，必须加以验证。我曾说过，灵感伴随着一种绝对确定的感觉；在我引用的那些情形中，这种感觉并没有欺骗我，而且多数时候都会如此。但是，我们必须小心，不要以为这是没有例外的规则。那种感觉常常会欺骗我们，尽管它并不会因此变得不鲜明；只有当我们试图建立证明时，才会察觉这一点。对于早晨或夜里我躺在床上、处于半睡半醒状态时来到我心中的想法，我尤其明显地观察到这一事实。\end{zhblock}

\begin{enblock}
Such are the facts of the case, and they suggest the following reflections. The result of all that precedes is to show that the unconscious ego, or, as it is called, the subliminal ego, plays a most important part in mathematical discovery. But the subliminal ego is generally thought of as purely automatic. Now we have seen that mathematical work is not a simple mechanical work, and that it could not be entrusted to any machine, whatever the degree of perfection we suppose it to have been brought to. It is not merely a question of applying certain rules, of manufacturing as many combinations as possible according to certain fixed laws. The combinations so obtained would be extremely numerous, useless, and encumbering. The real work of the discoverer consists in choosing between these combinations with a view to eliminating those that are useless, or rather not giving himself the trouble of making them at all. The rules which must guide this choice are extremely subtle and delicate, and it is practically impossible to state them in precise language; they must be felt rather than formulated. Under these conditions, how can we imagine a sieve capable of applying them mechanically?
\end{enblock}
\begin{zhblock}
事实就是这样，它们提示了以下思考。前面一切所得出的结果，是表明无意识自我，或者所谓潜意识自我，在数学发现中起着极其重要的作用。但是，潜意识自我通常被认为是纯粹自动的。现在我们已经看到，数学工作并不是简单的机械工作，也不能交给任何机器，无论我们假定这种机器已经达到怎样的完善程度。问题并不只是应用某些规则，按照某些固定法则制造出尽可能多的组合。这样得到的组合会极其众多、无用而碍事。发现者的真正工作在于在这些组合之间作出选择，以便排除无用的组合，或者更确切地说，根本不费力去制造它们。必须引导这种选择的规则极其微妙而精细，实际上不可能用精确语言表述；它们必须被感觉到，而不是被公式化。在这种条件下，我们怎能设想有一种筛子能够机械地应用这些规则呢？
\end{zhblock}

\begin{enblock}
The following, then, presents itself as a first hypothesis. The subliminal ego is in no way inferior to the conscious ego; it is not purely automatic; it is capable of discernment; it has tact and lightness of touch; it can select, and it can divine. More than that, it can divine better than the conscious ego, since it succeeds where the latter fails. In a word, is not the subliminal ego superior to the conscious ego? The importance of this question will be readily understood. In a recent lecture, M. Boutroux showed how it had arisen on entirely different occasions, and what consequences would be involved by an answer in the affirmative. (See also the same author's \emph{Science et religion}, pp. 313 et seq.)
\end{enblock}
\begin{zhblock}
于是，作为第一个假设，下面这种想法呈现出来：潜意识自我丝毫不低于意识自我；它并非纯粹自动；它有辨别力，有分寸感和轻盈的触感；它能够选择，也能够猜测。不止如此，它比意识自我更善于猜测，因为它在后者失败的地方取得成功。简言之，潜意识自我岂不是高于意识自我吗？这个问题的重要性很容易理解。在最近一次演讲中，布特鲁先生说明了它怎样在完全不同的场合中出现，以及若作肯定回答会牵涉哪些后果。（亦参见同一作者《科学与宗教》第313页及以下。）
\end{zhblock}

\begin{enblock}
Are we forced to give this affirmative answer by the facts I have just stated? I confess that, for my part, I should be loath to accept it. Let us, then, return to the facts, and see if they do not admit of some other explanation.
\end{enblock}
\begin{zhblock}
我刚才陈述的事实是否迫使我们作出这种肯定回答呢？我承认，就我而言，我很不愿接受它。那么，让我们回到事实，看看它们是否还能容许别的解释。\end{zhblock}

\begin{enblock}
It is certain that the combinations which present themselves to the mind in a kind of sudden illumination after a somewhat prolonged period of unconscious work are generally useful and fruitful combinations, which appear to be the result of a preliminary sifting. Does it follow from this that the subliminal ego, having divined by a delicate intuition that these combinations could be useful, has formed none but these, or has it formed a great many others which were devoid of interest, and remained unconscious?
\end{enblock}
\begin{zhblock}
可以肯定的是，在一段相当长的无意识工作之后，以某种突然照明的方式呈现在心中的组合，通常是有用且富有成果的组合；它们看起来像是一次初步筛选的结果。是否由此就能推出，潜意识自我凭借微妙直觉猜出这些组合可能有用，于是只形成了这些组合？还是说，它也形成了大量没有兴趣的其他组合，而这些组合始终停留在无意识中？\end{zhblock}

\begin{enblock}
Under this second aspect, all the combinations are formed as a result of the automatic action of the subliminal ego, but those only which are interesting find their way into the field of consciousness. This, too, is most mysterious. How can we explain the fact that, of the thousand products of our unconscious activity, some are invited to cross the threshold, while others remain outside? Is it mere chance that gives them this privilege? Evidently not. For instance, of all the excitements of our senses, it is only the most intense that retain our attention, unless it has been directed upon them by other causes. More commonly the privileged unconscious phenomena, those that are capable of becoming conscious, are those which, directly or indirectly, most deeply affect our sensibility.
\end{enblock}
\begin{zhblock}
从第二种角度看，一切组合都是潜意识自我自动活动的结果，但只有有趣的组合进入意识领域。这同样十分神秘。在我们无意识活动的上千个产物中，有些被邀请越过门槛，而另一些留在门外；这个事实如何解释？赋予它们这种特权的只是偶然吗？显然不是。例如，在我们所有感官刺激中，除非注意力因其他原因被引向它们，只有最强烈的刺激会保持我们的注意。更通常地说，那些享有特权的无意识现象，那些能够变成意识的现象，是直接或间接最深地影响我们感受性的现象。\end{zhblock}

\begin{enblock}
It may appear surprising that sensibility should be introduced in connection with mathematical demonstrations, which, it would seem, can only interest the intellect. But not if we bear in mind the feeling of mathematical beauty, of the harmony of numbers and forms and of geometric elegance. It is a real aesthetic feeling that all true mathematicians recognize, and this is truly sensibility.
\end{enblock}
\begin{zhblock}
把感受性引入数学证明的讨论，乍看可能令人惊讶，因为数学证明似乎只能使理智感兴趣。但如果我们记住数学之美的感觉，记住数与形的和谐以及几何的优雅，情况就不是这样了。这是一种真正的审美感觉，所有真正的数学家都承认它；而这确实就是感受性。\end{zhblock}

\begin{enblock}
Now, what are the mathematical entities to which we attribute this character of beauty and elegance, which are capable of developing in us a kind of aesthetic emotion? Those whose elements are harmoniously arranged so that the mind can, without effort, take in the whole without neglecting the details. This harmony is at once a satisfaction to our aesthetic requirements, and an assistance to the mind which it supports and guides. At the same time, by setting before our eyes a well-ordered whole, it gives us a presentiment of a mathematical law. Now, as I have said above, the only mathematical facts worthy of retaining our attention and capable of being useful are those which can make us acquainted with a mathematical law. Accordingly we arrive at the following conclusion. The useful combinations are precisely the most beautiful, I mean those that can most charm that special sensibility that all mathematicians know, but of which laymen are so ignorant that they are often tempted to smile at it.
\end{enblock}
\begin{zhblock}
那么，我们把这种美和优雅的性质归于哪些数学实体？哪些实体能够在我们心中发展出某种审美情感？是那些其元素和谐排列、使心智能够毫不费力地把握整体而又不忽略细节的实体。这种和谐既满足我们的审美要求，又帮助、支撑并引导心智。同时，它把一个秩序良好的整体呈现在我们眼前，使我们预感到一条数学定律。正如我上面所说，唯一值得保留我们注意、并且能够有用的数学事实，是那些能使我们认识一条数学定律的事实。因此我们得到如下结论：有用的组合正是最美的组合，我指的是那些最能迷住所有数学家都知道的特殊感受性的组合；外行人对这种感受性如此无知，以致常常忍不住嘲笑它。\end{zhblock}

\begin{enblock}
What follows, then? Of the very large number of combinations which the subliminal ego blindly forms, almost all are without interest and without utility. But, for that very reason, they are without action on the aesthetic sensibility; the consciousness will never know them. A few only are harmonious, and consequently at once useful and beautiful, and they will be capable of affecting the geometrician's special sensibility I have been speaking of; which, once aroused, will direct our attention upon them, and will thus give them the opportunity of becoming conscious.
\end{enblock}
\begin{zhblock}
那么，接下来会怎样？在潜意识自我盲目形成的极大量组合中，几乎所有都没有兴趣，也没有用处。但正因为如此，它们不会作用于审美感受性；意识永远不会知道它们。只有少数组合是和谐的，因而同时有用而美；它们能够影响我一直在谈论的几何学家的特殊感受性。感受性一旦被唤起，就会把我们的注意力引向它们，从而给它们成为意识内容的机会。\end{zhblock}

\begin{enblock}
This is only a hypothesis, and yet there is an observation which tends to confirm it. When a sudden illumination invades the mathematician's mind, it most frequently happens that it does not mislead him. But it also happens sometimes, as I have said, that it will not stand the test of verification. Well, it is to be observed almost always that this false idea, if it had been correct, would have flattered our natural instinct for mathematical elegance.
\end{enblock}
\begin{zhblock}
这只是一个假设；然而有一项观察倾向于证实它。当突然照明闯入数学家的心灵时，最常见的情况是它并不误导他。但正如我已经说过的，有时它也经不起验证的考验。值得注意的是，几乎总是这样：这个错误想法如果是真的，本来会迎合我们对数学优雅的自然本能。\end{zhblock}

\begin{enblock}
Thus it is this special aesthetic sensibility that plays the part of the delicate sieve of which I spoke above, and this makes it sufficiently clear why the man who has it not will never be a real discoverer.
\end{enblock}
\begin{zhblock}
因此，正是这种特殊的审美感受性扮演了我上面所说的精细筛子的角色；这也足以说明，为什么没有这种感受性的人永远不会成为真正的发现者。\end{zhblock}

\begin{enblock}
All the difficulties, however, have not disappeared. The conscious ego is strictly limited, but as regards the subliminal ego, we do not know its limitations, and that is why we are not too loath to suppose that in a brief space of time it can form more different combinations than could be comprised in the whole life of a conscious being. These limitations do exist, however. Is it conceivable that it can form all the possible combinations, whose number staggers the imagination? Nevertheless this would seem to be necessary, for if it produces only a small portion of the combinations, and that by chance, there will be very small likelihood of the right one, the one that must be selected, being found among them.
\end{enblock}
\begin{zhblock}
然而，所有困难并没有消失。意识自我受到严格限制；至于潜意识自我，我们不知道它的限制在哪里，因此我们不太抗拒这样的设想：它能在很短时间内形成比一个有意识存在者一生所能包含的还要多的不同组合。不过，这些限制确实存在。它能形成所有可能组合吗？这些组合的数量足以使想象力震惊。可是，这似乎又是必要的，因为如果它只是偶然地产生组合中的一小部分，那么正确组合，即必须被选中的那个组合，出现在其中的可能性就非常小。\end{zhblock}

\begin{enblock}
Perhaps we must look for the explanation in that period of preliminary conscious work which always precedes all fruitful unconscious work. If I may be permitted a crude comparison, let us represent the future elements of our combinations as something resembling Epicurus's hooked atoms. When the mind is in complete repose these atoms are immovable; they are, so to speak, attached to the wall. This complete repose may continue indefinitely without the atoms meeting, and, consequently, without the possibility of the formation of any combination.
\end{enblock}
\begin{zhblock}
也许我们必须在那段初步的有意识工作中寻找解释；一切富有成果的无意识工作之前，总有这样一段工作。如果允许我作一个粗略的比喻，让我们把未来组合的元素想象成类似伊壁鸠鲁的带钩原子。当心智完全静止时，这些原子不动；可以说，它们附着在墙上。这种完全静止可以无限持续下去，原子彼此不会相遇，因此也不可能形成任何组合。\end{zhblock}

\begin{enblock}
On the other hand, during a period of apparent repose, but of unconscious work, some of them are detached from the wall and set in motion. They plow through space in all directions, like a swarm of gnats, for instance, or, if we prefer a more learned comparison, like the gaseous molecules in the kinetic theory of gases. Their mutual collisions may then produce new combinations.
\end{enblock}
\begin{zhblock}
另一方面，在一段表面静止、实则无意识工作的时期，其中一些原子脱离墙面并开始运动。它们向各个方向穿过空间，例如像一群蚊蚋，或者，如果我们喜欢更学术的比喻，就像气体运动论中的气体分子。它们相互碰撞，便可能产生新的组合。\end{zhblock}

\begin{enblock}
What is the part to be played by the preliminary conscious work? Clearly it is to liberate some of these atoms, to detach them from the wall and set them in motion. We think we have accomplished nothing, when we have stirred up the elements in a thousand different ways to try to arrange them, and have not succeeded in finding a satisfactory arrangement. But after this agitation imparted to them by our will, they do not return to their original repose, but continue to circulate freely.
\end{enblock}
\begin{zhblock}
初步的有意识工作要扮演什么角色？显然，它的作用是释放其中一些原子，使它们脱离墙面并开始运动。当我们以千百种方式搅动这些元素，试图安排它们，却没有成功找到令人满意的排列时，我们以为自己什么也没有完成。但是，在我们的意志给予它们这种扰动之后，它们并不回到原先的静止状态，而是继续自由循环。\end{zhblock}

\begin{enblock}
Now our will did not select them at random, but in pursuit of a perfectly definite aim. Those it has liberated are not, therefore, chance atoms; they are those from which we may reasonably expect the desired solution. The liberated atoms will then experience collisions, either with each other, or with the atoms that have remained stationary, which they will run against in their course. I apologize once more. My comparison is very crude, but I cannot well see how I could explain my thought in any other way.
\end{enblock}
\begin{zhblock}
而我们的意志并不是随机选择它们，而是追求一个完全确定的目标。因此，被它释放的并不是偶然原子；它们是我们可以合理期待由之得到所需解答的那些原子。被释放的原子随后会发生碰撞，或者彼此碰撞，或者与那些仍保持静止、却在它们运动途中被撞上的原子碰撞。我再次道歉。这个比喻十分粗糙，但我实在看不出还能怎样解释我的想法。\end{zhblock}

\begin{enblock}
However it be, the only combinations that have any chance of being formed are those in which one at least of the elements is one of the atoms deliberately selected by our will. Now it is evidently among these that what I called just now the right combination is to be found. Perhaps there is here a means of modifying what was paradoxical in the original hypothesis.
\end{enblock}
\begin{zhblock}
无论如何，唯一有机会形成的组合，是那些至少有一个元素属于我们的意志有意选择的原子的组合。而我刚才所谓的正确组合，显然就在这些组合之中。也许这里有一种办法，可以修正原先假设中的悖论之处。\end{zhblock}

\begin{enblock}
Yet another observation. It never happens that unconscious work supplies ready-made the result of a lengthy calculation in which we have only to apply fixed rules. It might be supposed that the subliminal ego, purely automatic as it is, was peculiarly fitted for this kind of work, which is, in a sense, exclusively mechanical. It would seem that, by thinking overnight of the factors of a multiplication sum, we might hope to find the product ready-made for us on waking; or, again, that an algebraic calculation, for instance, or a verification could be made unconsciously. Observation proves that such is by no means the case. All that we can hope from these inspirations, which are the fruits of unconscious work, is to obtain points of departure for such calculations. As for the calculations themselves, they must be made in the second period of conscious work which follows the inspiration, and in which the results of the inspiration are verified and the consequences deduced. The rules of these calculations are strict and complicated; they demand discipline, attention, will, and consequently consciousness. In the subliminal ego, on the contrary, there reigns what I would call liberty, if one could give this name to the mere absence of discipline and to disorder born of chance. Only, this very disorder permits unexpected couplings.
\end{enblock}
\begin{zhblock}
还有一个观察。无意识工作从不会把一项漫长计算的现成结果提供给我们，而这种计算只需应用固定规则。有人也许会设想，潜意识自我既然纯粹自动，就特别适合这种在某种意义上完全机械的工作。看起来，只要夜里想着一道乘法题的各个因数，我们醒来时就有希望发现乘积已经现成摆在眼前；或者，例如一个代数计算、一项验证，也可以无意识地完成。观察证明，情况绝非如此。我们能从这些无意识工作果实即灵感中期望得到的一切，只是这类计算的出发点。至于计算本身，必须在灵感之后的第二段有意识工作中完成；在这段工作中，灵感的结果得到验证，后果被推出。这些计算的规则严格而复杂；它们要求纪律、注意、意志，因而要求意识。相反，在潜意识自我中占支配地位的，是我愿称为自由的东西，假如可以把这个名称给予纪律的单纯缺席以及偶然产生的混乱的话。只是，这种混乱本身允许意外的结合。\end{zhblock}

\begin{enblock}
I will make one last remark. When I related above some personal observations, I spoke of a night of excitement, on which I worked as though in spite of myself. The cases of this are frequent, and it is not necessary that the abnormal cerebral activity should be caused by a physical stimulant, as in the case quoted. Well, it appears that, in these cases, we are ourselves assisting at our own unconscious work, which becomes partly perceptible to the overexcited consciousness, but does not on that account change its nature. We then become vaguely aware of what distinguishes the two mechanisms, or, if you will, of the methods of working of the two egos. The psychological observations I have thus succeeded in making appear to me, in their general characteristics, to confirm the views I have been enunciating.
\end{enblock}
\begin{zhblock}
我再作最后一点说明。上面叙述个人观察时，我谈到一个兴奋的夜晚，那天夜里我仿佛违背自己意愿似地工作。这类情况很常见，而且异常的大脑活动并不一定要像我引用的例子那样由物理刺激物引起。在这些情形中，我们似乎是在亲自旁观自己的无意识工作；这种工作对过度兴奋的意识变得部分可感知，但并不因此改变其性质。于是我们会模糊地意识到两种机制的区别，或者也可以说，意识自我与潜意识自我的两种工作方法的区别。我得以作出的这些心理学观察，在总体特征上看来证实了我一直阐述的观点。\end{zhblock}

\begin{enblock}
Truly there is great need of this, for in spite of everything they are and remain largely hypothetical. The interest of the question is so great that I do not regret having submitted them to the reader.
\end{enblock}
\begin{zhblock}
确实很需要这种证实，因为尽管如此，这些观点过去并且仍然在很大程度上是假设性的。这个问题的意义如此重大，以至于我并不后悔把它们提交给读者。\end{zhblock}

\part*{Eric R. Kandel, \emph{In Search of Memory: The Emergence of a New Science of Mind}\\埃里克·R. 坎德尔《追寻记忆：新心智科学的兴起》}

\section*{Chapter 4. One Cell at a Time\\第四章 一次研究一个细胞}

\begin{enblock}
I entered Harry Grundfest's laboratory at Columbia University for a six-month elective period in the fall of 1955, hoping to learn something about higher brain functions. I did not anticipate embarking on a new career, a new way of life. But my very first conversation with Grundfest gave me reason to reflect. In that conversation I described my interest in psychoanalysis and my hope of learning something about where in the brain the ego, the id, and the superego might be located.
\end{enblock}
\begin{zhblock}
1955年秋，我进入哥伦比亚大学哈里·格伦德费斯特的实验室，进行为期六个月的选修学习，希望了解一些关于高级脑功能的知识。我并没有预料到自己会由此开始一种新的职业、一种新的生活方式。但我与格伦德费斯特的第一次谈话就让我开始反思。那次谈话中，我描述了自己对精神分析的兴趣，以及希望弄清自我、本我和超我在大脑中可能位于何处。\end{zhblock}

\begin{enblock}
My desire to find these three psychic agencies had been sparked by a diagram Freud published in the course of summarizing his new structural theory of mind, which he developed in the decade 1923 to 1933 (figure 4-1). That new theory maintained his earlier distinction between conscious and unconscious mental functions, but it added three interacting psychic agencies: the ego, the id, and the superego. Freud saw consciousness as the surface of the mental apparatus. Much of our mental function is submerged below that surface, Freud argued, just as the bulk of an iceberg is submerged below the surface of the ocean. The deeper a mental function lies below the surface, the less accessible it is to consciousness. Psychoanalysis provided a way of digging down to the buried mental strata, the preconscious and the unconscious components of the personality.
\end{enblock}
\begin{zhblock}
我想寻找这三个心理机构的愿望，源自弗洛伊德在概述其心智结构新理论时发表的一幅图示；这一理论是他在1923年至1933年这十年间发展出来的（图4-1）。新理论保留了他早先关于意识和无意识心理功能的区分，但又增加了三个相互作用的心理机构：自我、本我和超我。弗洛伊德把意识看作心理装置的表面。他认为，我们大部分心理功能都沉没在这个表面之下，正如冰山的大部分沉没在海面之下。心理功能位于表面之下越深，就越不容易被意识接近。精神分析提供了一种方式，可以向下挖掘到被埋藏的心理地层，即人格中的前意识和无意识成分。\end{zhblock}

\begin{figure}[htbp]
\centering
\safeincludeimage{a9b9ebf53506e9e0115c3e0ac3176d4a0e011804b72ed5e036c409bbfb033975.jpg}
\caption{Freud's structural theory. Freud conceived of three main psychic structures--the ego, the id, and the superego. The ego has a conscious component (perceptual consciousness, or pcpt.-cs.) that receives sensory information and is in direct contact with the outside world, as well as a preconscious component, an aspect of unconscious processing that has ready access to consciousness. The ego's unconscious components act through repression and other defenses to inhibit the instinctual urges of the id, the generator of sexual and aggressive instincts. The ego also responds to the pressures of the superego, the largely unconscious carrier of moral values. The dotted lines indicate the divisions between those processes that are accessible to consciousness and those that are completely unconscious. (From \emph{New Introductory Lectures on Psychoanalysis} [1933]).\\
弗洛伊德的结构理论。弗洛伊德设想了三个主要心理结构：自我、本我和超我。自我有一个意识成分（知觉意识，pcpt.-cs.），它接收感觉信息并与外部世界直接接触；自我还有一个前意识成分，即无意识加工中可较容易进入意识的一个方面。自我的无意识成分通过压抑及其他防御，抑制本我的本能冲动；本我是性本能和攻击本能的发生者。自我也回应超我的压力；超我是道德价值的、在很大程度上无意识的承载者。虚线标示了可进入意识的过程与完全无意识过程之间的分界。（引自《精神分析新引论》[1933]。）}
\end{figure}

\begin{enblock}
What gave Freud's new model a dramatic turn was the three interacting psychic agencies. Freud did not define the ego, the id, and the superego as either conscious or unconscious, but as differing in cognitive style, goal, and function.
\end{enblock}
\begin{zhblock}
使弗洛伊德新模型产生戏剧性转折的，是这三个相互作用的心理机构。弗洛伊德并没有把自我、本我和超我界定为意识的或无意识的，而是把它们界定为在认知风格、目标和功能上彼此不同。\end{zhblock}

\begin{enblock}
According to Freud's structural theory, the ego (the “I,” or autobiographical self) is the executive agency, and it has both a conscious and an unconscious component. The conscious component is in direct contact with the external world through the sensory apparatus for sight, sound, and touch; it is concerned with perception, reasoning, the planning of action, and the experiencing of pleasure and pain. In their work, Hartmann, Kris, and Lowenstein emphasized that this conflict-free component of the ego operates logically and is guided in its actions by the reality principle. The unconscious component of the ego is concerned with psychological defenses (repression, denial, sublimation), the mechanisms whereby the ego inhibits, channels, and redirects both the sexual and the aggressive instinctual drives of the id, the second psychic agency.
\end{enblock}
\begin{zhblock}
根据弗洛伊德的结构理论，自我（“我”，或自传式自我）是执行机构，并且同时具有意识成分和无意识成分。意识成分通过视觉、听觉和触觉的感觉装置与外部世界直接接触；它涉及知觉、推理、行动计划以及愉悦和痛苦的体验。哈特曼、克里斯和勒文斯坦在其工作中强调，自我这一没有冲突的成分以逻辑方式运行，并在行动中受现实原则引导。自我的无意识成分涉及心理防御（压抑、否认、升华），即自我用来抑制、疏导并重新定向第二个心理机构本我的性本能驱力和攻击本能驱力的机制。\end{zhblock}

\begin{enblock}
The id (the “it”, a term that Freud borrowed from Friedrich Nietzsche, is totally unconscious. It is not governed by logic or by reality but by the hedonistic principle of seeking pleasure and avoiding pain. The id, according to Freud, represents the primitive mind of the infant and is the only mental structure present at birth. The superego, the third governor, is the unconscious moral agency, the embodiment of our aspirations.
\end{enblock}
\begin{zhblock}
本我（“它”）这个术语，是弗洛伊德从弗里德里希·尼采那里借来的；本我是完全无意识的。它不受逻辑或现实支配，而受寻求快乐、避免痛苦的享乐原则支配。按照弗洛伊德的说法，本我代表婴儿的原始心智，是出生时唯一存在的心理结构。第三个统辖者超我，则是无意识的道德机构，是我们理想追求的体现。\end{zhblock}

\begin{enblock}
Although Freud did not intend his diagram to be a neuroanatomical map of mind, it stimulated me to wonder where in the elaborate folds of the human brain these psychic agencies might live, as it had earlier stimulated the curiosity of Kubie and Ostow. As I mentioned, these two psychoanalysts with a keen interest in biology had encouraged me to study with Grundfest.
\end{enblock}
\begin{zhblock}
尽管弗洛伊德并不打算把这幅图当作心智的神经解剖图，它仍激发我去想象，这些心理机构可能居住在人脑复杂褶皱中的什么地方；此前，它也曾激发库比和奥斯托的好奇心。正如我提到过的，这两位对生物学有浓厚兴趣的精神分析学家曾鼓励我跟随格伦德费斯特学习。\end{zhblock}

\begin{enblock}
Grundfest listened patiently as I told him of my rather grandiose ideas. Another biologist might well have dismissed me, wondering what to do with this naive and misguided medical student. But not Grundfest. He explained that my hope of understanding the biological basis of Freud's structural theory of mind was far beyond the grasp of contemporary brain science. Rather, he told me, to understand mind we needed to look at the brain one cell at a time.
\end{enblock}
\begin{zhblock}
我向格伦德费斯特讲述这些相当宏大的想法时，他耐心地听着。换作另一位生物学家，很可能已经把我打发走了，不知道该拿这个天真而误入歧途的医学生怎么办。但格伦德费斯特没有。他解释说，我想理解弗洛伊德心智结构理论的生物学基础，这远远超出了当时脑科学的能力范围。相反，他告诉我，要理解心智，我们需要一次研究大脑中的一个细胞。\end{zhblock}

\begin{enblock}
One cell at a time! I initially found those words demoralizing. How could one address psychoanalytic questions about the unconscious motivation of behavior, or the action of our conscious life, by studying the brain on the level of single nerve cells? But as we talked I suddenly remembered that in 1887, when Freud began his own career, he had sought to solve the hidden riddles of mental life by studying the brain one nerve cell at a time. Freud started out as an anatomist, studying single nerve cells, and had anticipated a key point of what later came to be called the neuron doctrine, the view that nerve cells are the building blocks of the brain. It was only later, after he began treating mentally ill patients in Vienna, that Freud made his monumental discoveries about unconscious mental processes.
\end{enblock}
\begin{zhblock}
一次研究一个细胞！起初我觉得这句话令人沮丧。通过在单个神经细胞层面研究大脑，怎么可能处理精神分析关于行为的无意识动机，或者关于我们意识生活作用的问题呢？但在谈话中，我突然想起，1887年弗洛伊德开始自己的职业生涯时，他正是试图通过一次研究一个神经细胞来解开心理生活的隐秘谜题。弗洛伊德最初是一名解剖学家，研究单个神经细胞；他还预见了后来所谓神经元学说的一个要点，即神经细胞是大脑的构件。只是后来，在维也纳开始治疗精神疾病患者之后，弗洛伊德才作出了关于无意识心理过程的重大并具有里程碑意义的发现。\end{zhblock}

\begin{enblock}
I found it ironic and remarkable that I was now being encouraged to take that journey in reverse, to move from an interest in the top-down structural theory of mind to the bottom-up study of the signaling elements of the nervous system, the intricate inner worlds of nerve cells. Harry Grundfest offered to guide me into this new world.
\end{enblock}
\begin{zhblock}
我觉得既讽刺又非凡的是，现在有人鼓励我反向踏上这段旅程：从对自上而下的心智结构理论的兴趣，转向自下而上地研究神经系统的信号元件，研究神经细胞复杂的内部世界。哈里·格伦德费斯特愿意引导我进入这个新世界。\end{zhblock}

\begin{center}
[\dots]
\end{center}

\section*{Chapter 28. Consciousness\\第二十八章 意识}

\begin{enblock}
Psychoanalysis introduced us to the unconscious in its several forms. Like many scientists now working on the brain, I have long been intrigued by the biggest question about the brain: the nature of consciousness and how various unconscious psychological processes relate to conscious thought. When I first talked with Harry Grundfest about Freud's structural theory of mind--the ego, the id, and the superego--the central focus of my thinking was: How do conscious and unconscious processes differ in their representation in the brain? But only recently has the new science of mind developed the tools for exploring this question experimentally.
\end{enblock}
\begin{zhblock}
精神分析使我们认识了无意识的多种形式。像今天许多研究大脑的科学家一样，我长期以来一直被关于大脑的最大问题所吸引：意识的本性，以及各种无意识心理过程如何与意识思维相关。当我第一次与哈里·格伦德费斯特谈论弗洛伊德的心智结构理论，即自我、本我和超我时，我思考的中心焦点是：意识过程和无意识过程在大脑中的表征有什么不同？但直到最近，新的心智科学才发展出实验探索这个问题的工具。\end{zhblock}

\begin{enblock}
To develop productive insights into consciousness, the new science of mind first had to settle on a working definition of consciousness as a state of perceptual awareness, or selective attention writ large. At its core, consciousness in people is an awareness of self, an awareness of being aware. Consciousness thus refers to our ability not simply to experience pleasure and pain but to attend to and reflect upon those experiences, and to do so in the context of our immediate lives and our life history. Conscious attention allows us to shut out extraneous experiences and focus on the critical event that confronts us, be it pleasure or pain, the blue of the sky, the cool northern light of a Vermeer painting, or the beauty and calm we experience at the seashore.
\end{enblock}
\begin{zhblock}
为了对意识形成有成果的洞见，新的心智科学首先必须确定一个工作性定义：意识是一种知觉觉知状态，或者说是放大了的选择性注意。就其核心而言，人的意识是自我觉知，是对“自己正在觉知”的觉知。因此，意识指的不只是我们体验愉悦和痛苦的能力，还包括我们关注并反思这些体验的能力，并且是在我们当下生活和生命史的语境中这样做的能力。意识注意使我们能够排除无关经验，聚焦于面对我们的关键事件，无论它是愉悦还是痛苦，是天空的蓝色、维米尔画作中清冷的北方光线，还是我们在海边体验到的美与宁静。\end{zhblock}

\begin{enblock}
Understanding consciousness is by far the most challenging task confronting science. The truth of this assertion can best be seen in the career of Francis Crick, perhaps the most creative and influential biologist of the second half of the twentieth century. When Crick first entered biology, after World War II, two great questions were thought to be beyond the capacities of science to answer: What distinguishes the living from the nonliving world? And what is the biological nature of consciousness? Crick turned first to the easier problem, distinguishing animate from inanimate matter, and explored the nature of the gene. By 1953, after just two years of collaboration, he and Jim Watson had helped solve that mystery. As Watson later described in \emph{The Double Helix}, “at lunch Francis winged into the Eagle [Pub] to tell everyone within hearing distance that we had found the secret of life.” In the next two decades, Crick helped crack the genetic code: how DNA makes RNA and RNA makes protein.
\end{enblock}
\begin{zhblock}
理解意识是迄今为止科学所面临的最具挑战性的任务。这一断言的真实性，最好可以从弗朗西斯·克里克的职业生涯中看出；他也许是二十世纪下半叶最富创造力、最有影响力的生物学家。第二次世界大战后，克里克刚进入生物学领域时，人们认为有两个大问题超出了科学能够回答的范围：什么把有生命世界与无生命世界区分开来？意识的生物学本性是什么？克里克首先转向较容易的问题，即区分有生命物质和无生命物质，并探索基因的本性。到1953年，经过仅仅两年的合作，他和吉姆·沃森已经帮助解开了这个谜。正如沃森后来在《双螺旋》中描述的：“午餐时，弗朗西斯飞奔进鹰酒馆，告诉所有听得见的人，我们找到了生命的秘密。”在接下来的二十年里，克里克帮助破解了遗传密码：DNA如何制造RNA，RNA如何制造蛋白质。\end{zhblock}

\begin{enblock}
In 1976, at age sixty, Crick turned to the remaining scientific mystery: the biological nature of consciousness. This he studied for the rest of his life in partnership with Christof Koch, a young computational neuroscientist. Crick brought his characteristic intelligence and optimism to bear on the question; moreover, he made consciousness a focus of the scientific community, which had previously ignored it. But, despite almost thirty years of continuous effort, Crick was able to budge the problem only a modest distance. Indeed, some scientists and philosophers of mind continue to find consciousness so inscrutable that they fear it can never be explained in physical terms. How can a biological system, a biological machine, they ask, feel anything? Even more doubtful, how can it think about itself?
\end{enblock}
\begin{zhblock}
1976年，六十岁的克里克转向剩下的那个科学谜题：意识的生物学本性。此后余生，他与年轻的计算神经科学家克里斯托夫·科赫合作研究这一问题。克里克把他特有的智慧和乐观精神用于这个问题；此外，他使意识成为科学共同体关注的焦点，而此前科学界一直忽视它。但是，尽管持续努力了近三十年，克里克也只把这个问题向前推进了一小段距离。确实，一些科学家和心灵哲学家仍然觉得意识如此不可理解，以至于担心它永远无法用物理术语解释。他们问：一个生物系统、一台生物机器，怎么会有任何感觉？更可疑的是，它怎么会思考自身？
\end{zhblock}

\begin{enblock}
These questions are not new. They were first posed in Western thought during the fifth century B.C. by Hippocrates and by the philosopher Plato, the founder of the Academy in Athens. Hippocrates was the first physician to cast superstition aside, basing his thinking on clinical observations and arguing that all mental processes emanate from the brain. Plato, who rejected observations and experiments, believed that the only reason we can think about ourselves and our mortal body is that we have a soul that is immaterial and immortal. The idea of an immortal soul was subsequently incorporated into Christian thought and elaborated upon by St. Thomas Aquinas in the thirteenth century. Aquinas and later religious thinkers held that the soul--the generator of consciousness--is not only distinct from the body, it is also of divine origin.
\end{enblock}
\begin{zhblock}
这些问题并不新。在西方思想中，它们最早由公元前五世纪的希波克拉底和哲学家柏拉图提出；柏拉图是雅典学园的创立者。希波克拉底是第一个把迷信抛在一边的医生，他把自己的思考建立在临床观察之上，并主张一切心理过程都源自大脑。拒绝观察和实验的柏拉图则认为，我们之所以能够思考自己以及自身会死的身体，唯一原因是我们有一个非物质且不朽的灵魂。不朽灵魂的观念随后被纳入基督教思想，并在十三世纪由圣托马斯·阿奎那加以阐发。阿奎那以及后来的宗教思想家认为，灵魂这个意识的产生者不仅不同于身体，而且具有神圣起源。\end{zhblock}

\begin{enblock}
In the seventeenth century, René Descartes developed the idea that human beings have a dual nature: they have a body, which is made up of material substance, and a mind, which derives from the spiritual nature of the soul. The soul receives signals from the body and can influence its actions but is itself made up of an immaterial substance that is unique to human beings. Descartes' thinking gave rise to the view that actions like eating and walking, as well as sensory perception, appetites, passions, and even simple forms of learning, are all mediated by the brain and can be studied scientifically. Mind, however, is sacred and as such is not a proper subject of science.
\end{enblock}
\begin{zhblock}
十七世纪，勒内·笛卡尔发展出这样一种观念：人具有双重本性。人有身体，由物质实体构成；也有心灵，源自灵魂的精神本性。灵魂接收来自身体的信号，并能影响身体的行动，但灵魂本身由一种人类所独有的非物质实体构成。笛卡尔的思想产生了一种看法，即吃饭、行走等行动，以及感觉知觉、欲望、激情，甚至简单形式的学习，都是由大脑中介的，并且可以被科学研究。然而心灵是神圣的，因此并不是科学的适当对象。\end{zhblock}

\begin{enblock}
It is remarkable to reflect that these seventeenth-century ideas were still current in the 1980s. Karl Popper, the Vienna-born philosopher of science, and John Eccles, the Nobel laureate neurobiologist, espoused dualism all of their lives. They agreed with Aquinas that the soul is immortal and independent of the brain. Gilbert Ryle, the British philosopher of science, referred to the notion of the soul as “the ghost in the machine.”\end{enblock}
\begin{zhblock}
令人惊讶的是，回想起来，这些十七世纪的观念在二十世纪八十年代仍然流行。出生于维也纳的科学哲学家卡尔·波普尔和诺贝尔奖得主、神经生物学家约翰·埃克尔斯，一生都拥护二元论。他们同意阿奎那的观点，认为灵魂不朽并且独立于大脑。英国科学哲学家吉尔伯特·赖尔把灵魂观念称为“机器中的幽灵”。\end{zhblock}

\begin{enblock}
Today, most philosophers of mind agree that what we call consciousness derives from the physical brain, but some disagree with Crick as to whether it can ever be approached scientifically. A few, such as Colin McGinn, believe that consciousness simply cannot be studied, because the architecture of the brain poses limitations on human cognitive capacities. In McGinn's view, the human mind may simply be incapable of solving certain problems. At the other extreme, philosophers such as Daniel Dennett deny that there is any problem at all. Dennett argues, much as neurologist John Hughlings Jackson did a century earlier, that consciousness is not a distinct operation of the brain; rather, it is the combined result of the computational workings of higher-order areas of the brain concerned with later stages of information processing.
\end{enblock}
\begin{zhblock}
今天，多数心灵哲学家都同意，我们所谓的意识来自物理大脑；但在意识是否终究能够以科学方式接近这一点上，有些人不同意克里克。少数人如科林·麦金认为，意识根本无法被研究，因为大脑的结构限制了人类认知能力。在麦金看来，人类心智也许就是不能解决某些问题。另一个极端上，丹尼尔·丹尼特等哲学家则否认这里有任何问题。丹尼特的论证与一个世纪前神经学家约翰·休林斯·杰克逊的论证十分相似：意识并不是大脑的一种独立运作；相反，它是大脑中与信息加工后期阶段有关的高级区域计算活动的综合结果。\end{zhblock}

\begin{enblock}
Finally, philosophers such as John Searle and Thomas Nagel take a middle position, holding that consciousness is a discrete set of biological processes. The processes are accessible to analysis, but we have made little headway in understanding them because they are very complex and represent more than the sum of their parts. Consciousness is therefore much more complicated than any property of the brain that we understand.
\end{enblock}
\begin{zhblock}
最后，约翰·塞尔和托马斯·内格尔等哲学家采取中间立场，认为意识是一组离散的生物过程。这些过程可以被分析，但由于它们非常复杂，并且不只是各部分之和，我们在理解它们方面进展很小。因此，意识比我们所理解的大脑任何性质都复杂得多。\end{zhblock}

\begin{enblock}
Searle and Nagel ascribe two characteristics to the conscious state: unity and subjectivity. The unitary nature of consciousness refers to the fact that our experiences come to us as a unified whole. All of the various sensory modalities are melded into a single, coherent, conscious experience. Thus when I approach a rosebush in the botanical garden at Wave Hill near my house in Riverdale, I sniff the exquisite fragrance of the blossoms at the same time that I see their beautiful red color--and I perceive this rosebush against the background of the Hudson River and the cliffs of the Palisade mountain ridge behind it. My perception is not only whole during the moment I experience it, it is also whole two weeks later, when I engage in mental time travel to recapture the moment. Despite the fact that there are different organs for smell and vision, and that each uses its own individual pathways, they converge in the brain in such a way that my perceptions are unified.
\end{enblock}
\begin{zhblock}
塞尔和内格尔把两个特征归于意识状态：统一性和主观性。意识的统一性是指，我们的经验以统一整体的形式来到我们这里。各种感觉通道都融合为单一、连贯的意识经验。因此，当我走近河谷区我家附近波浪山植物园里的一丛玫瑰时，我一边闻到花朵精致的芳香，一边看到它们美丽的红色；而且我是在以哈德逊河以及其后帕利塞兹山脊的悬崖为背景时知觉到这丛玫瑰的。我的知觉不仅在我体验它的那个瞬间是完整的；两周以后，当我进行心理时间旅行以重新捕捉那个瞬间时，它仍然是完整的。尽管嗅觉和视觉有不同器官，而且各自使用独立通路，它们仍以某种方式在大脑中汇合，使我的知觉成为统一的。\end{zhblock}

\begin{enblock}
The unitary nature of consciousness poses a difficult problem, but perhaps not an insurmountable one. This unitary nature can break down. In a surgical patient whose brain is severed between the two hemispheres, there are two conscious minds, each with its own unified percept.
\end{enblock}
\begin{zhblock}
意识的统一性提出了一个困难问题，但也许不是不可克服的问题。这种统一性可能破裂。在一名两半球之间连接被外科切断的患者那里，会有两个意识心灵，每一个都有自己的统一知觉。\end{zhblock}

\begin{enblock}
Subjectivity, the second characteristic of conscious awareness, poses the more formidable scientific challenge. Each of us experiences a world of private and unique sensations that is much more real to us than the experiences of others. We experience our own ideas, moods, and sensations directly, whereas we can only appreciate another person's experience indirectly, by observing or hearing about it. We therefore can ask, Is your response to the blue you see and the jasmine you smell--the meaning it has for you--identical to my response to the blue I see and the jasmine I smell and the meaning these have for me?
\end{enblock}
\begin{zhblock}
主观性是意识觉知的第二个特征，它提出了更艰巨的科学挑战。我们每个人都体验一个私人而独特的感觉世界；这个世界对我们来说比他人的经验真实得多。我们直接体验自己的观念、情绪和感觉，而只能通过观察或听别人讲述来间接理解另一个人的经验。因此我们可以问：你对你所看见的蓝色和所闻到的茉莉的反应，即它对你的意义，是否与我对我所看见的蓝色和所闻到的茉莉的反应以及它们对我的意义完全相同？
\end{zhblock}

\begin{enblock}
The issue here is not one of perception per se. It is not whether we each see a very similar shade of the same blue. That is relatively easy to establish by recording from single nerve cells in the visual system of different individuals. The brain does reconstruct our perception of an object, but the object perceived--the color blue or middle C on the piano--appears to correspond to the physical properties of the wavelength of the reflected light or the frequency of the emitted sound. Instead, the issue is the significance of that blue and that note for each of us. What we do not understand is how electrical activity in neurons gives rise to the meaning we ascribe to that color or that wavelength of sound. The fact that conscious experience is unique to each person raises the question of whether it is possible to determine objectively any characteristics of consciousness that are common to everyone. If the senses ultimately produce experiences that are completely and personally subjective, we cannot, the argument goes, arrive at a general definition of consciousness based on personal experience.
\end{enblock}
\begin{zhblock}
这里的问题并不是知觉本身的问题。问题不在于我们是否各自看见了同一种蓝色中非常相近的色调。通过记录不同个体视觉系统中单个神经细胞的活动，这一点相对容易确定。大脑确实会重构我们对一个对象的知觉，但被知觉的对象，即蓝色或钢琴上的中央C，似乎对应于反射光波长或发出声音频率的物理性质。相反，问题在于那种蓝色和那个音符对我们每个人的意义。我们不理解的是，神经元中的电活动如何产生我们赋予那种颜色或那一声波波长的意义。意识经验对每个人都是独特的这一事实提出了一个问题：是否可能客观地确定某些对每个人都共同的意识特征？如果感觉最终产生的经验是完全个人化、完全主观的，按照这种论证，我们就不能基于个人经验得出意识的一般定义。\end{zhblock}

\begin{enblock}
Nagel and Searle illustrate the difficulty of explaining the subjective nature of consciousness in physical terms as follows: Assume we succeed in recording the electrical activity of neurons in a region known to be important for consciousness while the person being studied carries out some task that requires conscious attention. For example, suppose we identified the cells that fire when I look at and become aware of a red image of the blossoms on a rosebush at Wave Hill. We have now taken a first step in studying consciousness--namely, we have found what Crick and Koch have called the neural correlate of consciousness for this one percept. For most of us, this would be a great advance because it pinpoints a material concomitant of conscious perception. From there we could go on to carry out experiments to determine whether these correlates also meld into a coherent whole, that is, the background of the Hudson River and the Palisades. But for Nagel and Searle, this is the easy problem of consciousness. The hard problem of consciousness is the second mystery, that of subjective experience.
\end{enblock}
\begin{zhblock}
内格尔和塞尔用如下方式说明以物理术语解释意识主观本性的困难：假定在被研究者执行某项需要意识注意的任务时，我们成功记录了已知对意识重要的某个脑区中神经元的电活动。例如，设想我看见并觉知到波浪山玫瑰丛花朵的红色图像时，我们识别出了放电的细胞。我们现在就迈出了研究意识的第一步，也就是找到了克里克和科赫所说的这一知觉的意识神经相关物。对我们大多数人来说，这会是一项巨大进展，因为它精确指出了意识知觉的一个物质伴随物。由此出发，我们可以继续做实验，确定这些相关物是否也融合为一个连贯整体，即哈德逊河和帕利塞兹的背景。但对内格尔和塞尔来说，这是意识的容易问题。意识的困难问题是第二个谜，即主观经验之谜。\end{zhblock}

\begin{enblock}
How is it that I respond to the red image of a rose with a feeling that is distinctive to me? To use another example, what grounds do we have for believing that when a mother looks at her child, the firing of cells in the region of the cortex concerned with face recognition accounts for the emotions she feels and for her ability to summon the memory of those emotions and that image of her child?
\end{enblock}
\begin{zhblock}
为什么我会以一种对我而言独特的感觉来回应玫瑰的红色图像？再举一个例子，我们有什么根据相信，当一位母亲看着自己的孩子时，皮层中与面孔识别有关区域的细胞放电，就能解释她感受到的情感，以及她唤起这些情感和孩子形象记忆的能力？\end{zhblock}

\begin{enblock}
As yet, we do not know how the firing of specific neurons leads to the subjective component of conscious perception, even in the simplest case. In fact, according to Searle and Nagel, we lack an adequate theory of how an objective phenomenon, such as electrical signals in the brain, can cause a subjective experience, such as pain. And because science as we currently practice it is a reductionist, analytical view of complicated events, while consciousness is irreducibly subjective, such a theory lies beyond our reach for now.
\end{enblock}
\begin{zhblock}
到目前为止，即便在最简单的情形中，我们也不知道特定神经元的放电如何导致意识知觉的主观成分。事实上，按照塞尔和内格尔的说法，我们缺少一种充分理论来解释，一个客观现象，例如大脑中的电信号，如何能够造成一种主观经验，例如疼痛。而且，由于我们当前实践的科学是一种还原论的、分析性的复杂事件观，而意识又不可还原地具有主观性，这样一种理论目前仍超出我们的能力范围。\end{zhblock}

\begin{enblock}
According to Nagel, science cannot take on consciousness without a significant change in methodology, a change that would enable scientists to identify and analyze the elements of subjective experience. Those elements are likely to be basic components of brain function, much as atoms and molecules are basic components of matter, but to exist in a form we cannot yet imagine. The reductions performed routinely in science are not problematic, Nagel holds. Biological science can readily explain how the properties of a particular type of matter arise from the objective properties of the molecules of which it is made. What science lacks are rules for explaining how subjective properties (consciousness) arise from the properties of objects (interconnected nerve cells).
\end{enblock}
\begin{zhblock}
内格尔认为，如果方法论没有重大改变，科学就不能处理意识；这种改变必须使科学家能够识别并分析主观经验的元素。这些元素很可能是大脑功能的基本成分，正如原子和分子是物质的基本成分一样；但它们以一种我们目前还无法想象的形式存在。内格尔认为，科学中常规进行的还原并没有问题。生物科学可以很容易地解释某一特定类型物质的性质如何来自构成它的分子的客观性质。科学所缺少的是一些规则，用来解释主观性质（意识）如何来自对象（相互连接的神经细胞）的性质。\end{zhblock}

\begin{enblock}
Nagel argues that our complete lack of insight into the elements of subjective experience should not prevent us from discovering the neural correlates of consciousness and the rules that relate conscious phenomena to cellular processes in the brain. In fact, it is only by accumulating such information that we will be in a position to think about the reduction of something subjective to something physical and objective. But to arrive at a theory that supports this reduction, we will first have to discover the elements of subjective consciousness. This discovery, says Nagel, will be enormous in its magnitude and its implications, requiring a revolution in biology and most likely a complete transformation of scientific thought.
\end{enblock}
\begin{zhblock}
内格尔主张，尽管我们对主观经验的元素完全缺乏洞见，这不应妨碍我们发现意识的神经相关物，以及把意识现象同大脑中的细胞过程联系起来的规则。事实上，只有通过积累这类信息，我们才会处在能够思考如何把某种主观的东西还原为某种物理而客观的东西的位置上。但是，为了得到一种支持这种还原的理论，我们首先必须发现主观意识的元素。内格尔说，这一发现的规模及其含义将极其巨大，它需要生物学发生一场革命，并且很可能需要科学思想的彻底转变。\end{zhblock}

\begin{enblock}
The aim of most neural scientists working on consciousness is much more modest than this grand perspective would imply. They are not deliberately working toward or anticipating a revolution in scientific thought. Although they must struggle with the difficulties of defining conscious phenomena experimentally, they do not see those difficulties as precluding all experimental study under existing paradigms. Neural scientists believe, and Searle for one agrees with them, that they have been able to make considerable progress in understanding the neurobiology of perception and memory without having to account for individual experience. For example, cognitive neural scientists have made advances in understanding the neural basis of the perception of the color blue without addressing the question of how each of us responds to the same blue.
\end{enblock}
\begin{zhblock}
大多数研究意识的神经科学家的目标，比这种宏大图景所暗示的要谦逊得多。他们并没有有意朝着科学思想革命工作，也没有预期这样的革命。虽然他们必须努力应对以实验方式界定意识现象的困难，但他们并不认为这些困难会排除在现有范式下进行一切实验研究。神经科学家相信，塞尔本人也同意他们，他们已经能够在不解释个人经验的情况下，在理解知觉和记忆的神经生物学方面取得相当进展。例如，认知神经科学家在理解蓝色知觉的神经基础方面取得了进展，而没有处理我们每个人如何回应同一种蓝色的问题。\end{zhblock}

\begin{enblock}
What we do not understand is the hard problem of consciousness--the mystery of how neural activity gives rise to subjective experience. Crick and Koch have argued that once we solve the easy problem of consciousness, the unity of consciousness, we will be able to manipulate those neural systems experimentally to solve the hard problem.
\end{enblock}
\begin{zhblock}
我们所不理解的，是意识的困难问题，即神经活动如何产生主观经验之谜。克里克和科赫认为，一旦我们解决了意识的容易问题，也就是意识的统一性，我们就能够通过实验操纵那些神经系统，从而解决困难问题。\end{zhblock}

\begin{enblock}
The unity of consciousness is a variant of the binding problem first identified in the study of visual perception. An intimate part of my experiencing the subjective pleasure of the moment at Wave Hill is how the look and the smell of roses in the gardens are bound together and unified with my view of the Hudson, the Palisades, and all the other component images of my perception. Each of these components of my subjective experience is mediated by different brain regions within my visual, olfactory, and emotional systems. The unity of my conscious experience implies that the binding process must somehow connect and integrate all of these separate areas in the brain.
\end{enblock}
\begin{zhblock}
意识的统一性是最初在视觉知觉研究中被识别出的绑定问题的一种变体。我在波浪山那个时刻体验到主观愉悦，其中一个亲密组成部分，就是花园中玫瑰的外观和气味如何与我对哈德逊河、帕利塞兹以及我知觉中所有其他组成图像的观看绑定在一起并统一起来。我的主观经验中每一个组成部分，都是由视觉、嗅觉和情感系统中的不同脑区中介的。我的意识经验的统一性意味着，绑定过程必须以某种方式把大脑中所有这些分离区域连接并整合起来。\end{zhblock}

\begin{enblock}
As a first step toward solving the easy problem of consciousness, we need to ask whether the unity of consciousness--a unity thought to be achieved by neural systems that mediate selective attention--is localized in one or just a few sites, which would enable us to manipulate them biologically. The answer to this question is by no means clear. Gerald Edelman, a leading theoretician on the brain and consciousness, has argued effectively that the neural machinery for the unity of consciousness is likely to be widely distributed throughout the cortex and thalamus. As a result, Edelman asserts, it is unlikely that we will be able to find consciousness through a simple set of neural correlates. Crick and Koch, on the other hand, believe that the unity of consciousness will have direct neural correlates because they most likely involve a specific set of neurons with specific molecular or neuroanatomical signatures. The neural correlates, they argue, probably require only a small set of neurons acting as a searchlight: the spotlight of attention. The initial task, they argue, is to locate within the brain that small set of neurons whose activity correlates best with the unity of conscious experience and then to determine the neural circuits to which they belong.
\end{enblock}
\begin{zhblock}
作为解决意识容易问题的第一步，我们需要问：意识的统一性，即被认为由中介选择性注意的神经系统实现的统一性，是否定位在一个或少数几个部位，从而使我们能够以生物学方式操纵它们？这个问题的答案远不清楚。杰拉尔德·埃德尔曼是大脑和意识领域的重要理论家，他有力地论证说，意识统一性的神经机制很可能广泛分布于整个皮层和丘脑。因此，埃德尔曼断言，我们不大可能通过一组简单的神经相关物找到意识。另一方面，克里克和科赫相信，意识的统一性会有直接的神经相关物，因为它们很可能涉及一组具有特定分子或神经解剖标志的神经元。他们认为，这些神经相关物也许只需要一小组像探照灯一样起作用的神经元：注意的聚光灯。他们认为，最初的任务是在大脑中定位那一小组神经元，其活动与意识经验的统一性最为相关，然后确定它们所属的神经回路。\end{zhblock}

\begin{enblock}
How are we to find this small population of nerve cells that could mediate the unity of consciousness? What criteria must they meet? In Crick and Koch's last paper (which Crick was still correcting on his way to the hospital a few hours before he died, on July 28, 2004), they focused on the claustrum, a sheet of brain tissue that is located below the cerebral cortex, as the site that mediates unity of experience. Little is known about the claustrum except that it connects to and exchanges information with almost all of the sensory and motor regions of the cortex as well as the amygdala, which plays an important role in emotion. Crick and Koch compare the claustrum to the conductor of an orchestra. Indeed, the neuroanatomical connections of the claustrum meet the requirements of a conductor; it can bind together and coordinate the various brain regions necessary for the unity of conscious awareness.
\end{enblock}
\begin{zhblock}
我们怎样找到这一小群可能中介意识统一性的神经细胞？它们必须满足什么标准？在克里克和科赫的最后一篇论文中（克里克在2004年7月28日去世前几小时、去医院的路上仍在修改这篇论文），他们把注意力集中在屏状核上。屏状核是位于大脑皮层下方的一片脑组织，被他们视为中介经验统一性的部位。关于屏状核，人们所知甚少；只知道它与皮层几乎所有感觉和运动区域以及杏仁核相连接并交换信息，而杏仁核在情绪中发挥重要作用。克里克和科赫把屏状核比作乐队指挥。确实，屏状核的神经解剖连接符合指挥的要求；它能够把意识觉知统一性所需的各个脑区绑定并协调起来。\end{zhblock}

\begin{enblock}
The idea that obsessed Crick at the end of his life--that the claustrum is the spotlight of attention, the site that binds the various components of any percept together--is the last in a series of important ideas he advanced. Crick's enormous contributions to biology (the double helical structure of DNA, the nature of the genetic code, the discovery of messenger RNA, the mechanisms of translating messenger RNA into the amino acid sequence of a protein, and the legitimizing of the biology of consciousness) put him in a class with Copernicus, Newton, Darwin, and Einstein. Yet his intense, lifelong focus on science, on the life of mind, is something he shares with many in the scientific community, and that obsession is symbolic of science at its best. The cognitive psychologist Vilayanur Ramachandran, a friend and colleague of Crick's, described Crick's focus on the claustrum during his last weeks:
\end{enblock}
\begin{zhblock}
克里克生命末期一直着迷的那个想法，即屏状核是注意的聚光灯，是把任何知觉的各个组成部分绑定在一起的部位，是他提出的一系列重要思想中的最后一个。克里克对生物学的巨大贡献，包括DNA双螺旋结构、遗传密码的本性、信使RNA的发现、把信使RNA翻译成蛋白质氨基酸序列的机制，以及使意识生物学获得合法地位，使他跻身哥白尼、牛顿、达尔文和爱因斯坦之列。然而，他对科学、对心智生命的强烈而终生的专注，是他与科学共同体中许多人共享的东西；这种执着象征着科学最好的样子。克里克的朋友和同事、认知心理学家维拉亚努尔·拉马钱德兰，描述了克里克生命最后几周对屏状核的关注。\end{zhblock}

\begin{enblock}
Three weeks prior to his death I visited him in his home in La Jolla. He was eighty-eight, had terminal cancer, was in pain, and was on chemotherapy; yet he had obviously been working away nonstop on his latest project. His very large desk--occupying half the room--was covered by articles, correspondence, envelopes, recent issues of \emph{Nature}, a laptop (despite his dislike of computers), and recent books on neuroanatomy. During the whole two hours that I was there, there was no mention of his illness--only a flight of ideas on the neural basis of consciousness. He was especially interested in a tiny structure called the claustrum which, he felt, had been largely ignored by mainstream pundits. As I was leaving he said: “Rama, I think the secret of consciousness lies in the claustrum--don't you? Why else would this tiny structure be connected to so many areas in the brain?” And gave me a sly, conspiratorial wink. It was the last time I saw him.
\end{enblock}
\begin{zhblock}
他去世前三周，我到拉霍亚的家中拜访他。他八十八岁，患有晚期癌症，忍受疼痛，并正在接受化疗；然而很明显，他一直不停地投入到最新项目中。他那张很大的书桌占了半个房间，上面堆满论文、信件、信封、近期的《自然》杂志、一台笔记本电脑（尽管他不喜欢电脑），以及新近出版的神经解剖学书籍。我在那里待了整整两个小时，期间没有提到他的病情，只有关于意识神经基础的思想飞翔。他特别感兴趣的是一个叫屏状核的微小结构，他觉得主流权威大多忽视了它。我离开时，他说：“拉玛，我认为意识的秘密就在屏状核里，你不这么想吗？否则为什么这个微小结构会连接到大脑这么多区域？”然后给了我一个狡黠、像同谋般的眨眼。那是我最后一次见到他。\end{zhblock}

\begin{enblock}
Since so little is known about the claustrum, Crick continued, he wanted to start an institute to focus on its function. In particular, he wanted to determine whether the claustrum is switched on when unconscious, subliminal perception of a given stimulus by a person's sensory organs turns into a conscious percept.
\end{enblock}
\begin{zhblock}
克里克接着说，由于关于屏状核所知太少，他想创办一个研究所，专门关注它的功能。尤其是，他想确定，当一个人的感觉器官对某个给定刺激的无意识、潜意识知觉转变为一个意识知觉时，屏状核是否被开启。\end{zhblock}

\begin{enblock}
One example of such switching that intrigued Crick and Koch is binocular rivalry. Here, two different images--say, vertical stripes and horizontal stripes--are presented to a person simultaneously in such a way that each eye sees only one set of stripes. The person may combine the two images and report seeing a plaid, but more commonly the person will see first one image, then the next, with horizontal and vertical stripes alternating back and forth spontaneously.
\end{enblock}
\begin{zhblock}
令克里克和科赫感兴趣的这种切换的一个例子是双眼竞争。在这种情形中，两种不同图像，例如竖条纹和横条纹，同时呈现给一个人，但呈现方式使每只眼睛只看见一组条纹。这个人可能把两幅图像合并起来，并报告说看见了格子图案；但更常见的是，他会先看见一幅图像，随后看见另一幅图像，横条纹和竖条纹自发地来回交替。\end{zhblock}

\begin{enblock}
Using MRI, Eric Lumer and his colleagues at University College, London have identified the frontal and parietal areas of the cortex as the regions of the brain that become active when a person's conscious attention switches from one image to another. These two regions have a special role in focusing conscious attention on objects in space. In turn, the prefrontal and posterior parietal regions of the cortex seem to relay the decision regarding which image is to be enhanced to the visual system, which then brings the image into consciousness. Indeed, people with damage to the prefrontal cortex have difficulty switching from one image to the other in situations of binocular rivalry. Crick and Koch might argue that the frontal and parietal areas of the cortex are recruited by the claustrum, which switches attention from one eye to the other and unifies the image presented to conscious awareness by each eye.
\end{enblock}
\begin{zhblock}
伦敦大学学院的埃里克·卢默及其同事利用MRI发现，当一个人的意识注意从一幅图像切换到另一幅图像时，皮层的额叶和顶叶区域会被激活。这两个区域在把意识注意聚焦于空间中的对象方面具有特殊作用。随后，皮层前额叶和后顶叶区域似乎把“哪一幅图像应当被增强”的决定传递给视觉系统，视觉系统再把该图像带入意识。事实上，前额叶皮层受损的人在双眼竞争情形中难以从一幅图像切换到另一幅图像。克里克和科赫也许会说，皮层额叶和顶叶区域是由屏状核招募的；屏状核把注意从一只眼切换到另一只眼，并统一每只眼呈现给意识觉知的图像。\end{zhblock}

\begin{enblock}
As these arguments make clear, consciousness remains an enormous problem. But through the efforts of Edelman on the one hand, and Crick and Koch on the other, we now have two specific and testable theories worthy of exploration.
\end{enblock}
\begin{zhblock}
正如这些论证所表明的，意识仍然是一个巨大的问题。但是，一方面通过埃德尔曼的努力，另一方面通过克里克和科赫的努力，我们现在拥有了两个具体且可检验、值得探索的理论。\end{zhblock}

\begin{enblock}
As someone interested in psychoanalysis, I wanted to take the Crick-Koch paradigm of comparing unconscious and conscious perception of the same stimulus to the next step: determining how visual perception becomes endowed with emotion. Unlike simple visual perception, emotionally charged visual perception is likely to differ between individuals. Therefore, a further question is, How and where are unconscious emotional perceptions processed?
\end{enblock}
\begin{zhblock}
作为一个对精神分析感兴趣的人，我想把克里克-科赫比较同一刺激的无意识知觉和意识知觉的范式推进到下一步：确定视觉知觉如何获得情绪色彩。不同于简单视觉知觉，带有情绪负荷的视觉知觉很可能因人而异。因此，进一步的问题是：无意识的情绪知觉在何处、以何种方式被加工？
\end{zhblock}

\begin{enblock}
Amit Etkin, a bold and creative M.D.-Ph.D. student, and I undertook a study in collaboration with Joy Hirsch, a brain imager at Columbia, in which we induced conscious and unconscious perceptions of emotional stimuli. Our approach paralleled in the emotional sphere that of Crick and Koch in the cognitive sphere. We explored how normal people respond consciously and unconsciously to pictures of people with a clearly neutral expression or an expression of fear on their faces. The pictures were provided by Peter Ekman at the University of California, San Francisco.
\end{enblock}
\begin{zhblock}
阿米特·埃特金是一名大胆而富有创造力的医学博士/哲学博士项目学生，我和他与哥伦比亚大学脑成像专家乔伊·赫希合作开展了一项研究，在其中诱发对情绪刺激的意识知觉和无意识知觉。我们在情绪领域采取的方法，平行于克里克和科赫在认知领域的方法。我们探索正常人如何以意识和无意识方式回应一些人的照片；这些人脸上带有明确的中性表情或恐惧表情。照片由加利福尼亚大学旧金山分校的保罗·艾克曼提供。\end{zhblock}

\begin{enblock}
Ekman, who has cataloged more than 100,000 human expressions, was able to show, as did Charles Darwin before him, that irrespective of sex or culture, conscious perceptions of seven facial expressions--happiness, fear, disgust, contempt, anger, surprise, and sadness--have virtually the same meaning to everyone (figure 28-1). We therefore argued that fearful faces should elicit a similar response from the healthy young medical and graduate student volunteers in our study, regardless of whether they perceived the stimulus consciously or unconsciously. We produced a conscious perception of fear by presenting the fearful faces for a long period, so people had time to reflect on them. We produced unconscious perception of fear by presenting the same faces so rapidly that the volunteers were unable to report which type of expression they had seen. Indeed, they were not even sure they had seen a face!
\end{enblock}
\begin{zhblock}
艾克曼已经编目了十万多种人类表情；他能够像此前的查尔斯·达尔文一样表明，不论性别或文化，七种面部表情，即快乐、恐惧、厌恶、轻蔑、愤怒、惊讶和悲伤，在意识知觉中对每个人实际上具有相同意义（图28-1）。因此我们认为，恐惧面孔应当在我们研究中的健康年轻医学生和研究生志愿者身上引发类似反应，无论他们是以意识方式还是无意识方式知觉到刺激。我们通过长时间呈现恐惧面孔来产生对恐惧的意识知觉，使人们有时间对它们进行反思。我们通过极快地呈现同样的面孔来产生对恐惧的无意识知觉，快到志愿者无法报告自己看见了哪一种表情。事实上，他们甚至不确定自己看见了一张脸。\end{zhblock}

\begin{figure}[htbp]
\centering
\safeincludeimage{983d9daa0542de65f607fe473e33e0b93324598b87b96544aac66a41c60daf77.jpg}
\caption{Ekman's seven universal facial expressions. (Courtesy of Paul Ekman.)\\艾克曼的七种普遍面部表情。（保罗·艾克曼惠允。）}
\end{figure}

\begin{enblock}
Since even normal people differ in their sensitivity to a threat, we gave all of the volunteers a questionnaire designed to measure background anxiety. In contrast to the momentary anxiety most people feel in a new situation, background anxiety reflects an enduring baseline trait.
\end{enblock}
\begin{zhblock}
由于即使正常人对威胁的敏感性也有差异，我们给所有志愿者发放了一份问卷，用来测量背景焦虑。与多数人在新情境中感到的瞬时焦虑不同，背景焦虑反映的是一种持久的基线特质。\end{zhblock}

\begin{enblock}
Not surprisingly, when we showed the volunteers pictures of faces with fearful expressions, we found prominent activity in the amygdala, the structure deep in the brain that mediates fear. What was surprising was that conscious and unconscious stimuli affected different regions of the amygdala, and they did so to differing degrees in different people, depending on their baseline anxiety.
\end{enblock}
\begin{zhblock}
不出所料，当我们向志愿者展示带有恐惧表情的面孔图片时，我们发现杏仁核有显著活动；杏仁核是位于大脑深部、中介恐惧的结构。令人惊讶的是，意识刺激和无意识刺激影响了杏仁核的不同区域，而且它们在不同人身上的影响程度也不同，这取决于他们的基线焦虑。\end{zhblock}

\begin{enblock}
Unconscious perception of fearful faces activated the basolateral nucleus. In people, as in mice, this area of the amygdala receives most of the incoming sensory information and is the primary means by which the amygdala communicates with the cortex. Activation of the basolateral nucleus by unconscious perception of fearful faces occurred in direct proportion to a person's background anxiety: the higher the measure of background anxiety, the greater the person's response. People with low background anxiety had no response at all. Conscious perception of fearful faces, in contrast, activated the dorsal region of the amygdala, which contains the central nucleus, and it did so regardless of a person's background anxiety. The central nucleus of the amygdala sends information to regions of the brain that are part of the autonomic nervous system--concerned with arousal and defensive responses. In sum, unconsciously perceived threats disproportionately affect people with high background anxiety, whereas consciously perceived threats activate the fight-or-flight response in all volunteers.
\end{enblock}
\begin{zhblock}
对恐惧面孔的无意识知觉激活了基底外侧核。在人类和小鼠中一样，杏仁核的这一区域接收大多数传入感觉信息，也是杏仁核与皮层交流的主要通道。由恐惧面孔无意识知觉引起的基底外侧核激活，与一个人的背景焦虑成正比：背景焦虑测量值越高，反应越强。背景焦虑低的人完全没有反应。相反，对恐惧面孔的意识知觉激活了杏仁核背侧区域，其中包含中央核；而且这种激活不受一个人背景焦虑的影响。杏仁核中央核把信息发送到大脑中属于自主神经系统的区域，这些区域与唤醒和防御反应有关。总之，无意识知觉到的威胁不成比例地影响背景焦虑高的人，而意识知觉到的威胁则在所有志愿者中激活战斗或逃跑反应。\end{zhblock}

\begin{enblock}
We also found that unconscious and conscious perception of fearful faces activates different neural networks outside the amygdala. Here again, the networks activated by unconsciously perceived threats were recruited only by the anxious volunteers. Surprisingly, even unconscious perception recruits participation of regions within the cerebral cortex.
\end{enblock}
\begin{zhblock}
我们还发现，对恐惧面孔的无意识知觉和意识知觉，会激活杏仁核之外的不同神经网络。在这里同样，由无意识知觉到的威胁激活的网络，只在焦虑志愿者中被招募。令人惊讶的是，即便无意识知觉也会招募大脑皮层内区域的参与。\end{zhblock}

\begin{enblock}
Thus viewing frightening stimuli activates two different brain systems, one that involves conscious, presumably top-down attention and one that involves unconscious, bottom-up attention, or vigilance, much as a signal of salience does in explicit and implicit memory in \emph{Aplysia} and in the mouse.
\end{enblock}
\begin{zhblock}
因此，观看令人恐惧的刺激会激活两个不同的大脑系统：一个涉及意识的、推测为自上而下的注意；另一个涉及无意识的、自下而上的注意或警觉。这很像显著性信号在海兔和小鼠的外显记忆与内隐记忆中所起的作用。\end{zhblock}

\begin{enblock}
These are fascinating results. First, they show that in the realm of emotion, as in the realm of perception, a stimulus can be perceived both unconsciously and consciously. They also support Crick and Koch's idea that in perception, distinct areas of the brain are correlated with conscious and unconscious awareness of a stimulus. Second, these studies confirm biologically the importance of the psychoanalytic idea of unconscious emotion. They suggest that the effects of anxiety are exerted most dramatically in the brain when the stimulus is left to the imagination rather than when it is perceived consciously. Once the image of a frightened face is confronted consciously, even anxious people can accurately appraise whether it truly poses a threat.
\end{enblock}
\begin{zhblock}
这些结果十分迷人。第一，它们表明，在情绪领域和知觉领域一样，一个刺激既可以被无意识地知觉，也可以被意识地知觉。它们还支持克里克和科赫的想法，即在知觉中，大脑不同区域与对刺激的意识觉知和无意识觉知相关。第二，这些研究从生物学上确认了精神分析中无意识情绪观念的重要性。它们提示，当刺激留给想象而不是被意识地知觉时，焦虑的影响在大脑中表现得最为剧烈。一旦恐惧面孔的图像被意识地面对，即使焦虑的人也能准确评估它是否真正构成威胁。\end{zhblock}

\begin{enblock}
A century after Freud suggested that psychopathology arises from conflict occurring on an unconscious level and that it can be regulated if the source of the conflict is confronted consciously, our imaging studies suggest ways in which such conflicting processes may be mediated in the brain. Moreover, the discovery of a correlation between volunteers' background anxiety and their unconscious neural processes validates biologically the Freudian idea that unconscious mental processes are part of the brain's system of information processing. While Freud's ideas have existed for more than one hundred years, no previous brain-imaging study had tried to account for how differences in people's behavior and interpretations of the world arise from differences in how they unconsciously process emotion. The finding that unconscious perception of fear lights up the basolateral nucleus of the amygdala in direct proportion to a person's baseline anxiety provides a biological marker for diagnosing an anxiety state and for evaluating the efficacy of various drugs and forms of psychotherapy.
\end{enblock}
\begin{zhblock}
在弗洛伊德提出精神病理源自无意识层面的冲突、而如果冲突源头被意识地面对就可以得到调节一个世纪之后，我们的成像研究提示了这类冲突过程可能如何在大脑中被中介。此外，志愿者背景焦虑与其无意识神经过程之间相关性的发现，从生物学上验证了弗洛伊德的思想：无意识心理过程是大脑信息加工系统的一部分。尽管弗洛伊德的思想已经存在一百多年，以前没有脑成像研究试图解释，人们行为和世界解释方式的差异如何来自他们无意识加工情绪方式的差异。无意识知觉恐惧会按一个人基线焦虑的比例点亮杏仁核基底外侧核，这一发现为诊断焦虑状态以及评价各种药物和心理治疗形式的疗效提供了一个生物学标志物。\end{zhblock}

\begin{enblock}
In discerning a correlation between the activity of a neural circuit and the unconscious and conscious perception of a threat, we are beginning to delineate the neural correlate of an emotion--fear. That description might well lead us to a scientific explanation of consciously perceived fear. It might give us an approximation of how neural events give rise to a mental event that enters our awareness. Thus, a half century after I left psychoanalysis for the biology of mind, the new biology of mind is getting ready to tackle some of the issues central to psychoanalysis and consciousness.
\end{enblock}
\begin{zhblock}
通过辨识一个神经回路的活动与对威胁的无意识知觉和意识知觉之间的相关性，我们正开始勾勒一种情绪，即恐惧的神经相关物。这种描述很可能引导我们对意识知觉到的恐惧作出科学解释。它也许能使我们近似地了解，神经事件如何产生一个进入我们觉知的心理事件。因此，在我离开精神分析、转向心智生物学半个世纪之后，新的心智生物学正准备处理某些对精神分析和意识而言居于中心地位的问题。\end{zhblock}

\begin{enblock}
One such issue is the nature of free will. Given Freud's discovery of psychic determinism--the fact that much of our cognitive and affective life is unconscious--what is left for personal choice, for freedom of action?
\end{enblock}
\begin{zhblock}
其中一个问题是自由意志的本性。鉴于弗洛伊德发现了心理决定论，即我们的认知生活和情感生活很大一部分是无意识的，那么个人选择、行动自由还剩下什么？
\end{zhblock}

\begin{enblock}
A critical set of experiments on this question was carried out in 1983 by Benjamin Libet at the University of California, San Francisco. Libet used as his starting point a discovery made by the German neuroscientist Hans Kornhuber. In his study, Kornhuber asked volunteers to move their right index finger. He then measured this voluntary movement with a strain gauge while at the same time recording the electrical activity of the brain by means of an electrode on the skull. After hundreds of trials, Kornhuber found that, invariably, each movement was preceded by a little blip in the electrical record from the brain, a spark of free will! He called this potential in the brain the “readiness potential” and found that it occurred one second before the voluntary movement.
\end{enblock}
\begin{zhblock}
关于这个问题的一组关键实验，是本杰明·利贝特于1983年在加利福尼亚大学旧金山分校完成的。利贝特以德国神经科学家汉斯·科恩胡伯的一项发现为出发点。在科恩胡伯的研究中，他要求志愿者移动右手食指。随后，他用应变计测量这种自主运动，同时通过头骨上的电极记录大脑电活动。经过数百次试验，科恩胡伯发现，每一次运动之前，大脑电记录中总会有一个小小的波动，一点自由意志的火花！他把大脑中的这种电位称为“准备电位”，并发现它发生在自主运动前一秒。\end{zhblock}

\begin{enblock}
Libet followed up on Kornhuber's finding with an experiment in which he asked volunteers to lift a finger whenever they felt the urge to do so. He placed an electrode on a volunteer's skull and confirmed a readiness potential about one second before the person lifted his or her finger. He then compared the time it took for the person to will the movement with the time of the readiness potential. Amazingly, Libet found that the readiness potential appeared not after, but 200 milliseconds before a person felt the urge to move his or her finger! Thus by merely observing the electrical activity of the brain, Libet could predict what a person would do before the person was actually aware of having decided to do it.
\end{enblock}
\begin{zhblock}
利贝特沿着科恩胡伯的发现继续研究，设计了一个实验：他要求志愿者在任何感到有冲动时抬起一根手指。他在志愿者头骨上放置电极，确认在这个人抬起手指前约一秒出现准备电位。随后，他把这个人产生行动意愿所需的时间与准备电位出现的时间进行比较。令人惊讶的是，利贝特发现，准备电位不是出现在一个人感到有冲动移动手指之后，而是出现在此前200毫秒！因此，仅仅通过观察大脑电活动，利贝特就能在一个人实际意识到自己已经决定要做某事之前，预测这个人会做什么。\end{zhblock}

\begin{enblock}
This finding has caused philosophers of mind to ask: If the choice is determined in the brain before we decide to act, where is free will? Is our sense of willing our movements only an illusion, a rationalization after the fact for what has happened? Or is the choice made freely, but not consciously? If so, choice in action, as in perception, may reflect the importance of unconscious inference. Libet proposes that the process of initiating a voluntary action occurs in an unconscious part of the brain, but that just before the action is initiated, consciousness is recruited to approve or veto the action. In the 200 milliseconds before a finger is lifted, consciousness determines whether it moves or not.
\end{enblock}
\begin{zhblock}
这一发现促使心灵哲学家发问：如果选择在我们决定行动之前就已经由大脑决定，那么自由意志在哪里？我们对自己意愿运动的感觉，是否只是一个幻觉，是对已经发生之事的事后合理化？或者，选择是自由作出的，但不是意识地作出的？如果是这样，行动中的选择就可能像知觉一样，反映无意识推理的重要性。利贝特提出，自主行动的启动过程发生在大脑的一个无意识部分，但在行动真正开始前，意识被招募来批准或否决该行动。在手指被抬起前200毫秒内，意识决定它是否移动。\end{zhblock}

\begin{enblock}
Whatever the reasons for the delay between decision and awareness, Libet's findings also raise the moral question: How can one be held responsible for decisions that are made without conscious awareness? The psychologists Richard Gregory and Vilayanur Ramachandran have drawn strict limits on that argument. They point out that “our conscious mind may not have free will, but it does have free won't.” Michael Gazzaniga, one of the pioneers in the development of cognitive neuroscience and a member of the American Council of Bioethics, has added, “Brains are automatic, but people are free.” One cannot infer the sum total of neural activity simply by looking at a few neural circuits in the brain.
\end{enblock}
\begin{zhblock}
无论决定与觉知之间延迟的原因是什么，利贝特的发现也提出了一个道德问题：一个人怎样能为没有意识觉知时作出的决定负责？心理学家理查德·格雷戈里和维拉亚努尔·拉马钱德兰为这种论证划定了严格界限。他们指出：“我们的意识心灵也许没有自由意志，但它确实有自由否决。”认知神经科学发展的先驱之一、美国生物伦理委员会成员迈克尔·加扎尼加又补充说：“大脑是自动的，但人是自由的。”不能仅仅通过观察大脑中的少数神经回路，就推断神经活动的总和。\end{zhblock}


% ===== 08_needham_sivin.tex =====
\section*{Needham: Chinese Scientific Terms, 五行, 阴阳, and 相关性思维}

\editorialnote{the chunk opens in the middle of a preceding discussion and with a partial Table 8. The table below keeps only the rows present in this file; ancient graph columns were too corrupt to restore safely.}

\begin{longtable}{>{\raggedright\arraybackslash}p{0.08\linewidth}
                  >{\raggedright\arraybackslash}p{0.23\linewidth}
                  >{\raggedright\arraybackslash}p{0.16\linewidth}
                  >{\raggedright\arraybackslash}p{0.09\linewidth}
                  >{\raggedright\arraybackslash}p{0.34\linewidth}}
\caption{Table 8 excerpt: ideographs relevant to early Chinese scientific terminology / 表8节选：与早期中国科学术语有关的字形}\\
\toprule
No. & Word & Modern Chinese & K no. & Cleaned remarks / 说明 \\
\midrule
\endfirsthead
\toprule
No. & Word & Modern Chinese & K no. & Cleaned remarks / 说明 \\
\midrule
\endhead
1a & affirmatory final particle & yeh / 也 & 4 & Possibly a cobra-like serpent, or the female external genitalia; the latter explanation was long accepted in Chinese tradition. / 可能取象于蛇，或取象于女阴；后说在中国传统中长期通行。\\
1b & affirmatory verb-noun, ``to be'', existence & shih / 是 & 866c & Sun plus foot and other strokes, probably related to cheng / 正, suggesting what is visible and not illusory. / 日与足等笔画组合，或与“正”相关，指日下可见、非虚妄者。\\
2a & negative particle, ``not'' & pu / 不 & 999 & A flower-head on a stalk with drooping leaves; probably a borrowed homophone. / 花蒂下垂之象，多半是假借同音字。\\
2b & negative verb-noun, non-existence & fei / 非 & 579 & Traditionally explained through the lower part of fei / 飞, two wings or birds back to back. / 传统上与“飞”的下部相释，作两翼或两鸟背向之象。\\
3 & different & i / 异 & 954 & A front-facing human figure with raised arms; perhaps linked with social difference and then otherness. / 正面人形举臂，或由社会差别引申为异、他者与差异。\\
4 & like, similar to & ju / 如 & 94g & Woman and mouth radicals; an early phonetic loan, with no clear archaic scientific significance. / 女、口合成；很早的假借或形声用法，古义与科学术语关系不明。\\
5 & if & jo / 若 & 777 & A kneeling person, perhaps gathering plants; likely a borrowed homophone. / 跪人采草之象，或许有关顺从义，但更可能是假借。\\
6 & change, permutation & i / 易 & 850 & A lizard, with meaning derived from colour change or quick shifts of position. / 蜥蜴之象，义或取其变色与迅速移动。\\
7 & gradual change, change of form & pien / 变 & 1780 & A later graph with silk hanks and a hand holding a stick, implying action and perhaps order from disorder. / 较晚之字，含丝与执杖之手，或寓由乱入治的动作。\\
8 & sudden change, change of substance & hua / 化 & 19 & Connected in Needham's table with exchange and the generalized idea of transformation. / 与交换和由此推广的变化观念相关。\\
9 & origin, first & yuan / 元 & 257 & A human figure in profile with emphasis on the head, hence first or beginning. / 侧身人形而突出头部，因而引申为首、始。\\
10 & cause, to rely on, following & yin / 因 & 370 & A woven mat, hence a basis or something relied upon; the meaning extends from static to temporal. / 席纹之象，义为凭借、根据，由空间的凭依引申到时间的因循。\\
11 & cause, reason, fact & ku / 故 & 49i & The ancient graph combines an old-object element with a hand holding a stick, suggesting prior action or precedent. / 古形含“古”意与执杖动作，指先事、先例。\\
12 & make, do, act & wei / 为 & 27 & An elephant with a human hand on its trunk, symbolizing dexterity. / 象与手相连，寓灵巧操作。\\
13 & begin & shih / 始 & 976 & A foetus and a woman; a beginning of beginnings. / 胎与女相合，表示发生之始。\\
14 & go, move & hsing / 行 & 748 & A diagram of crossroads. / 十字道路之象。\\
15 & go away, deprive, send away & chhü / 去 & 642 & A covered rice-basket; homophone borrowed for the present meaning. / 带盖食器之象，后假借为去。\\
16 & come to, reach, attain & chih / 至 & 413 & An arrow reaching target or ground. / 箭至其的或至地之象。\\
17 & stop & chih / 止 & 961 & A human foot. / 足形。\\
18 & end, finished, exhausted & chin / 尽 & 381 & A hand cleaning out a vessel with a brush. / 手以刷清器之象，故有尽、竭之义。\\
19 & true, truth, real & chen / 真 & 375 & A seal form probably implying full, filled up, or solid, opposed to empty or unreal. / 篆形义或为充实、坚实，与空虚不实相对。\\
20 & above, ascend, hand up & shang / 上 & 726 & Geometrical pictograph. / 几何式指事。\\
21 & below, descend, hand down & hsia / 下 & 35 & Geometrical pictograph. / 几何式指事。\\
\bottomrule
\end{longtable}

\begin{enblock}
There is a copious graphic vocabulary for study in the characters inscribed on Shang and Chou bronze vessels. Not all the characters have been identified, but enough exist to select ideographs that throw light on the origins of Chinese scientific terms. The ancient words probably had little influence on the thinking of the exponents of the proto-sciences in Ch'in and Han times, but for us they help us understand the ancient Chinese approach to science.
\end{enblock}
\begin{zhblock}
商、周青铜器铭文也提供了大量可供研究的图形词汇。并非所有文字都已释读，但其中已经足以选出若干能够说明中国科学术语起源的表意字。那些古字很可能并未明显影响秦汉时期原始科学诸家的思想；可是对于今天的研究者来说，它们本身有助于理解古代中国接近科学问题的方式。
\end{zhblock}

\begin{enblock}
A study of the table shows that the fundamental terms necessary for the beginnings of science were formed as one might expect from the principle of the ideograph. Only two of the terms here can be regarded as purely geometrical symbols; the rest are drawings of one kind or another. Some are borrowed homophones, and a few concern abstract ideas, but otherwise they depict natural objects, the human body and its parts, and human activities. What is most remarkable is that characters concerned with technology and communication are especially numerous. They demonstrate how ideographs developed from everyday life could later acquire abstract meanings, and how technical terminology for everyday thinking and experimentation actually arose.
\end{enblock}
\begin{zhblock}
考察这张表可以看出，科学开端所需的基本术语，正如表意文字的原则所预示的那样形成。这里真正可以看作纯几何符号的只有少数，其余大多是各种摹写性的图形。有些是假借同音字，有几项涉及抽象概念；其余则描绘自然物、人体及其部位，以及人的活动。尤其值得注意的是，与技术和交流有关的字形数量最多。这说明日常生活中的图像文字如何逐渐获得抽象意义，也说明日常思考和实验活动中的技术术语如何实际产生。
\end{zhblock}

\editorialnote{the printed source omits intervening pages here.}

\subsection*{The Stabilised Five-Element Theory / 定型化的五行学说}

\begin{enblock}
We are now in a position to consider the Five-Element theory as it was finalised in Han times and handed down to later ages. Two aspects merit special attention: the Enumeration Orders and the Symbolic Correlations. The Enumeration Orders are the orders in which the five elements were named in ancient and medieval presentations of the subject. They were not always the same, but the four most important were as follows.
\end{enblock}
\begin{zhblock}
现在可以考察汉代定型、并传之后世的五行学说。尤其值得注意的是两个方面：五行的排列次序，以及象征性相关系统。所谓排列次序，是古代和中古关于五行的论述中列举五行的顺序。它们并不总是相同，但最重要的四种如下。
\end{zhblock}

\begin{longtable}{>{\raggedright\arraybackslash}p{0.1\linewidth}
                  >{\raggedright\arraybackslash}p{0.34\linewidth}
                  >{\centering\arraybackslash}p{0.08\linewidth}
                  >{\centering\arraybackslash}p{0.08\linewidth}
                  >{\centering\arraybackslash}p{0.08\linewidth}
                  >{\centering\arraybackslash}p{0.08\linewidth}
                  >{\centering\arraybackslash}p{0.08\linewidth}}
\caption{Main Five-Element orders / 五行主要次序}\\
\toprule
i & Cosmogonic Order / 宇宙生成序 & Water / 水 & Fire / 火 & Wood / 木 & Metal / 金 & Earth / 土\\
ii & Mutual Production Order / 相生序 & Wood / 木 & Fire / 火 & Earth / 土 & Metal / 金 & Water / 水\\
iii & Mutual Conquest Order / 相胜（相克）序 & Wood / 木 & Metal / 金 & Fire / 火 & Water / 水 & Earth / 土\\
iv & ``Modern'' Order / “现代”口传序 & Metal / 金 & Wood / 木 & Water / 水 & Fire / 火 & Earth / 土\\
\bottomrule
\end{longtable}

\begin{enblock}
The Cosmogonic Order was the order in which the elements were supposed to have come into being. It begins with Water, echoing the recurrent Chinese emphasis on Water as the primeval element. The Mutual Production Order gave the seasons in correct order, beginning with Wood for spring and Water for winter, with Earth corresponding to an interval between summer and autumn.
\end{enblock}
\begin{zhblock}
宇宙生成序是设想五行产生的顺序。它以水为首，呼应了中国文献中不断出现的、以水为太初元素的强调。相生序则以正确的季节顺序呈现五行：木为春，水为冬，土对应夏秋之间的过渡月份。
\end{zhblock}

\begin{enblock}
The Mutual Conquest Order described the series in which each element was supposed to conquer its predecessor. It was associated with the teaching of Tsou Yen, and was based on everyday observations: Wood conquers Earth because a wooden spade digs earth; Metal conquers Wood because it cuts and carves it; Fire conquers Metal because it melts it; Water conquers Fire because it extinguishes it; Earth conquers Water because it dams and contains it. This order was also politically significant, offering an explanation of historical succession and therefore a tool for prediction. The ``Modern'' Order is obscure in significance, but it survives in colloquial Chinese as the familiar sequence Metal, Wood, Water, Fire, Earth.
\end{enblock}
\begin{zhblock}
相胜序描述每一行克制前一行的系列。它与邹衍的学说相联系，并以日常经验为根据：木胜土，因为木锹可以掘土；金胜木，因为金属可以砍削木材；火胜金，因为火能熔化金属；水胜火，因为水能灭火；土胜水，因为土能筑坝蓄水。对于重视灌溉和水利工程的人们来说，这是一种十分自然的比喻。这个次序还具有政治意义，被用来解释历史进程，并暗示它将继续适用于未来，因而可供预测。至于“现代”序，其意义较为晦暗，但“金、木、水、火、土”这一顺序已经进入汉语口语和童谣。
\end{zhblock}

\begin{enblock}
Two secondary principles in the Enumeration Orders concern the rate of change: the Principles of Control and Masking. The Principle of Control derives from the Mutual Conquest Order. A process of conquest is said to be controlled by the element that conquers the conqueror. Thus Metal conquers Wood but Fire controls the process; Fire conquers Metal but Water controls the process, and so on. Although used in fate-calculation, the Chinese were following logical paths of thought that modern experimental science has found applicable in fields such as ecology.
\end{enblock}
\begin{zhblock}
这些排列次序还牵涉两个关于变化速度的次级原则：控制原则和遮蔽原则。控制原则来自相胜序：某一克制过程由能够克制“克制者”的那一行所控制。例如金克木，而火控制这个过程；火克金，而水控制这个过程，如此类推。这个观念曾用于命算，但中国思想家所循的逻辑道路，在现代实验科学的若干领域中确有可比的应用，例如生态平衡中捕食者与被捕食者的数量关系。
\end{zhblock}

\begin{enblock}
One might criticize the Chinese principle of control by saying that in a cyclic process nothing could ever happen, since every element always inhibits another. But the Chinese did not suppose that all elements were effectively present everywhere at the same time. When the Mutual Production Order is related to the Mutual Conquest Order, the controlling element is always the one produced by the conquered element, and the system can continue. Wood conquers Earth in a process controlled by Metal, but Metal is produced by Earth, so feedback is built in.
\end{enblock}
\begin{zhblock}
有人或许会批评中国的控制原则：如果这是一个循环过程，那么一切都不可能发生，因为每一行总会抑制另一行。可是中国人并不认为所有五行在任何地方、任何时候都同样有效地同时在场。若把相生序与相胜序联系起来，就会发现控制者总是由被克者所生，系统因而能够继续运行。木克土而由金控制，但金又为土所生，于是反馈关系已经包含其中。
\end{zhblock}

\begin{enblock}
This idea could have social consequences, as when Confucians used it to prove that a son had the right to take revenge on his father's enemy. But its primary connotations were scientific. The idea that something acting on something else destroys it and is itself changed or destroyed in the process was known in ancient alchemy. It is also familiar in modern chemistry, for example when a reaction stops because reaction products accumulate; biochemistry offers even stronger examples.
\end{enblock}
\begin{zhblock}
这一观念也可能产生社会后果，例如儒家曾用它证明儿子有权为父复仇。但它的主要含义仍属科学性的。某物作用于他物、毁坏他物，同时又因这种作用而自身发生变化乃至毁坏，这一想法在古代炼金术中已经存在。现代化学同样熟悉这种情形，例如反应产物积累导致化学反应停止；生物化学中还有更有力的例证。
\end{zhblock}

\begin{enblock}
The Principle of Masking depends on both the Mutual Production and the Mutual Conquest Orders. It refers to the masking of a process of change by another process that creates more of the material than is being destroyed, or creates it faster. Wood destroys Earth, but Fire masks the process, since Fire destroys Wood and makes Earth, ash, faster than Wood can destroy Earth. Modern biological and ecological examples are easy to find.
\end{enblock}
\begin{zhblock}
遮蔽原则同时依赖相生序和相胜序。它指一种变化过程被另一种过程所遮蔽：后者生成的材料多于前者所毁坏的材料，或生成速度更快。木毁土，但火遮蔽这个过程，因为火能毁木并产生土，也就是灰，而且其速度快于木毁土。现代生物学和生态学中也可以找到类似例子。
\end{zhblock}

\begin{enblock}
In both principles a strong quantitative element lurks: the conclusions depend on quantities, speeds, and rates. This may have arisen from questions raised by the Enumeration Orders themselves. The early Chinese thinkers were not simply satisfied with them. The Mo Ching preserves a fragment of late Mohist criticism of the Naturalists: the Five Elements do not perpetually overcome one another. Fire naturally melts Metal if there is enough Fire; Metal can pulverise a burning fire to cinders if there is enough Metal. Metal stores Water but does not produce it; Fire attaches itself to Wood but is not produced from it.
\end{enblock}
\begin{zhblock}
在这两个原则中都潜藏着强烈的数量因素：结论取决于数量、速度和速率。这也许正来自排列次序本身所提出的问题，因为早期中国思想家并未满足于这些次序。《墨经》保存了后期墨家批评阴阳家的片段：五行并非永远相胜。只要火足够多，火自然熔金；只要金足够多，金也可把燃烧的火压成灰烬。金可以盛水，却不生水；火附着于木，却不由木生。
\end{zhblock}

\begin{enblock}
This attack on Tsou Yen's Mutual Conquest theory may have been a retort to his attitude toward the logical studies of the Mohist schools, but it remains interesting as a demonstration of the quantitative approach in Mohist scientific thinking.
\end{enblock}
\begin{zhblock}
这种对邹衍相胜说的攻击，也许是对他轻视墨家逻辑研究的一种反击；但无论如何，它作为墨家科学思维中数量化取向的证据，仍然十分有趣。
\end{zhblock}

\subsection*{The Symbolic Correlations / 象征性相关}

\begin{enblock}
The Five Elements gradually came to be associated with every conceivable category of things in the universe that could be classified in fives. Table 9 sets forth a selection.
\end{enblock}
\begin{zhblock}
五行逐渐同宇宙中一切可以分成五类的事物类别相联系。表9只列出其中一部分。
\end{zhblock}

\begin{longtable}{>{\raggedright\arraybackslash}p{0.12\linewidth}
                  >{\raggedright\arraybackslash}p{0.12\linewidth}
                  >{\raggedright\arraybackslash}p{0.12\linewidth}
                  >{\raggedright\arraybackslash}p{0.12\linewidth}
                  >{\raggedright\arraybackslash}p{0.12\linewidth}
                  >{\raggedright\arraybackslash}p{0.15\linewidth}
                  >{\raggedright\arraybackslash}p{0.16\linewidth}
                  >{\raggedright\arraybackslash}p{0.06\linewidth}}
\caption{Table 9a: symbolic correlations / 表9a：五行象征性相关}\\
\toprule
Element / 行 & Season / 时 & Direction / 方 & Taste / 味 & Smell / 臭 & Stems / 干 & Branches and animals / 支与动物 & Number / 数\\
\midrule
Wood / 木 & spring / 春 & east / 东 & sour / 酸 & goatish / 膻 & jia, yi / 甲乙 & yin tiger, mao hare / 寅虎、卯兔 & 8\\
Fire / 火 & summer / 夏 & south / 南 & bitter / 苦 & burning / 焦 & bing, ding / 丙丁 & wu horse, si serpent / 午马、巳蛇 & 7\\
Earth / 土 & sixth month or centre / 季夏或中央 & centre / 中 & sweet / 甘 & fragrant / 香 & wu, ji / 戊己 & xu dog, chou ox, wei sheep, chen dragon / 戌犬、丑牛、未羊、辰龙 & 5\\
Metal / 金 & autumn / 秋 & west / 西 & acrid / 辛 & rank / 腥 & geng, xin / 庚辛 & you cock, shen monkey / 酉鸡、申猴 & 9\\
Water / 水 & winter / 冬 & north / 北 & salty / 咸 & rotten / 腐 & ren, gui / 壬癸 & hai boar, zi rat / 亥豕、子鼠 & 6\\
\bottomrule
\end{longtable}

\begin{longtable}{>{\raggedright\arraybackslash}p{0.12\linewidth}
                  >{\raggedright\arraybackslash}p{0.14\linewidth}
                  >{\raggedright\arraybackslash}p{0.12\linewidth}
                  >{\raggedright\arraybackslash}p{0.15\linewidth}
                  >{\raggedright\arraybackslash}p{0.13\linewidth}
                  >{\raggedright\arraybackslash}p{0.12\linewidth}
                  >{\raggedright\arraybackslash}p{0.11\linewidth}
                  >{\raggedright\arraybackslash}p{0.09\linewidth}}
\caption{Table 9b: astronomical and political correlations / 表9b：天文与政治相关}\\
\toprule
Element / 行 & Musical note / 音 & Mansions / 宿 & Star palace / 宫 & Heavenly body / 天体 & Planet / 星 & Weather / 气象 & Colour / 色\\
\midrule
Wood / 木 & jiao / 角 & 1--7 & Azure Dragon / 青龙 & stars / 星 & Jupiter / 木星 & wind / 风 & green / 青\\
Fire / 火 & zhi / 徵 & 22--28 & Vermilion Bird / 朱雀 & sun / 日 & Mars / 火星 & heat / 热 & red / 赤\\
Earth / 土 & gong / 宫 & -- & Yellow Dragon / 黄龙 & earth / 地 & Saturn / 土星 & thunder / 雷 & yellow / 黄\\
Metal / 金 & shang / 商 & 15--21 & White Tiger / 白虎 & constellations / 宿 & Venus / 金星 & cold / 寒 & white / 白\\
Water / 水 & yu / 羽 & 8--14 & Sombre Warrior / 玄武 & moon / 月 & Mercury / 水星 & rain / 雨 & black / 黑\\
\bottomrule
\end{longtable}

\begin{longtable}{>{\raggedright\arraybackslash}p{0.11\linewidth}
                  >{\raggedright\arraybackslash}p{0.15\linewidth}
                  >{\raggedright\arraybackslash}p{0.13\linewidth}
                  >{\raggedright\arraybackslash}p{0.12\linewidth}
                  >{\raggedright\arraybackslash}p{0.12\linewidth}
                  >{\raggedright\arraybackslash}p{0.12\linewidth}
                  >{\raggedright\arraybackslash}p{0.13\linewidth}
                  >{\raggedright\arraybackslash}p{0.09\linewidth}}
\caption{Table 9c: bodily, agricultural, and ritual correlations / 表9c：身体、农事与礼仪相关}\\
\toprule
Element / 行 & Living class / 虫类 & Domestic animal / 牲 & Grain / 谷 & Sacrifice / 祭处 & Viscus / 脏 & Body part / 体 & Affect / 志\\
\midrule
Wood / 木 & scaly, fishes / 鳞类 & sheep / 羊 & wheat / 麦 & inner door / 户内 & spleen / 脾 & muscles / 筋 & anger / 怒\\
Fire / 火 & feathered, birds / 羽类 & fowl / 鸡 & beans / 豆 & hearth / 灶 & lungs / 肺 & pulse, blood / 脉血 & joy / 喜\\
Earth / 土 & naked, man / 倮类 & ox / 牛 & panicled millet / 稷 & inner court / 中庭 & heart / 心 & flesh / 肉 & desire / 欲\\
Metal / 金 & hairy, mammals / 毛类 & dog / 犬 & hemp / 麻 & outer door / 门外 & kidney / 肾 & skin and hair / 皮毛 & sorrow / 悲\\
Water / 水 & shell-covered / 介类 & pig / 豕 & millet / 黍 & well / 井 & liver / 肝 & bones, marrow / 骨髓 & fear / 恐\\
\bottomrule
\end{longtable}

\begin{enblock}
Careful analysis suggests that different groups of correlations were compiled by different groups of scholars. The astronomical group may go back to the ninth century B.C., though it shows the hand of the fourth-century astronomer Kan Te and probably the later astrologer Kan Chung-k'o. Then there were Naturalist groups associated with Tsou Yen, including the Yin-Yang group and the human psycho-physical functions. Finally, two groups are of particular scientific interest, one agricultural and one medical.
\end{enblock}
\begin{zhblock}
细加分析可以看出，不同组的相关系统很可能由不同学者群体编成。天文组也许早到公元前九世纪，尽管它显示出公元前四世纪天文学家甘德的影响，也可能受到后来的星占家甘忠可影响。随后是与邹衍相关的阴阳家配属，包括阴阳组以及人的心理-生理功能组。最后还有两组尤其具有科学史意义：一组主要关乎农业，一组主要关乎医学。
\end{zhblock}

\begin{enblock}
A striking point about the agricultural groups is that rice is absent from their lists of grains, although it appears in a medical group not reproduced here. This suggests that the agricultural correlations arose in northern China, or at an earlier date. Medical philosophers from very early times also had sixfold series complementing or replacing the fivefold series of scientific thinkers: six Yang organs, six Yin organs, and many analogous classifications.
\end{enblock}
\begin{zhblock}
农业组一个引人注意之处，是谷物表中没有稻，尽管稻出现在这里未列出的医学组中。这大概说明农业相关系统起源于中国北方，或者年代较早。医学思想家也很早就有六分系列，用来补充或替代科学思想家的五分系列：人体有六个阳腑、六个阴脏，医学中还可见许多类似分类。
\end{zhblock}

\begin{enblock}
These correlations were criticized because they led to absurdities. Wang Ch'ung pointed out in the first century A.D. that if yin corresponds to Wood and the tiger, while xu, chou, and wei correspond to Earth and to dog, ox, and sheep, then Wood conquering Earth would imply that the tiger overcomes the dog, ox, and sheep. If Water conquers Fire, why do rats not normally attack horses and drive them away? If Metal conquers Wood, why do cocks not devour hares?
\end{enblock}
\begin{zhblock}
这些相关系统曾遭到批评，因为它们会导出荒谬结论。公元一世纪的王充指出：若寅属木而配虎，戌、丑、未属土而配犬、牛、羊，那么木胜土就意味着虎胜犬、牛、羊。若水胜火，为什么鼠不常常攻击并赶走马？若金胜木，为什么鸡不吞食兔？
\end{zhblock}

\begin{enblock}
This criticism belonged to the Chinese skeptical tradition. Yet despite such attacks, these correlations seem at first to have helped scientific thought in China. They were no worse than the Greek theory of elements that dominated medieval Europe. They became harmful only when they grew over-elaborate and fanciful, too far removed from the observation of Nature.
\end{enblock}
\begin{zhblock}
这种批评属于中国怀疑论传统。可是尽管有这类攻击，在早期这些相关系统似乎确曾有助于中国的科学思维。它们并不比支配欧洲中世纪思想的希腊元素说更糟。只有当它们变得过度繁复、流于幻想、并且离自然观察太远时，才产生真正的危害。
\end{zhblock}

\begin{enblock}
As a positive example, Shen Kua's Dream Pool Essays records that in Qianshan, Xinzhou, a bitter spring forms a stream at the bottom of a gorge. When the water is heated it becomes tan fan, bitter alum, probably impure copper sulphate. When heated, this gives copper; if the alum is heated for a long time in an iron pan, the pan changes to copper. Shen Kua concluded that Water can be transformed into Metal.
\end{enblock}
\begin{zhblock}
作为正面例子，沈括《梦溪笔谈》记载：信州铅山县有苦泉，自峡谷底成溪。其水加热后成为胆矾，即苦矾，很可能是不纯的硫酸铜；胆矾再经加热可得铜；若在铁锅中久煎，铁锅会变成铜色。沈括由此认为水可以转化为金，这是一种异常的物质变化。
\end{zhblock}

\begin{enblock}
Shen Kua did not understand chemical change in the modern sense; like others of his time, he was dominated by Five-Element theory. But his observation is probably the earliest record in any language of metallic copper precipitated by iron, with the formation of iron sulphate.
\end{enblock}
\begin{zhblock}
沈括当然没有现代意义上的化学变化概念；和同时代人一样，他受五行理论支配。然而他的描述显示他是一位卓越的观察者，并且很可能是任何语言中关于铁置换出金属铜、并形成硫酸亚铁的最早记录。
\end{zhblock}

\subsection*{Numerology and Scientific Thinking / 数术与科学思维}

\begin{enblock}
Before leaving the Five-Element system, two examples of Chinese naturalistic thinking must be mentioned. Both are in the Ta Tai Li Chi, compiled between A.D. 85 and 105, though the quoted passages probably date from the second century B.C. In the first, Tseng Tzu explains that Heaven's Way is round and Earth's Way square. The bright radiates ch'i and the dark imbibes ch'i; Fire and the Sun have external brightness, while Metal and Water have internal brightness. The Yang is active and the Yin reactive.
\end{enblock}
\begin{zhblock}
在离开五行系统之前，还须提到中国自然主义思维的两个重要例子。二者都见于《大戴礼记》；该书成编于公元85至105年间，但所引文字大概可追溯到公元前二世纪。第一段中，曾子解释说天道圆、地道方；明者放射气，暗者吸纳气；火与日有外在之明，金与水有内在之明。因此阳主动，阴反应。
\end{zhblock}

\begin{enblock}
The passage continues: the seminal essence of Yang is called shen, and the germinal essence of Yin is called ling. Shen and ling are the roots of all living creatures and the ancestors of ritual, music, human-heartedness, righteousness, good and evil, order and disorder. When Yin and Yang keep their proper positions, there is quiet and peace. Hairy and feathered animals are born of Yang; animals with shells and scales are born of Yin; man alone is born naked because he has the balanced essences of both Yang and Yin.
\end{enblock}
\begin{zhblock}
这段文字继续说：阳之精称为神，阴之精称为灵。神与灵是一切生物的根本，也是礼乐、仁义、善恶、治乱等高级发展的祖先。阴阳各安其位，则安静太平。毛类与羽类动物由阳而生，介类与鳞类动物由阴而生；唯有人生而裸露，因为人兼具阴阳二气的平衡精华。
\end{zhblock}

\begin{enblock}
Needham stresses that calling the Sage the highest representative of naked animals is a supreme example of Chinese thought refusing to separate man from Nature, or individual man from social man. The passage also expresses the view that the basic forces at work in lower creatures are the same forces that, at higher levels, develop the highest forms of human social and ethical life.
\end{enblock}
\begin{zhblock}
李约瑟强调，把圣人称为裸虫的最高代表，极好地说明中国思想拒绝把人与自然分离，也拒绝把个人之人与社会之人分离。此段还表达了一种观点：在最低等生物中起作用的基本力量，与在较高层次上发展成人类社会生活和伦理生活最高形态的力量，是同一类力量。
\end{zhblock}

\begin{enblock}
The second passage is almost entirely biological, but its observations are fitted into number mysticism. Heaven is 1, Earth 2, Man 3; \(3 \times 3 = 9\), \(9 \times 9 = 81\). The Sun's number is 10, so man is born in the tenth month. Other multiplications are used to explain gestation periods of horses and dogs. The passage also classifies birds, fishes, silkworms, cicadas, gnats, scaled animals, shell-covered animals, birds, and mammals by habits, openings of the body, egg-laying, and nourishment in the womb.
\end{enblock}
\begin{zhblock}
第二段几乎完全是生物学性的，但其观察被纳入数字神秘主义框架。天为一，地为二，人为三；三三为九，九九八十一。日之数为十，所以人十月而生。其他乘数又被用来解释马、犬等动物的妊娠期。该段还按习性、身体孔窍、卵生与胎养等标准，分类叙述鸟、鱼、蚕、蝉、蜉蝣、鳞介动物、鸟类与哺乳动物。
\end{zhblock}

\begin{enblock}
The Naturalists were close observers of Nature, but their observations were fitted into number mysticism. Such mysticism appears already in the symbolic correlations and remained active as late as the twelfth century A.D.; it must be taken into account in any assessment of Chinese scientific ideas.
\end{enblock}
\begin{zhblock}
阴阳家，或者写下这些文字的人，显然善于观察自然；但他们把观察纳入数术框架。这种数字神秘主义已经见于象征性相关表，并且直到公元十二世纪仍有影响；评价中国科学观念时，必须把这一点纳入考虑。
\end{zhblock}

\subsection*{The Theory of the Two Fundamental Forces / 阴阳二气理论}

\begin{enblock}
Needham turns next to the two fundamental forces, Yin and Yang. We know less about them than about the Five Elements. They do not appear in surviving fragments of Tsou Yen, though his school was called the Yin-Yang Chia and later books usually credited their discussion to him. Their philosophical use probably began at the start of the fourth century B.C.; older passages that use them are likely interpolations.
\end{enblock}
\begin{zhblock}
李约瑟随后转向两种基本力量，即阴阳。相对于五行，我们对阴阳早期形态所知较少。现存邹衍残篇并不见阴阳二词，尽管他的学派被称为阴阳家，后世典籍也常把阴阳论归于他。阴阳作为哲学术语的使用大概始于公元前四世纪初；更早文本中出现的相关段落很可能是后人插入。
\end{zhblock}

\begin{enblock}
The Chinese characters for Yin and Yang are connected with darkness and light. Yin evokes cold, cloud, rain, femaleness, interiority, and darkness. Yang evokes sunshine, warmth, spring and summer, maleness, and brightness. They also had factual meanings: Yin was the shady side of a mountain or valley; Yang was the sunny side.
\end{enblock}
\begin{zhblock}
阴、阳二字与暗和明有关。阴唤起寒冷、云、雨、雌性、内部以及幽暗等观念；阳唤起日光、温暖、春夏、雄性与明亮等观念。它们也有更具事实性的意义：阴指山或谷的背阴面，阳指向阳面。
\end{zhblock}

\begin{enblock}
As philosophical terms they appear explicitly in the third-century B.C. appendices to the I Ching. There the universe is said to contain two fundamental forces or operations, one now dominating and then the other, in wave-like succession. The Mo Tzu and the Tao Te Ching also contain early references to the harmony of Yin and Yang.
\end{enblock}
\begin{zhblock}
作为哲学术语，阴阳在公元前三世纪《易经》附传中有明确用法。其意大致是：宇宙中只有两种基本力量或作用，它们此起彼伏，像波动一样交替占优。《墨子》和《道德经》也有关于阴阳和合的早期表述。
\end{zhblock}

\begin{enblock}
In the second century B.C., Tung Chung-shu wrote that Heaven has Yin and Yang, and so does man. When the Yin ch'i of Heaven and Earth begins to dominate, the Yin ch'i of man responds; if the Yin ch'i of man advances, the Yin ch'i of Heaven and Earth must also respond. Their Tao is one. To bring rain, Yin must be activated; to stop rain, Yang must be activated.
\end{enblock}
\begin{zhblock}
到公元前二世纪，董仲舒写道：天有阴阳，人亦有阴阳。天地之阴气开始占优，人之阴气也相应为先；若人之阴气先行，天地之阴气也应当随之兴起。其道为一。欲求雨，须发动阴；欲止雨，须发动阳。
\end{zhblock}

\begin{enblock}
The I Ching contains sixty-four hexagrams, each composed of six whole or broken lines corresponding to Yang and Yin. Figure 23 shows how Yang splits into two, then four, then eight, until sixty-four alternations exist. The process could continue indefinitely. Yin and Yang never become completely separated; at each stage, in any fragment, only one is manifest. Needham notes a scientific interest here: the repeated splitting of two factors, with one dominant and one recessive, has parallels in modern genetics.
\end{enblock}
\begin{zhblock}
《易经》有六十四卦，每卦由六爻组成，爻或整或断，分别对应阳与阴。图23显示阳如何分为二、再分为四、再分为八，直至形成六十四种交替。这个过程原则上可以无限继续。阴阳从未完全分离；在每一阶段、每一片段中，只有一方显现。李约瑟指出，这里具有科学史意义：两个因素反复分裂，其中一方显性、一方隐性，这与现代遗传学中的若干思路有相似之处。
\end{zhblock}

\begin{figure}[htbp]
\centering
\safeincludeimage{cd03d8e58d264dd5693c5d23747232e6341e7bbc70ec126811e0481af5645590.jpg}
\caption{Segregation table of the symbols of the Book of Changes, from Zhu Xi's twelfth-century \emph{Zhouyi benyi tushuo}. / 《易》象分化图，出自朱熹十二世纪《周易本义图说》。}
\end{figure}

\subsection*{``Associative'' Thinking and Its Significance / 相关性思维及其意义}

\begin{enblock}
Chinese scientific ideas involved two fundamental principles: the Two Forces, Yin and Yang, and the Five Elements. The Two Forces were derived from negative and positive projections of human sexual experience, while the Five Elements were believed to lie behind every substance and process. Everything in the universe that could be arranged in fives was correlated with the Five Elements; whatever could not fit that pattern belonged to other numerical orders such as fours, nines, or twenty-eights. This wider approach gave rise to number mysticism.
\end{enblock}
\begin{zhblock}
中国的科学观念包含两项基本原则：阴阳二气与五行。阴阳二气在根本上来自人类性经验的负、正投射；五行则被认为居于每一种物质与每一类过程的背后。宇宙中凡可分成五类者，都被配属到五行之中；凡不能归入五分框架者，则进入四、九、二十八等其他数目秩序。这个更广的取向产生了数字神秘主义。
\end{zhblock}

\begin{enblock}
Many European observers dismissed this number mysticism as superstition and claimed that it prevented true scientific thinking in China. Needham asks instead whether the traditional thought-system was merely superstition, merely primitive thought, or something characteristic of the civilization that produced it, and whether it contributed something positive to other civilizations.
\end{enblock}
\begin{zhblock}
许多欧洲观察者把这种数术视为纯粹迷信，并声称它阻碍了中国真正科学思维的兴起。李约瑟的问题则不同：传统思想体系究竟只是迷信、只是原始思维，还是产生它的文明所特有的某种思维形式？如果是后者，它是否也为其他文明提供了积极贡献？
\end{zhblock}

\begin{enblock}
Primitive magic has often been said to operate through two ``laws'': similarity, by which like produces like, and contagion, by which things once in contact continue to act on one another. These are precisely the kinds of laws behind Chinese correlations or associations. Their motive was magical, but that need not disturb us, since magic nourished early science. Magic tends toward the concrete and factual; it deals with reality and serves ordinary people through operations that resemble techniques.
\end{enblock}
\begin{zhblock}
对原始巫术的研究常说，它依据两条“法则”运作：相似律，即同类产生同类；接触律，即曾经接触过的事物即使分离后仍继续互相作用。中国的相关或配属系统背后，正是这类“法则”。其动机带有巫术性，但这并不必使我们不安，因为早期正是巫术滋养了科学。巫术趋向具体、事实，处理现实，并以类似技术操作的方式服务普通人。
\end{zhblock}

\begin{enblock}
Because of this practical emphasis, symbolic correlations were what a magician needed: they brought order to things when no one could know which conditions would bring success in magical or experimental procedures. If one experimented with water, it seemed logical not to wear red, the colour of fire. Such associations may have been intuitive, but at the time what else could they be?
\end{enblock}
\begin{zhblock}
正因这种实践取向，象征性相关系统正是术士所需要的东西：在无人确知何种条件会使命术或实验程序成功时，它为事物带来秩序。如果某人用水做实验或施术，那么不穿火之色的红衣就显得合乎逻辑。这类联想也许更多出于直觉，但在当时又还能是什么呢？
\end{zhblock}

\begin{enblock}
Modern scholars have called this mental approach ``associative thinking'' or ``co-ordinative thinking.'' It works by association and intuition, and has its own logic and laws of cause and effect. It is not mere superstition, though it differs from modern scientific thinking, where emphasis falls on external causes. It classifies ideas not in hierarchical ranks but side by side in a pattern. Things influence one another not mechanically, but by a kind of induction effect.
\end{enblock}
\begin{zhblock}
现代学者把这种心智取向称为“相关性思维”或“协调性思维”。它凭联想和直觉运作，并有自身的逻辑与因果法则。它并非单纯迷信，尽管它不同于强调外在原因的现代科学思维。它不是把观念排列成等级序列，而是把它们并置于一种图式之中。事物彼此影响，并不是通过机械原因，而是通过某种感应作用。
\end{zhblock}

\begin{enblock}
In this Chinese thought, the key words are Order and Pattern, or perhaps Organism. Symbolic correlations, correspondences, and the hexagrams of the I Ching formed part of one gigantic whole. Things behaved as they did not necessarily because of prior actions, but because their position in the cyclical universe endowed them with intrinsic natures. Their existence depended on the whole world-organism, and they reacted on one another by mysterious resonance.
\end{enblock}
\begin{zhblock}
在这里所讨论的中国思想中，关键词是秩序与图式，或者几乎可以说只有一个关键词：有机体。象征性相关、对应关系以及《易经》卦象，都是一个巨大整体的组成部分。事物以特定方式行动，并不一定因为此前其他事物的作用，而首先因为它们在不断变化的循环宇宙中所处的位置，使它们具有内在本性。它们的存在依赖整个世界有机体，并通过神秘的共振相互反应。
\end{zhblock}

\begin{enblock}
Tung Chung-shu expressed this in the Ch'un Ch'iu Fan Lu, in a chapter titled ``Things of the Same Genus Energise Each Other.'' Water poured on level ground avoids dry places and moves toward wet ones; fire avoids damp wood and ignites dry wood. Things reject what differs and follow what is akin. If ch'i are similar, they coalesce; if musical notes correspond, they resonate. The gong or shang note struck on one lute is answered by the same note on another instrument.
\end{enblock}
\begin{zhblock}
董仲舒在《春秋繁露》“同类相动”篇中最能表达这一点。水倾于平地，会避开干处而趋向湿处；火遇两块相同木柴，会避开湿者而燃着干者。万物排斥异类而趋附同类。气同则合，声应则鸣。一张琴上弹出宫声或商声，另一件弦乐器上相应的宫声或商声便会应和。
\end{zhblock}

\begin{enblock}
For Tung Chung-shu, causation was special and selective, acting in a stratified pattern rather than at random. The regularity of natural processes was conceived not as government by law, but as mutual adaptation to community life: a kind of mutual courtesy rather than strife among powers and processes.
\end{enblock}
\begin{zhblock}
对董仲舒而言，因果关系是特殊而有选择性的，它以分层图式起作用，而非随机发生。自然过程的规律性不是被设想为法则的统治，而是对共同体生活的相互适应：在无生命的力量和过程中，也有一种相互礼让，而非斗争。
\end{zhblock}

\begin{enblock}
If this expresses something deeply true about the Chinese world-picture, the fivefold correlations are an abstract chart of the whole thought-system. In true primitive thought anything may cause anything else; but once things are categorized in the fivefold system, anything can no longer be the cause of anything else.
\end{enblock}
\begin{zhblock}
如果这确实表达了中国世界图景中的深层真实，那么五分相关就是整个思想体系的一张抽象图表。在真正的原始思维中，任何事物都可能成为任何其他事物的原因；但一旦事物被纳入五分系统，任何事物便不再能够随意成为任何其他事物的原因。
\end{zhblock}

\begin{enblock}
Needham concludes that there are two ways of advancing from primitive truth. One, taken by some Greeks, refined causation toward a mechanical explanation of the universe, as in Democritus's atoms. The other systematized the universe of things and events into a structural pattern that conditioned mutual influences. In the Greek view, a particle was where it was because another particle pushed it there. In the other view, its behaviour was governed by its place in a field of force alongside similarly responsive particles; causation was environmental rather than mechanical.
\end{enblock}
\begin{zhblock}
李约瑟由此得出结论：从原始真理出发可以有两条进路。一条为若干希腊人所取，把因果观念精炼为宇宙的机械解释，如德谟克利特的原子论。另一条则把事物与事件的宇宙系统化为一种结构图式，由它制约各部分的相互影响。在希腊图景中，一个质点在某时某地，是因为另一个质点把它推到那里；在另一种图景中，它的行为由它在力场中的位置所支配，并与同样有反应能力的其他质点并列；这里的因果不是机械性的，而是环境性的。
\end{zhblock}

\begin{enblock}
The Democritean approach may have been necessary for modern science, but that does not make the Chinese view mere superstition. Aristotle too spoke of like attracting like, like nourishing like, and like affecting like. Greek thought as a whole moved toward mechanical cause and effect; Chinese thought developed the organic concept. The Chinese universe was not ruled by a creator-lawgiver or by atomistic clashes, but by a harmony of wills, spontaneous and patterned.
\end{enblock}
\begin{zhblock}
德谟克利特式进路也许是现代科学的必要前奏，但这并不意味着中国观点只是迷信。亚里士多德也曾说同类相吸、同类相养、同类相感。希腊思想总体上走向机械因果概念；中国思想则发展出有机体概念。中国的宇宙既不是由最高创造者兼立法者统治，也不是由原子的无情碰撞支配，而是由一种意志的和谐支配：自发而有图式。
\end{zhblock}

\begin{enblock}
Numbers reveal the contrast between the traditional Chinese view and modern science. China produced much creditable mathematics, but its numerology, connected with associative thinking, contributed little scientific value. Needham argues that it also did little harm; even the most exaggerated numerical correlations had validity in their own way, playing a part in Chinese scientific thinking as European extravagances foreshadowed later conceptions of laws of nature.
\end{enblock}
\begin{zhblock}
数字的使用清楚显示了传统中国观点与现代科学之间的对比。中国当然产生了许多优秀数学成果，但与相关性思维相连的数术似乎并未贡献多少科学价值。李约瑟同时认为，它也没有造成太大坏处；甚至最夸张的五行数字相关也以自己的方式具有有效性，并在中国科学思维发展中发挥了作用，正如欧洲某些荒诞做法也曾预示后来的自然法则概念。
\end{zhblock}

\begin{enblock}
Time and space were also conceived differently. For ancient Chinese thinkers, time was divided into seasons and subdivisions, though it flowed in one direction and was not the cyclical recurrent time of Indian philosophy. Space was not abstractly uniform, but divided into south, north, east, west, and centre, each connected with time and with the Five Elements. This compartmented world resembled medieval Europe before Galileo and Newton extended geometrical space and gravitation to the whole cosmos.
\end{enblock}
\begin{zhblock}
时间和空间也以不同方式被理解。对古代中国人来说，时间不是纯粹抽象量，而是分为不同季节及其细分；不过时间仍只向一个方向流动，中国并未采用印度哲学那种循环复归的时间观。空间也不是向各方向均匀延展的抽象物，而是分成南、北、东、西、中五方，每一方都与时间和五行相连。这个分隔化世界，很类似伽利略和牛顿把几何空间与万有引力推广到整个宇宙之前的中世纪欧洲。
\end{zhblock}

\begin{enblock}
For ancient Chinese thinkers, things were connected rather than caused. Tung Chung-shu wrote that opposing things cannot both arise simultaneously: Yin and Yang move parallel but not along the same road; they meet and take turns as controller. The universe is a vast organism, with one component after another taking the lead while all parts cooperate in mutual service, which is perfect freedom.
\end{enblock}
\begin{zhblock}
对古代中国思想家来说，事物之间更多是相连，而不是被造成。董仲舒说，自然常道是相反之物不能同时并兴：阴阳并行而不同路，相遇而轮流为主。宇宙是一个巨大的有机体，此时此地由这一部分领先，彼时彼地由另一部分领先，所有部分在相互服务中合作，而这正是完全的自由。
\end{zhblock}

\begin{enblock}
In such a system, causality is less a chain of events than what a modern biologist might call the endocrine orchestra of mammals, where all glands are working but it is not easy to say which leads at any moment. Chinese thought was dominated by a concept of causality in which succession was subordinated to interdependence.
\end{enblock}
\begin{zhblock}
在这样的系统中，因果性不像事件链条，而更像现代生物学家所谓哺乳动物的“内分泌乐队”：所有腺体都在工作，却不容易判断某一时刻究竟哪一项居于主导。中国思想中的因果观由这样一种观念支配：继起关系从属于相互依存关系。
\end{zhblock}

\editorialnote{the source chunk then omits intervening material and resumes with Sivin's essay.}

\section*{Nathan Sivin: Why the Scientific Revolution Did Not Take Place in China -- or Didn't It? / 席文：为什么科学革命没有在中国发生，或者说它没有发生吗？}

\begin{enblock}
This essay makes the case for two conclusions. First, why the Scientific Revolution did not take place in China is not a question that historical research can answer. It becomes useful mainly when one locates the fallacies that lead people to ask it. Second, by the criteria historians of science use, a scientific revolution did take place in China in the seventeenth century. It did not have the social consequences that we assume a scientific revolution will have. The most obvious conclusion is that those assumptions are mistaken.
\end{enblock}
\begin{zhblock}
本文论证两个结论。第一，“为什么科学革命没有在中国发生”并不是历史研究能够回答的问题；只有当我们辨认出促使人们提出这个问题的种种谬误时，它才有用。第二，按照科学史家通常使用的标准，中国在十七世纪确实发生过一场科学革命。只是它并未产生我们通常设想科学革命必然带来的社会后果。最明显的结论是：那些设想本身有误。
\end{zhblock}

\begin{enblock}
Anyone who has looked into the history of science, technology, and medicine over the last generation has known that all the great ancient civilizations had sophisticated traditions. The Chinese traditions are especially fascinating because they are fully recorded and were more independent of European influence than Islamic and Indian traditions. Beginning in the 1920s, Chinese and Japanese historians explained what the Chinese knew and did. In the 1950s Joseph Needham brought their work to educated readers in the West and encouraged further research. By now the study of China is one of the flourishing fields in the history of science.
\end{enblock}
\begin{zhblock}
过去一代以来，凡研究科学、技术和医学史的人都知道，古代各大文明都有自身复杂精致的传统。中国传统尤其引人入胜，因为其记录极为丰富，并且相较伊斯兰和印度传统，更少受欧洲影响。自二十世纪二十年代起，中日历史学家开始说明中国人知道什么、做了什么。二十世纪五十年代，李约瑟把他们的工作介绍给西方受过教育的读者，并鼓励进一步研究。如今，中国科学史已成为科学史中最繁荣的领域之一。
\end{zhblock}

\begin{enblock}
When people learn what scholars have uncovered, they usually wonder why the transition to modern science first happened where it did. Needham gave the classic formulation of the ``Scientific Revolution problem'' in 1969: Why did modern science, the mathematization of hypotheses about Nature with implications for advanced technology, rise meteorically only in the West at the time of Galileo? He added the question why Chinese civilization, from the first century B.C. to the fifteenth century A.D., was much more efficient than the West in applying natural knowledge to practical human needs.
\end{enblock}
\begin{zhblock}
当人们了解研究者已经揭示的材料后，通常会追问：为什么通向现代科学的转变首先在它实际发生的地方发生？李约瑟在1969年给出了“科学革命问题”的经典表述：为什么现代科学，即把关于自然的假说数学化、并对高级技术具有全部含义的那种科学，只在伽利略时代的西方迅猛兴起？他还提出第二个问题：为什么从公元前一世纪到公元十五世纪，中国文明在把自然知识应用于人类实际需要方面远比西方有效？
\end{zhblock}

\begin{enblock}
For much of that millennium and a half, European civilization was collapsing and then slowly recovering. If we want to explain European technological inferiority over fourteen centuries, we should look first at the western end of Eurasia, not the eastern. Even the claim of Chinese superiority is problematic, because the natural knowledge applied to human needs was not usually what we would call Chinese science.
\end{enblock}
\begin{zhblock}
在这一千五百年的大部分时间里，欧洲文明先经历缓慢而普遍的崩塌，然后又更缓慢地恢复。若要解释欧洲长达十四个世纪的技术落后，我们应当首先考察欧亚大陆西端，而不是远东。即使“中国优越”的说法也有问题，因为被应用于人类需要的自然知识，通常并不是我们所谓的中国科学。
\end{zhblock}

\begin{enblock}
Early technology did not succeed or fail according to how well it applied early science. In China, science was mainly practiced by members of a small educated class and transmitted in books. Technology was a matter of craft and manufacturing skills privately passed from artisans to children and apprentices. Most artisans could not read scientific books. They depended on practical and aesthetic knowledge that we can reconstruct only from artifacts and scattered literate testimony.
\end{enblock}
\begin{zhblock}
早期技术的成败并不取决于它在多大程度上应用了早期科学。在中国，科学主要由少数受教育阶层从事，并以书籍传承。技术则是工艺与制造技能，由工匠私下传给子女和学徒。多数工匠读不了科学家的书，只能依赖自身的实践知识与审美知识；这些知识今天只能通过遗物和少量文人记载来重建。
\end{zhblock}

\begin{enblock}
Comparing all scientific and engineering activity of one civilization with all that of another in a single generalization conceals more than it reveals. Only in modern times did these kinds of work become closely connected. A Chinese visitor to Europe before about 1400 would have found it technologically backward in many respects, but Chinese and European medicine before about 1850 may not have differed greatly in therapeutic effectiveness. Chinese mathematical astronomy by its last high point about 1300 never quite reached Ptolemy's predictive accuracy.
\end{enblock}
\begin{zhblock}
把一个文明的全部科学和工程活动与另一文明的全部活动用一句概括来比较，遮蔽的东西多于揭示的东西。只有在现代，这些工作才变得紧密相连。约1400年前到访欧洲的中国人，会在许多方面发现欧洲技术落后；另一方面，约1850年前中欧医学实践在疗效上大概并无巨大差距。中国数理天文学在约1300年的最后高峰，也并未完全达到托勒密一千一百年前已经掌握的预报精度。
\end{zhblock}

\begin{enblock}
Such comparisons tell us little about what we can learn from either culture. We do not stop studying European alchemy because Chinese alchemical literature is richer in chemical knowledge. What matters is that we can begin comparing several strong traditions of science and technology based on different ideas and social arrangements. Without such comparison we remain trapped in parochial viewpoints. Historians have more urgent work than proving that every other culture was inferior to the one they study.
\end{enblock}
\begin{zhblock}
这类比较并不能告诉我们可从任何一方文化学到什么。我们不会因为中国炼丹文献包含更丰富的化学知识，就停止研究欧洲炼金术。重要的是，我们现在能够开始比较若干强大的科学技术传统，它们建立在不同观念和社会安排之上。没有这种比较，我们就会被困在狭隘视角中。历史学家的更紧迫工作，绝不是证明其他一切文化都低于自己专攻的文化。
\end{zhblock}

\begin{enblock}
Shen Kua offers a clue to the character of early science in general. He is famous for discussing magnetic declination and movable type, applying permutations in traditional Chinese mathematics, proposing daily records of lunar and planetary positions, suggesting a purely solar calendar in East Asia, explaining land formation by silt deposition and erosion, and writing on medicine. He also contributed to archaeology, music, art criticism, literary criticism, economic theory, diplomacy, land reclamation, and political reform.
\end{enblock}
\begin{zhblock}
沈括提供了理解早期科学一般性质的线索。他以诸多成就著称：讨论磁偏角和活字印刷，在传统中国数学中应用排列组合，建议每日记录月亮和行星位置，在东亚首先提出纯太阳历，解释淤积和侵蚀造成的陆地形成，并撰写医学著作。他还对考古、音乐、艺术批评、文学批评、经济理论、外交、治水垦田和政治改革都有贡献。
\end{zhblock}

\begin{enblock}
Sivin studied Shen Kua to ask how traditional Chinese scientists explained to themselves what they were doing: how they understood nature, how they related to it as conscious individuals in society, and how the insights of various sciences fit together. Shen Kua seemed the obvious person to study because he was involved in so many fields.
\end{enblock}
\begin{zhblock}
席文研究沈括，是为了追问传统中国科学家如何向自己解释其所作所为：他们如何理解自然，作为生活在社会中的自觉个体又如何与自然发生关系，各门科学的洞见如何结合为总体理解。沈括涉猎诸多领域，因此似乎是最合适的研究对象。
\end{zhblock}

\begin{enblock}
The pattern that emerged was that there was no systematic connection among all the sciences in the minds of those who practiced them. They were not integrated under philosophy, as in European and Islamic schools and universities. China had sciences but no single ``science,'' no overarching conception or word for the sum of them.
\end{enblock}
\begin{zhblock}
研究显示，在从事这些学问的人心中，各门科学之间并不存在系统联系。它们没有像欧洲和伊斯兰的学校、大学中那样，被哲学统摄起来。中国有诸多“科学”，却没有一个单一的“科学”概念，也没有一个总括它们的词。
\end{zhblock}

\begin{enblock}
In Shen Kua's memoirs, physical and numerological aspects of astronomy, astrology, cosmology, and divination appear under ``regularities underlying the phenomena.'' Medicine, engineering, and mathematics appear under ``technical skills,'' alongside architecture and games. His chapter on ``strange occurrences'' includes plant fossils, a tornado, an experiment on rainbows, hearsay, and curiosities. What makes us think of Shen as a scientist was widely scattered through his own scheme of knowledge.
\end{enblock}
\begin{zhblock}
在沈括的笔记分类中，天文学、星占、宇宙论和占卜中物理与数术的方面，被归入“象数”或“物理之常”一类；医学、工程和数学则归入“技艺”，与建筑和游戏并列。他的“异事”一章同时包含植物化石、龙卷风、彩虹实验、传闻和奇谈。我们今天视为沈括“科学家”特征的内容，在他自己的知识体系中分散得很广。
\end{zhblock}

\begin{enblock}
Shen Kua's scheme cohered not at the level of science, but at a more general level. In his writing there are no clear boundaries between what fits the modern conception of science and what does not. That modern conception does not help us understand what Shen Kua was trying to do.
\end{enblock}
\begin{zhblock}
沈括的知识图式并不是在“科学”层面上保持一致，而是在更一般的层面上相互贯通。在他的著作中，符合现代科学概念的材料与不符合者之间没有清晰边界。用现代科学概念去理解沈括的意图，反而会造成障碍。
\end{zhblock}

\begin{enblock}
Shen lived in the late eleventh century, when social mobility was broadening the group that ruled China. Many new men were interested in practical affairs that earlier elites would have considered beneath them. Officials were expected to be competent and versatile; merit could be measured quantitatively in taxes collected, land reclaimed, and similar results. At leisure, members of Shen's group could indulge curiosity about anything, including technical matters once left to clerks or artisans.
\end{enblock}
\begin{zhblock}
沈括活跃于十一世纪后半叶，当时社会流动正在扩大中国统治集团。许多新进士人关心早先贵族会认为低贱的实际事务。官员被期望既能干又多才；政绩可用税收、垦田等量化指标衡量。闲暇时，沈括所属的大群体可以以业余方式追索宇宙万物，包括过去只适合吏胥或工匠的技术问题。
\end{zhblock}

\begin{enblock}
Only after Shen's lifetime did the amateur ideal settle again on philosophy, the arts, and literature, leaving earth and sky largely to technicians. In the eleventh century Shen was one of several polymaths whose scientific and technological interests emerged from official responsibilities. What connected his research interests were the diverse commitments of his civil service career.
\end{enblock}
\begin{zhblock}
直到沈括之后，士人的业余博雅理想才再次集中于哲学、艺术和文学，把天地之学大体留给技术人员。十一世纪的沈括只是若干博学者之一，他们的科学和技术兴趣都与多样的官职责任有关。把沈括各项研究兴趣联系起来的，正是他仕宦生涯中异常多样的职责与承诺。
\end{zhblock}

\begin{enblock}
Court astronomers, physicians, and alchemists had no reason to relate their arts to each other. Philosophers were not in a position to define a common discipline for all of them, as Aristotle and his successors had done in Europe. If anyone was going to seek common ground for the sciences in China, it was people like Shen Kua, who mastered many of them. But Shen connected his understanding in ways that did not directly link the fields of Chinese science and that intimately associated what we would call scientific with what we would call superstitious.
\end{enblock}
\begin{zhblock}
宫廷历算家、医生、炼丹者没有理由把自己的技艺彼此联系。哲学家也无力像亚里士多德及其后继者在欧洲所做的那样，为它们界定共同学科。若有人会寻找中国各门科学的共同基础，那必然是沈括这样通晓多门学问的人。然而沈括组织理解的方式，并未直接把中国科学各领域相连，而且把今天所谓科学与所谓迷信紧密结合在一起。
\end{zhblock}

\begin{enblock}
Sivin concludes that he failed to find the internal unity of Chinese science in Shen Kua's mind. Instead he learned the importance of relations between the sciences and other kinds of knowledge.
\end{enblock}
\begin{zhblock}
席文因此说，他未能在沈括的心智中找到中国科学的内在统一性。作为补偿，他学到了另一问题的重要性：各门科学与其他知识类型之间的关系。
\end{zhblock}

\subsection*{Back to the Scientific Revolution Problem / 回到科学革命问题}

\begin{enblock}
The question ``Why didn't the Chinese beat Europeans to the Scientific Revolution?'' is one of the few public questions people ask about why something did not happen in history. It is like asking why your name did not appear on page 3 of today's newspaper. Historians do not build research programs around such questions because they have no direct answers. They translate into questions about what did happen elsewhere: under what circumstances did the Scientific Revolution take place in seventeenth- and eighteenth-century Western Europe?
\end{enblock}
\begin{zhblock}
“为什么中国人没有抢在欧洲人之前完成科学革命？”是少数常被公开提出的、关于历史上某事为何没有发生的问题之一。它类似于问：为什么你的名字没有出现在今天报纸第三版？历史学家不会围绕这类问题建立研究计划，因为它们没有直接答案。它们只能转化为关于别处实际发生了什么的问题：科学革命是在怎样的条件下于十七、十八世纪的西欧发生的？
\end{zhblock}

\begin{enblock}
People keep asking the China question because it is heuristic: it encourages exploration and gives order to thought. Heuristic questions are useful at the beginning of inquiry, but as we understand complex patterns they grow murky and lose interest compared with the emerging clarity of what did happen.
\end{enblock}
\begin{zhblock}
人们之所以不断追问中国问题，是因为它有启发性：它鼓励探索，并为思考提供秩序。启发性问题在研究开端有用；但当我们逐渐理解复杂图式时，它们会变得含混，并且相较于“实际发生了什么”所呈现的清晰性而失去吸引力。
\end{zhblock}

\begin{enblock}
The urgency of the Scientific Revolution problem rests on shaky Western assumptions. Above all, we assume that the Scientific Revolution is what everyone ought to have had. It is not clear that everyone wanted it before modern survival under violent change made it urgent. We know little about how Europeans originally came to want that revolution because historians have focused on how it took place.
\end{enblock}
\begin{zhblock}
科学革命问题之所以显得迫切，是因为它依赖若干不稳固的西方假设。最重要的是，我们通常假定科学革命是每个文明都应当拥有的东西。然而在近代剧烈变动使其成为生存所需之前，并不清楚每个文明都曾想要它。我们对欧洲人最初如何在一个又一个国家中想要这场革命所知甚少，因为历史学家主要关注它如何发生。
\end{zhblock}

\begin{enblock}
These assumptions often link to the faith that European civilization, especially in its current American form, was somehow in touch with reality in a way no other civilization could be, and that its wealth and power came from an intrinsic fitness to inherit the earth. Historical study suggests instead that Western privilege came from a head start in technological exploitation of nature and political exploitation of societies not technologically equipped to defend themselves.
\end{enblock}
\begin{zhblock}
这些假设常常同一种信念相连：欧洲文明，尤其是其当代美国形态，以某种其他文明不可能具备的方式接触了现实；其财富和权力来自一种内在的、继承地球的适应性。历史研究并不支持这种看法。它更提示我们，西方的特权地位来自它在技术性开发自然和政治性剥削那些无技术能力自卫的社会方面取得了先发优势。
\end{zhblock}

\begin{enblock}
Scientists also tend to believe that because science has rapidly become international, it transcends European historical and philosophical biases and is as universal, objective, and value-free as the Nature it studies. This common sense does not survive examination. Modern science remains too marked by the special circumstances of its European development to be considered universal in that strong sense.
\end{enblock}
\begin{zhblock}
科学家还常相信：既然科学已经迅速国际化，它就超越了欧洲的历史和哲学偏见，并且像它所研究的自然一样普遍、客观、价值中立。这种常识经不起仔细检验。现代科学仍然深受其在欧洲发展时特殊情境的印记，不能在强意义上被视为普遍。
\end{zhblock}

\begin{enblock}
Chinese science got along without dichotomies between mind and body, objective and subjective, even wave and particle. In the West, mind-body and objective-subjective distinctions were entrenched by Plato and carried into modern times by Galileo, Descartes, and others. They helped early modern scientists claim authority over the physical world by separating natural knowledge from the realm of the soul and established religion.
\end{enblock}
\begin{zhblock}
中国科学并不依赖身心、客观与主观、甚至波与粒子之间的二分。在西方，身心和主客二分在柏拉图时代已根深蒂固，伽利略、笛卡尔等人又把它们带入近代。这些区分帮助早期现代科学家主张对物理世界的权威，因为纯粹自然知识被划离灵魂领域和既有宗教权威。
\end{zhblock}

\begin{enblock}
Science and religion have long learned to coexist, but we still live with sharp distinctions of this kind. If they are European peculiarities and sources of trouble, why has modern science not discarded them? It is not simple to root them out. Until we do, we should admit a certain parochialism in the foundations of science. Mathematical equations may be universal, but the allocation of human effort among possibilities of natural knowledge is not.
\end{enblock}
\begin{zhblock}
科学与宗教早已学会共存，但我们仍生活在这些尖锐区分之中。如果它们是欧洲特殊性，而且不断制造麻烦，为什么现代科学还没有摆脱它们？显然，连根拔除并不简单。在做到这一点之前，我们应当坦率承认科学基础中存在某种地方性。数学方程也许是普遍的，但人类把精力分配给哪些自然知识可能性，并不是普遍的。
\end{zhblock}

\begin{enblock}
Science and technology have spread worldwide, but that has not made them universal in the sense of transcending European patterns of thought. Encounters between old and new ideas in many societies have often been resolved by social change and political legislation, with traditional ideas excluded as backward, superstitious, or fit only for lower classes. True universality would require modern technology to coexist with and serve cultural diversity rather than standardizing it out of existence.
\end{enblock}
\begin{zhblock}
科学和技术已传播到全世界，但这并不意味着它们在超越欧洲思维模式的意义上成为普遍者。在许多社会中，新旧观念的相遇常常通过社会变革和政治立法来解决，传统观念则被以落后、迷信、倒退、只适合下层阶级等理由排除。真正的普遍性要求现代技术与文化多样性共存并服务于它，而不是把它标准化到消失。
\end{zhblock}

\begin{enblock}
Modern science has attained verifiability, internal consistency, taxonomic grasp, precision, and predictive accuracy unmatched by earlier activities. But the rigor that makes these possible quickly disappears when a law or theory is translated from equations, matrices, or technical models into ordinary cultural discourse. Such translation into value-laden analogies and metaphors precedes all public discussion of science and much philosophical reflection.
\end{enblock}
\begin{zhblock}
现代科学确已达到早期知识活动所不能及的可验证性、内在一致性、分类把握、精确说明和预测准确性。但一旦定律或理论从方程、范畴矩阵或严格技术模型转化为普通文化话语，使这些特征成为可能的严密性便迅速消失。这种转化为充满价值的类比和隐喻，是一切关于科学的公共讨论以及大部分哲学反思的前提。
\end{zhblock}

\begin{enblock}
Beyond the narrow abstract realm of exactitude, values and subjective judgments shape every activity situated in society. Contemporary science in China and the United States differs in convictions about basic and applied research, relations to general culture, scientists' roles in defining programs, planning and support, social aims, professional status, politics and knowledge, and allocation of resources. Invariant equations matter, but so does the ubiquity of opposable thumbs. Words like international and universal are out of place when applied beyond the narrow technical realm.
\end{enblock}
\begin{zhblock}
在精确性可能成立的狭窄抽象领域之外，价值和主观判断会进入社会中的每一项活动。当代中国和美国的科学活动，在基础与应用研究关系、科学与一般文化关系、科学家制定计划的角色、规划和资助程序、社会目标、职业组织与地位、政治观念与科学知识的联系、资源分配等方面，都有深刻差异。不变的方程固然是因素，但拇指对掌的普遍存在也是因素。只要讨论超出狭窄技术领域，“国际”与“普遍”这样的词就并不合适。
\end{zhblock}

\begin{enblock}
Nor should we accept uncritically the idea that modern science is in every essential respect European in origin. Any account of early science is lopsided if it restricts itself to discoveries at the western end of Eurasia, ignores movement of ideas among civilizations, neglects what Europeans had learned by 1600 about Islamic, Indian, and Chinese science, or ignores the impact of exotic technologies and materials on European experience.
\end{enblock}
\begin{zhblock}
同样，我们也不能不加批判地接受现代科学在一切本质方面都起源于欧洲的说法。任何早期科学史叙述，如果把自己局限于欧亚大陆西端的发现，忽视各文明之间观念的流动，未充分考虑1600年前欧洲人已经从伊斯兰、印度和中国科学中学到什么，或忽略异域技术和材料对欧洲经验的影响，就会在根本问题上偏斜误导。
\end{zhblock}

\subsection*{Fallacies of Historical Reasoning / 历史推理谬误}

\begin{enblock}
Growing awareness of ancient Chinese science and technology has produced many hypotheses about factors that inhibited modern science in China, or characteristics unique to the West that made a major scientific revolution possible. These often contain elementary fallacies of historical reasoning.
\end{enblock}
\begin{zhblock}
随着人们越来越意识到古代中国科学技术的高水平，各种关于“中国何以阻碍现代科学演化”的因素，或“西方何以独具促成重大科学革命的特征”的假说纷纷出现。它们常常包含基本的历史推理谬误。
\end{zhblock}

\begin{enblock}
For decades historians argued that Ch'ing thinkers took the world as observable, nominalistic fact, like Francis Bacon, but unlike Bacon did not develop a scientific methodology. They did not ask whether Bacon's method survived in contemporary science. Bacon's method was largely Scholastic in origin, concerned with taxonomies rather than theories of natural phenomena, and unconcerned with mathematical measurement. It was probably among the most sterile early modern attempts to define how physical science should proceed.
\end{enblock}
\begin{zhblock}
几十年来，历史学家常说清代思想家像弗朗西斯·培根一样，把世界看作可观察的、唯名论式的事实；但与培根不同，他们没有发展出科学方法论。这些论者甚至没有追问培根的科学方法是否还保存在当代科学实践中。培根的方法很大程度上源自经院传统，关心分类而非自然现象理论，并且不关心数学测量。在近代早期各种试图界定物理科学应如何前进的方案中，它大概是最贫瘠的之一。
\end{zhblock}

\begin{enblock}
This pattern faults the Chinese for not developing a method that later proved abortive in the West. Another example claims that late Warring States astronomers failed to develop a unified scientific system because they were not interested in applied technical sciences such as tools for cannon flight or navigation. This blames a civilization for lacking interests more than a millennium before cannon existed and ignores comparable European theorists whose lack of application did not prevent later mechanics.
\end{enblock}
\begin{zhblock}
这种思路责备中国人没有发展出一种后来在西方证明并无结果的方法。另一例认为，战国晚期的中国天文学家未能发展统一科学体系，是因为他们不关心可用于控制炮弹飞行或航海的应用技术科学。这实际上在火炮出现一千多年前责备一个文明缺乏相关兴趣，同时忽略欧洲冲力理论家同样缺乏应用兴趣却并未阻止后来的力学发展。
\end{zhblock}

\begin{enblock}
The general fallacy is to claim that if an important aspect of the European Scientific Revolution is absent in another civilization, the whole set of changes could not have happened there. This rests on the arbitrary, usually unspoken assumption that the circumstance is a necessary condition. But if we trace the prehistory of the actual Scientific Revolution far enough back, the same circumstance is often absent in Europe too. On what grounds, then, is it necessary?
\end{enblock}
\begin{zhblock}
一般说来，这种谬误就是声称：如果欧洲科学革命的某个重要方面在另一文明中不存在，那么那整组根本变化就不可能在那里发生。其基础是一个任意而通常未明说的假设：这个因素是必要条件。可是如果把实际科学革命的史前史向前追溯得足够远，在多数情况下也能找到欧洲尚不具备这个因素的时点。那么，它凭什么被视为必要条件？
\end{zhblock}

\begin{enblock}
The mirror image appears in estimates of the Book of Changes as a deterrent to science. Needham once wrote that while Five-Element and Yin-Yang theories favoured scientific thought in China, the elaborate symbolic system of the Book of Changes was a mischievous handicap, tempting those interested in Nature to rest in explanations that were no explanations at all. Later scholars similarly called it an inhibiting factor.
\end{enblock}
\begin{zhblock}
这个谬误的镜像，见于把《周易》视为科学阻碍的评价。李约瑟曾写道，五行和阴阳理论对中国科学思维有利而非有害，但《易经》精巧的象征系统几乎从一开始就是有害障碍，引诱关心自然的人停留在根本不是解释的解释上。后来学者也常把它称为抑制因素。
\end{zhblock}

\begin{enblock}
But natural philosophers often used the Book of Changes to build dynamic explanations of change, not merely static classifications. If one cannot show that Chinese scientists were steadily moving toward mathematical formulations and experimental verification, and if there is no evidence that without the Book of Changes they would have gone further, there is no warrant for calling it inhibition.
\end{enblock}
\begin{zhblock}
可是自然哲学家常用《易》建构关于变化的动态解释，而不只是静态分类。若不能证明早期中国科学家正在稳定走向数学公式化与实验验证，也没有具体证据说明没有《易》他们就会“更进一步”，那么就没有理由把《易》称为“抑制”。
\end{zhblock}

\begin{enblock}
The same confusion appears when scholars call the scholar-bureaucrat class an inhibiting factor: bookish, past-oriented, and focused on human institutions rather than Nature. But at the onset of the European Scientific Revolution, universities were dominated by Schoolmen and dons who were also bookish, past-oriented, and focused on human institutions. They did not prevent the great changes in Europe.
\end{enblock}
\begin{zhblock}
同样的混乱也出现于把士大夫官僚阶层称为“抑制因素”时：他们沉浸于书本，面向过去，关注人类制度而非自然。可是欧洲科学革命开始时，大学中占支配地位的是经院学者和学院教师，他们同样沉浸于书本，面向过去，关注人类制度。他们并没有阻止席卷欧洲的重大变化。
\end{zhblock}

\begin{enblock}
One might as well call Euclidean geometry an inhibiting factor for non-Euclidean geometry, since people satisfied with it did not move to the next step. But can we argue that non-Euclidean geometry would have developed sooner without Euclid? It is unfortunate to dismiss the technical language of the Book of Changes as an obstacle before its dimensions have been understood.
\end{enblock}
\begin{zhblock}
按这种逻辑，也可以把欧几里得几何称为非欧几里得几何发展的抑制因素，因为人们满足于前者时便没有走向下一步。但我们能说没有欧几里得，非欧几里得几何就会更早出现吗？在尚未理解《易》那套强有力地系统关联广泛人类经验的技术语言之前，就把它斥为障碍，是很可惜的。
\end{zhblock}

\begin{enblock}
The first fallacy confuses a culture's earlier state, or its way of doing something, with a cause or necessary condition. The second confuses the absence of the later state with an inhibitor. Both confound continuity with stasis. They are bad history because they are bad reasoning.
\end{enblock}
\begin{zhblock}
第一种谬误把一种文化的较早状态，或它做事的方式，误当成原因或必要条件。第二种谬误则把后来状态的缺失误当成抑制因素。二者都把连续性与停滞混为一谈。它们之所以是坏历史，是因为它们首先是坏推理。
\end{zhblock}

\begin{enblock}
Together these fallacies turn the history of world science into a saga of Europe's success and everyone else's failure, until redemption through modernization. If Europe had the horse and buggy, fallacy one makes its absence in China proof that an automobile analogue was impossible. If China had something analogous, fallacy two makes it an inhibiting factor. Thus European abstraction can be read as a stage toward inertial guidance, while Chinese abstraction is read as proof that such guidance could never arise in East Asia.
\end{enblock}
\begin{zhblock}
这两种谬误合用时，会把世界科学史变成欧洲成功、他者失败，直到现代化救赎降临的传奇。如果欧洲有马车，第一种谬误就会把中国没有马车读作汽车类似物不可能出现的证据；如果中国也有某种类似马车的东西，第二种谬误又会把它说成抑制因素。于是，欧洲抽象理论可被读作通向惯性制导的一阶段，而中国思想者同样不关心应用，则被读作东亚不可能产生这种技术的证据。
\end{zhblock}

\begin{enblock}
This is an infallible formula for reading the strength of modern science into Europe's past alone. Other civilizations are tested only by anticipation of, or approximation to, early European or modern science. European science is exempt because it is assumed to have given birth to the Scientific Revolution. Other civilizations shine only as they reflect European light.
\end{enblock}
\begin{zhblock}
这是一套万无一失的公式：把现代科学的力量读回历史过去，但只读回欧洲的过去。其他文明则只按其是否预示或接近早期欧洲科学或现代科学来检验。欧洲科学为什么无需检验？因为人们预设唯有它孕育了科学革命。于是其他文明只有在反射欧洲传统之光时才会发亮。
\end{zhblock}

\begin{enblock}
Such assumptions are disastrous because they assure us there is no point in comprehending non-Western technical inquiries on their own terms. They have been accepted not only in Europe but to some extent wherever Chinese history is studied. Sivin cites Nishijima Sadao's analysis of the ``static character'' hypothesis: it served to validate the self-consciousness of modern Western Europe and, in Japan, shaped ways of viewing China through modernization and backwardness.
\end{enblock}
\begin{zhblock}
这些假设是灾难性的，因为它们让我们确信没有必要按非西方技术探究自身的术语去理解它们。它们不仅在欧洲被接受，也在研究中国史的各国不同程度地被接受。席文引用西嶋定生对“中国社会停滞性”假说的分析：这一假说原本服务于确认现代西欧社会自我意识的价值；在日本，它也塑造了通过现代化与落后来观看中国的方式。
\end{zhblock}

\begin{enblock}
One more fallacy appears when historians choose aspects of the European experience for comparison: the assumption that the evolution of science can be explained by intellectual factors alone, or socioeconomic factors alone. People who think of science as a quest for truth see China's failure as a failure of thought; those who see science as social or economic see the defeat as backwardness. Neither exclusive approach is adequate. The distinction between internal and external factors lies in historians' habits and division of labor, not in the events themselves.
\end{enblock}
\begin{zhblock}
还有一种谬误，出现在历史学家选择欧洲经验的哪些方面来比较时：他们假定科学演化可以只用思想因素，或只用社会经济因素来解释。把科学主要视为追求自然真理的人，倾向于把中国未能先于英国进入现代科学看作思想失败；把科学主要视为社会或经济现象的人，则倾向于把这种失败看作社会经济落后。二者都不充分。内在因素与外在因素的区分，并不在被研究事件本身之中，而在历史学家的心智习惯、职业社群和劳动分工之中。
\end{zhblock}

\subsection*{Dimensions of the Scientific Revolution / 科学革命的维度}

\begin{enblock}
The Scientific Revolution and its consequences cut across historical specializations. Intellectually, it transformed knowledge of the external world: the questions asked, the means used, and the character of answers. It established for the first time the dominion of number and measure over every physical phenomenon.
\end{enblock}
\begin{zhblock}
科学革命及其后果跨越了历史专业分工的边界。从思想维度看，它转变了人们关于外部世界的知识：改变了所提问题、探索手段和答案性质。它首次确立了数与量度对一切物理现象的支配。
\end{zhblock}

\begin{enblock}
Ernst Gellner noted that the European Scientific Revolution was more than a new form of knowing. Earlier sciences asked several questions, not only ``is it true?'' but also whether it was beautiful, conventional, morally improving, led to the Good, or conformed to aesthetic patterns. Modern science displaced most of these with the test of truth, an appeal to public, verifiable, morally neutral fact. It created such facts for the first time: knowledge with no value except truth value.
\end{enblock}
\begin{zhblock}
厄内斯特·盖尔纳指出，欧洲科学革命不仅是跃迁到一种新的认识形式。早期科学提出的不只是“它是真的吗？”还包括它是否美、是否合乎传统、是否有道德教化作用、是否通向善、是否符合真理应有的美学图式。现代科学以真理检验取代了其中大部分要求，诉诸原则上公共、可验证、道德中立的事实。它首次创造出这类事实：除真值之外不具其他价值的知识。
\end{zhblock}

\begin{enblock}
Seventeenth-century China did not take the same leap. People there considered objective knowledge without wisdom, without moral or aesthetic significance, grotesque.
\end{enblock}
\begin{zhblock}
十七世纪的中国没有完成同样的跃迁。中国人会认为没有智慧、没有道德或审美意义的客观知识是怪诞的。
\end{zhblock}

\begin{enblock}
In Europe the Scientific Revolution also redefined natural philosophy's connections to other knowledge, man's orientation toward past and future, the authority determining uses of knowledge, what knowledge of nature was socially desirable, and how knowledge should relate to individuality and to active relations between man and nature.
\end{enblock}
\begin{zhblock}
在欧洲，科学革命还重新定义了自然哲学与其他知识的关系，重新定义了人对过去和未来的取向，重新定义了何种权威应决定知识用途，何种自然知识在社会上可欲或不可欲，以及知识应如何关联人的个体性和人与自然的能动关系。
\end{zhblock}

\begin{enblock}
Galileo and his allies could not overcome Church authority by ideas alone. But they were constructing a new intellectual community outside the old establishment. The Counter-Reformation made the Church defensive and obsessed with thought control, while new careers emerged, among them the profession of scientist. This profession did not provide salaried specialist careers until about 1800, but from the start its amateurs and devotees assumed independent authority to formulate laws of nature. Secular learning remade universities and displaced ancient institutions over several centuries.
\end{enblock}
\begin{zhblock}
伽利略及其同道不能单靠观念绕过教会权威。但他们已经在旧建制之外建构新的知识共同体。反宗教改革使教会变得防御性强并沉迷于思想控制；与此同时，新职业开始出现，其中包括“科学家”这一职业。直到约1800年，这一职业才提供有薪专家生涯；但从一开始，它的业余者、信奉者和热心人就主张自己拥有制定自然法则的独立权威。世俗学问在数百年演化与革命中改造大学，并取代古老制度。
\end{zhblock}

\begin{enblock}
This outline shows how much we miss if we look only at social factors or only at intellectual factors in China. Until recently, students of the topic, including Sivin, had overlooked an important part of the Chinese picture.
\end{enblock}
\begin{zhblock}
这番概述说明，如果考察中国情形时只关注社会因素或只关注思想因素，我们会错过很多。直到不久以前，关心这一主题的人，包括席文本人，都忽略了中国图景中的一个重要部分。
\end{zhblock}

\subsection*{Scientific Revolution in China / 中国的科学革命}

\begin{enblock}
By conventional intellectual criteria, China had its own scientific revolution in the seventeenth century. Western mathematics and mathematical astronomy were introduced shortly after 1600, in a form that would soon become obsolete in parts of Europe where current knowledge was accessible. Chinese scholars quickly responded and reshaped astronomy. Geometry and trigonometry largely replaced traditional numerical or algebraic procedures; planetary rotation and relative distance from Earth became important; Chinese astronomers came to believe mathematical models could explain as well as predict phenomena. These changes amount to a conceptual revolution in astronomy.
\end{enblock}
\begin{zhblock}
按通常的思想史标准，中国在十七世纪拥有自己的科学革命。1600年稍后，西方数学和数理天文学传入中国；这种形态在欧洲若干能够接触最新知识的地区不久便会过时。中国学者迅速回应，并重塑天文学实践。几何学和三角学在很大程度上取代了传统数值或代数程序；行星绝对旋转方向及其相对地球距离首次变得重要；中国天文学家开始相信数学模型不仅能预报现象，也能解释现象。这些变化构成了一场天文学概念革命。
\end{zhblock}

\begin{enblock}
That revolution did not generate the same tension as Europe's. It did not produce as fundamental a reorientation of thought about Nature, did not cast doubt on all traditional ideas of astronomical problems, and did not narrow views of what astronomical prediction meant for understanding Nature and man's relation to it.
\end{enblock}
\begin{zhblock}
这场革命没有产生与欧洲同时期革命同等强度的张力。它没有导致对自然思想同样根本的重新定向，没有使所有传统的天文学问题观念受到怀疑，也没有缩窄人们关于天文预报对理解自然及人与自然关系究竟意味着什么的看法。
\end{zhblock}

\begin{enblock}
The most striking long-term outcome was a revival of traditional Chinese astronomy and a rediscovery of forgotten methods studied alongside new ideas, supporting a new classicism. Rather than replacing traditional values, the new values implicit in foreign astronomical writings were used to renovate traditional values.
\end{enblock}
\begin{zhblock}
这场中西相遇最引人注目的长期结果，反而是传统中国天文学的复兴，以及被遗忘方法的重新发现；这些方法与新观念一起被研究，并支撑了一种新古典主义。外来天文学著作中隐含的新价值，并没有取代传统价值，而是被用来更新传统价值。
\end{zhblock}

\begin{enblock}
Why did this conceptual revolution not have the social consequences historians of Western science expect? Old and new astronomy were not in antagonistic competition once Chinese scholars acknowledged that European techniques gave more reliable predictions. By the mid seventeenth century European civilization had little political or social impact in China, and astronomy made its way on its own merits.
\end{enblock}
\begin{zhblock}
为什么这场概念革命没有产生西方科学史家预期的社会后果？一旦中国学者承认欧洲技术能给出可靠得多的预报，新旧天文学就并非处在敌对竞争中。到十七世纪中叶，欧洲文明在中国还没有明显政治或社会影响，天文学是凭自身优点进入中国的。
\end{zhblock}

\begin{enblock}
The later rooting of Western astronomy in China may have been the last major face-to-face encounter between non-Western and European science in world history. By the eighteenth century modern science crossed boundaries on the coattails of empire, and competition among sciences, literatures, or religions on their merits became a thing of the past. Even in the mid seventeenth century, despite eclipse-prediction contests at the late Ming court, European computational techniques triumphed through an imperial decision to hand operational control of the Astronomical Bureau to Jesuit missionaries.
\end{enblock}
\begin{zhblock}
西方天文学后来在中国扎根的过程，也许是世界史上非西方科学与欧洲科学最后一次重大的面对面相遇。到十八世纪，现代科学已借帝国势力跨越国界；不同科学、文学或宗教依据各自优点相互竞争，已经成为过去。即使在十七世纪中叶，尽管晚明宫廷上演了引人注目的日食预报竞赛，欧洲计算技术的胜利也不是通过伟大心灵的共识取得，而是通过皇帝决定把钦天监的实际运行控制权交给耶稣会士取得。
\end{zhblock}

\begin{enblock}
Foreign techniques, however powerful, offered Chinese students no alternative route to security and fame. The civil service examination system left no room for one. The only astronomers who could respond to Jesuit writings were members of the old intellectual elite, who evaluated innovations in light of established ideals they felt responsible to strengthen and pass on.
\end{enblock}
\begin{zhblock}
外来技术无论多么强大，都没有给中国学生提供通向安全和名望的替代道路。科举制度没有为此留下空间。能够回应耶稣会著作的天文学家，只有旧有知识精英成员；他们必然根据既有理想来评价创新，而这些理想正是他们自认为有责任加强并传给下一代的东西。
\end{zhblock}

\begin{enblock}
Revolutions in science, like those in politics, take place at society's margins. But the people who made the seventeenth-century Chinese revolution were firmly attached to dominant cultural values. There were no students of astronomy motivated to cast off tradition, no alienated intellectual groups willing to follow ideas wherever they led even if society fell apart.
\end{enblock}
\begin{zhblock}
科学革命和政治革命一样，常发生在社会边缘。但在十七世纪中国推动这场革命的人，却牢固依附于其文化的主导价值。当时没有天文学学生以摆脱传统价值为动机，也没有足够疏离的知识群体愿意追随观念走向任何地方，即使周围社会因此瓦解。
\end{zhblock}

\begin{enblock}
The most influential first-generation champions of Western astronomy were lower Yangtze men who lived through the Manchu invasion of the 1640s. They adopted the loyalist role of refusing service to a new, non-Chinese dynasty. Having rejected conventional careers, they spent their lives studying and teaching new mathematics and astronomy while using them to master neglected techniques of their own tradition. They rejected the Ch'ing present not for a modernist future, but to keep alive the lost Ming cause for another generation.
\end{enblock}
\begin{zhblock}
第一代最有影响力的西方天文学倡导者，是经历1640年代满洲入侵的江南人士。他们采取遗民角色，拒绝服务于他们眼中的非中国新王朝。既然拒绝在已经崩坏的社会中追求常规仕途，他们便把一生用于研究和教授新数学、新天文学，同时借此掌握本国传统中被忽视的技术。他们拒绝清代现实，并不是为了现代主义未来，而是为了让明代失落事业再延续一代。
\end{zhblock}

\begin{enblock}
If we seek in China people for whom science was not a means to conservative ends, or for whom a proven fact outweighed values evolved over millennia, we do not find them until the late nineteenth century. By then foreigners exempt from Chinese law and backed by gunboats could create new institutions and careers that attracted talented young people with few other prospects. The old astronomy had nothing left to do. It became rare for modern scientists to know that their country had its own scientific tradition, until that awareness became general only in recent generations.
\end{enblock}
\begin{zhblock}
如果我们要在中国寻找那些不把科学作为保守目的之手段的人，寻找那些认为一个已证事实重于数千年演化价值的人，那么直到十九世纪后期才找得到。那时，不受中国法律约束并有炮舰支持的外国人，已经可以在中国建立新制度和新职业路径，吸引别无前途的有才青年。此时不再能谈旧天文学与新天文学的相遇；社会政治变化已使旧者无事可为。随着时间推移，现代科学家很少知道本国曾有自己的科学传统；这种意识直到最近几代才重新普及。
\end{zhblock}

\subsection*{Conclusion / 结论}

\begin{enblock}
Sivin's frustrations in making sense of science in China arise partly from the many levels of human activity that must be encompassed over a vast sweep of time, and partly because the European Scientific Revolution itself calls for broader and deeper understanding than historians have often required. Once the many dimensions of scientific change and their complex relations are kept in mind, it becomes less surprising that the Scientific Revolution happened only when and where it did. The process resembles any historical evolution: a sum of human decisions and acts, some arbitrary and many wrongheaded. We need not appeal to fate, determinism, teleology, cultural superiority, inner logic, or a hidden World Spirit.
\end{enblock}
\begin{zhblock}
席文说，他试图理解中国科学时的挫折，部分来自必须把漫长时间跨度中的许多人类活动层次纳入考察，部分来自欧洲科学革命本身也要求比其历史学家通常坚持的更宽更深的理解。一旦牢记科学变化的许多维度及其复杂关系，科学革命只在它实际发生的时间和地点发生，就不再那么令人惊讶。这个过程越来越像任何历史演化：许多人类决定和行动的总和，其中有些任意，许多糊涂。我们不必诉诸命运、决定论、目的论、文化优越性、不可抗拒展开的内在逻辑，或某种隐藏的世界精神。
\end{zhblock}

\begin{enblock}
Looking at the two scientific revolutions--the familiar one in Europe and the unexpected one in seventeenth-century China--suggests that much remains to be learned about the specific circumstances of each before we can tell the world why it could not have happened elsewhere.
\end{enblock}
\begin{zhblock}
同时考察这两场科学革命，一场是我们自以为熟悉的欧洲科学革命，另一场是十七世纪中国那场不符合我们预期的科学革命，会提示我们：在准备告诉世界它为什么不可能发生在其他时间或地点之前，我们还需要深入理解每一场革命的具体情境及其全部维度。
\end{zhblock}

\begin{enblock}
Future breakthroughs in the study of Chinese science will concern the circumstances of the people who did science and technology: how technical ideas related to the rest of their thought; who formed scientific communities; what phenomena they considered problematic; what answers they considered legitimate; how communities related to society; how they were supported; how responsibilities to colleagues and to society were reconciled; and what larger ends the sciences served.
\end{enblock}
\begin{zhblock}
中国科学研究未来的突破将属于另一种类型：它们将关乎从事科学技术的人们的具体情境。技术观念如何同他们思想的其他部分相关？科学共同体由谁构成？谁形成共识，认为某些现象成问题，某些答案合法？这些共同体如何与社会相连？如何获得支持？有知识者对科学同行的责任如何同对社会的责任调和？各门科学服务于哪些更大的目的，以至于其法则能与中国绘画法则和基本道德原则相一致？
\end{zhblock}

\begin{enblock}
These issues remain poorly understood for both China and Europe. Comparative history of science will require much more study and reflection on both sides. By then, Sivin predicts, we will no longer ask why the transition to modern science did not first take place in China.
\end{enblock}
\begin{zhblock}
无论就中国还是欧洲而言，我们对这些问题都了解甚少。比较科学史还需要双方更多研究和反思。席文预言，到那时，我们将不再追问为什么向现代科学的转变没有首先在中国发生。
\end{zhblock}

\subsection*{Appendix: Chapter Headings in Shen Kua, \emph{Brush Talks from Dream Brook} / 附录：沈括《梦溪笔谈》篇目}

\begin{longtable}{>{\raggedright\arraybackslash}p{0.12\linewidth}
                  >{\raggedright\arraybackslash}p{0.58\linewidth}
                  >{\raggedright\arraybackslash}p{0.18\linewidth}}
\toprule
No. & Heading / 篇目 & Item / 条目 \\
\midrule
1 & Ancient Usages / 古制 & 1\\
3 & Philological Criticism / 辨证、校勘 & 42\\
5 & Music and Mathematical Harmonics / 乐律与数理声学 & 82\\
7 & Regularities Underlying the Phenomena / 象数、物理之常 & 116\\
9 & Human Affairs / 人事 & 151\\
11 & Civil Service / 官政 & 189\\
13 & Wisdom in Emergencies / 权智 & 224\\
14 & Literature / 文艺 & 245\\
17 & Calligraphy and Painting / 书画 & 277\\
18 & Technical Skills / 技艺 & 298\\
20 & The Supernormal / 神奇 & 338\\
21 & Strange Occurrences / 异事 & 357\\
22 & Errors / 谬误 & 388\\
23 & Wit and Satire / 讥谑 & 401\\
24 & Miscellaneous / 杂志 & 420\\
26 & Materia Medica / 药议 & 480\\
\bottomrule
\end{longtable}

\begin{enblock}
There are 507 jottings in Brush Talks from Dream Brook, or 609 including its two sequels. For a classification of the book's contents by field of knowledge, see Joseph Needham, \emph{Science and Civilisation in China}, volume 1, page 136.
\end{enblock}
\begin{zhblock}
《梦溪笔谈》共有507则札记；若包括两种续编，则为609则。按知识领域对全书内容所作的分类，可参见李约瑟《中国科学技术史》第一卷第136页。
\end{zhblock}

\begin{figure}[htbp]
\centering
\safeincludeimage{ab8db89f822f852b2811695f7d04179e4928c24f1184703305320c6d28dbdfd8.jpg}
\caption{Appendix image preserved from source chunk. / 保留原块中的附录图像。}
\end{figure}

\begin{figure}[htbp]
\centering
\safeincludeimage{39fb42c881e228e97ba0885f8316c52c8fc1fac4d2aab457bf7362ce851e1440.jpg}
\caption{Appendix image preserved from source chunk. / 保留原块中的附录图像。}
\end{figure}

\begin{figure}[htbp]
\centering
\safeincludeimage{916233d2f39dd52cfca28f6dd4858c99111a70e4ad0a2e974036e9fc4873485b.jpg}
\caption{Appendix image preserved from source chunk. / 保留原块中的附录图像。}
\end{figure}


% ===== 09_shen_brush_talks.tex =====
% Cleaned OCR chunk: chunks/09_shen_brush_talks.md
% Scope: Shen Kuo, Mengxi Bitan excerpts only.


\section*{夢溪筆談選段}

\subsection*{卷十八　技藝}

\subsubsection*{圍棋局數}

\paragraph{英譯補譯}
說書人相傳，唐代僧人一行曾經推算圍棋所有可能局面的總數，並認為可以窮盡。我細想之後，覺得此事本不難；只是數量極大，已非世間常用的數名可以表示。姑舉其大數：二路方盤、四枚棋子，可成八十一種局面；三路方盤、九枚棋子，可成一萬九千六百八十三種；四路方盤、十六枚棋子，可成四千三百零四萬六千七百二十一種；五路方盤、二十五枚棋子，可成八千四百七十二億八千八百六十萬九千四百四十三種。七路以上，數量已大到沒有通行名稱可記。至於十九路全盤三百六十一點，局面總數約為「萬」字連書四十三次的大數。

\subsubsection*{活版印刷}

\paragraph{英譯補譯}
慶曆年間，平民畢昇創製活版印刷。他的方法是用膠泥刻字，字模薄如銅錢邊緣，每一字各成一枚泥活字，再用火燒使之堅硬。先備一塊鐵板，在板上鋪松脂、蠟和紙灰之類的混合藥料。要印刷時，把鐵範放在鐵板上，在範中密排活字；排滿一範，就成為一塊穩固的印版。隨後近火烘烤，等背面的藥料稍稍熔化，再用平板壓平字面，使整版如磨石般平整。

若只印兩三本，這方法並不算簡便；若印數百、數千本，就極為神速。通常備兩版交替使用：一版正在刷印，另一版便同時排字；前一版印完，後一版也已備好，如此輪換，印刷便十分迅捷。

每個字都有數枚活字；常用字則各備二十枚以上，以應付同一頁中重複出現的字。活字不用時，按韻部分類，用紙簽標明，收在木格中。若遇到平日未備的生僻字，就臨時刻製，用草火燒成，片刻即可使用。

不用木活字，是因為木紋有粗有細，遇水後又會高低不平；排版時還會黏在藥料上，不易取下。因此不如用燒成的陶土活字。印完後，把印版再靠近火邊，使藥料熔化，用手一拂，活字便自行脫落，絲毫不被弄髒。

畢昇去世後，他的整副活字由沈括的族侄輩收藏，至當時仍視為珍寶。

\subsection*{卷二十一　異事（異疾附）}

\subsubsection*{虹}

\paragraph{原文}
世傳虹能入溪澗飲水，信然。

熙寧中，予使契丹，至其極北黑水境永安山下卓帳。是時新雨霽，見虹下帳前澗中。予與同職扣澗觀之，虹兩頭皆垂澗中。使人過澗，隔虹對立，相去數丈，中間如隔綃縠。自西望東則見；蓋夕虹也。立澗之東西望，則為日所鑠，都無所覩。久之稍稍正東，踰山而去。次日行一程，又復見之。孫彥先云：「虹乃雨中日影也，日照雨則有之。」

\paragraph{今譯}
世人傳說彩虹能進入溪澗飲水，這話確有其事。

熙寧年間，我出使契丹，到達它極北方的黑水境，在永安山下扎營。當時雨後剛剛放晴，看見一道彩虹垂到帳前的澗水中。我和同僚走近澗邊觀看，彩虹兩端都垂入澗水。又讓人渡過溪澗，隔著彩虹與我們相對而立，相距數丈，中間好像隔著薄紗。從西向東望去就能看見，因為那是傍晚的虹；若站在溪澗東邊向西看，便被日光照得耀眼，什麼也看不見。過了很久，彩虹漸漸向正東移動，越過山嶺而去。第二天走了一程，又再次看見它。孫彥先說：「虹是雨中的日影；日光照在雨點上，就會出現虹。」

\subsection*{卷二十四　雜誌一}

\subsubsection*{海陸變遷}

\paragraph{原文}
予奉使河北，遵太行而北，山崖之間，往往銜螺蚌殼及石子如鳥卵者，橫亙石壁如帶。此乃昔之海濱，今東距海已近千里。所謂大陸者，皆濁泥所湮耳。堯殛鯀於羽山，舊說在東海中，今乃在平陸。凡大河、漳水、滹沱、涿水、桑乾之類，悉是濁流。今關、陝以西，水行地中，不減百餘尺，其泥歲東流，皆為大陸之土，此理必然。

\paragraph{今譯}
我奉命出使河北，沿太行山向北而行，看見山崖之間常常夾有螺、蚌的殼，以及像鳥卵一樣的石子，橫貫石壁，形如帶狀地層。這裡從前應是海濱，如今向東距海已接近千里。所謂大陸，其實都是渾濁水流所挾帶的泥沙沉積、淤塞而成。堯誅殺鯀於羽山，舊說羽山在東海中，如今卻已在平陸。凡黃河、漳水、滹沱河、涿水、桑乾河一類，都是含沙量很高的濁流。如今潼關、陝州以西，河水下切地中，深處不下百餘尺；其中泥沙年年向東搬運，最終成為大陸之土，這在地貌變遷上是必然的道理。

\subsubsection*{指南針}

\paragraph{原文}
方家以磁石磨針鋒，則能指南，然常微偏東，不全南也。水浮多蕩搖。指爪及盌唇上皆可為之，運轉尤速，但堅滑易墜，不若縷懸為最善。其法取新纊中獨繭縷，以芥子許蠟，綴於針腰，無風處懸之，則針常指南。其中有磨而指北者。予家指南、北者皆有之。磁石之指南，猶柏之指西，莫可原其理。

\paragraph{今譯}
方術家用磁石摩擦針尖，針便能指南；但它常常稍微偏東，並不完全指向正南，這就是磁針的偏角現象。把磁針浮在水面上，常會飄蕩搖晃；放在指甲或碗口邊緣上也可以使它轉動，而且轉得更快，但這些支點堅硬光滑，磁針容易墜落，不如用絲縷懸掛最為穩妥。方法是從新絲綿中取一根獨繭絲，用芥子大小的一點蠟黏在針腰，把針懸在無風處，針便常常指南。其中也有摩擦後指北的。我家中指南、指北的磁針都有。磁石能指南，就像柏木枝幹趨西一樣，無法推究其原理。

\paragraph{原圖保留}
\safeincludeimage{d4a8886ce638fa3464f558b8c9a1932a570a557f66d4de747a34c69ade79d9e6.jpg}
\safeincludeimage{2f6c0487a50f0282edbf4c85b5d5c7d8454054cd248a8c5db183c048c711c25b.jpg}
\safeincludeimage{9ba1872eba396ba3fcb14170391f8c57060d494d202e9ea51e79d8c0ad1c9f3e.jpg}
\safeincludeimage{7b6c8c59583d67e076070ca60124f5b97f7e1f46d61eae8062b44497cf0ee5b2.jpg}
\safeincludeimage{a93d6ec79ea3faf441e39e04d4ffc9250f86fb9a33dc11d9399658e062718e4b.jpg}
\safeincludeimage{fbb24748342b96f2ec55f761c5b08597fd7647ff4d74cc96aca3f9d03f35358b.jpg}


% ===== 10_dunham_euclid.tex =====
% Processed from chunks/10_dunham_euclid.md
% Repaired bilingual LaTeX fragment. Image paths are preserved from the source Markdown.

\section*{Dunham on Greek Geometry and Euclid / 邓纳姆论希腊几何与欧几里得}

\noindent\textbf{原文} Yet Euclid included no postulate to justify transferring lengths in this fashion. Where we expect an axiomatic dispensation allowing this procedure, we find nothing. Although his compass could draw a circle, he did not explicitly permit it to be locked into position and moved about. Euclid's is thus somewhat facetiously called a ``collapsible compass,'' a device that falls shut the instant it is lifted from the page.

\noindent\textbf{译文} 然而，欧几里得并没有列入任何公设来说明可以用这种方式转移长度。我们本以为会看到一条公理化的许可，允许这种作图操作，可是在文本中却找不到。虽然他的圆规能够画圆，他并没有明说圆规可以固定开口后被移到别处。因此，人们多少带着玩笑意味地把欧几里得的圆规称为``可塌合圆规''：这种工具一离开纸面就会立刻合拢。

\begin{figure}[htbp]
\centering
\safeincludeimage{7eaf4a037eb4c823f9b908dbe1242bbb0c9ea9ef550ceb0c57f8993e27459236.jpg}
\caption{Euclid portrait / 欧几里得肖像}
\end{figure}

\noindent\textbf{原文} The portrait caption identifies Euclid as the mathematician of Alexandria who taught under Ptolemy Lagus, and notes that his books of the \textit{Elements} were the most celebrated among works connected with music and geometry. The caption also says that the last two books were attributed to Hypsicles of Alexandria rather than to Euclid.

\noindent\textbf{译文} 肖像题注称，欧几里得是亚历山大里亚的数学家，曾在托勒密·拉古斯统治时期任教；在与音乐和几何有关的著作中，他的《几何原本》最受赞誉。题注还说，最后两卷被归于亚历山大里亚的许普西克勒斯，而不是欧几里得本人。

\noindent\textbf{原文} This raises a serious question of logic: Did the esteemed Greek geometer forget to include a ``transfer of lengths'' postulate? Do we have here a Euclidean blunder?

\noindent\textbf{译文} 这就提出了一个严肃的逻辑问题：这位备受尊崇的希腊几何学家是不是忘了加入一条``长度转移''公设？这里难道有一个欧几里得式的失误吗？

\noindent\textbf{原文} Not at all. As we shall see in a moment, Euclid had his reasons, reasons logically sound and very Greek in nature, for not including such a postulate. Rather than a blunder, this omission is evidence of his geometric acumen and organizational ability.

\noindent\textbf{译文} 完全不是。我们马上会看到，欧几里得不加入这条公设自有其理由；这些理由在逻辑上站得住脚，也非常具有希腊数学的气质。与其说这是疏漏，不如说这一省略体现了他的几何洞察力和组织能力。

\noindent\textbf{原文} With postulates in place, Euclid introduced a few ``common notions,'' self-evident statements of a more general and less geometric nature. For instance, here we accept without proof that ``things which are equal to the same thing are also equal to one another,'' that ``if equals be added to equals, the wholes are equal,'' and that ``the whole is greater than the part.'' These are statements with which few will quibble.

\noindent\textbf{译文} 在列出公设之后，欧几里得又引入若干``共同概念''，也就是一些更一般、几何色彩较少的不证自明的陈述。例如，我们在这里无需证明便承认：``等于同一事物的事物彼此相等''，``等量加等量，其和相等''，以及``整体大于部分''。这些说法很少会有人争辩。

\noindent\textbf{原文} Then he was ready to plunge in. With a huge body of geometry to deduce from a tiny collection of definitions, postulates, and common notions, where does one start? This is the sort of initial challenge known to freeze mathematicians and authors in their tracks. But if, as the Chinese tell us, a journey of a thousand miles begins with a single step, then Euclid's journey through geometry began with an equilateral triangle. The very first proposition of the \textit{Elements} was the construction, upon a given line segment, of just such a figure.

\noindent\textbf{译文} 于是他便准备正式开始。要从少量定义、公设和共同概念出发，推出庞大的几何体系，究竟该从哪里下手？这种开端问题足以让数学家和作者都停在原地。但如果像中国谚语所说，千里之行始于足下，那么欧几里得的几何之旅便始于一个等边三角形。《几何原本》的第一条命题，正是在给定线段上作出这样一个图形。

\begin{figure}[htbp]
\centering
\safeincludeimage{84f5db85c3b186e63463411260378b177a27bfe2d50cb11f9cecead81e35076b.jpg}
\caption{Figure G.1 / 图 G.1}
\end{figure}

\noindent\textbf{原文} The argument is easy. Starting with given segment $AB$ in Figure G.1, we construct a circle with center at $A$ and radius $AB$ as permitted by postulate 3. Next, with center $B$ and radius $AB$, we invoke the same postulate to construct a second circle.

\noindent\textbf{译文} 论证很简单。从图 G.1 中给定的线段 $AB$ 出发，按照公设 3，以 $A$ 为圆心、$AB$ 为半径作圆。接着仍用同一公设，以 $B$ 为圆心、$AB$ 为半径作第二个圆。

\noindent\textbf{原文} Let $C$ be the point where the circular arcs intersect. We then draw lines $AC$ and $BC$ by postulate 1, forming $\Delta ABC$. In this triangle, $AB$ and $AC$ are the same length because they are radii of the first circle; and $AB$ and $BC$ are the same length since they are radii of the second. Because things equal to the same thing are equal to each other, all three sides are mutually equal. By Euclid's definition, the triangle is equilateral and the proof is complete.

\noindent\textbf{译文} 令 $C$ 为两段圆弧的交点。再由公设 1 连接 $AC$ 和 $BC$，得到 $\Delta ABC$。在这个三角形中，$AB$ 与 $AC$ 等长，因为它们同为第一个圆的半径；$AB$ 与 $BC$ 也等长，因为它们同为第二个圆的半径。由于等于同一事物的事物彼此相等，三条边便两两相等。按照欧几里得的定义，这个三角形是等边三角形，证明完成。

\noindent\textbf{原文} It is critical to observe that in using the compass for this construction, Euclid never needed to move it rigidly. After each arc was drawn, the compass can fall to pieces without affecting the proof in the least.

\noindent\textbf{译文} 关键在于注意：在这个作图中，欧几里得从不需要把圆规保持开口不变地搬动。每画完一道圆弧，圆规即便散开，也丝毫不会影响证明。

\noindent\textbf{原文} But in the next two propositions of Book I, Euclid showed how a length could be transferred even with a compass that collapses. This means that length transfer was implied by the postulates already on the table. A new postulate to this end would have been unnecessary baggage, and Euclid was sharp enough to realize this.

\noindent\textbf{译文} 但在第一卷随后的两个命题中，欧几里得说明了即使用可塌合圆规，也可以转移长度。这意味着``长度转移''其实已经由现有公设蕴含出来。为此另加一条公设只会成为多余的负担，而欧几里得敏锐地看出了这一点。

\begin{figure}[htbp]
\centering
\safeincludeimage{fa5ce8f9fb874791de9f1c517fc1d12dd85942158f0f12f8aae1e298b3723d0a.jpg}
\caption{Figure G.2 / 图 G.2}
\end{figure}

\noindent\textbf{原文} His proofs, here combined into a single argument, were quite elegant. Suppose we have segment $AB$ as shown in Figure G.2 and wish to transfer its length onto the segment $CD$ emanating from the point $C$. First, use the straightedge and apply postulate 1 to draw the segment connecting $B$ and $C$. Then upon segment $BC$ construct equilateral triangle $BCE$; the legitimacy of this construction, of course, is exactly what was established in the previous proposition.

\noindent\textbf{译文} 他的证明相当优雅；这里把两个证明合并成一个论证。假设我们有图 G.2 所示的线段 $AB$，并希望把它的长度转移到从点 $C$ 引出的线段 $CD$ 上。先用直尺并根据公设 1 作出连接 $B$ 与 $C$ 的线段。然后在线段 $BC$ 上作等边三角形 $BCE$；当然，这种作图的合法性正是上一命题已经建立的。

\noindent\textbf{原文} We then engage in a round of circle drawing. With center $B$ and radius $AB$, construct a circle meeting $BE$ at $F$; after the compass is picked up, suppose it collapses. With center $E$ and radius $EF$, draw a circle meeting $CE$ at $G$; again, the compass may collapse as we lift it from the page. With center $C$ and radius $CG$, draw a final circle meeting $CD$ at $H$. All of these constructions are permitted by Euclid's third postulate; none require a rigid compass.

\noindent\textbf{译文} 接下来进行一连串画圆。以 $B$ 为圆心、$AB$ 为半径作圆，交 $BE$ 于 $F$；拿起圆规后，假定它塌合也无妨。再以 $E$ 为圆心、$EF$ 为半径作圆，交 $CE$ 于 $G$；离开纸面时圆规仍可塌合。最后以 $C$ 为圆心、$CG$ 为半径作圆，交 $CD$ 于 $H$。所有这些作图都由欧几里得的第三公设允许，没有一个需要刚性圆规。

\noindent\textbf{原文} Now we merely generate a string of equalities. For ease of notation, denote the length of segment $XY$ by $\overline{XY}$:
\[
\begin{array}{rl}
\overline{AB} &= \overline{BF} \quad \text{because they are radii of the same circle}\\
&= \overline{BE}-\overline{EF}\\
&= \overline{BE}-\overline{EG} \quad \text{because } EF \text{ and } EG \text{ are radii of the same circle}\\
&= \overline{CE}-\overline{EG} \quad \text{because the three sides of } \Delta BCE \text{ are equally long}\\
&= \overline{CG}\\
&= \overline{CH} \quad \text{because again we have radii of the same circle.}
\end{array}
\]

\noindent\textbf{译文} 现在只需列出一串等式。为记号方便，用 $\overline{XY}$ 表示线段 $XY$ 的长度：
\[
\begin{array}{rl}
\overline{AB} &= \overline{BF} \quad \text{因为它们是同一圆的半径}\\
&= \overline{BE}-\overline{EF}\\
&= \overline{BE}-\overline{EG} \quad \text{因为 } EF \text{ 与 } EG \text{ 是同一圆的半径}\\
&= \overline{CE}-\overline{EG} \quad \text{因为 } \Delta BCE \text{ 的三边等长}\\
&= \overline{CG}\\
&= \overline{CH} \quad \text{因为这里又是同一圆的半径。}
\end{array}
\]

\noindent\textbf{原文} The beginning and end of this chain reveal that $\overline{AB}=\overline{CH}$. Therefore the initial length of $AB$ has been transferred onto segment $CD$ as required, yet nowhere did we have to pick up a compass and move it rigidly.

\noindent\textbf{译文} 这条等式链的开头和结尾表明 $\overline{AB}=\overline{CH}$。因此，原来 $AB$ 的长度已经按要求转移到了线段 $CD$ 上，而在整个过程中，我们从未需要拿起圆规并保持开口不变地移动它。

\noindent\textbf{原文} The surprising conclusion of this proof is that constructions seeming to require a noncollapsing compass can actually be accomplished with a collapsible one. As he subsequently developed his geometry, Euclid could therefore legitimately transfer a length from one place to another as if with a rigid compass, his rationale being the just-proved theorem. By getting it out of the way so early, and so simply, he was free to use it in all that followed.

\noindent\textbf{译文} 这个证明令人惊讶的结论是：那些看起来需要不塌合圆规的作图，实际上可以用可塌合圆规完成。因此，在随后发展几何时，欧几里得可以正当地像使用刚性圆规那样把长度从一处转到另一处，其依据正是刚刚证明的定理。他如此早、如此简洁地处理掉这一点，便能在后续所有论证中自由使用它。

\noindent\textbf{原文} At this point some readers may be stifling a yawn, regarding the whole business as much ado about nothing. After all, everyone knows that stationery stores sell cheap metal compasses that are made to stay open, and surely it would have done no great damage for Euclid to have included an extra postulate to that effect.

\noindent\textbf{译文} 读到这里，有些读者也许已经强忍哈欠，觉得整件事不过是小题大做。毕竟人人都知道，文具店出售的廉价金属圆规本来就能保持开口；欧几里得即便为此多添一条公设，似乎也不会造成多大损害。

\noindent\textbf{原文} Adherents of this position, we believe, have not yet gotten into the spirit of formal Greek geometry. First, the real-world existence of rigid compasses has no bearing on the development of ideal concepts. Second, stationery stores had not yet been invented. Third, and most crucially, Euclid would not have wanted to add an unnecessary postulate to his list. Why assume something that could be derived from other assumptions? It would make his postulates less pure, less streamlined, and less perfect, thereby violating an aesthetic principle rather than a mathematical one.

\noindent\textbf{译文} 我们认为，持这种看法的人还没有进入希腊形式几何的精神。第一，现实世界中是否存在刚性圆规，与理想概念的发展无关。第二，当时还没有文具店。第三，也是最关键的一点，欧几里得不会愿意在自己的清单中加入不必要的公设。既然某件事可以由其他假设推出，为什么还要把它当作假设？那会使他的公设不够纯粹、不够简洁、不够完美；它违反的不是数学原则，而是审美原则。

\noindent\textbf{原文} That aesthetic considerations were critical for the Greek mathematician is obvious. In Euclid's proof above we glimpse what Ivor Thomas meant when he wrote:

\begin{quote}
\noindent\textbf{原文} A feature which cannot fail to impress a modern mathematician is the perfection of form in the work of the great Greek geometers. This perfection of form, which is another expression of the same genius that gave us the Parthenon and the plays of Sophocles, is found equally in the proof of individual propositions and in the ordering of those separate propositions into books; it reaches its height, perhaps, in the \textit{Elements} of Euclid.

\noindent\textbf{译文} 现代数学家必定会深受触动的一个特征，是伟大希腊几何学家作品中形式的完美。这种形式的完美，是创造出帕特农神庙和索福克勒斯戏剧的同一天才的另一种表现；它既见于单个命题的证明，也见于把各个命题编排成卷册的秩序之中；也许，它在欧几里得的《几何原本》中达到了顶峰。
\end{quote}

\noindent\textbf{原文} We now move deeper into Book I to see further evidence of Euclid's genius. After disposing of the collapsible compass in propositions 2 and 3, Euclid proved the so-called side-angle-side, or SAS, congruence scheme in proposition 4. That is, if we have triangles $ABC$ and $DEF$ in which $\overline{AB}=\overline{DE}$, $\overline{AC}=\overline{DF}$, and the included angles $\alpha=\delta$, then the triangles are congruent, meaning they have precisely the same size and shape. In other words, if $\Delta DEF$ were picked up and placed atop $\Delta ABC$, the two would coincide perfectly, line for line, angle for angle, point for point.

\noindent\textbf{译文} 现在我们深入第一卷，看看欧几里得天才的更多证据。在命题 2 和命题 3 中处理完``可塌合圆规''之后，欧几里得在命题 4 中证明了所谓的边角边，即 SAS 全等判定。也就是说，若有三角形 $ABC$ 和 $DEF$，其中 $\overline{AB}=\overline{DE}$，$\overline{AC}=\overline{DF}$，且夹角 $\alpha=\delta$，那么这两个三角形全等，意思是它们有完全相同的大小和形状。换言之，如果把 $\Delta DEF$ 拿起来放在 $\Delta ABC$ 上，两者会线对线、角对角、点对点地完全重合。

\begin{figure}[htbp]
\centering
\safeincludeimagepair{3e123e38e33ce41d4eced941a453c9d23cc1fb1f3823f6bee06e8fcfcf595e74.jpg}{2bd811776486f854a6498e3c7308b4cd40dca32e62e6d03eaf1aef000381d89b.jpg}
\caption{Figure G.3 / 图 G.3}
\end{figure}

\noindent\textbf{原文} In Euclid's hands, triangle congruence was the key to proving geometric propositions. Later he established the additional congruence patterns of side-side-side (SSS) in proposition 8, and angle-side-angle (ASA) and angle-angle-side (AAS) in proposition 26.

\noindent\textbf{译文} 在欧几里得手中，三角形全等是证明几何命题的关键。后来他又在命题 8 中建立了边边边（SSS）全等判定，并在命题 26 中建立了角边角（ASA）和角角边（AAS）全等判定。

\noindent\textbf{原文} Proposition 5 of Book I proved that the base angles of an isosceles triangle are congruent. As noted, this result has been attributed to Thales, but the proof in the \textit{Elements} is probably Euclid's own. Although we shall not give it here, we note that it was accompanied by the diagram in Figure G.4. This configuration, which suggests a bridge to vivid imaginations, may account for calling proposition 5 the \textit{pons asinorum}, or bridge of asses. According to tradition, dullards, that is, asses, find the proof beyond them and thus cannot cross this logical bridge into the geometric promised land of the \textit{Elements}.

\noindent\textbf{译文} 第一卷命题 5 证明了等腰三角形的底角全等。如前所述，这一结果曾被归于泰勒斯，但《几何原本》中的证明很可能是欧几里得自己的。虽然我们这里不列出该证明，但要注意它配有图 G.4。这个图形在想象力丰富的人看来像一座桥，这也许解释了为什么命题 5 被称为 \textit{pons asinorum}，即``愚驴之桥''。按照传统，迟钝的人，也就是所谓的``驴''，理解不了这个证明，因此无法跨过这座逻辑之桥，进入《几何原本》的几何应许之地。

\begin{figure}[htbp]
\centering
\safeincludeimage{67343c2087e15ce1a36ae4fa638c9ba61463f662a18f6217e0c85e2596f8f2a3.jpg}
\caption{Figure G.4 / 图 G.4}
\end{figure}

\noindent\textbf{原文} If weak students were likened to asses, Euclid himself met a similar fate at the hands of the Epicureans for his proof of proposition 20. To see why, we must describe a few of the intervening theorems of Book I.

\noindent\textbf{译文} 如果薄弱的学生会被比作驴，那么欧几里得本人也因为命题 20 的证明，在伊壁鸠鲁学派那里遭遇了类似的命运。要看清原因，我们必须先说明第一卷中间的几个定理。

\noindent\textbf{原文} After crossing the \textit{pons asinorum}, Euclid showed how to bisect angles and construct perpendiculars with compass and straightedge and soon arrived at one of the critical theorems of Book I, commonly known as the exterior angle theorem. This result, which appeared as proposition 16, guaranteed that the exterior angle of any triangle exceeds either of the opposite and interior angles. That is, if we begin with $\Delta ABC$ and extend side $BC$ rightward to $D$, then both $\alpha$ and $\beta$ are less than $\angle ACD$.

\noindent\textbf{译文} 跨过``愚驴之桥''后，欧几里得说明了如何用圆规和直尺作角平分线、作垂线，并很快到达第一卷的一个关键定理，通常称为外角定理。这个出现在命题 16 中的结果保证，任意三角形的外角都大于任一不相邻内角。也就是说，如果从 $\Delta ABC$ 出发，把边 $BC$ 向右延长到 $D$，那么 $\alpha$ 和 $\beta$ 都小于 $\angle ACD$。

\begin{figure}[htbp]
\centering
\safeincludeimage{9fe2f5e7a9593709c5c05a4446fb2f56e135c53e3ccd705c3040aeeb0e47476b.jpg}
\caption{Figure G.5 / 图 G.5}
\end{figure}

\noindent\textbf{原文} The exterior angle theorem was the first geometric inequality in the \textit{Elements}. Where Euclid had previously shown sides or angles to be equal, as in the \textit{pons asinorum}, he here proved certain angles to be unequal. This theorem would play a significant role in the remainder of Book I.

\noindent\textbf{译文} 外角定理是《几何原本》中的第一个几何不等式。此前欧几里得证明的是边或角相等，例如``愚驴之桥''；而在这里，他证明某些角不相等。这个定理将在第一卷余下部分发挥重要作用。

\noindent\textbf{原文} This brings us to another inequality, proposition 19, whose diagram appears in Figure G.6. Euclid stated it as, ``In any triangle the greater angle is subtended by the greater side,'' which in more modern notation becomes:

\noindent\textbf{译文} 这把我们带到另一个不等式，即命题 19，其图形见图 G.6。欧几里得把它表述为：``在任意三角形中，较大的角由较大的边所对。'' 用较现代的记号写就是：

\noindent\textbf{原文} \textbf{Proposition 19.} In $\Delta ABC$, if $\beta>\alpha$, then $\overline{AC}>\overline{BC}$, that is, $b>a$.

\noindent\textbf{译文} \textbf{命题 19。} 在 $\Delta ABC$ 中，若 $\beta>\alpha$，则 $\overline{AC}>\overline{BC}$，也就是 $b>a$。

\noindent\textbf{原文} \textbf{Proof.} Here we assume that $\beta>\alpha$. It is our job to prove that side $AC$ opposite $\angle ABC$ is longer than side $BC$ opposite $\angle BAC$.

\noindent\textbf{译文} \textbf{证明。} 这里假设 $\beta>\alpha$。我们的任务是证明，与 $\angle ABC$ 相对的边 $AC$ 长于与 $\angle BAC$ 相对的边 $BC$。

\noindent\textbf{原文} Euclid separately considered the three possible cases: $b=a$, $b<a$, and $b>a$. His strategy was to show that the first two are impossible and thereby conclude that the third case must hold, as the theorem asserted. Such a technique is called a double \textit{reductio ad absurdum}, or a proof by double contradiction. This powerful logical strategy was nowhere better employed than in Greek mathematics. Here is how Euclid handled it.

\noindent\textbf{译文} 欧几里得分别考察了三种可能情形：$b=a$、$b<a$ 和 $b>a$。他的策略是说明前两种不可能，从而得出第三种必然成立，正如定理所断言。这种技巧称为双重归谬，或双重反证。没有哪个地方比希腊数学更善于运用这种强有力的逻辑策略。欧几里得的处理如下。

\noindent\textbf{原文} \textbf{Case 1.} Suppose $b=a$.

\noindent\textbf{译文} \textbf{情形 1。} 假设 $b=a$。

\noindent\textbf{原文} Referring to Figure G.6, we have $\overline{BC}=a=b=\overline{AC}$. This makes $\Delta ABC$ isosceles and so, invoking the \textit{pons asinorum}, we conclude that the base angles are themselves congruent. That is, $\angle BAC=\angle ABC$, or equivalently $\alpha=\beta$. But this contradicts the initial assumption that $\beta>\alpha$. Hence we dismiss Case 1 as impossible.

\noindent\textbf{译文} 参照图 G.6，我们有 $\overline{BC}=a=b=\overline{AC}$。于是 $\Delta ABC$ 是等腰三角形；调用``愚驴之桥''，可知其底角全等。也就是说，$\angle BAC=\angle ABC$，等价地，$\alpha=\beta$。但这与初始假设 $\beta>\alpha$ 矛盾。因此情形 1 不可能。

\begin{figure}[htbp]
\centering
\safeincludeimage{f7e1b267dc666d8de877e4a633130d80b036cdacf94ddd4b5aea15c997f8992c.jpg}
\caption{Figure G.6 / 图 G.6}
\end{figure}

\noindent\textbf{原文} \textbf{Case 2.} Suppose $b<a$.

\noindent\textbf{译文} \textbf{情形 2。} 假设 $b<a$。

\noindent\textbf{原文} Here we have the situation depicted in Figure G.7. Because $AC$ is assumed to be shorter than $BC$, we can construct segment $CD$ of length $b$ where $D$ falls within the longer side $BC$. Then draw $AD$ to form $\Delta ADC$. This triangle, having two sides of length $b$, is isosceles and therefore has congruent base angles $\angle DAC$ and $\angle ADC$. But applying the exterior angle theorem to the narrow $\Delta ABD$, we deduce that
\[
\begin{array}{rl}
\beta &= \text{interior } \angle ABD\\
&< \text{exterior } \angle ADC \quad \text{by the exterior angle theorem}\\
&= \angle DAC \quad \text{because } \Delta DAC \text{ is isosceles}\\
&< \angle BAC \quad \text{because the whole is greater than the part}\\
&= \alpha.
\end{array}
\]

\noindent\textbf{译文} 这里的情形如图 G.7 所示。由于假设 $AC$ 短于 $BC$，我们可以在较长的边 $BC$ 内取一点 $D$，使线段 $CD$ 的长度为 $b$。再连接 $AD$，形成 $\Delta ADC$。这个三角形有两条边长为 $b$，因此是等腰三角形，从而底角 $\angle DAC$ 与 $\angle ADC$ 全等。把外角定理用于狭长的 $\Delta ABD$，得到
\[
\begin{array}{rl}
\beta &= \text{内角 } \angle ABD\\
&< \text{外角 } \angle ADC \quad \text{由外角定理}\\
&= \angle DAC \quad \text{因为 } \Delta DAC \text{ 是等腰三角形}\\
&< \angle BAC \quad \text{因为整体大于部分}\\
&= \alpha.
\end{array}
\]

\begin{figure}[htbp]
\centering
\safeincludeimage{886014fc2a084db7af8e4c02c433fe3726c033b2487e545890387c57e8b73565.jpg}
\caption{Figure G.7 / 图 G.7}
\end{figure}

\noindent\textbf{原文} In other words, $\beta<\alpha$, which contradicts the theorem's original stipulation that $\beta>\alpha$. Case 2, leading as it does to a contradiction, also bites the dust. We are left with:

\noindent\textbf{译文} 换言之，$\beta<\alpha$，这与定理原先的规定 $\beta>\alpha$ 相矛盾。情形 2 既然导致矛盾，也被排除。剩下的只有：

\noindent\textbf{原文} \textbf{Case 3.} $b>a$.

\noindent\textbf{译文} \textbf{情形 3。} $b>a$。

\noindent\textbf{原文} This must be true because no alternatives remain, and the theorem is proved.

\noindent\textbf{译文} 因为已经没有其他可能，这一情形必定为真，定理得证。

\noindent\textbf{原文} We now have reached the proposition that so troubled the Epicurean philosophers. On the surface it sounds harmless enough:

\noindent\textbf{译文} 现在我们到达了那个令伊壁鸠鲁派哲学家很不满的命题。表面上它听起来相当无害：

\noindent\textbf{原文} \textbf{Proposition 20.} In any triangle two sides taken together in any manner are greater than the remaining one.

\noindent\textbf{译文} \textbf{命题 20。} 在任意三角形中，任取两边之和都大于余下的一边。

\noindent\textbf{原文} Why the controversy? Why the derision? We quote the commentator Proclus:

\noindent\textbf{译文} 为什么会有争议？为什么会有嘲笑？我们引用评注家普罗克洛的话：

\begin{quote}
\noindent\textbf{原文} The Epicureans are wont to ridicule this theorem, saying it is evident even to an ass and needs no proof; it is as much the mark of an ignorant man, they say, to require persuasion of evident truths as to believe what is obscure without question. That the present theorem is known to an ass they make out from the observation that, if straw is placed at one extremity of the sides, an ass in quest of provender will make his way along the one side and not by way of the two others.

\noindent\textbf{译文} 伊壁鸠鲁派惯于嘲笑这个定理，说它连驴都看得出来，根本不需要证明；他们说，对显然的真理还要求别人说服自己，正如对晦暗之物不加追问就相信一样，都是无知者的标志。他们说连驴都知道这个定理，是因为观察到：如果把草放在边的一端，寻找饲料的驴会沿着一条边走过去，而不会绕另外两条边过去。
\end{quote}

\noindent\textbf{原文} In short, even a dumb animal knows to take the straight route from $C$ to $B$ in Figure G.8 rather than go the long way around via $A$. So why, the Epicureans asked, did Euclid bother to prove something so blatantly obvious? Proclus gave an answer:

\noindent\textbf{译文} 简言之，即便是不会说话的动物也知道，在图 G.8 中应走从 $C$ 到 $B$ 的直路，而不是经由 $A$ 绕远路。所以伊壁鸠鲁派问：欧几里得为什么要费心证明如此显而易见的事情？普罗克洛给出了回答：

\begin{quote}
\noindent\textbf{原文} It should be replied that, granting the theorem is evident to sense-perception, it is still not clear for scientific thought. Many things have this character; for example, that fire warms. This is clear to perception, but it is the task of science to find out how it warms.

\noindent\textbf{译文} 应当回答说，即便承认这个定理对感官知觉是明显的，它对科学思想仍然并不清楚。许多事情都有这种性质；例如火会使人感到温暖。对知觉而言这是清楚的，但科学的任务是弄清它怎样使人温暖。
\end{quote}

\begin{figure}[htbp]
\centering
\safeincludeimage{e2cd42c9578b762ce58bcb4261733970990aaec7d5d6b3e0dc01b22a793f970f.jpg}
\caption{Figure G.8 / 图 G.8}
\end{figure}

\noindent\textbf{原文} In the Euclidean spirit, the spirit that so typified Greek geometry, we must employ our faculties of reason to demonstrate that which instinct already knows. Even a seemingly self-evident proposition cries out for a proof, and this Euclid was only too happy to provide. Building upon previous results, he reasoned as follows:

\noindent\textbf{译文} 在欧几里得精神中，也就是希腊几何最典型的精神中，我们必须运用理性能力来证明本能已经知道的事情。即便一个看似不证自明的命题也呼唤证明，而欧几里得非常乐意提供这样的证明。基于先前结果，他这样论证：

\noindent\textbf{原文} \textbf{Proposition 20.} In $\Delta ABC$, $\overline{AC}+\overline{AB}>\overline{BC}$, that is, $b+c>a$.

\noindent\textbf{译文} \textbf{命题 20。} 在 $\Delta ABC$ 中，$\overline{AC}+\overline{AB}>\overline{BC}$，也就是 $b+c>a$。

\noindent\textbf{原文} \textbf{Proof.} In Figure G.9, extend side $BA$ to $D$ so that $\overline{AD}=\overline{AC}=b$ and therefore $\overline{BD}=b+c$. This construction generates $\Delta DAC$, which is isosceles because it has two sides of length $b$. Considering the large $\Delta BDC$, we note that
\[
\begin{array}{rl}
\angle BCD &> \angle ACD \quad \text{because the whole is greater than the part}\\
&= \angle BDC \quad \text{because these are base angles of isosceles } \Delta DAC.
\end{array}
\]

\noindent\textbf{译文} \textbf{证明。} 在图 G.9 中，把边 $BA$ 延长到 $D$，使 $\overline{AD}=\overline{AC}=b$，因而 $\overline{BD}=b+c$。这个作图生成 $\Delta DAC$，它有两条边长为 $b$，所以是等腰三角形。考察较大的 $\Delta BDC$，注意到
\[
\begin{array}{rl}
\angle BCD &> \angle ACD \quad \text{因为整体大于部分}\\
&= \angle BDC \quad \text{因为它们是等腰三角形 } \Delta DAC \text{ 的底角。}
\end{array}
\]

\noindent\textbf{原文} So $\angle BCD$ is greater than $\angle BDC$. As Euclid had just proved that the greater side is opposite the greater angle, it follows that $\overline{BD}>\overline{BC}$; in other words, $b+c>a$, which is exactly what was to be proved.

\noindent\textbf{译文} 因此 $\angle BCD$ 大于 $\angle BDC$。由于欧几里得刚刚证明较大角对较大边，所以可得 $\overline{BD}>\overline{BC}$；换言之，$b+c>a$，这正是所要证明的。

\begin{figure}[htbp]
\centering
\safeincludeimage{37496d2ca430aa89754c56d324b80bae2c289097df3e5919bca7dbbcce3db9de.jpg}
\caption{Figure G.9 / 图 G.9}
\end{figure}

\noindent\textbf{原文} This is a splendid little proof. It has its subtleties, its clever use of inequalities, its quiet elegance.

\noindent\textbf{译文} 这是一个精彩的小证明。它有细腻之处，有对不等式的巧妙运用，也有一种安静的优雅。

\noindent\textbf{原文} In Sir Arthur Conan Doyle's \textit{A Study in Scarlet}, Dr. Watson described the deductive powers of Sherlock Holmes in these words: ``His conclusions were as infallible as so many propositions of Euclid.'' Watson was not alone in his high opinion of the Greek geometer. Centuries ago, the Arabic scholar al-Qifti said of Euclid, ``nay, there was no one even of later date who did not walk in his footsteps,'' and the incomparable Albert Einstein added his own tribute: ``If Euclid failed to kindle your youthful enthusiasm, then you were not born to be a scientific thinker.''

\noindent\textbf{译文} 在阿瑟·柯南·道尔爵士的《血字的研究》中，华生医生这样描述夏洛克·福尔摩斯的演绎能力：``他的结论像欧几里得的诸多命题一样无懈可击。'' 并非只有华生如此推崇这位希腊几何学家。数百年前，阿拉伯学者阿尔-基夫提评价欧几里得说：``甚至在后世，也没有一个人不是沿着他的足迹前行。'' 无与伦比的阿尔伯特·爱因斯坦也献上自己的赞词：``如果欧几里得未能点燃你少年时代的热情，那么你天生就不是科学思想者。''

\noindent\textbf{原文} What we have considered thus far, of course, is just the tip of the iceberg, just a sampler of what historian Morris Kline calls the Greeks' ``grand exercise in logic.'' We must leave them at this point. But in a sense, no mathematician can leave behind the legacy of the classical geometers. They gave demonstrative mathematics its start, honed its logical tools, and aimed it in directions it has traveled ever since. We end with the words of the twentieth-century British mathematician G. H. Hardy: ``The Greeks ... spoke a language which modern mathematicians can understand; as Littlewood said to me once, they are not clever schoolboys or scholarship candidates, but Fellows of another college.''

\noindent\textbf{译文} 当然，我们迄今所考察的只是冰山一角，只是历史学家莫里斯·克莱因所谓希腊人``逻辑的宏大操练''的一个样本。我们必须在此暂别他们。但从某种意义上说，没有哪位数学家能够真正离开古典几何学家的遗产。他们开创了论证数学，磨砺了它的逻辑工具，并把它引向了此后一直行进的方向。我们以二十世纪英国数学家 G. H. 哈代的话作为结尾：``希腊人……说的是现代数学家能够理解的语言；正如利特尔伍德有一次对我说的，他们不是聪明的学生或奖学金候选人，而是另一所学院的研究员。''

\noindent\textbf{原文} There is no arguing with Hardy when he says, ``Greek mathematics is the real thing.''

\noindent\textbf{译文} 哈代说``希腊数学是真正的数学''时，我们无法反驳。

\begin{center}
\safeincludeimagepair{abcb3bf131696c0acfd75ec97f91d1522860409ebbea163aaff717510f61489b.jpg}{948c10988b3b8098710740c00c6c465fbc59bcfc37867c9b395adc7af73e16aa.jpg}\quad
\safeincludeimage{abda016953b0c39c0618f44dccbcfdfda7411c9a08e2266903d730906f97cb2e.jpg}
\end{center}

\section*{Text 11b: \textit{Elements}, Book I / 文本 11b：《几何原本》第一卷}

\noindent\textbf{原文} From \textit{Elements} by Euclid.

\noindent\textbf{译文} 选自欧几里得《几何原本》。

\subsection*{Definitions / 定义}

\begin{enumerate}
\item \textbf{原文} A point is that which has no part. \par\textbf{译文} 点是没有部分的东西。
\item \textbf{原文} A line is breadthless length. \par\textbf{译文} 线是无宽之长。
\item \textbf{原文} The extremities of a line are points. \par\textbf{译文} 线的端是点。
\item \textbf{原文} A straight line is a line which lies evenly with the points on itself. \par\textbf{译文} 直线是同其自身上的点均匀放置的线。
\item \textbf{原文} A surface is that which has length and breadth only. \par\textbf{译文} 面是只有长和宽的东西。
\item \textbf{原文} The extremities of a surface are lines. \par\textbf{译文} 面的端是线。
\item \textbf{原文} A plane surface is a surface which lies evenly with the straight lines on itself. \par\textbf{译文} 平面是同其自身上的直线均匀放置的面。
\item \textbf{原文} A plane angle is the inclination to one another of two lines in a plane which meet one another and do not lie in a straight line. \par\textbf{译文} 平面角是在同一平面内两条相交而不在同一直线上的线彼此之间的倾斜。
\item \textbf{原文} When the lines containing the angle are straight, the angle is called rectilineal. \par\textbf{译文} 当构成角的线是直线时，这个角称为直线角。
\item \textbf{原文} When a straight line set up on a straight line makes the adjacent angles equal to one another, each of the equal angles is right, and the straight line standing on the other is called a perpendicular to that on which it stands. \par\textbf{译文} 当一条直线立在另一条直线上并使相邻角彼此相等时，每一个相等的角都是直角，而立在另一条直线上的那条直线称为它所立之线的垂线。
\item \textbf{原文} An obtuse angle is an angle greater than a right angle. \par\textbf{译文} 钝角是大于直角的角。
\item \textbf{原文} An acute angle is an angle less than a right angle. \par\textbf{译文} 锐角是小于直角的角。
\item \textbf{原文} A boundary is that which is an extremity of anything. \par\textbf{译文} 界是某物的端。
\item \textbf{原文} A figure is that which is contained by any boundary or boundaries. \par\textbf{译文} 图形是由一个或多个边界所围成的东西。
\item \textbf{原文} A circle is a plane figure contained by one line such that all the straight lines falling upon it from one point among those lying within the figure are equal to one another. \par\textbf{译文} 圆是由一条线围成的平面图形；从图形内部某一点引到这条线上的所有直线彼此相等。
\item \textbf{原文} The point is called the centre of the circle. \par\textbf{译文} 这一点称为圆心。
\item \textbf{原文} A diameter of the circle is any straight line drawn through the centre and terminated in both directions by the circumference of the circle, and such a straight line also bisects the circle. \par\textbf{译文} 圆的直径是任一条通过圆心、两端都止于圆周的直线；这样的直线也平分圆。
\item \textbf{原文} A semicircle is the figure contained by the diameter and the circumference cut off by it. The centre of the semicircle is the same as that of the circle. \par\textbf{译文} 半圆是由直径及其所截取的圆周围成的图形。半圆的圆心与原圆相同。
\item \textbf{原文} Rectilineal figures are those which are contained by straight lines, trilateral figures being those contained by three, quadrilateral those contained by four, and multilateral those contained by more than four straight lines. \par\textbf{译文} 直线图形是由直线围成的图形；三边形由三条直线围成，四边形由四条直线围成，多边形由四条以上直线围成。
\item \textbf{原文} Of trilateral figures, an equilateral triangle is that which has its three sides equal, an isosceles triangle that which has two of its sides alone equal, and a scalene triangle that which has its three sides unequal. \par\textbf{译文} 在三边形中，等边三角形是三边相等的三角形；等腰三角形是只有两边相等的三角形；不等边三角形是三边都不相等的三角形。
\item \textbf{原文} Of trilateral figures, a right-angled triangle is that which has a right angle, an obtuse-angled triangle that which has an obtuse angle, and an acute-angled triangle that which has its three angles acute. \par\textbf{译文} 在三边形中，直角三角形是有一个直角的三角形；钝角三角形是有一个钝角的三角形；锐角三角形是三个角都是锐角的三角形。
\item \textbf{原文} Of quadrilateral figures, a square is that which is both equilateral and right-angled; an oblong that which is right-angled but not equilateral; a rhombus that which is equilateral but not right-angled; and a rhomboid that which has its opposite sides and angles equal to one another but is neither equilateral nor right-angled. Let quadrilaterals other than these be called trapezia. \par\textbf{译文} 在四边形中，正方形是等边且直角的四边形；长方形是直角但不等边的四边形；菱形是等边但非直角的四边形；菱面形是对边和对角彼此相等、但既非等边也非直角的四边形。除此之外的四边形称为梯形。
\item \textbf{原文} Parallel straight lines are straight lines which, being in the same plane and being produced indefinitely in both directions, do not meet one another in either direction. \par\textbf{译文} 平行直线是在同一平面内、向两个方向无限延长后在任一方向都不相交的直线。
\end{enumerate}

\subsection*{Postulates / 公设}

\noindent\textbf{原文} Let the following be postulated:

\noindent\textbf{译文} 设定以下公设：

\begin{enumerate}
\item \textbf{原文} To draw a straight line from any point to any point. \par\textbf{译文} 从任一点到任一点作一直线。
\item \textbf{原文} To produce a finite straight line continuously in a straight line. \par\textbf{译文} 将有限直线沿直线连续延长。
\item \textbf{原文} To describe a circle with any centre and distance. \par\textbf{译文} 以任意圆心和距离作圆。
\item \textbf{原文} That all right angles are equal to one another. \par\textbf{译文} 所有直角彼此相等。
\item \textbf{原文} That, if a straight line falling on two straight lines make the interior angles on the same side less than two right angles, the two straight lines, if produced indefinitely, meet on that side on which are the angles less than the two right angles. \par\textbf{译文} 若一条直线交两条直线，并使同侧内角之和小于两个直角，则这两条直线若无限延长，会在内角之和小于两个直角的那一侧相交。
\end{enumerate}

\subsection*{Common Notions / 公理（共同概念）}

\begin{enumerate}
\item \textbf{原文} Things which are equal to the same thing are also equal to one another. \par\textbf{译文} 等于同一事物的事物彼此相等。
\item \textbf{原文} If equals be added to equals, the wholes are equal. \par\textbf{译文} 等量加等量，其和相等。
\item \textbf{原文} If equals be subtracted from equals, the remainders are equal. \par\textbf{译文} 等量减等量，其余量相等。
\item \textbf{原文} Things which coincide with one another are equal to one another. \par\textbf{译文} 彼此重合的事物彼此相等。
\item \textbf{原文} The whole is greater than the part. \par\textbf{译文} 整体大于部分。
\end{enumerate}

\subsection*{Propositions / 命题}

\subsubsection*{Proposition 1 / 命题 1}

\noindent\textbf{原文} On a given finite straight line to construct an equilateral triangle.

\noindent\textbf{译文} 在给定有限直线上作一个等边三角形。

\noindent\textbf{原文} Let $AB$ be the given finite straight line. Thus it is required to construct an equilateral triangle on the straight line $AB$. With centre $A$ and distance $AB$ let the circle $BCD$ be described; [Post. 3] again, with centre $B$ and distance $BA$ let the circle $ACE$ be described; [Post. 3] and from the point $C$, in which the circles cut one another, to the points $A$, $B$ let the straight lines $CA$, $CB$ be joined. [Post. 1]

\noindent\textbf{译文} 设 $AB$ 为给定的有限直线。因此要求在直线 $AB$ 上作一个等边三角形。以 $A$ 为圆心、$AB$ 为距离作圆 $BCD$；[公设 3] 再以 $B$ 为圆心、$BA$ 为距离作圆 $ACE$；[公设 3] 两圆相交于点 $C$，从点 $C$ 到点 $A$、$B$ 连接直线 $CA$、$CB$。[公设 1]

\begin{figure}[htbp]
\centering
\safeincludeimage{7f061bad1266dd45dcb307a48af1a0707f44763a1fbfbd168a0d78906e1cf15d.jpg}
\caption{Proposition 1 diagram / 命题 1 图}
\end{figure}

\noindent\textbf{原文} Now, since the point $A$ is the centre of the circle $CDB$, $AC$ is equal to $AB$. [Def. 15] Again, since the point $B$ is the centre of the circle $CAE$, $BC$ is equal to $BA$. [Def. 15] But $CA$ was also proved equal to $AB$; therefore each of the straight lines $CA$, $CB$ is equal to $AB$. And things which are equal to the same thing are also equal to one another; [C.N. 1] therefore $CA$ is also equal to $CB$. Therefore the three straight lines $CA$, $AB$, $BC$ are equal to one another. Therefore the triangle $ABC$ is equilateral, and it has been constructed on the given finite straight line $AB$. Being what it was required to do.

\noindent\textbf{译文} 现在，因为点 $A$ 是圆 $CDB$ 的圆心，所以 $AC$ 等于 $AB$。[定义 15] 又因为点 $B$ 是圆 $CAE$ 的圆心，所以 $BC$ 等于 $BA$。[定义 15] 但已证明 $CA$ 也等于 $AB$；因此直线 $CA$、$CB$ 各自都等于 $AB$。而等于同一事物的事物彼此相等；[共同概念 1] 所以 $CA$ 也等于 $CB$。因此三条直线 $CA$、$AB$、$BC$ 彼此相等。于是三角形 $ABC$ 是等边三角形，并且它已经作在给定有限直线 $AB$ 上。作图完毕。

\subsubsection*{Proposition 2 / 命题 2}

\noindent\textbf{原文} To place at a given point as an extremity a straight line equal to a given straight line.

\noindent\textbf{译文} 在给定点处，以该点为端点，作一条等于给定直线的直线。

\begin{figure}[htbp]
\centering
\safeincludeimage{b4e7f322f523fe70254e116490df11c37d34cc0bd8d8545e455ccb1eac3c7c8b.jpg}
\caption{Proposition 2 diagram / 命题 2 图}
\end{figure}

\noindent\textbf{原文} Let $A$ be the given point, and $BC$ the given straight line. Thus it is required to place at the point $A$ as an extremity a straight line equal to the given straight line $BC$. From the point $A$ to the point $B$ let the straight line $AB$ be joined; [Post. 1] and on it let the equilateral triangle $DAB$ be constructed. [I. 1] Let the straight lines $AE$, $BF$ be produced in a straight line with $DA$, $DB$; [Post. 2] with centre $B$ and distance $BC$ let the circle $CGH$ be described; [Post. 3] and again, with centre $D$ and distance $DG$ let the circle $GKL$ be described. [Post. 3]

\noindent\textbf{译文} 设 $A$ 为给定点，$BC$ 为给定直线。因此要求在点 $A$ 处，以 $A$ 为端点，作一条等于给定直线 $BC$ 的直线。从点 $A$ 到点 $B$ 连接直线 $AB$；[公设 1] 并在其上作等边三角形 $DAB$。[I. 1] 将直线 $AE$、$BF$ 分别沿 $DA$、$DB$ 的方向延长成直线；[公设 2] 以 $B$ 为圆心、$BC$ 为距离作圆 $CGH$；[公设 3] 再以 $D$ 为圆心、$DG$ 为距离作圆 $GKL$。[公设 3]

\noindent\textbf{原文} Then, since the point $B$ is the centre of the circle $CGH$, $BC$ is equal to $BG$. Again, since the point $D$ is the centre of the circle $GKL$, $DL$ is equal to $DG$. And in these $DA$ is equal to $DB$; therefore the remainder $AL$ is equal to the remainder $BG$. [C.N. 3] But $BC$ was also proved equal to $BG$; therefore each of the straight lines $AL$, $BC$ is equal to $BG$. And things which are equal to the same thing are also equal to one another; [C.N. 1] therefore $AL$ is also equal to $BC$. Therefore at the given point $A$ the straight line $AL$ is placed equal to the given straight line $BC$. Being what it was required to do.

\noindent\textbf{译文} 于是，因为点 $B$ 是圆 $CGH$ 的圆心，所以 $BC$ 等于 $BG$。又因为点 $D$ 是圆 $GKL$ 的圆心，所以 $DL$ 等于 $DG$。而其中 $DA$ 等于 $DB$；因此余下的 $AL$ 等于余下的 $BG$。[共同概念 3] 但已证明 $BC$ 也等于 $BG$；所以直线 $AL$、$BC$ 各自都等于 $BG$。等于同一事物的事物彼此相等；[共同概念 1] 因此 $AL$ 也等于 $BC$。所以在给定点 $A$ 处，已经作出直线 $AL$，使它等于给定直线 $BC$。作图完毕。

\subsubsection*{Proposition 3 / 命题 3}

\noindent\textbf{原文} Given two unequal straight lines, to cut off from the greater a straight line equal to the less.

\noindent\textbf{译文} 给定两条不相等的直线，从较大者上截取一条等于较小者的直线。

\noindent\textbf{原文} Let $AB$, $C$ be the two given unequal straight lines, and let $AB$ be the greater of them. Thus it is required to cut off from $AB$ the greater a straight line equal to $C$ the less. At the point $A$ let $AD$ be placed equal to the straight line $C$; [I. 2] and with centre $A$ and distance $AD$ let the circle $DEF$ be described. [Post. 3]

\noindent\textbf{译文} 设 $AB$、$C$ 为给定的两条不相等直线，并设 $AB$ 为其中较大者。因此要求从较大的 $AB$ 上截取一条等于较小的 $C$ 的直线。在点 $A$ 处作 $AD$ 等于直线 $C$；[I. 2] 并以 $A$ 为圆心、$AD$ 为距离作圆 $DEF$。[公设 3]

\begin{figure}[htbp]
\centering
\safeincludeimage{b6842ec6cac3d03b97eeae0e44cf38e34f337eaf8e080ecf2e6add1f7e90b52d.jpg}
\caption{Proposition 3 diagram / 命题 3 图}
\end{figure}

\noindent\textbf{原文} Now, since the point $A$ is the centre of the circle $DEF$, $AE$ is equal to $AD$. [Def. 15] But $C$ is also equal to $AD$. Therefore each of the straight lines $AE$, $C$ is equal to $AD$; so that $AE$ is also equal to $C$. [C.N. 1] Therefore, given the two straight lines $AB$, $C$, from $AB$ the greater, $AE$ has been cut off equal to $C$ the less. Being what it was required to do.

\noindent\textbf{译文} 现在，因为点 $A$ 是圆 $DEF$ 的圆心，所以 $AE$ 等于 $AD$。[定义 15] 但 $C$ 也等于 $AD$。因此直线 $AE$ 与 $C$ 各自都等于 $AD$；所以 $AE$ 也等于 $C$。[共同概念 1] 因此，给定直线 $AB$、$C$ 时，已从较大的 $AB$ 上截取 $AE$，使其等于较小的 $C$。作图完毕。

\subsubsection*{Proposition 4 / 命题 4}

\noindent\textbf{原文} If two triangles have the two sides equal to two sides respectively, and have the angles contained by the equal straight lines equal, they will also have the base equal to the base, the triangle will be equal to the triangle, and the remaining angles will be equal to the remaining angles respectively, namely those which the equal sides subtend.

\noindent\textbf{译文} 若两个三角形有两边分别等于两边，且由这些相等直线所夹的角相等，则它们的底也等于底，三角形也等于三角形，其余各角也分别相等，也就是相等边所对的那些角相等。

\noindent\textbf{原文} Let $ABC$, $DEF$ be two triangles having the two sides $AB$, $AC$ equal to the two sides $DE$, $DF$ respectively, namely $AB$ to $DE$ and $AC$ to $DF$, and the angle $BAC$ equal to the angle $EDF$.

\noindent\textbf{译文} 设 $ABC$、$DEF$ 为两个三角形，其中两边 $AB$、$AC$ 分别等于两边 $DE$、$DF$，即 $AB$ 等于 $DE$、$AC$ 等于 $DF$，并且角 $BAC$ 等于角 $EDF$。

\begin{figure}[htbp]
\centering
\safeincludeimagepair{080368f507e9eeb9ca14bad79c3117321aca4e45fc4883ad9766bea1bbf85251.jpg}{3c3f07b8219daea331169214d575bfc595f88e160f9bab12d7f1619b66c2e404.jpg}
\caption{Proposition 4 diagram / 命题 4 图}
\end{figure}

\noindent\textbf{原文} I say that the base $BC$ is also equal to the base $EF$, the triangle $ABC$ will be equal to the triangle $DEF$, and the remaining angles will be equal to the remaining angles respectively, namely those which the equal sides subtend, that is, the angle $ABC$ to the angle $DEF$, and the angle $ACB$ to the angle $DFE$.

\noindent\textbf{译文} 我说，底 $BC$ 也等于底 $EF$，三角形 $ABC$ 等于三角形 $DEF$，其余各角也分别相等，即相等边所对的角相等，也就是角 $ABC$ 等于角 $DEF$，角 $ACB$ 等于角 $DFE$。

\noindent\textbf{原文} For, if the triangle $ABC$ be applied to the triangle $DEF$, and if the point $A$ be placed on the point $D$ and the straight line $AB$ on $DE$, then the point $B$ will also coincide with $E$, because $AB$ is equal to $DE$. Again, $AB$ coinciding with $DE$, the straight line $AC$ will also coincide with $DF$, because the angle $BAC$ is equal to the angle $EDF$; hence the point $C$ will also coincide with the point $F$, because $AC$ is again equal to $DF$.

\noindent\textbf{译文} 因为，如果把三角形 $ABC$ 叠合到三角形 $DEF$ 上，使点 $A$ 放在点 $D$ 上，直线 $AB$ 放在 $DE$ 上，那么点 $B$ 也将与 $E$ 重合，因为 $AB$ 等于 $DE$。再者，当 $AB$ 与 $DE$ 重合时，直线 $AC$ 也将与 $DF$ 重合，因为角 $BAC$ 等于角 $EDF$；于是点 $C$ 也将与点 $F$ 重合，因为 $AC$ 又等于 $DF$。

\noindent\textbf{原文} But $B$ also coincided with $E$; hence the base $BC$ will coincide with the base $EF$ and will be equal to it. [C.N. 4] Thus the whole triangle $ABC$ will coincide with the whole triangle $DEF$ and will be equal to it. And the remaining angles will also coincide with the remaining angles and will be equal to them, the angle $ABC$ to the angle $DEF$, and the angle $ACB$ to the angle $DFE$. Therefore etc. Being what it was required to prove.

\noindent\textbf{译文} 但 $B$ 也已经与 $E$ 重合；因此底 $BC$ 将与底 $EF$ 重合，并且等于它。[共同概念 4] 于是整个三角形 $ABC$ 将与整个三角形 $DEF$ 重合，并且等于它。其余各角也将与其余各角重合并相等，即角 $ABC$ 等于角 $DEF$，角 $ACB$ 等于角 $DFE$。所以等等。证毕。

\subsubsection*{Proposition 5 / 命题 5}

\noindent\textbf{原文} In isosceles triangles the angles at the base are equal to one another; and, if the equal straight lines be produced further, the angles under the base will be equal to one another.

\noindent\textbf{译文} 在等腰三角形中，底上的两角彼此相等；并且，如果相等的直线继续延长，则底下的两角也彼此相等。

\noindent\textbf{原文} Let $ABC$ be an isosceles triangle having the side $AB$ equal to the side $AC$; and let the straight lines $BD$, $CE$ be produced further in a straight line with $AB$, $AC$. [Post. 2] I say that the angle $ABC$ is equal to the angle $ACB$, and the angle $CBD$ to the angle $BCE$.

\noindent\textbf{译文} 设 $ABC$ 为等腰三角形，其中边 $AB$ 等于边 $AC$；并将直线 $BD$、$CE$ 分别沿 $AB$、$AC$ 继续延长成直线。[公设 2] 我说，角 $ABC$ 等于角 $ACB$，且角 $CBD$ 等于角 $BCE$。

\noindent\textbf{原文} Let a point $F$ be taken at random on $BD$; from $AE$ the greater let $AG$ be cut off equal to $AF$ the less; [I. 3] and let the straight lines $FC$, $GB$ be joined. [Post. 1] Then, since $AF$ is equal to $AG$ and $AB$ to $AC$, the two sides $FA$, $AC$ are equal to the two sides $GA$, $AB$, respectively; and they contain a common angle, the angle $FAG$.

\noindent\textbf{译文} 在 $BD$ 上任取一点 $F$；从较大的 $AE$ 上截取 $AG$，使其等于较小的 $AF$；[I. 3] 并连接直线 $FC$、$GB$。[公设 1] 于是，由于 $AF$ 等于 $AG$，且 $AB$ 等于 $AC$，两边 $FA$、$AC$ 分别等于两边 $GA$、$AB$；并且它们夹着公共角，即角 $FAG$。

\begin{figure}[htbp]
\centering
\safeincludeimage{6b61f2fdc0064be390f614c6c8bc31d5431a4c86b8c9c1ecf3f622e794bc177f.jpg}
\caption{Proposition 5 diagram / 命题 5 图}
\end{figure}

\noindent\textbf{原文} Therefore the base $FC$ is equal to the base $GB$, and the triangle $AFC$ is equal to the triangle $AGB$, and the remaining angles will be equal to the remaining angles respectively, namely those which the equal sides subtend, that is, the angle $ACF$ to the angle $ABG$, and the angle $AFC$ to the angle $AGB$. [I. 4]

\noindent\textbf{译文} 因此底 $FC$ 等于底 $GB$，三角形 $AFC$ 等于三角形 $AGB$，其余各角也分别相等，即相等边所对的角相等，也就是角 $ACF$ 等于角 $ABG$，角 $AFC$ 等于角 $AGB$。[I. 4]

\noindent\textbf{原文} And, since the whole $AF$ is equal to the whole $AG$, and in these $AB$ is equal to $AC$, the remainder $BF$ is equal to the remainder $CG$. But $FC$ was also proved equal to $GB$; therefore the two sides $BF$, $FC$ are equal to the two sides $CG$, $GB$ respectively; and the angle $BFC$ is equal to the angle $CGB$, while the base $BC$ is common to them; therefore the triangle $BFC$ is also equal to the triangle $CGB$, and the remaining angles will be equal to the remaining angles respectively, namely those which the equal sides subtend; therefore the angle $FBC$ is equal to the angle $GCB$, and the angle $BCF$ to the angle $CBG$.

\noindent\textbf{译文} 又因为整个 $AF$ 等于整个 $AG$，且其中 $AB$ 等于 $AC$，所以余下的 $BF$ 等于余下的 $CG$。但已证明 $FC$ 等于 $GB$；因此两边 $BF$、$FC$ 分别等于两边 $CG$、$GB$；并且角 $BFC$ 等于角 $CGB$，而底 $BC$ 为二者共有；所以三角形 $BFC$ 也等于三角形 $CGB$，其余各角也分别相等，即相等边所对的角相等；因此角 $FBC$ 等于角 $GCB$，角 $BCF$ 等于角 $CBG$。

\noindent\textbf{原文} Accordingly, since the whole angle $ABG$ was proved equal to the angle $ACF$, and in these the angle $CBG$ is equal to the angle $BCF$, the remaining angle $ABC$ is equal to the remaining angle $ACB$; and they are at the base of the triangle $ABC$. But the angle $FBC$ was also proved equal to the angle $GCB$; and they are under the base. Therefore etc. Q.E.D.

\noindent\textbf{译文} 因此，既然已证明整个角 $ABG$ 等于角 $ACF$，并且其中角 $CBG$ 等于角 $BCF$，则余下的角 $ABC$ 等于余下的角 $ACB$；它们是三角形 $ABC$ 的底角。但也已经证明角 $FBC$ 等于角 $GCB$；它们是底下的角。所以等等。证毕。

\noindent\textbf{原文} [\ldots]

\noindent\textbf{译文} [原文此处省略]

\subsubsection*{Proposition 7 / 命题 7}

\noindent\textbf{原文} Given two straight lines constructed on a straight line from its extremities and meeting in a point, there cannot be constructed on the same straight line from its extremities, and on the same side of it, two other straight lines meeting in another point and equal to the former two respectively, namely each to that which has the same extremity with it.

\noindent\textbf{译文} 若从一条直线的两个端点作两条直线并交于一点，则不能在同一条直线的同一侧、仍从其两个端点再作另外两条直线，使它们交于另一点，并且分别等于前两条直线，也就是每一条都等于与它有同一端点的那一条。

\begin{figure}[htbp]
\centering
\safeincludeimage{828e9536c5eaa73dcb439a1abef7a7741489914511ed1091710b81aac5f2cffe.jpg}
\caption{Proposition 7 diagram / 命题 7 图}
\end{figure}

\noindent\textbf{原文} For, if possible, given two straight lines $AC$, $CB$ constructed on the straight line $AB$ and meeting at the point $C$, let two other straight lines $AD$, $DB$ be constructed on the same straight line $AB$, on the same side of it, meeting in another point $D$ and equal to the former two respectively, namely each to that which has the same extremity with it, so that $CA$ is equal to $DA$ which has the same extremity $A$ with it, and $CB$ to $DB$ which has the same extremity $B$ with it; and let $CD$ be joined.

\noindent\textbf{译文} 因为，若可能，给定在直线 $AB$ 上从其端点作出的两条直线 $AC$、$CB$，它们交于点 $C$；设还能在同一条直线 $AB$ 的同一侧作另外两条直线 $AD$、$DB$，它们交于另一点 $D$，并分别等于前两条直线，也就是各自等于与它有同一端点的那一条：$CA$ 等于同以 $A$ 为端点的 $DA$，$CB$ 等于同以 $B$ 为端点的 $DB$；再连接 $CD$。

\noindent\textbf{原文} Then, since $AC$ is equal to $AD$, the angle $ACD$ is also equal to the angle $ADC$. [I. 5] Therefore the angle $ADC$ is greater than the angle $DCB$; therefore the angle $CDB$ is much greater than the angle $DCB$. Again, since $CB$ is equal to $DB$, the angle $CDB$ is also equal to the angle $DCB$. But it was also proved much greater than it, which is impossible. Therefore etc. Q.E.D.

\noindent\textbf{译文} 于是，由于 $AC$ 等于 $AD$，角 $ACD$ 也等于角 $ADC$。[I. 5] 因此角 $ADC$ 大于角 $DCB$；所以角 $CDB$ 更大于角 $DCB$。又由于 $CB$ 等于 $DB$，角 $CDB$ 也等于角 $DCB$。但又已证明它远大于后者，这是不可能的。所以等等。证毕。

\subsubsection*{Proposition 8 / 命题 8}

\noindent\textbf{原文} If two triangles have the two sides equal to two sides respectively, and have also the base equal to the base, they will also have the angles equal which are contained by the equal straight lines.

\noindent\textbf{译文} 若两个三角形有两边分别等于两边，并且底也等于底，则由这些相等直线所夹的角也相等。

\noindent\textbf{原文} Let $ABC$, $DEF$ be two triangles having the two sides $AB$, $AC$ equal to the two sides $DE$, $DF$ respectively, namely $AB$ to $DE$, and $AC$ to $DF$; and let them have the base $BC$ equal to the base $EF$. I say that the angle $BAC$ is also equal to the angle $EDF$.

\noindent\textbf{译文} 设 $ABC$、$DEF$ 为两个三角形，其中两边 $AB$、$AC$ 分别等于两边 $DE$、$DF$，即 $AB$ 等于 $DE$，$AC$ 等于 $DF$；并且它们的底 $BC$ 等于底 $EF$。我说，角 $BAC$ 也等于角 $EDF$。

\noindent\textbf{原文} For, if the triangle $ABC$ be applied to the triangle $DEF$, and if the point $B$ be placed on the point $E$ and the straight line $BC$ on $EF$, the point $C$ will also coincide with $F$, because $BC$ is equal to $EF$.

\noindent\textbf{译文} 因为，如果把三角形 $ABC$ 叠合到三角形 $DEF$ 上，使点 $B$ 放在点 $E$ 上，并使直线 $BC$ 放在 $EF$ 上，那么点 $C$ 也将与 $F$ 重合，因为 $BC$ 等于 $EF$。

\begin{figure}[htbp]
\centering
\safeincludeimage{cc70ef0f787dd782c274c6403b90172997502bd074a5e5a5febde01f717b00f6.jpg}
\caption{Proposition 8 diagram / 命题 8 图}
\end{figure}

\noindent\textbf{原文} Then, $BC$ coinciding with $EF$, $BA$, $AC$ will also coincide with $ED$, $DF$; for, if the base $BC$ coincides with the base $EF$, and the sides $BA$, $AC$ do not coincide with $ED$, $DF$ but fall beside them as $EG$, $GF$, then, given two straight lines constructed on a straight line from its extremities and meeting in a point, there will have been constructed on the same straight line from its extremities, and on the same side of it, two other straight lines meeting in another point and equal to the former two respectively, namely each to that which has the same extremity with it. But they cannot be so constructed. [I. 7]

\noindent\textbf{译文} 然后，当 $BC$ 与 $EF$ 重合时，$BA$、$AC$ 也将与 $ED$、$DF$ 重合；因为，如果底 $BC$ 与底 $EF$ 重合，而边 $BA$、$AC$ 不与 $ED$、$DF$ 重合，却落在旁边成为 $EG$、$GF$，那么就会出现这样的情况：从一条直线的两个端点作两条直线并交于一点之后，又在同一条直线的同一侧、从其两个端点作出另外两条直线，交于另一点，并分别等于前两条直线，也就是每条都等于与它有同一端点的那条。但这样的作图不可能。[I. 7]

\noindent\textbf{原文} Therefore it is not possible that, if the base $BC$ be applied to the base $EF$, the sides $BA$, $AC$ should not coincide with $ED$, $DF$; they will therefore coincide, so that the angle $BAC$ will also coincide with the angle $EDF$, and will be equal to it. If therefore etc. Q.E.D.

\noindent\textbf{译文} 因此，如果底 $BC$ 叠合到底 $EF$ 上，边 $BA$、$AC$ 不与 $ED$、$DF$ 重合是不可能的；所以它们将重合，从而角 $BAC$ 也将与角 $EDF$ 重合，并且等于它。所以等等。证毕。

\subsubsection*{Proposition 9 / 命题 9}

\noindent\textbf{原文} To bisect a given rectilineal angle.

\noindent\textbf{译文} 平分给定直线角。

\begin{figure}[htbp]
\centering
\safeincludeimage{c15463dc3ef129c034e9bda5835bcfe11a3503c0ed6d14430e248712df29c268.jpg}
\caption{Proposition 9 diagram / 命题 9 图}
\end{figure}

\noindent\textbf{原文} Let the angle $BAC$ be the given rectilineal angle. Thus it is required to bisect it. Let a point $D$ be taken at random on $AB$; let $AE$ be cut off from $AC$ equal to $AD$; [I. 3] let $DE$ be joined, and on $DE$ let the equilateral triangle $DEF$ be constructed; let $AF$ be joined. I say that the angle $BAC$ has been bisected by the straight line $AF$.

\noindent\textbf{译文} 设角 $BAC$ 为给定直线角。因此要求平分它。在 $AB$ 上任取一点 $D$；从 $AC$ 上截取 $AE$，使其等于 $AD$；[I. 3] 连接 $DE$，并在 $DE$ 上作等边三角形 $DEF$；连接 $AF$。我说，角 $BAC$ 已被直线 $AF$ 平分，$AF$ 即为角平分线。

\noindent\textbf{原文} For, since $AD$ is equal to $AE$, and $AF$ is common, the two sides $DA$, $AF$ are equal to the two sides $EA$, $AF$ respectively. And the base $DF$ is equal to the base $EF$; therefore the angle $DAF$ is equal to the angle $EAF$. [I. 8] Therefore the given rectilineal angle $BAC$ has been bisected by the straight line $AF$. Q.E.F.

\noindent\textbf{译文} 因为 $AD$ 等于 $AE$，且 $AF$ 为公共边，所以两边 $DA$、$AF$ 分别等于两边 $EA$、$AF$。又底 $DF$ 等于底 $EF$；因此角 $DAF$ 等于角 $EAF$。[I. 8] 所以给定直线角 $BAC$ 已由直线 $AF$ 平分。作图完毕。

\subsubsection*{Proposition 10 / 命题 10}

\noindent\textbf{原文} To bisect a given finite straight line.

\noindent\textbf{译文} 平分给定有限直线。

\begin{figure}[htbp]
\centering
\safeincludeimage{344770768dd08da63064ff9c20f95022735918a2610550ba8f0f158d481aafb2.jpg}
\caption{Proposition 10 diagram / 命题 10 图}
\end{figure}

\noindent\textbf{原文} Let $AB$ be the given finite straight line. Thus it is required to bisect the finite straight line $AB$. Let the equilateral triangle $ABC$ be constructed on it, [I. 1] and let the angle $ACB$ be bisected by the straight line $CD$; [I. 9] I say that the straight line $AB$ has been bisected at the point $D$.

\noindent\textbf{译文} 设 $AB$ 为给定有限直线。因此要求平分有限直线 $AB$。在其上作等边三角形 $ABC$，[I. 1] 并由直线 $CD$ 平分角 $ACB$；[I. 9] 我说，直线 $AB$ 已在点 $D$ 处被平分。

\noindent\textbf{原文} For, since $AC$ is equal to $CB$, and $CD$ is common, the two sides $AC$, $CD$ are equal to the two sides $BC$, $CD$ respectively; and the angle $ACD$ is equal to the angle $BCD$; [I. 4] therefore the base $AD$ is equal to the base $BD$. Therefore the given finite straight line $AB$ has been bisected at $D$. Q.E.F.

\noindent\textbf{译文} 因为 $AC$ 等于 $CB$，且 $CD$ 为公共边，所以两边 $AC$、$CD$ 分别等于两边 $BC$、$CD$；并且角 $ACD$ 等于角 $BCD$；[I. 4] 因此底 $AD$ 等于底 $BD$。所以给定有限直线 $AB$ 已在 $D$ 处被平分。作图完毕。

\subsubsection*{Proposition 11 / 命题 11}

\noindent\textbf{原文} To draw a straight line at right angles to a given straight line from a given point on it.

\noindent\textbf{译文} 从给定直线上的给定点，作一条与该给定直线成直角的直线。

\noindent\textbf{原文} Let $AB$ be the given straight line, and $C$ the given point on it. Thus it is required to draw from the point $C$ a straight line at right angles to the straight line $AB$. Let a point $D$ be taken at random on $AC$; let $CE$ be made equal to $CD$; [I. 3] on $DE$ let the equilateral triangle $FDE$ be constructed, [I. 1] and let $FC$ be joined.

\noindent\textbf{译文} 设 $AB$ 为给定直线，$C$ 为其上的给定点。因此要求从点 $C$ 作一条与直线 $AB$ 成直角的直线。在 $AC$ 上任取一点 $D$；令 $CE$ 等于 $CD$；[I. 3] 在 $DE$ 上作等边三角形 $FDE$，[I. 1] 并连接 $FC$。

\begin{figure}[htbp]
\centering
\safeincludeimage{b9a58f36da06f5759bed049e2bb3826aa61e46c8c908e86e50ce69ade03b3e96.jpg}
\caption{Proposition 11 diagram / 命题 11 图}
\end{figure}

\noindent\textbf{原文} I say that the straight line $FC$ has been drawn at right angles to the given straight line $AB$ from $C$, the given point on it. For, since $DC$ is equal to $CE$, and $CF$ is common, the two sides $DC$, $CF$ are equal to the two sides $EC$, $CF$ respectively; and the base $DF$ is equal to the base $FE$; therefore the angle $DCF$ is equal to the angle $ECF$; [I. 8] and they are adjacent angles.

\noindent\textbf{译文} 我说，已经从给定直线 $AB$ 上的给定点 $C$ 作出与 $AB$ 成直角的直线 $FC$。因为 $DC$ 等于 $CE$，且 $CF$ 为公共边，所以两边 $DC$、$CF$ 分别等于两边 $EC$、$CF$；并且底 $DF$ 等于底 $FE$；因此角 $DCF$ 等于角 $ECF$；[I. 8] 而它们是相邻角。

\noindent\textbf{原文} But, when a straight line set up on a straight line makes the adjacent angles equal to one another, each of the equal angles is right; [Def. 10] therefore each of the angles $DCF$, $FCE$ is right. Therefore the straight line $CF$ has been drawn at right angles to the given straight line $AB$ from the given point $C$ on it. Q.E.F.

\noindent\textbf{译文} 但当一条直线立在另一条直线上并使相邻角彼此相等时，每一个相等的角都是直角；[定义 10] 所以角 $DCF$、$FCE$ 都是直角。因此，已经从给定直线 $AB$ 上的给定点 $C$ 作出与之成直角的直线 $CF$，即一条垂线。作图完毕。

\noindent\textbf{原文} [\ldots]

\noindent\textbf{译文} [原文此处省略]

\subsubsection*{Proposition 13 / 命题 13}

\noindent\textbf{原文} If a straight line set up on a straight line make angles, it will make either two right angles or angles equal to two right angles.

\noindent\textbf{译文} 若一条直线立在另一条直线上而成角，则它所成的要么是两个直角，要么是和等于两个直角的角。

\begin{figure}[htbp]
\centering
\safeincludeimage{71267a012817a97b765d3c633d03d379115716b87ef23ca8d5302d010a8a8c11.jpg}
\caption{Proposition 13 diagram / 命题 13 图}
\end{figure}

\noindent\textbf{原文} For let any straight line $AB$ set up on the straight line $CD$ make the angles $CBA$, $ABD$; I say that the angles $CBA$, $ABD$ are either two right angles or equal to two right angles. Now, if the angle $CBA$ is equal to the angle $ABD$, they are two right angles. [Def. 10]

\noindent\textbf{译文} 设任一直线 $AB$ 立在直线 $CD$ 上，形成角 $CBA$、$ABD$；我说，角 $CBA$、$ABD$ 要么是两个直角，要么其和等于两个直角。现在，如果角 $CBA$ 等于角 $ABD$，它们就是两个直角。[定义 10]

\noindent\textbf{原文} But, if not, let $BE$ be drawn from the point $B$ at right angles to $CD$; [I. 11] therefore the angles $CBE$, $EBD$ are two right angles. Then, since the angle $CBE$ is equal to the two angles $CBA$, $ABE$, let the angle $EBD$ be added to each; therefore the angles $CBE$, $EBD$ are equal to the three angles $CBA$, $ABE$, $EBD$. [C.N. 2] Again, since the angle $DBA$ is equal to the two angles $DBE$, $EBA$, let the angle $ABC$ be added to each; therefore the angles $DBA$, $ABC$ are equal to the three angles $DBE$, $EBA$, $ABC$. [C.N. 2]

\noindent\textbf{译文} 若不相等，则从点 $B$ 作 $BE$ 垂直于 $CD$；[I. 11] 因此角 $CBE$、$EBD$ 是两个直角。于是，由于角 $CBE$ 等于两个角 $CBA$、$ABE$ 之和，把角 $EBD$ 分别加到两边；所以角 $CBE$、$EBD$ 之和等于三个角 $CBA$、$ABE$、$EBD$ 之和。[共同概念 2] 又由于角 $DBA$ 等于两个角 $DBE$、$EBA$ 之和，把角 $ABC$ 分别加到两边；所以角 $DBA$、$ABC$ 之和等于三个角 $DBE$、$EBA$、$ABC$ 之和。[共同概念 2]

\noindent\textbf{原文} But the angles $CBE$, $EBD$ were also proved equal to the same three angles; and things which are equal to the same thing are also equal to one another; [C.N. 1] therefore the angles $CBE$, $EBD$ are also equal to the angles $DBA$, $ABC$. But the angles $CBE$, $EBD$ are two right angles; therefore the angles $DBA$, $ABC$ are also equal to two right angles. Therefore etc. Q.E.D.

\noindent\textbf{译文} 但角 $CBE$、$EBD$ 也已证明等于同样的三个角；而等于同一事物的事物彼此相等；[共同概念 1] 因此角 $CBE$、$EBD$ 也等于角 $DBA$、$ABC$。但角 $CBE$、$EBD$ 是两个直角；所以角 $DBA$、$ABC$ 之和也等于两个直角。所以等等。证毕。

\noindent\textbf{原文} [\ldots]

\noindent\textbf{译文} [原文此处省略]

\subsubsection*{Proposition 15 / 命题 15}

\noindent\textbf{原文} If two straight lines cut one another, they make the vertical angles equal to one another.

\noindent\textbf{译文} 若两条直线相交，则它们所成的对顶角彼此相等。

\noindent\textbf{原文} For let the straight lines $AB$, $CD$ cut one another at the point $E$; I say that the angle $AEC$ is equal to the angle $DEB$, and the angle $CEB$ to the angle $AED$.

\noindent\textbf{译文} 设直线 $AB$、$CD$ 相交于点 $E$；我说，角 $AEC$ 等于角 $DEB$，角 $CEB$ 等于角 $AED$。

\begin{figure}[htbp]
\centering
\safeincludeimage{6246ab53667b2e8dc6015ec40cae51efb7530a8d663f20e0861ba3249b47d47c.jpg}
\caption{Proposition 15 diagram / 命题 15 图}
\end{figure}

\noindent\textbf{原文} For, since the straight line $AE$ stands on the straight line $CD$, making the angles $CEA$, $AED$, the angles $CEA$, $AED$ are equal to two right angles. [I. 13] Again, since the straight line $DE$ stands on the straight line $AB$, making the angles $AED$, $DEB$, the angles $AED$, $DEB$ are equal to two right angles. [I. 13]

\noindent\textbf{译文} 因为直线 $AE$ 立在直线 $CD$ 上，形成角 $CEA$、$AED$，所以角 $CEA$、$AED$ 之和等于两个直角。[I. 13] 又由于直线 $DE$ 立在直线 $AB$ 上，形成角 $AED$、$DEB$，所以角 $AED$、$DEB$ 之和等于两个直角。[I. 13]

\noindent\textbf{原文} But the angles $CEA$, $AED$ were also proved equal to two right angles; therefore the angles $CEA$, $AED$ are equal to the angles $AED$, $DEB$. [Post. 4 and C.N. 1] Let the angle $AED$ be subtracted from each; therefore the remaining angle $CEA$ is equal to the remaining angle $BED$. [C.N. 3] Similarly it can be proved that the angles $CEB$, $DEA$ are also equal. Therefore etc. Q.E.D.

\noindent\textbf{译文} 但角 $CEA$、$AED$ 也已证明等于两个直角；因此角 $CEA$、$AED$ 之和等于角 $AED$、$DEB$ 之和。[公设 4 与共同概念 1] 从两边减去角 $AED$；于是余下的角 $CEA$ 等于余下的角 $BED$。[共同概念 3] 同理可证，角 $CEB$、$DEA$ 也相等。所以等等。证毕。

\noindent\textbf{原文} [\ldots]

\noindent\textbf{译文} [原文此处省略]

\subsubsection*{Proposition 16 / 命题 16}

\noindent\textbf{原文} In any triangle, if one of the sides be produced, the exterior angle is greater than either of the interior and opposite angles.

\noindent\textbf{译文} 在任意三角形中，若延长一边，则外角大于任一不相邻内角。

\noindent\textbf{原文} Let $ABC$ be a triangle, and let one side of it $BC$ be produced to $D$; I say that the exterior angle $ACD$ is greater than either of the interior and opposite angles $CBA$, $BAC$. Let $AC$ be bisected at $E$, [I. 10] and let $BE$ be joined and produced in a straight line to $F$.

\noindent\textbf{译文} 设 $ABC$ 为三角形，并将其一边 $BC$ 延长到 $D$；我说，外角 $ACD$ 大于任一内对角 $CBA$、$BAC$。令 $AC$ 在 $E$ 处被平分，[I. 10] 并连接 $BE$，再沿直线延长到 $F$。

\begin{figure}[htbp]
\centering
\safeincludeimage{6b4a6abdc4714f5762c7b8e3be81141be8ef74db03fbf56119ba30fb31e25698.jpg}
\caption{Proposition 16 diagram / 命题 16 图}
\end{figure}

\noindent\textbf{原文} Let $EF$ be made equal to $BE$, [I. 3] let $FC$ be joined, [Post. 1] and let $AC$ be drawn through to $G$. [Post. 2] Then, since $AE$ is equal to $EC$, and $BE$ to $EF$, the two sides $AE$, $EB$ are equal to the two sides $CE$, $EF$ respectively; and the angle $AEB$ is equal to the angle $FEC$, for they are vertical angles. [I. 15] Therefore the base $AB$ is equal to the base $FC$, and the triangle $ABE$ is equal to the triangle $CFE$, and the remaining angles are equal to the remaining angles respectively, namely those which the equal sides subtend; [I. 4] therefore the angle $BAE$ is equal to the angle $ECF$.

\noindent\textbf{译文} 令 $EF$ 等于 $BE$，[I. 3] 连接 $FC$，[公设 1] 并将 $AC$ 延长到 $G$。[公设 2] 于是，由于 $AE$ 等于 $EC$，且 $BE$ 等于 $EF$，两边 $AE$、$EB$ 分别等于两边 $CE$、$EF$；而角 $AEB$ 等于角 $FEC$，因为它们是对顶角。[I. 15] 因此底 $AB$ 等于底 $FC$，三角形 $ABE$ 等于三角形 $CFE$，其余各角也分别相等，即相等边所对的角相等；[I. 4] 所以角 $BAE$ 等于角 $ECF$。

\noindent\textbf{原文} But the angle $ECD$ is greater than the angle $ECF$; [C.N. 5] therefore the angle $ACD$ is greater than the angle $BAE$. Similarly also, if $BC$ be bisected, the angle $BCG$, that is, the angle $ACD$, [I. 15] can be proved greater than the angle $ABC$ as well. Therefore etc. Q.E.D.

\noindent\textbf{译文} 但角 $ECD$ 大于角 $ECF$；[共同概念 5] 因此角 $ACD$ 大于角 $BAE$。同样，若平分 $BC$，也可证明角 $BCG$，也就是角 $ACD$，[I. 15] 大于角 $ABC$。所以等等。证毕。

\noindent\textbf{原文} [\ldots]

\noindent\textbf{译文} [原文此处省略]

\subsubsection*{Proposition 18 / 命题 18}

\noindent\textbf{原文} In any triangle the greater side subtends the greater angle.

\noindent\textbf{译文} 在任意三角形中，较大的边对较大的角。

\begin{figure}[htbp]
\centering
\safeincludeimage{d9da22780dfcf229f4edbdb0297898dcad5788013f8abf1785e53626f11eaccb.jpg}
\caption{Proposition 18 diagram / 命题 18 图}
\end{figure}

\noindent\textbf{原文} For let $ABC$ be a triangle having the side $AC$ greater than $AB$; I say that the angle $ABC$ is also greater than the angle $BCA$. For, since $AC$ is greater than $AB$, let $AD$ be made equal to $AB$; [I. 3] and let $BD$ be joined. Then, since the angle $ADB$ is an exterior angle of the triangle $BCD$, it is greater than the interior and opposite angle $DCB$. But the angle $ADB$ is equal to the angle $ABD$, since the side $AB$ is equal to $AD$; [I. 5] therefore the angle $ABD$ is also greater than the angle $ACB$; therefore the angle $ABC$ is much greater than the angle $ACB$. Therefore etc. Q.E.D.

\noindent\textbf{译文} 设 $ABC$ 为三角形，且边 $AC$ 大于 $AB$；我说，角 $ABC$ 也大于角 $BCA$。因为 $AC$ 大于 $AB$，令 $AD$ 等于 $AB$；[I. 3] 并连接 $BD$。于是，由于角 $ADB$ 是三角形 $BCD$ 的外角，它大于内对角 $DCB$。但角 $ADB$ 等于角 $ABD$，因为边 $AB$ 等于 $AD$；[I. 5] 因此角 $ABD$ 也大于角 $ACB$；所以角 $ABC$ 更大于角 $ACB$。所以等等。证毕。

\subsubsection*{Proposition 19 / 命题 19}

\noindent\textbf{原文} In any triangle the greater angle is subtended by the greater side.

\noindent\textbf{译文} 在任意三角形中，较大的角由较大的边所对。

\noindent\textbf{原文} Let $ABC$ be a triangle having the angle $ABC$ greater than the angle $BCA$; I say that the side $AC$ is also greater than the side $AB$. For, if not, $AC$ is either equal to $AB$ or less. Now $AC$ is not equal to $AB$; for then the angle $ABC$ would also have been equal to the angle $ACB$; [I. 5] but it is not.

\noindent\textbf{译文} 设 $ABC$ 为三角形，且角 $ABC$ 大于角 $BCA$；我说，边 $AC$ 也大于边 $AB$。因为若不然，$AC$ 要么等于 $AB$，要么小于 $AB$。现在 $AC$ 不等于 $AB$；否则角 $ABC$ 也会等于角 $ACB$；[I. 5] 但它并不相等。

\begin{figure}[htbp]
\centering
\safeincludeimage{302a6c393eebc8820de0cfd0af310f1bc81e9a0762c586c845292d64ed567da0.jpg}
\caption{Proposition 19 diagram / 命题 19 图}
\end{figure}

\noindent\textbf{原文} Therefore $AC$ is not equal to $AB$. Neither is $AC$ less than $AB$, for then the angle $ABC$ would also have been less than the angle $ACB$; [I. 18] but it is not; therefore $AC$ is not less than $AB$. And it was proved that it is not equal either. Therefore $AC$ is greater than $AB$. Therefore etc. Q.E.D.

\noindent\textbf{译文} 因此 $AC$ 不等于 $AB$。$AC$ 也不小于 $AB$，因为若小于，则角 $ABC$ 也会小于角 $ACB$；[I. 18] 但事实并非如此；因此 $AC$ 不小于 $AB$。又已证明它也不相等。所以 $AC$ 大于 $AB$。所以等等。证毕。

\subsubsection*{Proposition 20 / 命题 20}

\noindent\textbf{原文} In any triangle two sides taken together in any manner are greater than the remaining one.

\noindent\textbf{译文} 在任意三角形中，任取两边之和都大于余下的一边。

\noindent\textbf{原文} For let $ABC$ be a triangle; I say that in the triangle $ABC$ two sides taken together in any manner are greater than the remaining one, namely $BA$, $AC$ greater than $BC$, $AB$, $BC$ greater than $AC$, and $BC$, $CA$ greater than $AB$.

\noindent\textbf{译文} 设 $ABC$ 为三角形；我说，在三角形 $ABC$ 中，任取两边之和都大于余下的一边，即 $BA$、$AC$ 之和大于 $BC$，$AB$、$BC$ 之和大于 $AC$，$BC$、$CA$ 之和大于 $AB$。

\begin{figure}[htbp]
\centering
\safeincludeimage{4f02fa1bc68d5bb688d94c409625e29512d0f27b81961c35b4b9b5890f24a90d.jpg}
\caption{Proposition 20 diagram / 命题 20 图}
\end{figure}

\noindent\textbf{原文} For let $BA$ be drawn through to the point $D$, let $DA$ be made equal to $CA$, and let $DC$ be joined. Then, since $DA$ is equal to $AC$, the angle $ADC$ is also equal to the angle $ACD$; [I. 5] therefore the angle $BCD$ is greater than the angle $ADC$. [C.N. 5]

\noindent\textbf{译文} 将 $BA$ 延长到点 $D$，令 $DA$ 等于 $CA$，并连接 $DC$。于是，由于 $DA$ 等于 $AC$，角 $ADC$ 也等于角 $ACD$；[I. 5] 因此角 $BCD$ 大于角 $ADC$。[共同概念 5]

\noindent\textbf{原文} And, since $DCB$ is a triangle having the angle $BCD$ greater than the angle $BDC$, and the greater angle is subtended by the greater side, [I. 19] therefore $DB$ is greater than $BC$. But $DA$ is equal to $AC$; therefore $BA$, $AC$ are greater than $BC$. Similarly we can prove that $AB$, $BC$ are also greater than $CA$, and $BC$, $CA$ than $AB$. Therefore etc. Q.E.D.

\noindent\textbf{译文} 又由于 $DCB$ 是一个三角形，其中角 $BCD$ 大于角 $BDC$，而较大的角由较大的边所对，[I. 19] 所以 $DB$ 大于 $BC$。但 $DA$ 等于 $AC$；因此 $BA$、$AC$ 之和大于 $BC$。同理可证，$AB$、$BC$ 之和也大于 $CA$，$BC$、$CA$ 之和大于 $AB$。所以等等。证毕。

\end{document}
